\documentclass[10pt,a4paper]{article}
\usepackage[margin=20mm,headheight=14pt,footskip=10mm]{geometry}
\usepackage{fontspec}
\setmainfont{Arial}
\setmonofont{Consolas}[Scale=MatchLowercase]
\usepackage{amsmath,amssymb,mathtools}
\usepackage{needspace}
\usepackage{xcolor}
\usepackage{url}

\urlstyle{same}
\def\UrlBreaks{\do\/\do\-\do\_\do\.\do\:\do\?\do\=}
\setlength{\parindent}{0pt}
\setlength{\parskip}{6pt}
\setlength{\emergencystretch}{3em}
\tolerance=2000
\allowdisplaybreaks[1]
\newcommand{\entrylabel}[1]{\par\needspace{3\baselineskip}\textbf{\color{blue!45!black}#1}\par}
\begin{document}


{\Huge\bfseries Probability:\\[5pt]An open-problem working book}\par\vspace{10mm}

A twenty-hour campaign with one research worker


Checkpoint:\allowbreak{} 2026-09-07T08:\allowbreak{}33:\allowbreak{}47Z


Start:\allowbreak{} 6 September 2026,\allowbreak{} 15:\allowbreak{}13:\allowbreak{}26 UTC. Deadline:\allowbreak{} 7 September 2026,\allowbreak{} 11:\allowbreak{}13:\allowbreak{}26 UTC. This evolving checkpoint does not declare the timed mandate complete.


\entrylabel{Coverage and its limits}

435 distinct questions presented as open in their primary sources;\allowbreak{} 435 have a substantive individual attempt. 0 original full solutions are promoted. The owner target is 1,\allowbreak{}000 distinct open probability problems:\allowbreak{} catalogue shortfall is 565. Historical source evidence is not a certification that all these questions remain open today.


\entrylabel{Candidate results}

12 current full candidate records:\allowbreak{} P0032,\allowbreak{} P0033,\allowbreak{} P0174,\allowbreak{} P0189,\allowbreak{} P0193,\allowbreak{} P0208,\allowbreak{} P0209,\allowbreak{} P0266,\allowbreak{} P0267,\allowbreak{} P0276,\allowbreak{} P0288,\allowbreak{} P0340. Candidates have proof packets and explicit limits. They are not promoted original solutions. Source-matched independent checks,\allowbreak{} current-literature status,\allowbreak{} and novelty are separate questions.


\entrylabel{Partial proof and counterexample candidates}

P0289,\allowbreak{} P0341,\allowbreak{} P0342,\allowbreak{} P0347,\allowbreak{} P0355,\allowbreak{} P0389,\allowbreak{} P0390,\allowbreak{} P0411,\allowbreak{} P0429,\allowbreak{} P0458,\allowbreak{} P0459. These records retain unresolved parts of their source questions and are excluded from the full-candidate total. The ledger status covers both partial proofs and scope counterexamples.


\entrylabel{Literature exclusions}

29 candidates have an established literature-resolution disposition and are excluded from the target. Other claimed resolutions retain their weaker status. In particular,\allowbreak{} P0248 is a Pak theorem rederivation and P0249 is a known consequence of Huang-Mao variance comparison.


\entrylabel{Persistent probability math brain}

38 private unverified entries and 39 edges were actually ingested. Retrieval probes and negative controls are recorded separately. Store acceptance and retrieval do not provide independent mathematical review. Prepared notes are not counted as ingested entries.


\entrylabel{Operational limits}

The original worker reported a dream age of 30.5 hours against a 30.0-hour ceiling. The resumed thread completed its own live wake with identity intact and fresh,\allowbreak{} and no owner-alert finding. The signed-project warning remains unresolved. Private Git checkpoints are not contract recovery. This research used no GPU training,\allowbreak{} new paid service,\allowbreak{} public publication,\allowbreak{} deployment,\allowbreak{} or remote computation.


\clearpage

\section*{Method and status vocabulary}

\entrylabel{source\_\allowbreak{}says\_\allowbreak{}open}

The cited primary source poses the question as unresolved at its recorded writing date.


\entrylabel{unverified}

Present literature status has not been established by this campaign.


\entrylabel{no\_\allowbreak{}resolution\_\allowbreak{}found\_\allowbreak{}in\_\allowbreak{}documented\_\allowbreak{}search}

A dated bounded search found no matching resolution. This does not certify current openness.


\entrylabel{unsolved or attempted\_\allowbreak{}unsolved}

The recorded attempt does not settle the whole mathematical target.


\entrylabel{candidate\_\allowbreak{}full\_\allowbreak{}solution}

A full argument is written,\allowbreak{} but its stated proof,\allowbreak{} assumption,\allowbreak{} review,\allowbreak{} or literature checks remain incomplete.


\entrylabel{proved\_\allowbreak{}here\_\allowbreak{}author\_\allowbreak{}checked}

Reserved for a source-matched complete argument that survives repeated author checks;\allowbreak{} distinct from independent review.


\entrylabel{resolved\_\allowbreak{}in\_\allowbreak{}literature}

An earlier theorem or counterexample addresses the source question. It is not a new solution here.


Distinct numbered or explicitly posed questions are admitted after mathematical-content deduplication. Related parameter versions and bundled conjectures are grouped. Every admitted item has its source,\allowbreak{} locator,\allowbreak{} dated historical status,\allowbreak{} precise target,\allowbreak{} individual argument,\allowbreak{} check,\allowbreak{} and obstruction. A counterexample to an attempted proof method is never relabelled a counterexample to the actual conjecture. Earlier attempt records remain visible when later evidence changes their status;\allowbreak{} the current per-problem disposition governs the headline count.


\clearpage

\section*{Question index}

P0001  Negative correlation in uniform forests


P0002  Self-avoiding-walk counting exponents


P0003  Hexagonal self-avoiding-walk scaling limit


P0004  Existence of the quenched connective limit


P0005  Ehrenfest mirror escape


P0006  Poisson mirror escape threshold


P0007  Unbounded mirror component trapping every ray


P0008  Diffusive Poisson mirror trajectories


P0009  Manhattan pinball localization


P0010  Biased random square-lattice orientations


P0011  Brownian SIR survival in dimensions at least two


P0012  Gaussian tensor type-2 bound


P0013  Expected Lovasz theta of a dense independent-edge graph


P0014  Expected Lovasz theta of a random circulant graph


P0015  Random complex phase-retrieval injectivity at 4M-5 measurements


P0016  Worst square-submatrix conditioning in a Gaussian frame


P0017  Dimension-free log-concave Poincare or Cheeger constant


P0018  Sharp polynomial-logarithmic expected norm of graph matrices


P0019  Brownian exit approach directions in higher dimensions


P0020  Finite-crossing arcs through a planar Brownian trace


P0021  Interior remainder of planar Brownian trace


P0022  Percolation dimension of a planar Brownian trace


P0023  Efficient reflected-Brownian coupling in acute triangles


P0024  Failure of synchronous contraction in a bounded planar domain


P0025  Synchronous reflection outside a disk


P0026  Non-explosion of Fleming-Viot Brownian particles in arbitrary domains


P0027  Rigidity of shy reflected-Brownian couplings


P0028  Bounded independent Brownian pieces and two-sided Brownian motion


P0029  Separation of two peaks in random square-grid labelings


P0030  Compatibility of forward and backward Brownian decompositions


P0031  Sparse first-passage pieces in a forward Brownian concatenation


P0032  Embedding an equivalent Brownian law by adding a finite prefix


P0033  Ergodicity of stationary Galton-Watson flows


P0034  Recurrence of the planar horizontal-new vertical-old walk


P0035  Positive speed of inward-biased walks below the Cayley growth rate


P0036  Percolation threshold on arbitrary geodesic Cayley spanning trees


P0037  BK disjoint occurrence for weighted uniform spanning trees


P0038  Transience preserved in Cayley percolation clusters


P0039  Positive-speed equivalence for Cayley percolation clusters


P0040  Uniqueness threshold below one in one-ended transitive graphs


P0041  Nonuniqueness phase for nonamenable quasi-transitive graphs


P0042  Bootstrap thresholds as a characterization of amenable groups


P0043  Absence of infinite clusters at critical percolation


P0044  Invariant monotone coupling of free and wired uniform spanning forests


P0045  Rational expected free-forest degree for finitely presented groups


P0046  Connecting a free uniform spanning forest by arbitrarily sparse invariant additions


P0047  Deterministic component count of the free uniform spanning forest


P0048  One or infinitely many free-forest components without unimodularity


P0049  Infinitely many ends in all free-forest components without unimodularity


P0050  Singularity of distinct free and wired forest measures


P0051  Zero-one law for all forest components having one end


P0052  Limiting bushes along a wired uniform-forest ray


P0053  One-ended wired minimal forests without unimodularity


P0054  Infinitely many ends in free minimal forests distinct from wired


P0055  Minimal-forest connectivity and uniqueness at p\_\allowbreak{}u


P0056  One or infinitely many minimal-forest components


P0057  Uniqueness-critical percolation versus invariant spanning trees


P0058  Indistinguishability of minimal spanning-forest components


P0059  Limiting bushes along a wired minimal-forest ray


P0060  Dimensions with a connected Euclidean minimal spanning forest


P0061  Full scaling limit of the square-lattice minimal spanning tree


P0062  Connectivity of minimal forests on graphs roughly isometric to trees


P0063  Singularity of binary and ternary harmonic codings of the Cantor set


P0064  Dimension drop for every nonuniform Galton-Watson flow


P0065  Monotonicity of harmonic dimension under inward bias


P0066  Monotonicity of biased-walk speed on Galton-Watson trees


P0067  Global Kuramoto synchronization of random cubic graphs


P0068  Subspectral detection from a growing number of synchronization frequencies


P0069  Subexponential detection at fixed synchronization frequencies


P0070  Kikuchi spectral thresholds as the polynomial degree increases


P0071  Polynomial Glauber mixing throughout the SK high-temperature regime


P0072  Geodesic selection by reinforced ants


P0073  Deterministic limits for trace-reinforced ants


P0074  Maximal degree and its location in the infinite-path fitness regime


P0075  Empirical degree scaling for superlinear attachment with fitness


P0076  Multiple-attachment phase transition with superlinear fitness


P0077  Fluctuations of fitness-reinforced sequence sums


P0078  Boundary profile of the reinforced-sequence condensate


P0079  Growth of sums with unbounded reinforcement fitness


P0080  Weak reinforcement on infinite WARM networks


P0081  Critical WARM on the integer line


P0082  Sharp reinforcement threshold for WARM forests


P0083  Sharp reinforcement threshold for WARM whiskers


P0084  Finite surviving components in infinite superlinear WARM


P0085  WARM percolation with vanishing rates on lattices


P0086  Recurrence of sublinear vertex-reinforced walk on Z


P0087  Forced network urn convergence and diffusion scaling


P0088  Matching local limits of Bernoulli and Poisson-outdegree attachment


P0089  Local limits for nonaffine attachment with multiple edges


P0090  Multistate pivot-tree criterion for interpretation


P0091  Infinite-state tree-automaton interpretation criteria


P0092  Randomized tree automata and endogeny


P0093  Infinitely many fixed points of a finite-color automaton


P0094  Prescribing rogue and interpretable automaton fixed points


P0095  Automaton-level conditions for a unique fixed point


P0097  An explicit nontrivial FPP time constant and limit shape


P0098  Strict FPP comparison in higher dimensions at a large minimum atom


P0099  Strict convexity of iid FPP limit shapes


P0100  Excluding axial facets and cubical FPP shapes


P0101  Classification of flat FPP boundary pieces


P0102  Strong Busemann limits for subadditive arrays and FPP


P0103  Uniform curvature of continuous-law FPP shapes


P0104  Subdiffusive and KPZ passage-time fluctuation exponents


P0105  High-dimensional limit of the FPP fluctuation exponent


P0106  Polynomial overlap of independent planar FPP geodesics


P0107  Critical Bernoulli FPP growth in three dimensions


P0108  Sharp width threshold for the thin-cylinder FPP CLT


P0109  Existence of finite FPP geodesics without distributional assumptions


P0110  An iid FPP law with proved KPZ exponent relation


P0111  Infinitely many planar FPP geodesic rays for continuous laws


P0112  Full-sequence convergence of FPP geodesics to a ray


P0113  Almost-sure escape from a fixed midpoint in FPP


P0114  Uniqueness of planar directional geodesics in every fixed direction


P0115  Coalescence of directional FPP rays in higher dimensions


P0116  Growth of the infinite-descendant part of the FPP geodesic tree


P0117  Absence of bi-infinite geodesics in planar iid FPP


P0118  Noncircular low-dimensional Eden limit shape


P0119  Noncoexistence of unequal-speed Richardson infections


P0120  Strict interface speedup in half-plane inhomogeneous FPP


P0121  Uniform Lipschitz height models and their interfaces


P0122  SLE4 from spectral truncations of the Gaussian free field


P0123  High-temperature triangular Ising interfaces converge to SLE6


P0124  Critical planar random-cluster interfaces for general cluster weight


P0125  Conformal limits of the critical and dense O(n) loop model


P0126  Planar uniform lattice-tree diameter and scaling limit


P0127  Efficient approximately uniform sampling of lattice trees


P0128  Cardy conformal invariance for Voronoi percolation


P0129  Cardy conformal invariance for Boolean continuum percolation


P0130  Sharp planar one-arm probability beyond its exponent


P0131  Exact critical percolation chemical-distance exponent


P0132  Measurable loop erasure of planar Brownian motion


P0133  Torus spanning-tree dual homotopy distribution


P0134  A nontrivial exact Hausdorff gauge for SLE


P0135  Almost-sure fibers of chronological loop addition


P0136  Chronological Brownian loop addition in three dimensions


P0137  A loop-erased representation of a unit Brownian loop


P0138  Law of loop-decorated SLE beyond kappa2


P0139  Rough-path continuity of chronological loop attachment


P0141  First-order distance constants in finite-variance preferential attachment


P0142  Fluctuations of preferential-attachment distances and diameter


P0143  Second-order extremal hopcount in exponential complete-graph FPP


P0144  Conservative weighted-distance fluctuations in heavy-tailed configuration models


P0145  Hopcount universality in infinite-variance weighted configuration models


P0146  Weak-to-strong disorder universality in complete-graph FPP


P0147  Fluctuations of weighted flooding and diameter on configuration models


P0148  General edge-weight laws for random-graph flooding and diameter


P0149  Extreme hopcounts for general random-graph degree and weight laws


P0150  Minimal assumptions for heavy-tailed critical metric convergence


P0151  Critical window for fixed-outdegree preferential attachment


P0152  Kernel criteria for critical inhomogeneous random-graph universality


P0153  Vacant-set percolation threshold on general configuration models


P0154  Metric scaling of critical random-walk vacant components


P0155  Erdos-Renyi component-size limits for high-dimensional tori


P0156  Configuration-model minimum spanning-tree metric limits


P0157  Two competing largest-component scales in barely subcritical configuration models


P0158  Ordered-phase susceptibility exponent on random trees and graphs


P0159  Quenched non-Gaussian critical spin limits on heterogeneous random graphs


P0160  Fixed-temperature metastability for general configuration models


P0161  Potts free energy for general configuration-model degrees


P0162  Fourier entropy versus total influence


P0163  Polynomial Fourier concentration for DNF formulas


P0164  Gaussian lower bound for noise stability of unbiased thresholds


P0165  Optimal odd functions for many-player correlation distillation


P0166  Influential coordinates of bounded low-degree polynomials


P0167  Minimum correlation of symmetric Gaussian sets


P0168  Noise-optimal equal-volume Gaussian partitions


P0169  Linear Fourier mass versus Boolean degree


P0170  Optimal limited independence for Gaussian polynomial thresholds


P0171  Polynomial small bias sufficient to fool DNF events


P0172  Sharp erasure-distillation bound from the Gaussian limit


P0173  Asymptotic influence bounds for polynomial thresholds


P0174  Sharp support-diameter threshold for free Brownian unimodality


P0175  Symmetric unimodality under free addition


P0176  Existence of nontrivial freely strongly unimodal laws


P0177  Location of Brownian modes at large time


P0178  Critical moment condition for eventual stable-process unimodality


P0179  Uniform permutation representations of the Gieseking group


P0180  Optimal spectra of random symmetric-group Cayley graphs


P0181  Fixed-degree random Cayley graphs with optimal spectral gap


P0182  Strong convergence of high-dimensional representations of a fixed compact group


P0183  Deterministic strong convergence from LPS generators


P0184  Strong convergence of modular-group actions on finite projective lines


P0185  Ramanujan inequalities for growing families of matrix polynomials


P0186  Largest coefficient dimension compatible with strong convergence


P0187  Optimal edge-fluctuation scales for arbitrary random-matrix polynomials


P0188  Joint strong convergence for permutation and deterministic matrices


P0189  Brown-measure convergence from general strong convergence


P0190  Intrinsic freeness for non-Hermitian eigenvalues


P0191  Least singular values at the polynomial constant term


P0192  Brown-measure universality beyond complex entry densities


P0193  Eventual log-unimodality from arbitrary compact positive initial laws


P0194  Polynomial mixing with plus Ising boundary conditions


P0195  Monotonicity of the ferromagnetic Ising spectral gap


P0196  Sharp comparison of systematic and random Ising updates


P0197  Relation between Ising relaxation and mixing on transitive graphs


P0198  Cutoff on bounded-degree transitive expanders


P0199  Cutoff for cyclic-to-random transpositions


P0200  Cutoff for card-cyclic random insertion


P0201  Optimal Glauber mixing near the coloring degree threshold


P0202  Diameter-squared upper bound for transitive mixing


P0203  Mixing exponent of uniform binary row additions


P0204  Mixing on large monotone subsets of the hypercube


P0205  Metropolis relaxation for geometric-time Cayley targets


P0208  Extremal arc-transitive walk parameters away from cycles


P0209  Uniform stationary normal approximation from gap and bounded observable


P0210  Martingales detectable by computationally practical stopping rules


P0211  Relaxation of linear-polymer coagulation and fragmentation


P0212  Low-density genealogical limit of branching coalescing walks


P0213  Constant-factor online scheduling with subadditive batch durations


P0214  Fluctuations of shortest routes in Poisson proximity networks


P0215  Growth exponent for common subcladograms of random trees


P0216  Limit constant for common random two-dimensional partial orders


P0217  Percolation of merging planar empires


P0218  Growth exponents in the spatial city-formation model


P0219  High-density oriented percolation flow asymptotics


P0220  Explicit stationary law of the leftmost Hammersley drift-jump particles


P0221  Limit of the optimal online random spanning-tree cost


P0222  Near-unit stretch versus length in spatial networks


P0223  Sharp relaxation exponent for local tree rotations and triangulation flips


P0224  Feasible length and stretch statistics of scale-invariant networks


P0225  Attainment of the optimal SIRSN length-stretch tradeoff


P0226  Local finiteness of scale-invariant network traffic measures


P0227  Implications among SIRSN finiteness and geodesic properties


P0228  Uniqueness of the unbounded major-road component


P0229  Integrable supremum of routes to a bounded region


P0230  Asymptotic total length of finite SIRSN samples


P0231  A zero-one law for Poisson-coloring frontier dimension


P0232  Continuous curves in the Poisson-coloring frontier


P0233  Absence of double points on Poisson-coloring interface curves


P0234  A safe margin around extended initial color seeds


P0235  A Feller topology for limits of reversible Markov chains


P0236  Infinite mean route length at each large fixed Euclidean separation


P0242  Optimal stationary stopping under Markov-chain lumping


P0243  Sharp constant-one comparisons of reversible mixing and stopping parameters


P0244  Sharp extremal walk parameters among regular graphs


P0245  Exact minimum stationary cover time for reversible chains


P0246  Universal leading cover-time upper bound under transitivity


P0247  Asymptotic uniformity of the final discovered vertex


P0250  Kalai absorption bound from a directed cut parameter


P0251  Mixing of the random-pair binary crossover chain


P0252  Sharp harmonic-mean characterization of meeting times


P0253  Maximal stationary antivoter variance


P0254  Mixing of gap-jump exclusion on a discrete circle


P0255  Variation of nearest-unvisited walk length under a random start


P0256  Weak concentration of nearest-unvisited walks under exponential edge noise


P0257  Law of large numbers and constants for NUV lengths


P0258  Variance scale and Gaussian fluctuations of NUV lengths


P0259  Never-visited density for the planar Poisson greedy walk


P0266  Critical variational SPDE maximal second-moment bound


P0267  Endpoint energy-space continuity for critical variational SPDEs


P0268  Abstract stochastic maximal regularity with non-small multiplicative noise


P0269  Necessary and sufficient operator criterion for stochastic maximal regularity


P0270  Critical time regularization at Hilbert exponent two


P0271  Global stochastic reaction diffusion under a triangular reaction structure


P0272  Global stochastic reaction diffusion with quadratic reactions and mass control


P0273  Local stochastic Navier Stokes theory with no slip and non-small transport noise


P0274  Optimal critical spatial smoothness for stochastic Navier Stokes transport noise


P0275  Supremum endpoint Serrin criterion for stochastic Navier Stokes


P0276  Product condition and cutoff for unweighted Abelian Cayley graphs


P0277  Cutoff for conservative exclusion in dimension at least two


P0278  Cutoff for unit-rate zero range on higher dimensional lattices


P0279  Cutoff for bounded-degree Ising throughout the rapid-mixing phase


P0280  Cutoff for bounded-degree hardcore throughout the rapid-mixing phase


P0281  Log concavity of finite particle TASEP transition probabilities


P0282  Chen log concavity conjecture for fixed size random permutation LIS


P0283  Exact binary longest common subsequence time constant


P0284  Gaussian fluctuations for uniform binary longest common subsequences


P0285  Finitary ordering permutation of a KPZ line ensemble vertical slice


P0286  Stationary two color ASEP covariance decay on a large ring


P0287  Invariant measures for full line TASEP with killing at one site


P0288  Rank one polynomial asymptotics for lambda transient Markov chains


P0289  Ross modulation monotonicity and a criterion for Cox workload ordering


P0290  Ross workload ordering for an image of a doubly monotone Markov chain


P0291  A nonmonotone Markov intensity with supermodular time ordering


P0292  Characterize CTMC speed monotonicity in supermodular order


P0293  Order one over N accuracy for least workload mean field queues


P0294  Compute a first order bias correction for continuous workload mean field models


P0295  Fluid limit with dependent service and patience marks


P0296  Interchange stationary and many server limits with dependent marks


P0297  Exact Hausdorff dimension of Liouville quantum gravity


P0298  General Toda three-point structure constants


P0299  Probabilistic conformal bootstrap for Toda theory


P0300  Conformal bootstrap for the noncompact H3 WZW model


P0301  Full probabilistic noncompact imaginary Liouville theory


P0302  One critically loaded queue under general k-limited polling inputs


P0303  Simultaneous criticality in k-limited polling


P0305  Hausdorff dimensions of bulk LQG geodesic star sets


P0306  One-dimensional boundary two-star set in LQG


P0307  Exact dimension of the LQG metric net


P0308  Canonical and grand-canonical equivalence for long-range Riesz gases


P0309  Positive-temperature crystalline Gibbs phases for three-dimensional Riesz gases


P0310  Uniqueness for the one-dimensional short-range Riesz Gibbs gas


P0311  Optimal truncated pair-correlation decay in the Riesz fluid phase


P0312  Infinite-volume BBGKY equation for the general sine-beta process


P0313  Translation invariance in low-dimensional long-range Riesz gases


P0314  Uniqueness of the general sine-beta DLR Gibbs measure


P0315  Forrester-Haldane pointwise pair-correlation asymptotics for sine beta


P0316  Two phase transitions in negative-exponent one-dimensional Riesz gases


P0317  Hyperuniformity of general long-range Riesz Gibbs gases


P0318  Sharp number-variance growth exponents for long-range Riesz gases


P0320  Number rigidity in the remaining Riesz exponent band


P0321  Existence of a hexatic phase for two-dimensional Riesz gases


P0322  Critical gap-correlation decay in the one-dimensional s=\allowbreak{}1 Riesz gas


P0323  DLR characterization of the circular Riesz thermodynamic limit


P0324  Analyticity of the one-dimensional Riesz free energy


P0325  Uniqueness of the microscopic bulk limit of the two-dimensional Coulomb gas


P0326  High-temperature uniqueness and variational characterization of the two-dimensional plasma


P0327  Exponential clustering in the high-temperature two-dimensional plasma


P0328  Gravitational allocation for the two-dimensional one-component plasma


P0329  Jancovici-Lebowitz-Manificat discrepancy deviations for general beta


P0330  Higher-order free-energy expansion of the two-dimensional Coulomb gas


P0331  Sharpness of the Debye scale in Coulomb local laws


P0332  Regularity threshold for bounded Coulomb-gas linear-statistic fluctuations


P0333  Multivariate fluctuation limits for empirical transport with Euclidean cost


P0334  Necessary and sufficient conditions for an Lp quantile-process limit


P0335  Nondegenerate centered fluctuations of null quadratic empirical transport


P0336  Selected limit of empirical Euclidean transport between separated rectangles


P0337  Consistent estimation of the Sinkhorn null-limit spectrum


P0338  Minimax rates for conditional optimal-transport maps


P0339  Sample complexity of Gromov-Wasserstein plans and dual potentials


P0340  Growth of the minimal drawn length of critical beta-splitting trees


P0341  Uniform occupation estimates and the polynomial-test ansatz


P0342  Transferring a stationary fringe-process CLT to finite beta-splitting trees


P0343  Correlations of critical-tree leaves at mesoscopic separation


P0344  Heights of prescribed interval-labelled leaves in the critical tree


P0345  The leading constant for maximum discrete critical-tree height


P0346  Distinct and repeated clade shapes in critical beta-splitting trees


P0347  Distributions of balance indices for critical beta-splitting trees


P0348  The drawn width profile of a critical beta-splitting cladogram


P0349  Complex powers of clade sizes in the critical tree


P0350  A branchpoint point-process scaling limit of the critical fringe tree


P0351  Connections between the critical fragmentation and beta(2,\allowbreak{}1) coalescent


P0355  Endogeny of nonminimal discounted-tree fixed points


P0356  A general boundedness criterion for discounted tree sums


P0357  Rays with persistent extremal branching-random-walk local times


P0358  The extremal ray-height constant in a heavy-tailed branching random walk


P0359  The multifractal spectrum of extremal ray local times


P0360  Does frozen-percolation endogeny imply unrestricted almost sure uniqueness?


P0361  Uniqueness of the nondiagonal scale-invariant frozen-percolation coupling


P0362  Constructing frozen percolation on the PWIT as a finite-graph limit


P0363  A finite-graph approximation to frozen percolation on the three-regular tree


P0364  Uniqueness in law from the unoriented freezing rule and tree symmetry


P0365  Classifying endogenous laws among all frozen-percolation RDE solutions


P0366  The extra randomness in nonendogenous frozen percolation


P0367  Support of the conditional burning-time law


P0368  All bivariate frozen-percolation fixed points and their attraction domains


P0369  Nonendogeny from instability of the diagonal coupling


P0370  Endogeny of frozen percolation on higher-degree trees and the PWIT


P0371  Existence of infinite-cluster frozen percolation on higher-dimensional lattices


P0372  Spherical shape and uniform bulk density of ARW aggregates


P0373  Full microscopic limit of the ARW point-source aggregate


P0374  Quadrature-domain scaling limit from extended ARW sources


P0375  Independence of ARW bulk microstructure from the initial source shape


P0376  Universal critical limit obtained by sprinkling subcritical ARW states


P0377  Universal stationary ARW density in higher-dimensional growing domains


P0378  Exhaustion-independent infinite-volume wired ARW measure


P0379  Equality of the three infinite-volume critical ARW laws


P0380  Equality of bulk one-site density and wired global ARW density


P0381  Higher-dimensional hockey-stick density evolution for wired ARW


P0382  Critical-density total-variation cutoff for wired ARW in higher dimensions


P0383  Exponential incompressibility of the infinite-volume wired ARW state


P0384  Sharp first slow-stabilization threshold of the ARW free chain


P0385  Asymptotic equivalence of supercritical free and wake ARW states


P0386  Uniform loss of initial-state information in dense ARW stabilization


P0387  Slow-to-one-step mixing transition of the ARW wake chain


P0388  Hyperuniformity of the infinite-volume wired ARW state


P0389  Onset of mesoscopic hyperuniformity in the ARW free chain


P0390  Hyperuniformity transition in stationary ARW wake states


P0391  Diffusion limit for an arbitrary symmetric M/\allowbreak{}G/\allowbreak{}1 discipline


P0392  Infinite-variance scaling limits of symmetric queue lengths


P0393  Uniform busy-cycle maximum bounds for symmetric queues


P0394  Efficient input estimation from periodically observed queue workloads


P0395  Network decomposition with several bottlenecks and internal flow covariance


P0396  IDC-based decomposition for networks of multiserver queues


P0397  Multiclass and non-Markovian-routing robust queueing decomposition


P0398  Robust network approximations under time-varying arrivals


P0399  Identifiable and efficient probe inference in tandem queueing networks


P0400  Response-time laws with correlated redundancy replicas


P0401  Optimizing replication degree under correlated service requirements


P0402  Stability beyond processor sharing with correlated redundancy copies


P0403  Optimal scheduling for correlated redundancy systems


P0404  Tractable concurrent-replica policies with broad stability guarantees


P0405  Characterizing optimality of MAX-H in multiclass queues


P0406  Finite-system performance of size-based scheduling with impatience and noisy sizes


P0408  Recovering an ideal Voronoi tessellation from vertices or its dual


P0409  Joint law of neighboring ideal Voronoi vertices


P0410  Hole probabilities for lower-dimensional ideal Voronoi faces


P0411  Boundary-hole monotonicity and coefficient positivity for ideal cells


P0412  A general transfer theorem between horospherical and hyperbolic typical values


P0413  Low-intensity variance of typical hyperbolic Voronoi cell volume


P0414  Hausdorff dimension of exceptional rays in the planar ideal Voronoi tree


P0415  Canonical intensity and typical-cell laws for unbounded hyperbolic cells


P0416  Amenability of ideal Bernoulli Voronoi relations on nonexact groups


P0417  Regular exponential volume growth without full graph IPVT convergence


P0418  Non-Poisson low-intensity ideal Voronoi diagrams


P0419  Equal-growth L1 products beyond the divergent tilted-volume condition


P0420  Low-intensity Voronoi limits for products with the L2 metric


P0421  Convergence of edge-measured product-graph Voronoi diagrams


P0422  More than two distinguishable types of group IPVT cells


P0425  Optimal parallel allocation with unknown nonexponential job sizes


P0426  Optimal known-size parallel allocation with Poisson arrivals


P0427  Optimal allocation between two exponential parallelism classes


P0428  Mean-response optimization for stochastically arriving DAG jobs


P0429  Largest growing fork-join batch compatible with asymptotic independence


P0430  Response laws for limited fork-join shortest-queue dispatching


P0431  Straggler replication in queued limited fork-join systems


P0432  Optimal known-size multiserver scheduling beyond SRPT


P0433  Optimal dispatching to SRPT queues outside heavy traffic


P0434  Sharp approximation ratio for simple alternatives to Gittins


P0435  Optimal nonasymptotic scheduling with unknown sizes in M/\allowbreak{}G/\allowbreak{}k


P0436  Optimal finite-deadline success scheduling


P0437  Dropping probability and optimality gap of DTLH


P0438  Feasibility of class-specific response quantiles


P0439  Server provisioning for a finite response-time SLO


P0440  Waiting-time tails with deterministic multiserver setup


P0441  Deterministic DelayedOff waiting-power tradeoff


P0442  Optimal waiting-power Pareto frontier with deterministic setup


P0443  Closed-form performance bounds for general setup distributions


P0444  Computational complexity of the multiserver-job stability threshold


P0445  Monotone workload recurrences beyond FCFS multiserver jobs


P0446  Perfect multiserver-job samples with a stability-matching upper bound


P0447  Finite-size response characterization for classical fork-join queues


P0450  Conditional generation limits for the Erlang weighted tree


P0451  Unimodularity and reversibility of the Erlang branching spine


P0452  Giant-component density from Erlang weighted-tree survival


P0453  Weaker degree-tail hypotheses for Erlang weighted-tree growth


P0456  The fine connectivity window for bilateral preference graphs


P0457  Diameter near the newly established bilateral connectivity threshold


P0458  The standardized degree law in growing bilateral preference graphs


P0459  Simultaneous concentration of bilateral preference degrees


P0461  Entropy concavity for determinantal kernels


P0462  Disjoint coupling of orthogonal projection processes


P0463  Block factors for finite-range determinantal symbols


P0464  Geometric-mean product bounds for equivariant determinantal laws


P0465  Linear coupling distance for noncommuting equivariant kernels


\clearpage

\section*{P0001  Negative correlation in uniform forests}

\entrylabel{Primary source}

Selected problems in probability theory;\allowbreak{} Conjecture 4.2,\allowbreak{} PDF page 8


\url{https://arxiv.org/pdf/2205.07318}


Source statement date:\allowbreak{} not established. Archived access:\allowbreak{} 2026-09-06T15:\allowbreak{}29:\allowbreak{}11Z.


Historical:\allowbreak{} source\_\allowbreak{}says\_\allowbreak{}open | Current:\allowbreak{} no\_\allowbreak{}resolution\_\allowbreak{}found\_\allowbreak{}in\_\allowbreak{}documented\_\allowbreak{}search | Result:\allowbreak{} partial\_\allowbreak{}result


\entrylabel{Question}

For every finite simple graph,\allowbreak{} does its uniform acyclic edge subset have pairwise nonpositive edge covariance?


\entrylabel{Dated literature check}

\{"date":\allowbreak{} "2026-09-06",\allowbreak{} "queries":\allowbreak{} ["uniform forests negative correlation site:\allowbreak{}arxiv.org"],\allowbreak{} "evidence\_\allowbreak{}url":\allowbreak{} "\url{https://arxiv.org/html/2603.10738v1",} "evidence\_\allowbreak{}locator":\allowbreak{} "Conjecture 1.3",\allowbreak{} "scope":\allowbreak{} "March 2026 primary paper still states general conjecture;\allowbreak{} later complete-graph results do not settle arbitrary graphs"\}


\subsection*{P0001-A001 | Condition independent weighted Bernoulli edges on acyclicity;\allowbreak{} factor disjoint cycle constraints in cactus graphs}

\entrylabel{Argument or result}

Complete author-checked proof of negative association for all finite cactus graphs and finite positive edge weights;\allowbreak{} target on arbitrary graphs remains unsolved here


\entrylabel{Second check}

Direct subset counts for uniform cycles give covariance -2\textasciicircum{}(m-2)/\allowbreak{}(2\textasciicircum{}m-1)\textasciicircum{}2;\allowbreak{} triangle -2/\allowbreak{}49 and tree zero agree


\entrylabel{Remaining obstruction}

Overlapping cycle constraints do not factor;\allowbreak{} generic conditioning on a decreasing event fails. Need an inequality using graphic-matroid circuit elimination.


Verification:\allowbreak{} author checked | Independent review:\allowbreak{} none | Full target solved:\allowbreak{} False


\entrylabel{Full proof packet}

proofs/\allowbreak{}P0001\_\allowbreak{}cactus\_\allowbreak{}forests.md


\clearpage

\section*{P0002  Self-avoiding-walk counting exponents}

\entrylabel{Primary source}

Selected problems in probability theory;\allowbreak{} Conjecture 2.1,\allowbreak{} PDF page 3 and footnote


\url{https://arxiv.org/pdf/2205.07318}


Source statement date:\allowbreak{} not established. Archived access:\allowbreak{} 2026-09-06T15:\allowbreak{}29:\allowbreak{}11Z.


Historical:\allowbreak{} source\_\allowbreak{}says\_\allowbreak{}open | Current:\allowbreak{} unverified | Result:\allowbreak{} attempted\_\allowbreak{}unsolved


\entrylabel{Question}

Determine the asymptotic correction to c\_\allowbreak{}n=\allowbreak{}mu\textasciicircum{}n for lattice self-avoiding walks in dimensions two,\allowbreak{} three,\allowbreak{} and four.


\entrylabel{Scope}

The source predicts gamma\_\allowbreak{}2=\allowbreak{}43/\allowbreak{}32 and explicitly requires logarithmic corrections in dimension four. High-dimensional cases already proved are excluded from the unresolved target. This remains one source problem,\allowbreak{} not separate parameter rows.


\subsection*{P0002-A001 | Submultiplicativity and its limits}

\entrylabel{Argument or result}

On Z\textasciicircum{}d,\allowbreak{} splitting a walk gives c\_\allowbreak{}(m+\allowbreak{}n)<=\allowbreak{}c\_\allowbreak{}m c\_\allowbreak{}n. Positive-coordinate walks give d\textasciicircum{}n<=\allowbreak{}c\_\allowbreak{}n,\allowbreak{} and the nonbacktracking bound gives c\_\allowbreak{}n<=\allowbreak{}2d(2d-1)\textasciicircum{}(n-1). Fekete determines the exponential growth rate. However a\_\allowbreak{}n=\allowbreak{}mu\textasciicircum{}n exp(sqrt(n)) is submultiplicative and has no fixed polynomial prefactor.


\entrylabel{Second check}

The example is a real sequence,\allowbreak{} not a self-avoiding-walk counterexample. The square-root subadditivity checks the claimed property directly.


\entrylabel{Remaining obstruction}

Submultiplicativity cannot determine the critical exponent or the dimension-four logarithmic correction.


Verification:\allowbreak{} author checked limited argument | Independent review:\allowbreak{} none | Full target solved:\allowbreak{} False


\clearpage

\section*{P0003  Hexagonal self-avoiding-walk scaling limit}

\entrylabel{Primary source}

Selected problems in probability theory;\allowbreak{} Question 2.2,\allowbreak{} PDF page 3


\url{https://arxiv.org/pdf/2205.07318}


Source statement date:\allowbreak{} not established. Archived access:\allowbreak{} 2026-09-06T15:\allowbreak{}29:\allowbreak{}11Z.


Historical:\allowbreak{} source\_\allowbreak{}says\_\allowbreak{}open | Current:\allowbreak{} unverified | Result:\allowbreak{} attempted\_\allowbreak{}unsolved


\entrylabel{Question}

Does the rescaled uniform n-step hexagonal-lattice self-avoiding walk converge to SLE\_\allowbreak{}(8/\allowbreak{}3)?


\entrylabel{Scope}

The source leaves the precise rescaling,\allowbreak{} curve topology,\allowbreak{} parametrization and limiting SLE geometry implicit;\allowbreak{} these are proof obligations,\allowbreak{} not permissions to change the question.


\subsection*{P0003-A001 | Exact finite-domain conditional law}

\entrylabel{Argument or result}

For a finite-graph a-to-b self-avoiding path with weight z\textasciicircum{}length,\allowbreak{} condition on an initial simple prefix. The remaining path has the same fugacity law on the graph with earlier prefix vertices removed. The factor z\textasciicircum{}prefix\_\allowbreak{}length cancels between numerator and partition function.


\entrylabel{Second check}

The endpoint stays in the graph;\allowbreak{} deleting it would invalidate the argument. Fixed-length conditioning also fixes the remaining length and is a different ensemble.


\entrylabel{Remaining obstruction}

This domain Markov property supplies neither tightness nor conformal invariance and does not identify an SLE scaling limit.


Verification:\allowbreak{} author checked limited argument | Independent review:\allowbreak{} none | Full target solved:\allowbreak{} False


\clearpage

\section*{P0004  Existence of the quenched connective limit}

\entrylabel{Primary source}

Selected problems in probability theory;\allowbreak{} Question 2.3,\allowbreak{} PDF pages 3-4


\url{https://arxiv.org/pdf/2205.07318}


Source statement date:\allowbreak{} not established. Archived access:\allowbreak{} 2026-09-06T15:\allowbreak{}29:\allowbreak{}11Z.


Historical:\allowbreak{} source\_\allowbreak{}says\_\allowbreak{}open | Current:\allowbreak{} unverified | Result:\allowbreak{} partial\_\allowbreak{}result


\entrylabel{Question}

On the infinite supercritical bond-percolation cluster of Z\textasciicircum{}2,\allowbreak{} does c\_\allowbreak{}n(v)\textasciicircum{}(1/\allowbreak{}n) converge almost surely,\allowbreak{} independently of v?


\subsection*{P0004-A001 | Adjacent-root walk decomposition and a Markov/\allowbreak{}Borel-Cantelli bound}

\entrylabel{Argument or result}

Root-independent limsup on every infinite connected bounded-degree graph;\allowbreak{} annealed upper bound p mu


\entrylabel{Second check}

Walks avoiding the neighboring root inject by prepending;\allowbreak{} walks visiting it inject into reversed-prefix/\allowbreak{}suffix pairs. These are disjoint cases and preserve lengths.


\entrylabel{Remaining obstruction}

The cut-walk inequality has an environment-dependent endpoint maximum. No stationary scalar subadditive process or liminf lower bound was obtained.


Verification:\allowbreak{} author checked | Independent review:\allowbreak{} none | Full target solved:\allowbreak{} False


\entrylabel{Full proof packet}

proofs/\allowbreak{}P0004\_\allowbreak{}quenched\_\allowbreak{}limsup.md


\clearpage

\section*{P0005  Ehrenfest mirror escape}

\entrylabel{Primary source}

Selected problems in probability theory;\allowbreak{} Question 3.1,\allowbreak{} PDF pages 4-5


\url{https://arxiv.org/pdf/2205.07318}


Source statement date:\allowbreak{} not established. Archived access:\allowbreak{} 2026-09-06T15:\allowbreak{}29:\allowbreak{}11Z.


Historical:\allowbreak{} source\_\allowbreak{}says\_\allowbreak{}open | Current:\allowbreak{} unverified | Result:\allowbreak{} attempted\_\allowbreak{}unsolved


\entrylabel{Question}

Which mirror densities p give positive escape probability in the square-lattice model with equally likely NE/\allowbreak{}NW mirrors?


\subsection*{P0005-A001 | Invertible routing and finite-volume diagnostics}

\entrylabel{Argument or result}

The local mirror pairings define an invertible map on oriented light states. Its orbits are cycles or bi-infinite trajectories. On a finite periodic torus every orbit is a cycle. A mirror-free initial straight run of N sites has probability (1-p)\textasciicircum{}N.


\entrylabel{Second check}

The finite-state bijection proves the torus assertion independently of sampling. The straight-run bound tends to zero for each fixed p>0.


\entrylabel{Remaining obstruction}

Finite periodic loop counts cannot decide infinite-volume escape;\allowbreak{} the N-to-infinity and p-to-zero limits cannot be interchanged from this bound.


Verification:\allowbreak{} author checked limited argument | Independent review:\allowbreak{} none | Full target solved:\allowbreak{} False


\clearpage

\section*{P0006  Poisson mirror escape threshold}

\entrylabel{Primary source}

Selected problems in probability theory;\allowbreak{} Question 3.2(a-c),\allowbreak{} PDF pages 5-6


\url{https://arxiv.org/pdf/2205.07318}


Source statement date:\allowbreak{} not established. Archived access:\allowbreak{} 2026-09-06T15:\allowbreak{}29:\allowbreak{}11Z.


Historical:\allowbreak{} source\_\allowbreak{}says\_\allowbreak{}open | Current:\allowbreak{} unverified | Result:\allowbreak{} attempted\_\allowbreak{}unsolved


\entrylabel{Question}

For nondegenerate Poisson-mirror orientation laws,\allowbreak{} characterize escape below the vacant-set threshold,\allowbreak{} including uniform orientations.


\entrylabel{Scope}

All three source subquestions are grouped to avoid counting nested parameter variants separately. The direction quantifier for theta\_\allowbreak{}mu must be resolved from the model before a theorem claim.


\subsection*{P0006-A001 | Exact first-collision distribution for Poisson mirrors}

\entrylabel{Argument or result}

For a straight test ray segment of length L in direction alpha and mirrors of total length epsilon with orientation law mu,\allowbreak{} the centers whose mirror crosses the segment form a parallelogram of area epsilon L |sin(alpha-beta)| for angle beta. Unit-intensity marked Poisson counting gives no-collision probability exp(-epsilon L integral |sin(alpha-beta)| mu(dbeta)). Thus the first free flight is exponential whenever this rate is positive.


\entrylabel{Second check}

The factor is total mirror length,\allowbreak{} not half-length. Parallel mirrors contribute zero area and endpoint coincidences have zero intensity. The rate is positive for every alpha if mu is not concentrated on a single unoriented direction.


\entrylabel{Remaining obstruction}

After reflection,\allowbreak{} previously examined empty regions and previously hit mirrors make the flight history dependent. The first-flight computation gives no independent-flight construction or escape threshold.


Verification:\allowbreak{} author checked limited argument | Independent review:\allowbreak{} none | Full target solved:\allowbreak{} False


\clearpage

\section*{P0007  Unbounded mirror component trapping every ray}

\entrylabel{Primary source}

Selected problems in probability theory;\allowbreak{} Question 3.3,\allowbreak{} PDF page 6


\url{https://arxiv.org/pdf/2205.07318}


Source statement date:\allowbreak{} not established. Archived access:\allowbreak{} 2026-09-06T15:\allowbreak{}29:\allowbreak{}11Z.


Historical:\allowbreak{} source\_\allowbreak{}says\_\allowbreak{}open | Current:\allowbreak{} unverified | Result:\allowbreak{} attempted\_\allowbreak{}unsolved


\entrylabel{Question}

Can an unbounded component of a mirror complement contain an origin from which every reflected ray stays bounded?


\subsection*{P0007-A001 | Attempted reduction to straight visibility}

\entrylabel{Argument or result}

An unbounded unobstructed straight ray from the origin immediately rules out the proposed all-directions trapped configuration. However an unbounded connected planar open set need not contain any unbounded straight ray:\allowbreak{} a sufficiently narrow unbounded spiral corridor is a topological counterexample to that inference.


\entrylabel{Second check}

The corridor is only a counterexample to the topological inference. It is not asserted to be a configuration of the source mirror model.


\entrylabel{Remaining obstruction}

Unboundedness of the vacant component gives a continuous route,\allowbreak{} not a realizable billiard route. No mirror configuration trapping all angles,\allowbreak{} or theorem preventing one,\allowbreak{} was constructed.


Verification:\allowbreak{} author checked limited argument | Independent review:\allowbreak{} none | Full target solved:\allowbreak{} False


\clearpage

\section*{P0008  Diffusive Poisson mirror trajectories}

\entrylabel{Primary source}

Selected problems in probability theory;\allowbreak{} Question 3.4,\allowbreak{} PDF page 7


\url{https://arxiv.org/pdf/2205.07318}


Source statement date:\allowbreak{} not established. Archived access:\allowbreak{} 2026-09-06T15:\allowbreak{}29:\allowbreak{}11Z.


Historical:\allowbreak{} source\_\allowbreak{}says\_\allowbreak{}open | Current:\allowbreak{} unverified | Result:\allowbreak{} attempted\_\allowbreak{}unsolved


\entrylabel{Question}

For rationally supported nondegenerate mirror orientations,\allowbreak{} do escaping trajectories have finite positive diffusivity and a Gaussian limit?


\subsection*{P0008-A001 | Green-Kubo variance and a failed CLT shortcut}

\entrylabel{Argument or result}

For a stationary mean-zero scalar velocity with integrable covariance C(s),\allowbreak{} Var(integral\_\allowbreak{}0\textasciicircum{}t V\_\allowbreak{}s ds)/\allowbreak{}t =\allowbreak{} 2 integral\_\allowbreak{}0\textasciicircum{}t (1-s/\allowbreak{}t) C(s) ds tends to 2 integral\_\allowbreak{}0\textasciicircum{}infinity C(s) ds. Covariance control alone does not imply a Gaussian limit:\allowbreak{} let Y\_\allowbreak{}i=\allowbreak{}A epsilon\_\allowbreak{}i with independent fair signs and independent A equal to 1 or 2 with probability 1/\allowbreak{}2. Distinct increments are uncorrelated,\allowbreak{} but n\textasciicircum{}(-1/\allowbreak{}2) sum Y\_\allowbreak{}i converges to the equally weighted mixture of N(0,\allowbreak{}1) and N(0,\allowbreak{}4).


\entrylabel{Second check}

Conditioning on A proves the mixture limit;\allowbreak{} its fourth moment differs from that of a Gaussian with the same variance. This is not a unit-speed mirror-model counterexample.


\entrylabel{Remaining obstruction}

Need mixing or a martingale approximation for the billiard environment,\allowbreak{} including the effect of conditioning on an unbounded trajectory.


Verification:\allowbreak{} author checked limited argument | Independent review:\allowbreak{} none | Full target solved:\allowbreak{} False


\clearpage

\section*{P0009  Manhattan pinball localization}

\entrylabel{Primary source}

Selected problems in probability theory;\allowbreak{} Question 3.5,\allowbreak{} PDF page 7;\allowbreak{} model Figure 3.2


\url{https://arxiv.org/pdf/2205.07318}


Source statement date:\allowbreak{} not established. Archived access:\allowbreak{} 2026-09-06T15:\allowbreak{}29:\allowbreak{}11Z.


Historical:\allowbreak{} source\_\allowbreak{}says\_\allowbreak{}open | Current:\allowbreak{} unverified | Result:\allowbreak{} attempted\_\allowbreak{}unsolved


\entrylabel{Question}

Does every positive mirror density localize the Manhattan-oriented pinball trajectory?


\subsection*{P0009-A001 | Small-density straight paths and lack of an order argument}

\entrylabel{Argument or result}

At q=\allowbreak{}0 a ray moves indefinitely along its Manhattan street. The event that its first N potential collisions contain no mirror has probability (1-q)\textasciicircum{}N. This only controls a finite initial flight and gives no positive lower bound as N increases at fixed q.


\entrylabel{Second check}

The finite-run probability agrees with independence and has the correct q=\allowbreak{}0 endpoint. Adding a mirror changes routing rather than merely deleting traversable states,\allowbreak{} so an edgewise percolation monotonicity proof has not been supplied.


\entrylabel{Remaining obstruction}

The known localization interval near and above one-half cannot be extended to every q>0 using this finite-run bound or an unproved monotonicity assumption.


Verification:\allowbreak{} author checked limited argument | Independent review:\allowbreak{} none | Full target solved:\allowbreak{} False


\clearpage

\section*{P0010  Biased random square-lattice orientations}

\entrylabel{Primary source}

Selected problems in probability theory;\allowbreak{} Question 5.1,\allowbreak{} PDF page 9


\url{https://arxiv.org/pdf/2205.07318}


Source statement date:\allowbreak{} not established. Archived access:\allowbreak{} 2026-09-06T15:\allowbreak{}29:\allowbreak{}11Z.


Historical:\allowbreak{} source\_\allowbreak{}says\_\allowbreak{}open | Current:\allowbreak{} unverified | Result:\allowbreak{} attempted\_\allowbreak{}unsolved


\entrylabel{Question}

Does every nonzero orientation bias give positive probability of an infinite directed path from the origin?


\subsection*{P0010-A001 | Directed-percolation subgraph comparison}

\entrylabel{Argument or result}

Reflecting both coordinate axes maps p to 1-p,\allowbreak{} so the infinite-path probability is symmetric about one-half. North-east paths use independent edges pointing north/\allowbreak{}east with probability p;\allowbreak{} survival of the usual north-east directed percolation therefore implies survival in the full oriented graph. The p=\allowbreak{}0 and p=\allowbreak{}1 endpoints have deterministic infinite paths.


\entrylabel{Second check}

The comparison is inclusion of path events,\allowbreak{} not equality of the two models. It proves a region above the directed-percolation threshold and its reflected counterpart only.


\entrylabel{Remaining obstruction}

It does not settle the interval between one-half and that threshold. The critical p=\allowbreak{}1/\allowbreak{}2 result does not extend by continuity alone.


Verification:\allowbreak{} author checked limited argument | Independent review:\allowbreak{} none | Full target solved:\allowbreak{} False


\clearpage

\section*{P0011  Brownian SIR survival in dimensions at least two}

\entrylabel{Primary source}

Selected problems in probability theory;\allowbreak{} Question 6.1(ii),\allowbreak{} PDF pages 9-10


\url{https://arxiv.org/pdf/2205.07318}


Source statement date:\allowbreak{} not established. Archived access:\allowbreak{} 2026-09-06T15:\allowbreak{}29:\allowbreak{}11Z.


Historical:\allowbreak{} source\_\allowbreak{}says\_\allowbreak{}open | Current:\allowbreak{} unverified | Result:\allowbreak{} attempted\_\allowbreak{}unsolved


\entrylabel{Question}

Does the fully diffusing Poisson Brownian SIR model survive for sufficiently small positive removal rate when d>=\allowbreak{}2?


\subsection*{P0011-A001 | Counterexample to the usual pathwise lifetime monotonicity coupling}

\entrylabel{Argument or result}

Consider a deterministic contact schedule with A initially infected:\allowbreak{} A-C at time .1,\allowbreak{} A-B at 1,\allowbreak{} C-B at 9,\allowbreak{} B-D at 10,\allowbreak{} and no other contacts. Baseline infectious lifetimes are L\_\allowbreak{}A=\allowbreak{}.5,\allowbreak{} L\_\allowbreak{}C=\allowbreak{}12,\allowbreak{} L\_\allowbreak{}B=\allowbreak{}2. Then B is infected at 9 and infects D at 10. Multiply every lifetime by 3:\allowbreak{} A now infects B at 1;\allowbreak{} B is removed at 7,\allowbreak{} cannot be reinfected at 9,\allowbreak{} and D is never infected.


\entrylabel{Second check}

Check both timelines:\allowbreak{} C remains infectious at 9 in each case;\allowbreak{} A has already died at 1 only in the baseline. The example refutes infection-set inclusion under scaled lifetimes on a fixed contact schedule.


\entrylabel{Remaining obstruction}

This is not a counterexample to survival monotonicity or to the Brownian epidemic conjecture. It blocks a direct attractiveness proof;\allowbreak{} a Brownian contact-support realization was not proved.


Verification:\allowbreak{} author checked limited argument | Independent review:\allowbreak{} none | Full target solved:\allowbreak{} False


\clearpage

\section*{P0012  Gaussian tensor type-2 bound}

\entrylabel{Primary source}

Randomstrasse101:\allowbreak{} Open Problems of 2025;\allowbreak{} Conjecture 16,\allowbreak{} Entry 8


\url{https://arxiv.org/pdf/2603.29571}


Source statement date:\allowbreak{} not established. Archived access:\allowbreak{} 2026-09-06T15:\allowbreak{}32:\allowbreak{}10Z.


Historical:\allowbreak{} source\_\allowbreak{}says\_\allowbreak{}open | Current:\allowbreak{} unverified | Result:\allowbreak{} attempted\_\allowbreak{}unsolved


\entrylabel{Question}

For symmetric order-r tensors T\_\allowbreak{}i on R\textasciicircum{}d,\allowbreak{} p>=\allowbreak{}2 and independent standard Gaussians g\_\allowbreak{}i,\allowbreak{} is E||sum\_\allowbreak{}i g\_\allowbreak{}i T\_\allowbreak{}i||\_\allowbreak{}(I\_\allowbreak{}p) bounded by polylog(d,\allowbreak{}n) times O\_\allowbreak{}(r,\allowbreak{}p)(d\textasciicircum{}(1/\allowbreak{}2-1/\allowbreak{}p) sqrt(sum\_\allowbreak{}i ||T\_\allowbreak{}i||\_\allowbreak{}(I\_\allowbreak{}p)\textasciicircum{}2))?


\entrylabel{Scope}

The unresolved range is p<2r apart from known matrix and other special cases. The source says p>=\allowbreak{}2r is proved. All parameters remain one problem.


\subsection*{P0012-A001 | Diagonal-tensor Gaussian estimate}

\entrylabel{Argument or result}

For a diagonal homogeneous form T(x)=\allowbreak{}sum\_\allowbreak{}j a\_\allowbreak{}j x\_\allowbreak{}j\textasciicircum{}r and 2<=\allowbreak{}p<=\allowbreak{}r,\allowbreak{} its supremum over the l\_\allowbreak{}p unit ball is max\_\allowbreak{}j |a\_\allowbreak{}j|. Indeed sum |x\_\allowbreak{}j|\textasciicircum{}r<=\allowbreak{}1 and a coordinate vector attains the maximum. For independent standard g\_\allowbreak{}i,\allowbreak{} each diagonal coefficient of sum\_\allowbreak{}i g\_\allowbreak{}i T\_\allowbreak{}i has variance at most sum\_\allowbreak{}i ||T\_\allowbreak{}i||\textasciicircum{}2,\allowbreak{} so the Gaussian maximum bound gives E||sum\_\allowbreak{}i g\_\allowbreak{}i T\_\allowbreak{}i|| <=\allowbreak{} sqrt(2 log(2d)) (sum\_\allowbreak{}i ||T\_\allowbreak{}i||\textasciicircum{}2)\textasciicircum{}(1/\allowbreak{}2).


\entrylabel{Second check}

The upper bound follows by the exponential-moment union bound and minimizing its parameter;\allowbreak{} coordinate vectors prove the norm identity even for even r by taking absolute value outside the form.


\entrylabel{Remaining obstruction}

This verifies a restricted diagonal family. Off-diagonal terms and the full p range require the tensor geometry absent here.


Verification:\allowbreak{} author checked limited argument | Independent review:\allowbreak{} none | Full target solved:\allowbreak{} False


\clearpage

\section*{P0013  Expected Lovasz theta of a dense independent-edge graph}

\entrylabel{Primary source}

Randomstrasse101:\allowbreak{} Open Problems of 2025;\allowbreak{} Conjecture 17,\allowbreak{} Entry 9


\url{https://arxiv.org/pdf/2603.29571}


Source statement date:\allowbreak{} not established. Archived access:\allowbreak{} 2026-09-06T15:\allowbreak{}32:\allowbreak{}10Z.


Historical:\allowbreak{} source\_\allowbreak{}says\_\allowbreak{}open | Current:\allowbreak{} unverified | Result:\allowbreak{} attempted\_\allowbreak{}unsolved


\entrylabel{Question}

For G distributed as G(n,\allowbreak{}1/\allowbreak{}2),\allowbreak{} is E theta(G)=\allowbreak{}(1+\allowbreak{}o(1)) sqrt(n)?


\entrylabel{Scope}

Theta is the SDP with PSD matrix,\allowbreak{} trace one,\allowbreak{} and zero entries on graph edges.


\subsection*{P0013-A001 | Complement symmetry and uniform-integrability obstruction}

\entrylabel{Argument or result}

The theta product inequality theta(G) theta(complement G)>=\allowbreak{}n and complement symmetry of G(n,\allowbreak{}1/\allowbreak{}2) give E theta>=\allowbreak{}sqrt(n) by arithmetic-geometric mean. A proof of concentration on a logarithmic scale would still need tail control:\allowbreak{} set Z\_\allowbreak{}n to plus or minus (log n)/\allowbreak{}2,\allowbreak{} each with probability 1/\allowbreak{}(2 sqrt(n)),\allowbreak{} and zero otherwise. Then Z\_\allowbreak{}n tends to zero in probability and variance,\allowbreak{} but E exp(Z\_\allowbreak{}n) tends to 3/\allowbreak{}2.


\entrylabel{Second check}

The example has a symmetric law and exp(Z\_\allowbreak{}n) in [n\textasciicircum{}(-1/\allowbreak{}2),\allowbreak{}n\textasciicircum{}(1/\allowbreak{}2)]. Its expectation is 1-1/\allowbreak{}sqrt(n)+\allowbreak{}(1+\allowbreak{}n\textasciicircum{}(-1))/\allowbreak{}(2),\allowbreak{} which tends to 3/\allowbreak{}2. It is not a random-graph construction.


\entrylabel{Remaining obstruction}

Need a sharp upper bound and uniform integrability at the sqrt(n) scale;\allowbreak{} complement symmetry only yields the lower bound.


Verification:\allowbreak{} author checked limited argument | Independent review:\allowbreak{} none | Full target solved:\allowbreak{} False


\clearpage

\section*{P0014  Expected Lovasz theta of a random circulant graph}

\entrylabel{Primary source}

Randomstrasse101:\allowbreak{} Open Problems of 2025;\allowbreak{} Conjecture 18,\allowbreak{} Entry 9


\url{https://arxiv.org/pdf/2603.29571}


Source statement date:\allowbreak{} not established. Archived access:\allowbreak{} 2026-09-06T15:\allowbreak{}32:\allowbreak{}10Z.


Historical:\allowbreak{} source\_\allowbreak{}says\_\allowbreak{}open | Current:\allowbreak{} unverified | Result:\allowbreak{} attempted\_\allowbreak{}unsolved


\entrylabel{Question}

For the graph on Z/\allowbreak{}nZ formed by independently including each unordered nonzero generator pair \{k,\allowbreak{}-k\} with probability 1/\allowbreak{}2,\allowbreak{} is E theta(G)=\allowbreak{}(1+\allowbreak{}o(1)) sqrt(n)?


\entrylabel{Scope}

The inverse-pair convention includes the order-two generator once when n is even;\allowbreak{} this differs substantively from independent edges.


\subsection*{P0014-A001 | Vertex-transitive complement identity}

\entrylabel{Argument or result}

A circulant graph is vertex transitive,\allowbreak{} so theta(G) theta(complement G)=\allowbreak{}n. Its random model is complement invariant. For Z=\allowbreak{}log(theta(G)/\allowbreak{}sqrt(n)),\allowbreak{} Z has the law of -Z and E theta/\allowbreak{}sqrt(n)=\allowbreak{}E cosh Z. The conjecture is therefore equivalent to E(cosh Z-1) tending to zero.


\entrylabel{Second check}

The exact identity is stronger than the general product inequality. The symmetric rare-tail example in P0013 has Z tending to zero in probability yet E cosh Z tending to 3/\allowbreak{}2,\allowbreak{} so concentration alone is insufficient.


\entrylabel{Remaining obstruction}

Must estimate the exponential tails of the logarithmic theta fluctuation for the dependent circulant adjacency structure.


Verification:\allowbreak{} author checked limited argument | Independent review:\allowbreak{} none | Full target solved:\allowbreak{} False


\clearpage

\section*{P0015  Random complex phase-retrieval injectivity at 4M-5 measurements}

\entrylabel{Primary source}

Randomstrasse101:\allowbreak{} Open Problems of 2025;\allowbreak{} Conjecture 19(a,\allowbreak{}b),\allowbreak{} Entry 10


\url{https://arxiv.org/pdf/2603.29571}


Source statement date:\allowbreak{} not established. Archived access:\allowbreak{} 2026-09-06T15:\allowbreak{}32:\allowbreak{}10Z.


Historical:\allowbreak{} source\_\allowbreak{}says\_\allowbreak{}open | Current:\allowbreak{} unverified | Result:\allowbreak{} attempted\_\allowbreak{}unsolved


\entrylabel{Question}

Let p\_\allowbreak{}M be the probability that a complex Gaussian (4M-5)-by-M matrix gives an injective phaseless measurement map modulo global phase. Is p\_\allowbreak{}M<1 for each meaningful M,\allowbreak{} and does p\_\allowbreak{}M tend to zero?


\entrylabel{Scope}

M>=\allowbreak{}2 is needed for the matrix size to be meaningful. The two clauses are retained as one source conjecture,\allowbreak{} not counted twice.


\subsection*{P0015-A001 | Complete M=\allowbreak{}2 obstruction via Hermitian lifting}

\entrylabel{Argument or result}

For M=\allowbreak{}2 the matrix has three complex Gaussian rows. Almost surely these span C\textasciicircum{}2. The real-linear measurement map from the four-dimensional Hermitian 2-by-2 space to R\textasciicircum{}3 has a nonzero kernel element H. H cannot be positive or negative semidefinite:\allowbreak{} a\_\allowbreak{}i\textasciicircum{}* H a\_\allowbreak{}i=\allowbreak{}0 for a spanning set would force H=\allowbreak{}0. Thus H has one positive and one negative eigenvalue and can be written xx\textasciicircum{}*-yy\textasciicircum{}* with noncollinear x,\allowbreak{}y. Their three intensity measurements agree. Hence p\_\allowbreak{}2=\allowbreak{}0.


\entrylabel{Second check}

Dimension gives a nonzero H without assuming generic rank three;\allowbreak{} the spanning condition rules out semidefinite kernel elements. Spectral decomposition independently verifies equal squared magnitudes.


\entrylabel{Remaining obstruction}

For M>2 a nonzero kernel matrix need not have rank at most two with one eigenvalue of each sign. The requisite low-rank intersection probability remains unresolved.


Verification:\allowbreak{} author checked limited argument | Independent review:\allowbreak{} none | Full target solved:\allowbreak{} False


\clearpage

\section*{P0016  Worst square-submatrix conditioning in a Gaussian frame}

\entrylabel{Primary source}

Randomstrasse101:\allowbreak{} Open Problems of 2025;\allowbreak{} Open Problem 21 and Definition 10.1,\allowbreak{} Entry 10


\url{https://arxiv.org/pdf/2603.29571}


Source statement date:\allowbreak{} not established. Archived access:\allowbreak{} 2026-09-06T15:\allowbreak{}32:\allowbreak{}10Z.


Historical:\allowbreak{} source\_\allowbreak{}says\_\allowbreak{}open | Current:\allowbreak{} unverified | Result:\allowbreak{} attempted\_\allowbreak{}unsolved


\entrylabel{Question}

For a (2M-1)-by-M matrix of independent standard real Gaussians,\allowbreak{} determine the asymptotic size of omega(A),\allowbreak{} the minimum smallest singular value over row subsets whose complement has rank below M. In particular determine whether its decay is exponential and identify the exponential rate.


\entrylabel{Scope}

Almost surely every M-row subset is nonsingular,\allowbreak{} so the minimum reduces to M-row subsets. Interpret asymptotic size in probability;\allowbreak{} the source does not specify a convergence mode or a normalized logarithmic limit.


\subsection*{P0016-A001 | Union small-ball lower bound and one-normal upper bound}

\entrylabel{Argument or result}

For a standard M-by-M real Gaussian matrix A,\allowbreak{} sigma\_\allowbreak{}min(A)<=\allowbreak{}t implies that some column has distance at most sqrt(M)t from the span of the other columns:\allowbreak{} project a minimizing unit vector equation using a coordinate of magnitude at least 1/\allowbreak{}sqrt(M). Each such distance has the law |N(0,\allowbreak{}1)|,\allowbreak{} so P(sigma\_\allowbreak{}min<=\allowbreak{}t)<=\allowbreak{}sqrt(2/\allowbreak{}pi) M\textasciicircum{}(3/\allowbreak{}2)t. There are at most 4\textasciicircum{}M M-row submatrices of a (2M-1)-by-M frame. Therefore its minimum omega exceeds exp(-(log 4+\allowbreak{}epsilon)M) with probability tending to one. Conversely fix M-1 rows and a unit normal u. The M remaining row projections are independent N(0,\allowbreak{}1),\allowbreak{} and omega is at most their minimum absolute value,\allowbreak{} giving omega<=\allowbreak{}C log(M)/\allowbreak{}M with probability tending to one for sufficiently large C.


\entrylabel{Second check}

Full spark holds almost surely. Adding rows increases the Gram matrix,\allowbreak{} so the minimum relevant frame bound is attained among M-row subsets. Both bounds use union/\allowbreak{}inclusion inequalities and never assume independence among overlapping minors.


\entrylabel{Remaining obstruction}

The exponential lower bound and polynomial upper bound are far apart. Sharp extremes of the strongly dependent minors remain uncontrolled.


Verification:\allowbreak{} author checked limited argument | Independent review:\allowbreak{} none | Full target solved:\allowbreak{} False


\clearpage

\section*{P0017  Dimension-free log-concave Poincare or Cheeger constant}

\entrylabel{Primary source}

Randomstrasse101:\allowbreak{} Open Problems of 2025;\allowbreak{} Conjecture 30,\allowbreak{} Entry 13


\url{https://arxiv.org/pdf/2603.29571}


Source statement date:\allowbreak{} not established. Archived access:\allowbreak{} 2026-09-06T15:\allowbreak{}32:\allowbreak{}10Z.


Historical:\allowbreak{} source\_\allowbreak{}says\_\allowbreak{}open | Current:\allowbreak{} unverified | Result:\allowbreak{} attempted\_\allowbreak{}unsolved


\entrylabel{Question}

Is there a universal positive lower bound for the Euclidean Cheeger constant of every isotropic log-concave probability measure,\allowbreak{} in every dimension?


\entrylabel{Scope}

Use outer boundary liminf in the standard definition;\allowbreak{} the source's displayed halfspace formula has an S/\allowbreak{}H notational mismatch. Thin-shell and slicing results are separate and reported resolved by this source.


\subsection*{P0017-A001 | Gaussian special case by orthogonal expansion}

\entrylabel{Argument or result}

For standard Gaussian measure,\allowbreak{} expand a polynomial f in orthonormal multivariate Hermite polynomials with coefficients a\_\allowbreak{}alpha. Then Var(f)=\allowbreak{}sum\_\allowbreak{}(|alpha|>=\allowbreak{}1) a\_\allowbreak{}alpha\textasciicircum{}2 <=\allowbreak{} sum |alpha| a\_\allowbreak{}alpha\textasciicircum{}2 =\allowbreak{} E||gradient f||\textasciicircum{}2. Closure in Gaussian Sobolev norm extends the inequality to its natural domain.


\entrylabel{Second check}

The constant is attained by any nonconstant linear function,\allowbreak{} checking both normalization and dimension independence.


\entrylabel{Remaining obstruction}

A general isotropic log-concave measure lacks this diagonal Hermite generator. This special case does not prove its uniform Poincare or Cheeger bound.


Verification:\allowbreak{} author checked limited argument | Independent review:\allowbreak{} none | Full target solved:\allowbreak{} False


\clearpage

\section*{P0018  Sharp polynomial-logarithmic expected norm of graph matrices}

\entrylabel{Primary source}

Randomstrasse101:\allowbreak{} Open Problems of 2025;\allowbreak{} Conjecture 31 and Definition 14.2,\allowbreak{} Entry 14


\url{https://arxiv.org/pdf/2603.29571}


Source statement date:\allowbreak{} not established. Archived access:\allowbreak{} 2026-09-06T15:\allowbreak{}32:\allowbreak{}10Z.


Historical:\allowbreak{} source\_\allowbreak{}says\_\allowbreak{}open | Current:\allowbreak{} unverified | Result:\allowbreak{} attempted\_\allowbreak{}unsolved


\entrylabel{Question}

For each fixed graph shape alpha with specified left and right index sets,\allowbreak{} form M\_\allowbreak{}alpha by summing injective labelings with products of independent Rademacher edge variables. Must E||M\_\allowbreak{}alpha|| be Theta(n\textasciicircum{}f(alpha) (log n)\textasciicircum{}g(alpha)) for some shape-dependent exponents?


\entrylabel{Scope}

The target is the general existence of sharp exponents,\allowbreak{} not a separate problem for every shape. Definitions of overlapping or internal vertex sets must follow the full cited graph-matrix convention.


\subsection*{P0018-A001 | Moment expansion with a crude operator-norm bound}

\entrylabel{Argument or result}

For every finite random graph matrix,\allowbreak{} E||M|| <=\allowbreak{} (E||M||\_\allowbreak{}F\textasciicircum{}2)\textasciicircum{}(1/\allowbreak{}2). Expanding each entry as a sum of Rademacher edge monomials,\allowbreak{} cross terms vanish unless the two monomials use every edge with even combined multiplicity. For a fixed shape the remaining label identifications can be counted to obtain a polynomial upper bound in n.


\entrylabel{Second check}

Orthogonality concerns equality of edge monomials,\allowbreak{} not equality of vertex labelings. Distinct labelings can produce the same edge set and must not be discarded.


\entrylabel{Remaining obstruction}

The Frobenius bound can lose powers of n and cannot decide the sharp operator norm or the existence of one power-log asymptotic form.


Verification:\allowbreak{} author checked limited argument | Independent review:\allowbreak{} none | Full target solved:\allowbreak{} False


\clearpage

\section*{P0019  Brownian exit approach directions in higher dimensions}

\entrylabel{Primary source}

Krzysztof Burdzy:\allowbreak{} My favorite open problems;\allowbreak{} Problem 1


\url{https://sites.math.washington.edu/~burdzy/open_mathjax.php}


Source statement date:\allowbreak{} not established. Archived access:\allowbreak{} 2026-09-06T16:\allowbreak{}41:\allowbreak{}36Z.


Historical:\allowbreak{} source\_\allowbreak{}says\_\allowbreak{}open | Current:\allowbreak{} unverified | Result:\allowbreak{} attempted\_\allowbreak{}unsolved


\entrylabel{Question}

For Brownian motion exiting a domain in dimension d>2,\allowbreak{} is the set of limiting approach directions at the exit point almost surely a sphere or a hemisphere?


\entrylabel{Scope}

The question is meaningful on finite exit;\allowbreak{} either condition on finite exit or restrict to bounded domains. Do not call infinite exit a counterexample to the intended geometric question.


\subsection*{P0019-A001 | Smooth-boundary directional containment}

\entrylabel{Argument or result}

At an exit point z where the domain boundary is C\textasciicircum{}1 with outward unit normal n,\allowbreak{} local differentiability gives n dot (z-x)>=\allowbreak{}-o(|z-x|) for x in the domain approaching z. Thus every limit of (z-B\_\allowbreak{}s)/\allowbreak{}|z-B\_\allowbreak{}s| as s increases to the exit time lies in the closed hemisphere \{u:\allowbreak{}n dot u>=\allowbreak{}0\}.


\entrylabel{Second check}

The sign is checked on a half-space,\allowbreak{} where the vector from an interior point to the exit point has positive outward-normal component.


\entrylabel{Remaining obstruction}

Containment does not prove that all hemisphere directions occur,\allowbreak{} and there is no supporting tangent half-space at the arbitrary boundaries in the full question.


Verification:\allowbreak{} author checked limited argument | Independent review:\allowbreak{} none | Full target solved:\allowbreak{} False


\clearpage

\section*{P0020  Finite-crossing arcs through a planar Brownian trace}

\entrylabel{Primary source}

Krzysztof Burdzy:\allowbreak{} My favorite open problems;\allowbreak{} Problem 2(i)


\url{https://sites.math.washington.edu/~burdzy/open_mathjax.php}


Source statement date:\allowbreak{} not established. Archived access:\allowbreak{} 2026-09-06T16:\allowbreak{}41:\allowbreak{}36Z.


Historical:\allowbreak{} source\_\allowbreak{}says\_\allowbreak{}open | Current:\allowbreak{} unverified | Result:\allowbreak{} attempted\_\allowbreak{}unsolved


\entrylabel{Question}

Almost surely,\allowbreak{} can every two points outside the trace of planar Brownian motion on [0,\allowbreak{}1] be connected by a Jordan arc meeting the trace at only finitely many points?


\entrylabel{Scope}

The quantifier is simultaneous over all pairs of outside points,\allowbreak{} not a single deterministic pair.


\subsection*{P0020-A001 | Reduction to component adjacency and access}

\entrylabel{Argument or result}

If both endpoints lie in the same component of the complement of the Brownian trace,\allowbreak{} an open connected planar set is polygonally connected;\allowbreak{} removing loops gives a simple arc avoiding the trace entirely. If an arc crosses the trace finitely many times,\allowbreak{} its complementary intervals give a finite sequence of complement components with touching closures at successive crossing points.


\entrylabel{Second check}

A closure contact alone does not ensure an accessible common boundary point. This is why the necessary finite-chain condition was not promoted to an equivalence.


\entrylabel{Remaining obstruction}

Need finite accessible chains for every pair of components,\allowbreak{} simultaneously for the Brownian trace.


Verification:\allowbreak{} author checked limited argument | Independent review:\allowbreak{} none | Full target solved:\allowbreak{} False


\clearpage

\section*{P0021  Interior remainder of planar Brownian trace}

\entrylabel{Primary source}

Krzysztof Burdzy:\allowbreak{} My favorite open problems;\allowbreak{} Problem 2(ii)


\url{https://sites.math.washington.edu/~burdzy/open_mathjax.php}


Source statement date:\allowbreak{} not established. Archived access:\allowbreak{} 2026-09-06T16:\allowbreak{}41:\allowbreak{}36Z.


Historical:\allowbreak{} source\_\allowbreak{}says\_\allowbreak{}open | Current:\allowbreak{} unverified | Result:\allowbreak{} attempted\_\allowbreak{}unsolved


\entrylabel{Question}

Delete from a planar Brownian trace on [0,\allowbreak{}1] all boundaries of connected components of its complement. Is the remaining set totally disconnected almost surely?


\entrylabel{Scope}

This is a different topological property from Problem 2(i),\allowbreak{} although a negative answer is stated to imply a negative answer there.


\subsection*{P0021-A001 | Failure of transferring total disconnectedness through the path map}

\entrylabel{Argument or result}

A possible approach is to identify a totally disconnected set of Brownian times surviving boundary removal. That would not suffice:\allowbreak{} the Cantor function restricted to the middle-third Cantor set maps this totally disconnected compact set continuously onto [0,\allowbreak{}1]. Continuous images can be connected.


\entrylabel{Second check}

The restriction is onto because every value of the Cantor function is attained on the Cantor set,\allowbreak{} including plateau values at endpoints.


\entrylabel{Remaining obstruction}

The spatial remainder needs direct Brownian geometry. Showing it contains no simple arc would also be weaker than showing every connected component is a singleton.


Verification:\allowbreak{} author checked limited argument | Independent review:\allowbreak{} none | Full target solved:\allowbreak{} False


\clearpage

\section*{P0022  Percolation dimension of a planar Brownian trace}

\entrylabel{Primary source}

Krzysztof Burdzy:\allowbreak{} My favorite open problems;\allowbreak{} Problem 3


\url{https://sites.math.washington.edu/~burdzy/open_mathjax.php}


Source statement date:\allowbreak{} not established. Archived access:\allowbreak{} 2026-09-06T16:\allowbreak{}41:\allowbreak{}36Z.


Historical:\allowbreak{} source\_\allowbreak{}says\_\allowbreak{}open | Current:\allowbreak{} unverified | Result:\allowbreak{} attempted\_\allowbreak{}unsolved


\entrylabel{Question}

Is the infimum of Hausdorff dimensions of nontrivial Jordan arcs contained in a planar Brownian trace on [0,\allowbreak{}1] equal to one almost surely?


\entrylabel{Scope}

A nontrivial arc contains at least two distinct points. The lower bound one alone is not a solution.


\subsection*{P0022-A001 | Universal dimension lower bound for connecting sets}

\entrylabel{Argument or result}

Any connected set containing distinct x and y has Hausdorff dimension at least one. Orthogonal projection onto their joining line is Lipschitz and has connected image containing the interval between the projections,\allowbreak{} so its dimension is at least one.


\entrylabel{Second check}

This argument does not require rectifiability or a parametrized simple curve,\allowbreak{} and hence applies to the connectivity notion in a percolation-dimension lower bound.


\entrylabel{Remaining obstruction}

Polygonal approximations of a Brownian path are not subsets of its trace. The missing step is constructing sufficiently thin connecting subsets inside the trace itself.


Verification:\allowbreak{} author checked limited argument | Independent review:\allowbreak{} none | Full target solved:\allowbreak{} False


\clearpage

\section*{P0023  Efficient reflected-Brownian coupling in acute triangles}

\entrylabel{Primary source}

Krzysztof Burdzy:\allowbreak{} My favorite open problems;\allowbreak{} Problem 4


\url{https://sites.math.washington.edu/~burdzy/open_mathjax.php}


Source statement date:\allowbreak{} not established. Archived access:\allowbreak{} 2026-09-06T16:\allowbreak{}41:\allowbreak{}36Z.


Historical:\allowbreak{} source\_\allowbreak{}says\_\allowbreak{}open | Current:\allowbreak{} unverified | Result:\allowbreak{} attempted\_\allowbreak{}unsolved


\entrylabel{Question}

In any acute triangle,\allowbreak{} can reflected Brownian motions from distinct starting points be coupled so their meeting-time tail decays at every exponential rate strictly below the first positive Neumann spectral rate?


\entrylabel{Scope}

Check the generator convention in Burdzy-Kendall (2000):\allowbreak{} the page calls its rate a Laplacian eigenvalue. A factor-two convention discrepancy is not a solution of the coupling question.


\subsection*{P0023-A001 | Spectral lower bound on coupling tails}

\entrylabel{Argument or result}

Let phi be a bounded nonconstant Neumann eigenfunction with semigroup P\_\allowbreak{}t phi=\allowbreak{}e\textasciicircum{}(-lambda t) phi. For any coalescent coupling with meeting time T,\allowbreak{} |P\_\allowbreak{}t phi(x)-P\_\allowbreak{}t phi(y)|<=\allowbreak{}2||phi||\_\allowbreak{}infinity P(T>t). Choosing phi(x)!=\allowbreak{}phi(y) yields P(T>t)>=\allowbreak{}c e\textasciicircum{}(-lambda t).


\entrylabel{Second check}

The exponent is the eigenvalue of the actual reflected-Brownian generator:\allowbreak{} using Delta rather than Delta/\allowbreak{}2 changes it by a factor of two.


\entrylabel{Remaining obstruction}

This proves that the spectral rate is an upper limit on coupling efficiency,\allowbreak{} but gives no coupling attaining arbitrarily close to that rate in an acute triangle.


Verification:\allowbreak{} author checked limited argument | Independent review:\allowbreak{} none | Full target solved:\allowbreak{} False


\clearpage

\section*{P0024  Failure of synchronous contraction in a bounded planar domain}

\entrylabel{Primary source}

Krzysztof Burdzy:\allowbreak{} My favorite open problems;\allowbreak{} Problem 5(i)


\url{https://sites.math.washington.edu/~burdzy/open_mathjax.php}


Source statement date:\allowbreak{} not established. Archived access:\allowbreak{} 2026-09-06T16:\allowbreak{}41:\allowbreak{}36Z.


Historical:\allowbreak{} source\_\allowbreak{}says\_\allowbreak{}open | Current:\allowbreak{} unverified | Result:\allowbreak{} attempted\_\allowbreak{}unsolved


\entrylabel{Question}

Does some bounded connected smooth planar domain admit synchronously coupled reflected Brownian motions whose distance has positive limsup with positive probability?


\entrylabel{Scope}

The page states that any such domain must have at least two holes by existing results.


\subsection*{P0024-A001 | Synchronous contraction on convex domains}

\entrylabel{Argument or result}

For two reflected motions driven by the same Brownian path,\allowbreak{} the noise cancels in their difference. With inward normals,\allowbreak{} d|X-Y|\textasciicircum{}2=\allowbreak{}2<X-Y,\allowbreak{}n(X)>dL\_\allowbreak{}X-2<X-Y,\allowbreak{}n(Y)>dL\_\allowbreak{}Y. Convexity makes both contributions nonpositive,\allowbreak{} so the distance is nonincreasing.


\entrylabel{Second check}

Each sign follows from the supporting half-space inequality at the boundary point. The calculation uses finite-variation reflection terms and has no missing quadratic-variation term.


\entrylabel{Remaining obstruction}

Nonincrease alone does not prove convergence to zero. Nonconvex domains can have both signs and require control of the long-run reflection geometry.


Verification:\allowbreak{} author checked limited argument | Independent review:\allowbreak{} none | Full target solved:\allowbreak{} False


\clearpage

\section*{P0025  Synchronous reflection outside a disk}

\entrylabel{Primary source}

Krzysztof Burdzy:\allowbreak{} My favorite open problems;\allowbreak{} Problem 5(ii)


\url{https://sites.math.washington.edu/~burdzy/open_mathjax.php}


Source statement date:\allowbreak{} not established. Archived access:\allowbreak{} 2026-09-06T16:\allowbreak{}41:\allowbreak{}36Z.


Historical:\allowbreak{} source\_\allowbreak{}says\_\allowbreak{}open | Current:\allowbreak{} unverified | Result:\allowbreak{} attempted\_\allowbreak{}unsolved


\entrylabel{Question}

For synchronously coupled reflected Brownian motions in the exterior of a nondegenerate closed disk,\allowbreak{} can their distance have positive limsup with positive probability?


\entrylabel{Scope}

This unbounded exterior geometry differs from the bounded-domain existence question.


\subsection*{P0025-A001 | Exterior-disk reflection has both distance signs}

\entrylabel{Argument or result}

At x=\allowbreak{}(R,\allowbreak{}0) on the inner circular boundary of the exterior domain,\allowbreak{} its inward-to-domain normal is (1,\allowbreak{}0). For y=\allowbreak{}(r,\allowbreak{}0),\allowbreak{} r>R,\allowbreak{} <x-y,\allowbreak{}n>=\allowbreak{}R-r<0;\allowbreak{} for y=\allowbreak{}(-r,\allowbreak{}0) it is R+\allowbreak{}r>0. Hence a reflection at X can either shrink or enlarge the synchronous separation.


\entrylabel{Second check}

Both test locations lie in the exterior and the normal points into that same domain,\allowbreak{} fixing the otherwise easy sign error.


\entrylabel{Remaining obstruction}

A global monotone-distance proof is unavailable. The question turns on boundary occupation and angular correlations over indefinitely long times.


Verification:\allowbreak{} author checked limited argument | Independent review:\allowbreak{} none | Full target solved:\allowbreak{} False


\clearpage

\section*{P0026  Non-explosion of Fleming-Viot Brownian particles in arbitrary domains}

\entrylabel{Primary source}

Krzysztof Burdzy:\allowbreak{} My favorite open problems;\allowbreak{} Problem 6 with 2024 update


\url{https://sites.math.washington.edu/~burdzy/open_mathjax.php}


Source statement date:\allowbreak{} not established. Archived access:\allowbreak{} 2026-09-06T16:\allowbreak{}41:\allowbreak{}36Z.


Historical:\allowbreak{} source\_\allowbreak{}says\_\allowbreak{}open | Current:\allowbreak{} unverified | Result:\allowbreak{} attempted\_\allowbreak{}unsolved


\entrylabel{Question}

For N>=\allowbreak{}3 Brownian particles in any bounded open connected domain,\allowbreak{} killed at the boundary and immediately cloned from a uniformly selected survivor,\allowbreak{} are the successive death times almost surely non-accumulating?


\entrylabel{Scope}

N=\allowbreak{}2 is explicitly reported solved by Kwasnicki (2024). The earlier general proof had an irreparable error;\allowbreak{} it is not admissible evidence for the general result.


\subsection*{P0026-A001 | Attempted boundary-distance Lyapunov proof}

\entrylabel{Argument or result}

The candidate Lyapunov function -sum\_\allowbreak{}i log dist(X\_\allowbreak{}i,\allowbreak{}boundary D) becomes singular when particles approach the boundary,\allowbreak{} but an Ito argument needs derivative bounds. Even for a smooth domain,\allowbreak{} the distance is nonsmooth on its medial set;\allowbreak{} for arbitrary domains there may be no uniformly controlled smooth substitute. At a cloning event another particle may be still closer to the boundary,\allowbreak{} so the jump is not uniformly favorable.


\entrylabel{Second check}

The known N=\allowbreak{}2 result is excluded from this target. A pair comparison does not transfer automatically to N>=\allowbreak{}3 because a particle can clone the third particle rather than its proposed comparator.


\entrylabel{Remaining obstruction}

No Lyapunov drift bound uniform over arbitrary boundary geometry and all cloning configurations was obtained.


Verification:\allowbreak{} author checked limited argument | Independent review:\allowbreak{} none | Full target solved:\allowbreak{} False


\clearpage

\section*{P0027  Rigidity of shy reflected-Brownian couplings}

\entrylabel{Primary source}

Krzysztof Burdzy:\allowbreak{} My favorite open problems;\allowbreak{} Problem 7


\url{https://sites.math.washington.edu/~burdzy/open_mathjax.php}


Source statement date:\allowbreak{} not established. Archived access:\allowbreak{} 2026-09-06T16:\allowbreak{}41:\allowbreak{}36Z.


Historical:\allowbreak{} source\_\allowbreak{}says\_\allowbreak{}open | Current:\allowbreak{} unverified | Result:\allowbreak{} attempted\_\allowbreak{}unsolved


\entrylabel{Question}

If reflected Brownian motions in a bounded connected domain can remain a positive distance apart forever with positive probability,\allowbreak{} must there be such a coupling with one path a fixed deterministic function of the other at each time?


\entrylabel{Scope}

Regularity sufficient to define the intended reflected Brownian motion needs to be fixed from the cited background. A metric-graph counterexample on the page is not a Euclidean-domain counterexample.


\subsection*{P0027-A001 | Explicit deterministic shy coupling on an annulus}

\entrylabel{Argument or result}

If an isometry F preserves a reflected-Brownian domain and has inf\_\allowbreak{}x |x-F(x)|>0,\allowbreak{} take Y\_\allowbreak{}t=\allowbreak{}F(X\_\allowbreak{}t). Orthogonal transformation of Brownian noise and equivariance of inward reflection show Y has the correct reflected law,\allowbreak{} while separation stays positive. On R<|x|<R\_\allowbreak{}out,\allowbreak{} F(x)=\allowbreak{}-x yields |X\_\allowbreak{}t-Y\_\allowbreak{}t|>=\allowbreak{}2R.


\entrylabel{Second check}

The annulus excludes the origin;\allowbreak{} a disk would have zero minimum displacement and would not satisfy this check.


\entrylabel{Remaining obstruction}

This verifies a nonempty deterministic class. It does not turn an arbitrary shy coupling into a single state map with the required path consistency.


Verification:\allowbreak{} author checked limited argument | Independent review:\allowbreak{} none | Full target solved:\allowbreak{} False


\clearpage

\section*{P0028  Bounded independent Brownian pieces and two-sided Brownian motion}

\entrylabel{Primary source}

Krzysztof Burdzy:\allowbreak{} My favorite open problems;\allowbreak{} Problem 8;\allowbreak{} compare Burdzy-Scheutzow (2014) Problem 7.7


\url{https://sites.math.washington.edu/~burdzy/open_mathjax.php}


Source statement date:\allowbreak{} not established. Archived access:\allowbreak{} 2026-09-06T16:\allowbreak{}41:\allowbreak{}36Z.


Historical:\allowbreak{} source\_\allowbreak{}says\_\allowbreak{}open | Current:\allowbreak{} unverified | Result:\allowbreak{} attempted\_\allowbreak{}unsolved


\entrylabel{Question}

Concatenate independent Brownian path segments cut at stopping times indexed by all integers,\allowbreak{} with lengths uniformly bounded by a deterministic constant and divergent sums in both directions. Is the resulting path two-sided Brownian after some random time shift?


\entrylabel{Scope}

The paper's related backward-Brownian question and its almost-sure random bound are recorded as related variants,\allowbreak{} not an extra count. Independent lengths alone are not the same as stopping times depending on their Brownian paths.


\subsection*{P0028-A001 | Stationary Palm construction in the identically distributed finite-mean case}

\entrylabel{Argument or result}

For iid stopped Brownian segments with 0<E T<infinity,\allowbreak{} the cycle sequence is ergodic under shifting its index. A stationary occupation law is Q(F)=\allowbreak{}E[integral\_\allowbreak{}0\textasciicircum{}T F(theta\_\allowbreak{}s X) ds]/\allowbreak{}E T. Shifting by t changes the numerator by E[integral\_\allowbreak{}T\textasciicircum{}(T+\allowbreak{}t) F(theta\_\allowbreak{}s X) ds-integral\_\allowbreak{}0\textasciicircum{}t F(theta\_\allowbreak{}s X) ds]=\allowbreak{}0 because the concatenation rooted at the next cycle has the same Palm law. Tilting the iid cycle sequence by T\_\allowbreak{}0/\allowbreak{}E T and choosing a uniform time in that cycle realizes this law;\allowbreak{} the forward filling lemma supplies that tilt by a finite randomized index shift.


\entrylabel{Second check}

For each deterministic s,\allowbreak{} \{s<T\} is measurable from the segment up to s. Future Brownian increments followed by independent stopped segments retain the Brownian law and are independent of the past. Integrating over s gives a Brownian future under Q;\allowbreak{} stationarity then supplies a Brownian past and its independence. This is a proof outline pending full artifact and earlier-literature comparison,\allowbreak{} not a newly claimed theorem.


\entrylabel{Remaining obstruction}

The actual question allows independent non-identically distributed segments. There is then no invariant iid index law or common Palm mean to which this argument applies.


Verification:\allowbreak{} author checked limited argument | Independent review:\allowbreak{} none | Full target solved:\allowbreak{} False


\subsection*{P0028-A002 | Primary-source check of the iid finite-mean reduction}

\entrylabel{Argument or result}

Burdzy-Scheutzow (2014),\allowbreak{} Theorem 5.1 explicitly proves that strongly decomposable processes with finite mean segment length are two-sided Brownian motions. P0028-A001 rederives this known special case and is not a novel solution.


\entrylabel{Remaining obstruction}

The independent non-identically distributed bounded-length question remains unresolved here.


Verification:\allowbreak{} source theorem checked | Independent review:\allowbreak{} none | Full target solved:\allowbreak{} False


\clearpage

\section*{P0029  Separation of two peaks in random square-grid labelings}

\entrylabel{Primary source}

Krzysztof Burdzy:\allowbreak{} My favorite open problems;\allowbreak{} Problem 9


\url{https://sites.math.washington.edu/~burdzy/open_mathjax.php}


Source statement date:\allowbreak{} not established. Archived access:\allowbreak{} 2026-09-06T16:\allowbreak{}41:\allowbreak{}36Z.


Historical:\allowbreak{} source\_\allowbreak{}says\_\allowbreak{}open | Current:\allowbreak{} unverified | Result:\allowbreak{} attempted\_\allowbreak{}unsolved


\entrylabel{Question}

Uniformly label an n-by-n square grid and condition on exactly two local maxima. Does their graph distance divided by n converge in distribution to zero?


\entrylabel{Scope}

This is the literal displayed question,\allowbreak{} despite the heading about repulsion. Verify against Twin peaks (2020) before correcting the sign or reporting a counterexample.


\subsection*{P0029-A001 | Exact small-grid count and an orientation reformulation}

\entrylabel{Argument or result}

On the 2-by-2 grid (a four-cycle),\allowbreak{} exactly two peaks must be opposite and carry labels 3 and 4. There are two choices of diagonal,\allowbreak{} two orders of the top labels,\allowbreak{} and two orders of the bottom labels:\allowbreak{} eight of the 24 labelings. Their graph distance is always two. More generally orient each edge toward its larger label. Two peaks means precisely two sinks of the induced acyclic orientation;\allowbreak{} the number of labelings yielding a fixed orientation is its number of linear extensions.


\entrylabel{Second check}

Adjacent vertices cannot both be peaks. In the four-cycle,\allowbreak{} a proposed peak with label 2 would have to exceed both other nonpeak labels,\allowbreak{} which is impossible;\allowbreak{} this checks the count without enumeration software.


\entrylabel{Remaining obstruction}

The small-grid calculation gives no large-n scale. The orientation sum requires sharp weighted counts of acyclic orientations with two prescribed sinks. The source title says repel while its displayed question asks distance/\allowbreak{}n -> 0;\allowbreak{} that wording discrepancy is retained,\allowbreak{} not silently corrected.


Verification:\allowbreak{} author checked limited argument | Independent review:\allowbreak{} none | Full target solved:\allowbreak{} False


\clearpage

\section*{P0030  Compatibility of forward and backward Brownian decompositions}

\entrylabel{Primary source}

Forward Brownian motion;\allowbreak{} Problems 7.1-7.3,\allowbreak{} PDF pages 25-26


\url{https://math.washington.edu/sites/math/files/documents/research/1302.6958.pdf}


Source statement date:\allowbreak{} not established. Archived access:\allowbreak{} 2026-09-06T16:\allowbreak{}41:\allowbreak{}55Z.


Historical:\allowbreak{} source\_\allowbreak{}says\_\allowbreak{}open | Current:\allowbreak{} unverified | Result:\allowbreak{} attempted\_\allowbreak{}unsolved


\entrylabel{Question}

Does a decomposable forward Brownian motion which is also backward Brownian necessarily become two-sided Brownian after a random shift? The source also asks weaker related formulations using independence of past and future at a random time.


\entrylabel{Scope}

The three linked formulations are grouped as one compatibility problem. A general process that is both forward and backward Brownian without the extra assumptions is already counterexampled in the paper.


\subsection*{P0030-A001 | Common-splice compatibility obstruction}

\entrylabel{Argument or result}

Two independent standard Brownian half-paths attached at the same deterministic time give a two-sided Brownian motion. For a general process certified separately as forward and backward Brownian,\allowbreak{} the certifying random times can differ. Moving either certificate to another random time can bias its marginal law,\allowbreak{} as the selection-density calculation in P0032 demonstrates.


\entrylabel{Second check}

Independence of two halves at a proposed joining time is insufficient if their Brownian marginal laws have been lost there.


\entrylabel{Remaining obstruction}

No construction aligns both certificates at one time while preserving both marginal laws and their independence.


Verification:\allowbreak{} author checked limited argument | Independent review:\allowbreak{} none | Full target solved:\allowbreak{} False


\clearpage

\section*{P0031  Sparse first-passage pieces in a forward Brownian concatenation}

\entrylabel{Primary source}

Forward Brownian motion;\allowbreak{} Problem 7.4,\allowbreak{} PDF page 26


\url{https://math.washington.edu/sites/math/files/documents/research/1302.6958.pdf}


Source statement date:\allowbreak{} not established. Archived access:\allowbreak{} 2026-09-06T16:\allowbreak{}41:\allowbreak{}55Z.


Historical:\allowbreak{} source\_\allowbreak{}says\_\allowbreak{}open | Current:\allowbreak{} unverified | Result:\allowbreak{} attempted\_\allowbreak{}unsolved


\entrylabel{Question}

In a two-sided independent concatenation,\allowbreak{} use unit-length Brownian segments except at increasingly rare negative-time indices where segments stop on first hitting -1. Characterize frequencies giving backward or two-sided Brownian motion.


\entrylabel{Scope}

The exact deterministic or random schedule is part of the classification;\allowbreak{} do not invent a schedule and count it as another sourced question.


\subsection*{P0031-A001 | Sparse exceptional hitting-time clock estimate}

\entrylabel{Argument or result}

For a standard Brownian first hitting time tau of -1,\allowbreak{} reflection gives P(tau>t)=\allowbreak{}2 Phi(1/\allowbreak{}sqrt(t))-1 \textasciitilde{} sqrt(2/\allowbreak{}pi)t\textasciicircum{}(-1/\allowbreak{}2),\allowbreak{} hence E tau=\allowbreak{}infinity. If exceptional segment indices are 2\textasciicircum{}k,\allowbreak{} then sum\_\allowbreak{}k P(tau\_\allowbreak{}k>epsilon 2\textasciicircum{}k)<infinity. Borel-Cantelli yields tau\_\allowbreak{}k/\allowbreak{}2\textasciicircum{}k ->0 almost surely. Splitting the sum into finitely many early terms and a geometrically weighted tail gives sum\_\allowbreak{}(k<=\allowbreak{}K) tau\_\allowbreak{}k/\allowbreak{}2\textasciicircum{}K ->0.


\entrylabel{Second check}

Independence is not needed for this Borel-Cantelli upper bound. Infinite means do not contradict the almost-sure sparse-clock estimate.


\entrylabel{Remaining obstruction}

A negligible exceptional time fraction is an asymptotic statement;\allowbreak{} it does not construct an exact backward Brownian law after a finite random shift. The full frequency classification remains open here.


Verification:\allowbreak{} author checked limited argument | Independent review:\allowbreak{} none | Full target solved:\allowbreak{} False


\clearpage

\section*{P0032  Embedding an equivalent Brownian law by adding a finite prefix}

\entrylabel{Primary source}

Forward Brownian motion;\allowbreak{} Problem 7.5,\allowbreak{} PDF page 26


\url{https://math.washington.edu/sites/math/files/documents/research/1302.6958.pdf}


Source statement date:\allowbreak{} not established. Archived access:\allowbreak{} 2026-09-06T16:\allowbreak{}41:\allowbreak{}55Z.


Historical:\allowbreak{} source\_\allowbreak{}says\_\allowbreak{}open | Current:\allowbreak{} unverified | Result:\allowbreak{} candidate\_\allowbreak{}full\_\allowbreak{}solution


\entrylabel{Question}

Given a process whose law on nonnegative time is equivalent to Wiener measure,\allowbreak{} can a random finite prefix be added,\allowbreak{} on an enlarged probability space if necessary,\allowbreak{} to produce a process with Wiener law?


\entrylabel{Scope}

No independence or Brownian stopping-time restriction is explicit in the source. The candidate proof uses a finite nonnegative integer random time depending on entire paths and auxiliary randomness.


\subsection*{P0032-A001 | Forward density matching on the ergodic two-sided Wiener shift}

\entrylabel{Argument or result}

Complete candidate construction of a finite nonnegative integer prefix time for any suffix law absolutely continuous with respect to Wiener measure


\entrylabel{Second check}

Two-state cyclic example checks the pushforward direction and normalization;\allowbreak{} finite-window rejection sampling gives an independent special-case construction


\entrylabel{Remaining obstruction}

No mathematical gap currently identified under the stated interpretation;\allowbreak{} source interpretation,\allowbreak{} measurability details,\allowbreak{} and earlier filling/\allowbreak{}embedding literature still under review


Verification:\allowbreak{} candidate under author and literature checks | Independent review:\allowbreak{} none | Full target solved:\allowbreak{} False


\entrylabel{Full proof packet}

proofs/\allowbreak{}P0032\_\allowbreak{}brownian\_\allowbreak{}prefix\_\allowbreak{}candidate.md


\clearpage

\section*{P0033  Ergodicity of stationary Galton-Watson flows}

\entrylabel{Primary source}

Probability on Trees and Networks,\allowbreak{} paperback electronic version 19 August 2026;\allowbreak{} Question 17.18,\allowbreak{} printed p.575,\allowbreak{} PDF p.592


\url{https://rdlyons.pages.iu.edu/prbtree/book_pb.pdf}


Source statement date:\allowbreak{} not established. Archived access:\allowbreak{} 2026-09-06T15:\allowbreak{}30:\allowbreak{}43Z.


Historical:\allowbreak{} source\_\allowbreak{}says\_\allowbreak{}open | Current:\allowbreak{} unverified | Result:\allowbreak{} candidate\_\allowbreak{}full\_\allowbreak{}solution


\entrylabel{Question}

If a consistent equivariant Borel flow rule has a stationary probability measure equivalent to Galton-Watson law,\allowbreak{} must the associated subtree Markov chain be ergodic?


\entrylabel{Scope}

The definition permits zero child weights. Strict positivity is a separate sufficient hypothesis in Proposition 17.17. Exercise 17.22 only treats a merely absolutely continuous stationary measure.


\subsection*{P0033-A001 | Endogenous AND/\allowbreak{}OR tree coloring and conditional subtree selection}

\entrylabel{Argument or result}

Candidate counterexample:\allowbreak{} GW itself stationary,\allowbreak{} with two invariant color classes of mass one half


\entrylabel{Second check}

Exact second-moment recursion proves color measurability;\allowbreak{} full conditional subtree law proves stationarity;\allowbreak{} source definition expressly permits zero child weights


\entrylabel{Remaining obstruction}

No mathematical gap currently identified;\allowbreak{} full author recheck and prior tree-automaton literature pending


Verification:\allowbreak{} candidate under author and literature checks | Independent review:\allowbreak{} none | Full target solved:\allowbreak{} False


\entrylabel{Full proof packet}

proofs/\allowbreak{}P0033\_\allowbreak{}stationary\_\allowbreak{}flow\_\allowbreak{}counterexample\_\allowbreak{}candidate.md


\clearpage

\section*{P0034  Recurrence of the planar horizontal-new vertical-old walk}

\entrylabel{Primary source}

Probability on Trees and Networks,\allowbreak{} paperback electronic version 19 August 2026;\allowbreak{} Question 2.27,\allowbreak{} printed p.56


\url{https://rdlyons.pages.iu.edu/prbtree/book_pb.pdf}


Source statement date:\allowbreak{} not established. Archived access:\allowbreak{} 2026-09-06T15:\allowbreak{}30:\allowbreak{}43Z.


Historical:\allowbreak{} source\_\allowbreak{}says\_\allowbreak{}open | Current:\allowbreak{} unverified | Result:\allowbreak{} attempted\_\allowbreak{}unsolved


\entrylabel{Question}

Start at the origin in Z\textasciicircum{}2;\allowbreak{} after first visiting a site step horizontally with either sign equally likely,\allowbreak{} and after a repeat visit step vertically with either sign equally likely. Does any site receive infinitely many visits with positive probability?


\entrylabel{Scope}

The source indexing omits the first step;\allowbreak{} use its intended new-site rule also at time zero. The analogous four- and three-dimensional results do not decide this planar model.


\subsection*{P0034-A001 | Martingale estimates for the axis-switching walk}

\entrylabel{Argument or result}

The axis of the next step is predictable from the past,\allowbreak{} while its sign is fair. Both coordinates are martingales,\allowbreak{} their cross quadratic covariation is zero,\allowbreak{} and E|X\_\allowbreak{}n|\textasciicircum{}2=\allowbreak{}n. Bounded increments and Azuma give summable bounds for P(|X\_\allowbreak{}n|>epsilon n),\allowbreak{} so X\_\allowbreak{}n/\allowbreak{}n tends to zero almost surely.


\entrylabel{Second check}

The predictable quadratic variation of the horizontal coordinate counts first visits,\allowbreak{} and the vertical one counts repeat visits;\allowbreak{} their sum is exactly n.


\entrylabel{Remaining obstruction}

Zero speed and a diffusive second moment do not imply recurrence. The occupation-dependent balance of the two coordinate clocks remains uncontrolled.


Verification:\allowbreak{} author checked limited argument | Independent review:\allowbreak{} none | Full target solved:\allowbreak{} False


\clearpage

\section*{P0035  Positive speed of inward-biased walks below the Cayley growth rate}

\entrylabel{Primary source}

Probability on Trees and Networks,\allowbreak{} paperback electronic version 19 August 2026;\allowbreak{} Question 3.11,\allowbreak{} printed p.91


\url{https://rdlyons.pages.iu.edu/prbtree/book_pb.pdf}


Source statement date:\allowbreak{} not established. Archived access:\allowbreak{} 2026-09-06T15:\allowbreak{}30:\allowbreak{}43Z.


Historical:\allowbreak{} source\_\allowbreak{}says\_\allowbreak{}open | Current:\allowbreak{} unverified | Result:\allowbreak{} attempted\_\allowbreak{}unsolved


\entrylabel{Question}

For a Cayley graph of exponential growth rate b and 1<lambda<b,\allowbreak{} does the distance from the identity divided by time for RW\_\allowbreak{}lambda converge almost surely to a strictly positive speed?


\entrylabel{Scope}

RW\_\allowbreak{}lambda uses the source's radial conductances,\allowbreak{} favoring motion toward the identity for lambda>1. The case lambda=\allowbreak{}1 may have zero speed and is excluded.


\subsection*{P0035-A001 | Exact radial speed on a regular tree}

\entrylabel{Argument or result}

On the (b+\allowbreak{}1)-regular rooted tree,\allowbreak{} away from the root the radial walk steps out with probability b/\allowbreak{}(b+\allowbreak{}lambda) and in with probability lambda/\allowbreak{}(b+\allowbreak{}lambda). For lambda<b its speed is (b-lambda)/\allowbreak{}(b+\allowbreak{}lambda)>0:\allowbreak{} the root is visited finitely often and the bounded-difference martingale law of large numbers applies to the radial drift.


\entrylabel{Second check}

The speed tends to zero at lambda=\allowbreak{}b and to one at lambda=\allowbreak{}0,\allowbreak{} checking the normalization.


\entrylabel{Remaining obstruction}

General Cayley graphs have dead ends and nonuniform distance-level neighborhoods. Transience alone does not control time spent in traps.


Verification:\allowbreak{} author checked limited argument | Independent review:\allowbreak{} none | Full target solved:\allowbreak{} False


\clearpage

\section*{P0036  Percolation threshold on arbitrary geodesic Cayley spanning trees}

\entrylabel{Primary source}

Probability on Trees and Networks,\allowbreak{} paperback electronic version 19 August 2026;\allowbreak{} Question 3.12,\allowbreak{} printed p.91


\url{https://rdlyons.pages.iu.edu/prbtree/book_pb.pdf}


Source statement date:\allowbreak{} not established. Archived access:\allowbreak{} 2026-09-06T15:\allowbreak{}30:\allowbreak{}43Z.


Historical:\allowbreak{} source\_\allowbreak{}says\_\allowbreak{}open | Current:\allowbreak{} unverified | Result:\allowbreak{} attempted\_\allowbreak{}unsolved


\entrylabel{Question}

If T is any geodesic spanning tree rooted at the identity of a Cayley graph of growth rate b,\allowbreak{} must br(T)=\allowbreak{}b,\allowbreak{} equivalently must its Bernoulli bond percolation threshold be 1/\allowbreak{}b?


\entrylabel{Scope}

The branching-number/\allowbreak{}percolation equivalence is a theorem of the source. The question is stated there geometrically;\allowbreak{} its probability formulation is recorded transparently. Lexicographically minimal spanning trees satisfy it.


\subsection*{P0036-A001 | Equal level growth does not determine branching number}

\entrylabel{Argument or result}

A geodesic spanning tree has exactly the Cayley graph sphere sizes at each level. But growth can exceed branching number for a bounded-degree tree:\allowbreak{} take a single infinite spine and at its nth vertex attach a finite b-ary tree of height n. Its level sizes have exponential rate sqrt(b)>1,\allowbreak{} while its only infinite ray is the spine and arbitrarily deep singleton cutsets give branching number one.


\entrylabel{Second check}

At level m the largest attached contribution starts around n=\allowbreak{}m/\allowbreak{}2 and has order b\textasciicircum{}(m/\allowbreak{}2),\allowbreak{} up to constant factors and a geometric sum. Every attachment is finite,\allowbreak{} so no additional infinite ray appears.


\entrylabel{Remaining obstruction}

This is not a geodesic Cayley spanning-tree counterexample. It proves that matching sphere counts alone cannot establish the conjectured branching number.


Verification:\allowbreak{} author checked limited argument | Independent review:\allowbreak{} none | Full target solved:\allowbreak{} False


\clearpage

\section*{P0037  BK disjoint occurrence for weighted uniform spanning trees}

\entrylabel{Primary source}

Probability on Trees and Networks,\allowbreak{} paperback electronic version 19 August 2026;\allowbreak{} Conjecture 4.10,\allowbreak{} printed p.125


\url{https://rdlyons.pages.iu.edu/prbtree/book_pb.pdf}


Source statement date:\allowbreak{} not established. Archived access:\allowbreak{} 2026-09-06T15:\allowbreak{}30:\allowbreak{}43Z.


Historical:\allowbreak{} source\_\allowbreak{}says\_\allowbreak{}open | Current:\allowbreak{} unverified | Result:\allowbreak{} attempted\_\allowbreak{}unsolved


\entrylabel{Question}

For increasing events A,\allowbreak{}B on edge subsets,\allowbreak{} is the probability of their disjoint occurrence under any weighted uniform spanning-tree law at most P(A)P(B)?


\entrylabel{Scope}

Disjoint witnesses use disjoint edge-coordinate sets. Arbitrary events are excluded;\allowbreak{} negative association for fixed disjoint coordinates is weaker. Goldman 2023 proves a special class generated by simple points.


\subsection*{P0037-A001 | BK for fixed conjunctions and a nonincreasing-event obstruction}

\entrylabel{Argument or result}

For disjoint fixed edge sets F and H,\allowbreak{} the spanning-tree determinantal formula and Hadamard-Fischer inequality give P(F union H subset T)<=\allowbreak{}P(F subset T)P(H subset T). Thus fixed increasing conjunctions obey the desired inequality. The increasing condition matters:\allowbreak{} on a triangle,\allowbreak{} A=\allowbreak{}\{e in T\} and B=\allowbreak{}\{f not in T\},\allowbreak{} e!=\allowbreak{}f,\allowbreak{} have disjoint coordinate witnesses but P(A intersection B)=\allowbreak{}1/\allowbreak{}3>2/\allowbreak{}9=\allowbreak{}P(A)P(B).


\entrylabel{Second check}

There are three equiprobable spanning trees of a triangle,\allowbreak{} which checks the latter probabilities directly. The determinantal argument is for a positive semidefinite transfer-current kernel.


\entrylabel{Remaining obstruction}

General increasing disjoint occurrence is a union over adaptive witnesses. A fixed-coordinate negative-association inequality cannot simply be applied after that union.


Verification:\allowbreak{} author checked limited argument | Independent review:\allowbreak{} none | Full target solved:\allowbreak{} False


\clearpage

\section*{P0038  Transience preserved in Cayley percolation clusters}

\entrylabel{Primary source}

Probability on Trees and Networks,\allowbreak{} paperback electronic version 19 August 2026;\allowbreak{} Conjecture 6.44,\allowbreak{} printed p.216


\url{https://rdlyons.pages.iu.edu/prbtree/book_pb.pdf}


Source statement date:\allowbreak{} not established. Archived access:\allowbreak{} 2026-09-06T15:\allowbreak{}30:\allowbreak{}43Z.


Historical:\allowbreak{} source\_\allowbreak{}says\_\allowbreak{}open | Current:\allowbreak{} unverified | Result:\allowbreak{} attempted\_\allowbreak{}unsolved


\entrylabel{Question}

Is every infinite Bernoulli percolation cluster on a transient Cayley graph almost surely transient?


\entrylabel{Scope}

The claim concerns all infinite clusters simultaneously. Nonamenable and Euclidean cases have substantial existing results.


\subsection*{P0038-A001 | Percolation on regular trees}

\entrylabel{Argument or result}

For bond percolation on the (b+\allowbreak{}1)-regular tree,\allowbreak{} descendants of a nonroot vertex form a Galton-Watson process with Binomial(b,\allowbreak{}p) offspring. Conditional on survival when bp>1,\allowbreak{} its branching number is bp almost surely,\allowbreak{} so the electrical-network branching-number criterion gives transience.


\entrylabel{Second check}

The root has b+\allowbreak{}1 possible children but a finite root modification does not change transience. At bp<=\allowbreak{}1 the nondegenerate branching process dies out.


\entrylabel{Remaining obstruction}

The reduction uses independent descendant trees and does not extend to clusters with overlapping paths in a general transient Cayley graph.


Verification:\allowbreak{} author checked limited argument | Independent review:\allowbreak{} none | Full target solved:\allowbreak{} False


\clearpage

\section*{P0039  Positive-speed equivalence for Cayley percolation clusters}

\entrylabel{Primary source}

Probability on Trees and Networks,\allowbreak{} paperback electronic version 19 August 2026;\allowbreak{} Conjecture 6.45,\allowbreak{} printed p.216


\url{https://rdlyons.pages.iu.edu/prbtree/book_pb.pdf}


Source statement date:\allowbreak{} not established. Archived access:\allowbreak{} 2026-09-06T15:\allowbreak{}30:\allowbreak{}43Z.


Historical:\allowbreak{} source\_\allowbreak{}says\_\allowbreak{}open | Current:\allowbreak{} unverified | Result:\allowbreak{} attempted\_\allowbreak{}unsolved


\entrylabel{Question}

Does simple random walk on a Cayley graph have positive speed if and only if simple random walk on its infinite Bernoulli percolation clusters has positive speed almost surely?


\entrylabel{Scope}

Use graph distance in the cluster for the cluster speed as in the source context;\allowbreak{} compare with ambient speed carefully. Nonamenable cases are reported known.


\subsection*{P0039-A001 | Zero speed on every bounded-degree subexponential-growth subgraph}

\entrylabel{Argument or result}

Let H be a connected subgraph of a bounded-degree graph G with |B\_\allowbreak{}G(x,\allowbreak{}n)|=\allowbreak{}exp(o(n)). The Carne-Varopoulos bound gives P\_\allowbreak{}x(X\_\allowbreak{}n=\allowbreak{}y)<=\allowbreak{}2 sqrt(deg\_\allowbreak{}H(y)/\allowbreak{}deg\_\allowbreak{}H(x)) exp(-d\_\allowbreak{}H(x,\allowbreak{}y)\textasciicircum{}2/\allowbreak{}(2n)). Summing over reachable vertices yields P(d\_\allowbreak{}H(x,\allowbreak{}X\_\allowbreak{}n)>=\allowbreak{}epsilon n)<=\allowbreak{}C exp(o(n)-epsilon\textasciicircum{}2 n/\allowbreak{}2),\allowbreak{} which is summable. Thus the intrinsic speed on H is zero almost surely.


\entrylabel{Second check}

Reachability in H implies membership in B\_\allowbreak{}G(x,\allowbreak{}n),\allowbreak{} but distances in the exponential bound remain intrinsic to H. Bounded degrees control the prefactor. This is a consequence of the classical heat-kernel bound,\allowbreak{} not a novelty claim.


\entrylabel{Remaining obstruction}

It settles the zero-speed equivalence for subexponential-growth Cayley graphs and their clusters only. Exponential-growth cases need different information.


Verification:\allowbreak{} author checked limited argument | Independent review:\allowbreak{} none | Full target solved:\allowbreak{} False


\clearpage

\section*{P0040  Uniqueness threshold below one in one-ended transitive graphs}

\entrylabel{Primary source}

Probability on Trees and Networks,\allowbreak{} paperback electronic version 19 August 2026;\allowbreak{} Conjecture 7.27,\allowbreak{} printed p.258


\url{https://rdlyons.pages.iu.edu/prbtree/book_pb.pdf}


Source statement date:\allowbreak{} not established. Archived access:\allowbreak{} 2026-09-06T15:\allowbreak{}30:\allowbreak{}43Z.


Historical:\allowbreak{} source\_\allowbreak{}says\_\allowbreak{}open | Current:\allowbreak{} unverified | Result:\allowbreak{} attempted\_\allowbreak{}unsolved


\entrylabel{Question}

Does every locally finite quasi-transitive graph with exactly one end have Bernoulli percolation uniqueness threshold p\_\allowbreak{}u<1?


\entrylabel{Scope}

A one-ended finite-presentation class is known. The assertion p\_\allowbreak{}c<1 alone is weaker and now resolved in the transitive setting.


\subsection*{P0040-A001 | Amenable Cayley special case and a cutset bottleneck}

\entrylabel{Argument or result}

For an amenable Cayley graph with one end,\allowbreak{} uniqueness gives p\_\allowbreak{}u=\allowbreak{}p\_\allowbreak{}c. The known strict p\_\allowbreak{}c<1 theorem for transitive graphs not roughly isometric to Z therefore gives p\_\allowbreak{}u<1 in this special case. One-endedness excludes the two-ended roughly-Z case.


\entrylabel{Second check}

This use is restricted to Cayley graphs;\allowbreak{} no unverified extension of the strict-critical-threshold theorem to every quasi-transitive graph is presumed.


\entrylabel{Remaining obstruction}

The remaining nonamenable quasi-transitive setting needs bounds on separating cutsets or cycle structure. A bounded cycle-basis argument is unavailable without its hypotheses.


Verification:\allowbreak{} author checked limited argument | Independent review:\allowbreak{} none | Full target solved:\allowbreak{} False


\clearpage

\section*{P0041  Nonuniqueness phase for nonamenable quasi-transitive graphs}

\entrylabel{Primary source}

Probability on Trees and Networks,\allowbreak{} paperback electronic version 19 August 2026;\allowbreak{} Conjectures 7.31 and 7.41,\allowbreak{} printed pp.260 and 265


\url{https://rdlyons.pages.iu.edu/prbtree/book_pb.pdf}


Source statement date:\allowbreak{} not established. Archived access:\allowbreak{} 2026-09-06T15:\allowbreak{}30:\allowbreak{}43Z.


Historical:\allowbreak{} source\_\allowbreak{}says\_\allowbreak{}open | Current:\allowbreak{} unverified | Result:\allowbreak{} attempted\_\allowbreak{}unsolved


\entrylabel{Question}

Must p\_\allowbreak{}c(G)<p\_\allowbreak{}u(G) for every nonamenable quasi-transitive graph? The source's related formulation asks whether the property p\_\allowbreak{}c<p\_\allowbreak{}u is preserved by rough isometries among quasi-transitive graphs.


\entrylabel{Scope}

The closely linked rough-isometry conjecture is grouped here to avoid counting a reformulation separately. Nonunimodular and many other cases are established.


\subsection*{P0041-A001 | Strict gap on regular trees}

\entrylabel{Argument or result}

On a (b+\allowbreak{}1)-regular tree with b>=\allowbreak{}2,\allowbreak{} branching-process survival gives p\_\allowbreak{}c=\allowbreak{}1/\allowbreak{}b. For every p in (1/\allowbreak{}b,\allowbreak{}1),\allowbreak{} infinitely many separated descendant subtrees contain infinite open clusters while their attachment edge is closed,\allowbreak{} using independent positive-probability events. Thus there are infinitely many clusters and p\_\allowbreak{}u=\allowbreak{}1.


\entrylabel{Second check}

The attachment events can be placed on an infinite collection of disjoint descendant cones. For p<=\allowbreak{}1/\allowbreak{}b there is no infinite cluster.


\entrylabel{Remaining obstruction}

This proves the strict gap in a fundamental nonamenable family. Nonamenability by itself does not provide the same independent cone decomposition.


Verification:\allowbreak{} author checked limited argument | Independent review:\allowbreak{} none | Full target solved:\allowbreak{} False


\clearpage

\section*{P0042  Bootstrap thresholds as a characterization of amenable groups}

\entrylabel{Primary source}

Probability on Trees and Networks,\allowbreak{} paperback electronic version 19 August 2026;\allowbreak{} Question 7.42,\allowbreak{} printed p.266


\url{https://rdlyons.pages.iu.edu/prbtree/book_pb.pdf}


Source statement date:\allowbreak{} not established. Archived access:\allowbreak{} 2026-09-06T15:\allowbreak{}30:\allowbreak{}43Z.


Historical:\allowbreak{} source\_\allowbreak{}says\_\allowbreak{}open | Current:\allowbreak{} unverified | Result:\allowbreak{} attempted\_\allowbreak{}unsolved


\entrylabel{Question}

Is a finitely generated group amenable exactly when every Cayley graph from every finite generating set has bootstrap complete-occupation critical probability p(G,\allowbreak{}k) in \{0,\allowbreak{}1\} for every threshold k?


\entrylabel{Scope}

Initially each vertex is occupied independently with probability p;\allowbreak{} vacant vertices with at least k occupied neighbors become occupied. Criticality is for positive probability of eventual occupation of all vertices.


\subsection*{P0042-A001 | Bootstrap thresholds on the line and the three-regular tree}

\entrylabel{Argument or result}

On Z the one-neighbor bootstrap threshold is zero,\allowbreak{} whereas the two-neighbor threshold is one because any initially vacant adjacent pair stays vacant. On the three-regular tree with threshold two,\allowbreak{} the final vacant set is the two-core of the initial vacant percolation. If its density q<=\allowbreak{}1/\allowbreak{}2,\allowbreak{} all vacant clusters are finite almost surely and leaf pruning empties them. If q>1/\allowbreak{}2,\allowbreak{} with positive probability a vertex has two surviving vacant branches,\allowbreak{} giving a bi-infinite vacant path that remains vacant. Hence the complete-occupation threshold is p=\allowbreak{}1/\allowbreak{}2.


\entrylabel{Second check}

An infinite bi-infinite path supplies a stable set since each vertex has two vacant neighbors. Finite trees have an exposed leaf and are exhausted by the pruning process. Critical Binomial(2,\allowbreak{}1/\allowbreak{}2) clusters die out.


\entrylabel{Remaining obstruction}

This illustrates both the amenable endpoint behavior and a nonamenable interior threshold. It does not classify every generating set or every group.


Verification:\allowbreak{} author checked limited argument | Independent review:\allowbreak{} none | Full target solved:\allowbreak{} False


\clearpage

\section*{P0043  Absence of infinite clusters at critical percolation}

\entrylabel{Primary source}

Probability on Trees and Networks,\allowbreak{} paperback electronic version 19 August 2026;\allowbreak{} Conjecture 8.15,\allowbreak{} printed p.285


\url{https://rdlyons.pages.iu.edu/prbtree/book_pb.pdf}


Source statement date:\allowbreak{} not established. Archived access:\allowbreak{} 2026-09-06T15:\allowbreak{}30:\allowbreak{}43Z.


Historical:\allowbreak{} source\_\allowbreak{}says\_\allowbreak{}open | Current:\allowbreak{} unverified | Result:\allowbreak{} attempted\_\allowbreak{}unsolved


\entrylabel{Question}

If G is quasi-transitive and p\_\allowbreak{}c(G)<1,\allowbreak{} is there almost surely no infinite Bernoulli cluster at p\_\allowbreak{}c(G)?


\entrylabel{Scope}

Do not count each lattice dimension as a separate problem. Known nonamenable and high-dimensional cases do not settle the entire assertion.


\subsection*{P0043-A001 | Critical branching-process extinction}

\entrylabel{Argument or result}

For the regular tree at p=\allowbreak{}1/\allowbreak{}b,\allowbreak{} an infinite descendant cluster would have survival probability x satisfying x=\allowbreak{}1-(1-px)\textasciicircum{}b. For b>=\allowbreak{}2 strict concavity gives 1-(1-x/\allowbreak{}b)\textasciicircum{}b<x for every x>0,\allowbreak{} so only x=\allowbreak{}0 is possible.


\entrylabel{Second check}

The strict inequality can also be obtained from the binomial generating function lying strictly above its tangent at one. This excludes the deterministic b=\allowbreak{}1 exception,\allowbreak{} whose p\_\allowbreak{}c is one.


\entrylabel{Remaining obstruction}

Critical clusters in a general quasi-transitive graph are not independent branching trees. Their loop dependence is not captured by this fixed-point proof.


Verification:\allowbreak{} author checked limited argument | Independent review:\allowbreak{} none | Full target solved:\allowbreak{} False


\clearpage

\section*{P0044  Invariant monotone coupling of free and wired uniform spanning forests}

\entrylabel{Primary source}

Probability on Trees and Networks,\allowbreak{} paperback electronic version 19 August 2026;\allowbreak{} Question 10.6,\allowbreak{} printed pp.348-349


\url{https://rdlyons.pages.iu.edu/prbtree/book_pb.pdf}


Source statement date:\allowbreak{} not established. Archived access:\allowbreak{} 2026-09-06T15:\allowbreak{}30:\allowbreak{}43Z.


Historical:\allowbreak{} source\_\allowbreak{}says\_\allowbreak{}open | Current:\allowbreak{} unverified | Result:\allowbreak{} attempted\_\allowbreak{}unsolved


\entrylabel{Question}

For an infinite network,\allowbreak{} is there a coupling (W,\allowbreak{}F) of its wired and free uniform spanning forests with W subset F almost surely whose joint law is invariant under all network automorphisms?


\entrylabel{Scope}

Strassen gives an unconstrained monotone coupling. Sofic Cayley groups admit invariant couplings by Lyons-Thom 2016;\allowbreak{} arbitrary invariant percolations can fail this property.


\subsection*{P0044-A001 | Invariant monotone coupling under an amenable action}

\entrylabel{Argument or result}

Start with a monotone coupling of the two invariant forest laws under a countable amenable action. Average its translates over a Folner sequence. The product space of two edge configurations is compact,\allowbreak{} the prescribed marginals are unchanged,\allowbreak{} and the order relation is closed. A subsequential weak limit is therefore still a monotone coupling and is invariant because translation changes each average by at most its Folner boundary ratio.


\entrylabel{Second check}

The argument averages probability measures,\allowbreak{} not individual configurations,\allowbreak{} so convexity preserves the marginals and the monotone support. Compactness is available for a countable locally finite edge set.


\entrylabel{Remaining obstruction}

For a nonamenable acting group there is no Folner averaging sequence. Existence of an arbitrary monotone coupling does not by itself produce an invariant one.


Verification:\allowbreak{} author checked limited argument | Independent review:\allowbreak{} none | Full target solved:\allowbreak{} False


\clearpage

\section*{P0045  Rational expected free-forest degree for finitely presented groups}

\entrylabel{Primary source}

Probability on Trees and Networks,\allowbreak{} paperback electronic version 19 August 2026;\allowbreak{} Question 10.11,\allowbreak{} printed p.351


\url{https://rdlyons.pages.iu.edu/prbtree/book_pb.pdf}


Source statement date:\allowbreak{} not established. Archived access:\allowbreak{} 2026-09-06T15:\allowbreak{}30:\allowbreak{}43Z.


Historical:\allowbreak{} source\_\allowbreak{}says\_\allowbreak{}open | Current:\allowbreak{} unverified | Result:\allowbreak{} attempted\_\allowbreak{}unsolved


\entrylabel{Question}

For a Cayley graph of a finitely presented group,\allowbreak{} is the expected degree at the identity in its free uniform spanning forest rational?


\entrylabel{Scope}

This is the first l2-Betti-number question via E degree=\allowbreak{}2+\allowbreak{}2 beta\_\allowbreak{}1. Counterexamples to broader Atiyah statements must be checked for both first degree and finite presentation before exclusion.


\subsection*{P0045-A001 | Rational finite-volume values do not force a rational limit}

\entrylabel{Argument or result}

For the standard free-group Cayley tree the free spanning forest is the entire tree,\allowbreak{} so its expected degree is the integer 2r. Finite exhaustion spanning-tree expectations are rational under unit conductances,\allowbreak{} but rational approximants can converge to an irrational number,\allowbreak{} as floor(n sqrt(2))/\allowbreak{}n demonstrates.


\entrylabel{Second check}

A finite tree has only one spanning tree,\allowbreak{} so the free-group calculation is exact and requires no limiting interchange.


\entrylabel{Remaining obstruction}

The general finitely presented-group question needs the algebraic structure behind the first l2 Betti number;\allowbreak{} rationality of finite combinatorial counts is insufficient.


Verification:\allowbreak{} author checked limited argument | Independent review:\allowbreak{} none | Full target solved:\allowbreak{} False


\clearpage

\section*{P0046  Connecting a free uniform spanning forest by arbitrarily sparse invariant additions}

\entrylabel{Primary source}

Probability on Trees and Networks,\allowbreak{} paperback electronic version 19 August 2026;\allowbreak{} Question 10.12,\allowbreak{} printed p.351


\url{https://rdlyons.pages.iu.edu/prbtree/book_pb.pdf}


Source statement date:\allowbreak{} not established. Archived access:\allowbreak{} 2026-09-06T15:\allowbreak{}30:\allowbreak{}43Z.


Historical:\allowbreak{} source\_\allowbreak{}says\_\allowbreak{}open | Current:\allowbreak{} unverified | Result:\allowbreak{} attempted\_\allowbreak{}unsolved


\entrylabel{Question}

For any Cayley graph,\allowbreak{} its free uniform spanning forest F,\allowbreak{} and epsilon>0,\allowbreak{} is there an invariant connected random spanning subgraph omega containing F with P(e in omega minus F)<epsilon for every edge e?


\entrylabel{Scope}

Added edges may depend on F. Independent Bernoulli sprinkling is explicitly known to fail for some nonamenable graphs.


\subsection*{P0046-A001 | Connection to the cost equality obstruction}

\entrylabel{Argument or result}

If the free forest is already connected,\allowbreak{} taking its own configuration gives the requested extension with no added edges. Otherwise,\allowbreak{} an invariant connected extension with each edge added with probability at most epsilon would have expected root degree at most E deg\_\allowbreak{}F(o)+\allowbreak{}deg\_\allowbreak{}G(o) epsilon. The standard invariant-graphing bound would then imply cost(Gamma)<=\allowbreak{}E deg\_\allowbreak{}F(o)/\allowbreak{}2=\allowbreak{}1+\allowbreak{}beta\_\allowbreak{}1(Gamma) after epsilon tends to zero.


\entrylabel{Second check}

The root-degree estimate is the finite sum of per-edge addition probabilities. Combined with the known opposite cost inequality,\allowbreak{} the proposed construction would prove equality,\allowbreak{} not merely a weak approximation statement.


\entrylabel{Remaining obstruction}

No dependent sprinkling construction was found. This reduction explains why independent sprinkling or a generic connectivity argument does not settle the question.


Verification:\allowbreak{} author checked limited argument | Independent review:\allowbreak{} none | Full target solved:\allowbreak{} False


\clearpage

\section*{P0047  Deterministic component count of the free uniform spanning forest}

\entrylabel{Primary source}

Probability on Trees and Networks,\allowbreak{} paperback electronic version 19 August 2026;\allowbreak{} Question 10.28,\allowbreak{} printed p.359


\url{https://rdlyons.pages.iu.edu/prbtree/book_pb.pdf}


Source statement date:\allowbreak{} not established. Archived access:\allowbreak{} 2026-09-06T15:\allowbreak{}30:\allowbreak{}43Z.


Historical:\allowbreak{} source\_\allowbreak{}says\_\allowbreak{}open | Current:\allowbreak{} unverified | Result:\allowbreak{} attempted\_\allowbreak{}unsolved


\entrylabel{Question}

On every infinite connected locally finite network,\allowbreak{} is the number of free uniform spanning-forest components almost surely a constant?


\entrylabel{Scope}

Automorphisms with an infinite orbit give a positive answer. The remaining arbitrary-network setting lacks that ergodic shortcut.


\subsection*{P0047-A001 | Why tail triviality alone does not fix the component count}

\entrylabel{Argument or result}

An ergodic action under which the forest law is invariant makes its component count deterministic. On an arbitrary network there need not be such an action. Trying to substitute tail triviality fails at the logical step:\allowbreak{} adding one edge between two components can change the component count,\allowbreak{} so this statistic has not been shown to be measurable with respect to the edge tail sigma-field.


\entrylabel{Second check}

The finite edge change can preserve acyclicity when it joins two different tree components,\allowbreak{} so acyclicity does not remove this obstruction.


\entrylabel{Remaining obstruction}

A forest-law-specific argument making the count constant is still needed in the absence of a useful invariant ergodic action.


Verification:\allowbreak{} author checked limited argument | Independent review:\allowbreak{} none | Full target solved:\allowbreak{} False


\clearpage

\section*{P0048  One or infinitely many free-forest components without unimodularity}

\entrylabel{Primary source}

Probability on Trees and Networks,\allowbreak{} paperback electronic version 19 August 2026;\allowbreak{} Question 10.29,\allowbreak{} printed p.360


\url{https://rdlyons.pages.iu.edu/prbtree/book_pb.pdf}


Source statement date:\allowbreak{} not established. Archived access:\allowbreak{} 2026-09-06T15:\allowbreak{}30:\allowbreak{}43Z.


Historical:\allowbreak{} source\_\allowbreak{}says\_\allowbreak{}open | Current:\allowbreak{} unverified | Result:\allowbreak{} attempted\_\allowbreak{}unsolved


\entrylabel{Question}

On every transitive network,\allowbreak{} is the free uniform spanning forest almost surely connected or composed of infinitely many components?


\entrylabel{Scope}

The unimodular case was solved by Hutchcroft-Nachmias 2017 and Timar 2018. The unresolved scope retained here is nonunimodular.


\subsection*{P0048-A001 | Invariant ergodicity permits a finite component count larger than one}

\entrylabel{Argument or result}

On the Cayley graph of Z with generators +\allowbreak{}/\allowbreak{}-1 and +\allowbreak{}/\allowbreak{}-2,\allowbreak{} retain exactly the edges of length two. The resulting deterministic spanning forest has two bi-infinite components,\allowbreak{} is translation invariant,\allowbreak{} and its point-mass law is ergodic. Thus those two properties alone cannot exclude a finite count of two.


\entrylabel{Second check}

The two components are the even and odd integers;\allowbreak{} every retained vertex has degree two and neither component contains a cycle.


\entrylabel{Remaining obstruction}

This forest is not asserted to have a uniform-spanning-forest law. The target requires extra free-forest structure in the nonunimodular case.


Verification:\allowbreak{} author checked limited argument | Independent review:\allowbreak{} none | Full target solved:\allowbreak{} False


\clearpage

\section*{P0049  Infinitely many ends in all free-forest components without unimodularity}

\entrylabel{Primary source}

Probability on Trees and Networks,\allowbreak{} paperback electronic version 19 August 2026;\allowbreak{} Question 10.52,\allowbreak{} printed p.377


\url{https://rdlyons.pages.iu.edu/prbtree/book_pb.pdf}


Source statement date:\allowbreak{} not established. Archived access:\allowbreak{} 2026-09-06T15:\allowbreak{}30:\allowbreak{}43Z.


Historical:\allowbreak{} source\_\allowbreak{}says\_\allowbreak{}open | Current:\allowbreak{} unverified | Result:\allowbreak{} attempted\_\allowbreak{}unsolved


\entrylabel{Question}

On a transitive network with distinct free and wired uniform spanning-forest laws,\allowbreak{} must every free-forest component have infinitely many ends almost surely?


\entrylabel{Scope}

The source's subsequent indistinguishability theorem settles the unimodular case. Haar-light components alone do not automatically refute an ends statement.


\subsection*{P0049-A001 | Missing mass-transport step in the nonunimodular case}

\entrylabel{Argument or result}

In the unimodular setting,\allowbreak{} an expected-degree argument can establish the presence of infinitely ended components,\allowbreak{} and indistinguishability can promote a component property to all components. On a nonunimodular graph the unweighted root mass balance used in that argument is invalid;\allowbreak{} modular weights enter and light components need not receive the same balance.


\entrylabel{Second check}

An existence statement about some components would remain weaker than the source assertion about every component even if the first step survived.


\entrylabel{Remaining obstruction}

No replacement weighted transport and no counterexample with the required FUSF/\allowbreak{}WUSF inequality were obtained. This is a failed proof route,\allowbreak{} not a new theorem.


Verification:\allowbreak{} author checked limited argument | Independent review:\allowbreak{} none | Full target solved:\allowbreak{} False


\clearpage

\section*{P0050  Singularity of distinct free and wired forest measures}

\entrylabel{Primary source}

Probability on Trees and Networks,\allowbreak{} paperback electronic version 19 August 2026;\allowbreak{} Question 10.58,\allowbreak{} printed p.385


\url{https://rdlyons.pages.iu.edu/prbtree/book_pb.pdf}


Source statement date:\allowbreak{} not established. Archived access:\allowbreak{} 2026-09-06T15:\allowbreak{}30:\allowbreak{}43Z.


Historical:\allowbreak{} source\_\allowbreak{}says\_\allowbreak{}open | Current:\allowbreak{} unverified | Result:\allowbreak{} attempted\_\allowbreak{}unsolved


\entrylabel{Question}

If the wired and free uniform spanning-forest measures differ on an infinite network,\allowbreak{} must they be mutually singular?


\entrylabel{Scope}

Known for trees and networks with an infinite automorphism orbit. Distinct probability measures need not be singular in general.


\subsection*{P0050-A001 | Distinct invariant ergodic laws are mutually singular}

\entrylabel{Argument or result}

For a fixed countable group action,\allowbreak{} two distinct invariant ergodic probability measures are mutually singular. In the Lebesgue decomposition of one relative to the other,\allowbreak{} the absolutely continuous part is invariant. Its density is invariant and hence constant under the second ergodic law;\allowbreak{} ergodicity of the first forces its absolutely continuous mass to be either zero or one. Nonzero overlap therefore forces equality.


\entrylabel{Second check}

The singular and absolutely continuous parts are preserved by the action by uniqueness of Lebesgue decomposition;\allowbreak{} this is the load-bearing invariance step.


\entrylabel{Remaining obstruction}

This handles forest laws admitting such an ergodic common action. Arbitrary networks may lack it;\allowbreak{} two distinct probability laws need not be singular in general.


Verification:\allowbreak{} author checked limited argument | Independent review:\allowbreak{} none | Full target solved:\allowbreak{} False


\clearpage

\section*{P0051  Zero-one law for all forest components having one end}

\entrylabel{Primary source}

Probability on Trees and Networks,\allowbreak{} paperback electronic version 19 August 2026;\allowbreak{} Question 10.59,\allowbreak{} printed p.386


\url{https://rdlyons.pages.iu.edu/prbtree/book_pb.pdf}


Source statement date:\allowbreak{} not established. Archived access:\allowbreak{} 2026-09-06T15:\allowbreak{}30:\allowbreak{}43Z.


Historical:\allowbreak{} source\_\allowbreak{}says\_\allowbreak{}open | Current:\allowbreak{} unverified | Result:\allowbreak{} attempted\_\allowbreak{}unsolved


\entrylabel{Question}

For both free and wired uniform spanning forests of an arbitrary infinite network,\allowbreak{} is the event that every component has one end of probability zero or one?


\entrylabel{Scope}

Both boundary conditions are grouped. Changing finitely many forest edges can merge components,\allowbreak{} so tail measurability needs a proof.


\subsection*{P0051-A001 | The all-one-ended event is not automatically a tail event}

\entrylabel{Argument or result}

Take two disjoint one-ended infinite trees and add one edge joining them. Before the addition every component has one end;\allowbreak{} afterwards the connected tree has two ends,\allowbreak{} since deleting the joining edge separates two infinite pieces.


\entrylabel{Second check}

The added edge creates no cycle and the change is finite. Thus the property can change within the class of spanning forests under a finite edge modification.


\entrylabel{Remaining obstruction}

Uniform-forest tail triviality cannot be invoked until a separate tail representation of the event is proved. The example is a logical obstruction,\allowbreak{} not a uniform-forest counterexample.


Verification:\allowbreak{} author checked limited argument | Independent review:\allowbreak{} none | Full target solved:\allowbreak{} False


\clearpage

\section*{P0052  Limiting bushes along a wired uniform-forest ray}

\entrylabel{Primary source}

Probability on Trees and Networks,\allowbreak{} paperback electronic version 19 August 2026;\allowbreak{} Conjecture 10.60,\allowbreak{} printed p.386


\url{https://rdlyons.pages.iu.edu/prbtree/book_pb.pdf}


Source statement date:\allowbreak{} not established. Archived access:\allowbreak{} 2026-09-06T15:\allowbreak{}30:\allowbreak{}43Z.


Historical:\allowbreak{} source\_\allowbreak{}says\_\allowbreak{}open | Current:\allowbreak{} unverified | Result:\allowbreak{} attempted\_\allowbreak{}unsolved


\entrylabel{Question}

Along the unique ray from the identity in its wired uniform spanning-forest component,\allowbreak{} do the finite bushes at successive ray vertices,\allowbreak{} translated back to the identity,\allowbreak{} converge in distribution?


\entrylabel{Scope}

Restrict to Cayley graphs where the stated unique ray exists. The source records solutions for growth at least five and standard Z\textasciicircum{}2,\allowbreak{}Z\textasciicircum{}3,\allowbreak{}Z\textasciicircum{}4;\allowbreak{} the remaining low-growth Cayley cases must be checked individually before any promotion.


\subsection*{P0052-A001 | Stationarity along a chosen spine is an additional requirement}

\entrylabel{Argument or result}

The WUSF viewed from an ordinary invariant root and the WUSF viewed from the nth vertex on its distinguished spine are different rooted laws. Moving along the spine is a configuration-dependent selection and need not preserve the original root law. Compactness or tightness,\allowbreak{} even if proved for the bush sequence,\allowbreak{} would only give subsequential limits.


\entrylabel{Second check}

The elementary sequence of probability laws delta\_\allowbreak{}0,\allowbreak{}delta\_\allowbreak{}1,\allowbreak{}delta\_\allowbreak{}0,\allowbreak{}delta\_\allowbreak{}1 is tight but has no full limit,\allowbreak{} illustrating precisely the missing implication rather than modeling WUSF bushes.


\entrylabel{Remaining obstruction}

Need regeneration or asymptotic stationarity along the loop-erased spine. Known high-growth and standard-lattice cases in the source must be removed before claiming a new remaining case.


Verification:\allowbreak{} author checked limited argument | Independent review:\allowbreak{} none | Full target solved:\allowbreak{} False


\clearpage

\section*{P0053  One-ended wired minimal forests without unimodularity}

\entrylabel{Primary source}

Probability on Trees and Networks,\allowbreak{} paperback electronic version 19 August 2026;\allowbreak{} Question 11.13,\allowbreak{} printed p.401


\url{https://rdlyons.pages.iu.edu/prbtree/book_pb.pdf}


Source statement date:\allowbreak{} not established. Archived access:\allowbreak{} 2026-09-06T15:\allowbreak{}30:\allowbreak{}43Z.


Historical:\allowbreak{} source\_\allowbreak{}says\_\allowbreak{}open | Current:\allowbreak{} unverified | Result:\allowbreak{} attempted\_\allowbreak{}unsolved


\entrylabel{Question}

Does each wired minimal spanning-forest component of a nonunimodular transitive graph have one end almost surely?


\entrylabel{Scope}

The accompanying request for useful nonhomogeneous sufficient conditions is retained as a related extension,\allowbreak{} not another count.


\subsection*{P0053-A001 | A failed largest-edge argument on an infinite path}

\entrylabel{Argument or result}

A finite-cycle minimum-spanning-tree proof deletes the largest edge of a cycle. A proposed extension to a bi-infinite path fails because its supremum need not be attained:\allowbreak{} labels 1-1/\allowbreak{}(|k|+\allowbreak{}2) along k in Z approach one and have no maximum.


\entrylabel{Second check}

The deterministic sequence only disproves the logical inference that every infinite path has a largest label. It is not an iid weighted-graph realization or a counterexample to the wired forest question.


\entrylabel{Remaining obstruction}

One-endedness needs an invasion-percolation or mass-transport argument that treats unattained suprema and the nonunimodular setting.


Verification:\allowbreak{} author checked limited argument | Independent review:\allowbreak{} none | Full target solved:\allowbreak{} False


\clearpage

\section*{P0054  Infinitely many ends in free minimal forests distinct from wired}

\entrylabel{Primary source}

Probability on Trees and Networks,\allowbreak{} paperback electronic version 19 August 2026;\allowbreak{} Conjecture 11.16,\allowbreak{} printed p.402


\url{https://rdlyons.pages.iu.edu/prbtree/book_pb.pdf}


Source statement date:\allowbreak{} not established. Archived access:\allowbreak{} 2026-09-06T15:\allowbreak{}30:\allowbreak{}43Z.


Historical:\allowbreak{} source\_\allowbreak{}says\_\allowbreak{}open | Current:\allowbreak{} unverified | Result:\allowbreak{} attempted\_\allowbreak{}unsolved


\entrylabel{Question}

If a quasi-transitive graph has distinct wired and free minimal spanning-forest laws,\allowbreak{} does every free minimal-forest component have infinitely many ends almost surely?


\entrylabel{Scope}

The source and later Timar results establish unimodular cases;\allowbreak{} retain the genuinely remaining scope. All edges have independent continuous labels.


\subsection*{P0054-A001 | Regular-tree special case and resampling obstruction}

\entrylabel{Argument or result}

On a regular tree every edge is a bridge and there are no finite cycles,\allowbreak{} so the free minimal spanning forest is the whole tree. It has infinitely many ends,\allowbreak{} providing a direct special case where free and wired forests differ. On a graph with cycles,\allowbreak{} inserting an edge can force deletion of another edge by the cycle-maximum rule.


\entrylabel{Second check}

The free forest statement follows directly from its finite-cycle definition;\allowbreak{} it is not the wired definition,\allowbreak{} which also sees connections at infinity.


\entrylabel{Remaining obstruction}

A minimal forest is not a Bernoulli insertion-tolerant process. Arguments proving end properties for percolation do not transfer without controlling the forced deletions.


Verification:\allowbreak{} author checked limited argument | Independent review:\allowbreak{} none | Full target solved:\allowbreak{} False


\clearpage

\section*{P0055  Minimal-forest connectivity and uniqueness at p\_\allowbreak{}u}

\entrylabel{Primary source}

Probability on Trees and Networks,\allowbreak{} paperback electronic version 19 August 2026;\allowbreak{} Conjecture 11.18,\allowbreak{} printed p.403


\url{https://rdlyons.pages.iu.edu/prbtree/book_pb.pdf}


Source statement date:\allowbreak{} not established. Archived access:\allowbreak{} 2026-09-06T15:\allowbreak{}30:\allowbreak{}43Z.


Historical:\allowbreak{} source\_\allowbreak{}says\_\allowbreak{}open | Current:\allowbreak{} unverified | Result:\allowbreak{} attempted\_\allowbreak{}unsolved


\entrylabel{Question}

On a nonamenable quasi-transitive graph with p\_\allowbreak{}u<1,\allowbreak{} is the free minimal spanning forest connected almost surely exactly when Bernoulli percolation has a unique infinite cluster at p\_\allowbreak{}u?


\entrylabel{Scope}

The source's stronger bijection between forest components and p\_\allowbreak{}u clusters is grouped as a strengthening. One direction is already proved there.


\subsection*{P0055-A001 | The easy direction of the uniqueness connection}

\entrylabel{Argument or result}

Under the p\_\allowbreak{}u>p\_\allowbreak{}c hypothesis in the source proposition,\allowbreak{} each FMSF component meets exactly one infinite p\_\allowbreak{}u cluster. Consequently a connected FMSF can meet only one such cluster,\allowbreak{} so connectivity implies uniqueness at p\_\allowbreak{}u in that setting.


\entrylabel{Second check}

The condition p\_\allowbreak{}u>p\_\allowbreak{}c is retained;\allowbreak{} assuming it for every nonamenable graph would silently solve another admitted conjecture. The proposition does not say that one p\_\allowbreak{}u cluster meets only one forest component.


\entrylabel{Remaining obstruction}

The reverse direction needs to exclude several free-forest components meeting the same infinite p\_\allowbreak{}u cluster. The stronger component-bijection formulation remains unproved here.


Verification:\allowbreak{} author checked limited argument | Independent review:\allowbreak{} none | Full target solved:\allowbreak{} False


\clearpage

\section*{P0056  One or infinitely many minimal-forest components}

\entrylabel{Primary source}

Probability on Trees and Networks,\allowbreak{} paperback electronic version 19 August 2026;\allowbreak{} Question 11.19,\allowbreak{} printed p.403


\url{https://rdlyons.pages.iu.edu/prbtree/book_pb.pdf}


Source statement date:\allowbreak{} not established. Archived access:\allowbreak{} 2026-09-06T15:\allowbreak{}30:\allowbreak{}43Z.


Historical:\allowbreak{} source\_\allowbreak{}says\_\allowbreak{}open | Current:\allowbreak{} unverified | Result:\allowbreak{} attempted\_\allowbreak{}unsolved


\entrylabel{Question}

Must both free and wired minimal spanning forests on a quasi-transitive graph have either one or infinitely many components almost surely?


\entrylabel{Scope}

Free and wired are grouped. Distinct-law unimodular cases are reported resolved by Timar 2018.


\subsection*{P0056-A001 | Ergodic factor laws give constancy but not the value of the count}

\entrylabel{Argument or result}

Minimal spanning forests are equivariant functions of iid continuous edge labels,\allowbreak{} so an ergodic label action makes their component count constant. That alone permits a finite constant larger than one:\allowbreak{} the invariant deterministic two-component forest on Z with generators +\allowbreak{}/\allowbreak{}-1,\allowbreak{}+\allowbreak{}/\allowbreak{}-2 from P0048 is a counterexample to the purely ergodic inference.


\entrylabel{Second check}

The example is explicitly not a minimal-spanning-forest law. It separates the zero-one statement for invariant events from a classification of possible constant values.


\entrylabel{Remaining obstruction}

The usual insertion-tolerance route must account for cycle-breaking deletions. No argument excluding a finite count between two and infinity was obtained in the full quasi-transitive scope.


Verification:\allowbreak{} author checked limited argument | Independent review:\allowbreak{} none | Full target solved:\allowbreak{} False


\clearpage

\section*{P0057  Uniqueness-critical percolation versus invariant spanning trees}

\entrylabel{Primary source}

Probability on Trees and Networks,\allowbreak{} paperback electronic version 19 August 2026;\allowbreak{} Question 11.31,\allowbreak{} printed p.408


\url{https://rdlyons.pages.iu.edu/prbtree/book_pb.pdf}


Source statement date:\allowbreak{} not established. Archived access:\allowbreak{} 2026-09-06T15:\allowbreak{}30:\allowbreak{}43Z.


Historical:\allowbreak{} source\_\allowbreak{}says\_\allowbreak{}open | Current:\allowbreak{} unverified | Result:\allowbreak{} attempted\_\allowbreak{}unsolved


\entrylabel{Question}

For a Cayley graph with p\_\allowbreak{}u<1,\allowbreak{} is there an implication in either direction between multiple infinite clusters at p\_\allowbreak{}u and the nonexistence of an invariant random spanning tree?


\entrylabel{Scope}

Each condition implies the free minimal spanning forest is disconnected,\allowbreak{} but a common consequence does not establish either implication.


\subsection*{P0057-A001 | Conditional implication through the component-bijection conjecture}

\entrylabel{Argument or result}

If there is no invariant random spanning tree,\allowbreak{} the invariant FMSF cannot be connected,\allowbreak{} because a connected spanning forest would itself be such a tree. If the stronger conjecture in P0055 gave a bijection between free-forest components and infinite p\_\allowbreak{}u clusters,\allowbreak{} that disconnection would imply multiple clusters at p\_\allowbreak{}u.


\entrylabel{Second check}

The first implication uses only the definition of an invariant spanning tree. The second is explicitly conditional on a separate unproved component correspondence.


\entrylabel{Remaining obstruction}

Two properties sharing FMSF disconnection as a consequence do not thereby imply one another. Neither unconditional implication in the source question is established here.


Verification:\allowbreak{} author checked limited argument | Independent review:\allowbreak{} none | Full target solved:\allowbreak{} False


\clearpage

\section*{P0058  Indistinguishability of minimal spanning-forest components}

\entrylabel{Primary source}

Probability on Trees and Networks,\allowbreak{} paperback electronic version 19 August 2026;\allowbreak{} Conjecture 11.33,\allowbreak{} printed p.409


\url{https://rdlyons.pages.iu.edu/prbtree/book_pb.pdf}


Source statement date:\allowbreak{} not established. Archived access:\allowbreak{} 2026-09-06T15:\allowbreak{}30:\allowbreak{}43Z.


Historical:\allowbreak{} source\_\allowbreak{}says\_\allowbreak{}open | Current:\allowbreak{} unverified | Result:\allowbreak{} attempted\_\allowbreak{}unsolved


\entrylabel{Question}

Are the components of both free and wired minimal spanning forests on a unimodular transitive graph indistinguishable by every invariant Borel component property?


\entrylabel{Scope}

Timar 2018 settles the free case when free and wired differ. The remaining equal-law or wired scope is the target.


\subsection*{P0058-A001 | Invariant ergodicity does not imply component indistinguishability}

\entrylabel{Argument or result}

Construct a periodic forest on Z\textasciicircum{}2 using blocks of four horizontal rows. The first row is a bi-infinite horizontal path. The next three rows form a comb:\allowbreak{} the middle row is a horizontal spine,\allowbreak{} with each vertex in the two outer rows attached vertically to that spine. Repeat and choose a uniformly random vertical phase. The law is translation invariant and ergodic on its finite phase orbit,\allowbreak{} yet its components are distinguishable as lines and combs by the existence of degree-four vertices.


\entrylabel{Second check}

Every vertex belongs to exactly one component and there are no cycles. If full square-lattice automorphism invariance is desired,\allowbreak{} also choose a uniform horizontal/\allowbreak{}vertical orientation;\allowbreak{} the full automorphism group is transitive on the resulting finite orbit.


\entrylabel{Remaining obstruction}

This is not a minimal-spanning-forest construction. It disproves only the shortcut from invariance and ergodicity to indistinguishability.


Verification:\allowbreak{} author checked limited argument | Independent review:\allowbreak{} none | Full target solved:\allowbreak{} False


\clearpage

\section*{P0059  Limiting bushes along a wired minimal-forest ray}

\entrylabel{Primary source}

Probability on Trees and Networks,\allowbreak{} paperback electronic version 19 August 2026;\allowbreak{} Conjecture 11.34,\allowbreak{} printed p.409


\url{https://rdlyons.pages.iu.edu/prbtree/book_pb.pdf}


Source statement date:\allowbreak{} not established. Archived access:\allowbreak{} 2026-09-06T15:\allowbreak{}30:\allowbreak{}43Z.


Historical:\allowbreak{} source\_\allowbreak{}says\_\allowbreak{}open | Current:\allowbreak{} unverified | Result:\allowbreak{} attempted\_\allowbreak{}unsolved


\entrylabel{Question}

For a Cayley graph not roughly isometric to Z,\allowbreak{} does the sequence of translated bushes along a ray in the identity's wired minimal spanning-forest component converge in distribution?


\entrylabel{Scope}

Uniqueness of that ray is itself tied to critical percolation. Keep the conditional structural assumptions explicit rather than selecting an unmentioned canonical ray.


\subsection*{P0059-A001 | Unique-ray and limiting-law issues must be separated}

\entrylabel{Argument or result}

A bush decomposition first needs an almost surely unique infinite ray in the relevant component. Assuming that ray before proving the applicable one-endedness result would import another open assertion. Even when uniqueness is available,\allowbreak{} rooting successively along the ray is a configuration-dependent operation and does not automatically preserve the original rooted law.


\entrylabel{Second check}

A convergent subsequence of bush laws would not establish a full limit. The period-two sequence example in P0052 shows the logical distinction.


\entrylabel{Remaining obstruction}

No regeneration theorem or coupling proving convergence along the wired minimal-forest ray was obtained. The preliminary one-endedness assumptions must be checked for each claimed case.


Verification:\allowbreak{} author checked limited argument | Independent review:\allowbreak{} none | Full target solved:\allowbreak{} False


\clearpage

\section*{P0060  Dimensions with a connected Euclidean minimal spanning forest}

\entrylabel{Primary source}

Probability on Trees and Networks,\allowbreak{} paperback electronic version 19 August 2026;\allowbreak{} Question 11.35,\allowbreak{} printed p.409


\url{https://rdlyons.pages.iu.edu/prbtree/book_pb.pdf}


Source statement date:\allowbreak{} not established. Archived access:\allowbreak{} 2026-09-06T15:\allowbreak{}30:\allowbreak{}43Z.


Historical:\allowbreak{} source\_\allowbreak{}says\_\allowbreak{}open | Current:\allowbreak{} unverified | Result:\allowbreak{} attempted\_\allowbreak{}unsolved


\entrylabel{Question}

For which dimensions d is the minimal spanning forest on Z\textasciicircum{}d connected almost surely?


\entrylabel{Scope}

Competing predictions put a transition near dimension six or eight. These dimension variants form one classification problem.


\subsection*{P0060-A001 | Exact one-dimensional connectivity}

\entrylabel{Argument or result}

On Z every edge is a bridge,\allowbreak{} so the free minimal spanning forest retains all edges. For an edge with label u<1 to be excluded by a wired connection at infinity,\allowbreak{} both infinite half-lines would have to supply paths of edges with labels below u. Each such event has probability zero for iid uniform labels,\allowbreak{} since the probability for its first N edges is u\textasciicircum{}N. Thus the wired forest also retains every edge almost surely.


\entrylabel{Second check}

There are only countably many edges,\allowbreak{} allowing the probability-one statement to hold simultaneously for all of them. The argument uses the unique paths on Z.


\entrylabel{Remaining obstruction}

The higher-dimensional connectivity classification depends on intersections of invasion structures and is not determined by the one-dimensional calculation.


Verification:\allowbreak{} author checked limited argument | Independent review:\allowbreak{} none | Full target solved:\allowbreak{} False


\clearpage

\section*{P0061  Full scaling limit of the square-lattice minimal spanning tree}

\entrylabel{Primary source}

Probability on Trees and Networks,\allowbreak{} paperback electronic version 19 August 2026;\allowbreak{} Question 11.36,\allowbreak{} printed p.410


\url{https://rdlyons.pages.iu.edu/prbtree/book_pb.pdf}


Source statement date:\allowbreak{} not established. Archived access:\allowbreak{} 2026-09-06T15:\allowbreak{}30:\allowbreak{}43Z.


Historical:\allowbreak{} source\_\allowbreak{}says\_\allowbreak{}open | Current:\allowbreak{} unverified | Result:\allowbreak{} attempted\_\allowbreak{}unsolved


\entrylabel{Question}

Does the minimal spanning tree on epsilon Z\textasciicircum{}2 have a full distributional scaling limit as epsilon tends to zero,\allowbreak{} in the embedded-tree topology used for its known subsequential limits?


\entrylabel{Scope}

The source leaves the exact topology to the cited Aizenman-Burchard-Newman-Wilson framework. A triangular-lattice theorem does not settle square-lattice universality.


\subsection*{P0061-A001 | Exact universality under a monotone change of edge labels}

\entrylabel{Argument or result}

On every finite graph with distinct edge weights,\allowbreak{} a strictly increasing transform preserves their order and hence preserves Kruskal selection and the minimum spanning tree. Independently,\allowbreak{} the cycle characterization retains exactly the same nonmaximal edges under that transform. For iid continuous distributions this identifies the law through ranks,\allowbreak{} or through the probability integral transform almost surely.


\entrylabel{Second check}

For a continuous CDF that is not strictly increasing everywhere,\allowbreak{} flat intervals have zero probability of a sampled weight,\allowbreak{} so the rank argument still applies almost surely.


\entrylabel{Remaining obstruction}

This is label-distribution universality on a fixed graph. It does not compare triangular and square lattice geometries or prove tightness and uniqueness of a continuum scaling limit.


Verification:\allowbreak{} author checked limited argument | Independent review:\allowbreak{} none | Full target solved:\allowbreak{} False


\clearpage

\section*{P0062  Connectivity of minimal forests on graphs roughly isometric to trees}

\entrylabel{Primary source}

Probability on Trees and Networks,\allowbreak{} paperback electronic version 19 August 2026;\allowbreak{} Question 11.37,\allowbreak{} printed p.410


\url{https://rdlyons.pages.iu.edu/prbtree/book_pb.pdf}


Source statement date:\allowbreak{} not established. Archived access:\allowbreak{} 2026-09-06T15:\allowbreak{}30:\allowbreak{}43Z.


Historical:\allowbreak{} source\_\allowbreak{}says\_\allowbreak{}open | Current:\allowbreak{} unverified | Result:\allowbreak{} attempted\_\allowbreak{}unsolved


\entrylabel{Question}

If a connected locally finite graph is roughly isometric to a tree,\allowbreak{} must its free minimal spanning forest be connected almost surely?


\entrylabel{Scope}

This concerns minimal forests,\allowbreak{} not uniform forests;\allowbreak{} Pete-Timar's uniform-forest counterexample cannot be substituted.


\subsection*{P0062-A001 | Connected FMSF on trees of finite cycle blocks}

\entrylabel{Argument or result}

Suppose an infinite cactus graph is built from bridges and edge-disjoint finite cycles. Its free minimal spanning forest retains every bridge and removes exactly the unique largest-weight edge of each cycle. Each resulting cycle block remains connected. A finite block path between any two vertices remains connected after these deletions,\allowbreak{} so the forest is globally connected and acyclic.


\entrylabel{Second check}

A simple cycle in a cactus is one of its cycle blocks;\allowbreak{} there are no additional cycle constraints coupling two blocks. Continuous iid labels make every finite maximum unique almost surely.


\entrylabel{Remaining obstruction}

Uniformly bounded block diameters give a class roughly isometric to trees. General roughly-tree graphs can have overlapping cycles and need not admit this block decomposition.


Verification:\allowbreak{} author checked limited argument | Independent review:\allowbreak{} none | Full target solved:\allowbreak{} False


\clearpage

\section*{P0063  Singularity of binary and ternary harmonic codings of the Cantor set}

\entrylabel{Primary source}

Probability on Trees and Networks,\allowbreak{} paperback electronic version 19 August 2026;\allowbreak{} Question 16.7,\allowbreak{} printed p.542


\url{https://rdlyons.pages.iu.edu/prbtree/book_pb.pdf}


Source statement date:\allowbreak{} not established. Archived access:\allowbreak{} 2026-09-06T15:\allowbreak{}30:\allowbreak{}43Z.


Historical:\allowbreak{} source\_\allowbreak{}says\_\allowbreak{}open | Current:\allowbreak{} unverified | Result:\allowbreak{} attempted\_\allowbreak{}unsolved


\entrylabel{Question}

For the middle-thirds Cantor set E,\allowbreak{} let nu\_\allowbreak{}b be harmonic measure on the source's base-b coding tree T[b](E),\allowbreak{} pushed to E. Is nu\_\allowbreak{}2 mutually singular with the Cantor-Lebesgue measure nu\_\allowbreak{}3?


\entrylabel{Scope}

Use the coding-tree construction and boundary map in Section 16.2;\allowbreak{} distinct harmonic measures or dimensions alone do not prove singularity.


\subsection*{P0063-A001 | Likelihood-ratio strategy and its exact missing input}

\entrylabel{Argument or result}

For product Bernoulli measures with different parameters,\allowbreak{} the strong law gives disjoint full-measure sets of limiting digit frequencies and hence mutual singularity. A possible analogous strategy for the two Cantor measures is a likelihood-ratio martingale along a common binary coding,\allowbreak{} with cumulative conditional relative entropy forcing its limit to vanish.


\entrylabel{Second check}

Merely proving that two measures are unequal is insufficient:\allowbreak{} two different positive bounded densities relative to the same measure are still equivalent. The binary digits of Cantor-Lebesgue measure are not assumed iid.


\entrylabel{Remaining obstruction}

No divergent conditional-entropy estimate or martingale tail control was established for the actual harmonic measure versus base-three Cantor measure.


Verification:\allowbreak{} author checked limited argument | Independent review:\allowbreak{} none | Full target solved:\allowbreak{} False


\clearpage

\section*{P0064  Dimension drop for every nonuniform Galton-Watson flow}

\entrylabel{Primary source}

Probability on Trees and Networks,\allowbreak{} paperback electronic version 19 August 2026;\allowbreak{} Conjecture 17.23,\allowbreak{} printed p.580


\url{https://rdlyons.pages.iu.edu/prbtree/book_pb.pdf}


Source statement date:\allowbreak{} not established. Archived access:\allowbreak{} 2026-09-06T15:\allowbreak{}30:\allowbreak{}43Z.


Historical:\allowbreak{} source\_\allowbreak{}says\_\allowbreak{}open | Current:\allowbreak{} unverified | Result:\allowbreak{} attempted\_\allowbreak{}unsolved


\entrylabel{Question}

For a nondegenerate supercritical Galton-Watson law of mean m under the chapter's moment assumptions,\allowbreak{} must every consistent equivariant flow rule different almost surely from UNIF have boundary-measure Hausdorff dimension below log m almost surely?


\entrylabel{Scope}

The source proves a stationary-measure sufficient condition. Difference on a GW-null set is not an admissible counterexample.


\subsection*{P0064-A001 | Exact dimension drop for the nonergodic same-color flow candidate}

\entrylabel{Argument or result}

For the AND/\allowbreak{}OR color construction of P0033,\allowbreak{} let K be the number of children of the parent color and choose uniformly among them. Conditional on root color zero,\allowbreak{} P(K=\allowbreak{}1,\allowbreak{}2,\allowbreak{}3)=\allowbreak{}(3/\allowbreak{}5,\allowbreak{}3/\allowbreak{}10,\allowbreak{}1/\allowbreak{}10);\allowbreak{} conditional on color one it is (3/\allowbreak{}10,\allowbreak{}3/\allowbreak{}5,\allowbreak{}1/\allowbreak{}10). The full conditional-child selection lemma implies these K variables are iid along the selected ray within each color class. Hence the local dimension is h\_\allowbreak{}0=\allowbreak{}(3/\allowbreak{}10)log 2+\allowbreak{}(1/\allowbreak{}10)log 3 or h\_\allowbreak{}1=\allowbreak{}(3/\allowbreak{}5)log 2+\allowbreak{}(1/\allowbreak{}10)log 3 by the strong law. Both are strictly less than log(12/\allowbreak{}5),\allowbreak{} the Galton-Watson dimension.


\entrylabel{Second check}

For color zero the unconditioned masses are .30 for two mixed children (K=\allowbreak{}1),\allowbreak{} .15 for two zero children (K=\allowbreak{}2),\allowbreak{} and .05 for three zero children (K=\allowbreak{}3);\allowbreak{} divide by color probability .5. For color one the corresponding masses are .15,\allowbreak{}.30,\allowbreak{}.05. Cylinder mass is exactly the product of reciprocal K values;\allowbreak{} Fubini transfers the ray statement to almost every tree and almost every flow ray.


\entrylabel{Remaining obstruction}

This author-checked special construction supports dimension drop and is not a counterexample to the full question. It inherits the P0033 coloring proof and remains without independent review.


Verification:\allowbreak{} author checked limited argument | Independent review:\allowbreak{} none | Full target solved:\allowbreak{} False


\clearpage

\section*{P0065  Monotonicity of harmonic dimension under inward bias}

\entrylabel{Primary source}

Probability on Trees and Networks,\allowbreak{} paperback electronic version 19 August 2026;\allowbreak{} Question 17.29,\allowbreak{} printed p.582


\url{https://rdlyons.pages.iu.edu/prbtree/book_pb.pdf}


Source statement date:\allowbreak{} not established. Archived access:\allowbreak{} 2026-09-06T15:\allowbreak{}30:\allowbreak{}43Z.


Historical:\allowbreak{} source\_\allowbreak{}says\_\allowbreak{}open | Current:\allowbreak{} unverified | Result:\allowbreak{} attempted\_\allowbreak{}unsolved


\entrylabel{Question}

For a nondegenerate supercritical Galton-Watson tree without leaves and 0<=\allowbreak{}lambda<m,\allowbreak{} does harmonic-measure dimension of RW\_\allowbreak{}lambda increase with lambda,\allowbreak{} and is the increase strict?


\entrylabel{Scope}

On a deterministic regular tree the dimension is constant,\allowbreak{} so retain the chapter's nondegeneracy assumptions for strictness. Both monotonic and strict clauses are grouped.


\subsection*{P0065-A001 | Regular-tree equality and a local-entropy countercheck}

\entrylabel{Argument or result}

On a deterministic b-ary tree,\allowbreak{} symmetry makes harmonic measure uniform on children for every transient bias lambda<b,\allowbreak{} so its dimension is log b,\allowbreak{} independent of lambda. Thus strict monotonicity must exclude deterministic offspring. A local root-entropy proof also fails:\allowbreak{} attach regular trees of branching 2 and 4 as two root subtrees. Their escape probabilities are 1-lambda/\allowbreak{}2 and 1-lambda/\allowbreak{}4 for lambda<2;\allowbreak{} last-exit child probabilities are proportional to these values and become more unequal as lambda increases toward 2,\allowbreak{} so their entropy decreases.


\entrylabel{Second check}

The two-child example is not an iid Galton-Watson tree. Its role is to disprove a pointwise local entropy monotonicity claim,\allowbreak{} not the global dimension conjecture.


\entrylabel{Remaining obstruction}

Global dimension averages local entropy under a bias-dependent law of the environment seen along the ray. Controlling that changing law remains necessary.


Verification:\allowbreak{} author checked limited argument | Independent review:\allowbreak{} none | Full target solved:\allowbreak{} False


\clearpage

\section*{P0066  Monotonicity of biased-walk speed on Galton-Watson trees}

\entrylabel{Primary source}

Probability on Trees and Networks,\allowbreak{} paperback electronic version 19 August 2026;\allowbreak{} Question 17.33,\allowbreak{} printed p.592


\url{https://rdlyons.pages.iu.edu/prbtree/book_pb.pdf}


Source statement date:\allowbreak{} not established. Archived access:\allowbreak{} 2026-09-06T15:\allowbreak{}30:\allowbreak{}43Z.


Historical:\allowbreak{} source\_\allowbreak{}says\_\allowbreak{}open | Current:\allowbreak{} unverified | Result:\allowbreak{} attempted\_\allowbreak{}unsolved


\entrylabel{Question}

For Galton-Watson trees with p\_\allowbreak{}0=\allowbreak{}0,\allowbreak{} is the speed of RW\_\allowbreak{}lambda nonincreasing as the bias parameter lambda increases?


\entrylabel{Scope}

The transient range is 0<lambda<m;\allowbreak{} beyond criticality the speed is zero. Positive-regeneration and large-bias partial results are not the complete monotonicity claim.


\subsection*{P0066-A001 | The changing stationary environment blocks a pointwise drift proof}

\entrylabel{Argument or result}

For fixed offspring D,\allowbreak{} the radial drift (D-lambda)/\allowbreak{}(D+\allowbreak{}lambda) decreases with lambda,\allowbreak{} and on a regular tree it gives the exact speed. In a nondegenerate Galton-Watson tree the stationary distribution of the environment seen by the walker also changes with lambda,\allowbreak{} so differentiating the speed has an additional measure-change term.


\entrylabel{Second check}

An elementary comparison shows the logical gap:\allowbreak{} drift at lambda=\allowbreak{}1,\allowbreak{}D=\allowbreak{}2 is 1/\allowbreak{}3,\allowbreak{} while at lambda=\allowbreak{}2,\allowbreak{}D=\allowbreak{}100 it is 98/\allowbreak{}102. Pointwise monotonicity in lambda cannot compare averages taken under different offspring laws.


\entrylabel{Remaining obstruction}

The example is not a Galton-Watson speed counterexample. No sign bound for the measure-change term,\allowbreak{} or alternative monotone coupling of speeds,\allowbreak{} was proved.


Verification:\allowbreak{} author checked limited argument | Independent review:\allowbreak{} none | Full target solved:\allowbreak{} False


\clearpage

\section*{P0067  Global Kuramoto synchronization of random cubic graphs}

\entrylabel{Primary source}

Randomstrasse101:\allowbreak{} Open Problems of 2024;\allowbreak{} Conjecture3,\allowbreak{} PDF p.3


\url{https://arxiv.org/pdf/2504.20539v1}


Source statement date:\allowbreak{} not established. Archived access:\allowbreak{} 2026-09-06T16:\allowbreak{}42:\allowbreak{}33Z.


Historical:\allowbreak{} source\_\allowbreak{}says\_\allowbreak{}open | Current:\allowbreak{} unverified | Result:\allowbreak{} attempted\_\allowbreak{}unsolved


\entrylabel{Question}

Is a uniformly random simple 3-regular graph globally synchronizing with probability tending to one? Global synchronization means that every local minimum of sum\_\allowbreak{}(ij in E)(1-cos(theta\_\allowbreak{}i-theta\_\allowbreak{}j)) is a consensus configuration.


\entrylabel{Scope}

Use even n. Connectivity is necessary but not sufficient;\allowbreak{} degree-three deterministic examples do not refute an asymptotic random-graph assertion.


\subsection*{P0067-A001 | Global Kuramoto synchronization of random cubic graphs special-case or structural calculation}

\entrylabel{Argument or result}

On the cubic prism C\_\allowbreak{}m times K\_\allowbreak{}2,\allowbreak{} m>=\allowbreak{}5,\allowbreak{} put theta\_\allowbreak{}(j,\allowbreak{}a)=\allowbreak{}2 pi j/\allowbreak{}m for both layers. The sine forces along the two cycle neighbors cancel and rung forces vanish. The Hessian is a weighted Laplacian with cycle weights cos(2 pi/\allowbreak{}m)>0 and rung weights1,\allowbreak{} so it is positive definite modulo the global rotation. This is a nonconsensus strict local minimum modulo rotation.


\entrylabel{Second check}

Both the vanishing gradient and Hessian signs were checked;\allowbreak{} the graph is connected and cubic. This is an exact deterministic obstruction to a connectivity-only proof.


\entrylabel{Remaining obstruction}

Must show that uniformly random cubic graphs avoid all such stable phase arrangements with probability tending to one;\allowbreak{} the prism family itself has vanishing probability.


Verification:\allowbreak{} author checked limited argument | Independent review:\allowbreak{} none | Full target solved:\allowbreak{} False


\clearpage

\section*{P0068  Subspectral detection from a growing number of synchronization frequencies}

\entrylabel{Primary source}

Randomstrasse101:\allowbreak{} Open Problems of 2024;\allowbreak{} OpenProblem7,\allowbreak{} PDF p.7


\url{https://arxiv.org/pdf/2504.20539v1}


Source statement date:\allowbreak{} not established. Archived access:\allowbreak{} 2026-09-06T16:\allowbreak{}42:\allowbreak{}33Z.


Historical:\allowbreak{} source\_\allowbreak{}says\_\allowbreak{}open | Current:\allowbreak{} unverified | Result:\allowbreak{} attempted\_\allowbreak{}unsolved


\entrylabel{Question}

For uniform U(1) signal x in C\textasciicircum{}n with |x\_\allowbreak{}i|=\allowbreak{}1 and independent Wigner noises,\allowbreak{} observe Y\_\allowbreak{}k=\allowbreak{}(lambda/\allowbreak{}n)x\textasciicircum{}(k)(x\textasciicircum{}(k))*+\allowbreak{}W\_\allowbreak{}k/\allowbreak{}sqrt(n),\allowbreak{} k=\allowbreak{}1,\allowbreak{}...,\allowbreak{}L\_\allowbreak{}n. Find a scaling of distinct frequencies L\_\allowbreak{}n for which a polynomial-time algorithm reliably detects signal at some fixed lambda<1.


\entrylabel{Scope}

The source also names finite groups. For a fixed finite group the number of distinct irreducible frequencies is bounded;\allowbreak{} a growing group or repeated-observation convention must be specified. The unambiguous U(1) case is the admitted target.


\subsection*{P0068-A001 | Subspectral detection from a growing number of synchronization frequencies special-case or structural calculation}

\entrylabel{Argument or result}

If L independent observations repeated the same rank-one signal,\allowbreak{} averaging would divide noise standard deviation by sqrt(L) and reduce the effective spectral threshold. In the actual model the vectors x\textasciicircum{}(k) differ with k,\allowbreak{} so averaging Y\_\allowbreak{}k yields a sum of different rank-one projectors. It is not a repeated observation of xx*.


\entrylabel{Second check}

The noise-average variance calculation is exact. The signal identity fails already for phases whose squares are not a common multiple of their first powers.


\entrylabel{Remaining obstruction}

Need an algorithm exploiting algebraic relations between distinct powers. Repetition-based variance reduction solves a different experiment.


Verification:\allowbreak{} author checked limited argument | Independent review:\allowbreak{} none | Full target solved:\allowbreak{} False


\clearpage

\section*{P0069  Subexponential detection at fixed synchronization frequencies}

\entrylabel{Primary source}

Randomstrasse101:\allowbreak{} Open Problems of 2024;\allowbreak{} OpenProblem8,\allowbreak{} PDF p.7


\url{https://arxiv.org/pdf/2504.20539v1}


Source statement date:\allowbreak{} not established. Archived access:\allowbreak{} 2026-09-06T16:\allowbreak{}42:\allowbreak{}33Z.


Historical:\allowbreak{} source\_\allowbreak{}says\_\allowbreak{}open | Current:\allowbreak{} unverified | Result:\allowbreak{} attempted\_\allowbreak{}unsolved


\entrylabel{Question}

In the same multifrequency U(1) or specified finite-group Gaussian synchronization model,\allowbreak{} fix L independently of n. For signal strength between its statistical threshold and1,\allowbreak{} is reliable detection possible in time exp(n\textasciicircum{}delta) for some delta<1?


\entrylabel{Scope}

The statistical threshold and group must refer to the fixed model. This is a time-accuracy question distinct from growing the observation model in OpenProblem7.


\subsection*{P0069-A001 | Subexponential detection at fixed synchronization frequencies special-case or structural calculation}

\entrylabel{Argument or result}

Try enumerating an oracle seed of m=\allowbreak{}o(n) signal coordinates and use their edges to estimate a new coordinate in one channel. After multiplying by known seed phases,\allowbreak{} the sum has signal amplitude lambda m/\allowbreak{}n and noise standard deviation sqrt(m/\allowbreak{}n),\allowbreak{} so signal-to-noise ratio is lambda sqrt(m/\allowbreak{}n),\allowbreak{} which tends to zero.


\entrylabel{Second check}

The calculation uses the displayed matrix normalization;\allowbreak{} using unnormalized Wigner noise gives the same ratio. An oracle seed is stronger than an enumerated guessed seed.


\entrylabel{Remaining obstruction}

This only defeats a one-step seed-and-average method. Iterative amplification or a different subexponential algorithm is not ruled out.


Verification:\allowbreak{} author checked limited argument | Independent review:\allowbreak{} none | Full target solved:\allowbreak{} False


\clearpage

\section*{P0070  Kikuchi spectral thresholds as the polynomial degree increases}

\entrylabel{Primary source}

Randomstrasse101:\allowbreak{} Open Problems of 2024;\allowbreak{} Conjecture9,\allowbreak{} PDF p.8


\url{https://arxiv.org/pdf/2504.20539v1}


Source statement date:\allowbreak{} not established. Archived access:\allowbreak{} 2026-09-06T16:\allowbreak{}42:\allowbreak{}33Z.


Historical:\allowbreak{} source\_\allowbreak{}says\_\allowbreak{}open | Current:\allowbreak{} unverified | Result:\allowbreak{} attempted\_\allowbreak{}unsolved


\entrylabel{Question}

For even tensor order r and fixed l>=\allowbreak{}r/\allowbreak{}2,\allowbreak{} index M by l-subsets of [n] and set M\_\allowbreak{}IJ=\allowbreak{}lambda+\allowbreak{}Z\_\allowbreak{}(I symmetric\_\allowbreak{}difference J) when that difference has size r,\allowbreak{} and zero otherwise,\allowbreak{} with independent standard Gaussian Z\_\allowbreak{}S. Does the normalized spectral-outlier threshold n\textasciicircum{}(r/\allowbreak{}4) lambda\_\allowbreak{}critical(r,\allowbreak{}l) tend to zero as l tends to infinity after n tends to infinity?


\entrylabel{Scope}

The source suppresses the n-dependence of its threshold notation. Its literal inclusion of r=\allowbreak{}2 needs a separate check against matrix-PCA limits;\allowbreak{} no silent restriction to r>=\allowbreak{}4 is made and no r=\allowbreak{}2 counterexample is yet claimed.


\subsection*{P0070-A001 | Kikuchi spectral thresholds as the polynomial degree increases special-case or structural calculation}

\entrylabel{Argument or result}

Each row has k=\allowbreak{}binomial(l,\allowbreak{}r/\allowbreak{}2)binomial(n-l,\allowbreak{}r/\allowbreak{}2) nonzero entries. For the normalized all-ones vector v,\allowbreak{} v*M(lambda)*v=\allowbreak{}lambda k+\allowbreak{}(k/\allowbreak{}binomial(n,\allowbreak{}r))sum\_\allowbreak{}(|S|=\allowbreak{}r)Z\_\allowbreak{}S. Thus its noise variance is k\textasciicircum{}2/\allowbreak{}binomial(n,\allowbreak{}r). Write M(0)=\allowbreak{}sum\_\allowbreak{}S Z\_\allowbreak{}S A\_\allowbreak{}S;\allowbreak{} A\_\allowbreak{}S\textasciicircum{}2 is the diagonal indicator of |I intersection S|=\allowbreak{}r/\allowbreak{}2,\allowbreak{} hence sum\_\allowbreak{}S A\_\allowbreak{}S\textasciicircum{}2=\allowbreak{}k I.


\entrylabel{Second check}

Double-counting pairs (I,\allowbreak{}J) with symmetric difference S gives the Rayleigh coefficient. Squaring A\_\allowbreak{}S toggles the same subset S twice and returns I exactly when that toggle preserves cardinality.


\entrylabel{Remaining obstruction}

A Gaussian matrix-series bound introduces sqrt(log binomial(n,\allowbreak{}l)) into the norm estimate. Removing that loss and controlling the actual outlier transition remain necessary;\allowbreak{} the r=\allowbreak{}2 scope caveat needs resolution.


Verification:\allowbreak{} author checked limited argument | Independent review:\allowbreak{} none | Full target solved:\allowbreak{} False


\clearpage

\section*{P0071  Polynomial Glauber mixing throughout the SK high-temperature regime}

\entrylabel{Primary source}

Randomstrasse101:\allowbreak{} Open Problems of 2024;\allowbreak{} OpenProblem15,\allowbreak{} PDF p.12


\url{https://arxiv.org/pdf/2504.20539v1}


Source statement date:\allowbreak{} not established. Archived access:\allowbreak{} 2026-09-06T16:\allowbreak{}42:\allowbreak{}33Z.


Historical:\allowbreak{} source\_\allowbreak{}says\_\allowbreak{}open | Current:\allowbreak{} unverified | Result:\allowbreak{} attempted\_\allowbreak{}unsolved


\entrylabel{Question}

For the Sherrington-Kirkpatrick Ising model with arbitrary external field,\allowbreak{} does single-site heat-bath Glauber dynamics mix in polynomial total-variation time for every fixed beta<1,\allowbreak{} with high probability over the Gaussian disorder?


\entrylabel{Scope}

Keep SK Hamiltonian normalization consistent. A July2026 primary paper proves results for sufficiently large external fields,\allowbreak{} and an August2026 preprint claims optimal mixing for beta<1/\allowbreak{}2 uniformly over fields;\allowbreak{} neither states the full beta<1 arbitrary-field conclusion.


\entrylabel{Dated literature check}

[\{"finding":\allowbreak{} "Large-field result,\allowbreak{} not arbitrary-field full regime",\allowbreak{} "url":\allowbreak{} "\url{https://arxiv.org/abs/2607.06813"%7D,} \{"finding":\allowbreak{} "Preprint claims beta<1/\allowbreak{}2 uniform fields;\allowbreak{} full beta<1 remainder not resolved by its statement",\allowbreak{} "url":\allowbreak{} "\url{https://arxiv.org/abs/2608.22159"%7D]}


\subsection*{P0071-A001 | Polynomial Glauber mixing throughout the SK high-temperature regime special-case or structural calculation}

\entrylabel{Argument or result}

At beta=\allowbreak{}0,\allowbreak{} spins are independent under the target product law. Couple two random-scan heat-bath chains using the same update coordinate and uniform draw. Once a coordinate has been updated it agrees forever. The probability some coordinate remains unupdated by t is at most n(1-1/\allowbreak{}n)\textasciicircum{}t<=\allowbreak{}n exp(-t/\allowbreak{}n),\allowbreak{} giving mixing time at most n log(n/\allowbreak{}epsilon) for every external field.


\entrylabel{Second check}

The update probability of each spin may depend on its field,\allowbreak{} but at beta=\allowbreak{}0 it does not depend on neighboring spins,\allowbreak{} so agreement is preserved. This is a complete endpoint proof,\allowbreak{} with no novelty claim.


\entrylabel{Remaining obstruction}

For beta>0 a later update can destroy agreement through neighboring disagreements. The beta=\allowbreak{}0 coupon-collector argument does not control the whole replica-symmetric interval.


Verification:\allowbreak{} author checked limited argument | Independent review:\allowbreak{} none | Full target solved:\allowbreak{} False


\clearpage

\section*{P0072  Geodesic selection by reinforced ants}

\entrylabel{Primary source}

Aletti et al.:\allowbreak{} Open problems in stochastic reinforcement models;\allowbreak{} March2023 workshop,\allowbreak{} published2025;\allowbreak{} Section2,\allowbreak{} main conjecture,\allowbreak{} PDF pp.1-2


\url{https://www.matrix-inst.org.au/wp_Matrix2016/wp-content/uploads/2023/Aletti.pdf}


Source statement date:\allowbreak{} not established. Archived access:\allowbreak{} 2026-09-06T18:\allowbreak{}00:\allowbreak{}32Z.


Historical:\allowbreak{} source\_\allowbreak{}says\_\allowbreak{}open | Current:\allowbreak{} unverified | Result:\allowbreak{} attempted\_\allowbreak{}unsolved


\entrylabel{Question}

On a finite graph with nest N and food F,\allowbreak{} ants take edge-weighted random walks until hitting F. Reinforce either a uniformly chosen shortest N-F path in the visited trace,\allowbreak{} or the loop erasure of the reversed walk. Are all edges whose weights grow linearly contained in shortest N-F paths of the original graph?


\entrylabel{Scope}

The two reinforcement algorithms form one conjecture family here. The source proves series-parallel and losange cases,\allowbreak{} not all graphs.


\subsection*{P0072-A001 | Geodesic selection by reinforced ants structural argument}

\entrylabel{Argument or result}

On a finite tree with one food vertex,\allowbreak{} every N-F walk has the same simple path after loop erasure,\allowbreak{} and the shortest path in its trace is that path too. Each edge of that path therefore receives exactly one increment per ant;\allowbreak{} every other edge receives zero. The normalized weights converge to its indicator without a probabilistic limit theorem.


\entrylabel{Second check}

Uniqueness of the simple path verifies both algorithms separately. This endpoint is included in known special cases,\allowbreak{} not a new solution.


\entrylabel{Remaining obstruction}

Cycles permit different simple paths and histories affect their weights. The deterministic tree reduction gives no control of their competition.


Verification:\allowbreak{} author checked limited argument | Independent review:\allowbreak{} none | Full target solved:\allowbreak{} False


\clearpage

\section*{P0073  Deterministic limits for trace-reinforced ants}

\entrylabel{Primary source}

Aletti et al.:\allowbreak{} Open problems in stochastic reinforcement models;\allowbreak{} March2023 workshop,\allowbreak{} published2025;\allowbreak{} OpenProblems1-2,\allowbreak{} PDF p.2;\allowbreak{} underlying Kious-Mailler-Schapira Conjecture1.1


\url{https://www.matrix-inst.org.au/wp_Matrix2016/wp-content/uploads/2023/Aletti.pdf}


Source statement date:\allowbreak{} not established. Archived access:\allowbreak{} 2026-09-06T18:\allowbreak{}00:\allowbreak{}32Z.


Historical:\allowbreak{} source\_\allowbreak{}says\_\allowbreak{}open | Current:\allowbreak{} unverified | Result:\allowbreak{} attempted\_\allowbreak{}unsolved


\entrylabel{Question}

When each ant reinforces every edge it visited at least once,\allowbreak{} do all edge weights divided by the number of ants converge almost surely to deterministic constants on general finite graphs,\allowbreak{} and does an appropriate edgewise statement hold on Z\textasciicircum{}2?


\entrylabel{Scope}

The finite and Z\textasciicircum{}2 versions are grouped. The underlying paper excludes multiple edges incident to food;\allowbreak{} the brief workshop statement omits this load-bearing restriction. Its wording main conjecture is interpreted using the immediately preceding trace-process paragraph.


\subsection*{P0073-A001 | Deterministic limits for trace-reinforced ants structural argument}

\entrylabel{Argument or result}

If two parallel edges connect N directly to F,\allowbreak{} an ant crosses exactly one and reinforces it. Starting at weights a,\allowbreak{}b gives the classical Polya urn:\allowbreak{} the first normalized weight has Beta(a,\allowbreak{}b) limit with variance ab/\allowbreak{}((a+\allowbreak{}b)\textasciicircum{}2(a+\allowbreak{}b+\allowbreak{}1))>0. Thus the unrestricted multigraph formulation is false,\allowbreak{} and the underlying no-parallel-food assumption is necessary.


\entrylabel{Second check}

The transition probability and one-increment update coincide exactly with the urn. This is a scope counterexample,\allowbreak{} not a refutation under the original restriction.


\entrylabel{Remaining obstruction}

For simple finite graphs the nonlinear trace-hitting probabilities still require a uniqueness/\allowbreak{}attraction argument. On Z\textasciicircum{}2,\allowbreak{} changing conductances additionally complicate the guarantee that every ant reaches food.


Verification:\allowbreak{} author checked limited argument | Independent review:\allowbreak{} none | Full target solved:\allowbreak{} False


\clearpage

\section*{P0074  Maximal degree and its location in the infinite-path fitness regime}

\entrylabel{Primary source}

Aletti et al.:\allowbreak{} Open problems in stochastic reinforcement models;\allowbreak{} March2023 workshop,\allowbreak{} published2025;\allowbreak{} OpenProblem6,\allowbreak{} PDF p.5


\url{https://www.matrix-inst.org.au/wp_Matrix2016/wp-content/uploads/2023/Aletti.pdf}


Source statement date:\allowbreak{} not established. Archived access:\allowbreak{} 2026-09-06T18:\allowbreak{}00:\allowbreak{}32Z.


Historical:\allowbreak{} source\_\allowbreak{}says\_\allowbreak{}open | Current:\allowbreak{} unverified | Result:\allowbreak{} attempted\_\allowbreak{}unsolved


\entrylabel{Question}

For superlinear preferential attachment with iid positive fitness and attachment f(k,\allowbreak{}w)=\allowbreak{}w k\textasciicircum{}p or k\textasciicircum{}p+\allowbreak{}w,\allowbreak{} determine the growth of the maximum degree and the label attaining it in the regime where the limiting tree is locally finite and has one infinite ray.


\entrylabel{Scope}

Use a positive initial-degree convention;\allowbreak{} the displayed one-vertex degree-zero multiplicative formula otherwise has a zero denominator. The source model must be initialized consistently.


\subsection*{P0074-A001 | Maximal degree and its location in the infinite-path fitness regime structural argument}

\entrylabel{Argument or result}

Let Delta\_\allowbreak{}n be the maximum degree and I\_\allowbreak{}n a maximizing label. If Delta\_\allowbreak{}n tends to infinity and each fixed vertex has finite limiting degree,\allowbreak{} then I\_\allowbreak{}n tends to infinity:\allowbreak{} for every fixed K,\allowbreak{} the maximum of the K finite limiting degrees is finite,\allowbreak{} and eventually Delta\_\allowbreak{}n exceeds it.


\entrylabel{Second check}

This is a pathwise implication and does not assume independence. It deliberately leaves divergence of Delta\_\allowbreak{}n as an extra hypothesis.


\entrylabel{Remaining obstruction}

Neither the rate of Delta\_\allowbreak{}n nor even its divergence follows from local finiteness alone. A deterministic infinite path has maximum degree2.


Verification:\allowbreak{} author checked limited argument | Independent review:\allowbreak{} none | Full target solved:\allowbreak{} False


\clearpage

\section*{P0075  Empirical degree scaling for superlinear attachment with fitness}

\entrylabel{Primary source}

Aletti et al.:\allowbreak{} Open problems in stochastic reinforcement models;\allowbreak{} March2023 workshop,\allowbreak{} published2025;\allowbreak{} OpenProblem7,\allowbreak{} PDF p.5


\url{https://www.matrix-inst.org.au/wp_Matrix2016/wp-content/uploads/2023/Aletti.pdf}


Source statement date:\allowbreak{} not established. Archived access:\allowbreak{} 2026-09-06T18:\allowbreak{}00:\allowbreak{}32Z.


Historical:\allowbreak{} source\_\allowbreak{}says\_\allowbreak{}open | Current:\allowbreak{} unverified | Result:\allowbreak{} attempted\_\allowbreak{}unsolved


\entrylabel{Question}

Determine the scaling and limiting empirical degree distribution of the superlinear preferential-attachment tree with fitness described in Section3,\allowbreak{} including the infinite-path regime.


\entrylabel{Scope}

A degree-count measure and a fitness-degree joint measure are different objects;\allowbreak{} the admitted target is the empirical degree distribution explicitly named in the source.


\subsection*{P0075-A001 | Empirical degree scaling for superlinear attachment with fitness structural argument}

\entrylabel{Argument or result}

For every n-vertex tree,\allowbreak{} the empirical degree law pi\_\allowbreak{}n satisfies sum\_\allowbreak{}k k pi\_\allowbreak{}n(k)=\allowbreak{}2-2/\allowbreak{}n,\allowbreak{} hence pi\_\allowbreak{}n([K,\allowbreak{}infinity))<=\allowbreak{}2/\allowbreak{}K. Thus the unscaled empirical laws are tight on the positive integers. Any weak subsequential limit is a probability law with mean at most2;\allowbreak{} its mean can be smaller because degrees are not uniformly integrable.


\entrylabel{Second check}

The handshaking identity proves the first statement. A star has empirical law tending to delta\_\allowbreak{}1 while its empirical mean tends to2,\allowbreak{} confirming the moment caveat.


\entrylabel{Remaining obstruction}

Tightness supplies no unique limit or model-specific tail scaling. Escape of the first moment must be tracked separately.


Verification:\allowbreak{} author checked limited argument | Independent review:\allowbreak{} none | Full target solved:\allowbreak{} False


\clearpage

\section*{P0076  Multiple-attachment phase transition with superlinear fitness}

\entrylabel{Primary source}

Aletti et al.:\allowbreak{} Open problems in stochastic reinforcement models;\allowbreak{} March2023 workshop,\allowbreak{} published2025;\allowbreak{} OpenProblem8,\allowbreak{} PDF p.5


\url{https://www.matrix-inst.org.au/wp_Matrix2016/wp-content/uploads/2023/Aletti.pdf}


Source statement date:\allowbreak{} not established. Archived access:\allowbreak{} 2026-09-06T18:\allowbreak{}00:\allowbreak{}32Z.


Historical:\allowbreak{} source\_\allowbreak{}says\_\allowbreak{}open | Current:\allowbreak{} unverified | Result:\allowbreak{} attempted\_\allowbreak{}unsolved


\entrylabel{Question}

In the Section3 fitness model,\allowbreak{} replace one attachment per new vertex by m>=\allowbreak{}2 choices and determine how the transition between degree condensation and an infinite-path structure changes.


\entrylabel{Scope}

Sampling with or without replacement and sequential updating must be fixed;\allowbreak{} the source proposes a multigraph extension. One family is counted for all m,\allowbreak{} rather than parameter copies.


\subsection*{P0076-A001 | Multiple-attachment phase transition with superlinear fitness structural argument}

\entrylabel{Argument or result}

After n arrivals with exactly m outgoing attachments each,\allowbreak{} total degree is 2m(n-1) up to the finite seed. Consequently no collection of vertices can each have degree fraction greater than a if it contains more than 2m/\allowbreak{}a vertices. This bounds the number of macroscopic hubs but permits several of them,\allowbreak{} unlike a single-star conclusion.


\entrylabel{Second check}

The bound follows by summing nonnegative degrees;\allowbreak{} it remains valid with parallel edges and sequential updates.


\entrylabel{Remaining obstruction}

A tree embedding into independent genealogical birth processes is lost when a vertex has several parents. The degree budget does not identify the phase boundary.


Verification:\allowbreak{} author checked limited argument | Independent review:\allowbreak{} none | Full target solved:\allowbreak{} False


\clearpage

\section*{P0077  Fluctuations of fitness-reinforced sequence sums}

\entrylabel{Primary source}

Aletti et al.:\allowbreak{} Open problems in stochastic reinforcement models;\allowbreak{} March2023 workshop,\allowbreak{} published2025;\allowbreak{} OpenProblem9,\allowbreak{} PDF p.7


\url{https://www.matrix-inst.org.au/wp_Matrix2016/wp-content/uploads/2023/Aletti.pdf}


Source statement date:\allowbreak{} not established. Archived access:\allowbreak{} 2026-09-06T18:\allowbreak{}00:\allowbreak{}32Z.


Historical:\allowbreak{} source\_\allowbreak{}says\_\allowbreak{}open | Current:\allowbreak{} unverified | Result:\allowbreak{} attempted\_\allowbreak{}unsolved


\entrylabel{Question}

Let a new positive value be sampled from mu with probability beta,\allowbreak{} and otherwise copy a previous value with probability proportional to value times multiplicity. With ess sup mu=\allowbreak{}1 and mu(\{1\})=\allowbreak{}0,\allowbreak{} find the fluctuation limit of sum\_\allowbreak{}(i<=\allowbreak{}n)(X\_\allowbreak{}i-lambda\_\allowbreak{}star),\allowbreak{} where lambda\_\allowbreak{}star is the Section4 law-of-large-numbers constant.


\entrylabel{Scope}

The source equation beta integral x/\allowbreak{}(lambda\_\allowbreak{}star-(1-beta)x) dmu=\allowbreak{}1 defines lambda\_\allowbreak{}star when possible;\allowbreak{} otherwise lambda\_\allowbreak{}star=\allowbreak{}1-beta. A universal square-root normalization is not assumed.


\subsection*{P0077-A001 | Fluctuations of fitness-reinforced sequence sums structural argument}

\entrylabel{Argument or result}

Writing S\_\allowbreak{}n=\allowbreak{}sum X\_\allowbreak{}i and Q\_\allowbreak{}n=\allowbreak{}sum X\_\allowbreak{}i\textasciicircum{}2 gives E[X\_\allowbreak{}(n+\allowbreak{}1)|F\_\allowbreak{}n]=\allowbreak{}beta m\_\allowbreak{}1+\allowbreak{}(1-beta)Q\_\allowbreak{}n/\allowbreak{}S\_\allowbreak{}n. Subtracting this predictable mean yields bounded martingale differences,\allowbreak{} but the centered sum also contains sum[beta m\_\allowbreak{}1+\allowbreak{}(1-beta)Q\_\allowbreak{}i/\allowbreak{}S\_\allowbreak{}i-lambda\_\allowbreak{}star].


\entrylabel{Second check}

The conditional mean is obtained by summing x times its copying probability. Boundedness controls the martingale part without controlling the accumulated drift.


\entrylabel{Remaining obstruction}

A martingale CLT alone cannot settle the question:\allowbreak{} fluctuations of Q\_\allowbreak{}n/\allowbreak{}S\_\allowbreak{}n and the slowly relaxing drift may contribute at the same or larger scale.


Verification:\allowbreak{} author checked limited argument | Independent review:\allowbreak{} none | Full target solved:\allowbreak{} False


\clearpage

\section*{P0078  Boundary profile of the reinforced-sequence condensate}

\entrylabel{Primary source}

Aletti et al.:\allowbreak{} Open problems in stochastic reinforcement models;\allowbreak{} March2023 workshop,\allowbreak{} published2025;\allowbreak{} OpenProblem10 and displayed conjecture,\allowbreak{} PDF pp.7-8


\url{https://www.matrix-inst.org.au/wp_Matrix2016/wp-content/uploads/2023/Aletti.pdf}


Source statement date:\allowbreak{} not established. Archived access:\allowbreak{} 2026-09-06T18:\allowbreak{}00:\allowbreak{}32Z.


Historical:\allowbreak{} source\_\allowbreak{}says\_\allowbreak{}open | Current:\allowbreak{} unverified | Result:\allowbreak{} attempted\_\allowbreak{}unsolved


\entrylabel{Question}

When beta integral (1-x)\textasciicircum{}(-1) dmu<1 and mu(1-x,\allowbreak{}1) is regularly varying at zero with index alpha>1,\allowbreak{} determine a\_\allowbreak{}n tending to infinity and c>0 for which n\textasciicircum{}(-1) sum f(a\_\allowbreak{}n(1-X\_\allowbreak{}i)) converges in probability to c integral\_\allowbreak{}0\textasciicircum{}infinity f(x)e\textasciicircum{}(-x)x\textasciicircum{}alpha dx for compactly supported continuous f.


\entrylabel{Scope}

Convergence against compactly supported test functions does not by itself identify total condensate mass;\allowbreak{} mass at zero or infinity must be separately controlled.


\subsection*{P0078-A001 | Boundary profile of the reinforced-sequence condensate structural argument}

\entrylabel{Argument or result}

Let rho=\allowbreak{}1-beta integral (1-x)\textasciicircum{}(-1)dmu be the mass missing from the bulk limit. If the claimed rescaled measure is tight on (0,\allowbreak{}infinity) and carries all that missing mass,\allowbreak{} integrating its density forces c=\allowbreak{}rho/\allowbreak{}Gamma(alpha+\allowbreak{}1). Without those two additional controls this identification is invalid.


\entrylabel{Second check}

The integral of e\textasciicircum{}(-x)x\textasciicircum{}alpha equals Gamma(alpha+\allowbreak{}1);\allowbreak{} the deduction is explicitly conditional on no escaped mass.


\entrylabel{Remaining obstruction}

The source provides only vague convergence away from the boundary. Finding a\_\allowbreak{}n and proving boundary tightness are the load-bearing missing steps.


Verification:\allowbreak{} author checked limited argument | Independent review:\allowbreak{} none | Full target solved:\allowbreak{} False


\clearpage

\section*{P0079  Growth of sums with unbounded reinforcement fitness}

\entrylabel{Primary source}

Aletti et al.:\allowbreak{} Open problems in stochastic reinforcement models;\allowbreak{} March2023 workshop,\allowbreak{} published2025;\allowbreak{} OpenProblem11,\allowbreak{} PDF p.8


\url{https://www.matrix-inst.org.au/wp_Matrix2016/wp-content/uploads/2023/Aletti.pdf}


Source statement date:\allowbreak{} not established. Archived access:\allowbreak{} 2026-09-06T18:\allowbreak{}00:\allowbreak{}32Z.


Historical:\allowbreak{} source\_\allowbreak{}says\_\allowbreak{}open | Current:\allowbreak{} unverified | Result:\allowbreak{} attempted\_\allowbreak{}unsolved


\entrylabel{Question}

For the same innovation/\allowbreak{}copying scheme with unbounded positive innovation law mu,\allowbreak{} determine the superlinear growth rate of S\_\allowbreak{}n=\allowbreak{}sum\_\allowbreak{}(i<=\allowbreak{}n) X\_\allowbreak{}i as a function of the tail of mu.


\entrylabel{Scope}

No single rate can be asserted for every unbounded mu without describing its tail;\allowbreak{} the question is a distribution-dependent classification.


\subsection*{P0079-A001 | Growth of sums with unbounded reinforcement fitness structural argument}

\entrylabel{Argument or result}

Let K\_\allowbreak{}n count innovation draws and let Y\_\allowbreak{}j be their iid values. Every innovation appears at least once,\allowbreak{} so S\_\allowbreak{}n>=\allowbreak{}sum\_\allowbreak{}(j<=\allowbreak{}K\_\allowbreak{}n)Y\_\allowbreak{}j and S\_\allowbreak{}n>=\allowbreak{}max\_\allowbreak{}(j<=\allowbreak{}K\_\allowbreak{}n)Y\_\allowbreak{}j,\allowbreak{} while K\_\allowbreak{}n/\allowbreak{}n tends to beta almost surely. If E\_\allowbreak{}mu Y=\allowbreak{}infinity,\allowbreak{} the truncated strong law and then monotone convergence give S\_\allowbreak{}n/\allowbreak{}n tending to infinity from this lower bound alone.


\entrylabel{Second check}

For a truncation L,\allowbreak{} divide the innovation sum by n and use the iid strong law;\allowbreak{} let L increase afterward.


\entrylabel{Remaining obstruction}

Finite-mean unbounded mu requires amplification by reinforcement;\allowbreak{} innovation lower bounds alone remain linear and cannot give the requested rate.


Verification:\allowbreak{} author checked limited argument | Independent review:\allowbreak{} none | Full target solved:\allowbreak{} False


\clearpage

\section*{P0080  Weak reinforcement on infinite WARM networks}

\entrylabel{Primary source}

Aletti et al.:\allowbreak{} Open problems in stochastic reinforcement models;\allowbreak{} March2023 workshop,\allowbreak{} published2025;\allowbreak{} Conjecture1,\allowbreak{} PDF p.10


\url{https://www.matrix-inst.org.au/wp_Matrix2016/wp-content/uploads/2023/Aletti.pdf}


Source statement date:\allowbreak{} not established. Archived access:\allowbreak{} 2026-09-06T18:\allowbreak{}00:\allowbreak{}32Z.


Historical:\allowbreak{} source\_\allowbreak{}says\_\allowbreak{}open | Current:\allowbreak{} unverified | Result:\allowbreak{} attempted\_\allowbreak{}unsolved


\entrylabel{Question}

For a bounded-degree infinite graph,\allowbreak{} bounded positive vertex firing rates,\allowbreak{} and WARM exponent 0<alpha<1,\allowbreak{} establish convergence of every normalized edge count to the intended unique positive equilibrium.


\entrylabel{Scope}

The source literally says unique equilibrium while allowing nonnegative coordinates. Boundary equilibria make that literal uniqueness false;\allowbreak{} the meaningful positive-attractor formulation is recorded explicitly,\allowbreak{} not silently substituted.


\subsection*{P0080-A001 | Weak reinforcement on infinite WARM networks structural argument}

\entrylabel{Argument or result}

On Z with all firing rates1,\allowbreak{} set edge values alternately2 and0. At each vertex the only positive incident edge is selected with probability1;\allowbreak{} its two endpoints supply total rate2,\allowbreak{} while zero edges get rate0. This is an equilibrium for every alpha>0. Shifting gives a second one,\allowbreak{} and the uniform edge value1 is a third. Thus uniqueness among all nonnegative equilibria is false.


\entrylabel{Second check}

All denominators equal2\textasciicircum{}alpha,\allowbreak{} so the example has no undefined 0/\allowbreak{}0 term. It does not contradict convergence from strictly positive initial counts.


\entrylabel{Remaining obstruction}

The positive-equilibrium uniqueness and stochastic attraction problem on general infinite graphs remains;\allowbreak{} no finite total Lyapunov function is available.


Verification:\allowbreak{} author checked limited argument | Independent review:\allowbreak{} none | Full target solved:\allowbreak{} False


\clearpage

\section*{P0081  Critical WARM on the integer line}

\entrylabel{Primary source}

Aletti et al.:\allowbreak{} Open problems in stochastic reinforcement models;\allowbreak{} March2023 workshop,\allowbreak{} published2025;\allowbreak{} Conjecture2,\allowbreak{} PDF p.10


\url{https://www.matrix-inst.org.au/wp_Matrix2016/wp-content/uploads/2023/Aletti.pdf}


Source statement date:\allowbreak{} not established. Archived access:\allowbreak{} 2026-09-06T18:\allowbreak{}00:\allowbreak{}32Z.


Historical:\allowbreak{} source\_\allowbreak{}says\_\allowbreak{}open | Current:\allowbreak{} unverified | Result:\allowbreak{} attempted\_\allowbreak{}unsolved


\entrylabel{Question}

For WARM on Z with alpha=\allowbreak{}1 and all vertex firing rates1,\allowbreak{} prove that each edge count divided by time converges almost surely to1.


\entrylabel{Scope}

Keep the stochastic count process distinct from its deterministic equilibrium equation.


\subsection*{P0081-A001 | Critical WARM on the integer line structural argument}

\entrylabel{Argument or result}

For any a,\allowbreak{}b>0 with a+\allowbreak{}b=\allowbreak{}2,\allowbreak{} the alternating edge configuration a,\allowbreak{}b,\allowbreak{}a,\allowbreak{}b,\allowbreak{}... solves the equilibrium equation:\allowbreak{} an a-edge receives a/\allowbreak{}(a+\allowbreak{}b) from each endpoint,\allowbreak{} totaling a,\allowbreak{} and similarly for b. Thus the ODE has a continuum of positive equilibria on the same graph.


\entrylabel{Second check}

Substitution in the local rate formula verifies every edge. This extends the source even-cycle observation but concerns deterministic equilibria only.


\entrylabel{Remaining obstruction}

Convergence to1 cannot follow merely from uniqueness of ODE equilibria,\allowbreak{} which is false. One must exploit stochastic dynamics or spatial structure to select the uniform equilibrium.


Verification:\allowbreak{} author checked limited argument | Independent review:\allowbreak{} none | Full target solved:\allowbreak{} False


\clearpage

\section*{P0082  Sharp reinforcement threshold for WARM forests}

\entrylabel{Primary source}

Aletti et al.:\allowbreak{} Open problems in stochastic reinforcement models;\allowbreak{} March2023 workshop,\allowbreak{} published2025;\allowbreak{} OpenProblem12,\allowbreak{} PDF pp.10-11


\url{https://www.matrix-inst.org.au/wp_Matrix2016/wp-content/uploads/2023/Aletti.pdf}


Source statement date:\allowbreak{} not established. Archived access:\allowbreak{} 2026-09-06T18:\allowbreak{}00:\allowbreak{}32Z.


Historical:\allowbreak{} source\_\allowbreak{}says\_\allowbreak{}open | Current:\allowbreak{} unverified | Result:\allowbreak{} attempted\_\allowbreak{}unsolved


\entrylabel{Question}

Improve the known sufficient threshold alpha>2 ensuring that every connected component of a stable finite-graph WARM equilibrium is a tree,\allowbreak{} allowing arbitrary positive firing rates.


\entrylabel{Scope}

The source identifies the uniform triangle threshold4/\allowbreak{}3 as a lower obstruction. A rederivation of it is not a new threshold theorem.


\subsection*{P0082-A001 | Sharp reinforcement threshold for WARM forests structural argument}

\entrylabel{Argument or result}

At a positive equilibrium x,\allowbreak{} put s\_\allowbreak{}v=\allowbreak{}sum\_\allowbreak{}(e\textasciitilde{}v)x\_\allowbreak{}e\textasciicircum{}alpha. The Hessian of the correctly rate-weighted Lyapunov function satisfies H[delta,\allowbreak{}delta]=\allowbreak{}(alpha-1)sum\_\allowbreak{}e delta\_\allowbreak{}e\textasciicircum{}2/\allowbreak{}x\_\allowbreak{}e-alpha sum\_\allowbreak{}v p\_\allowbreak{}v/\allowbreak{}s\_\allowbreak{}v\textasciicircum{}2 (sum\_\allowbreak{}(e\textasciitilde{}v)x\_\allowbreak{}e\textasciicircum{}(alpha-1)delta\_\allowbreak{}e)\textasciicircum{}2. Along an even cycle choose alternating delta\_\allowbreak{}e=\allowbreak{}+\allowbreak{}/\allowbreak{}-x\_\allowbreak{}e\textasciicircum{}(1-alpha),\allowbreak{} zero elsewhere. Every vertex sum vanishes,\allowbreak{} giving a positive direction for alpha>1 and hence linear instability.


\entrylabel{Second check}

The equilibrium identity x\_\allowbreak{}e=\allowbreak{}sum\_\allowbreak{}(v\textasciitilde{}e)p\_\allowbreak{}v x\_\allowbreak{}e\textasciicircum{}alpha/\allowbreak{}s\_\allowbreak{}v reduces the diagonal term. Multiplication by the positive diagonal x preserves the Hessian inertia of the linearized gradient flow.


\entrylabel{Remaining obstruction}

This eliminates even cycles,\allowbreak{} not odd cycles or all odd unicyclic supports. It does not improve the universal forest threshold by itself.


Verification:\allowbreak{} author checked limited argument | Independent review:\allowbreak{} none | Full target solved:\allowbreak{} False


\clearpage

\section*{P0083  Sharp reinforcement threshold for WARM whiskers}

\entrylabel{Primary source}

Aletti et al.:\allowbreak{} Open problems in stochastic reinforcement models;\allowbreak{} March2023 workshop,\allowbreak{} published2025;\allowbreak{} OpenProblem13,\allowbreak{} PDF pp.10-11


\url{https://www.matrix-inst.org.au/wp_Matrix2016/wp-content/uploads/2023/Aletti.pdf}


Source statement date:\allowbreak{} not established. Archived access:\allowbreak{} 2026-09-06T18:\allowbreak{}00:\allowbreak{}32Z.


Historical:\allowbreak{} source\_\allowbreak{}says\_\allowbreak{}open | Current:\allowbreak{} unverified | Result:\allowbreak{} attempted\_\allowbreak{}unsolved


\entrylabel{Question}

For finite WARM graphs with uniform vertex firing rates,\allowbreak{} improve the sufficient threshold alpha>25 for every surviving stable component to have diameter at most3.


\entrylabel{Scope}

This is a geometry question beyond the forest conclusion;\allowbreak{} all alpha values are grouped as one threshold problem.


\subsection*{P0083-A001 | Sharp reinforcement threshold for WARM whiskers structural argument}

\entrylabel{Argument or result}

Consider a support path of four edges with reflection-symmetric values a,\allowbreak{}b,\allowbreak{}b,\allowbreak{}a and unit vertex rates. The endpoint equation is a=\allowbreak{}1+\allowbreak{}a\textasciicircum{}alpha/\allowbreak{}(a\textasciicircum{}alpha+\allowbreak{}b\textasciicircum{}alpha),\allowbreak{} so 1<a<2. The neighboring inner-edge equation is b=\allowbreak{}b\textasciicircum{}alpha/\allowbreak{}(a\textasciicircum{}alpha+\allowbreak{}b\textasciicircum{}alpha)+\allowbreak{}1/\allowbreak{}2. Summing gives a+\allowbreak{}b=\allowbreak{}5/\allowbreak{}2. Thus any such equilibrium reduces to a scalar equation with b=\allowbreak{}5/\allowbreak{}2-a.


\entrylabel{Second check}

The total rate of the five vertices is5,\allowbreak{} equal to2a+\allowbreak{}2b. The central vertex contributes1/\allowbreak{}2 to each equal inner edge.


\entrylabel{Remaining obstruction}

Existence and stability of this scalar branch must still be analyzed,\allowbreak{} and excluding this symmetric path alone would not exclude all larger asymmetric trees.


Verification:\allowbreak{} author checked limited argument | Independent review:\allowbreak{} none | Full target solved:\allowbreak{} False


\clearpage

\section*{P0084  Finite surviving components in infinite superlinear WARM}

\entrylabel{Primary source}

Aletti et al.:\allowbreak{} Open problems in stochastic reinforcement models;\allowbreak{} March2023 workshop,\allowbreak{} published2025;\allowbreak{} Conjecture3,\allowbreak{} PDF p.11


\url{https://www.matrix-inst.org.au/wp_Matrix2016/wp-content/uploads/2023/Aletti.pdf}


Source statement date:\allowbreak{} not established. Archived access:\allowbreak{} 2026-09-06T18:\allowbreak{}00:\allowbreak{}32Z.


Historical:\allowbreak{} source\_\allowbreak{}says\_\allowbreak{}open | Current:\allowbreak{} unverified | Result:\allowbreak{} attempted\_\allowbreak{}unsolved


\entrylabel{Question}

With bounded-degree graph,\allowbreak{} firing rates bounded above and bounded away from zero,\allowbreak{} and alpha>1,\allowbreak{} prove edgewise convergence of normalized WARM counts and finiteness of every component of surviving edges.


\entrylabel{Scope}

The lower bound on firing rates is essential in the source:\allowbreak{} examples with decaying rates on trees have infinitely many infinite components.


\subsection*{P0084-A001 | Finite surviving components in infinite superlinear WARM structural argument}

\entrylabel{Argument or result}

If every edge had survival probability at most q,\allowbreak{} this would not imply domination by iid Bernoulli percolation. For example,\allowbreak{} activate all edges simultaneously with probability q and none otherwise. Every marginal is q,\allowbreak{} yet there is an infinite component with probability q for every q>0 on an infinite connected graph.


\entrylabel{Second check}

The example computes the marginals and connectivity exactly. It is a counterexample to a marginal-only proof method,\allowbreak{} not to WARM.


\entrylabel{Remaining obstruction}

A valid comparison needs conditional probability bounds or explicit dependence control,\allowbreak{} together with convergence of the edge weights;\allowbreak{} small one-edge probabilities are insufficient.


Verification:\allowbreak{} author checked limited argument | Independent review:\allowbreak{} none | Full target solved:\allowbreak{} False


\clearpage

\section*{P0085  WARM percolation with vanishing rates on lattices}

\entrylabel{Primary source}

Aletti et al.:\allowbreak{} Open problems in stochastic reinforcement models;\allowbreak{} March2023 workshop,\allowbreak{} published2025;\allowbreak{} OpenProblem14,\allowbreak{} PDF p.11


\url{https://www.matrix-inst.org.au/wp_Matrix2016/wp-content/uploads/2023/Aletti.pdf}


Source statement date:\allowbreak{} not established. Archived access:\allowbreak{} 2026-09-06T18:\allowbreak{}00:\allowbreak{}32Z.


Historical:\allowbreak{} source\_\allowbreak{}says\_\allowbreak{}open | Current:\allowbreak{} unverified | Result:\allowbreak{} attempted\_\allowbreak{}unsolved


\entrylabel{Question}

Do bounded positive firing rates on Z\textasciicircum{}d,\allowbreak{} d>=\allowbreak{}2,\allowbreak{} which are not bounded away from zero,\allowbreak{} and some alpha>1 yield a positive probability of an infinite component of surviving WARM edges? If so,\allowbreak{} determine uniqueness.


\entrylabel{Scope}

The existence and uniqueness parts are grouped. Tree examples do not directly embed as bounded-length lattice edges.


\subsection*{P0085-A001 | WARM percolation with vanishing rates on lattices structural argument}

\entrylabel{Argument or result}

There is no bounded-length embedding of a b-ary tree,\allowbreak{} b>=\allowbreak{}2,\allowbreak{} into Z\textasciicircum{}d with distinct vertices:\allowbreak{} its depth-r ball contains order b\textasciicircum{}r vertices,\allowbreak{} whereas images lie in a lattice ball of radius Lr with only O(r\textasciicircum{}d) vertices. Therefore transferring the known regular-tree construction requires unbounded subdivisions or a different geometry.


\entrylabel{Second check}

This is the exponential-versus-polynomial volume comparison,\allowbreak{} independent of reinforcement details.


\entrylabel{Remaining obstruction}

Unbounded subdivisions change the local competition and rates along paths. The volume obstruction explains why a direct bounded-distortion transfer fails but does not exclude lattice WARM percolation.


Verification:\allowbreak{} author checked limited argument | Independent review:\allowbreak{} none | Full target solved:\allowbreak{} False


\clearpage

\section*{P0086  Recurrence of sublinear vertex-reinforced walk on Z}

\entrylabel{Primary source}

Aletti et al.:\allowbreak{} Open problems in stochastic reinforcement models;\allowbreak{} March2023 workshop,\allowbreak{} published2025;\allowbreak{} OpenProblem15,\allowbreak{} PDF p.12


\url{https://www.matrix-inst.org.au/wp_Matrix2016/wp-content/uploads/2023/Aletti.pdf}


Source statement date:\allowbreak{} not established. Archived access:\allowbreak{} 2026-09-06T18:\allowbreak{}00:\allowbreak{}32Z.


Historical:\allowbreak{} source\_\allowbreak{}says\_\allowbreak{}open | Current:\allowbreak{} unverified | Result:\allowbreak{} attempted\_\allowbreak{}unsolved


\entrylabel{Question}

For nearest-neighbor vertex-reinforced random walk on Z,\allowbreak{} with transition weights w(1+\allowbreak{}local\_\allowbreak{}time) and positive w satisfying w(x)=\allowbreak{}o(x\textasciicircum{}alpha) for some alpha<1,\allowbreak{} is the walk recurrent?


\entrylabel{Scope}

The source states positivity and the growth upper bound,\allowbreak{} without an explicit monotonicity hypothesis. Any use of monotonicity must be flagged.


\subsection*{P0086-A001 | Recurrence of sublinear vertex-reinforced walk on Z structural argument}

\entrylabel{Argument or result}

For constant w the walk is simple symmetric random walk. Starting at1 and stopping at0 or m,\allowbreak{} the harmonic function h(k)=\allowbreak{}k/\allowbreak{}m gives probability1/\allowbreak{}m of hitting m before0. Letting m tend to infinity shows that the probability of avoiding0 forever is0;\allowbreak{} by the strong Markov property returns recur indefinitely.


\entrylabel{Second check}

This is the standard gambler-ruin proof of the constant-weight case,\allowbreak{} not a new result. Constant weights satisfy the stated sublinear condition for any alpha in(0,\allowbreak{}1).


\entrylabel{Remaining obstruction}

For nonconstant reinforcement,\allowbreak{} position alone is not Markov and the gambler-ruin harmonic function no longer solves the conditional transition equation.


Verification:\allowbreak{} author checked limited argument | Independent review:\allowbreak{} none | Full target solved:\allowbreak{} False


\clearpage

\section*{P0087  Forced network urn convergence and diffusion scaling}

\entrylabel{Primary source}

Aletti et al.:\allowbreak{} Open problems in stochastic reinforcement models;\allowbreak{} March2023 workshop,\allowbreak{} published2025;\allowbreak{} OpenProblem16,\allowbreak{} PDF pp.14-16


\url{https://www.matrix-inst.org.au/wp_Matrix2016/wp-content/uploads/2023/Aletti.pdf}


Source statement date:\allowbreak{} not established. Archived access:\allowbreak{} 2026-09-06T18:\allowbreak{}00:\allowbreak{}32Z.


Historical:\allowbreak{} source\_\allowbreak{}says\_\allowbreak{}open | Current:\allowbreak{} unverified | Result:\allowbreak{} attempted\_\allowbreak{}unsolved


\entrylabel{Question}

For Z\_\allowbreak{}(n+\allowbreak{}1)=\allowbreak{}(1-r-f)Z\_\allowbreak{}n+\allowbreak{}r X\_\allowbreak{}(n+\allowbreak{}1)+\allowbreak{}f p with E[X\_\allowbreak{}(n+\allowbreak{}1)|F\_\allowbreak{}n]=\allowbreak{}Q Z\_\allowbreak{}n,\allowbreak{} Q=\allowbreak{}A\textasciicircum{}T irreducible row-stochastic and X binary coordinatewise,\allowbreak{} determine nonconvergence and the continuous-time limit when r tends to0 and f is asymptotic to c r\textasciicircum{}2.


\entrylabel{Scope}

Assume r,\allowbreak{}f>0 and r+\allowbreak{}f<=\allowbreak{}1. The source permits 0<=\allowbreak{}p<=\allowbreak{}1,\allowbreak{} including extreme vectors;\allowbreak{} those have a different convergence answer. The displayed conditional mean alone does not specify the cross-coordinate noise covariance.


\subsection*{P0087-A001 | Forced network urn convergence and diffusion scaling structural argument}

\entrylabel{Argument or result}

A complete fixed-step classification follows. On any event Z\_\allowbreak{}n->z,\allowbreak{} the update forces X\_\allowbreak{}n to converge to a binary vector x and hence eventually equal it. The bounded martingale-difference strong law gives x=\allowbreak{}Qz. The update limit gives z=\allowbreak{}(r x+\allowbreak{}f p)/\allowbreak{}(r+\allowbreak{}f). A binary coordinate of Qz can equal0 or1 only when all its positive-weight predecessors have the same x-value and matching p-value. Irreducibility then forces x and p to be the all-zero or all-one vector. Conversely if p=\allowbreak{}0,\allowbreak{} the positive stationary vector pi gives E[pi.Z\_\allowbreak{}n]=\allowbreak{}(1-f)\textasciicircum{}n E[pi.Z\_\allowbreak{}0],\allowbreak{} and Markov plus Borel-Cantelli yields Z\_\allowbreak{}n->0;\allowbreak{} p=\allowbreak{}1 follows by complementing.


\entrylabel{Second check}

The proof needs no conditional independence. The strong law follows from summable variances of martingale increments divided by n,\allowbreak{} then Kronecker. The maximum principle can also be checked by propagating the set of zero coordinates backwards through irreducibility.


\entrylabel{Remaining obstruction}

This resolves the fixed-step convergence subquestion under the stated assumptions,\allowbreak{} but not the whole diffusion question. Fast synchronization,\allowbreak{} an initial layer,\allowbreak{} and the covariance specification must be treated before a full solution label.


Verification:\allowbreak{} author checked limited argument | Independent review:\allowbreak{} none | Full target solved:\allowbreak{} False


\clearpage

\section*{P0088  Matching local limits of Bernoulli and Poisson-outdegree attachment}

\entrylabel{Primary source}

Aletti et al.:\allowbreak{} Open problems in stochastic reinforcement models;\allowbreak{} March2023 workshop,\allowbreak{} published2025;\allowbreak{} OpenProblem23,\allowbreak{} PDF p.21


\url{https://www.matrix-inst.org.au/wp_Matrix2016/wp-content/uploads/2023/Aletti.pdf}


Source statement date:\allowbreak{} not established. Archived access:\allowbreak{} 2026-09-06T18:\allowbreak{}00:\allowbreak{}32Z.


Historical:\allowbreak{} source\_\allowbreak{}says\_\allowbreak{}open | Current:\allowbreak{} unverified | Result:\allowbreak{} attempted\_\allowbreak{}unsolved


\entrylabel{Question}

Relate the local weak limit of the Dereich-Morters Bernoulli preferential-attachment model,\allowbreak{} whose older vertices are chosen independently with probability f(in\_\allowbreak{}degree)/\allowbreak{}(n-1),\allowbreak{} to an appropriately matched model with iid Poisson outgoing degree and attachment based on in-degree.


\entrylabel{Scope}

The matching of the affine rule and Poisson parameter is part of the question. Early probabilities exceeding1 need a consistent initialization or cap.


\subsection*{P0088-A001 | Matching local limits of Bernoulli and Poisson-outdegree attachment structural argument}

\entrylabel{Argument or result}

Conditional on current degrees,\allowbreak{} the Bernoulli outgoing count has mean Lambda\_\allowbreak{}n=\allowbreak{}sum\_\allowbreak{}i p\_\allowbreak{}ni. Its total variation distance to Poisson(Lambda\_\allowbreak{}n) is at most sum\_\allowbreak{}i p\_\allowbreak{}ni\textasciicircum{}2 by the elementary independent Bernoulli-to-Poisson coupling. If max\_\allowbreak{}i p\_\allowbreak{}ni->0 and Lambda\_\allowbreak{}n stays bounded,\allowbreak{} this error tends to0.


\entrylabel{Second check}

Use sum p\_\allowbreak{}i\textasciicircum{}2<=\allowbreak{}max p\_\allowbreak{}i sum p\_\allowbreak{}i. This verifies the count approximation only;\allowbreak{} the vertex targets retain degree-dependent marks.


\entrylabel{Remaining obstruction}

A matching outgoing count does not prove local graph convergence. Target choices,\allowbreak{} ancestor types,\allowbreak{} and dependencies across successive births must be coupled over an explored neighborhood.


Verification:\allowbreak{} author checked limited argument | Independent review:\allowbreak{} none | Full target solved:\allowbreak{} False


\clearpage

\section*{P0089  Local limits for nonaffine attachment with multiple edges}

\entrylabel{Primary source}

Aletti et al.:\allowbreak{} Open problems in stochastic reinforcement models;\allowbreak{} March2023 workshop,\allowbreak{} published2025;\allowbreak{} OpenProblem24,\allowbreak{} PDF p.21


\url{https://www.matrix-inst.org.au/wp_Matrix2016/wp-content/uploads/2023/Aletti.pdf}


Source statement date:\allowbreak{} not established. Archived access:\allowbreak{} 2026-09-06T18:\allowbreak{}00:\allowbreak{}32Z.


Historical:\allowbreak{} source\_\allowbreak{}says\_\allowbreak{}open | Current:\allowbreak{} unverified | Result:\allowbreak{} attempted\_\allowbreak{}unsolved


\entrylabel{Question}

Determine local weak limits of preferential-attachment graphs with fixed outgoing degree m>=\allowbreak{}2 and nonaffine attachment function f,\allowbreak{} extending the known tree and affine cases.


\entrylabel{Scope}

State a nonexplosive/\allowbreak{}noncondensed class of f before expecting a locally finite limit;\allowbreak{} the source asks for a theory beyond the affine case,\allowbreak{} not one arbitrary universal formula.


\subsection*{P0089-A001 | Local limits for nonaffine attachment with multiple edges structural argument}

\entrylabel{Argument or result}

With m=\allowbreak{}1 and equal attachment weights,\allowbreak{} the graph is a uniform recursive tree and its empirical degree law is tight by the tree degree identity. With m>=\allowbreak{}2 the degree budget still gives uniform-root degree tightness,\allowbreak{} since its mean is at most2m plus a vanishing seed term.


\entrylabel{Second check}

Markov gives P(deg(U\_\allowbreak{}n)>K)<=\allowbreak{}2m/\allowbreak{}K+\allowbreak{}o(1) without assumptions on f.


\entrylabel{Remaining obstruction}

Tightness of the root degree alone does not ensure tightness of radius-two neighborhoods:\allowbreak{} a uniformly selected leaf can attach to a hub of diverging degree. Condensation and size bias must be controlled.


Verification:\allowbreak{} author checked limited argument | Independent review:\allowbreak{} none | Full target solved:\allowbreak{} False


\clearpage

\section*{P0090  Multistate pivot-tree criterion for interpretation}

\entrylabel{Primary source}

Random tree recursions:\allowbreak{} which fixed points correspond to tangible sets of trees?;\allowbreak{} Section6,\allowbreak{} PDF p.39,\allowbreak{} first paragraph


\url{https://arxiv.org/pdf/1808.03019v2}


Source statement date:\allowbreak{} not established. Archived access:\allowbreak{} 2026-09-06T16:\allowbreak{}42:\allowbreak{}18Z.


Historical:\allowbreak{} source\_\allowbreak{}says\_\allowbreak{}open | Current:\allowbreak{} unverified | Result:\allowbreak{} attempted\_\allowbreak{}unsolved


\entrylabel{Question}

Extend the two-state interpretation criterion of Theorem1.7 to finite color sets of size at least3,\allowbreak{} covering critical and supercritical pivot trees even without positive regularity.


\entrylabel{Scope}

A deterministic automaton acts on child-color counts. Interpretability is measurability of root color from the uncolored Galton-Watson tree;\allowbreak{} positive-mass colors and the standing offspring assumptions are retained.


\subsection*{P0090-A001 | Multistate pivot-tree criterion for interpretation structural argument}

\entrylabel{Argument or result}

For a finite-type pivot branching process with mean matrix M and spectral radius less than1,\allowbreak{} E[population at level n] tends to0,\allowbreak{} hence P(level n nonempty)<=\allowbreak{}E[size] tends to0. This proves extinction and explains the subcritical direction already covered by Theorem1.8.


\entrylabel{Second check}

Jordan normal form shows M\textasciicircum{}n->0 when its spectral radius is below1,\allowbreak{} including reducible matrices.


\entrylabel{Remaining obstruction}

At spectral radius1 the first moment need not decay,\allowbreak{} and reducible type classes can behave differently. The established subcritical argument does not settle critical/\allowbreak{}supercritical interpretation.


Verification:\allowbreak{} author checked limited argument | Independent review:\allowbreak{} none | Full target solved:\allowbreak{} False


\clearpage

\section*{P0091  Infinite-state tree-automaton interpretation criteria}

\entrylabel{Primary source}

Random tree recursions:\allowbreak{} which fixed points correspond to tangible sets of trees?;\allowbreak{} Section6,\allowbreak{} PDF pp.39-40,\allowbreak{} second paragraph


\url{https://arxiv.org/pdf/1808.03019v2}


Source statement date:\allowbreak{} not established. Archived access:\allowbreak{} 2026-09-06T16:\allowbreak{}42:\allowbreak{}18Z.


Historical:\allowbreak{} source\_\allowbreak{}says\_\allowbreak{}open | Current:\allowbreak{} unverified | Result:\allowbreak{} attempted\_\allowbreak{}unsolved


\entrylabel{Question}

Extend the random-tree interpretation and pivot-process theory to infinite state spaces,\allowbreak{} so it can handle recursive distributional models beyond finite colors.


\entrylabel{Scope}

The source asks for an extension framework. Finite-type spectral statements cannot be copied unchanged to a countable type space.


\subsection*{P0091-A001 | Infinite-state tree-automaton interpretation criteria structural argument}

\entrylabel{Argument or result}

Consider a countable-type branching process in which each type k individual has deterministically one child of type k+\allowbreak{}1. It survives forever,\allowbreak{} although every finite truncation of its mean matrix is nilpotent and every fixed type eventually disappears.


\entrylabel{Second check}

The population is exactly one at every level. The top-left K by K shift matrix has Kth power zero.


\entrylabel{Remaining obstruction}

Finite truncations falsely suggest extinction because population escapes through types. Any infinite-state criterion needs tightness or a suitable operator topology,\allowbreak{} not only finite truncation eigenvalues.


Verification:\allowbreak{} author checked limited argument | Independent review:\allowbreak{} none | Full target solved:\allowbreak{} False


\clearpage

\section*{P0092  Randomized tree automata and endogeny}

\entrylabel{Primary source}

Random tree recursions:\allowbreak{} which fixed points correspond to tangible sets of trees?;\allowbreak{} Section6,\allowbreak{} PDF pp.39-40,\allowbreak{} second paragraph


\url{https://arxiv.org/pdf/1808.03019v2}


Source statement date:\allowbreak{} not established. Archived access:\allowbreak{} 2026-09-06T16:\allowbreak{}42:\allowbreak{}18Z.


Historical:\allowbreak{} source\_\allowbreak{}says\_\allowbreak{}open | Current:\allowbreak{} unverified | Result:\allowbreak{} attempted\_\allowbreak{}unsolved


\entrylabel{Question}

Extend the interpretation/\allowbreak{}pivot theory when each vertex has iid auxiliary randomness and its state is a function of its child states and that local randomness.


\entrylabel{Scope}

Endogeny must be relative to the tree plus all vertex innovations;\allowbreak{} it is a different sigma-field from the bare uncolored tree.


\subsection*{P0092-A001 | Randomized tree automata and endogeny structural argument}

\entrylabel{Argument or result}

Let each vertex state be its own independent Bernoulli(p) innovation,\allowbreak{} ignoring its children. The pivot tree is empty and the state is measurable from the marked tree,\allowbreak{} but for 0<p<1 it is independent of the unmarked tree and cannot be a measurable function of that tree.


\entrylabel{Second check}

Conditional variance given the bare tree is p(1-p)>0;\allowbreak{} conditional variance given innovations is0.


\entrylabel{Remaining obstruction}

A literal extension retaining only the unmarked-tree sigma-field is false. A useful theorem must specify the innovation sigma-field and then handle nontrivial child dependence.


Verification:\allowbreak{} author checked limited argument | Independent review:\allowbreak{} none | Full target solved:\allowbreak{} False


\clearpage

\section*{P0093  Infinitely many fixed points of a finite-color automaton}

\entrylabel{Primary source}

Random tree recursions:\allowbreak{} which fixed points correspond to tangible sets of trees?;\allowbreak{} Section6,\allowbreak{} PDF p.40,\allowbreak{} third paragraph


\url{https://arxiv.org/pdf/1808.03019v2}


Source statement date:\allowbreak{} not established. Archived access:\allowbreak{} 2026-09-06T16:\allowbreak{}42:\allowbreak{}18Z.


Historical:\allowbreak{} source\_\allowbreak{}says\_\allowbreak{}open | Current:\allowbreak{} unverified | Result:\allowbreak{} attempted\_\allowbreak{}unsolved


\entrylabel{Question}

For a finite color set of size at least3 and a deterministic child-count automaton under the paper standing offspring assumptions,\allowbreak{} can the distributional map have infinitely many fixed points?


\entrylabel{Scope}

The one-child identity example is excluded by the genuine-branching assumption. Empty-color embeddings are not used to manufacture examples.


\subsection*{P0093-A001 | Infinitely many fixed points of a finite-color automaton structural argument}

\entrylabel{Argument or result}

If the offspring distribution has finite support bounded by D,\allowbreak{} each coordinate of the distributional map is a polynomial of degree at most D in the color probabilities. Fixed points are therefore a compact real algebraic subset of the probability simplex. In contrast with one variable,\allowbreak{} polynomial equations may have curves of zeros,\allowbreak{} so analyticity alone cannot prove finiteness.


\entrylabel{Second check}

For example,\allowbreak{} the algebraic system x-y=\allowbreak{}0 in a square has a whole line of zeros;\allowbreak{} this illustrates the logical gap without asserting that this system comes from an automaton.


\entrylabel{Remaining obstruction}

The nonnegative multinomial coefficients and deterministic count assignments constrain which algebraic sets can occur. Neither an admissible positive-dimensional example nor a general exclusion is obtained.


Verification:\allowbreak{} author checked limited argument | Independent review:\allowbreak{} none | Full target solved:\allowbreak{} False


\clearpage

\section*{P0094  Prescribing rogue and interpretable automaton fixed points}

\entrylabel{Primary source}

Random tree recursions:\allowbreak{} which fixed points correspond to tangible sets of trees?;\allowbreak{} Section6,\allowbreak{} PDF p.40,\allowbreak{} third paragraph


\url{https://arxiv.org/pdf/1808.03019v2}


Source statement date:\allowbreak{} not established. Archived access:\allowbreak{} 2026-09-06T16:\allowbreak{}42:\allowbreak{}18Z.


Historical:\allowbreak{} source\_\allowbreak{}says\_\allowbreak{}open | Current:\allowbreak{} unverified | Result:\allowbreak{} attempted\_\allowbreak{}unsolved


\entrylabel{Question}

Given a finite prescribed configuration of two-state fixed points with rogue/\allowbreak{}interpretable labels,\allowbreak{} determine whether a deterministic tree automaton and offspring distribution can realize it,\allowbreak{} including the monotone-automaton restriction.


\entrylabel{Scope}

Endpoint and multiplicity constraints are part of realizability. The request is not permission to assign arbitrary labels that contradict the pivot criterion.


\subsection*{P0094-A001 | Prescribing rogue and interpretable automaton fixed points structural argument}

\entrylabel{Argument or result}

For a differentiable increasing distributional map Psi,\allowbreak{} let g(p)=\allowbreak{}Psi(p)-p. At an isolated simple fixed point,\allowbreak{} crossing from g>0 to g<0 gives Psi prime<1,\allowbreak{} while the opposite crossing gives Psi prime>1. Consecutive simple interior crossings must alternate signs. Thus a prescribed list of simple attractive/\allowbreak{}repulsive crossings cannot have two consecutive crossings of the same orientation without an intervening zero.


\entrylabel{Second check}

This is the intermediate value theorem and the derivative sign at a simple root;\allowbreak{} it makes no claim for tangencies or nondifferentiable maps.


\entrylabel{Remaining obstruction}

The derivative restriction is necessary,\allowbreak{} not sufficient for an automaton realization or rogue classification in degenerate cases. Constructing compatible multinomial maps remains open here.


Verification:\allowbreak{} author checked limited argument | Independent review:\allowbreak{} none | Full target solved:\allowbreak{} False


\clearpage

\section*{P0095  Automaton-level conditions for a unique fixed point}

\entrylabel{Primary source}

Random tree recursions:\allowbreak{} which fixed points correspond to tangible sets of trees?;\allowbreak{} Section6,\allowbreak{} PDF p.40,\allowbreak{} fourth paragraph


\url{https://arxiv.org/pdf/1808.03019v2}


Source statement date:\allowbreak{} not established. Archived access:\allowbreak{} 2026-09-06T16:\allowbreak{}42:\allowbreak{}18Z.


Historical:\allowbreak{} source\_\allowbreak{}says\_\allowbreak{}open | Current:\allowbreak{} unverified | Result:\allowbreak{} attempted\_\allowbreak{}unsolved


\entrylabel{Question}

Find structural conditions on a tree automaton itself that force uniqueness of its distributional fixed point,\allowbreak{} complementing the source condition on a first-order interpretation.


\entrylabel{Scope}

The condition must involve the offspring law as appropriate;\allowbreak{} the same rule can undergo bifurcations when that law changes.


\subsection*{P0095-A001 | Automaton-level conditions for a unique fixed point structural argument}

\entrylabel{Argument or result}

Suppose that for each offspring count D the rule depends on a fixed set of at most k(D) child positions,\allowbreak{} and E k(D)=\allowbreak{}c<1. Couple child colors optimally with disagreement probability dTV(mu,\allowbreak{}nu);\allowbreak{} the union bound over that fixed relevant set gives dTV(Psi(mu),\allowbreak{}Psi(nu))<=\allowbreak{}c dTV(mu,\allowbreak{}nu). Iteration is a contraction and yields a unique fixed point on the finite simplex.


\entrylabel{Second check}

The relevant set must be fixed independently of the child configuration;\allowbreak{} a pointwise bound on pivotal children alone would not justify this global coupling. For symmetric count automata the crude k(D)=\allowbreak{}D criterion reduces to E D<1.


\entrylabel{Remaining obstruction}

For supercritical offspring the crude contraction constant exceeds1. A useful structural criterion must exploit much finer properties than the number of children.


Verification:\allowbreak{} author checked limited argument | Independent review:\allowbreak{} none | Full target solved:\allowbreak{} False


\clearpage

\section*{P0097  An explicit nontrivial FPP time constant and limit shape}

\entrylabel{Primary source}

Auffinger,\allowbreak{} Damron and Hanson:\allowbreak{} 50 years of first-passage percolation;\allowbreak{} Questions1 and37,\allowbreak{} PDF pp.11 and141


\url{https://arxiv.org/pdf/1511.03262}


Source statement date:\allowbreak{} 2016-09-26. Archived access:\allowbreak{} 2026-09-06T18:\allowbreak{}34:\allowbreak{}26Z.


Historical:\allowbreak{} source\_\allowbreak{}says\_\allowbreak{}open | Current:\allowbreak{} unverified | Result:\allowbreak{} attempted\_\allowbreak{}unsolved


\entrylabel{Question}

Identify an explicit positive axial time constant for a nondegenerate iid nearest-neighbor passage-time law on Z\textasciicircum{}d,\allowbreak{} or explicitly determine a nontrivial compact limit shape.


\entrylabel{Scope}

The survey itself calls Question37 a rephrasing of Question1;\allowbreak{} only one problem is counted. A zero time constant from a supercritical zero-edge atom or the deterministic l1 ball is excluded.


\subsection*{P0097-A001 | An explicit nontrivial FPP time constant and limit shape structural argument}

\entrylabel{Argument or result}

If every edge time is multiplied by c>0,\allowbreak{} every finite path cost and every passage time is multiplied by c,\allowbreak{} so the time constant is multiplied by c and the limit shape is dilated by1/\allowbreak{}c. Thus scaling a known degenerate example cannot produce the missing nontrivial shape.


\entrylabel{Second check}

The equality holds configuration by configuration before taking any limit. Positive-time and nondegeneracy restrictions exclude the zero-time percolation examples.


\entrylabel{Remaining obstruction}

Homogeneity supplies a transformation law,\allowbreak{} not a computable nondegenerate base model. Integrable directed last-passage formulas apply to a different optimization problem.


Verification:\allowbreak{} author checked limited argument | Independent review:\allowbreak{} none | Full target solved:\allowbreak{} False


\clearpage

\section*{P0098  Strict FPP comparison in higher dimensions at a large minimum atom}

\entrylabel{Primary source}

Auffinger,\allowbreak{} Damron and Hanson:\allowbreak{} 50 years of first-passage percolation;\allowbreak{} Question3,\allowbreak{} PDF p.13;\allowbreak{} Definition2.8 and Theorem2.12,\allowbreak{} pp.12-13


\url{https://arxiv.org/pdf/1511.03262}


Source statement date:\allowbreak{} 2016-09-26. Archived access:\allowbreak{} 2026-09-06T18:\allowbreak{}34:\allowbreak{}26Z.


Historical:\allowbreak{} source\_\allowbreak{}says\_\allowbreak{}open | Current:\allowbreak{} unverified | Result:\allowbreak{} attempted\_\allowbreak{}unsolved


\entrylabel{Question}

Extend strict comparison of axial time constants for strictly ordered passage-time laws in increasing-concave order to d>2 when the minimum support value r>0 has mass at least the oriented-percolation threshold.


\entrylabel{Scope}

More variable means integral phi dF <=\allowbreak{} integral phi dG for every increasing concave phi. The claim is axial;\allowbreak{} equality inside an oriented percolation cone cannot be ignored in statements about arbitrary directions.


\subsection*{P0098-A001 | Strict FPP comparison in higher dimensions at a large minimum atom structural argument}

\entrylabel{Argument or result}

For the special coupling tau'\_\allowbreak{}e=\allowbreak{}tau\_\allowbreak{}e+\allowbreak{}delta with delta>0,\allowbreak{} every path from0 to ne1 has at least n edges. Its primed cost is at least its original cost plus delta n. Taking infima and limits gives mu'(e1)>=\allowbreak{}mu(e1)+\allowbreak{}delta.


\entrylabel{Second check}

The bound does not depend on a minimizing path or an upper bound on its length. It proves an actual strict gap for this special stochastic ordering.


\entrylabel{Remaining obstruction}

General increasing-concave order,\allowbreak{} including mean-preserving spreads,\allowbreak{} has no pointwise additive coupling. The argument does not handle the remaining comparison class.


Verification:\allowbreak{} author checked limited argument | Independent review:\allowbreak{} none | Full target solved:\allowbreak{} False


\clearpage

\section*{P0099  Strict convexity of iid FPP limit shapes}

\entrylabel{Primary source}

Auffinger,\allowbreak{} Damron and Hanson:\allowbreak{} 50 years of first-passage percolation;\allowbreak{} Question4,\allowbreak{} PDF p.21;\allowbreak{} Section2.8


\url{https://arxiv.org/pdf/1511.03262}


Source statement date:\allowbreak{} 2016-09-26. Archived access:\allowbreak{} 2026-09-06T18:\allowbreak{}34:\allowbreak{}26Z.


Historical:\allowbreak{} source\_\allowbreak{}says\_\allowbreak{}open | Current:\allowbreak{} unverified | Result:\allowbreak{} attempted\_\allowbreak{}unsolved


\entrylabel{Question}

For continuous iid passage-time distributions satisfying the shape-theorem assumptions,\allowbreak{} prove that the deterministic limit shape is strictly convex.


\entrylabel{Scope}

This qualitative absence of boundary segments is distinct from a quantitative uniform-curvature estimate. Continuity means no atoms;\allowbreak{} finite moment and positive-speed assumptions must accompany a compact shape.


\subsection*{P0099-A001 | Strict convexity of iid FPP limit shapes structural argument}

\entrylabel{Argument or result}

Finite geodesic uniqueness does not imply strict convexity of an asymptotic norm. For example,\allowbreak{} the strictly convex norms ||x||\_\allowbreak{}p converge uniformly on compact sets to ||x||\_\allowbreak{}infinity as p tends to infinity,\allowbreak{} and the latter unit ball has flat faces. Even stronger finite approximations can lose strictness in the limit.


\entrylabel{Second check}

The inequality ||x||\_\allowbreak{}infinity <=\allowbreak{} ||x||\_\allowbreak{}p <=\allowbreak{} d\textasciicircum{}(1/\allowbreak{}p)||x||\_\allowbreak{}infinity verifies the convergence directly.


\entrylabel{Remaining obstruction}

This is a counterexample to an inference about limits,\allowbreak{} not an iid FPP counterexample. A uniform quantitative gap in the triangle inequality is missing.


Verification:\allowbreak{} author checked limited argument | Independent review:\allowbreak{} none | Full target solved:\allowbreak{} False


\clearpage

\section*{P0100  Excluding axial facets and cubical FPP shapes}

\entrylabel{Primary source}

Auffinger,\allowbreak{} Damron and Hanson:\allowbreak{} 50 years of first-passage percolation;\allowbreak{} Questions5,\allowbreak{}8,\allowbreak{}9,\allowbreak{} PDF pp.21 and24


\url{https://arxiv.org/pdf/1511.03262}


Source statement date:\allowbreak{} 2016-09-26. Archived access:\allowbreak{} 2026-09-06T18:\allowbreak{}34:\allowbreak{}26Z.


Historical:\allowbreak{} source\_\allowbreak{}says\_\allowbreak{}open | Current:\allowbreak{} unverified | Result:\allowbreak{} attempted\_\allowbreak{}unsolved


\entrylabel{Question}

Rule out a cube as an iid FPP limit shape and,\allowbreak{} more strongly,\allowbreak{} rule out a boundary segment perpendicular to e1 at the axial boundary point;\allowbreak{} exhibit a nondegenerate law with no such axial segment.


\entrylabel{Scope}

These nested axial-flatness formulations are grouped. Convexity,\allowbreak{} coordinate-reflection symmetry,\allowbreak{} and a nonempty interior are known but do not settle them.


\subsection*{P0100-A001 | Excluding axial facets and cubical FPP shapes structural argument}

\entrylabel{Argument or result}

The l\_\allowbreak{}infinity unit ball satisfies every elementary symmetry and convexity condition of an FPP limit shape and has axial facets. The l1 unit ball also satisfies those conditions but has an axial corner. Thus the shape theorem's convexity and lattice symmetries permit both geometries.


\entrylabel{Second check}

Both balls are compact with nonempty interior and invariant under coordinate permutations and sign changes. Their intersections with the supporting plane x1=\allowbreak{}1 distinguish the facet from the corner.


\entrylabel{Remaining obstruction}

One needs a property specific to iid optimization that excludes the cubical candidate;\allowbreak{} deterministic convex geometry alone supplies none.


Verification:\allowbreak{} author checked limited argument | Independent review:\allowbreak{} none | Full target solved:\allowbreak{} False


\clearpage

\section*{P0101  Classification of flat FPP boundary pieces}

\entrylabel{Primary source}

Auffinger,\allowbreak{} Damron and Hanson:\allowbreak{} 50 years of first-passage percolation;\allowbreak{} Questions6-7,\allowbreak{} PDF p.24;\allowbreak{} definition of M\_\allowbreak{}p,\allowbreak{} p.22


\url{https://arxiv.org/pdf/1511.03262}


Source statement date:\allowbreak{} 2016-09-26. Archived access:\allowbreak{} 2026-09-06T18:\allowbreak{}34:\allowbreak{}26Z.


Historical:\allowbreak{} source\_\allowbreak{}says\_\allowbreak{}open | Current:\allowbreak{} unverified | Result:\allowbreak{} attempted\_\allowbreak{}unsolved


\entrylabel{Question}

For planar iid FPP,\allowbreak{} show that the only flat boundary portions are the oriented-percolation-cone portions associated,\allowbreak{} after rescaling,\allowbreak{} with a positive minimum edge weight carrying a sufficiently large atom.


\entrylabel{Scope}

M\_\allowbreak{}p in the source normalizes the minimum support value to1 and its atom to p>=\allowbreak{}oriented p\_\allowbreak{}c. The two source questions concerning flats outside the cone and converse classification are grouped.


\subsection*{P0101-A001 | Classification of flat FPP boundary pieces structural argument}

\entrylabel{Argument or result}

If tau\_\allowbreak{}e>=\allowbreak{}1,\allowbreak{} then every path from0 to x costs at least ||x||\_\allowbreak{}1,\allowbreak{} so the limit shape lies inside the l1 ball. An oriented path consisting entirely of unit-time edges attains this lower bound exactly. This explains the mechanism generating the known cone flats.


\entrylabel{Second check}

The lower bound is pathwise. Along an oriented path,\allowbreak{} its length equals the l1 displacement,\allowbreak{} so equality is exact.


\entrylabel{Remaining obstruction}

The mechanism proves why these flats occur,\allowbreak{} not why no alternative mechanism creates a flat segment away from that cone.


Verification:\allowbreak{} author checked limited argument | Independent review:\allowbreak{} none | Full target solved:\allowbreak{} False


\clearpage

\section*{P0102  Strong Busemann limits for subadditive arrays and FPP}

\entrylabel{Primary source}

Auffinger,\allowbreak{} Damron and Hanson:\allowbreak{} 50 years of first-passage percolation;\allowbreak{} Questions10 and32,\allowbreak{} PDF pp.26 and113


\url{https://arxiv.org/pdf/1511.03262}


Source statement date:\allowbreak{} 2016-09-26. Archived access:\allowbreak{} 2026-09-06T18:\allowbreak{}34:\allowbreak{}26Z.


Historical:\allowbreak{} source\_\allowbreak{}says\_\allowbreak{}open | Current:\allowbreak{} unverified | Result:\allowbreak{} attempted\_\allowbreak{}unsolved


\entrylabel{Question}

Find useful conditions ensuring almost-sure limits of X\_\allowbreak{}(m,\allowbreak{}N)-X\_\allowbreak{}(n,\allowbreak{}N),\allowbreak{} including FPP differences of passage times to receding hyperplanes.


\entrylabel{Scope}

The survey explicitly translates the FPP limit question into the general-array question. These are one family;\allowbreak{} a weak subsequential or averaged limit is not the requested pathwise limit.


\subsection*{P0102-A001 | Strong Busemann limits for subadditive arrays and FPP structural argument}

\entrylabel{Argument or result}

Let S be simple symmetric random walk and X\_\allowbreak{}(m,\allowbreak{}n)=\allowbreak{}|S\_\allowbreak{}n-S\_\allowbreak{}m|. This is an integrable stationary ergodic subadditive array. For N with S\_\allowbreak{}N>=\allowbreak{}2,\allowbreak{} X\_\allowbreak{}(0,\allowbreak{}N)-X\_\allowbreak{}(1,\allowbreak{}N)=\allowbreak{}S\_\allowbreak{}1;\allowbreak{} for S\_\allowbreak{}N<=\allowbreak{}-2 it equals -S\_\allowbreak{}1. Both events occur infinitely often almost surely,\allowbreak{} so this Busemann difference does not converge.


\entrylabel{Second check}

Subadditivity is the triangle inequality;\allowbreak{} stationarity and ergodicity come from the iid increments. Recurrence in both directions proves the two subsequential values,\allowbreak{} which differ because S\_\allowbreak{}1 is+\allowbreak{}1 or-1.


\entrylabel{Remaining obstruction}

Kingman's usual hypotheses are insufficient. FPP triangle bounds yield compact subsequences,\allowbreak{} but an additional uniqueness or directional structure is needed for a full limit.


Verification:\allowbreak{} author checked limited argument | Independent review:\allowbreak{} none | Full target solved:\allowbreak{} False


\clearpage

\section*{P0103  Uniform curvature of continuous-law FPP shapes}

\entrylabel{Primary source}

Auffinger,\allowbreak{} Damron and Hanson:\allowbreak{} 50 years of first-passage percolation;\allowbreak{} Question11,\allowbreak{} PDF p.31


\url{https://arxiv.org/pdf/1511.03262}


Source statement date:\allowbreak{} 2016-09-26. Archived access:\allowbreak{} 2026-09-06T18:\allowbreak{}34:\allowbreak{}26Z.


Historical:\allowbreak{} source\_\allowbreak{}says\_\allowbreak{}open | Current:\allowbreak{} unverified | Result:\allowbreak{} attempted\_\allowbreak{}unsolved


\entrylabel{Question}

Prove that every continuous iid passage-time law with a compact limit shape has a uniform constant C>0 such that 1-mu(z)>=\allowbreak{}C min(mu(z-z1),\allowbreak{}mu(z-z2))\textasciicircum{}2 for boundary points z1,\allowbreak{}z2 and z=\allowbreak{}(1-lambda)z1+\allowbreak{}lambda z2.


\entrylabel{Scope}

This is the source's geometric definition,\allowbreak{} avoiding ambiguity in printed curvature-exponent conventions. It is quantitatively stronger than strict convexity.


\subsection*{P0103-A001 | Uniform curvature of continuous-law FPP shapes structural argument}

\entrylabel{Argument or result}

Strict convexity alone cannot imply this inequality. On the l\_\allowbreak{}p unit ball for p>2,\allowbreak{} take z1=\allowbreak{}((1-t\textasciicircum{}p)\textasciicircum{}(1/\allowbreak{}p),\allowbreak{}t),\allowbreak{} z2=\allowbreak{}((1-t\textasciicircum{}p)\textasciicircum{}(1/\allowbreak{}p),\allowbreak{}-t),\allowbreak{} and their midpoint z. Then 1-||z||\_\allowbreak{}p is asymptotic to t\textasciicircum{}p/\allowbreak{}p,\allowbreak{} whereas ||z-z\_\allowbreak{}i||\_\allowbreak{}p\textasciicircum{}2=\allowbreak{}t\textasciicircum{}2. Their ratio tends to0.


\entrylabel{Second check}

The expansion (1-u)\textasciicircum{}(1/\allowbreak{}p)=\allowbreak{}1-u/\allowbreak{}p+\allowbreak{}o(u) gives the asserted ratio. The l\_\allowbreak{}p ball is strictly convex for p>1.


\entrylabel{Remaining obstruction}

This is a norm counterexample to a proof strategy,\allowbreak{} not an FPP counterexample. Uniform lower curvature must be proved from the random model.


Verification:\allowbreak{} author checked limited argument | Independent review:\allowbreak{} none | Full target solved:\allowbreak{} False


\clearpage

\section*{P0104  Subdiffusive and KPZ passage-time fluctuation exponents}

\entrylabel{Primary source}

Auffinger,\allowbreak{} Damron and Hanson:\allowbreak{} 50 years of first-passage percolation;\allowbreak{} Questions13 and15,\allowbreak{} PDF pp.34 and40;\allowbreak{} Section3


\url{https://arxiv.org/pdf/1511.03262}


Source statement date:\allowbreak{} 2016-09-26. Archived access:\allowbreak{} 2026-09-06T18:\allowbreak{}34:\allowbreak{}26Z.


Historical:\allowbreak{} source\_\allowbreak{}says\_\allowbreak{}open | Current:\allowbreak{} unverified | Result:\allowbreak{} attempted\_\allowbreak{}unsolved


\entrylabel{Question}

Under suitable nondegenerate iid passage-time assumptions,\allowbreak{} prove a fluctuation exponent strictly below1/\allowbreak{}2,\allowbreak{} identify the planar exponent as1/\allowbreak{}3,\allowbreak{} and establish a positive exponent in a suitable dimension.


\entrylabel{Scope}

These closely related exponent-range targets are grouped. Different definitions using variances or stronger concentration must be specified before claiming equivalence.


\subsection*{P0104-A001 | Subdiffusive and KPZ passage-time fluctuation exponents structural argument}

\entrylabel{Argument or result}

A variance upper bound of order n/\allowbreak{}log n does not establish an exponent below1/\allowbreak{}2. Indeed log(n/\allowbreak{}log n)/\allowbreak{}(2log n) tends to1/\allowbreak{}2. Sums of floor(n/\allowbreak{}log n) independent mean-zero variance-one variables realize that variance order in an elementary model.


\entrylabel{Second check}

The logarithm calculation is exact in the limit;\allowbreak{} the example has additive variances. It is not asserted to be FPP.


\entrylabel{Remaining obstruction}

Known sublinear-variance estimates leave a polynomial-exponent gap. An additional power saving,\allowbreak{} and a matching lower bound for identification,\allowbreak{} are required.


Verification:\allowbreak{} author checked limited argument | Independent review:\allowbreak{} none | Full target solved:\allowbreak{} False


\clearpage

\section*{P0105  High-dimensional limit of the FPP fluctuation exponent}

\entrylabel{Primary source}

Auffinger,\allowbreak{} Damron and Hanson:\allowbreak{} 50 years of first-passage percolation;\allowbreak{} Question14,\allowbreak{} PDF p.34


\url{https://arxiv.org/pdf/1511.03262}


Source statement date:\allowbreak{} 2016-09-26. Archived access:\allowbreak{} 2026-09-06T18:\allowbreak{}34:\allowbreak{}26Z.


Historical:\allowbreak{} source\_\allowbreak{}says\_\allowbreak{}open | Current:\allowbreak{} unverified | Result:\allowbreak{} attempted\_\allowbreak{}unsolved


\entrylabel{Question}

Determine whether the dimension-dependent passage-time fluctuation exponent chi(d) tends to zero as d tends to infinity.


\entrylabel{Scope}

Existence and definition of chi(d) are themselves part of the scope;\allowbreak{} decay of the time constant is not decay of its fluctuation exponent.


\subsection*{P0105-A001 | High-dimensional limit of the FPP fluctuation exponent structural argument}

\entrylabel{Argument or result}

Embedding Z\textasciicircum{}d into Z\textasciicircum{}(d+\allowbreak{}1) with shared weights gives T\_\allowbreak{}(d+\allowbreak{}1)<=\allowbreak{}T\_\allowbreak{}d and hence mu\_\allowbreak{}(d+\allowbreak{}1)<=\allowbreak{}mu\_\allowbreak{}d. This ordering says nothing comparable about variance:\allowbreak{} a Bernoulli random variable is bounded above by the constant1 but has larger variance.


\entrylabel{Second check}

The passage-time inequality follows because the higher-dimensional infimum has more admissible paths. The variance counterexample has values1/\allowbreak{}4 and0 at Bernoulli parameter1/\allowbreak{}2.


\entrylabel{Remaining obstruction}

A distributional comparison controlling centered fluctuations across dimensions is missing. Monotonicity of means cannot supply it.


Verification:\allowbreak{} author checked limited argument | Independent review:\allowbreak{} none | Full target solved:\allowbreak{} False


\clearpage

\section*{P0106  Polynomial overlap of independent planar FPP geodesics}

\entrylabel{Primary source}

Auffinger,\allowbreak{} Damron and Hanson:\allowbreak{} 50 years of first-passage percolation;\allowbreak{} Question16,\allowbreak{} PDF p.42


\url{https://arxiv.org/pdf/1511.03262}


Source statement date:\allowbreak{} 2016-09-26. Archived access:\allowbreak{} 2026-09-06T18:\allowbreak{}34:\allowbreak{}26Z.


Historical:\allowbreak{} source\_\allowbreak{}says\_\allowbreak{}open | Current:\allowbreak{} unverified | Result:\allowbreak{} attempted\_\allowbreak{}unsolved


\entrylabel{Question}

Show that two independent geodesics from0 to ne1,\allowbreak{} sampled in independent iid planar environments with the same law,\allowbreak{} have expected edge overlap at least n\textasciicircum{}a for some a>0.


\entrylabel{Scope}

Independent environments and a fixed measurable geodesic-selection rule are intended;\allowbreak{} two samples in one unique-geodesic environment would be identical.


\subsection*{P0106-A001 | Polynomial overlap of independent planar FPP geodesics structural argument}

\entrylabel{Argument or result}

Let E\_\allowbreak{}r be the 8r+\allowbreak{}4 edges leaving the square[-r,\allowbreak{}r]\textasciicircum{}2,\allowbreak{} for r=\allowbreak{}0,\allowbreak{}...,\allowbreak{}n-1,\allowbreak{} and p\_\allowbreak{}e=\allowbreak{}P(e belongs to the selected geodesic). Every path to ne1 crosses E\_\allowbreak{}r,\allowbreak{} so sum\_\allowbreak{}(e in E\_\allowbreak{}r)p\_\allowbreak{}e>=\allowbreak{}1. Independence and Cauchy-Schwarz give E|Gamma intersect Gamma'|=\allowbreak{}sum\_\allowbreak{}e p\_\allowbreak{}e\textasciicircum{}2 >=\allowbreak{} sum\_\allowbreak{}(r=\allowbreak{}0)\textasciicircum{}(n-1)1/\allowbreak{}(8r+\allowbreak{}4) >=\allowbreak{}(1/\allowbreak{}8)log(2n+\allowbreak{}1).


\entrylabel{Second check}

The sets E\_\allowbreak{}r are disjoint;\allowbreak{} corners contribute two outward edges and total8r+\allowbreak{}4 is exact. Tonelli justifies the infinite sum. Integrating1/\allowbreak{}(8x+\allowbreak{}4) gives the final lower bound.


\entrylabel{Remaining obstruction}

This establishes only a logarithmic overlap lower bound. Improving shell-level spreading to a power requires model-specific concentration of geodesic locations.


Verification:\allowbreak{} author checked limited argument | Independent review:\allowbreak{} none | Full target solved:\allowbreak{} False


\clearpage

\section*{P0107  Critical Bernoulli FPP growth in three dimensions}

\entrylabel{Primary source}

Auffinger,\allowbreak{} Damron and Hanson:\allowbreak{} 50 years of first-passage percolation;\allowbreak{} Question19,\allowbreak{} PDF p.71


\url{https://arxiv.org/pdf/1511.03262}


Source statement date:\allowbreak{} 2016-09-26. Archived access:\allowbreak{} 2026-09-06T18:\allowbreak{}34:\allowbreak{}26Z.


Historical:\allowbreak{} source\_\allowbreak{}says\_\allowbreak{}open | Current:\allowbreak{} unverified | Result:\allowbreak{} attempted\_\allowbreak{}unsolved


\entrylabel{Question}

For iid edge times0 with probability p\_\allowbreak{}c(Z\textasciicircum{}3) and1 otherwise,\allowbreak{} determine whether E T(0,\allowbreak{}boundary B(n)) is comparable to log n.


\entrylabel{Scope}

The target is undirected edge FPP at the three-dimensional bond threshold;\allowbreak{} planar critical exponents cannot be transferred.


\subsection*{P0107-A001 | Critical Bernoulli FPP growth in three dimensions structural argument}

\entrylabel{Argument or result}

Collapse every connected zero-edge cluster to a vertex and retain unit-cost edges between the resulting vertices. Passage time is precisely graph distance in this random quotient graph,\allowbreak{} counting the fewest positive edges needed to reach the target. In the original graph it is bounded above by the number of positive edges on a fixed n-step coordinate path.


\entrylabel{Second check}

Every original path projects to a quotient path with the same number of unit-cost crossings,\allowbreak{} and every finite quotient path lifts through finite zero-edge connections. This proves equality even when cluster sizes are unbounded.


\entrylabel{Remaining obstruction}

The logarithmic estimate requires quantitative geometry of the cluster-adjacency quotient. The planar argument based on nested dual circuits has no direct three-dimensional counterpart.


Verification:\allowbreak{} author checked limited argument | Independent review:\allowbreak{} none | Full target solved:\allowbreak{} False


\clearpage

\section*{P0108  Sharp width threshold for the thin-cylinder FPP CLT}

\entrylabel{Primary source}

Auffinger,\allowbreak{} Damron and Hanson:\allowbreak{} 50 years of first-passage percolation;\allowbreak{} Question20,\allowbreak{} PDF p.72


\url{https://arxiv.org/pdf/1511.03262}


Source statement date:\allowbreak{} 2016-09-26. Archived access:\allowbreak{} 2026-09-06T18:\allowbreak{}34:\allowbreak{}26Z.


Historical:\allowbreak{} source\_\allowbreak{}says\_\allowbreak{}open | Current:\allowbreak{} unverified | Result:\allowbreak{} attempted\_\allowbreak{}unsolved


\entrylabel{Question}

Determine the largest power alpha for which standardized passage times constrained to Z times[-n\textasciicircum{}alpha,\allowbreak{}n\textasciicircum{}alpha]\textasciicircum{}(d-1) converge to a standard normal law.


\entrylabel{Scope}

The source defines gamma\_\allowbreak{}F(d) by this threshold under a finite2+\allowbreak{}delta edge-time moment. A sufficient exponent from the source theorem is not an optimality result.


\subsection*{P0108-A001 | Sharp width threshold for the thin-cylinder FPP CLT structural argument}

\entrylabel{Argument or result}

At width0 the graph is a line,\allowbreak{} and the passage time to ne1 is exactly the sum of n iid edge weights. For a nondegenerate finite-variance law the classical iid CLT gives normalization sqrt(n Var(tau)) and mean n E(tau).


\entrylabel{Second check}

Any line path containing a backtracking excursion can have that excursion removed without increasing its nonnegative cost,\allowbreak{} leaving the unique coordinate segment.


\entrylabel{Remaining obstruction}

For growing width,\allowbreak{} competing paths introduce correlations and varying block sizes. The width0 reduction does not control how rapidly an approximately independent block decomposition deteriorates.


Verification:\allowbreak{} author checked limited argument | Independent review:\allowbreak{} none | Full target solved:\allowbreak{} False


\clearpage

\section*{P0109  Existence of finite FPP geodesics without distributional assumptions}

\entrylabel{Primary source}

Auffinger,\allowbreak{} Damron and Hanson:\allowbreak{} 50 years of first-passage percolation;\allowbreak{} Question21,\allowbreak{} PDF p.73


\url{https://arxiv.org/pdf/1511.03262}


Source statement date:\allowbreak{} 2016-09-26. Archived access:\allowbreak{} 2026-09-06T18:\allowbreak{}34:\allowbreak{}26Z.


Historical:\allowbreak{} source\_\allowbreak{}says\_\allowbreak{}open | Current:\allowbreak{} unverified | Result:\allowbreak{} attempted\_\allowbreak{}unsolved


\entrylabel{Question}

For every iid finite nonnegative passage-time distribution in d>=\allowbreak{}3,\allowbreak{} prove that the infimum over finite paths between any two vertices is attained.


\entrylabel{Scope}

Infinite passage-time edges are not included in this formulation. Non-iid weighted graphs can fail attainment and therefore cannot settle the iid question by analogy.


\subsection*{P0109-A001 | Existence of finite FPP geodesics without distributional assumptions structural argument}

\entrylabel{Argument or result}

If all edge times are at least a>0,\allowbreak{} a fixed connecting path of cost M bounds any competing path length by M/\allowbreak{}a;\allowbreak{} only finitely many simple paths of that length start at a fixed vertex,\allowbreak{} so a minimum exists. Without a positive lower bound,\allowbreak{} a deterministic ladder can fail:\allowbreak{} give each vertical edge from height j-1 to j on both rails weight2\textasciicircum{}(-j),\allowbreak{} and the rung at height k weight3*2\textasciicircum{}(-k). A route across at height k costs2+\allowbreak{}2\textasciicircum{}(-k),\allowbreak{} with infimum2 unattained.


\entrylabel{Second check}

For a finite ladder route attaining maximum height k,\allowbreak{} vertical cost is at least2(1-2\textasciicircum{}(-k)) and at least one rung has cost at least3*2\textasciicircum{}(-k);\allowbreak{} this gives the same strict lower bound. Embed the ladder into Z\textasciicircum{}d with all other edges cost100 to preserve failure.


\entrylabel{Remaining obstruction}

The ladder environment is deterministic and spatially inhomogeneous,\allowbreak{} so it is only a counterexample to a general compactness inference. The iid critical case requires additional probabilistic structure.


Verification:\allowbreak{} author checked limited argument | Independent review:\allowbreak{} none | Full target solved:\allowbreak{} False


\clearpage

\section*{P0110  An iid FPP law with proved KPZ exponent relation}

\entrylabel{Primary source}

Auffinger,\allowbreak{} Damron and Hanson:\allowbreak{} 50 years of first-passage percolation;\allowbreak{} Question22,\allowbreak{} PDF p.89;\allowbreak{} equation4.19


\url{https://arxiv.org/pdf/1511.03262}


Source statement date:\allowbreak{} 2016-09-26. Archived access:\allowbreak{} 2026-09-06T18:\allowbreak{}34:\allowbreak{}26Z.


Historical:\allowbreak{} source\_\allowbreak{}says\_\allowbreak{}open | Current:\allowbreak{} unverified | Result:\allowbreak{} attempted\_\allowbreak{}unsolved


\entrylabel{Question}

Exhibit a nondegenerate iid undirected FPP law with well-defined fluctuation and wandering exponents satisfying chi=\allowbreak{}2xi-1.


\entrylabel{Scope}

Conditional proofs assuming suitably strong exponents exist do not exhibit a law for which those assumptions hold. Exactly solvable directed last-passage models are outside the target.


\subsection*{P0110-A001 | An iid FPP law with proved KPZ exponent relation structural argument}

\entrylabel{Argument or result}

If a geodesic deviates transversely by n\textasciicircum{}xi and the norm has quadratic curvature,\allowbreak{} the deterministic excess norm over a displacement of length n has order n*(n\textasciicircum{}(xi-1))\textasciicircum{}2=\allowbreak{}n\textasciicircum{}(2xi-1). This yields the proposed balance only after assuming random fluctuations are of the same order.


\entrylabel{Second check}

The power follows from norm homogeneity and the square of the relative transverse displacement. A curvature exponent kappa would instead give kappa*xi-(kappa-1).


\entrylabel{Remaining obstruction}

Both the required exponent existence and error control when replacing random passage times by the norm remain unproved for a concrete iid law here. The balance calculation is conditional,\allowbreak{} not a proof.


Verification:\allowbreak{} author checked limited argument | Independent review:\allowbreak{} none | Full target solved:\allowbreak{} False


\clearpage

\section*{P0111  Infinitely many planar FPP geodesic rays for continuous laws}

\entrylabel{Primary source}

Auffinger,\allowbreak{} Damron and Hanson:\allowbreak{} 50 years of first-passage percolation;\allowbreak{} Questions23-24,\allowbreak{} PDF p.91


\url{https://arxiv.org/pdf/1511.03262}


Source statement date:\allowbreak{} 2016-09-26. Archived access:\allowbreak{} 2026-09-06T18:\allowbreak{}34:\allowbreak{}26Z.


Historical:\allowbreak{} source\_\allowbreak{}says\_\allowbreak{}open | Current:\allowbreak{} unverified | Result:\allowbreak{} attempted\_\allowbreak{}unsolved


\entrylabel{Question}

Prove that continuous iid planar passage-time distributions have infinitely many pairwise eventually disjoint geodesic rays;\allowbreak{} the source also asks for an example outside the percolation-cone class.


\entrylabel{Scope}

The universal continuous-law question and its weaker example target are grouped. Distinct rays mean sharing at most finitely many vertices,\allowbreak{} not merely having different first edges.


\subsection*{P0111-A001 | Infinitely many planar FPP geodesic rays for continuous laws structural argument}

\entrylabel{Argument or result}

Finite geodesic existence and finite lattice degree yield at least one ray:\allowbreak{} among geodesics to points tending to infinity,\allowbreak{} choose an infinite subsequence with a common first edge,\allowbreak{} then a common second edge,\allowbreak{} and continue. Every finite segment of the diagonal limit is minimizing.


\entrylabel{Second check}

At each stage finitely many choices ensure an infinite subsequence. Path lengths tend to infinity,\allowbreak{} and any shortcut of a limiting finite segment would contradict minimization in all sufficiently long selected paths.


\entrylabel{Remaining obstruction}

This compactness proof gives one ray,\allowbreak{} not many eventually disjoint rays. Endpoint sequences in different directions could still share the same limiting tail without additional shape information.


Verification:\allowbreak{} author checked limited argument | Independent review:\allowbreak{} none | Full target solved:\allowbreak{} False


\clearpage

\section*{P0112  Full-sequence convergence of FPP geodesics to a ray}

\entrylabel{Primary source}

Auffinger,\allowbreak{} Damron and Hanson:\allowbreak{} 50 years of first-passage percolation;\allowbreak{} Question25,\allowbreak{} PDF p.92


\url{https://arxiv.org/pdf/1511.03262}


Source statement date:\allowbreak{} 2016-09-26. Archived access:\allowbreak{} 2026-09-06T18:\allowbreak{}34:\allowbreak{}26Z.


Historical:\allowbreak{} source\_\allowbreak{}says\_\allowbreak{}open | Current:\allowbreak{} unverified | Result:\allowbreak{} attempted\_\allowbreak{}unsolved


\entrylabel{Question}

For iid passage-time laws with unique finite geodesics,\allowbreak{} determine whether the geodesics from0 to ne1 converge locally to a single infinite ray as n tends to infinity.


\entrylabel{Scope}

Convergence means stabilization of every finite initial segment. Existence of some subsequential ray is strictly weaker.


\subsection*{P0112-A001 | Full-sequence convergence of FPP geodesics to a ray structural argument}

\entrylabel{Argument or result}

The space of infinite rooted paths in a locally finite graph has a compact product topology after adjoining finite-path endpoints. Thus long finite geodesics admit subsequential limits. But the abstract sequence of paths alternating between two disjoint infinite branches has two subsequential limits and no full limit.


\entrylabel{Second check}

On a rooted binary tree,\allowbreak{} alternating the all-left and all-right depth-n paths supplies an explicit sequence illustrating the failure of compactness to imply uniqueness.


\entrylabel{Remaining obstruction}

The example is not an iid lattice FPP counterexample. A uniqueness-of-limit or directional-coalescence argument is missing for the specified endpoint sequence.


Verification:\allowbreak{} author checked limited argument | Independent review:\allowbreak{} none | Full target solved:\allowbreak{} False


\clearpage

\section*{P0113  Almost-sure escape from a fixed midpoint in FPP}

\entrylabel{Primary source}

Auffinger,\allowbreak{} Damron and Hanson:\allowbreak{} 50 years of first-passage percolation;\allowbreak{} Question26,\allowbreak{} PDF p.92


\url{https://arxiv.org/pdf/1511.03262}


Source statement date:\allowbreak{} 2016-09-26. Archived access:\allowbreak{} 2026-09-06T18:\allowbreak{}34:\allowbreak{}26Z.


Historical:\allowbreak{} source\_\allowbreak{}says\_\allowbreak{}open | Current:\allowbreak{} unverified | Result:\allowbreak{} attempted\_\allowbreak{}unsolved


\entrylabel{Question}

In uniquely geodesic planar iid FPP,\allowbreak{} prove that no fixed vertex belongs to geodesics from-ne1 to ne1 for infinitely many n.


\entrylabel{Scope}

This is the stronger almost-sure formulation printed in the survey. The known BKS result that the probability of midpoint membership tends to zero is not by itself this statement.


\subsection*{P0113-A001 | Almost-sure escape from a fixed midpoint in FPP structural argument}

\entrylabel{Argument or result}

If a fixed vertex v belongs to infinitely many such geodesics,\allowbreak{} both endpoint distances from v tend to infinity. A diagonal subsequence of the two oriented halves produces a doubly infinite geodesic through v. Therefore absence of all bigeodesics would imply the target.


\entrylabel{Second check}

Finite degree supplies the diagonal subsequence and every limiting finite segment is a minimizing segment of a selected geodesic. Conversely,\allowbreak{} probabilities tending to0 cannot justify a Borel-Cantelli conclusion:\allowbreak{} independent events with probability1/\allowbreak{}n occur infinitely often.


\entrylabel{Remaining obstruction}

Neither a summable midpoint bound nor unconditional absence of bigeodesics is obtained. Later midpoint papers must be compared against this exact almost-sure quantifier.


Verification:\allowbreak{} author checked limited argument | Independent review:\allowbreak{} none | Full target solved:\allowbreak{} False


\clearpage

\section*{P0114  Uniqueness of planar directional geodesics in every fixed direction}

\entrylabel{Primary source}

Auffinger,\allowbreak{} Damron and Hanson:\allowbreak{} 50 years of first-passage percolation;\allowbreak{} Questions27-28,\allowbreak{} PDF p.95


\url{https://arxiv.org/pdf/1511.03262}


Source statement date:\allowbreak{} 2016-09-26. Archived access:\allowbreak{} 2026-09-06T18:\allowbreak{}34:\allowbreak{}26Z.


Historical:\allowbreak{} source\_\allowbreak{}says\_\allowbreak{}open | Current:\allowbreak{} unverified | Result:\allowbreak{} attempted\_\allowbreak{}unsolved


\entrylabel{Question}

Show that every deterministic direction belongs to the Licea-Newman full-measure set D:\allowbreak{} from a fixed lattice vertex there are almost surely not two distinct geodesic rays with that direction.


\entrylabel{Scope}

The axial case and all-direction case are grouped. This is at-most-one,\allowbreak{} not an assertion that a ray in that direction exists. The null event may depend on the fixed direction.


\subsection*{P0114-A001 | Uniqueness of planar directional geodesics in every fixed direction structural argument}

\entrylabel{Argument or result}

A full-measure deterministic set may omit the axis,\allowbreak{} for example[0,\allowbreak{}2pi) without0. Also,\allowbreak{} a random exceptional set can be nonempty almost surely while excluding each fixed angle almost surely:\allowbreak{} a singleton containing an independent uniform angle has this property.


\entrylabel{Second check}

The first example checks the deterministic exceptional-set issue. In the second,\allowbreak{} P(U=\allowbreak{}theta)=\allowbreak{}0 for every fixed theta but the singleton is never empty.


\entrylabel{Remaining obstruction}

Fubini-based almost-every-direction arguments do not settle even a prescribed exceptional angle. A geometric argument sensitive to that direction is needed.


Verification:\allowbreak{} author checked limited argument | Independent review:\allowbreak{} none | Full target solved:\allowbreak{} False


\clearpage

\section*{P0115  Coalescence of directional FPP rays in higher dimensions}

\entrylabel{Primary source}

Auffinger,\allowbreak{} Damron and Hanson:\allowbreak{} 50 years of first-passage percolation;\allowbreak{} Question29,\allowbreak{} PDF p.97;\allowbreak{} Theorem4.22,\allowbreak{} p.96


\url{https://arxiv.org/pdf/1511.03262}


Source statement date:\allowbreak{} 2016-09-26. Archived access:\allowbreak{} 2026-09-06T18:\allowbreak{}34:\allowbreak{}26Z.


Historical:\allowbreak{} source\_\allowbreak{}says\_\allowbreak{}open | Current:\allowbreak{} unverified | Result:\allowbreak{} attempted\_\allowbreak{}unsolved


\entrylabel{Question}

Extend the planar coalescence conclusion to higher dimensions:\allowbreak{} for a fixed direction with almost-sure uniqueness of a directional ray from each vertex,\allowbreak{} must directional rays from different vertices coalesce?


\entrylabel{Scope}

The source asks whether its two-dimensional theorem extends. The relevant directional-uniqueness hypothesis must be stated rather than silently assumed in higher dimensions.


\subsection*{P0115-A001 | Coalescence of directional FPP rays in higher dimensions structural argument}

\entrylabel{Argument or result}

Translation invariance and one outgoing edge at every vertex do not force coalescence. Orient every edge line in the positive e1 direction in Z\textasciicircum{}d. There is exactly one forward ray from each vertex,\allowbreak{} but rays on different parallel lines never meet.


\entrylabel{Second check}

The construction is deterministic and translation invariant. All rays have the same asymptotic direction,\allowbreak{} so directional agreement alone adds no coalescence.


\entrylabel{Remaining obstruction}

This is a forest counterexample to a proof strategy,\allowbreak{} not the iid FPP model. Planar shielding and ordering used in the source proof do not extend automatically.


Verification:\allowbreak{} author checked limited argument | Independent review:\allowbreak{} none | Full target solved:\allowbreak{} False


\clearpage

\section*{P0116  Growth of the infinite-descendant part of the FPP geodesic tree}

\entrylabel{Primary source}

Auffinger,\allowbreak{} Damron and Hanson:\allowbreak{} 50 years of first-passage percolation;\allowbreak{} Question30,\allowbreak{} PDF p.97


\url{https://arxiv.org/pdf/1511.03262}


Source statement date:\allowbreak{} 2016-09-26. Archived access:\allowbreak{} 2026-09-06T18:\allowbreak{}34:\allowbreak{}26Z.


Historical:\allowbreak{} source\_\allowbreak{}says\_\allowbreak{}open | Current:\allowbreak{} unverified | Result:\allowbreak{} attempted\_\allowbreak{}unsolved


\entrylabel{Question}

Let f(r) count edges meeting the Euclidean circle of radius r that belong to infinitely many finite geodesics from0. Determine whether f(r) tends to infinity,\allowbreak{} and at what rate.


\entrylabel{Scope}

Assume a unique rooted finite-geodesic tree. Source questions about how many rays exist and this radial growth function are related but the latter requests additional spatial rate information.


\subsection*{P0116-A001 | Growth of the infinite-descendant part of the FPP geodesic tree structural argument}

\entrylabel{Argument or result}

In a locally finite rooted tree,\allowbreak{} an edge has infinitely many descendants exactly when it lies on an infinite ray,\allowbreak{} by the finite-branching compactness lemma. If the tree has infinitely many ends,\allowbreak{} then for every k one can choose k rays with disjoint tails outside some fixed ball;\allowbreak{} each tail crosses every sufficiently large circle. Consequently f(r)>=\allowbreak{}k eventually.


\entrylabel{Second check}

The forward implication is Konig's lemma applied below the edge;\allowbreak{} the reverse is immediate. A finite selection of distinct ends separates beyond finitely many branch vertices.


\entrylabel{Remaining obstruction}

This proves divergence conditional on infinitely many ends,\allowbreak{} without a rate. It supplies no unconditional ray multiplicity and no bound on where k rays separate.


Verification:\allowbreak{} author checked limited argument | Independent review:\allowbreak{} none | Full target solved:\allowbreak{} False


\clearpage

\section*{P0117  Absence of bi-infinite geodesics in planar iid FPP}

\entrylabel{Primary source}

Auffinger,\allowbreak{} Damron and Hanson:\allowbreak{} 50 years of first-passage percolation;\allowbreak{} Question31,\allowbreak{} PDF p.98


\url{https://arxiv.org/pdf/1511.03262}


Source statement date:\allowbreak{} 2016-09-26. Archived access:\allowbreak{} 2026-09-06T18:\allowbreak{}34:\allowbreak{}26Z.


Historical:\allowbreak{} source\_\allowbreak{}says\_\allowbreak{}open | Current:\allowbreak{} unverified | Result:\allowbreak{} attempted\_\allowbreak{}unsolved


\entrylabel{Question}

For nondegenerate continuous iid planar passage-time laws in the usual positive-speed regime,\allowbreak{} determine whether a bi-infinite geodesic exists.


\entrylabel{Scope}

The survey's short question sits in the uniquely geodesic random setting. Deterministic equal weights and the one-dimensional lattice have obvious geodesic lines and do not resolve the intended planar question.


\subsection*{P0117-A001 | Absence of bi-infinite geodesics in planar iid FPP structural argument}

\entrylabel{Argument or result}

Taking geodesics between-ne1 and ne1 gives candidate long paths,\allowbreak{} but a compactness argument produces a bi-infinite limit through a bounded region only if their distance from that region is tight along a subsequence. Paths can escape every fixed ball while connecting endpoints on opposite sides.


\entrylabel{Second check}

An explicit lattice path from(-n,\allowbreak{}0) up to(-n,\allowbreak{}n),\allowbreak{} across to(n,\allowbreak{}n),\allowbreak{} then down to(n,\allowbreak{}0) avoids every fixed ball for all large n;\allowbreak{} it illustrates the missing compactness anchor,\allowbreak{} without being asserted to be a geodesic.


\entrylabel{Remaining obstruction}

Finite-geodesic existence does not give the required anchoring. Proving or ruling out such anchoring under iid passage weights remains unresolved here.


Verification:\allowbreak{} author checked limited argument | Independent review:\allowbreak{} none | Full target solved:\allowbreak{} False


\clearpage

\section*{P0118  Noncircular low-dimensional Eden limit shape}

\entrylabel{Primary source}

Auffinger,\allowbreak{} Damron and Hanson:\allowbreak{} 50 years of first-passage percolation;\allowbreak{} Question33,\allowbreak{} PDF p.130;\allowbreak{} model definition pp.128-129


\url{https://arxiv.org/pdf/1511.03262}


Source statement date:\allowbreak{} 2016-09-26. Archived access:\allowbreak{} 2026-09-06T18:\allowbreak{}34:\allowbreak{}26Z.


Historical:\allowbreak{} source\_\allowbreak{}says\_\allowbreak{}open | Current:\allowbreak{} unverified | Result:\allowbreak{} attempted\_\allowbreak{}unsolved


\entrylabel{Question}

Prove that the Eden model's limit shape is not a Euclidean ball in the dimensions left open by existing high-dimensional results,\allowbreak{} in particular dimension2.


\entrylabel{Scope}

The source defines an Eden version choosing a boundary vertex uniformly. Edge-exponential infection generally chooses vertices at rates proportional to infected neighbors;\allowbreak{} these versions must not be identified without checking the time change.


\subsection*{P0118-A001 | Noncircular low-dimensional Eden limit shape structural argument}

\entrylabel{Argument or result}

With independent rate-one clocks on currently available boundary vertices,\allowbreak{} the next added vertex is uniform and the holding time is exponential with rate equal to their number. In edge-clock growth,\allowbreak{} a vertex with k infected neighbors has k active edge clocks and hence total rate k. A boundary vertex with two infected neighbors is selected twice as often as one with a single neighbor in that version.


\entrylabel{Second check}

The minimum of m independent rate-one exponentials has rate m and a uniform minimizing index. Summing k independent clock hazards gives rate k.


\entrylabel{Remaining obstruction}

This identifies the correct finite-state generator and prevents a model substitution,\allowbreak{} but the angular growth rate does not close as a function of the current radius. No anisotropic limit-shape inequality is obtained.


Verification:\allowbreak{} author checked limited argument | Independent review:\allowbreak{} none | Full target solved:\allowbreak{} False


\clearpage

\section*{P0119  Noncoexistence of unequal-speed Richardson infections}

\entrylabel{Primary source}

Auffinger,\allowbreak{} Damron and Hanson:\allowbreak{} 50 years of first-passage percolation;\allowbreak{} Questions34-35,\allowbreak{} PDF pp.134-135


\url{https://arxiv.org/pdf/1511.03262}


Source statement date:\allowbreak{} 2016-09-26. Archived access:\allowbreak{} 2026-09-06T18:\allowbreak{}34:\allowbreak{}26Z.


Historical:\allowbreak{} source\_\allowbreak{}says\_\allowbreak{}open | Current:\allowbreak{} unverified | Result:\allowbreak{} attempted\_\allowbreak{}unsolved


\entrylabel{Question}

Show that two-type Richardson growth on Z\textasciicircum{}d,\allowbreak{} d>=\allowbreak{}2,\allowbreak{} from distinct finite seeds cannot coexist indefinitely at unequal infection speeds;\allowbreak{} establish the proposed monotonicity in relative speed if that route is valid.


\entrylabel{Scope}

The source assumes an ordering of the two parameters before asking monotonicity. Speed versus mean-passage-time conventions must be fixed. The universal noncoexistence and monotonicity route are grouped.


\subsection*{P0119-A001 | Noncoexistence of unequal-speed Richardson infections structural argument}

\entrylabel{Argument or result}

A common-clock monotone comparison of one infection's absolute size does not imply monotonicity of coexistence,\allowbreak{} because the latter requires both species to remain infinite. On Z with seeds at0 and1,\allowbreak{} every site at most0 can only be reached first through the left seed and every site at least1 through the right seed;\allowbreak{} with positive finite passage times both species survive for every speed ratio.


\entrylabel{Second check}

The seeds block any nearest-neighbor path across the interface once colors are permanent. Both outward rays are infinite,\allowbreak{} independent of relative rates.


\entrylabel{Remaining obstruction}

The one-dimensional example shows the essential role of bypass geometry and is excluded from the intended d>=\allowbreak{}2 conclusion. A coupling controlling the joint survival event in higher dimensions remains missing.


Verification:\allowbreak{} author checked limited argument | Independent review:\allowbreak{} none | Full target solved:\allowbreak{} False


\clearpage

\section*{P0120  Strict interface speedup in half-plane inhomogeneous FPP}

\entrylabel{Primary source}

Auffinger,\allowbreak{} Damron and Hanson:\allowbreak{} 50 years of first-passage percolation;\allowbreak{} Question38,\allowbreak{} PDF p.142


\url{https://arxiv.org/pdf/1511.03262}


Source statement date:\allowbreak{} 2016-09-26. Archived access:\allowbreak{} 2026-09-06T18:\allowbreak{}34:\allowbreak{}26Z.


Historical:\allowbreak{} source\_\allowbreak{}says\_\allowbreak{}open | Current:\allowbreak{} unverified | Result:\allowbreak{} attempted\_\allowbreak{}unsolved


\entrylabel{Question}

Characterize the pairs of iid half-plane edge laws F\_\allowbreak{}minus,\allowbreak{}F\_\allowbreak{}plus for which vertical growth along their interface is strictly faster than in either homogeneous model.


\entrylabel{Scope}

Greater speed means a smaller directional time constant. The exact convention for edges on the interface must be fixed;\allowbreak{} finite detours in the constant-law check below do not affect its conclusion.


\subsection*{P0120-A001 | Strict interface speedup in half-plane inhomogeneous FPP structural argument}

\entrylabel{Argument or result}

For constant times a on one half-plane and b on the other,\allowbreak{} let c=\allowbreak{}min(a,\allowbreak{}b)>0. Every path with vertical displacement n costs at least cn. A path taking a bounded horizontal detour into the cheaper half-plane,\allowbreak{} moving n steps vertically,\allowbreak{} then returning costs cn+\allowbreak{}O(1). Thus the interface time constant is exactly c,\allowbreak{} so there is no strict speedup beyond both homogeneous models.


\entrylabel{Second check}

The lower bound uses path length at least n;\allowbreak{} the explicit path gives a matching upper bound after division by n.


\entrylabel{Remaining obstruction}

Random switching between regions may exploit complementary favorable segments. The deterministic comparison cannot characterize when that advantage is strictly positive.


Verification:\allowbreak{} author checked limited argument | Independent review:\allowbreak{} none | Full target solved:\allowbreak{} False


\clearpage

\section*{P0121  Uniform Lipschitz height models and their interfaces}

\entrylabel{Primary source}

Schramm:\allowbreak{} Conformally invariant scaling limits,\allowbreak{} overview and collection of problems,\allowbreak{} 2006v2;\allowbreak{} Problems2.2-2.3,\allowbreak{} PDF pp.15-16


\url{https://arxiv.org/pdf/math/0602151}


Source statement date:\allowbreak{} 2006. Archived access:\allowbreak{} 2026-09-06T18:\allowbreak{}33:\allowbreak{}45Z.


Historical:\allowbreak{} source\_\allowbreak{}says\_\allowbreak{}open | Current:\allowbreak{} unverified | Result:\allowbreak{} attempted\_\allowbreak{}unsolved


\entrylabel{Question}

Establish Gaussian-free-field scaling and SLE4 interface limits for the source's uniform odd-integer triangular-lattice heights,\allowbreak{} and for its uniform real-valued Lipschitz heights with the appropriate boundary amplitude.


\entrylabel{Scope}

The two genuinely different height ensembles are conservatively grouped as one universality family. Integer heights have nearest-neighbor difference0 or2 and boundary values plus/\allowbreak{}minus1;\allowbreak{} real heights have difference at most1 and tuned boundary values plus/\allowbreak{}minus lambda.


\subsection*{P0121-A001 | Uniform Lipschitz height models and their interfaces structural reduction or obstruction}

\entrylabel{Argument or result}

Conditioning an interior height on all other heights gives a uniform law on its feasible interval,\allowbreak{} intersected with the odd integers in the discrete model. The interval endpoints are the maximum of neighbor heights minus the slope bound and the minimum plus the bound. Its conditional mean is the midpoint of these extrema,\allowbreak{} not generally the arithmetic mean of neighboring heights.


\entrylabel{Second check}

For real slope bound1 and six neighboring heights0,\allowbreak{}0,\allowbreak{}0,\allowbreak{}0,\allowbreak{}0,\allowbreak{}1,\allowbreak{} the feasible interval is[0,\allowbreak{}1],\allowbreak{} so the conditional mean is1/\allowbreak{}2 while the neighbor average is1/\allowbreak{}6. These neighbor values satisfy their mutual triangular-edge constraints.


\entrylabel{Remaining obstruction}

The discrete Gaussian harmonic-observable proof therefore does not transfer by replacing a Gaussian field with a uniform Lipschitz field. One must first control the effective large-scale field and then the discontinuous operation of extracting its zero interface.


Verification:\allowbreak{} author checked limited argument | Independent review:\allowbreak{} none | Full target solved:\allowbreak{} False


\clearpage

\section*{P0122  SLE4 from spectral truncations of the Gaussian free field}

\entrylabel{Primary source}

Schramm:\allowbreak{} Conformally invariant scaling limits,\allowbreak{} overview and collection of problems,\allowbreak{} 2006v2;\allowbreak{} Problem2.4,\allowbreak{} PDF p.16


\url{https://arxiv.org/pdf/math/0602151}


Source statement date:\allowbreak{} 2006. Archived access:\allowbreak{} 2026-09-06T18:\allowbreak{}33:\allowbreak{}45Z.


Historical:\allowbreak{} source\_\allowbreak{}says\_\allowbreak{}open | Current:\allowbreak{} unverified | Result:\allowbreak{} attempted\_\allowbreak{}unsolved


\entrylabel{Question}

Project a planar Dirichlet Gaussian free field onto Laplace eigenmodes with eigenvalue magnitude at most R,\allowbreak{} add the harmonic extension of suitably chosen two-arc boundary values,\allowbreak{} and prove that the connecting zero interface converges to SLE4 as R tends to infinity.


\entrylabel{Scope}

The cutoff is spectral and nonlocal. Distributional convergence of the field itself is already built into the construction and is not the requested interface convergence.


\subsection*{P0122-A001 | SLE4 from spectral truncations of the Gaussian free field structural reduction or obstruction}

\entrylabel{Argument or result}

Even uniform convergence of smooth functions does not make zero sets continuous. On a fixed square,\allowbreak{} h\_\allowbreak{}n(x,\allowbreak{}y)=\allowbreak{}1/\allowbreak{}n has no zero set,\allowbreak{} whereas its uniform limit0 vanishes everywhere. With changing signs,\allowbreak{} n\textasciicircum{}(-1)sin(nx) also converges uniformly to0 but has many macroscopic vertical zero components.


\entrylabel{Second check}

The sup norms are at most1/\allowbreak{}n. For the sinusoidal example the zero lines are x=\allowbreak{}k*pi/\allowbreak{}n wherever they lie in the square;\allowbreak{} their number diverges.


\entrylabel{Remaining obstruction}

These examples do not satisfy the prescribed random-field law or boundary data;\allowbreak{} they invalidate only the proposed continuous-mapping argument. Uniform control of near-zero topology and a characterization of the selected interface are missing.


Verification:\allowbreak{} author checked limited argument | Independent review:\allowbreak{} none | Full target solved:\allowbreak{} False


\clearpage

\section*{P0123  High-temperature triangular Ising interfaces converge to SLE6}

\entrylabel{Primary source}

Schramm:\allowbreak{} Conformally invariant scaling limits,\allowbreak{} overview and collection of problems,\allowbreak{} 2006v2;\allowbreak{} Problem2.5,\allowbreak{} PDF p.17,\allowbreak{} high-temperature portion


\url{https://arxiv.org/pdf/math/0602151}


Source statement date:\allowbreak{} 2006. Archived access:\allowbreak{} 2026-09-06T18:\allowbreak{}33:\allowbreak{}45Z.


Historical:\allowbreak{} source\_\allowbreak{}says\_\allowbreak{}open | Current:\allowbreak{} unverified | Result:\allowbreak{} attempted\_\allowbreak{}unsolved


\entrylabel{Question}

For every fixed inverse temperature strictly between0 and the triangular Ising critical inverse temperature,\allowbreak{} prove convergence of the two-boundary-arc spin interface to chordal SLE6.


\entrylabel{Scope}

The critical-temperature SLE3 component has been resolved in later literature and is not admitted here. The independent-spin beta=\allowbreak{}0 endpoint is also known. Fixed positive high-temperature correlations are the target.


\subsection*{P0123-A001 | High-temperature triangular Ising interfaces converge to SLE6 structural reduction or obstruction}

\entrylabel{Argument or result}

The high-temperature Gibbs law is not a small total-variation perturbation of independent spins uniformly over a growing domain. For n independent two-spin blocks with coupling beta>0,\allowbreak{} the average within-block product concentrates at tanh(beta),\allowbreak{} whereas at beta=\allowbreak{}0 it concentrates at0. Choosing a threshold halfway between these means distinguishes the two product laws with probability tending to1.


\entrylabel{Second check}

Each block product is plus/\allowbreak{}minus1 with mean tanh(beta);\allowbreak{} its variance is at most1. Chebyshev applied to n blocks proves asymptotic separation without a central limit approximation.


\entrylabel{Remaining obstruction}

Thus finite-range or exponentially decaying local dependence does not justify global total-variation replacement. A renormalization or observable argument controlling crossing geometry is still needed for the actual triangular model.


Verification:\allowbreak{} author checked limited argument | Independent review:\allowbreak{} none | Full target solved:\allowbreak{} False


\clearpage

\section*{P0124  Critical planar random-cluster interfaces for general cluster weight}

\entrylabel{Primary source}

Schramm:\allowbreak{} Conformally invariant scaling limits,\allowbreak{} overview and collection of problems,\allowbreak{} 2006v2;\allowbreak{} Problems2.6 and2.11,\allowbreak{} PDF pp.17-18


\url{https://arxiv.org/pdf/math/0602151}


Source statement date:\allowbreak{} 2006. Archived access:\allowbreak{} 2026-09-06T18:\allowbreak{}33:\allowbreak{}45Z.


Historical:\allowbreak{} source\_\allowbreak{}says\_\allowbreak{}open | Current:\allowbreak{} unverified | Result:\allowbreak{} attempted\_\allowbreak{}unsolved


\entrylabel{Question}

For the square-lattice FK model at p=\allowbreak{}sqrt(q)/\allowbreak{}(1+\allowbreak{}sqrt(q)),\allowbreak{} q in(0,\allowbreak{}4),\allowbreak{} establish the predicted chordal SLE limit with kappa=\allowbreak{}4*pi/\allowbreak{}arccos(-sqrt(q)/\allowbreak{}2).


\entrylabel{Scope}

The q=\allowbreak{}2 FK-Ising case is already resolved and excluded from the unresolved scope. Square bond percolation is q=\allowbreak{}1 and is grouped here instead of being counted again as Problem2.11.


\subsection*{P0124-A001 | Critical planar random-cluster interfaces for general cluster weight structural reduction or obstruction}

\entrylabel{Argument or result}

Given all edges except e,\allowbreak{} the FK conditional probability that e is open is p if its endpoints are already connected,\allowbreak{} and p/\allowbreak{}[p+\allowbreak{}q(1-p)] otherwise. At the proposed self-dual p these become sqrt(q)/\allowbreak{}(1+\allowbreak{}sqrt(q)) and1/\allowbreak{}(1+\allowbreak{}sqrt(q)). They straddle1/\allowbreak{}2 for q>1.


\entrylabel{Second check}

Opening e changes its Bernoulli odds by p/\allowbreak{}(1-p);\allowbreak{} if it merges components it also changes the cluster-count weight by1/\allowbreak{}q. Normalizing the two possible weights gives the displayed conditional probabilities.


\entrylabel{Remaining obstruction}

Self-duality supplies symmetry but not an independent percolation exploration law. The dependence on long-range connectivity prevents a direct reuse of the q=\allowbreak{}1 or q=\allowbreak{}2 observable argument.


Verification:\allowbreak{} author checked limited argument | Independent review:\allowbreak{} none | Full target solved:\allowbreak{} False


\clearpage

\section*{P0125  Conformal limits of the critical and dense O(n) loop model}

\entrylabel{Primary source}

Schramm:\allowbreak{} Conformally invariant scaling limits,\allowbreak{} overview and collection of problems,\allowbreak{} 2006v2;\allowbreak{} Problem2.7,\allowbreak{} PDF p.17


\url{https://arxiv.org/pdf/math/0602151}


Source statement date:\allowbreak{} 2006. Archived access:\allowbreak{} 2026-09-06T18:\allowbreak{}33:\allowbreak{}45Z.


Historical:\allowbreak{} source\_\allowbreak{}says\_\allowbreak{}open | Current:\allowbreak{} unverified | Result:\allowbreak{} attempted\_\allowbreak{}unsolved


\entrylabel{Question}

For the hexagonal O(n) loop model,\allowbreak{} establish the predicted SLE limits of a boundary-to-boundary strand at the critical edge weight and in the dense phase.


\entrylabel{Scope}

The source's whole parameter family is one problem. Previously admitted self-avoiding walk at n=\allowbreak{}0 is cross-referenced,\allowbreak{} and known Ising cases at n=\allowbreak{}1 are not new open items. The remaining interacting-loop ensembles form this entry.


\subsection*{P0125-A001 | Conformal limits of the critical and dense O(n) loop model structural reduction or obstruction}

\entrylabel{Argument or result}

Deleting all closed loops from a configuration does not give a self-avoiding walk with weights depending only on its length. The marginal weight of a strand gamma is proportional to x\textasciicircum{}|gamma| times the partition function of loops compatible with gamma. Two equal-length strands may leave different complementary regions and therefore different loop partition functions.


\entrylabel{Second check}

Summing the joint weight x\textasciicircum{}(|gamma|+\allowbreak{}total loop length)*n\textasciicircum{}(number of loops) over compatible loop families factors out x\textasciicircum{}|gamma| and leaves exactly the stated complementary partition function.


\entrylabel{Remaining obstruction}

The environmental partition-function factor is the missing ingredient in a reduction to self-avoiding walk. Its conformal asymptotics have not been derived here.


Verification:\allowbreak{} author checked limited argument | Independent review:\allowbreak{} none | Full target solved:\allowbreak{} False


\clearpage

\section*{P0126  Planar uniform lattice-tree diameter and scaling limit}

\entrylabel{Primary source}

Schramm:\allowbreak{} Conformally invariant scaling limits,\allowbreak{} overview and collection of problems,\allowbreak{} 2006v2;\allowbreak{} Problem2.8,\allowbreak{} PDF p.18


\url{https://arxiv.org/pdf/math/0602151}


Source statement date:\allowbreak{} 2006. Archived access:\allowbreak{} 2026-09-06T18:\allowbreak{}33:\allowbreak{}45Z.


Historical:\allowbreak{} source\_\allowbreak{}says\_\allowbreak{}open | Current:\allowbreak{} unverified | Result:\allowbreak{} attempted\_\allowbreak{}unsolved


\entrylabel{Question}

Determine the expected Euclidean diameter growth and the rescaled random-set limit of the uniform n-vertex lattice tree containing the origin,\allowbreak{} in particular on the planar square lattice.


\entrylabel{Scope}

A uniform embedded n-vertex tree is not a uniform spanning tree of a fixed region or a critical Galton-Watson tree. High-dimensional spread-out results do not settle the planar model.


\subsection*{P0126-A001 | Planar uniform lattice-tree diameter and scaling limit structural reduction or obstruction}

\entrylabel{Argument or result}

Every such planar tree has Euclidean diameter at most n-1. If its diameter is D,\allowbreak{} all its vertices lie in an axis-aligned square of side at most2D around any chosen vertex,\allowbreak{} containing at most(2D+\allowbreak{}1)\textasciicircum{}2 lattice points. Hence D is at least(sqrt(n)-1)/\allowbreak{}2,\allowbreak{} deterministically.


\entrylabel{Second check}

The upper bound follows from connectedness and at most n-1 edges on any simple tree path. The lower bound counts the containing square and applies to every connected n-vertex subset,\allowbreak{} independently of the measure.


\entrylabel{Remaining obstruction}

The interval between square-root and linear growth is too wide to identify the exponent. Expected-diameter normalization also gives no tightness of the full tree geometry by itself.


Verification:\allowbreak{} author checked limited argument | Independent review:\allowbreak{} none | Full target solved:\allowbreak{} False


\clearpage

\section*{P0127  Efficient approximately uniform sampling of lattice trees}

\entrylabel{Primary source}

Schramm:\allowbreak{} Conformally invariant scaling limits,\allowbreak{} overview and collection of problems,\allowbreak{} 2006v2;\allowbreak{} Problem2.9,\allowbreak{} PDF p.18


\url{https://arxiv.org/pdf/math/0602151}


Source statement date:\allowbreak{} 2006. Archived access:\allowbreak{} 2026-09-06T18:\allowbreak{}33:\allowbreak{}45Z.


Historical:\allowbreak{} source\_\allowbreak{}says\_\allowbreak{}open | Current:\allowbreak{} unverified | Result:\allowbreak{} attempted\_\allowbreak{}unsolved


\entrylabel{Question}

Give a provably efficient algorithm to sample approximately uniformly from the n-vertex lattice trees containing the origin,\allowbreak{} or prove an appropriate computational obstruction.


\entrylabel{Scope}

An ergodic chain with uniform stationary distribution is not automatically a polynomial-time sampler. The meaning of approximation must specify,\allowbreak{} for example,\allowbreak{} total variation error.


\subsection*{P0127-A001 | Efficient approximately uniform sampling of lattice trees structural reduction or obstruction}

\entrylabel{Argument or result}

All eligible trees lie in the radius n-1 lattice ball. Enumerating all subsets of its edges,\allowbreak{} retaining exactly the connected acyclic n-vertex subsets containing0,\allowbreak{} and choosing a uniformly random retained index gives an exact finite sampler. On Z\textasciicircum{}2 the ambient edge count is O(n\textasciicircum{}2),\allowbreak{} so the naive running time is at most exponential in n\textasciicircum{}2.


\entrylabel{Second check}

Every tree vertex is reached by at most n-1 unit steps from0,\allowbreak{} proving the finite containment. Acyclicity plus connectedness and n vertices characterizes the required objects,\allowbreak{} with each edge set counted once.


\entrylabel{Remaining obstruction}

This supplies existence and exactness but fails efficiency. The later Fredes-Marckert survey constructs uniform-stationary chains;\allowbreak{} their existence alone does not provide a polynomial mixing bound.


Verification:\allowbreak{} author checked limited argument | Independent review:\allowbreak{} none | Full target solved:\allowbreak{} False


\clearpage

\section*{P0128  Cardy conformal invariance for Voronoi percolation}

\entrylabel{Primary source}

Schramm:\allowbreak{} Conformally invariant scaling limits,\allowbreak{} overview and collection of problems,\allowbreak{} 2006v2;\allowbreak{} Problem2.12,\allowbreak{} PDF p.19


\url{https://arxiv.org/pdf/math/0602151}


Source statement date:\allowbreak{} 2006. Archived access:\allowbreak{} 2026-09-06T18:\allowbreak{}33:\allowbreak{}45Z.


Historical:\allowbreak{} source\_\allowbreak{}says\_\allowbreak{}open | Current:\allowbreak{} unverified | Result:\allowbreak{} attempted\_\allowbreak{}unsolved


\entrylabel{Question}

Prove the analogue of Smirnov's crossing-probability theorem for critical planar Voronoi percolation,\allowbreak{} with the corresponding conformally invariant interface limit.


\entrylabel{Scope}

Annealed Poisson geometry and independent fair cell colors are intended. Invariance under changing a background metric,\allowbreak{} known earlier,\allowbreak{} is a different statement from Cardy's formula.


\subsection*{P0128-A001 | Cardy conformal invariance for Voronoi percolation structural reduction or obstruction}

\entrylabel{Argument or result}

A planar conformal map f sends a homogeneous Poisson process of intensity lambda to an inhomogeneous process whose intensity at w is lambda/\allowbreak{}|f'(f\textasciicircum{}(-1)(w))|\textasciicircum{}2. Thus exact rotation and translation symmetry of the original process do not produce exact invariance under arbitrary conformal maps.


\entrylabel{Second check}

The expected image count in a Borel set A equals lambda times the area of f\textasciicircum{}(-1)(A);\allowbreak{} the change-of-variables formula gives the stated Jacobian. Disjoint-set independence is preserved,\allowbreak{} so this identifies the transformed process.


\entrylabel{Remaining obstruction}

One would need to prove that a smoothly varying local intensity does not affect the limiting critical crossing law. The Jacobian computation exposes that missing universality step without proving it.


Verification:\allowbreak{} author checked limited argument | Independent review:\allowbreak{} none | Full target solved:\allowbreak{} False


\clearpage

\section*{P0129  Cardy conformal invariance for Boolean continuum percolation}

\entrylabel{Primary source}

Schramm:\allowbreak{} Conformally invariant scaling limits,\allowbreak{} overview and collection of problems,\allowbreak{} 2006v2;\allowbreak{} Problem2.13,\allowbreak{} PDF p.19


\url{https://arxiv.org/pdf/math/0602151}


Source statement date:\allowbreak{} 2006. Archived access:\allowbreak{} 2026-09-06T18:\allowbreak{}33:\allowbreak{}45Z.


Historical:\allowbreak{} source\_\allowbreak{}says\_\allowbreak{}open | Current:\allowbreak{} unverified | Result:\allowbreak{} attempted\_\allowbreak{}unsolved


\entrylabel{Question}

Establish Smirnov-type conformally invariant critical crossing limits for the planar Poisson Boolean model of equal-radius disks.


\entrylabel{Scope}

Varying the disk radius at intensity1 and varying intensity at radius1 are related by a Euclidean scaling and are one problem,\allowbreak{} not two corpus entries.


\subsection*{P0129-A001 | Cardy conformal invariance for Boolean continuum percolation structural reduction or obstruction}

\entrylabel{Argument or result}

Scaling space by1/\allowbreak{}r maps intensity1 centers with disk radius r to radius-one disks whose centers have intensity r\textasciicircum{}2. Therefore all Euclidean-scale behavior depends on the dimensionless product lambda*r\textasciicircum{}2 for general intensity lambda. Under a nonaffine conformal map,\allowbreak{} however,\allowbreak{} each disk is deformed and its local size depends on position.


\entrylabel{Second check}

An area A in scaled coordinates corresponds to original area r\textasciicircum{}2 A,\allowbreak{} yielding mean center count r\textasciicircum{}2 A. This proves the exact parameter equivalence.


\entrylabel{Remaining obstruction}

The scalar scaling relation does not control spatially inhomogeneous radii and noncircular images under conformal maps. Proving their irrelevance at criticality remains the obstruction.


Verification:\allowbreak{} author checked limited argument | Independent review:\allowbreak{} none | Full target solved:\allowbreak{} False


\clearpage

\section*{P0130  Sharp planar one-arm probability beyond its exponent}

\entrylabel{Primary source}

Schramm:\allowbreak{} Conformally invariant scaling limits,\allowbreak{} overview and collection of problems,\allowbreak{} 2006v2;\allowbreak{} Problem3.1,\allowbreak{} PDF p.20,\allowbreak{} one-arm remainder


\url{https://arxiv.org/pdf/math/0602151}


Source statement date:\allowbreak{} 2006. Archived access:\allowbreak{} 2026-09-06T18:\allowbreak{}33:\allowbreak{}45Z.


Historical:\allowbreak{} source\_\allowbreak{}says\_\allowbreak{}open | Current:\allowbreak{} unverified | Result:\allowbreak{} attempted\_\allowbreak{}unsolved


\entrylabel{Question}

Improve the critical triangular-percolation one-arm probability R\textasciicircum{}(-5/\allowbreak{}48+\allowbreak{}o(1)) to sharp multiplicative asymptotics,\allowbreak{} at least bounds by constant multiples of R\textasciicircum{}(-5/\allowbreak{}48).


\entrylabel{Scope}

Du-Gao-Li-Zhuang resolve sharp half-plane and planar polychromatic multi-arm asymptotics;\allowbreak{} those resolved portions of the 2006 question are excluded. Their abstract does not include the planar one-arm event.


\subsection*{P0130-A001 | Sharp planar one-arm probability beyond its exponent structural reduction or obstruction}

\entrylabel{Argument or result}

Quasi-multiplicative information and a logarithmic exponent need not imply a constant asymptotic amplitude. For alpha>0 and0<a<alpha,\allowbreak{} define f(R)=\allowbreak{}R\textasciicircum{}(-alpha)*exp(a*sin(log R)),\allowbreak{} R>=\allowbreak{}1. Then f is decreasing,\allowbreak{} has logarithmic exponent-alpha,\allowbreak{} and obeys exp(-3a)<=\allowbreak{}f(RS)/\allowbreak{}(f(R)f(S))<=\allowbreak{}exp(3a),\allowbreak{} but R\textasciicircum{}alpha f(R) has no limit.


\entrylabel{Second check}

The derivative of log f with respect to log R is-alpha+\allowbreak{}a*cos(log R)<0. The quasi-multiplicative bound follows from three sine terms each bounded by1;\allowbreak{} subsequences with sine plus/\allowbreak{}minus1 give different amplitudes.


\entrylabel{Remaining obstruction}

This is a scalar obstruction to extracting an asymptotic constant from coarse multiplicativity. It is not a percolation counterexample,\allowbreak{} and even constant-factor bounds for the specified probability are not proved by it.


Verification:\allowbreak{} author checked limited argument | Independent review:\allowbreak{} none | Full target solved:\allowbreak{} False


\clearpage

\section*{P0131  Exact critical percolation chemical-distance exponent}

\entrylabel{Primary source}

Schramm:\allowbreak{} Conformally invariant scaling limits,\allowbreak{} overview and collection of problems,\allowbreak{} 2006v2;\allowbreak{} Problem3.3,\allowbreak{} PDF p.20


\url{https://arxiv.org/pdf/math/0602151}


Source statement date:\allowbreak{} 2006. Archived access:\allowbreak{} 2026-09-06T18:\allowbreak{}33:\allowbreak{}45Z.


Historical:\allowbreak{} source\_\allowbreak{}says\_\allowbreak{}open | Current:\allowbreak{} unverified | Result:\allowbreak{} attempted\_\allowbreak{}unsolved


\entrylabel{Question}

Determine the exponent alpha for the expected length of the shortest left-right critical percolation crossing of a square of side R,\allowbreak{} conditional on a crossing existing.


\entrylabel{Scope}

The Euclidean connection event,\allowbreak{} the number of edges in a shortest open path,\allowbreak{} and its conditional expectation are distinct quantities. The source requests the exponent itself,\allowbreak{} not merely improved upper or lower bounds.


\subsection*{P0131-A001 | Exact critical percolation chemical-distance exponent structural reduction or obstruction}

\entrylabel{Argument or result}

A shortest crossing may be chosen self-avoiding by deleting loops. Its length is at least a constant times R and at most the number of sites in the square,\allowbreak{} which is O(R\textasciicircum{}2). Thus any exponent defined by the stated expectation must lie in[1,\allowbreak{}2].


\entrylabel{Second check}

The lower bound is net horizontal displacement divided by the maximum horizontal displacement per edge. The upper bound counts sites on a self-avoiding path,\allowbreak{} so it holds on every configuration with a crossing.


\entrylabel{Remaining obstruction}

These geometric bounds do not quantify the critical detours around closed clusters. Interface dimensions are not automatically dimensions of shortest open crossings.


Verification:\allowbreak{} author checked limited argument | Independent review:\allowbreak{} none | Full target solved:\allowbreak{} False


\clearpage

\section*{P0132  Measurable loop erasure of planar Brownian motion}

\entrylabel{Primary source}

Schramm:\allowbreak{} Conformally invariant scaling limits,\allowbreak{} overview and collection of problems,\allowbreak{} 2006v2;\allowbreak{} Problem6.1,\allowbreak{} PDF p.24;\allowbreak{} Berestycki-Sauzedde2026,\allowbreak{} Section7 first question


\url{https://arxiv.org/pdf/math/0602151}


Source statement date:\allowbreak{} 2006. Archived access:\allowbreak{} 2026-09-06T18:\allowbreak{}33:\allowbreak{}45Z.


Historical:\allowbreak{} source\_\allowbreak{}says\_\allowbreak{}open | Current:\allowbreak{} unverified | Result:\allowbreak{} attempted\_\allowbreak{}unsolved


\entrylabel{Question}

In the joint scaling limit of a planar random walk and its loop erasure,\allowbreak{} determine whether the SLE2 path is a measurable function of the Brownian path.


\entrylabel{Scope}

The 2026 loop-addition theorem establishes the coupling and its chronological Brownian law;\allowbreak{} the authors explicitly leave inverse measurability open. Existence of a coupling does not imply a deterministic factor.


\subsection*{P0132-A001 | Measurable loop erasure of planar Brownian motion structural reduction or obstruction}

\entrylabel{Argument or result}

For random elements(G,\allowbreak{}W) in standard Borel spaces,\allowbreak{} G is measurable from W exactly when the regular conditional distribution Law(G|W) is a point mass almost surely. An arbitrary measurable choice of a point in each fiber is insufficient:\allowbreak{} if G is standard normal and W=\allowbreak{}|G|,\allowbreak{} choosing the positive preimage changes the joint law,\allowbreak{} whose conditional signs remain fair.


\entrylabel{Second check}

The forward implication follows by conditioning G=\allowbreak{}g(W). For the reverse,\allowbreak{} a probability kernel supported on Dirac measures identifies a measurable atom through a countable separating family. In the normal example symmetry gives conditional sign probability1/\allowbreak{}2 for W>0.


\entrylabel{Remaining obstruction}

No degeneration of the conditional SLE law given the complete Brownian chronology has been shown. The new forward map does not provide this missing conditional-law argument.


Verification:\allowbreak{} author checked limited argument | Independent review:\allowbreak{} none | Full target solved:\allowbreak{} False


\clearpage

\section*{P0133  Torus spanning-tree dual homotopy distribution}

\entrylabel{Primary source}

Schramm:\allowbreak{} Conformally invariant scaling limits,\allowbreak{} overview and collection of problems,\allowbreak{} 2006v2;\allowbreak{} Problem6.2,\allowbreak{} PDF p.24


\url{https://arxiv.org/pdf/math/0602151}


Source statement date:\allowbreak{} 2006. Archived access:\allowbreak{} 2026-09-06T18:\allowbreak{}33:\allowbreak{}45Z.


Historical:\allowbreak{} source\_\allowbreak{}says\_\allowbreak{}open | Current:\allowbreak{} unverified | Result:\allowbreak{} attempted\_\allowbreak{}unsolved


\entrylabel{Question}

Determine the limiting distribution of the homotopy classes of the three edge-simple noncontractible closed paths in the dual complement of a uniform spanning tree on a fine toroidal square grid.


\entrylabel{Scope}

The question concerns the joint winding information,\allowbreak{} not only the number of dual cycles. The geometry of the torus is an input. Later line-bundle determinant methods require a separate exact-scope comparison.


\subsection*{P0133-A001 | Torus spanning-tree dual homotopy distribution structural reduction or obstruction}

\entrylabel{Argument or result}

On a toroidal N-vertex square grid,\allowbreak{} E=\allowbreak{}2N and there are N faces. Removing the N-1 primal tree edges leaves N+\allowbreak{}1 dual edges. The dual complement is connected,\allowbreak{} so its first Betti number is(N+\allowbreak{}1)-N+\allowbreak{}1=\allowbreak{}2. This explains two independent topological cycles but does not itself specify the triple of simple-loop homotopy classes.


\entrylabel{Second check}

Euler's formula on the torus gives V-E+\allowbreak{}F=\allowbreak{}0. Dual connectivity follows because a separating primal-dual cut would correspond to a cycle contained in the primal tree,\allowbreak{} which is impossible. The cycle-rank formula is E-V+\allowbreak{}1 for a connected finite graph.


\entrylabel{Remaining obstruction}

Rank2 determines neither winding slopes nor their probabilities. A scalar Laplacian determinant counts all spanning trees without resolving this joint topological information.


Verification:\allowbreak{} author checked limited argument | Independent review:\allowbreak{} none | Full target solved:\allowbreak{} False


\clearpage

\section*{P0134  A nontrivial exact Hausdorff gauge for SLE}

\entrylabel{Primary source}

Schramm:\allowbreak{} Conformally invariant scaling limits,\allowbreak{} overview and collection of problems,\allowbreak{} 2006v2;\allowbreak{} Problem7.1,\allowbreak{} PDF p.25


\url{https://arxiv.org/pdf/math/0602151}


Source statement date:\allowbreak{} 2006. Archived access:\allowbreak{} 2026-09-06T18:\allowbreak{}33:\allowbreak{}45Z.


Historical:\allowbreak{} source\_\allowbreak{}says\_\allowbreak{}open | Current:\allowbreak{} unverified | Result:\allowbreak{} attempted\_\allowbreak{}unsolved


\entrylabel{Question}

Find a nontrivial Hausdorff gauge measure for SLE\_\allowbreak{}kappa,\allowbreak{} kappa<8,\allowbreak{} which is sigma-finite and preferably finite on compact subarcs,\allowbreak{} as intended by the source's natural-parametrization question.


\entrylabel{Scope}

Literally,\allowbreak{} the zero measure is sigma-finite;\allowbreak{} the source context requires a useful nonzero measure. Rezaei proves that the critical power gauge r\textasciicircum{}(1+\allowbreak{}kappa/\allowbreak{}8) gives zero Hausdorff measure. Minkowski content is not automatically a Hausdorff measure.


\subsection*{P0134-A001 | A nontrivial exact Hausdorff gauge for SLE structural reduction or obstruction}

\entrylabel{Argument or result}

Hausdorff dimension alone cannot select the exact gauge. For any set of dimension d,\allowbreak{} a power exponent s>d gives zero Hausdorff measure,\allowbreak{} while s<d gives infinite measure;\allowbreak{} at s=\allowbreak{}d either zero,\allowbreak{} finite positive,\allowbreak{} or infinite behavior can occur. Multiplying r\textasciicircum{}d by a slowly varying factor can change the critical measure while leaving the dimension unchanged.


\entrylabel{Second check}

The two off-critical conclusions are the defining threshold property of Hausdorff dimension. At the critical exponent,\allowbreak{} a smooth interval has finite positive H\textasciicircum{}1,\allowbreak{} while a countable union of disjoint intervals with divergent total length has infinite H\textasciicircum{}1;\allowbreak{} critical zero examples exist among dimension-one Cantor sets.


\entrylabel{Remaining obstruction}

One needs fine local occupation estimates,\allowbreak{} not just dimension or existence of Minkowski content. The predicted logarithmic correction has not been proved here.


Verification:\allowbreak{} author checked limited argument | Independent review:\allowbreak{} none | Full target solved:\allowbreak{} False


\clearpage

\section*{P0135  Almost-sure fibers of chronological loop addition}

\entrylabel{Primary source}

Berestycki and Sauzedde:\allowbreak{} Brownian paths as loop-decorated SLEs,\allowbreak{}2026;\allowbreak{} Section7,\allowbreak{} Almost sure(non-)injectivity


\url{https://arxiv.org/pdf/2602.04673}


Source statement date:\allowbreak{} 2026-08-24. Archived access:\allowbreak{} 2026-09-06T19:\allowbreak{}42:\allowbreak{}01Z.


Historical:\allowbreak{} source\_\allowbreak{}says\_\allowbreak{}open | Current:\allowbreak{} unverified | Result:\allowbreak{} attempted\_\allowbreak{}unsolved


\entrylabel{Question}

Determine whether the chronological attachment map from a typical SLE2 path and the intersected Brownian loops to the resulting Brownian path has a unique preimage almost surely,\allowbreak{} or instead infinitely many.


\entrylabel{Scope}

Global noninjectivity is already noted by the authors. The question concerns typical fibers in the specified support. It is related to,\allowbreak{} but stronger than,\allowbreak{} recovering only the SLE path measurably.


\subsection*{P0135-A001 | Almost-sure fibers of chronological loop addition structural reduction or obstruction}

\entrylabel{Argument or result}

Global noninjectivity cannot settle almost-sure invertibility. The map f:\allowbreak{}R->R with f(0)=\allowbreak{}f(1)=\allowbreak{}0 and f(x)=\allowbreak{}x otherwise is noninjective,\allowbreak{} yet it is injective off the two-point set and therefore injective on a full-measure set for a continuous input law. Conversely,\allowbreak{} x->|x| has two typical preimages for a symmetric continuous law.


\entrylabel{Second check}

A finite set has probability0 under an atomless law. For the absolute-value map,\allowbreak{} every y>0 has two preimages of equal conditional probability under a standard normal input.


\entrylabel{Remaining obstruction}

One must exhibit an ambiguity mechanism occurring on a positive-probability or full-probability family of loop configurations. A deterministic exceptional pair or a measurable selection theorem does not do this.


Verification:\allowbreak{} author checked limited argument | Independent review:\allowbreak{} none | Full target solved:\allowbreak{} False


\clearpage

\section*{P0136  Chronological Brownian loop addition in three dimensions}

\entrylabel{Primary source}

Berestycki and Sauzedde:\allowbreak{} Brownian paths as loop-decorated SLEs,\allowbreak{}2026;\allowbreak{} Section7,\allowbreak{} Dimension3


\url{https://arxiv.org/pdf/2602.04673}


Source statement date:\allowbreak{} 2026-08-24. Archived access:\allowbreak{} 2026-09-06T19:\allowbreak{}42:\allowbreak{}01Z.


Historical:\allowbreak{} source\_\allowbreak{}says\_\allowbreak{}open | Current:\allowbreak{} unverified | Result:\allowbreak{} attempted\_\allowbreak{}unsolved


\entrylabel{Question}

Extend the chronological loop-decoration representation to the three-dimensional scaling limit of loop-erased random walk,\allowbreak{} yielding the full parametrized Brownian path rather than only equality of random ranges.


\entrylabel{Scope}

The source explains that its planar deterministic continuity hypothesis fails in dimension3. A probabilistic continuity statement at typical configurations could still suffice.


\subsection*{P0136-A001 | Chronological Brownian loop addition in three dimensions structural reduction or obstruction}

\entrylabel{Argument or result}

A fixed planar circle intersecting a line in R\textasciicircum{}3 becomes disjoint from that line after an arbitrarily small translation in the perpendicular direction. For example,\allowbreak{} gamma is the x-axis segment[-2,\allowbreak{}2] and ell is the unit circle in the xy-plane;\allowbreak{} ell+\allowbreak{}(0,\allowbreak{}0,\allowbreak{}epsilon) avoids gamma for every epsilon nonzero while converging uniformly to ell.


\entrylabel{Second check}

Every point of the translated loop has z-coordinate epsilon and every point of gamma has z-coordinate0. At epsilon0 the circle meets gamma at(plus/\allowbreak{}minus1,\allowbreak{}0,\allowbreak{}0).


\entrylabel{Remaining obstruction}

Attaching or omitting a loop of fixed diameter changes the resulting path macroscopically. Typical random intersections need a stability argument beyond deterministic uniform closeness;\allowbreak{} the example does not refute such an almost-sure statement.


Verification:\allowbreak{} author checked limited argument | Independent review:\allowbreak{} none | Full target solved:\allowbreak{} False


\clearpage

\section*{P0137  A loop-erased representation of a unit Brownian loop}

\entrylabel{Primary source}

Berestycki and Sauzedde:\allowbreak{} Brownian paths as loop-decorated SLEs,\allowbreak{}2026;\allowbreak{} Section7,\allowbreak{} SLE2 loop analogue


\url{https://arxiv.org/pdf/2602.04673}


Source statement date:\allowbreak{} 2026-08-24. Archived access:\allowbreak{} 2026-09-06T19:\allowbreak{}42:\allowbreak{}01Z.


Historical:\allowbreak{} source\_\allowbreak{}says\_\allowbreak{}open | Current:\allowbreak{} unverified | Result:\allowbreak{} attempted\_\allowbreak{}unsolved


\entrylabel{Question}

Describe a unit Brownian loop as an SLE2-type loop decorated by an independent Brownian loop soup,\allowbreak{} with a precise law and chronology.


\entrylabel{Scope}

A closed loop has no distinguished endpoint ordering. Rooting conventions,\allowbreak{} duration conditioning,\allowbreak{} and the law of the proposed SLE loop must be fixed together.


\subsection*{P0137-A001 | A loop-erased representation of a unit Brownian loop structural reduction or obstruction}

\entrylabel{Argument or result}

Even in a finite closed walk,\allowbreak{} rooting before chronological loop erasure can discard all information:\allowbreak{} the standard loop erasure of a finite walk ending at its starting vertex is the single starting vertex. When the terminal return is processed,\allowbreak{} it erases the whole currently retained path.


\entrylabel{Second check}

The standard stack algorithm retains distinct vertices and,\allowbreak{} upon revisiting a vertex already on the stack,\allowbreak{} removes everything after its previous occurrence. The final return to the first vertex therefore leaves only that vertex.


\entrylabel{Remaining obstruction}

A nontrivial loop analogue cannot simply apply the open-path loop-erasure algorithm to an arbitrarily rooted closed Brownian path. A cyclic or marked-edge construction and its resulting sampling bias remain to be established.


Verification:\allowbreak{} author checked limited argument | Independent review:\allowbreak{} none | Full target solved:\allowbreak{} False


\clearpage

\section*{P0138  Law of loop-decorated SLE beyond kappa2}

\entrylabel{Primary source}

Berestycki and Sauzedde:\allowbreak{} Brownian paths as loop-decorated SLEs,\allowbreak{}2026;\allowbreak{} Section7,\allowbreak{} Other values of kappa


\url{https://arxiv.org/pdf/2602.04673}


Source statement date:\allowbreak{} 2026-08-24. Archived access:\allowbreak{} 2026-09-06T19:\allowbreak{}42:\allowbreak{}01Z.


Historical:\allowbreak{} source\_\allowbreak{}says\_\allowbreak{}open | Current:\allowbreak{} unverified | Result:\allowbreak{} attempted\_\allowbreak{}unsolved


\entrylabel{Question}

Identify the chronological path law obtained by decorating chordal SLE\_\allowbreak{}kappa with Brownian loops at restriction intensity lambda=\allowbreak{}(8-3*kappa)*(6-kappa)/\allowbreak{}(2*kappa),\allowbreak{} including its proposed interpretation as partially erased Brownian motion.


\entrylabel{Scope}

Only parameters giving a nonnegative soup intensity and an applicable attachment construction are meaningful. The interval2<=\allowbreak{}kappa<=\allowbreak{}8/\allowbreak{}3 interpolates the Brownian kappa2 endpoint and undecorated SLE8/\allowbreak{}3 endpoint.


\subsection*{P0138-A001 | Law of loop-decorated SLE beyond kappa2 structural reduction or obstruction}

\entrylabel{Argument or result}

At kappa2 the displayed intensity is1;\allowbreak{} at kappa8/\allowbreak{}3 it is0. Agreement of the filled hull with a restriction law cannot identify the ordered path:\allowbreak{} a circle traversed clockwise or counterclockwise has the same trace and hull,\allowbreak{} and arbitrary monotone time changes also preserve the hull.


\entrylabel{Second check}

Substitution gives lambda2=\allowbreak{}(8-6)*(6-2)/\allowbreak{}4=\allowbreak{}1 and lambda8/\allowbreak{}3=\allowbreak{}0. The two circle trajectories differ as parametrized maps while having identical ranges.


\entrylabel{Remaining obstruction}

Restriction formulas control avoidance of sets,\allowbreak{} not chronological finite-dimensional distributions or quadratic variation. A path-level stochastic characterization is missing.


Verification:\allowbreak{} author checked limited argument | Independent review:\allowbreak{} none | Full target solved:\allowbreak{} False


\clearpage

\section*{P0139  Rough-path continuity of chronological loop attachment}

\entrylabel{Primary source}

Berestycki and Sauzedde:\allowbreak{} Brownian paths as loop-decorated SLEs,\allowbreak{}2026;\allowbreak{} Section7,\allowbreak{} Continuity in rough path topologies


\url{https://arxiv.org/pdf/2602.04673}


Source statement date:\allowbreak{} 2026-08-24. Archived access:\allowbreak{} 2026-09-06T19:\allowbreak{}42:\allowbreak{}01Z.


Historical:\allowbreak{} source\_\allowbreak{}says\_\allowbreak{}open | Current:\allowbreak{} unverified | Result:\allowbreak{} attempted\_\allowbreak{}unsolved


\entrylabel{Question}

Prove appropriate Holder,\allowbreak{} p-variation,\allowbreak{} or rough-path continuity estimates for the chronological loop-attachment map,\allowbreak{} possibly restricted to the cusp conditions proposed by the authors.


\entrylabel{Scope}

The precise topology must include all necessary time and iterated-integral information. Continuity in the source's existing path topology does not imply continuity of rough-path lifts.


\subsection*{P0139-A001 | Rough-path continuity of chronological loop attachment structural reduction or obstruction}

\entrylabel{Argument or result}

Concatenate m positively oriented circles of radius m\textasciicircum{}(-1/\allowbreak{}2),\allowbreak{} each based at0,\allowbreak{} over a unit time interval. Their sup norm tends to0 and,\allowbreak{} for p>2,\allowbreak{} their first-level p-variation is at most C\_\allowbreak{}p*m\textasciicircum{}(1/\allowbreak{}p-1/\allowbreak{}2),\allowbreak{} which also tends to0. But their total signed area is m*pi*m\textasciicircum{}(-1)=\allowbreak{}pi.


\entrylabel{Second check}

Each circle contributes O(r\textasciicircum{}p) to the p-variation power,\allowbreak{} and splitting a partition at returns to0 changes the bound by only a p-dependent factor. Areas add over closed concatenated loops and equal pi*r\textasciicircum{}2 per circle.


\entrylabel{Remaining obstruction}

This proves that first-level convergence alone cannot control the second iterated integral. It does not rule out continuity with properly enhanced input loops or the source's additional cusp restrictions.


Verification:\allowbreak{} author checked limited argument | Independent review:\allowbreak{} none | Full target solved:\allowbreak{} False


\clearpage

\section*{P0141  First-order distance constants in finite-variance preferential attachment}

\entrylabel{Primary source}

Van der Hofstad:\allowbreak{} Stochastic processes on random graphs,\allowbreak{} Saint-Flour2017;\allowbreak{} Open Problem2.3,\allowbreak{} PDF p.74


\url{https://www.win.tue.nl/~rhofstad/SaintFlour_SPoRG_fin_Saint_Flour.pdf}


Source statement date:\allowbreak{} 2017-07-12. Archived access:\allowbreak{} 2026-09-06T19:\allowbreak{}31:\allowbreak{}10Z.


Historical:\allowbreak{} source\_\allowbreak{}says\_\allowbreak{}open | Current:\allowbreak{} unverified | Result:\allowbreak{} attempted\_\allowbreak{}unsolved


\entrylabel{Question}

For affine preferential attachment with m>=\allowbreak{}2 and delta>0,\allowbreak{} prove convergence of typical graph distance/\allowbreak{}log n and diameter/\allowbreak{}log n to explicit positive constants a<b.


\entrylabel{Scope}

Two-sided logarithmic bounds are weaker than convergence to constants. The positive-delta regime is the finite-variance degree regime;\allowbreak{} logarithmic-logarithmic formulas from negative delta do not apply.


\subsection*{P0141-A001 | First-order distance constants in finite-variance preferential attachment structural reduction}

\entrylabel{Argument or result}

For every connected graph,\allowbreak{} the distance D(U,\allowbreak{}V) between two uniform vertices is bounded by its diameter,\allowbreak{} so if the two normalized quantities converge in probability to constants a,\allowbreak{}b then necessarily a<=\allowbreak{}b. Strict inequality does not follow:\allowbreak{} a complete graph has typical distance and diameter both1 with probability tending to1.


\entrylabel{Second check}

If a>b,\allowbreak{} choose epsilon<(a-b)/\allowbreak{}2;\allowbreak{} simultaneous convergence would force D(U,\allowbreak{}V)>diameter with probability tending to1,\allowbreak{} a contradiction. The complete-graph calculation checks the strictness issue separately.


\entrylabel{Remaining obstruction}

This gives only the necessary weak ordering. The preferential-attachment age correlations and extreme peripheral vertices must be quantified to identify constants and show a strict gap.


Verification:\allowbreak{} author checked limited argument | Independent review:\allowbreak{} none | Full target solved:\allowbreak{} False


\clearpage

\section*{P0142  Fluctuations of preferential-attachment distances and diameter}

\entrylabel{Primary source}

Van der Hofstad:\allowbreak{} Stochastic processes on random graphs,\allowbreak{} Saint-Flour2017;\allowbreak{} Open Problem2.4,\allowbreak{} PDF p.76


\url{https://www.win.tue.nl/~rhofstad/SaintFlour_SPoRG_fin_Saint_Flour.pdf}


Source statement date:\allowbreak{} 2017-07-12. Archived access:\allowbreak{} 2026-09-06T19:\allowbreak{}31:\allowbreak{}10Z.


Historical:\allowbreak{} source\_\allowbreak{}says\_\allowbreak{}open | Current:\allowbreak{} unverified | Result:\allowbreak{} attempted\_\allowbreak{}unsolved


\entrylabel{Question}

Determine the centered fluctuation scales and limiting laws for typical distances and diameter in affine preferential-attachment graphs with m>=\allowbreak{}2.


\entrylabel{Scope}

Parameter regimes are one source problem,\allowbreak{} not separate corpus variants. Tightness and full-sequence convergence must be distinguished because graph distances are integer-valued.


\subsection*{P0142-A001 | Fluctuations of preferential-attachment distances and diameter structural reduction}

\entrylabel{Argument or result}

If integer-valued X\_\allowbreak{}n is centered by real c\_\allowbreak{}n,\allowbreak{} every centered value lies in Z-c\_\allowbreak{}n. Along a subsequence with fractional parts c\_\allowbreak{}n tending to r,\allowbreak{} any weak limit is supported on Z-r. Two subsequences whose fractional parts have different limits can therefore have incompatible translated-lattice limits.


\entrylabel{Second check}

For a closed set staying a positive distance from Z-r,\allowbreak{} its probability is eventually0;\allowbreak{} the portmanteau theorem puts any weak limit on Z-r. The deterministic example X\_\allowbreak{}n=\allowbreak{}floor(c\_\allowbreak{}n) gives centered value minus the fractional part.


\entrylabel{Remaining obstruction}

A claimed continuous full-sequence limit without additional scaling would require checking this lattice obstruction. No centered distance law for the actual attachment model is obtained.


Verification:\allowbreak{} author checked limited argument | Independent review:\allowbreak{} none | Full target solved:\allowbreak{} False


\clearpage

\section*{P0143  Second-order extremal hopcount in exponential complete-graph FPP}

\entrylabel{Primary source}

Van der Hofstad:\allowbreak{} Stochastic processes on random graphs,\allowbreak{} Saint-Flour2017;\allowbreak{} Open Problem3.1,\allowbreak{} PDF pp.99-100


\url{https://www.win.tue.nl/~rhofstad/SaintFlour_SPoRG_fin_Saint_Flour.pdf}


Source statement date:\allowbreak{} 2017-07-12. Archived access:\allowbreak{} 2026-09-06T19:\allowbreak{}31:\allowbreak{}10Z.


Historical:\allowbreak{} source\_\allowbreak{}says\_\allowbreak{}open | Current:\allowbreak{} unverified | Result:\allowbreak{} attempted\_\allowbreak{}unsolved


\entrylabel{Question}

Determine fluctuations of the longest minimum-weight path,\allowbreak{} both from a fixed root and over all vertex pairs,\allowbreak{} beyond the first-order e log n and alpha\_\allowbreak{}star log n scales.


\entrylabel{Scope}

The rooted shortest-path tree is a random recursive tree,\allowbreak{} for which substantial height refinements already exist. The all-pairs fluctuation target is retained;\allowbreak{} known rooted results are not new solutions.


\subsection*{P0143-A001 | Second-order extremal hopcount in exponential complete-graph FPP structural reduction}

\entrylabel{Argument or result}

For a random recursive tree,\allowbreak{} F\_\allowbreak{}n(z)=\allowbreak{}E sum\_\allowbreak{}v z\textasciicircum{}(depth(v)) equals product\_\allowbreak{}(j=\allowbreak{}1)\textasciicircum{}(n-1)(1+\allowbreak{}z/\allowbreak{}j). For z>1,\allowbreak{} P(height>=\allowbreak{}k)<=\allowbreak{}z\textasciicircum{}(-k)F\_\allowbreak{}n(z)<=\allowbreak{}exp(z H\_\allowbreak{}(n-1)-k log z). Taking k=\allowbreak{}a log n,\allowbreak{}z=\allowbreak{}a gives an upper exponent a-a log a. A union bound over n roots adds1,\allowbreak{} explaining alpha\_\allowbreak{}star log(alpha\_\allowbreak{}star)-alpha\_\allowbreak{}star=\allowbreak{}1.


\entrylabel{Second check}

When vertex n is added,\allowbreak{} its parent is uniform among n-1 vertices,\allowbreak{} so the expected generating function is multiplied by1+\allowbreak{}z/\allowbreak{}(n-1). The bound log(1+\allowbreak{}z/\allowbreak{}j)<=\allowbreak{}z/\allowbreak{}j yields the display.


\entrylabel{Remaining obstruction}

This recovers a first-order upper bound only. It loses dependence between extreme paths and the logarithmic corrections needed for a fluctuation law.


Verification:\allowbreak{} author checked limited argument | Independent review:\allowbreak{} none | Full target solved:\allowbreak{} False


\clearpage

\section*{P0144  Conservative weighted-distance fluctuations in heavy-tailed configuration models}

\entrylabel{Primary source}

Van der Hofstad:\allowbreak{} Stochastic processes on random graphs,\allowbreak{} Saint-Flour2017;\allowbreak{} Open Problem3.2,\allowbreak{} PDF p.142


\url{https://www.win.tue.nl/~rhofstad/SaintFlour_SPoRG_fin_Saint_Flour.pdf}


Source statement date:\allowbreak{} 2017-07-12. Archived access:\allowbreak{} 2026-09-06T19:\allowbreak{}31:\allowbreak{}10Z.


Historical:\allowbreak{} source\_\allowbreak{}says\_\allowbreak{}open | Current:\allowbreak{} unverified | Result:\allowbreak{} attempted\_\allowbreak{}unsolved


\entrylabel{Question}

Describe the fluctuations of the minimal total weight between two typical vertices of an infinite-variance configuration model when the associated weighted branching process is conservative.


\entrylabel{Scope}

The source explicitly says the question concerns random fluctuations beyond the leading quantile-sum scale. Later leading-order weighted-distance results must not be mistaken for a complete centered limiting law.


\subsection*{P0144-A001 | Conservative weighted-distance fluctuations in heavy-tailed configuration models structural reduction}

\entrylabel{Argument or result}

For nonnegative weights,\allowbreak{} a path sum is at least its maximum edge,\allowbreak{} so the FPP distance dominates the minimax path value. The latter is exactly the smallest level t at which the two vertices are connected using edges of weight at most t. This follows by comparing the same path family under the two objectives.


\entrylabel{Second check}

A connection at level t supplies a path with every edge at most t,\allowbreak{} and any path with maximum edge at most t supplies such a connection;\allowbreak{} taking infima proves the minimax/\allowbreak{}percolation identity.


\entrylabel{Remaining obstruction}

The lower bound tracks the largest obstruction but discards accumulated small costs. In the conservative regime infinitely many layers can contribute,\allowbreak{} so the minimax comparison does not determine centered sum fluctuations.


Verification:\allowbreak{} author checked limited argument | Independent review:\allowbreak{} none | Full target solved:\allowbreak{} False


\clearpage

\section*{P0145  Hopcount universality in infinite-variance weighted configuration models}

\entrylabel{Primary source}

Van der Hofstad:\allowbreak{} Stochastic processes on random graphs,\allowbreak{} Saint-Flour2017;\allowbreak{} Open Problem3.3,\allowbreak{} PDF p.142


\url{https://www.win.tue.nl/~rhofstad/SaintFlour_SPoRG_fin_Saint_Flour.pdf}


Source statement date:\allowbreak{} 2017-07-12. Archived access:\allowbreak{} 2026-09-06T19:\allowbreak{}31:\allowbreak{}10Z.


Historical:\allowbreak{} source\_\allowbreak{}says\_\allowbreak{}open | Current:\allowbreak{} unverified | Result:\allowbreak{} attempted\_\allowbreak{}unsolved


\entrylabel{Question}

Determine the scaling and possible central limit laws for the number of edges on a minimum-weight path between typical vertices of an infinite-variance configuration model.


\entrylabel{Scope}

The target is hopcount,\allowbreak{} not total weight or unweighted distance. Separate edge-weight laws may lead to different behaviors but are one classification problem.


\subsection*{P0145-A001 | Hopcount universality in infinite-variance weighted configuration models structural reduction}

\entrylabel{Argument or result}

If all weights lie in[a,\allowbreak{}b] with0<a<=\allowbreak{}b,\allowbreak{} the hopcount H of a minimizing path and graph distance d satisfy d<=\allowbreak{}H<=\allowbreak{}(b/\allowbreak{}a)d:\allowbreak{} a shortest unweighted path costs at most bd,\allowbreak{} whereas the chosen path costs at least aH. Monotone transformations do not preserve hopcount:\allowbreak{} paths with weights(1,\allowbreak{}4) and(1.8,\allowbreak{}1.8,\allowbreak{}1.8) reverse their cost ordering when weights are squared.


\entrylabel{Second check}

The original costs are5 and5.4;\allowbreak{} the squared costs are17 and9.72. Both paths can be realized as the two branches of a finite cycle,\allowbreak{} so the comparison is exact.


\entrylabel{Remaining obstruction}

A bounded-factor comparison gives no fluctuation law,\allowbreak{} and it disappears for weights approaching0. The exploration process must jointly control accumulated cost and generation depth.


Verification:\allowbreak{} author checked limited argument | Independent review:\allowbreak{} none | Full target solved:\allowbreak{} False


\clearpage

\section*{P0146  Weak-to-strong disorder universality in complete-graph FPP}

\entrylabel{Primary source}

Van der Hofstad:\allowbreak{} Stochastic processes on random graphs,\allowbreak{} Saint-Flour2017;\allowbreak{} Open Problem3.4,\allowbreak{} PDF p.149 and discussion p.154


\url{https://www.win.tue.nl/~rhofstad/SaintFlour_SPoRG_fin_Saint_Flour.pdf}


Source statement date:\allowbreak{} 2017-07-12. Archived access:\allowbreak{} 2026-09-06T19:\allowbreak{}31:\allowbreak{}10Z.


Historical:\allowbreak{} source\_\allowbreak{}says\_\allowbreak{}open | Current:\allowbreak{} unverified | Result:\allowbreak{} attempted\_\allowbreak{}unsolved


\entrylabel{Question}

Complete the classification of complete-graph first-passage percolation with varying edge-weight distributions,\allowbreak{} including the strong-disorder regime related to the minimum spanning tree.


\entrylabel{Scope}

The source highlights weights E\_\allowbreak{}e raised to s\_\allowbreak{}n and the regimes not covered by its theorem. A value comparison with a bottleneck problem does not prove convergence of path geometry.


\subsection*{P0146-A001 | Weak-to-strong disorder universality in complete-graph FPP structural reduction}

\entrylabel{Argument or result}

For any positive edge labels a\_\allowbreak{}e and any simple path of at most n-1 edges,\allowbreak{} max a\_\allowbreak{}e\textasciicircum{}s <=\allowbreak{} sum a\_\allowbreak{}e\textasciicircum{}s <=\allowbreak{}(n-1)max a\_\allowbreak{}e\textasciicircum{}s. Taking infima over paths and s-th roots shows that the FPP value lies between the minimax value M and(n-1)\textasciicircum{}(1/\allowbreak{}s)M. If s/\allowbreak{}log n tends to infinity,\allowbreak{} this ratio bound tends to1.


\entrylabel{Second check}

The lower bound holds for every path;\allowbreak{} applying the upper bound to a minimax path gives the reverse comparison. No exchange of a random limit and an infimum is used.


\entrylabel{Remaining obstruction}

Nearly equal objective values may be attained by very different paths. This deterministic estimate does not identify the hopcount crossover or a metric-space scaling limit.


Verification:\allowbreak{} author checked limited argument | Independent review:\allowbreak{} none | Full target solved:\allowbreak{} False


\clearpage

\section*{P0147  Fluctuations of weighted flooding and diameter on configuration models}

\entrylabel{Primary source}

Van der Hofstad:\allowbreak{} Stochastic processes on random graphs,\allowbreak{} Saint-Flour2017;\allowbreak{} Open Problem3.5,\allowbreak{} PDF pp.150-151


\url{https://www.win.tue.nl/~rhofstad/SaintFlour_SPoRG_fin_Saint_Flour.pdf}


Source statement date:\allowbreak{} 2017-07-12. Archived access:\allowbreak{} 2026-09-06T19:\allowbreak{}31:\allowbreak{}10Z.


Historical:\allowbreak{} source\_\allowbreak{}says\_\allowbreak{}open | Current:\allowbreak{} unverified | Result:\allowbreak{} attempted\_\allowbreak{}unsolved


\entrylabel{Question}

Identify centered fluctuation laws for exponential-weight flooding time and weighted diameter in the configuration-model regimes of Theorem3.37.


\entrylabel{Scope}

The source assumes a convergent degree law with a uniformly bounded2+\allowbreak{}epsilon moment and positive mass at the minimum degree. If the graph is disconnected,\allowbreak{} distances must be restricted to the relevant giant component;\allowbreak{} infinite distances are not silently included.


\subsection*{P0147-A001 | Fluctuations of weighted flooding and diameter on configuration models structural reduction}

\entrylabel{Argument or result}

On any fixed graph,\allowbreak{} the minimum incident exponential edge weight at a degree-d vertex is Exp(d). For a collection of m pairwise nonadjacent degree-d vertices these minima use disjoint edge sets and are independent. Hence P(d max X\_\allowbreak{}v-log m<=\allowbreak{}x)=\allowbreak{}(1-exp(-x)/\allowbreak{}m)\textasciicircum{}m tends to exp(-exp(-x)).


\entrylabel{Second check}

The minimum of d independent rate-one exponentials has tail exp(-dt). Nonadjacent vertices cannot share an incident edge. Multiplying their distribution functions proves the extreme-value limit.


\entrylabel{Remaining obstruction}

This identifies one endpoint-delay component of the extremes. The bulk connecting distances and their dependence on the extreme endpoints are not controlled,\allowbreak{} so no full flooding or diameter limit is claimed.


Verification:\allowbreak{} author checked limited argument | Independent review:\allowbreak{} none | Full target solved:\allowbreak{} False


\clearpage

\section*{P0148  General edge-weight laws for random-graph flooding and diameter}

\entrylabel{Primary source}

Van der Hofstad:\allowbreak{} Stochastic processes on random graphs,\allowbreak{} Saint-Flour2017;\allowbreak{} Open Problem3.6,\allowbreak{} PDF p.152


\url{https://www.win.tue.nl/~rhofstad/SaintFlour_SPoRG_fin_Saint_Flour.pdf}


Source statement date:\allowbreak{} 2017-07-12. Archived access:\allowbreak{} 2026-09-06T19:\allowbreak{}31:\allowbreak{}10Z.


Historical:\allowbreak{} source\_\allowbreak{}says\_\allowbreak{}open | Current:\allowbreak{} unverified | Result:\allowbreak{} attempted\_\allowbreak{}unsolved


\entrylabel{Question}

Extend weighted flooding-time and diameter asymptotics from exponential weights to general iid nonnegative edge-weight laws on the configuration model.


\entrylabel{Scope}

Both the lower tail used by fast paths and the upper tail of the minimum weight leaving a peripheral vertex can matter. This is one law-classification problem.


\subsection*{P0148-A001 | General edge-weight laws for random-graph flooding and diameter structural reduction}

\entrylabel{Argument or result}

At a degree-d vertex the minimum incident weight X has tail P(X>t)=\allowbreak{}(1-F(t))\textasciicircum{}d. Over m nonadjacent such vertices,\allowbreak{} P(max X\_\allowbreak{}v<=\allowbreak{}t)=\allowbreak{}[1-(1-F(t))\textasciicircum{}d]\textasciicircum{}m. For weights E\textasciicircum{}s this gives max X\_\allowbreak{}v asymptotic to(log m/\allowbreak{}d)\textasciicircum{}s in probability.


\entrylabel{Second check}

For E\textasciicircum{}s,\allowbreak{} P(X>t)=\allowbreak{}exp(-d*t\textasciicircum{}(1/\allowbreak{}s));\allowbreak{} setting t=\allowbreak{}((1 plus/\allowbreak{}minus epsilon)log m/\allowbreak{}d)\textasciicircum{}s and using the displayed maximum distribution gives matching upper and lower probability bounds.


\entrylabel{Remaining obstruction}

Endpoint delays alone can change scale,\allowbreak{} but the fastest route through the bulk is a different optimization. Matching these two mechanisms for general F is still missing.


Verification:\allowbreak{} author checked limited argument | Independent review:\allowbreak{} none | Full target solved:\allowbreak{} False


\clearpage

\section*{P0149  Extreme hopcounts for general random-graph degree and weight laws}

\entrylabel{Primary source}

Van der Hofstad:\allowbreak{} Stochastic processes on random graphs,\allowbreak{} Saint-Flour2017;\allowbreak{} Open Problem3.7,\allowbreak{} PDF p.152


\url{https://www.win.tue.nl/~rhofstad/SaintFlour_SPoRG_fin_Saint_Flour.pdf}


Source statement date:\allowbreak{} 2017-07-12. Archived access:\allowbreak{} 2026-09-06T19:\allowbreak{}31:\allowbreak{}10Z.


Historical:\allowbreak{} source\_\allowbreak{}says\_\allowbreak{}open | Current:\allowbreak{} unverified | Result:\allowbreak{} attempted\_\allowbreak{}unsolved


\entrylabel{Question}

Determine the first-order extreme hopcount from one root and over all pairs for general configuration-model degree laws,\allowbreak{} and extend beyond exponential edge weights.


\entrylabel{Scope}

The known regular-graph constants use a branching-generation large-deviation calculation. Replacing the degree by its mean is not justified for extremes.


\subsection*{P0149-A001 | Extreme hopcounts for general random-graph degree and weight laws structural reduction}

\entrylabel{Argument or result}

On a graph with maximum degree Delta,\allowbreak{} there are at most n*Delta*(Delta-1)\textasciicircum{}(k-1) oriented simple k-edge paths. With iid weights Y,\allowbreak{} the chance any such path has total weight at most t is at most this count times exp(a*t)*(E exp(-aY))\textasciicircum{}k for every a>0.


\entrylabel{Second check}

After the first edge there are at most Delta-1 nonbacktracking choices. Independence along a simple path and exponential Markov inequality bound its lower-tail cost. The union bound needs no independence between different paths.


\entrylabel{Remaining obstruction}

The bound is too crude when a few hubs make Delta large and does not automatically constrain minimizing paths with unknown total cost. Degree-biased branching and extreme endpoint selection require sharper control.


Verification:\allowbreak{} author checked limited argument | Independent review:\allowbreak{} none | Full target solved:\allowbreak{} False


\clearpage

\section*{P0150  Minimal assumptions for heavy-tailed critical metric convergence}

\entrylabel{Primary source}

Van der Hofstad:\allowbreak{} Stochastic processes on random graphs,\allowbreak{} Saint-Flour2017;\allowbreak{} Open Problem4.2,\allowbreak{} PDF p.218;\allowbreak{} Theorem4.39


\url{https://www.win.tue.nl/~rhofstad/SaintFlour_SPoRG_fin_Saint_Flour.pdf}


Source statement date:\allowbreak{} 2017-07-12. Archived access:\allowbreak{} 2026-09-06T19:\allowbreak{}31:\allowbreak{}10Z.


Historical:\allowbreak{} source\_\allowbreak{}says\_\allowbreak{}open | Current:\allowbreak{} unverified | Result:\allowbreak{} attempted\_\allowbreak{}unsolved


\entrylabel{Question}

Remove the extra degree restrictions used to upgrade heavy-tailed critical configuration-model components from Gromov-weak convergence to Gromov-Hausdorff-Prokhorov convergence.


\entrylabel{Scope}

The printed problem heading says finite third moment,\allowbreak{} but its referenced theorem and surrounding text concern infinite third moment,\allowbreak{} tau in(3,\allowbreak{}4). Bhamidi et al2020 improve the global lower-mass bound under slightly stronger assumptions;\allowbreak{} the minimal-assumption remainder requires exact comparison.


\subsection*{P0150-A001 | Minimal assumptions for heavy-tailed critical metric convergence structural reduction}

\entrylabel{Argument or result}

Let X\_\allowbreak{}n=\allowbreak{}[0,\allowbreak{}1] with the ordinary metric and probability measure(1-1/\allowbreak{}n)delta\_\allowbreak{}0+\allowbreak{}(1/\allowbreak{}n)Lebesgue. Every fixed finite iid sample from this measure lies at0 with probability tending to1,\allowbreak{} so its distance matrices converge to those of a point. Yet the spaces all have diameter1 and cannot converge in Gromov-Hausdorff distance to a point.


\entrylabel{Second check}

The probability that k sampled points all select the delta\_\allowbreak{}0 component is(1-1/\allowbreak{}n)\textasciicircum{}k. The Gromov-Hausdorff distance from a compact space to a point is half its diameter. Each X\_\allowbreak{}n still has full-support measure.


\entrylabel{Remaining obstruction}

This is a precise failure of weak measured convergence to control thin appendages. A uniform lower bound on masses of fixed-radius balls is the missing type of estimate.


Verification:\allowbreak{} author checked limited argument | Independent review:\allowbreak{} none | Full target solved:\allowbreak{} False


\clearpage

\section*{P0151  Critical window for fixed-outdegree preferential attachment}

\entrylabel{Primary source}

Van der Hofstad:\allowbreak{} Stochastic processes on random graphs,\allowbreak{} Saint-Flour2017;\allowbreak{} Open Problem4.4,\allowbreak{} PDF p.220


\url{https://www.win.tue.nl/~rhofstad/SaintFlour_SPoRG_fin_Saint_Flour.pdf}


Source statement date:\allowbreak{} 2017-07-12. Archived access:\allowbreak{} 2026-09-06T19:\allowbreak{}31:\allowbreak{}10Z.


Historical:\allowbreak{} source\_\allowbreak{}says\_\allowbreak{}open | Current:\allowbreak{} unverified | Result:\allowbreak{} attempted\_\allowbreak{}unsolved


\entrylabel{Question}

For affine preferential attachment with a fixed m>=\allowbreak{}2 edges per arriving vertex and delta>0,\allowbreak{} identify the percolation scaling window and largest critical component sizes.


\entrylabel{Scope}

The percolation threshold itself is covered by later work2312.14085. The2025 growing-graph window theorem concerns models with independently attached edges;\allowbreak{} it is not automatically the fixed-outdegree model asked here.


\subsection*{P0151-A001 | Critical window for fixed-outdegree preferential attachment structural reduction}

\entrylabel{Argument or result}

Local weak convergence alone cannot determine macroscopic cluster sizes. A single cycle on n vertices and a disjoint union of cycles whose lengths are about sqrt(n) have the same local weak limit,\allowbreak{} the two-sided infinite path,\allowbreak{} but their largest components have orders n and sqrt(n).


\entrylabel{Second check}

For every fixed radius r,\allowbreak{} once each cycle length exceeds2r+\allowbreak{}1 all rooted radius-r neighborhoods are identical. Their total component sizes remain explicitly different.


\entrylabel{Remaining obstruction}

This prevents deriving the critical window from the Polya-point local limit alone. A growing-scale exploration argument must retain age and fixed-outdegree dependencies.


Verification:\allowbreak{} author checked limited argument | Independent review:\allowbreak{} none | Full target solved:\allowbreak{} False


\clearpage

\section*{P0152  Kernel criteria for critical inhomogeneous random-graph universality}

\entrylabel{Primary source}

Van der Hofstad:\allowbreak{} Stochastic processes on random graphs,\allowbreak{} Saint-Flour2017;\allowbreak{} Open Problem4.5-4.6,\allowbreak{} PDF pp.222-223


\url{https://www.win.tue.nl/~rhofstad/SaintFlour_SPoRG_fin_Saint_Flour.pdf}


Source statement date:\allowbreak{} 2017-07-12. Archived access:\allowbreak{} 2026-09-06T19:\allowbreak{}31:\allowbreak{}10Z.


Historical:\allowbreak{} source\_\allowbreak{}says\_\allowbreak{}open | Current:\allowbreak{} unverified | Result:\allowbreak{} attempted\_\allowbreak{}unsolved


\entrylabel{Question}

Characterize critical component universality for general inhomogeneous random-graph kernels,\allowbreak{} including the proposed role of the positive principal eigenfunction belonging to L3.


\entrylabel{Scope}

The general question and the eigenfunction criterion are grouped. Rank-one and finite-type theorems do not cover every irreducible unbounded kernel.


\subsection*{P0152-A001 | Kernel criteria for critical inhomogeneous random-graph universality structural reduction}

\entrylabel{Argument or result}

L3 integrability of the principal eigenfunction does not itself give a spectral gap. Partition[0,\allowbreak{}1] into blocks of masses p\_\allowbreak{}i tending to0,\allowbreak{} choose epsilon\_\allowbreak{}i>0 tending to0,\allowbreak{} and let m=\allowbreak{}sum p\_\allowbreak{}i epsilon\_\allowbreak{}i. On blocks i,\allowbreak{}j set kappa=\allowbreak{}(1-epsilon\_\allowbreak{}i)1\{i=\allowbreak{}j\}/\allowbreak{}p\_\allowbreak{}i+\allowbreak{}epsilon\_\allowbreak{}i epsilon\_\allowbreak{}j/\allowbreak{}m. This symmetric positive kernel has row integral1,\allowbreak{} so its principal eigenfunction is the constant1 in L3. For the centered indicator of block i,\allowbreak{} the Rayleigh quotient tends to1 as epsilon\_\allowbreak{}i tends to0.


\entrylabel{Second check}

The row integral is1-epsilon\_\allowbreak{}i+\allowbreak{}epsilon\_\allowbreak{}i=\allowbreak{}1. The Dirichlet energy of the block indicator is at most p\_\allowbreak{}i epsilon\_\allowbreak{}i,\allowbreak{} while its centered L2 norm squared is p\_\allowbreak{}i(1-p\_\allowbreak{}i). Thus the spectral gap is at most epsilon\_\allowbreak{}i/\allowbreak{}(1-p\_\allowbreak{}i),\allowbreak{} tending to0. Positivity of the kernel makes the constant eigenfunction unique up to scale.


\entrylabel{Remaining obstruction}

The example is an obstruction to using an unstated spectral-gap consequence of the conjectured criterion. It is not asserted to refute random-graph universality;\allowbreak{} the component scaling for this explicit kernel remains to be analyzed.


Verification:\allowbreak{} author checked limited argument | Independent review:\allowbreak{} none | Full target solved:\allowbreak{} False


\clearpage

\section*{P0153  Vacant-set percolation threshold on general configuration models}

\entrylabel{Primary source}

Van der Hofstad:\allowbreak{} Stochastic processes on random graphs,\allowbreak{} Saint-Flour2017;\allowbreak{} Open Problem4.8,\allowbreak{} PDF p.225


\url{https://www.win.tue.nl/~rhofstad/SaintFlour_SPoRG_fin_Saint_Flour.pdf}


Source statement date:\allowbreak{} 2017-07-12. Archived access:\allowbreak{} 2026-09-06T19:\allowbreak{}31:\allowbreak{}10Z.


Historical:\allowbreak{} source\_\allowbreak{}says\_\allowbreak{}open | Current:\allowbreak{} unverified | Result:\allowbreak{} attempted\_\allowbreak{}unsolved


\entrylabel{Question}

Identify the critical time for percolation of the vacant set left after a random walk runs for order n time on a configuration model,\allowbreak{} including degree laws with diverging second moment.


\entrylabel{Scope}

Time normalization,\allowbreak{} the giant component,\allowbreak{} and initial distribution must match the source. The suggestion that the threshold vanishes for infinite second moment is a question,\allowbreak{} not an established consequence of ordinary percolation.


\subsection*{P0153-A001 | Vacant-set percolation threshold on general configuration models structural reduction}

\entrylabel{Argument or result}

For simple random walk on a finite connected graph,\allowbreak{} stationary mass at v is d\_\allowbreak{}v/\allowbreak{}sum\_\allowbreak{}w d\_\allowbreak{}w. This controls expected occupation counts in stationarity,\allowbreak{} but it does not identify the probability of ever visiting v:\allowbreak{} repeated returns can account for many occupation times. In particular P(visit v by time t)<=\allowbreak{}E[number of visits by time t],\allowbreak{} with no matching lower bound from this mean alone.


\entrylabel{Second check}

The stationary law follows from detailed balance pi\_\allowbreak{}u/\allowbreak{}d\_\allowbreak{}u=\allowbreak{}pi\_\allowbreak{}v/\allowbreak{}d\_\allowbreak{}v on each edge. The displayed visit bound is Markov's inequality for the nonnegative integer occupation count.


\entrylabel{Remaining obstruction}

Degree-biased occupation cannot be substituted for independent vertex deletion. One needs escape probabilities and joint avoidance estimates to determine the critical remaining branching structure.


Verification:\allowbreak{} author checked limited argument | Independent review:\allowbreak{} none | Full target solved:\allowbreak{} False


\clearpage

\section*{P0154  Metric scaling of critical random-walk vacant components}

\entrylabel{Primary source}

Van der Hofstad:\allowbreak{} Stochastic processes on random graphs,\allowbreak{} Saint-Flour2017;\allowbreak{} Open Problem4.9,\allowbreak{} PDF p.225


\url{https://www.win.tue.nl/~rhofstad/SaintFlour_SPoRG_fin_Saint_Flour.pdf}


Source statement date:\allowbreak{} 2017-07-12. Archived access:\allowbreak{} 2026-09-06T19:\allowbreak{}31:\allowbreak{}10Z.


Historical:\allowbreak{} source\_\allowbreak{}says\_\allowbreak{}open | Current:\allowbreak{} unverified | Result:\allowbreak{} attempted\_\allowbreak{}unsolved


\entrylabel{Question}

Obtain metric-space scaling limits for the largest critical vacant-set components of general configuration models,\allowbreak{} including heavy-tailed degree sequences.


\entrylabel{Scope}

Random-regular vacant-set metric limits are already established. The general degree law and possible suppression of high-degree hubs form the unresolved extension.


\subsection*{P0154-A001 | Metric scaling of critical random-walk vacant components structural reduction}

\entrylabel{Argument or result}

A high degree need not imply large escape capacity. Attach K leaves to a vertex v that has just one additional edge leading into a fixed transient infinite tree. Every walk step from v into a leaf returns to v with probability1;\allowbreak{} only the remaining edge can contribute to escape. Thus d\_\allowbreak{}v*P\_\allowbreak{}v(no return after time0) is at most1,\allowbreak{} uniformly in K.


\entrylabel{Second check}

Condition on the first step from v. Each leaf contributes0 to the sum of neighbor escape probabilities,\allowbreak{} and the one transient branch contributes at most1. This is the usual unit-conductance singleton capacity calculation.


\entrylabel{Remaining obstruction}

The example is not a configuration-model limit theorem. It shows why deleting hubs based only on degree can fail:\allowbreak{} escape geometry,\allowbreak{} not stationary occupation alone,\allowbreak{} controls avoidance.


Verification:\allowbreak{} author checked limited argument | Independent review:\allowbreak{} none | Full target solved:\allowbreak{} False


\clearpage

\section*{P0155  Erdos-Renyi component-size limits for high-dimensional tori}

\entrylabel{Primary source}

Van der Hofstad:\allowbreak{} Stochastic processes on random graphs,\allowbreak{} Saint-Flour2017;\allowbreak{} Open Problem4.10,\allowbreak{} PDF p.226


\url{https://www.win.tue.nl/~rhofstad/SaintFlour_SPoRG_fin_Saint_Flour.pdf}


Source statement date:\allowbreak{} 2017-07-12. Archived access:\allowbreak{} 2026-09-06T19:\allowbreak{}31:\allowbreak{}10Z.


Historical:\allowbreak{} source\_\allowbreak{}says\_\allowbreak{}open | Current:\allowbreak{} unverified | Result:\allowbreak{} attempted\_\allowbreak{}unsolved


\entrylabel{Question}

For nearest-neighbor percolation on growing d-dimensional tori with fixed d>6 in the critical window around p\_\allowbreak{}c(Z\textasciicircum{}d),\allowbreak{} prove convergence of the ordered component sizes to the appropriate shifted and rescaled multiplicative-coalescent law.


\entrylabel{Scope}

Correct window centering and constants are part of the source question. Results for a growing-dimensional hypercube,\allowbreak{} or order-of-magnitude estimates,\allowbreak{} are not automatically this fixed-d torus theorem.


\subsection*{P0155-A001 | Erdos-Renyi component-size limits for high-dimensional tori structural reduction}

\entrylabel{Argument or result}

Two-sided order bounds cannot identify a scaling law:\allowbreak{} random variables X\_\allowbreak{}n=\allowbreak{}n\textasciicircum{}(2/\allowbreak{}3)A\_\allowbreak{}n,\allowbreak{} where A\_\allowbreak{}n alternates between two different bounded positive distributions,\allowbreak{} always have the predicted magnitude and are tight after scaling,\allowbreak{} but have different subsequential limits.


\entrylabel{Second check}

Take A\_\allowbreak{}n identically1 for even n and identically2 for odd n. The normalized variables have no weak limit despite uniform upper and lower deterministic bounds.


\entrylabel{Remaining obstruction}

For the percolation model one must identify the joint exploration process and the window shift,\allowbreak{} rather than only establish component-size tightness. No such identification is proved here.


Verification:\allowbreak{} author checked limited argument | Independent review:\allowbreak{} none | Full target solved:\allowbreak{} False


\clearpage

\section*{P0156  Configuration-model minimum spanning-tree metric limits}

\entrylabel{Primary source}

Van der Hofstad:\allowbreak{} Stochastic processes on random graphs,\allowbreak{} Saint-Flour2017;\allowbreak{} Open Problem4.11-4.12,\allowbreak{} PDF p.231


\url{https://www.win.tue.nl/~rhofstad/SaintFlour_SPoRG_fin_Saint_Flour.pdf}


Source statement date:\allowbreak{} 2017-07-12. Archived access:\allowbreak{} 2026-09-06T19:\allowbreak{}31:\allowbreak{}10Z.


Historical:\allowbreak{} source\_\allowbreak{}says\_\allowbreak{}open | Current:\allowbreak{} unverified | Result:\allowbreak{} attempted\_\allowbreak{}unsolved


\entrylabel{Question}

Identify and prove the metric-measure scaling limits of minimum spanning trees on configuration models with minimum limiting degree at least3,\allowbreak{} across finite-third-moment and heavy-tailed degree regimes.


\entrylabel{Scope}

The two source regimes are grouped. The random3-regular case is resolved by Addario-Berry and Sen1810.03802;\allowbreak{} it is excluded from the open remainder. The law of iid continuous edge weights is not an extra family of distinct problems.


\subsection*{P0156-A001 | Configuration-model minimum spanning-tree metric limits structural reduction}

\entrylabel{Argument or result}

For a fixed finite graph with distinct edge weights,\allowbreak{} applying any strictly increasing function to every weight leaves the minimum spanning tree unchanged. Kruskal's algorithm inspects edges in the same order and makes the same cycle decisions. Consequently,\allowbreak{} with iid continuous weights,\allowbreak{} the distribution of the unweighted MST geometry depends only on a uniform random edge ordering.


\entrylabel{Second check}

An iid continuous law gives no ties almost surely and every ordering has probability1/\allowbreak{}|E|!. The greedy cycle criterion depends only on the ordering,\allowbreak{} which independently verifies distributional invariance.


\entrylabel{Remaining obstruction}

The remaining difficulty is the graph ensemble and the global near-critical component-merger geometry. Rank invariance does not establish tightness,\allowbreak{} scaling exponents,\allowbreak{} or the limiting tree.


Verification:\allowbreak{} author checked limited argument | Independent review:\allowbreak{} none | Full target solved:\allowbreak{} False


\clearpage

\section*{P0157  Two competing largest-component scales in barely subcritical configuration models}

\entrylabel{Primary source}

Van der Hofstad:\allowbreak{} Stochastic processes on random graphs,\allowbreak{} Saint-Flour2017;\allowbreak{} Open Problem4.13,\allowbreak{} PDF p.234


\url{https://www.win.tue.nl/~rhofstad/SaintFlour_SPoRG_fin_Saint_Flour.pdf}


Source statement date:\allowbreak{} 2017-07-12. Archived access:\allowbreak{} 2026-09-06T19:\allowbreak{}31:\allowbreak{}10Z.


Historical:\allowbreak{} source\_\allowbreak{}says\_\allowbreak{}open | Current:\allowbreak{} unverified | Result:\allowbreak{} attempted\_\allowbreak{}unsolved


\entrylabel{Question}

Under convergence of third degree moments,\allowbreak{} prove that the largest barely subcritical configuration-model component has order max\{d\_\allowbreak{}max/\allowbreak{}epsilon\_\allowbreak{}n,\allowbreak{} epsilon\_\allowbreak{}n\textasciicircum{}(-2)log(n epsilon\_\allowbreak{}n\textasciicircum{}3)\} in the regime specified by the source.


\entrylabel{Scope}

The asymptotic regime must have epsilon\_\allowbreak{}n tending to0 with n epsilon\_\allowbreak{}n\textasciicircum{}3 large;\allowbreak{} otherwise the displayed logarithm is not the intended barely-subcritical scale. Infinite-third-moment results are a different regime.


\subsection*{P0157-A001 | Two competing largest-component scales in barely subcritical configuration models structural reduction}

\entrylabel{Argument or result}

A branching-process approximation already explains the high-degree candidate scale. If the root has d children and every later vertex has mean offspring1-epsilon,\allowbreak{} then expected total progeny is1+\allowbreak{}d/\allowbreak{}epsilon:\allowbreak{} each child starts a subcritical tree of expected size1/\allowbreak{}epsilon.


\entrylabel{Second check}

Generation expectations are d(1-epsilon)\textasciicircum{}(k-1) for k>=\allowbreak{}1;\allowbreak{} Tonelli and the geometric series yield1+\allowbreak{}d/\allowbreak{}epsilon.


\entrylabel{Remaining obstruction}

Expectation at a high-degree root is not a high-probability largest-component estimate. Uniform concentration,\allowbreak{} collision errors,\allowbreak{} and the competing extreme size of many ordinary components remain to be controlled.


Verification:\allowbreak{} author checked limited argument | Independent review:\allowbreak{} none | Full target solved:\allowbreak{} False


\clearpage

\section*{P0158  Ordered-phase susceptibility exponent on random trees and graphs}

\entrylabel{Primary source}

Van der Hofstad:\allowbreak{} Stochastic processes on random graphs,\allowbreak{} Saint-Flour2017;\allowbreak{} Open Problem5.1,\allowbreak{} PDF pp.286-287


\url{https://www.win.tue.nl/~rhofstad/SaintFlour_SPoRG_fin_Saint_Flour.pdf}


Source statement date:\allowbreak{} 2017-07-12. Archived access:\allowbreak{} 2026-09-06T19:\allowbreak{}31:\allowbreak{}10Z.


Historical:\allowbreak{} source\_\allowbreak{}says\_\allowbreak{}open | Current:\allowbreak{} unverified | Result:\allowbreak{} attempted\_\allowbreak{}unsolved


\entrylabel{Question}

Prove that the susceptibility critical exponent gamma\_\allowbreak{}prime on the ordered side of the locally tree-like ferromagnetic Ising transition equals1 under the source's degree assumptions.


\entrylabel{Scope}

The high-temperature exponent gamma=\allowbreak{}1 and the random-regular ordered-phase case are already known. The source seeks the missing upper bound for general degree laws with a nontrivial finite critical inverse temperature.


\subsection*{P0158-A001 | Ordered-phase susceptibility exponent on random trees and graphs structural reduction}

\entrylabel{Argument or result}

At zero cavity field,\allowbreak{} the expected generation-k contribution to susceptibility is proportional to nu\textasciicircum{}(k-1)*(tanh beta)\textasciicircum{}k,\allowbreak{} yielding a geometric series only when nu*tanh beta<1. Above criticality that same bound diverges,\allowbreak{} so it cannot prove the desired finite susceptibility asymptotic in the plus state.


\entrylabel{Second check}

The ratio of consecutive generation terms is nu*tanh beta,\allowbreak{} which equals1 at criticality and exceeds1 on the ordered side. Nonzero random cavity fields supply extra attenuation factors in the exact source formula.


\entrylabel{Remaining obstruction}

A sharp lower bound on that attenuation,\allowbreak{} averaged along the size-biased spine and uniformly near criticality,\allowbreak{} is missing. Analytically continuing the high-temperature geometric sum would give an invalid negative denominator.


Verification:\allowbreak{} author checked limited argument | Independent review:\allowbreak{} none | Full target solved:\allowbreak{} False


\clearpage

\section*{P0159  Quenched non-Gaussian critical spin limits on heterogeneous random graphs}

\entrylabel{Primary source}

Van der Hofstad:\allowbreak{} Stochastic processes on random graphs,\allowbreak{} Saint-Flour2017;\allowbreak{} Open Problem5.2,\allowbreak{} PDF p.305


\url{https://www.win.tue.nl/~rhofstad/SaintFlour_SPoRG_fin_Saint_Flour.pdf}


Source statement date:\allowbreak{} 2017-07-12. Archived access:\allowbreak{} 2026-09-06T19:\allowbreak{}31:\allowbreak{}10Z.


Historical:\allowbreak{} source\_\allowbreak{}says\_\allowbreak{}open | Current:\allowbreak{} unverified | Result:\allowbreak{} attempted\_\allowbreak{}unsolved


\entrylabel{Question}

Extend the critical non-classical total-spin limit theorem from annealed generalized random graphs to quenched configuration models or quenched generalized random graphs.


\entrylabel{Scope}

Quenched means fixing the sampled graph before taking its spin law. Matching critical exponents or annealed pressure does not establish a quenched distributional limit.


\subsection*{P0159-A001 | Quenched non-Gaussian critical spin limits on heterogeneous random graphs structural reduction}

\entrylabel{Argument or result}

For environment G and partition functions Z\_\allowbreak{}G(h),\allowbreak{} the averaged quenched moment generating function is E[Z\_\allowbreak{}G(h)/\allowbreak{}Z\_\allowbreak{}G(0)],\allowbreak{} whereas the annealed one is E Z\_\allowbreak{}G(h)/\allowbreak{}E Z\_\allowbreak{}G(0). These are generally different. If two equally likely environments have(Z0,\allowbreak{}Zh)=\allowbreak{}(1,\allowbreak{}1) and(2,\allowbreak{}2*cosh h),\allowbreak{} the expressions are(1+\allowbreak{}cosh h)/\allowbreak{}2 and(1+\allowbreak{}2*cosh h)/\allowbreak{}3.


\entrylabel{Second check}

The two finite Boltzmann systems can be realized by one zero-spin state versus two states with spins plus/\allowbreak{}minus1. The formulas then follow directly from normalization. This is a normalization example,\allowbreak{} not a graph-model counterexample.


\entrylabel{Remaining obstruction}

The annealed integral identity weights graph environments by their partition functions. Removing this tilt while preserving critical fluctuations is the unresolved step.


Verification:\allowbreak{} author checked limited argument | Independent review:\allowbreak{} none | Full target solved:\allowbreak{} False


\clearpage

\section*{P0160  Fixed-temperature metastability for general configuration models}

\entrylabel{Primary source}

Van der Hofstad:\allowbreak{} Stochastic processes on random graphs,\allowbreak{} Saint-Flour2017;\allowbreak{} Open Problem5.4,\allowbreak{} PDF p.313


\url{https://www.win.tue.nl/~rhofstad/SaintFlour_SPoRG_fin_Saint_Flour.pdf}


Source statement date:\allowbreak{} 2017-07-12. Archived access:\allowbreak{} 2026-09-06T19:\allowbreak{}31:\allowbreak{}10Z.


Historical:\allowbreak{} source\_\allowbreak{}says\_\allowbreak{}open | Current:\allowbreak{} unverified | Result:\allowbreak{} attempted\_\allowbreak{}unsolved


\entrylabel{Question}

Prove metastable behavior of low-temperature ferromagnetic Ising dynamics on configuration models as n tends to infinity with inverse temperature fixed and large.


\entrylabel{Scope}

The source contrasts this with taking temperature to0 at fixed n. Later random-regular metastability results are known;\allowbreak{} the full general configuration-model extension is retained.


\subsection*{P0160-A001 | Fixed-temperature metastability for general configuration models structural reduction}

\entrylabel{Argument or result}

A deterministic expansion condition gives an energy barrier. If every vertex set S of size at most n/\allowbreak{}2 has at least h|S| boundary edges,\allowbreak{} and the Hamiltonian is-J sum\_\allowbreak{}edges sigma\_\allowbreak{}u sigma\_\allowbreak{}v-B sum\_\allowbreak{}v sigma\_\allowbreak{}v,\allowbreak{} then changing S from minus to plus costs2J|boundary S|-2B|S|,\allowbreak{} at least2(Jh-B)|S|. Every single-spin path from all minus to all plus crosses a set of size floor(n/\allowbreak{}2).


\entrylabel{Second check}

Each boundary edge changes its interaction contribution by2J and each flipped spin changes its field contribution by-2B. The number of plus spins changes by1 per step,\allowbreak{} forcing the intermediate set size.


\entrylabel{Remaining obstruction}

A positive energy barrier is not by itself a fixed-temperature exit-time theorem because exponentially many configurations contribute entropy. One must compare beta times the barrier with entropy and control the metastable basin,\allowbreak{} with expansion estimates valid for the sampled graph.


Verification:\allowbreak{} author checked limited argument | Independent review:\allowbreak{} none | Full target solved:\allowbreak{} False


\subsection*{P0160-A002 | Energy plus entropy conductance bound}

\entrylabel{Argument or result}

Let G have odd n and edge expansion |boundary S|>=\allowbreak{}h|S| for all |S|<=\allowbreak{}n/\allowbreak{}2. Use H(sigma)=\allowbreak{}-J sum\_\allowbreak{}edges sigma\_\allowbreak{}u sigma\_\allowbreak{}v-B sum\_\allowbreak{}v sigma\_\allowbreak{}v with J>0,\allowbreak{}B>=\allowbreak{}0,\allowbreak{}Jh>B,\allowbreak{} and Gibbs law pi proportional to exp(-beta H). Consider any reversible single-spin continuous-time chain with each individual spin-flip rate at most1. Put A=\allowbreak{}\{number of plus spins< n/\allowbreak{}2\}. Global spin reversal maps A bijectively onto its complement and does not decrease Gibbs weight,\allowbreak{} so pi(A)<=\allowbreak{}1/\allowbreak{}2. An A-to-complement jump must start with(n-1)/\allowbreak{}2 plus spins,\allowbreak{} whose energy exceeds the all-minus energy by at least(Jh-B)(n-1). There are at most2\textasciicircum{}n such states and total exit rate at most n per state. Since pi(A) is at least the all-minus weight,\allowbreak{} the conductance Phi=\allowbreak{}Q(A,\allowbreak{}A-complement)/\allowbreak{}pi(A) is at most n*2\textasciicircum{}n*exp[-beta(Jh-B)(n-1)]. Starting with pi conditioned on A,\allowbreak{} the time-s law is bounded above by pi/\allowbreak{}pi(A),\allowbreak{} by Markov-kernel positivity and stationarity. Consequently the expected number of A-to-complement jumps by time t is at most t*Phi. The probability of still being in A at time t is at least1-t*Phi,\allowbreak{} giving total variation distance at least1/\allowbreak{}2-t*Phi. Convexity of total variation implies the same lower bound for the worst starting state. Therefore t\_\allowbreak{}mix(1/\allowbreak{}4)>=\allowbreak{}1/\allowbreak{}(4*Phi),\allowbreak{} at least exp[beta(Jh-B)(n-1)]/\allowbreak{}(4*n*2\textasciicircum{}n). This is exponential in n when beta(Jh-B)>log2. The calculation accounts explicitly for entropy. It is a deterministic expansion-based mixing lower bound;\allowbreak{} it does not prove a metastable exit-time distribution,\allowbreak{} or expansion for every configuration-model degree law.


\entrylabel{Second check}

Boundary-layer counting and stationary domination give separate energy and time estimates.


\entrylabel{Remaining obstruction}

No general random-graph expansion theorem or metastable exit-time distribution is proved.


Verification:\allowbreak{} author checked limited argument | Independent review:\allowbreak{} none | Full target solved:\allowbreak{} False


\clearpage

\section*{P0161  Potts free energy for general configuration-model degrees}

\entrylabel{Primary source}

Van der Hofstad:\allowbreak{} Stochastic processes on random graphs,\allowbreak{} Saint-Flour2017;\allowbreak{} Open Problem5.5,\allowbreak{} PDF p.314


\url{https://www.win.tue.nl/~rhofstad/SaintFlour_SPoRG_fin_Saint_Flour.pdf}


Source statement date:\allowbreak{} 2017-07-12. Archived access:\allowbreak{} 2026-09-06T19:\allowbreak{}31:\allowbreak{}10Z.


Historical:\allowbreak{} source\_\allowbreak{}says\_\allowbreak{}open | Current:\allowbreak{} unverified | Result:\allowbreak{} attempted\_\allowbreak{}unsolved


\entrylabel{Question}

Extend the ferromagnetic Potts Bethe free-energy and phase-diagram results from regular locally tree-like graphs to general configuration-model degree laws.


\entrylabel{Scope}

Can and van der Hofstad2503.17636 establish additional regular-graph and parameter-region results. A lower Bethe bound is not equality throughout the general heterogeneous phase diagram.


\subsection*{P0161-A001 | Potts free energy for general configuration-model degrees structural reduction}

\entrylabel{Argument or result}

For a finite forest at zero field,\allowbreak{} Z=\allowbreak{}q\textasciicircum{}(number of components)*(exp(beta)+\allowbreak{}q-1)\textasciicircum{}(number of edges),\allowbreak{} by successively summing leaf colors. Applying that formula to a d-regular locally tree-like graph as if it were a forest predicts pressure log q+\allowbreak{}(d/\allowbreak{}2)log[(exp(beta)+\allowbreak{}q-1)/\allowbreak{}q]. For d>2 and large beta this is below beta*d/\allowbreak{}2,\allowbreak{} while the q monochromatic states alone imply pressure at least beta*d/\allowbreak{}2+\allowbreak{}log(q)/\allowbreak{}n.


\entrylabel{Second check}

Leaf elimination gives one factor exp(beta)+\allowbreak{}q-1 per edge and q for each remaining component root. In a connected graph the monochromatic states have weight exp(beta*|E|),\allowbreak{} so their contribution is an exact lower bound.


\entrylabel{Remaining obstruction}

Local tree neighborhoods do not erase the linear number of global cycle constraints. Correct Bethe fixed-point selection,\allowbreak{} beyond the symmetric forest expression,\allowbreak{} remains essential.


Verification:\allowbreak{} author checked limited argument | Independent review:\allowbreak{} none | Full target solved:\allowbreak{} False


\clearpage

\section*{P0162  Fourier entropy versus total influence}

\entrylabel{Primary source}

O Donnell:\allowbreak{} Open problems in analysis of Boolean functions;\allowbreak{} PDF p.5,\allowbreak{} Fourier Entropy-Influence Conjecture


\url{https://arxiv.org/pdf/1204.6447}


Source statement date:\allowbreak{} 2012-04-29. Archived access:\allowbreak{} 2026-09-06T18:\allowbreak{}55:\allowbreak{}14Z.


Historical:\allowbreak{} source\_\allowbreak{}says\_\allowbreak{}open | Current:\allowbreak{} unverified | Result:\allowbreak{} attempted\_\allowbreak{}unsolved


\entrylabel{Question}

Bound the Shannon entropy of the squared Fourier coefficients of every Boolean function by a universal constant times its total influence.


\entrylabel{Scope}

The weaker min-entropy formulation is grouped here. Gonzalez,\allowbreak{}MacManus,\allowbreak{}Pereyra2606.00246 still calls the full statement a conjecture and proves additional classes.


\subsection*{P0162-A001 | Fourier entropy versus total influence exact calculation or structural obstruction}

\entrylabel{Argument or result}

Discarding the Boolean constraint destroys the claim:\allowbreak{} the probability measure uniform on the n singleton subsets has entropy log2(n) but mean subset size1. It cannot be the squared spectrum of a Boolean function when n>=\allowbreak{}2.


\entrylabel{Second check}

If all Fourier mass were on singletons,\allowbreak{} f=\allowbreak{}sum\_\allowbreak{}i a\_\allowbreak{}i*x\_\allowbreak{}i with a\_\allowbreak{}i=\allowbreak{}plus/\allowbreak{}minus1/\allowbreak{}sqrt(n). In f\textasciicircum{}2 the coefficient of x\_\allowbreak{}i*x\_\allowbreak{}j is2*a\_\allowbreak{}i*a\_\allowbreak{}j,\allowbreak{} nonzero for distinct i,\allowbreak{}j,\allowbreak{} whereas Boolean f requires f\textasciicircum{}2=\allowbreak{}1. This independently exhibits the violated Fourier convolution identity.


\entrylabel{Remaining obstruction}

Entropy versus mean cardinality alone is insufficient. Any proof must use the nonlinear identities imposed by f\textasciicircum{}2=\allowbreak{}1;\allowbreak{} no dimension-free inequality exploiting all of them is obtained.


Verification:\allowbreak{} author checked limited argument | Independent review:\allowbreak{} none | Full target solved:\allowbreak{} False


\clearpage

\section*{P0163  Polynomial Fourier concentration for DNF formulas}

\entrylabel{Primary source}

O Donnell:\allowbreak{} Open problems in analysis of Boolean functions;\allowbreak{} PDF p.4,\allowbreak{} Mansour Conjecture


\url{https://arxiv.org/pdf/1204.6447}


Source statement date:\allowbreak{} 2012-04-29. Archived access:\allowbreak{} 2026-09-06T18:\allowbreak{}55:\allowbreak{}14Z.


Historical:\allowbreak{} source\_\allowbreak{}says\_\allowbreak{}open | Current:\allowbreak{} unverified | Result:\allowbreak{} attempted\_\allowbreak{}unsolved


\entrylabel{Question}

For a size-s DNF,\allowbreak{} concentrate all but epsilon of its squared Fourier mass on at most s raised to O(log(1/\allowbreak{}epsilon)) coefficients.


\entrylabel{Scope}

The weaker polynomial-in-s formulation is grouped. Lecomte-Tan2109.04525 improve read-dependent bounds,\allowbreak{} not the full assertion.


\subsection*{P0163-A001 | Polynomial Fourier concentration for DNF formulas exact calculation or structural obstruction}

\entrylabel{Argument or result}

Let C\_\allowbreak{}j be the indicator of clause j,\allowbreak{} and encode the DNF by f=\allowbreak{}1-2*product\_\allowbreak{}j(1-C\_\allowbreak{}j). Each consistent conjunction has Fourier l1 norm1,\allowbreak{} and any product of conjunction indicators is either another conjunction or zero. Inclusion-exclusion gives L=\allowbreak{}sum\_\allowbreak{}S |fhat(S)|<=\allowbreak{}2\textasciicircum{}(s+\allowbreak{}1)-1. Keeping coefficients of magnitude at least epsilon/\allowbreak{}L leaves squared tail at most epsilon and uses at most L\textasciicircum{}2/\allowbreak{}epsilon coefficients.


\entrylabel{Second check}

For discarded coefficients,\allowbreak{} sum a\_\allowbreak{}S\textasciicircum{}2 <=\allowbreak{}(epsilon/\allowbreak{}L)*sum |a\_\allowbreak{}S|<=\allowbreak{}epsilon. For retained coefficients,\allowbreak{} their number times epsilon/\allowbreak{}L is at most L. Thus the explicit exponential bound is valid.


\entrylabel{Remaining obstruction}

This handles every DNF but is exponentially larger than the requested polynomial bound. The expansion ignores cancellations among overlapping clauses.


Verification:\allowbreak{} author checked limited argument | Independent review:\allowbreak{} none | Full target solved:\allowbreak{} False


\clearpage

\section*{P0164  Gaussian lower bound for noise stability of unbiased thresholds}

\entrylabel{Primary source}

Kamath,\allowbreak{}Kumar,\allowbreak{}Manurangsi:\allowbreak{} When Majority Fails,\allowbreak{}2026;\allowbreak{} Section1.4,\allowbreak{} Conjecture1;\allowbreak{} O-Donnell2012 p.5 weaker consequences


\url{https://arxiv.org/pdf/2604.06590}


Source statement date:\allowbreak{} 2026-04-08. Archived access:\allowbreak{} 2026-09-06T20:\allowbreak{}45:\allowbreak{}41Z.


Historical:\allowbreak{} source\_\allowbreak{}says\_\allowbreak{}open | Current:\allowbreak{} unverified | Result:\allowbreak{} attempted\_\allowbreak{}unsolved


\entrylabel{Question}

For an unbiased linear threshold f and rho in(0,\allowbreak{}1),\allowbreak{} prove Stab\_\allowbreak{}rho(f)>=\allowbreak{}(2/\allowbreak{}pi)*arcsin(rho).


\entrylabel{Scope}

This is the authors refined conjecture,\allowbreak{} not the false finite-majority comparison. The degree-one Fourier lower bound2/\allowbreak{}pi and small-noise consequence are grouped as parts of this family.


\subsection*{P0164-A001 | Gaussian lower bound for noise stability of unbiased thresholds exact calculation or structural obstruction}

\entrylabel{Argument or result}

For f=\allowbreak{}sign(S),\allowbreak{} S=\allowbreak{}sum\_\allowbreak{}i w\_\allowbreak{}i X\_\allowbreak{}i with sum w\_\allowbreak{}i\textasciicircum{}2=\allowbreak{}1 and no zero ties,\allowbreak{} Holder gives E|S|>=\allowbreak{}(E S\textasciicircum{}2)\textasciicircum{}(3/\allowbreak{}2)/\allowbreak{}(E S\textasciicircum{}4)\textasciicircum{}(1/\allowbreak{}2)>=\allowbreak{}1/\allowbreak{}sqrt(3),\allowbreak{} since E S\textasciicircum{}4=\allowbreak{}3-2*sum w\_\allowbreak{}i\textasciicircum{}4<=\allowbreak{}3. Cauchy-Schwarz applied to E[fS]=\allowbreak{}sum\_\allowbreak{}i w\_\allowbreak{}i*fhat(i) yields sum\_\allowbreak{}i fhat(i)\textasciicircum{}2>=\allowbreak{}1/\allowbreak{}3. Therefore Stab\_\allowbreak{}rho(f)>=\allowbreak{}rho/\allowbreak{}3.


\entrylabel{Second check}

Interpolate L2 between L1 and L4:\allowbreak{} 1=\allowbreak{}||S||2 <=\allowbreak{}||S||1\textasciicircum{}(1/\allowbreak{}3)||S||4\textasciicircum{}(2/\allowbreak{}3). Fourier noise stability is sum\_\allowbreak{}S rho\textasciicircum{}|S|*fhat(S)\textasciicircum{}2,\allowbreak{} with nonnegative terms for rho>0. This verifies the bound without assuming regular weights.


\entrylabel{Remaining obstruction}

The elementary bound is weaker even than the established robust Khintchine bounds and far below the proposed arcsine curve. A central limit argument gives the arcsine law only when all normalized weights vanish;\allowbreak{} it does not compare arbitrary nonregular thresholds.


Verification:\allowbreak{} author checked limited argument | Independent review:\allowbreak{} none | Full target solved:\allowbreak{} False


\clearpage

\section*{P0165  Optimal odd functions for many-player correlation distillation}

\entrylabel{Primary source}

O Donnell:\allowbreak{} Open problems in analysis of Boolean functions;\allowbreak{} PDF p.6,\allowbreak{} Optimality of Majorities for Non-Interactive Correlation Distillation


\url{https://arxiv.org/pdf/1204.6447}


Source statement date:\allowbreak{} 2012-04-29. Archived access:\allowbreak{} 2026-09-06T18:\allowbreak{}55:\allowbreak{}14Z.


Historical:\allowbreak{} source\_\allowbreak{}says\_\allowbreak{}open | Current:\allowbreak{} unverified | Result:\allowbreak{} attempted\_\allowbreak{}unsolved


\entrylabel{Question}

When r independent noisy copies of a common uniform bit string use the same odd Boolean function,\allowbreak{} is the agreement probability maximized by majority on some odd subset of coordinates?


\entrylabel{Scope}

The number of players and the noise parameter are one optimization problem. This differs from two-party erasure distillation and from minimizing threshold stability.


\subsection*{P0165-A001 | Optimal odd functions for many-player correlation distillation exact calculation or structural obstruction}

\entrylabel{Argument or result}

Write rho=\allowbreak{}1-2*epsilon and m(x)=\allowbreak{}T\_\allowbreak{}rho f(x). Conditional independence gives agreement probability E[((1+\allowbreak{}m)/\allowbreak{}2)\textasciicircum{}r+\allowbreak{}((1-m)/\allowbreak{}2)\textasciicircum{}r]. For r=\allowbreak{}2 this is(1+\allowbreak{}E m\textasciicircum{}2)/\allowbreak{}2,\allowbreak{} and for r=\allowbreak{}3 it is(1+\allowbreak{}3E m\textasciicircum{}2)/\allowbreak{}4. Since odd f is unbiased,\allowbreak{} Parseval gives E m\textasciicircum{}2=\allowbreak{}sum\_\allowbreak{}nonemptyS rho\textasciicircum{}(2|S|)*fhat(S)\textasciicircum{}2<=\allowbreak{}rho\textasciicircum{}2. Dictators attain both bounds.


\entrylabel{Second check}

Expand the two binomials;\allowbreak{} odd powers cancel. A dictator has m=\allowbreak{}rho*x\_\allowbreak{}1 and equality in the Fourier bound. For r=\allowbreak{}4 the expression becomes(1+\allowbreak{}6E m\textasciicircum{}2+\allowbreak{}E m\textasciicircum{}4)/\allowbreak{}8,\allowbreak{} displaying a genuinely additional moment.


\entrylabel{Remaining obstruction}

The exact r=\allowbreak{}2,\allowbreak{}3 calculation does not optimize the higher moments jointly for r>=\allowbreak{}4. No argument shows that a subset-majority maximizes all those combinations.


Verification:\allowbreak{} author checked limited argument | Independent review:\allowbreak{} none | Full target solved:\allowbreak{} False


\clearpage

\section*{P0166  Influential coordinates of bounded low-degree polynomials}

\entrylabel{Primary source}

O Donnell:\allowbreak{} Open problems in analysis of Boolean functions;\allowbreak{} PDF p.7,\allowbreak{} Aaronson-Ambainis Conjecture


\url{https://arxiv.org/pdf/1204.6447}


Source statement date:\allowbreak{} 2012-04-29. Archived access:\allowbreak{} 2026-09-06T18:\allowbreak{}55:\allowbreak{}14Z.


Historical:\allowbreak{} source\_\allowbreak{}says\_\allowbreak{}open | Current:\allowbreak{} unverified | Result:\allowbreak{} attempted\_\allowbreak{}unsolved


\entrylabel{Question}

For f on the uniform discrete cube taking values in[-1,\allowbreak{}1] and having degree at most d,\allowbreak{} prove some L2 influence is at least(Var(f)/\allowbreak{}d)\textasciicircum{}C for a universal C.


\entrylabel{Scope}

Bhattacharya2402.13952 and its March2026 revision concern random restrictions. That does not settle every original polynomial.


\subsection*{P0166-A001 | Influential coordinates of bounded low-degree polynomials exact calculation or structural obstruction}

\entrylabel{Argument or result}

For d=\allowbreak{}1,\allowbreak{} write f=\allowbreak{}a0+\allowbreak{}sum\_\allowbreak{}i a\_\allowbreak{}i*x\_\allowbreak{}i. Evaluating at coordinate signs that align with a\_\allowbreak{}i and at their opposites gives |a0|+\allowbreak{}sum\_\allowbreak{}i|a\_\allowbreak{}i|<=\allowbreak{}1. Hence Var(f)=\allowbreak{}sum\_\allowbreak{}i a\_\allowbreak{}i\textasciicircum{}2<=\allowbreak{}max\_\allowbreak{}i|a\_\allowbreak{}i|,\allowbreak{} so max\_\allowbreak{}i Inf\_\allowbreak{}i(f)=\allowbreak{}max\_\allowbreak{}i a\_\allowbreak{}i\textasciicircum{}2>=\allowbreak{}Var(f)\textasciicircum{}2.


\entrylabel{Second check}

The two extreme cube values are a0 plus/\allowbreak{}minus sum|a\_\allowbreak{}i|. L2 influence for an affine polynomial is exactly a\_\allowbreak{}i\textasciicircum{}2. The proof is dimension-free and uses the full range constraint.


\entrylabel{Remaining obstruction}

In higher degree,\allowbreak{} cancellation among monomials prevents this direct l1-coefficient bound. The argument gives no polynomial-in-d control of the largest influence.


Verification:\allowbreak{} author checked limited argument | Independent review:\allowbreak{} none | Full target solved:\allowbreak{} False


\clearpage

\section*{P0167  Minimum correlation of symmetric Gaussian sets}

\entrylabel{Primary source}

O Donnell:\allowbreak{} Open problems in analysis of Boolean functions;\allowbreak{} PDF p.8,\allowbreak{} Symmetric Gaussian Problem


\url{https://arxiv.org/pdf/1204.6447}


Source statement date:\allowbreak{} 2012-04-29. Archived access:\allowbreak{} 2026-09-06T18:\allowbreak{}55:\allowbreak{}14Z.


Historical:\allowbreak{} source\_\allowbreak{}says\_\allowbreak{}open | Current:\allowbreak{} unverified | Result:\allowbreak{} attempted\_\allowbreak{}unsolved


\entrylabel{Question}

For correlated standard Gaussian vectors with prescribed rho and set measures mu,\allowbreak{}nu,\allowbreak{} minimize P(X in A,\allowbreak{}Y in B) under central symmetry A=\allowbreak{}-A.


\entrylabel{Scope}

The proposed centered-ball optimizer is not the general optimization question. Heilman1511.00382 gives negative results for ball/\allowbreak{}complement maximizers in dimension at least2;\allowbreak{} this does not identify every optimum.


\subsection*{P0167-A001 | Minimum correlation of symmetric Gaussian sets exact calculation or structural obstruction}

\entrylabel{Argument or result}

For fixed A,\allowbreak{} put h=\allowbreak{}T\_\allowbreak{}rho 1\_\allowbreak{}A. The objective is integral\_\allowbreak{}B h dgamma. Its minimum over sets of measure nu is attained by a lower-level set of h,\allowbreak{} with a suitable fraction of a level set if necessary. Since A is symmetric,\allowbreak{} h is even,\allowbreak{} so B can also be chosen symmetric.


\entrylabel{Second check}

If B selects positive measure at a higher value of h while omitting equal measure at a lower value,\allowbreak{} exchanging the two decreases the integral. Applying the threshold at the nu-quantile proves the characterization. Gaussian measure is nonatomic,\allowbreak{} including after passing to antipodal pairs.


\entrylabel{Remaining obstruction}

This eliminates one set but leaves a nonlinear optimization over A. Symmetry alone does not make its smoothed indicator radial or force a ball.


Verification:\allowbreak{} author checked limited argument | Independent review:\allowbreak{} none | Full target solved:\allowbreak{} False


\clearpage

\section*{P0168  Noise-optimal equal-volume Gaussian partitions}

\entrylabel{Primary source}

O Donnell:\allowbreak{} Open problems in analysis of Boolean functions;\allowbreak{} PDF p.9,\allowbreak{} Standard Simplex Conjecture


\url{https://arxiv.org/pdf/1204.6447}


Source statement date:\allowbreak{} 2012-04-29. Archived access:\allowbreak{} 2026-09-06T18:\allowbreak{}55:\allowbreak{}14Z.


Historical:\allowbreak{} source\_\allowbreak{}says\_\allowbreak{}open | Current:\allowbreak{} unverified | Result:\allowbreak{} attempted\_\allowbreak{}unsolved


\entrylabel{Question}

Prove that equal-volume regular-simplex partitions maximize Gaussian noise stability for positive correlation and minimize it for negative correlation,\allowbreak{} for every admissible number of cells.


\entrylabel{Scope}

The signs of correlation and cell numbers form one source conjecture. Heilman2608.00333 claims a substantial three-cell correlation range;\allowbreak{} it does not cover all cell counts or every positive correlation.


\subsection*{P0168-A001 | Noise-optimal equal-volume Gaussian partitions exact calculation or structural obstruction}

\entrylabel{Argument or result}

For four cells in dimension at least3,\allowbreak{} partition by the signs of the first two coordinates. Its stability is[1/\allowbreak{}2+\allowbreak{}arcsin(rho)/\allowbreak{}pi]\textasciicircum{}2,\allowbreak{} because the two Gaussian coordinate pairs are independent and each pair has sign-agreement probability1/\allowbreak{}2+\allowbreak{}arcsin(rho)/\allowbreak{}pi.


\entrylabel{Second check}

Represent a correlated pair as(G1,\allowbreak{}rho*G1+\allowbreak{}sqrt(1-rho\textasciicircum{}2)*G2). Rotational symmetry in the plane gives sign-disagreement angle arccos(rho)/\allowbreak{}pi. The product of the two coordinate agreement events supplies the displayed four-cell benchmark.


\entrylabel{Remaining obstruction}

This computes a non-simplex competitor exactly but does not order it against every partition. At rho=\allowbreak{}0 every equal-volume partition has stability1/\allowbreak{}q,\allowbreak{} so independence is not a uniqueness argument.


Verification:\allowbreak{} author checked limited argument | Independent review:\allowbreak{} none | Full target solved:\allowbreak{} False


\clearpage

\section*{P0169  Linear Fourier mass versus Boolean degree}

\entrylabel{Primary source}

O Donnell:\allowbreak{} Open problems in analysis of Boolean functions;\allowbreak{} PDF p.9,\allowbreak{} Linear Coefficients versus Total Degree


\url{https://arxiv.org/pdf/1204.6447}


Source statement date:\allowbreak{} 2012-04-29. Archived access:\allowbreak{} 2026-09-06T18:\allowbreak{}55:\allowbreak{}14Z.


Historical:\allowbreak{} source\_\allowbreak{}says\_\allowbreak{}open | Current:\allowbreak{} unverified | Result:\allowbreak{} attempted\_\allowbreak{}unsolved


\entrylabel{Question}

Bound the sum of degree-one Fourier coefficients of a degree-d Boolean function by sqrt(d),\allowbreak{} or by the sharper majority-on-d-bits value proposed in the source.


\entrylabel{Scope}

The stronger and weaker proposed constants are grouped. Jha1611.01029 and Wang2020/\allowbreak{}002 discuss partial cases and reformulations;\allowbreak{} neither title is treated as a complete resolution.


\subsection*{P0169-A001 | Linear Fourier mass versus Boolean degree exact calculation or structural obstruction}

\entrylabel{Argument or result}

Under independent input bits with common mean t,\allowbreak{} F(t)=\allowbreak{}E\_\allowbreak{}t f is a polynomial of degree at most d bounded by1 in absolute value for t in[-1,\allowbreak{}1]. Its derivative at0 is sum\_\allowbreak{}i fhat(i). However,\allowbreak{} boundedness and degree alone cannot force the desired sqrt(d) bound:\allowbreak{} for odd d the Chebyshev polynomial T\_\allowbreak{}d is bounded by1 on this interval and |T\_\allowbreak{}d prime(0)|=\allowbreak{}d.


\entrylabel{Second check}

The biased expectation of the character on S is t\textasciicircum{}|S|. Also T\_\allowbreak{}d(t)=\allowbreak{}cos(d arccos t),\allowbreak{} whose derivative at0 has magnitude d for odd d.


\entrylabel{Remaining obstruction}

A proof through one-variable biased expectations must exploit additional constraints inherited from a Boolean function. Chebyshev is an obstruction to the relaxed method,\allowbreak{} not a Boolean counterexample.


Verification:\allowbreak{} author checked limited argument | Independent review:\allowbreak{} none | Full target solved:\allowbreak{} False


\clearpage

\section*{P0170  Optimal limited independence for Gaussian polynomial thresholds}

\entrylabel{Primary source}

O Donnell:\allowbreak{} Open problems in analysis of Boolean functions;\allowbreak{} PDF p.9,\allowbreak{} k-wise Independence for PTFs


\url{https://arxiv.org/pdf/1204.6447}


Source statement date:\allowbreak{} 2012-04-29. Archived access:\allowbreak{} 2026-09-06T18:\allowbreak{}55:\allowbreak{}14Z.


Historical:\allowbreak{} source\_\allowbreak{}says\_\allowbreak{}open | Current:\allowbreak{} unverified | Result:\allowbreak{} attempted\_\allowbreak{}unsolved


\entrylabel{Question}

Determine the smallest k(d,\allowbreak{}epsilon) for which every k-wise independent standard-Gaussian vector fools every degree-d polynomial threshold to error epsilon.


\entrylabel{Scope}

The joint distribution of every k coordinates must be product Gaussian. Matching finitely many mixed moments alone is a weaker property and is not substituted for this assumption.


\subsection*{P0170-A001 | Optimal limited independence for Gaussian polynomial thresholds exact calculation or structural obstruction}

\entrylabel{Argument or result}

Take n=\allowbreak{}k+\allowbreak{}1 signs uniformly distributed subject to product\_\allowbreak{}i sign\_\allowbreak{}i=\allowbreak{}1,\allowbreak{} independent of n independent absolute standard Gaussians R\_\allowbreak{}i. Then X\_\allowbreak{}i=\allowbreak{}sign\_\allowbreak{}i*R\_\allowbreak{}i has k-wise independent standard-Gaussian coordinates. Yet product\_\allowbreak{}i X\_\allowbreak{}i>0 almost surely,\allowbreak{} while the same product for fully independent Gaussians is positive with probability1/\allowbreak{}2.


\entrylabel{Second check}

Every assignment of any k signs has exactly one extension satisfying the parity constraint,\allowbreak{} so its probability is2\textasciicircum{}-k. Independent magnitudes then yield independent Gaussian marginals on that subset. The distinguishing polynomial has degree k+\allowbreak{}1.


\entrylabel{Remaining obstruction}

For epsilon<1/\allowbreak{}2 this proves the necessary bound k>=\allowbreak{}d. It is weaker than known quantitative lower bounds and supplies no sharp dependence on epsilon.


Verification:\allowbreak{} author checked limited argument | Independent review:\allowbreak{} none | Full target solved:\allowbreak{} False


\clearpage

\section*{P0171  Polynomial small bias sufficient to fool DNF events}

\entrylabel{Primary source}

O Donnell:\allowbreak{} Open problems in analysis of Boolean functions;\allowbreak{} PDF pp.9-10,\allowbreak{} epsilon-biased Sets for DNFs


\url{https://arxiv.org/pdf/1204.6447}


Source statement date:\allowbreak{} 2012-04-29. Archived access:\allowbreak{} 2026-09-06T18:\allowbreak{}55:\allowbreak{}14Z.


Historical:\allowbreak{} source\_\allowbreak{}says\_\allowbreak{}open | Current:\allowbreak{} unverified | Result:\allowbreak{} attempted\_\allowbreak{}unsolved


\entrylabel{Question}

For each fixed error delta>0,\allowbreak{} determine whether every distribution whose nonconstant parity biases are at most s\textasciicircum{}(-C\_\allowbreak{}delta) fools every size-s DNF to error delta.


\entrylabel{Scope}

The2025 Generalised Linial-Nisan counterexample controls marginals of bounded-size coordinate sets;\allowbreak{} its theorem alone does not assert small bias for every parity. It is not treated as a resolution of this different source statement.


\subsection*{P0171-A001 | Polynomial small bias sufficient to fool DNF events exact calculation or structural obstruction}

\entrylabel{Argument or result}

For a single conjunction C of k literals,\allowbreak{} its Fourier coefficients have magnitude2\textasciicircum{}-k on subsets of those k coordinates. An epsilon-biased distribution therefore changes E C by at most epsilon*(1-2\textasciicircum{}-k). Inclusion-exclusion for a union of s clauses gives the valid but crude error at most(2\textasciicircum{}s-1)*epsilon,\allowbreak{} since each nonempty intersection of consistent clauses is another conjunction.


\entrylabel{Second check}

The constant coefficient contributes no expectation difference,\allowbreak{} while all other coefficients are paired with parity biases of magnitude at most epsilon. Expanding the union into its2\textasciicircum{}s-1 nonempty clause intersections proves the stated global bound without assuming independence of clauses.


\entrylabel{Remaining obstruction}

To obtain polynomial rather than exponential bias requirements,\allowbreak{} one needs a substantially more economical approximation or cancellation argument. Small errors for individual clause probabilities do not by themselves control their union.


Verification:\allowbreak{} author checked limited argument | Independent review:\allowbreak{} none | Full target solved:\allowbreak{} False


\clearpage

\section*{P0172  Sharp erasure-distillation bound from the Gaussian limit}

\entrylabel{Primary source}

Kamath,\allowbreak{}Kumar,\allowbreak{}Manurangsi:\allowbreak{} When Majority Fails,\allowbreak{}2026;\allowbreak{} Section1.4,\allowbreak{} Conjecture2


\url{https://arxiv.org/pdf/2604.06590}


Source statement date:\allowbreak{} 2026-04-08. Archived access:\allowbreak{} 2026-09-06T20:\allowbreak{}45:\allowbreak{}41Z.


Historical:\allowbreak{} source\_\allowbreak{}says\_\allowbreak{}open | Current:\allowbreak{} unverified | Result:\allowbreak{} attempted\_\allowbreak{}unsolved


\entrylabel{Question}

For unbiased Boolean f and coordinate-retention probability p in(0,\allowbreak{}1/\allowbreak{}2),\allowbreak{} prove E|E[f(X) given the retained coordinates]|<=\allowbreak{}(2/\allowbreak{}pi)*arcsin(sqrt(p)).


\entrylabel{Scope}

The false finite-majority comparator is excluded. This is the2026 authors refined conjecture. Requiring the optimizing second output to be unbiased is a distinct constraint;\allowbreak{} the source formulation here uses the absolute conditional expectation.


\subsection*{P0172-A001 | Sharp erasure-distillation bound from the Gaussian limit exact calculation or structural obstruction}

\entrylabel{Argument or result}

Identify the conditional expectation with the multilinear extension f(Z),\allowbreak{} where each Z\_\allowbreak{}i is0 with probability1-p and otherwise a fair sign. Orthogonality gives E f(Z)\textasciicircum{}2=\allowbreak{}sum\_\allowbreak{}S p\textasciicircum{}|S|*fhat(S)\textasciicircum{}2<=\allowbreak{}p for unbiased f. Thus Cauchy-Schwarz yields E|f(Z)|<=\allowbreak{}sqrt(p). A parity on m coordinates instead has exact value p\textasciicircum{}m.


\entrylabel{Second check}

If S and T differ,\allowbreak{} the expectation of Z\_\allowbreak{}S Z\_\allowbreak{}T vanishes because a differing coordinate has mean0. For S=\allowbreak{}T it equals p\textasciicircum{}|S|. Parseval and the absent constant coefficient give the bound. A parity survives only when all m relevant coordinates are retained.


\entrylabel{Remaining obstruction}

The L2 estimate loses the proposed arcsine constant. Conditional distributions of the smoothed Boolean function,\allowbreak{} beyond variance alone,\allowbreak{} must be controlled.


Verification:\allowbreak{} author checked limited argument | Independent review:\allowbreak{} none | Full target solved:\allowbreak{} False


\clearpage

\section*{P0173  Asymptotic influence bounds for polynomial thresholds}

\entrylabel{Primary source}

O Donnell:\allowbreak{} Open problems in analysis of Boolean functions;\allowbreak{} PDF p.3,\allowbreak{} weaker Gotsman-Linial formulations


\url{https://arxiv.org/pdf/1204.6447}


Source statement date:\allowbreak{} 2012-04-29. Archived access:\allowbreak{} 2026-09-06T18:\allowbreak{}55:\allowbreak{}14Z.


Historical:\allowbreak{} source\_\allowbreak{}says\_\allowbreak{}open | Current:\allowbreak{} unverified | Result:\allowbreak{} attempted\_\allowbreak{}unsolved


\entrylabel{Question}

Prove total influence at most C*d*sqrt(n),\allowbreak{} or at least C\_\allowbreak{}d*sqrt(n),\allowbreak{} for every degree-d polynomial threshold function on n uniform bits.


\entrylabel{Scope}

The exact symmetric-maximizer conjecture is false,\allowbreak{} with Chapman2108.02288 giving later counterexamples. The weaker asymptotic bounds explicitly posed in the source are grouped here.


\subsection*{P0173-A001 | Asymptotic influence bounds for polynomial thresholds exact calculation or structural obstruction}

\entrylabel{Argument or result}

The desired C*d*sqrt(n) bound holds when the polynomial is a product of d affine linear forms with no cube zeros. Indeed its sign is a product of d threshold functions. If that product changes on a coordinate flip,\allowbreak{} at least one factor changes,\allowbreak{} so its influence is at most the sum of the factor influences. For a threshold g,\allowbreak{} unateness gives Inf\_\allowbreak{}i(g)=\allowbreak{}|ghat(i)| and hence I(g)<=\allowbreak{}sqrt(n)*sqrt(sum\_\allowbreak{}i ghat(i)\textasciicircum{}2)<=\allowbreak{}sqrt(n).


\entrylabel{Second check}

For a monotone coordinate of a Boolean function,\allowbreak{} its discrete derivative is either0 or1 after the conventional half-difference,\allowbreak{} so its expectation equals its flip probability. Cauchy-Schwarz and Parseval then give the threshold bound. A union bound gives the product estimate.


\entrylabel{Remaining obstruction}

A general multivariate polynomial need not factor into real affine forms. The argument supplies a genuine subclass result,\allowbreak{} not a reduction for arbitrary polynomial thresholds.


Verification:\allowbreak{} author checked limited argument | Independent review:\allowbreak{} none | Full target solved:\allowbreak{} False


\clearpage

\section*{P0174  Sharp support-diameter threshold for free Brownian unimodality}

\entrylabel{Primary source}

Hasebe,\allowbreak{}Ueda:\allowbreak{} Large time unimodality for classical and free Brownian motions;\allowbreak{} Problem 3.3,\allowbreak{} arXiv PDF p.6;\allowbreak{} journal p.359


\url{https://arxiv.org/pdf/1710.08240}


Source statement date:\allowbreak{} 2018-09-14. Archived access:\allowbreak{} 2026-09-06T20:\allowbreak{}48:\allowbreak{}42Z.


Historical:\allowbreak{} source\_\allowbreak{}says\_\allowbreak{}open | Current:\allowbreak{} unverified | Result:\allowbreak{} candidate\_\allowbreak{}full\_\allowbreak{}solution


\entrylabel{Question}

Find the optimal universal C such that every compactly supported law mu becomes unimodal after free convolution with the variance-t semicircle law for all t >=\allowbreak{} C times the square of its support diameter.


\entrylabel{Scope}

The source prints 2<=\allowbreak{}C<=\allowbreak{}4,\allowbreak{} but its own symmetric-Bernoulli threshold t=\allowbreak{}4 and diameter 2 imply only C>=\allowbreak{}1. The candidate proof derives both bounds independently. Manuscript header 2021 differs from the arXiv revision date 2018.


\subsection*{P0174-A001 | Centered fourth-moment diameter inequality and reweighted Cauchy-profile critical points}

\entrylabel{Argument or result}

Candidate sharp answer C=\allowbreak{}1. For any centered law Q of support diameter D,\allowbreak{} with a=\allowbreak{}E X\textasciicircum{}2 and b=\allowbreak{}E X\textasciicircum{}4,\allowbreak{} a<=\allowbreak{}D\textasciicircum{}2/\allowbreak{}4 and b+\allowbreak{}3a\textasciicircum{}2<=\allowbreak{}D\textasciicircum{}2*a. Reweighting mu by ((x-u)\textasciicircum{}2+\allowbreak{}R\textasciicircum{}2)\textasciicircum{}(-2) at a stationary Cauchy profile gives y=\allowbreak{}R\textasciicircum{}2<=\allowbreak{}3a when the second derivative is nonnegative,\allowbreak{} and 1/\allowbreak{}F=\allowbreak{}(b+\allowbreak{}2ay+\allowbreak{}y\textasciicircum{}2)/\allowbreak{}(a+\allowbreak{}y)<=\allowbreak{}D\textasciicircum{}2. Equality forces symmetric endpoint Bernoulli and a strict fourth-order maximum. Thus no local minimum has F<=\allowbreak{}1/\allowbreak{}D\textasciicircum{}2;\allowbreak{} each level 1/\allowbreak{}t for t>=\allowbreak{}D\textasciicircum{}2 has at most two roots. Biane parametrization and the source level-set lemma imply unimodality. Symmetric Bernoulli gives sharpness.


\entrylabel{Second check}

Second analytic route differentiates v\_\allowbreak{}t directly at critical points. For symmetric Bernoulli on \{-c,\allowbreak{}c\},\allowbreak{} v\_\allowbreak{}t second derivative at 0 equals (4*c\textasciicircum{}2-t)/\allowbreak{}(t*sqrt(t-c\textasciicircum{}2)),\allowbreak{} so t<4*c\textasciicircum{}2=\allowbreak{}D\textasciicircum{}2 gives a positive density local minimum. The centered moment inequality follows independently from E[(X-Xprime)\textasciicircum{}4]<=\allowbreak{}D\textasciicircum{}2*E[(X-Xprime)\textasciicircum{}2]. These are same-author checks,\allowbreak{} not independent review.


\entrylabel{Remaining obstruction}

No mathematical gap currently identified in the full written candidate. Later-literature novelty is unverified and there is no independent proof review. Check the reweighting normalization,\allowbreak{} Jensen direction,\allowbreak{} equality case,\allowbreak{} and exact source density assumptions before promotion.


Verification:\allowbreak{} author checked complete candidate | Independent review:\allowbreak{} none | Full target solved:\allowbreak{} False


\entrylabel{Full proof packet}

proofs/\allowbreak{}P0174\_\allowbreak{}free\_\allowbreak{}brownian\_\allowbreak{}unimodality\_\allowbreak{}constant\_\allowbreak{}candidate.md


\clearpage

\section*{P0175  Symmetric unimodality under free addition}

\entrylabel{Primary source}

Hasebe-Ueda:\allowbreak{} Large time unimodality,\allowbreak{} published version;\allowbreak{} Conjecture 3.5


\url{https://alea.math.cnrs.fr/articles/v15/15-15.pdf}


Source statement date:\allowbreak{} 2018-03-23. Archived access:\allowbreak{} 2026-09-06T21:\allowbreak{}27:\allowbreak{}11Z.


Historical:\allowbreak{} source\_\allowbreak{}says\_\allowbreak{}open | Current:\allowbreak{} unverified | Result:\allowbreak{} attempted\_\allowbreak{}unsolved


\entrylabel{Question}

Is the free additive convolution of two symmetric unimodal probability laws unimodal?


\entrylabel{Scope}

The older AIM free-analysis question is the same family and is not counted again. A semicircle factor and some freely infinitely divisible subclasses are established cases.


\subsection*{P0175-A001 | Symmetric unimodality under free addition structural calculation}

\entrylabel{Argument or result}

The classical proof by mixing symmetric unimodal components does not transfer by linearity:\allowbreak{} free convolution is not affine even on symmetric unimodal laws. Let mu=\allowbreak{}(delta\_\allowbreak{}0+\allowbreak{}S(0,\allowbreak{}2))/\allowbreak{}2 and nu=\allowbreak{}S(0,\allowbreak{}1). The sixth moment of mu free-added to nu is67,\allowbreak{} whereas the mixture of the two component free convolutions has sixth moment70.


\entrylabel{Second check}

For mu,\allowbreak{} m2=\allowbreak{}1,\allowbreak{}m4=\allowbreak{}4,\allowbreak{}m6=\allowbreak{}20,\allowbreak{} so its free cumulants are k2=\allowbreak{}1,\allowbreak{}k4=\allowbreak{}2,\allowbreak{}k6=\allowbreak{}3. Adding nu changes k2 to2 only,\allowbreak{} giving m6=\allowbreak{}3+\allowbreak{}6*2*2+\allowbreak{}5*2\textasciicircum{}3=\allowbreak{}67. The affine mixture would be(S(0,\allowbreak{}1)+\allowbreak{}S(0,\allowbreak{}3))/\allowbreak{}2,\allowbreak{} with m6=\allowbreak{}(5+\allowbreak{}135)/\allowbreak{}2=\allowbreak{}70. Both mixture components and mu are symmetric unimodal.


\entrylabel{Remaining obstruction}

This disproves a proposed proof mechanism,\allowbreak{} not the unimodality conjecture. The actual free convolution in this example is covered by the known semicircle-factor theorem.


Verification:\allowbreak{} author checked limited argument | Independent review:\allowbreak{} none | Full target solved:\allowbreak{} False


\clearpage

\section*{P0176  Existence of nontrivial freely strongly unimodal laws}

\entrylabel{Primary source}

Hasebe-Ueda:\allowbreak{} Large time unimodality,\allowbreak{} published version;\allowbreak{} Problem 3.11


\url{https://alea.math.cnrs.fr/articles/v15/15-15.pdf}


Source statement date:\allowbreak{} 2018-03-23. Archived access:\allowbreak{} 2026-09-06T21:\allowbreak{}27:\allowbreak{}11Z.


Historical:\allowbreak{} source\_\allowbreak{}says\_\allowbreak{}open | Current:\allowbreak{} unverified | Result:\allowbreak{} attempted\_\allowbreak{}unsolved


\entrylabel{Question}

Does a non-point-mass law exist whose free additive convolution with every unimodal law is unimodal?


\entrylabel{Scope}

The source already excludes all nondegenerate finite-variance laws. The remaining possible examples must have infinite variance.


\subsection*{P0176-A001 | Existence of nontrivial freely strongly unimodal laws structural calculation}

\entrylabel{Argument or result}

An explicit infinite-variance candidate also fails:\allowbreak{} the standard Cauchy law. Let nu=\allowbreak{}(1/\allowbreak{}20)*delta\_\allowbreak{}0+\allowbreak{}(19/\allowbreak{}20)\emph{Uniform[0,\allowbreak{}1000],\allowbreak{} a unimodal law with mode0. Cauchy free convolution equals classical convolution,\allowbreak{} whose density derivative,\allowbreak{} multiplied by pi,\allowbreak{} is -x/\allowbreak{}[10}(1+\allowbreak{}x\textasciicircum{}2)\textasciicircum{}2]+\allowbreak{}19/\allowbreak{}20000*[1/\allowbreak{}(1+\allowbreak{}x\textasciicircum{}2)-1/\allowbreak{}(1+\allowbreak{}(x-1000)\textasciicircum{}2)]. Its signs at0,\allowbreak{}1,\allowbreak{}250,\allowbreak{}1000 are respectively positive,\allowbreak{}negative,\allowbreak{}positive,\allowbreak{}negative,\allowbreak{} so the density has at least two peaks.


\entrylabel{Second check}

At1 the negative term is-1/\allowbreak{}40 and the positive term is less than19/\allowbreak{}40000. At250,\allowbreak{} writing A=\allowbreak{}62501 and B=\allowbreak{}562501,\allowbreak{} the derivative is -25/\allowbreak{}A\textasciicircum{}2+\allowbreak{}475/\allowbreak{}(A*B),\allowbreak{} positive because19*A>B. At1000 both displayed contributions are negative. The initial law has a single atom at its mode and a constant density to its right.


\entrylabel{Remaining obstruction}

This gives a concrete Cauchy counterexample,\allowbreak{} consistent with the known strong-unimodality criterion,\allowbreak{} but does not exclude every infinite-variance law.


Verification:\allowbreak{} author checked limited argument | Independent review:\allowbreak{} none | Full target solved:\allowbreak{} False


\clearpage

\section*{P0177  Location of Brownian modes at large time}

\entrylabel{Primary source}

Hasebe-Ueda:\allowbreak{} Large time unimodality,\allowbreak{} published version;\allowbreak{} Problem 4.6


\url{https://alea.math.cnrs.fr/articles/v15/15-15.pdf}


Source statement date:\allowbreak{} 2018-03-23. Archived access:\allowbreak{} 2026-09-06T21:\allowbreak{}27:\allowbreak{}11Z.


Historical:\allowbreak{} source\_\allowbreak{}says\_\allowbreak{}open | Current:\allowbreak{} unverified | Result:\allowbreak{} attempted\_\allowbreak{}unsolved


\entrylabel{Question}

Estimate the mode of classical Brownian motion started from a law with a finite exponential quadratic moment,\allowbreak{} and determine analogous estimates for free Brownian motion.


\entrylabel{Scope}

The source gives a growing interval of radius sqrt(t)/\allowbreak{}2 once unimodality holds. Classical and free versions of this mode-location request are grouped.


\subsection*{P0177-A001 | Location of Brownian modes at large time structural calculation}

\entrylabel{Argument or result}

For the classical mode,\allowbreak{} let m be the initial mean and mu\_\allowbreak{}k=\allowbreak{}E[(X-m)\textasciicircum{}k]. The source exponential-moment assumption yields the asymptotic expansion mode(t)=\allowbreak{}m-mu\_\allowbreak{}3/\allowbreak{}(2*t)+\allowbreak{}(mu\_\allowbreak{}5-6*mu\_\allowbreak{}2*mu\_\allowbreak{}3)/\allowbreak{}(8*t\textasciicircum{}2)+\allowbreak{}O(t\textasciicircum{}-3). In particular the mode tends to the initial mean,\allowbreak{} and the leading displacement has sign opposite the third central moment.


\entrylabel{Second check}

Set s=\allowbreak{}1/\allowbreak{}t and H(h,\allowbreak{}s)=\allowbreak{}E[(X-m-h)*exp(-s*(X-m-h)\textasciicircum{}2/\allowbreak{}2)]. The exponential quadratic moment makes H analytic near(0,\allowbreak{}0),\allowbreak{} with H\_\allowbreak{}h=\allowbreak{}-1 there. The implicit function theorem gives h=\allowbreak{}a*s+\allowbreak{}b*s\textasciicircum{}2+\allowbreak{}O(s\textasciicircum{}3). Expanding H gives a=\allowbreak{}-mu\_\allowbreak{}3/\allowbreak{}2 and b=\allowbreak{}mu\_\allowbreak{}5/\allowbreak{}8-3*mu\_\allowbreak{}2*mu\_\allowbreak{}3/\allowbreak{}4. The associated density stationary point is a strict local maximum since H\_\allowbreak{}h remains negative;\allowbreak{} the source large-time unimodality identifies it with the unique global mode.


\entrylabel{Remaining obstruction}

This is an author-checked asymptotic partial answer,\allowbreak{} with no novelty claim. It does not provide a uniform explicit finite-time error constant or solve the free-Brownian half.


Verification:\allowbreak{} author checked limited argument | Independent review:\allowbreak{} none | Full target solved:\allowbreak{} False


\clearpage

\section*{P0178  Critical moment condition for eventual stable-process unimodality}

\entrylabel{Primary source}

Hasebe-Ueda:\allowbreak{} Large time unimodality,\allowbreak{} published version;\allowbreak{} Conjecture 5.5,\allowbreak{} journal p.373


\url{https://alea.math.cnrs.fr/articles/v15/15-15.pdf}


Source statement date:\allowbreak{} 2018-03-23. Archived access:\allowbreak{} 2026-09-06T21:\allowbreak{}27:\allowbreak{}11Z.


Historical:\allowbreak{} source\_\allowbreak{}says\_\allowbreak{}open | Current:\allowbreak{} unverified | Result:\allowbreak{} attempted\_\allowbreak{}unsolved


\entrylabel{Question}

For a classical or free strictly alpha-stable process,\allowbreak{} alpha in(0,\allowbreak{}2),\allowbreak{} does a finite absolute(2+\allowbreak{}alpha)-moment of the initial law ensure eventual unimodality under the corresponding convolution?


\entrylabel{Scope}

Use the published wording:\allowbreak{} classical convolution with classical stable laws,\allowbreak{} free convolution with free stable laws. The archived preprint wording differs. All stability indices are one source conjecture;\allowbreak{} the Cauchy case is already established.


\subsection*{P0178-A001 | Critical moment condition for eventual stable-process unimodality structural calculation}

\entrylabel{Argument or result}

A separated-atom counterexample of the source type cannot retain the critical moment while making its usual tail mechanism arbitrarily strong. If mu=\allowbreak{}sum\_\allowbreak{}k w\_\allowbreak{}k*delta\_\allowbreak{}(a\_\allowbreak{}k) and sum\_\allowbreak{}k w\_\allowbreak{}k*|a\_\allowbreak{}k|\textasciicircum{}(2+\allowbreak{}alpha)<infinity,\allowbreak{} then w\_\allowbreak{}k*|a\_\allowbreak{}k|\textasciicircum{}(2+\allowbreak{}alpha) tends to0. Thus it cannot diverge along a subsequence,\allowbreak{} which is the condition used to make an isolated atom dominate the far-tail derivative of the main stable density.


\entrylabel{Second check}

Every term of a convergent nonnegative series tends to0. In the Cauchy case alpha=\allowbreak{}1,\allowbreak{} the derivative of the bulk kernel at a\_\allowbreak{}k has order |a\_\allowbreak{}k|\textasciicircum{}-3,\allowbreak{} while the local atom contributes order w\_\allowbreak{}k at any fixed smoothing time;\allowbreak{} their ratio is of order w\_\allowbreak{}k*|a\_\allowbreak{}k|\textasciicircum{}3. This matches the critical exponent in the exact known construction.


\entrylabel{Remaining obstruction}

Failure of that counterexample mechanism is not a uniform derivative estimate for arbitrary initial laws. Asymmetric stable laws and the free cases still require global control.


Verification:\allowbreak{} author checked limited argument | Independent review:\allowbreak{} none | Full target solved:\allowbreak{} False


\clearpage

\section*{P0179  Uniform permutation representations of the Gieseking group}

\entrylabel{Primary source}

The strong convergence phenomenon,\allowbreak{}2025 survey;\allowbreak{} Section 6.1,\allowbreak{} p.64


\url{https://arxiv.org/pdf/2507.00346}


Source statement date:\allowbreak{} 2026-02-11. Archived access:\allowbreak{} 2026-09-06T20:\allowbreak{}51:\allowbreak{}12Z.


Historical:\allowbreak{} source\_\allowbreak{}says\_\allowbreak{}open | Current:\allowbreak{} unverified | Result:\allowbreak{} attempted\_\allowbreak{}unsolved


\entrylabel{Question}

Do uniformly sampled permutation pairs satisfying U\textasciicircum{}2*V\textasciicircum{}2=\allowbreak{}V*U converge strongly,\allowbreak{} after removing the common constant vector,\allowbreak{} to the regular representation of the Gieseking group?


\entrylabel{Scope}

Surface groups are already covered by the source. The later2602.11905 free-product theorem concerns a different relation and does not settle this hyperbolic3-manifold group.


\subsection*{P0179-A001 | Uniform permutation representations of the Gieseking group structural calculation}

\entrylabel{Argument or result}

Independent-permutation moment formulas cannot be inserted unchanged. Taking determinants in the relation forces sign(U)=\allowbreak{}sign(V). More fundamentally,\allowbreak{} the relator U\textasciicircum{}2*V\textasciicircum{}2*U\textasciicircum{}-1*V\textasciicircum{}-1 is identically I,\allowbreak{} so its normalized trace is1,\allowbreak{} whereas the corresponding nontrivial reduced free-group word has trace0 in the free regular representation.


\entrylabel{Second check}

The determinant equation is1=\allowbreak{}sign(V)*sign(U). The relator identity follows by multiplying U\textasciicircum{}2*V\textasciicircum{}2=\allowbreak{}VU on the right by U\textasciicircum{}-1*V\textasciicircum{}-1. Removing the constant vector leaves the identity on the remaining space,\allowbreak{} so the trace discrepancy persists.


\entrylabel{Remaining obstruction}

The correct nonfree group trace respects the relator,\allowbreak{} but no asymptotic enumeration or outlier estimate for uniform constrained pairs is obtained.


Verification:\allowbreak{} author checked limited argument | Independent review:\allowbreak{} none | Full target solved:\allowbreak{} False


\clearpage

\section*{P0180  Optimal spectra of random symmetric-group Cayley graphs}

\entrylabel{Primary source}

The strong convergence phenomenon,\allowbreak{}2025 survey;\allowbreak{} Section 6.2,\allowbreak{} p.65


\url{https://arxiv.org/pdf/2507.00346}


Source statement date:\allowbreak{} 2026-02-11. Archived access:\allowbreak{} 2026-09-06T20:\allowbreak{}51:\allowbreak{}12Z.


Historical:\allowbreak{} source\_\allowbreak{}says\_\allowbreak{}open | Current:\allowbreak{} unverified | Result:\allowbreak{} attempted\_\allowbreak{}unsolved


\entrylabel{Question}

For a fixed number of uniform random generators of S\_\allowbreak{}N,\allowbreak{} do all nontrivial eigenvalues apart from the sign representation approach the free-group spectral bound? Develop the associated high-dimensional representation convergence.


\entrylabel{Scope}

The weaker expansion question and the stronger arbitrary-irrep formulation are grouped with the optimal-gap family. The trivial and sign representations must both be removed.


\subsection*{P0180-A001 | Optimal spectra of random symmetric-group Cayley graphs structural calculation}

\entrylabel{Argument or result}

The sign obstruction has an exact positive probability:\allowbreak{} when all r generator signs agree,\allowbreak{} the corresponding adjacency eigenvalue has absolute value2r,\allowbreak{} above2*sqrt(2r-1) for r>=\allowbreak{}2. This event has probability2\textasciicircum{}(1-r). After removing it,\allowbreak{} pointwise convergence for each selected irrep still cannot justify the supremum over all irreps:\allowbreak{} in a triangular array with one uniformly chosen bad coordinate per row,\allowbreak{} every coordinate is bad with probability1/\allowbreak{}m\_\allowbreak{}N tending to0,\allowbreak{} yet the row maximum is bad surely.


\entrylabel{Second check}

The sign variables are independent fair signs for uniform permutations. The array example verifies the quantifier failure exactly and indicates the need for uniform probability bounds,\allowbreak{} rather than just increasing the range of an individual representation theorem.


\entrylabel{Remaining obstruction}

No uniform tail estimate over the high-dimensional irreducible representations is proved.


Verification:\allowbreak{} author checked limited argument | Independent review:\allowbreak{} none | Full target solved:\allowbreak{} False


\clearpage

\section*{P0181  Fixed-degree random Cayley graphs with optimal spectral gap}

\entrylabel{Primary source}

The strong convergence phenomenon,\allowbreak{}2025 survey;\allowbreak{} Section 6.2,\allowbreak{} pp.65-66


\url{https://arxiv.org/pdf/2507.00346}


Source statement date:\allowbreak{} 2026-02-11. Archived access:\allowbreak{} 2026-09-06T20:\allowbreak{}51:\allowbreak{}12Z.


Historical:\allowbreak{} source\_\allowbreak{}says\_\allowbreak{}open | Current:\allowbreak{} unverified | Result:\allowbreak{} attempted\_\allowbreak{}unsolved


\entrylabel{Question}

Find finite-group sequences whose fixed-number uniform random Cayley generators have asymptotically optimal nontrivial spectral gap,\allowbreak{} particularly finite special linear groups.


\entrylabel{Scope}

Different field and rank regimes are grouped;\allowbreak{} growing-number-of-generator theorems do not answer the fixed-degree question.


\subsection*{P0181-A001 | Fixed-degree random Cayley graphs with optimal spectral gap structural calculation}

\entrylabel{Argument or result}

Abelian groups cannot supply even a uniform fixed-degree gap in a basic test case. For arbitrary r generators a\_\allowbreak{}i of Z/\allowbreak{}pZ,\allowbreak{} divide the r-dimensional phase cube into M\textasciicircum{}r boxes with M\textasciicircum{}r<p. Two of the p points k*(a\_\allowbreak{}1,\allowbreak{}...,\allowbreak{}a\_\allowbreak{}r)/\allowbreak{}p fall in one box. Their nonzero difference k gives a nontrivial character for which all generator phases have distance at most2*pi/\allowbreak{}M from zero. Hence its adjacency eigenvalue is at least2r-4*pi\textasciicircum{}2*r/\allowbreak{}M\textasciicircum{}2.


\entrylabel{Second check}

Use cos(x)>=\allowbreak{}1-x\textasciicircum{}2/\allowbreak{}2 in the eigenvalue2*sum\_\allowbreak{}i cos(2*pi*k*a\_\allowbreak{}i/\allowbreak{}p). Taking M tending to infinity at order p\textasciicircum{}(1/\allowbreak{}r) forces the gap to0,\allowbreak{} regardless of how the generators were chosen.


\entrylabel{Remaining obstruction}

The argument identifies a structural obstruction to abelian candidates but supplies no optimal bound for the nonabelian special-linear families.


Verification:\allowbreak{} author checked limited argument | Independent review:\allowbreak{} none | Full target solved:\allowbreak{} False


\clearpage

\section*{P0182  Strong convergence of high-dimensional representations of a fixed compact group}

\entrylabel{Primary source}

The strong convergence phenomenon,\allowbreak{}2025 survey;\allowbreak{} Section 6.3,\allowbreak{} p.66


\url{https://arxiv.org/pdf/2507.00346}


Source statement date:\allowbreak{} 2026-02-11. Archived access:\allowbreak{} 2026-09-06T20:\allowbreak{}51:\allowbreak{}12Z.


Historical:\allowbreak{} source\_\allowbreak{}says\_\allowbreak{}open | Current:\allowbreak{} unverified | Result:\allowbreak{} attempted\_\allowbreak{}unsolved


\entrylabel{Question}

For independent Haar elements of a fixed compact semisimple group,\allowbreak{} do their growing irreducible representations converge strongly to free Haar generators,\allowbreak{} starting with symmetric powers of SU(2)?


\entrylabel{Scope}

The source already establishes weak convergence in broad generality. The SU(2) test case and broader fixed-group question are one family.


\subsection*{P0182-A001 | Strong convergence of high-dimensional representations of a fixed compact group structural calculation}

\entrylabel{Argument or result}

The one-generator SU(2) case has an elementary complete answer. A Haar element has eigenvalues exp(plus/\allowbreak{}minus i*theta),\allowbreak{} with theta/\allowbreak{}pi irrational almost surely. Its N-th symmetric power has eigenvalues exp(i*(N-2k)*theta),\allowbreak{}0<=\allowbreak{}k<=\allowbreak{}N. These become dense on the unit circle,\allowbreak{} so every Laurent-polynomial norm converges to the supremum on the circle,\allowbreak{} the regular Z norm.


\entrylabel{Second check}

Irrational rotation makes the first N+\allowbreak{}1 multiples of2*theta asymptotically dense;\allowbreak{} multiplying the set by exp(i*N*theta) preserves its maximal circular gap. Uniform continuity then proves norm convergence for each Laurent polynomial. Geometric sums give convergence of every nonzero normalized trace moment to0.


\entrylabel{Remaining obstruction}

For two or more independently sampled group elements,\allowbreak{} their representation eigenbases interact;\allowbreak{} the scalar circle-density argument no longer controls joint polynomial norms.


Verification:\allowbreak{} author checked limited argument | Independent review:\allowbreak{} none | Full target solved:\allowbreak{} False


\clearpage

\section*{P0183  Deterministic strong convergence from LPS generators}

\entrylabel{Primary source}

The strong convergence phenomenon,\allowbreak{}2025 survey;\allowbreak{} Section 6.4,\allowbreak{} pp.66-67


\url{https://arxiv.org/pdf/2507.00346}


Source statement date:\allowbreak{} 2026-02-11. Archived access:\allowbreak{} 2026-09-06T20:\allowbreak{}51:\allowbreak{}12Z.


Historical:\allowbreak{} source\_\allowbreak{}says\_\allowbreak{}open | Current:\allowbreak{} unverified | Result:\allowbreak{} attempted\_\allowbreak{}unsolved


\entrylabel{Question}

Do the explicit LPS-type finite-group generators yield strong convergence to free generators,\allowbreak{} beyond their established adjacency spectral gap?


\entrylabel{Scope}

This is the explicit deterministic model in the source,\allowbreak{} not a request to rediscover deterministic Ramanujan graphs.


\subsection*{P0183-A001 | Deterministic strong convergence from LPS generators structural calculation}

\entrylabel{Argument or result}

Weak freeness plus the usual adjacency norm is insufficient in general. Start with any strongly convergent r-unitary model and append a one-dimensional block in which every generator equals z=\allowbreak{}exp(i*theta),\allowbreak{} with cos(theta)=\allowbreak{}sqrt(2r-1)/\allowbreak{}r. This block does not change limiting traces and its adjacency eigenvalue is exactly2*sqrt(2r-1). Yet the phase-twisted adjacency polynomial has value2r on that block,\allowbreak{} above its free norm2*sqrt(2r-1).


\entrylabel{Second check}

Normalized trace weight of the added block tends to0. Multiplying each free Haar generator by a fixed phase preserves the free law and norm. Thus both asserted facts hold while a different linear polynomial has a persistent outlier.


\entrylabel{Remaining obstruction}

This is an obstruction to inferring strong convergence from one spectral-gap theorem;\allowbreak{} the constructed scalar blocks are not the actual LPS representation,\allowbreak{} so the source conjecture is not refuted.


Verification:\allowbreak{} author checked limited argument | Independent review:\allowbreak{} none | Full target solved:\allowbreak{} False


\clearpage

\section*{P0184  Strong convergence of modular-group actions on finite projective lines}

\entrylabel{Primary source}

The strong convergence phenomenon,\allowbreak{}2025 survey;\allowbreak{} Section 6.4,\allowbreak{} p.67


\url{https://arxiv.org/pdf/2507.00346}


Source statement date:\allowbreak{} 2026-02-11. Archived access:\allowbreak{} 2026-09-06T20:\allowbreak{}51:\allowbreak{}12Z.


Historical:\allowbreak{} source\_\allowbreak{}says\_\allowbreak{}open | Current:\allowbreak{} unverified | Result:\allowbreak{} attempted\_\allowbreak{}unsolved


\entrylabel{Question}

Does the PSL\_\allowbreak{}2(Z) action on P\textasciicircum{}1(F\_\allowbreak{}q),\allowbreak{} restricted to constants-orthogonal vectors,\allowbreak{} converge strongly to the regular group representation for every fixed tuple of integer matrices?


\entrylabel{Scope}

The diffusion-operator formulation and arbitrary-polynomial formulation are equivalent in the source and are counted once. Random homomorphisms from C2*C3 are a different model.


\subsection*{P0184-A001 | Strong convergence of modular-group actions on finite projective lines structural calculation}

\entrylabel{Argument or result}

Weak convergence follows directly from fixed points. For a fixed nonidentity matrix X in PSL\_\allowbreak{}2(Z),\allowbreak{} its fixed points on P\textasciicircum{}1(F\_\allowbreak{}q) solve a projective quadratic equation. Except for finitely many primes where X becomes scalar,\allowbreak{} there are at most two fixed points. The normalized trace in the q-dimensional constants-orthogonal representation is therefore(fix(X)-1)/\allowbreak{}q,\allowbreak{} which tends to0.


\entrylabel{Second check}

For X represented by entries a,\allowbreak{}b,\allowbreak{}c,\allowbreak{}d,\allowbreak{} finite fixed points satisfy c*z\textasciicircum{}2+\allowbreak{}(d-a)*z-b=\allowbreak{}0,\allowbreak{} with the point at infinity handled projectively. A nonzero quadratic on the projective line has at most two zeros. Applying this to each word yields the regular trace moments. Faithfulness of the regular trace then gives the lower polynomial-norm bound by even moments.


\entrylabel{Remaining obstruction}

Moment convergence gives no matching upper bound on operator norms;\allowbreak{} sparse exceptional eigenvalues remain the exact obstruction.


Verification:\allowbreak{} author checked limited argument | Independent review:\allowbreak{} none | Full target solved:\allowbreak{} False


\clearpage

\section*{P0185  Ramanujan inequalities for growing families of matrix polynomials}

\entrylabel{Primary source}

The strong convergence phenomenon,\allowbreak{}2025 survey;\allowbreak{} Section 6.5,\allowbreak{} pp.67-68


\url{https://arxiv.org/pdf/2507.00346}


Source statement date:\allowbreak{} 2026-02-11. Archived access:\allowbreak{} 2026-09-06T20:\allowbreak{}51:\allowbreak{}12Z.


Historical:\allowbreak{} source\_\allowbreak{}says\_\allowbreak{}open | Current:\allowbreak{} unverified | Result:\allowbreak{} attempted\_\allowbreak{}unsolved


\entrylabel{Question}

Can permutation models satisfy the free polynomial norm upper bound with no error for every polynomial of degree at most q(N) tending to infinity,\allowbreak{} or for each fixed polynomial along a chosen model?


\entrylabel{Scope}

The per-polynomial weaker question and simultaneous-degree version are grouped. The source excludes a bound simultaneously for every degree at fixed dimension.


\subsection*{P0185-A001 | Ramanujan inequalities for growing families of matrix polynomials structural calculation}

\entrylabel{Argument or result}

The one-generator amenable case satisfies the exact inequality at every degree:\allowbreak{} take an N-cycle permutation and remove its constant eigenvector. Its remaining eigenvalues lie on the unit circle,\allowbreak{} so the norm of any Laurent polynomial is at most its supremum on that circle,\allowbreak{} exactly the norm in the regular representation of Z.


\entrylabel{Second check}

A normal finite matrix is unitarily diagonalizable,\allowbreak{} and Laurent-polynomial norms are maxima over its eigenvalues. This proves the bound without an asymptotic error or probability qualifier.


\entrylabel{Remaining obstruction}

For r>=\allowbreak{}2 the free reduced group algebra is noncommutative and has no finite-dimensional faithful representation. The circle argument gives no construction for any diverging degree window in that setting.


Verification:\allowbreak{} author checked limited argument | Independent review:\allowbreak{} none | Full target solved:\allowbreak{} False


\clearpage

\section*{P0186  Largest coefficient dimension compatible with strong convergence}

\entrylabel{Primary source}

The strong convergence phenomenon,\allowbreak{}2025 survey;\allowbreak{} Section 6.6,\allowbreak{} p.68


\url{https://arxiv.org/pdf/2507.00346}


Source statement date:\allowbreak{} 2026-02-11. Archived access:\allowbreak{} 2026-09-06T20:\allowbreak{}51:\allowbreak{}12Z.


Historical:\allowbreak{} source\_\allowbreak{}says\_\allowbreak{}open | Current:\allowbreak{} unverified | Result:\allowbreak{} attempted\_\allowbreak{}unsolved


\entrylabel{Question}

Determine the optimal growth of matrix-coefficient dimension D\_\allowbreak{}N for strong convergence of polynomials in Gaussian or Haar random matrices.


\entrylabel{Scope}

Known subexponential sufficiency and failure at sufficiently large exp(C*N\textasciicircum{}2) leave a gap;\allowbreak{} Gaussian and Haar versions are grouped. New q-Gaussian model statements are not automatically the same fixed-model question.


\subsection*{P0186-A001 | Largest coefficient dimension compatible with strong convergence structural calculation}

\entrylabel{Argument or result}

Expected outlier counts can fail as a method even when the extreme eigenvalue converges in probability. Let E\_\allowbreak{}N have probability p\_\allowbreak{}N tending to0,\allowbreak{} and take X\_\allowbreak{}N=\allowbreak{}1\_\allowbreak{}(E\_\allowbreak{}N)*I\_\allowbreak{}(D\_\allowbreak{}N). Then its norm tends to0 in probability,\allowbreak{} but the expected number of eigenvalues above1/\allowbreak{}2 equals D\_\allowbreak{}N*p\_\allowbreak{}N. With D\_\allowbreak{}N=\allowbreak{}exp(N) and p\_\allowbreak{}N=\allowbreak{}exp(-N/\allowbreak{}2),\allowbreak{} that count diverges.


\entrylabel{Second check}

The event norm(X\_\allowbreak{}N)>1/\allowbreak{}2 is exactly E\_\allowbreak{}N,\allowbreak{} whereas the count on that event is exactly D\_\allowbreak{}N. This isolates multiplicity rather than a failure of norm convergence.


\entrylabel{Remaining obstruction}

The example does not refute strong convergence in the requested models. It explains why a trace-statistic proof needs a different mechanism once coefficient dimension creates large multiplicities.


Verification:\allowbreak{} author checked limited argument | Independent review:\allowbreak{} none | Full target solved:\allowbreak{} False


\clearpage

\section*{P0187  Optimal edge-fluctuation scales for arbitrary random-matrix polynomials}

\entrylabel{Primary source}

The strong convergence phenomenon,\allowbreak{}2025 survey;\allowbreak{} Section 6.7,\allowbreak{} pp.68-69


\url{https://arxiv.org/pdf/2507.00346}


Source statement date:\allowbreak{} 2026-02-11. Archived access:\allowbreak{} 2026-09-06T20:\allowbreak{}51:\allowbreak{}12Z.


Historical:\allowbreak{} source\_\allowbreak{}says\_\allowbreak{}open | Current:\allowbreak{} unverified | Result:\allowbreak{} attempted\_\allowbreak{}unsolved


\entrylabel{Question}

Under suitable nondegeneracy assumptions,\allowbreak{} establish the N\textasciicircum{}(-2/\allowbreak{}3) fluctuation scale and edge law for norms of arbitrary-degree polynomials of random matrices.


\entrylabel{Scope}

The source explicitly excludes degenerate situations and separates deterministic spikes,\allowbreak{} which may produce N\textasciicircum{}(-1/\allowbreak{}2) fluctuations.


\subsection*{P0187-A001 | Optimal edge-fluctuation scales for arbitrary random-matrix polynomials structural calculation}

\entrylabel{Argument or result}

The nondegeneracy qualifier is necessary even for a one-variable quadratic polynomial. For a standard GUE matrix G\_\allowbreak{}N let P(x)=\allowbreak{}20-(x-2)\textasciicircum{}2. With probability tending to1 its spectrum lies in a fixed neighborhood of[-2,\allowbreak{}2] where P is positive. Then norm(P(G\_\allowbreak{}N))=\allowbreak{}20-min\_\allowbreak{}i(lambda\_\allowbreak{}i-2)\textasciicircum{}2. Since the largest GUE eigenvalue is2+\allowbreak{}O\_\allowbreak{}p(N\textasciicircum{}-2/\allowbreak{}3),\allowbreak{} one has0<=\allowbreak{}20-norm(P(G\_\allowbreak{}N))<=\allowbreak{}O\_\allowbreak{}p(N\textasciicircum{}-4/\allowbreak{}3).


\entrylabel{Second check}

The upper bound uses the largest eigenvalue as one candidate in the minimum. Spectral mapping and positivity give the norm identity. Consequently N\textasciicircum{}(2/\allowbreak{}3)*(norm(P(G\_\allowbreak{}N))-20) tends to0,\allowbreak{} rather than a nondegenerate Tracy-Widom variable.


\entrylabel{Remaining obstruction}

The example is deliberately degenerate and does not contradict the stated research direction. A correct general theorem must identify regularity of the polynomial edge before asserting the universal scale.


Verification:\allowbreak{} author checked limited argument | Independent review:\allowbreak{} none | Full target solved:\allowbreak{} False


\clearpage

\section*{P0188  Joint strong convergence for permutation and deterministic matrices}

\entrylabel{Primary source}

The strong convergence phenomenon,\allowbreak{}2025 survey;\allowbreak{} Section 6.8,\allowbreak{} p.69


\url{https://arxiv.org/pdf/2507.00346}


Source statement date:\allowbreak{} 2026-02-11. Archived access:\allowbreak{} 2026-09-06T20:\allowbreak{}51:\allowbreak{}12Z.


Historical:\allowbreak{} source\_\allowbreak{}says\_\allowbreak{}open | Current:\allowbreak{} unverified | Result:\allowbreak{} attempted\_\allowbreak{}unsolved


\entrylabel{Question}

Develop a permutation-matrix counterpart of joint strong convergence with strongly convergent deterministic matrices.


\entrylabel{Scope}

Removing the common constant vector of permutation matrices requires an explicit convention for deterministic matrices that do not preserve its orthogonal complement.


\subsection*{P0188-A001 | Joint strong convergence for permutation and deterministic matrices structural calculation}

\entrylabel{Argument or result}

Naively compressing arbitrary deterministic matrices creates persistent defects. Let e be the normalized all-ones vector,\allowbreak{} P=\allowbreak{}I-ee*,\allowbreak{} and B a balanced diagonal matrix with entries plus/\allowbreak{}minus1. B strongly converges to the symmetric Bernoulli law. But on e-perp,\allowbreak{} PBP has the zero eigenvector Be,\allowbreak{} while all its other eigenvalues are plus/\allowbreak{}minus1. Thus norm(I\_\allowbreak{}(e-perp)-(PBP)\textasciicircum{}2)=\allowbreak{}1,\allowbreak{} whereas the corresponding Bernoulli-limit polynomial is zero.


\entrylabel{Second check}

Because B\textasciicircum{}2=\allowbreak{}I and e*Be=\allowbreak{}0,\allowbreak{} PBP(Be)=\allowbreak{}PB(e)=\allowbreak{}Pe=\allowbreak{}0. More generally (PBP)\textasciicircum{}2-PB\textasciicircum{}2P=\allowbreak{}-PBee*BP,\allowbreak{} whose norm here is1. Its normalized trace tends to0 but its operator norm does not.


\entrylabel{Remaining obstruction}

This refutes the naive unchanged-deterministic-limit formulation after compression,\allowbreak{} not every possible counterpart requested by the survey. Appropriate invariant subspaces or additional deterministic hypotheses remain to be formulated and proved.


Verification:\allowbreak{} author checked limited argument | Independent review:\allowbreak{} none | Full target solved:\allowbreak{} False


\clearpage

\section*{P0189  Brown-measure convergence from general strong convergence}

\entrylabel{Primary source}

The strong convergence phenomenon,\allowbreak{}2025 survey;\allowbreak{} Section 6.9,\allowbreak{} pp.69-70


\url{https://arxiv.org/pdf/2507.00346}


Source statement date:\allowbreak{} 2026-02-11. Archived access:\allowbreak{} 2026-09-06T20:\allowbreak{}51:\allowbreak{}12Z.


Historical:\allowbreak{} source\_\allowbreak{}says\_\allowbreak{}open | Current:\allowbreak{} unverified | Result:\allowbreak{} candidate\_\allowbreak{}full\_\allowbreak{}solution


\entrylabel{Question}

Does strong convergence of a family of random matrices imply convergence of empirical complex-eigenvalue measures of every noncommutative polynomial to its limiting Brown measure?


\entrylabel{Scope}

This is the literal general statement. It does not impose independent entries,\allowbreak{} independence of the matrix family,\allowbreak{} or a free limiting family. The narrower Ginibre-with-adjoints question is recorded as a remaining subclass,\allowbreak{} not counted again.


\subsection*{P0189-A001 | Brown-measure convergence from general strong convergence structural calculation}

\entrylabel{Argument or result}

Complete candidate counterexample:\allowbreak{} rational clock/\allowbreak{}shift unitary pairs U\_\allowbreak{}j,\allowbreak{}V\_\allowbreak{}j approach an irrational rotation pair u,\allowbreak{}v strongly,\allowbreak{} while U\_\allowbreak{}j\textasciicircum{}q=\allowbreak{}I and V\_\allowbreak{}j\textasciicircum{}q=\allowbreak{}-I force det(zI-U\_\allowbreak{}j-V\_\allowbreak{}j)=\allowbreak{}z\textasciicircum{}q. Hence every finite sum is nilpotent with empirical measure delta\_\allowbreak{}0. The established Brown measure of u+\allowbreak{}v is Haar measure on the unit circle. Simultaneous Haar conjugation makes the finite family random without changing any of these statements.


\entrylabel{Second check}

The strong-convergence proof maps the simple irrational rotation algebra into the product of matrix algebras modulo norm-null sequences;\allowbreak{} injectivity gives every subsequential limsup polynomial norm and therefore the full limit. Normalized word traces are checked directly. The determinant calculation independently checks nilpotency. Full proof and all source dependencies are in the proof packet.


\entrylabel{Remaining obstruction}

No gap presently identified in the literal-general-scope construction,\allowbreak{} but there is no independent review and no novelty claim. It does not resolve independent Ginibre polynomials with adjoints;\allowbreak{} the June2026 holomorphic-Ginibre theorem is compatible with this example.


Verification:\allowbreak{} author checked limited argument | Independent review:\allowbreak{} none | Full target solved:\allowbreak{} False


\entrylabel{Full proof packet}

proofs/\allowbreak{}P0189\_\allowbreak{}brown\_\allowbreak{}measure\_\allowbreak{}strong\_\allowbreak{}convergence\_\allowbreak{}counterexample\_\allowbreak{}candidate.md


\clearpage

\section*{P0190  Intrinsic freeness for non-Hermitian eigenvalues}

\entrylabel{Primary source}

The strong convergence phenomenon,\allowbreak{}2025 survey;\allowbreak{} Section 6.9,\allowbreak{} p.70


\url{https://arxiv.org/pdf/2507.00346}


Source statement date:\allowbreak{} 2026-02-11. Archived access:\allowbreak{} 2026-09-06T20:\allowbreak{}51:\allowbreak{}12Z.


Historical:\allowbreak{} source\_\allowbreak{}says\_\allowbreak{}open | Current:\allowbreak{} unverified | Result:\allowbreak{} attempted\_\allowbreak{}unsolved


\entrylabel{Question}

Find a non-Hermitian complex-eigenvalue counterpart of intrinsic freeness for general nonhomogeneous random matrices,\allowbreak{} and identify the appropriate comparison objects.


\entrylabel{Scope}

This concerns nonasymptotic comparison of broad covariance models,\allowbreak{} distinct from the abstract continuity questionP0189.


\subsection*{P0190-A001 | Intrinsic freeness for non-Hermitian eigenvalues structural calculation}

\entrylabel{Argument or result}

Hermitization exposes a missing hypothesis. Weak singular-value convergence and bounded largest singular value do not control the logarithmic potential:\allowbreak{} nu\_\allowbreak{}N=\allowbreak{}(1-1/\allowbreak{}N)*delta\_\allowbreak{}1+\allowbreak{}(1/\allowbreak{}N)*delta\_\allowbreak{}(exp(-N)) converges weakly to delta\_\allowbreak{}1 and has bounded support,\allowbreak{} but integral log(s) dnu\_\allowbreak{}N=\allowbreak{}-1 rather than0. Thus a comparison based only on bounded polynomial spectral statistics cannot justify passage through the logarithm.


\entrylabel{Second check}

All positive moments converge to1,\allowbreak{} while the single tiny singular value contributes(1/\allowbreak{}N)*(-N)=\allowbreak{}-1. Uniform integrability of the negative logarithm would remove this particular obstruction.


\entrylabel{Remaining obstruction}

The scalar example is an obstruction to a proof step,\allowbreak{} not itself a counterexample to empirical eigenvalue convergence. Quantitative least-singular-value control adapted to intrinsic covariance models is still missing.


Verification:\allowbreak{} author checked limited argument | Independent review:\allowbreak{} none | Full target solved:\allowbreak{} False


\clearpage

\section*{P0191  Least singular values at the polynomial constant term}

\entrylabel{Primary source}

Yi Han:\allowbreak{} Brown measure convergence for the spectrum of polynomials in Ginibre matrices;\allowbreak{} Question 1.6


\url{https://arxiv.org/pdf/2606.01664}


Source statement date:\allowbreak{} 2026-06-01. Archived access:\allowbreak{} 2026-09-06T21:\allowbreak{}27:\allowbreak{}13Z.


Historical:\allowbreak{} source\_\allowbreak{}says\_\allowbreak{}open | Current:\allowbreak{} unverified | Result:\allowbreak{} attempted\_\allowbreak{}unsolved


\entrylabel{Question}

Extend the polynomial least-singular-value bound for a nonzero holomorphic polynomial of independent complex Ginibre matrices to the shift equal to its constant term.


\entrylabel{Scope}

The source theorem excludes that single shift. Qualitative invertibility and a quantitative polynomial lower-tail estimate are kept distinct.


\subsection*{P0191-A001 | Least singular values at the polynomial constant term structural calculation}

\entrylabel{Argument or result}

For a monomial P=\allowbreak{}X\_\allowbreak{}(i1)...X\_\allowbreak{}(id),\allowbreak{} the missing shift is0 and the estimate reduces to the individual matrices. Since smin(P)>=\allowbreak{}product\_\allowbreak{}k smin(X\_\allowbreak{}(ik)),\allowbreak{} the event smin(P)<=\allowbreak{}epsilon implies at least one factor has smin<=\allowbreak{}epsilon\textasciicircum{}(1/\allowbreak{}d). The standard complex-Ginibre bound then gives P(smin(P)<=\allowbreak{}epsilon)<=\allowbreak{}C*d*N\textasciicircum{}2*epsilon\textasciicircum{}(2/\allowbreak{}d). Repeated factors require no independence in this union bound.


\entrylabel{Second check}

For normalized complex Ginibre,\allowbreak{} N*smin has the usual bounded quadratic lower tail. Multiply the deterministic singular-value inequalities successively,\allowbreak{} and apply the union bound over the d factor positions.


\entrylabel{Remaining obstruction}

Sums of monomials can cancel,\allowbreak{} so this argument gives no least-singular-value estimate for an arbitrary noncommutative polynomial at its constant term.


Verification:\allowbreak{} author checked limited argument | Independent review:\allowbreak{} none | Full target solved:\allowbreak{} False


\clearpage

\section*{P0192  Brown-measure universality beyond complex entry densities}

\entrylabel{Primary source}

Yi Han:\allowbreak{} Brown measure convergence for the spectrum of polynomials in Ginibre matrices;\allowbreak{} Question 1.7 and first part of Question 1.8


\url{https://arxiv.org/pdf/2606.01664}


Source statement date:\allowbreak{} 2026-06-01. Archived access:\allowbreak{} 2026-09-06T21:\allowbreak{}27:\allowbreak{}13Z.


Historical:\allowbreak{} source\_\allowbreak{}says\_\allowbreak{}open | Current:\allowbreak{} unverified | Result:\allowbreak{} attempted\_\allowbreak{}unsolved


\entrylabel{Question}

Extend Brown-measure convergence for holomorphic noncommutative polynomials to real Gaussian or discrete entry laws such as Bernoulli,\allowbreak{} allowing complex polynomial coefficients.


\entrylabel{Scope}

Real-valued and discrete entry versions are grouped as one source universality family. The source bounded-density-on-C theorem does not cover either automatically.


\subsection*{P0192-A001 | Brown-measure universality beyond complex entry densities structural calculation}

\entrylabel{Argument or result}

Every one-matrix holomorphic polynomial is already reducible to the circular law. If X\_\allowbreak{}N has a circular-law limit and P is a fixed complex polynomial,\allowbreak{} spectral mapping gives mu\_\allowbreak{}(P(X\_\allowbreak{}N))=\allowbreak{}P\_\allowbreak{}pushforward(mu\_\allowbreak{}(X\_\allowbreak{}N)) exactly,\allowbreak{} including algebraic multiplicities. The continuous mapping theorem therefore yields the pushforward of the circular law. The Brown measure of P(c) is the same pushforward,\allowbreak{} by factorizing P(z)-lambda and using multiplicativity of the Fuglede-Kadison determinant.


\entrylabel{Second check}

For each test function h bounded and continuous,\allowbreak{} h composed with P is bounded and continuous,\allowbreak{} so weak convergence passes through P. At the logarithmic-potential level,\allowbreak{} P(c)-lambda factors into commuting(c-root\_\allowbreak{}j) terms,\allowbreak{} and determinant multiplicativity gives the same transformed Brown measure. Existing real-Gaussian and Bernoulli circular laws supply the initial one-matrix convergence.


\entrylabel{Remaining obstruction}

For several noncommuting input matrices there is no joint eigenvalue spectral mapping. Arithmetic and structural dependence in least-singular-value estimates remain unresolved here.


Verification:\allowbreak{} author checked limited argument | Independent review:\allowbreak{} none | Full target solved:\allowbreak{} False


\clearpage

\section*{P0193  Eventual log-unimodality from arbitrary compact positive initial laws}

\entrylabel{Primary source}

Hasebe-Ueda-Wang:\allowbreak{} Log-unimodality for free positive multiplicative Brownian motion;\allowbreak{} Problem 4.7,\allowbreak{} PDF p.11;\allowbreak{} Lemma 4.1,\allowbreak{} p.8


\url{https://arxiv.org/pdf/2009.13848}


Source statement date:\allowbreak{} 2020-09-29. Archived access:\allowbreak{} 2026-09-06T21:\allowbreak{}33:\allowbreak{}40Z.


Historical:\allowbreak{} source\_\allowbreak{}says\_\allowbreak{}open | Current:\allowbreak{} unverified | Result:\allowbreak{} candidate\_\allowbreak{}full\_\allowbreak{}solution


\entrylabel{Question}

If nu is supported on a compact interval [alpha,\allowbreak{}beta] with 0<alpha<beta,\allowbreak{} must sigma\_\allowbreak{}t free-multiplied by nu be log-unimodal for all sufficiently large times?


\entrylabel{Scope}

This asks whether the support-ratio restriction in Theorem 4.4 can be removed. Log-unimodality is unimodality after taking logarithms. The known positive theorem assumes a multiplicatively symmetric log-unimodal initial law;\allowbreak{} the question here allows non-log-unimodal initial laws.


\entrylabel{Dated literature check}

Focused primary-source searches2026-09-06 found no verified later resolution. The arXiv page has one2020version and a2022journal reference. Author publication pages and Auer2026 concern related results,\allowbreak{} not a verified resolution here. Search absence is not current-open evidence.


\subsection*{P0193-A001 | Four exact crossings of the multiplicative Cauchy-profile criterion}

\entrylabel{Argument or result}

Candidate negative answer:\allowbreak{} nu=\allowbreak{}(delta\_\allowbreak{}(1/\allowbreak{}10)+\allowbreak{}delta\_\allowbreak{}10)/\allowbreak{}2 yields failure of log-unimodality for every t>=\allowbreak{}5*pi. Set a=\allowbreak{}log10 and K\_\allowbreak{}R(s)=\allowbreak{}1/\allowbreak{}[2*(cosh(s)-cosR)],\allowbreak{} F\_\allowbreak{}R(u)=\allowbreak{}(K\_\allowbreak{}R(u-a)+\allowbreak{}K\_\allowbreak{}R(u+\allowbreak{}a))/\allowbreak{}2. Choose the unique R in[pi/\allowbreak{}2,\allowbreak{}pi) such that R/\allowbreak{}sinR=\allowbreak{}t/\allowbreak{}10. Lemma4.1 then asks whether F\_\allowbreak{}R(u)=\allowbreak{}1/\allowbreak{}10 has at most two roots. But F\_\allowbreak{}R(0)<=\allowbreak{}10/\allowbreak{}101<1/\allowbreak{}10,\allowbreak{} F\_\allowbreak{}R(+\allowbreak{}-a)>=\allowbreak{}1/\allowbreak{}8>1/\allowbreak{}10,\allowbreak{} and F\_\allowbreak{}R tends to0 at both infinities. Continuity gives four distinct roots,\allowbreak{} contradicting the necessary condition for log-unimodality.


\entrylabel{Second check}

Same-author second algebra route in r:\allowbreak{} at r=\allowbreak{}1 the integral is10/\allowbreak{}(101-20*cosR)<=\allowbreak{}10/\allowbreak{}101;\allowbreak{} at r=\allowbreak{}10 and1/\allowbreak{}10 the resonant atom alone contributes1/\allowbreak{}[4*(1-cosR)]>=\allowbreak{}1/\allowbreak{}8;\allowbreak{} endpoint limits are0. The four roots lie in(0,\allowbreak{}1/\allowbreak{}10),\allowbreak{}(1/\allowbreak{}10,\allowbreak{}1),\allowbreak{}(1,\allowbreak{}10),\allowbreak{}(10,\allowbreak{}infinity). Monotonicity of R/\allowbreak{}sinR follows from(sinR-R*cosR)prime=\allowbreak{}R*sinR>0.


\entrylabel{Remaining obstruction}

No gap currently identified in the full written argument. Later-literature novelty and independent review remain unverified. Check the exact necessity direction in Lemma4.1,\allowbreak{} the time normalization,\allowbreak{} and whether the published question changed from the archived version before promotion.


Verification:\allowbreak{} author checked complete candidate | Independent review:\allowbreak{} none | Full target solved:\allowbreak{} False


\entrylabel{Full proof packet}

proofs/\allowbreak{}P0193\_\allowbreak{}compact\_\allowbreak{}multiplicative\_\allowbreak{}nonunimodality\_\allowbreak{}candidate.md


\clearpage

\section*{P0194  Polynomial mixing with plus Ising boundary conditions}

\entrylabel{Primary source}

Levin-Peres:\allowbreak{} Markov Chains and Mixing Times,\allowbreak{} second edition;\allowbreak{} Chapter26 Question1,\allowbreak{} PDFp375


\url{https://pages.uoregon.edu/dlevin/MARKOV/mcmt2e.pdf}


Source statement date:\allowbreak{} 2017. Archived access:\allowbreak{} 2026-09-06T21:\allowbreak{}43:\allowbreak{}30Z.


Historical:\allowbreak{} source\_\allowbreak{}says\_\allowbreak{}open | Current:\allowbreak{} unverified | Result:\allowbreak{} attempted\_\allowbreak{}unsolved


\entrylabel{Question}

Is Glauber mixing on the n-by-n square with plus boundary polynomial in n at every fixed positive temperature?


\entrylabel{Scope}

The2026 phase-ordering paper still uses a quasi-polynomial plus-boundary estimate. Its biased random initialization theorem is not a worst-start plus-boundary polynomial bound. Critical and high-temperature established cases are not counted separately.


\subsection*{P0194-A001 | Polynomial mixing with plus Ising boundary conditions explicit structural calculation}

\entrylabel{Argument or result}

The direct disagreement coupling proves the high-temperature part. Write N=\allowbreak{}n\textasciicircum{}2 and c=\allowbreak{}4*tanh(beta). A single disagreement vanishes if its site is updated,\allowbreak{} while at most four neighboring updates create disagreement with probability at most tanh(beta). Thus the expected Hamming distance contracts by1-(1-c)/\allowbreak{}N when c<1. Iteration and the coupling inequality give tmix(1/\allowbreak{}4)<=\allowbreak{}N*log(4*N)/\allowbreak{}(1-c),\allowbreak{} up to integer rounding.


\entrylabel{Second check}

Changing one neighboring spin changes the heat-bath plus probability by at most tanh(beta),\allowbreak{} since the conditional probability is(1+\allowbreak{}tanh(h))/\allowbreak{}2. A path coupling along single-spin differences extends the bound to arbitrary initial pairs. Boundary sites have at most four variable neighbors,\allowbreak{} so the estimate remains valid with fixed plus boundary.


\entrylabel{Remaining obstruction}

For4*tanh(beta)>=\allowbreak{}1 the coefficient ceases to contract. This is a failure of the local metric argument,\allowbreak{} not a lower bound showing slow mixing;\allowbreak{} the low-temperature interface regime remains unresolved here.


Verification:\allowbreak{} author checked limited argument | Independent review:\allowbreak{} none | Full target solved:\allowbreak{} False


\clearpage

\section*{P0195  Monotonicity of the ferromagnetic Ising spectral gap}

\entrylabel{Primary source}

Levin-Peres:\allowbreak{} Markov Chains and Mixing Times,\allowbreak{} second edition;\allowbreak{} Chapter26 Question2,\allowbreak{} PDFp375


\url{https://pages.uoregon.edu/dlevin/MARKOV/mcmt2e.pdf}


Source statement date:\allowbreak{} 2017. Archived access:\allowbreak{} 2026-09-06T21:\allowbreak{}43:\allowbreak{}30Z.


Historical:\allowbreak{} source\_\allowbreak{}says\_\allowbreak{}open | Current:\allowbreak{} unverified | Result:\allowbreak{} attempted\_\allowbreak{}unsolved


\entrylabel{Question}

Does increasing any ferromagnetic interaction decrease the Glauber spectral gap? What stronger time-dependent total-variation monotonicity holds?


\entrylabel{Scope}

Temperature,\allowbreak{} edge addition and individual-coupling versions are grouped. The source has zero field in its general formula. Cycle and mean-field special cases are established in the cited literature.


\subsection*{P0195-A001 | Monotonicity of the ferromagnetic Ising spectral gap explicit structural calculation}

\entrylabel{Argument or result}

The two-spin calculation supports the conjecture exactly. With interaction J>=\allowbreak{}0,\allowbreak{} let c=\allowbreak{}tanh(J) and update either spin with probability1/\allowbreak{}2. The eigenfunctions1,\allowbreak{} sigma1+\allowbreak{}sigma2,\allowbreak{} sigma1-sigma2,\allowbreak{} sigma1*sigma2-c have eigenvalues1,\allowbreak{}(1+\allowbreak{}c)/\allowbreak{}2,\allowbreak{}(1-c)/\allowbreak{}2,\allowbreak{}0. Hence the spectral gap is(1-tanh(J))/\allowbreak{}2,\allowbreak{} strictly decreasing for finite J.


\entrylabel{Second check}

Conditionally E[sigma1|sigma2]=\allowbreak{}c*sigma2. Applying the update kernel to the four displayed basis functions gives the eigenvalues directly. The stationary mean of sigma1*sigma2 is c,\allowbreak{} so the last function is centered.


\entrylabel{Remaining obstruction}

For a general graph,\allowbreak{} increasing J changes both the reversible measure and the Dirichlet form. Pointwise comparison of transition probabilities alone does not order the Rayleigh quotient because its variance denominator also changes.


Verification:\allowbreak{} author checked limited argument | Independent review:\allowbreak{} none | Full target solved:\allowbreak{} False


\clearpage

\section*{P0196  Sharp comparison of systematic and random Ising updates}

\entrylabel{Primary source}

Levin-Peres:\allowbreak{} Markov Chains and Mixing Times,\allowbreak{} second edition;\allowbreak{} Chapter26 Question3,\allowbreak{} PDFpp375-376


\url{https://pages.uoregon.edu/dlevin/MARKOV/mcmt2e.pdf}


Source statement date:\allowbreak{} 2017. Archived access:\allowbreak{} 2026-09-06T21:\allowbreak{}43:\allowbreak{}30Z.


Historical:\allowbreak{} source\_\allowbreak{}says\_\allowbreak{}open | Current:\allowbreak{} unverified | Result:\allowbreak{} attempted\_\allowbreak{}unsolved


\entrylabel{Question}

For ferromagnetic Ising dynamics,\allowbreak{} compare the number of updates in a systematic sweep order with random-scan mixing,\allowbreak{} with universal upper constant and lower logarithmic loss as posed in the source.


\entrylabel{Scope}

The two inequalities are one comparison problem. General Gibbs-system polynomial comparison is established by Gaitonde-Mossel2025;\allowbreak{} their counterexamples for arbitrary target laws do not by themselves refute the ferromagnetic Ising scope.


\subsection*{P0196-A001 | Sharp comparison of systematic and random Ising updates explicit structural calculation}

\entrylabel{Argument or result}

At infinite temperature one deterministic sweep gives an exact independent fair-spin sample. Random scan has mixing time asymptotic to(N/\allowbreak{}2)*log N. The upper bound follows from Fourier eigenvalues1-|S|/\allowbreak{}N:\allowbreak{} chi-square distance is sum\_\allowbreak{}(k>=\allowbreak{}1) binom(N,\allowbreak{}k)*(1-k/\allowbreak{}N)\textasciicircum{}(2t) <=\allowbreak{}(1+\allowbreak{}exp(-2t/\allowbreak{}N))\textasciicircum{}N-1. For the lower bound,\allowbreak{} magnetization from all plus has mean N*(1-1/\allowbreak{}N)\textasciicircum{}t and variance at most N,\allowbreak{} whereas stationary magnetization has variance N. Thus at t=\allowbreak{}(N/\allowbreak{}2)*(log N-C) the means are separated by order exp(C/\allowbreak{}2)*sqrt(N).


\entrylabel{Second check}

For distinct spins,\allowbreak{} E[sigma\_\allowbreak{}i*sigma\_\allowbreak{}j]=\allowbreak{}(1-2/\allowbreak{}N)\textasciicircum{}t <=\allowbreak{}(1-1/\allowbreak{}N)\textasciicircum{}(2t),\allowbreak{} so their covariance is nonpositive and the variance bound follows. Chebyshev separates the two magnetization laws as C grows. This independently recovers the factor log N between random updates and one sweep.


\entrylabel{Remaining obstruction}

The product case shows why the logarithmic allowance is necessary. Interacting conditional-expectation projections fail to commute,\allowbreak{} so their product cannot be replaced by the power of their average.


Verification:\allowbreak{} author checked limited argument | Independent review:\allowbreak{} none | Full target solved:\allowbreak{} False


\clearpage

\section*{P0197  Relation between Ising relaxation and mixing on transitive graphs}

\entrylabel{Primary source}

Levin-Peres:\allowbreak{} Markov Chains and Mixing Times,\allowbreak{} second edition;\allowbreak{} Chapter26 Question4,\allowbreak{} PDFp376


\url{https://pages.uoregon.edu/dlevin/MARKOV/mcmt2e.pdf}


Source statement date:\allowbreak{} 2017. Archived access:\allowbreak{} 2026-09-06T21:\allowbreak{}43:\allowbreak{}30Z.


Historical:\allowbreak{} source\_\allowbreak{}says\_\allowbreak{}open | Current:\allowbreak{} unverified | Result:\allowbreak{} attempted\_\allowbreak{}unsolved


\entrylabel{Question}

On transitive graph families,\allowbreak{} is Glauber relaxation of order N equivalent to mixing of order N*log N as temperature varies?


\entrylabel{Scope}

Time counts individual random-site updates. This concerns a relation of orders,\allowbreak{} distinct from the separate question of whether there is cutoff.


\subsection*{P0197-A001 | Relation between Ising relaxation and mixing on transitive graphs explicit structural calculation}

\entrylabel{Argument or result}

At zero interaction the heat-bath operator has eigenvalues1-k/\allowbreak{}N with multiplicity binom(N,\allowbreak{}k),\allowbreak{} so trel=\allowbreak{}N and tmix is asymptotic to(N/\allowbreak{}2)*log N. This exhibits the proposed relation,\allowbreak{} and also shows that the logarithm comes from many coordinate modes rather than a single slow eigenfunction. A gap estimate alone gives only tmix<=\allowbreak{}trel*log(4/\allowbreak{}pi\_\allowbreak{}min);\allowbreak{} at beta0,\allowbreak{} pi\_\allowbreak{}min=\allowbreak{}2\textasciicircum{}-N,\allowbreak{} which overestimates by a factor of order N/\allowbreak{}log N.


\entrylabel{Second check}

The largest nonconstant eigenvalue is1-1/\allowbreak{}N. Summing the full spectral multiplicities yields the sharper chi-square expression(1+\allowbreak{}exp(-2t/\allowbreak{}N))\textasciicircum{}N-1 as an upper bound;\allowbreak{} retaining only the gap loses the binomial structure.


\entrylabel{Remaining obstruction}

Transitivity of vertices does not itself identify the multiplicities or nonlinear slow observables of the interacting spin configuration chain. A comparison preserving this spectral structure is missing.


Verification:\allowbreak{} author checked limited argument | Independent review:\allowbreak{} none | Full target solved:\allowbreak{} False


\clearpage

\section*{P0198  Cutoff on bounded-degree transitive expanders}

\entrylabel{Primary source}

Levin-Peres:\allowbreak{} Markov Chains and Mixing Times,\allowbreak{} second edition;\allowbreak{} Chapter26 Question5,\allowbreak{} PDFp376


\url{https://pages.uoregon.edu/dlevin/MARKOV/mcmt2e.pdf}


Source statement date:\allowbreak{} 2017. Archived access:\allowbreak{} 2026-09-06T21:\allowbreak{}43:\allowbreak{}30Z.


Historical:\allowbreak{} source\_\allowbreak{}says\_\allowbreak{}open | Current:\allowbreak{} unverified | Result:\allowbreak{} attempted\_\allowbreak{}unsolved


\entrylabel{Question}

Must lazy random walk on every bounded-degree vertex-transitive expander sequence exhibit total-variation cutoff?


\entrylabel{Scope}

Salez2025 repeats this as a current challenge. Optimal Ramanujan expansion is an established special case. This is not the random-Cayley spectral-gap question P0181.


\subsection*{P0198-A001 | Cutoff on bounded-degree transitive expanders explicit structural calculation}

\entrylabel{Argument or result}

The assumptions already imply mixing of order log N and the usual product condition. From any vertex,\allowbreak{} after t steps the support lies in a ball of at most(Delta+\allowbreak{}1)\textasciicircum{}t vertices,\allowbreak{} so total variation is at least1-(Delta+\allowbreak{}1)\textasciicircum{}t/\allowbreak{}N. Thus tmix(1/\allowbreak{}4)>=\allowbreak{}log(3*N/\allowbreak{}4)/\allowbreak{}log(Delta+\allowbreak{}1). A uniform positive lazy spectral gap gamma gives tmix(1/\allowbreak{}4)<=\allowbreak{}gamma\textasciicircum{}-1*log(4*N),\allowbreak{} while trel<=\allowbreak{}gamma\textasciicircum{}-1.


\entrylabel{Second check}

The lower bound tests the reachable set against uniform stationarity;\allowbreak{} transitivity makes the graph regular. The upper bound is the standard reversible L2 contraction. Their ratio is bounded and trel/\allowbreak{}tmix tends to0.


\entrylabel{Remaining obstruction}

These bounds establish only a bounded transition-time ratio,\allowbreak{} not a ratio tending to1. They provide no concentration of the relevant entropic or hitting-time clock,\allowbreak{} which is the unresolved step.


Verification:\allowbreak{} author checked limited argument | Independent review:\allowbreak{} none | Full target solved:\allowbreak{} False


\clearpage

\section*{P0199  Cutoff for cyclic-to-random transpositions}

\entrylabel{Primary source}

Levin-Peres:\allowbreak{} Markov Chains and Mixing Times,\allowbreak{} second edition;\allowbreak{} Chapter26 Question6(a),\allowbreak{} PDFp376


\url{https://pages.uoregon.edu/dlevin/MARKOV/mcmt2e.pdf}


Source statement date:\allowbreak{} 2017. Archived access:\allowbreak{} 2026-09-06T21:\allowbreak{}43:\allowbreak{}30Z.


Historical:\allowbreak{} source\_\allowbreak{}says\_\allowbreak{}open | Current:\allowbreak{} unverified | Result:\allowbreak{} attempted\_\allowbreak{}unsolved


\entrylabel{Question}

Does the shuffle that swaps the cyclically chosen position with a uniform random position have cutoff?


\entrylabel{Scope}

This is a time-inhomogeneous deck chain or a chain sampled by full sweeps. It differs from the cyclic adjacent-transposition shuffle,\allowbreak{} whose cutoff was proved by Nam-Nestoridi.


\subsection*{P0199-A001 | Cutoff for cyclic-to-random transpositions explicit structural calculation}

\entrylabel{Argument or result}

The one-card marginal admits an exact update that may identify candidate slow modes. If position i is chosen deterministically and the tracked-card probabilities are p\_\allowbreak{}j,\allowbreak{} then p\_\allowbreak{}i becomes1/\allowbreak{}N,\allowbreak{} while p\_\allowbreak{}j becomes(1-1/\allowbreak{}N)*p\_\allowbreak{}j+\allowbreak{}p\_\allowbreak{}i/\allowbreak{}N for j different from i. Thus a centered vector z has new z\_\allowbreak{}i=\allowbreak{}0 and new z\_\allowbreak{}j=\allowbreak{}(1-1/\allowbreak{}N)*z\_\allowbreak{}j+\allowbreak{}z\_\allowbreak{}i/\allowbreak{}N. A sweep is the explicit product of these N linear maps.


\entrylabel{Second check}

If the card is at i it moves uniformly;\allowbreak{} if it is at j different from i,\allowbreak{} it stays unless the random partner is j. Summing the two cases gives the formula,\allowbreak{} which preserves the uniform distribution. Adding a deterministic phase as a state coordinate would give a periodic phase marginal and cannot justify ordinary mixing of the augmented chain.


\entrylabel{Remaining obstruction}

The one-card product can supply lower-bound observables,\allowbreak{} but does not control correlations of all cards or produce a matching sharp upper bound for the full permutation.


Verification:\allowbreak{} author checked limited argument | Independent review:\allowbreak{} none | Full target solved:\allowbreak{} False


\clearpage

\section*{P0200  Cutoff for card-cyclic random insertion}

\entrylabel{Primary source}

Levin-Peres:\allowbreak{} Markov Chains and Mixing Times,\allowbreak{} second edition;\allowbreak{} Chapter26 Question6(c),\allowbreak{} PDFp376


\url{https://pages.uoregon.edu/dlevin/MARKOV/mcmt2e.pdf}


Source statement date:\allowbreak{} 2017. Archived access:\allowbreak{} 2026-09-06T21:\allowbreak{}43:\allowbreak{}30Z.


Historical:\allowbreak{} source\_\allowbreak{}says\_\allowbreak{}open | Current:\allowbreak{} unverified | Result:\allowbreak{} attempted\_\allowbreak{}unsolved


\entrylabel{Question}

Does the shuffle that cyclically chooses a card label,\allowbreak{} removes that card and reinserts it uniformly have cutoff?


\entrylabel{Scope}

Random-to-random insertion is a different shuffle and is resolved. Here the selected label,\allowbreak{} rather than its present position,\allowbreak{} advances cyclically.


\subsection*{P0200-A001 | Cutoff for card-cyclic random insertion explicit structural calculation}

\entrylabel{Argument or result}

A coupon-collector strong-stationary-time argument fails even after every label has been updated. For three cards initially123,\allowbreak{} update labels1,\allowbreak{}2,\allowbreak{}3 once each. The final probabilities of123,\allowbreak{}132,\allowbreak{}312 are5/\allowbreak{}27 each,\allowbreak{} and those of213,\allowbreak{}231,\allowbreak{}321 are4/\allowbreak{}27 each,\allowbreak{} so the total-variation distance from uniform is1/\allowbreak{}18 after a complete sweep.


\entrylabel{Second check}

After labels1 and2,\allowbreak{} the states213,\allowbreak{}123,\allowbreak{}132 each have probability2/\allowbreak{}9 and231,\allowbreak{}321,\allowbreak{}312 each have probability1/\allowbreak{}9. Deleting card3 leaves relative order12 with probability5/\allowbreak{}9 and21 with probability4/\allowbreak{}9. Uniform insertion of3 gives the six stated probabilities,\allowbreak{} summing to1.


\entrylabel{Remaining obstruction}

Complete update coverage is not exact stationarity because the updated card positions remain dependent. A large-N contraction or spectral mechanism with a sharp window is not obtained.


Verification:\allowbreak{} author checked limited argument | Independent review:\allowbreak{} none | Full target solved:\allowbreak{} False


\clearpage

\section*{P0201  Optimal Glauber mixing near the coloring degree threshold}

\entrylabel{Primary source}

Levin-Peres:\allowbreak{} Markov Chains and Mixing Times,\allowbreak{} second edition;\allowbreak{} Chapter26 Question7,\allowbreak{} PDFp376


\url{https://pages.uoregon.edu/dlevin/MARKOV/mcmt2e.pdf}


Source statement date:\allowbreak{} 2017. Archived access:\allowbreak{} 2026-09-06T21:\allowbreak{}43:\allowbreak{}30Z.


Historical:\allowbreak{} source\_\allowbreak{}says\_\allowbreak{}open | Current:\allowbreak{} unverified | Result:\allowbreak{} attempted\_\allowbreak{}unsolved


\entrylabel{Question}

Does proper-coloring Glauber dynamics mix in order N*log N whenever the color count exceeds Delta+\allowbreak{}2 on graphs of maximum degree Delta?


\entrylabel{Scope}

The book asks strictly more than Delta+\allowbreak{}2. Nearby formulations at Delta+\allowbreak{}2 are grouped. Later2026 results for high-girth graphs do not settle unrestricted graphs;\allowbreak{} vertex coloring is different from edge coloring.


\subsection*{P0201-A001 | Optimal Glauber mixing near the coloring degree threshold explicit structural calculation}

\entrylabel{Argument or result}

The elementary Hamming path coupling stops at q>2*Delta. For configurations differing at one vertex,\allowbreak{} updating that vertex removes the disagreement. Updating one of its at most Delta neighbors creates disagreement with probability at most1/\allowbreak{}(q-Delta),\allowbreak{} by maximally coupling the two uniform available-color sets. Therefore E[H\_\allowbreak{}next]<=\allowbreak{}1-1/\allowbreak{}N+\allowbreak{}Delta/\allowbreak{}[N*(q-Delta)],\allowbreak{} which contracts precisely when q>2*Delta.


\entrylabel{Second check}

The two available-color sets differ by at most one exchanged color,\allowbreak{} each has at least q-Delta elements,\allowbreak{} and their uniform laws have total variation at most1/\allowbreak{}(q-Delta). Iterating the path coupling yields an N*log N bound with constant proportional to(q-Delta)/\allowbreak{}(q-2*Delta).


\entrylabel{Remaining obstruction}

The denominator degenerates before the source threshold. Exploiting local color availability or multi-vertex moves is necessary;\allowbreak{} the failed bound is not a counterexample to rapid mixing.


Verification:\allowbreak{} author checked limited argument | Independent review:\allowbreak{} none | Full target solved:\allowbreak{} False


\clearpage

\section*{P0202  Diameter-squared upper bound for transitive mixing}

\entrylabel{Primary source}

Levin-Peres:\allowbreak{} Markov Chains and Mixing Times,\allowbreak{} second edition;\allowbreak{} Chapter26 Question8,\allowbreak{} PDFp377


\url{https://pages.uoregon.edu/dlevin/MARKOV/mcmt2e.pdf}


Source statement date:\allowbreak{} 2017. Archived access:\allowbreak{} 2026-09-06T21:\allowbreak{}43:\allowbreak{}30Z.


Historical:\allowbreak{} source\_\allowbreak{}says\_\allowbreak{}open | Current:\allowbreak{} unverified | Result:\allowbreak{} attempted\_\allowbreak{}unsolved


\entrylabel{Question}

Is lazy random-walk mixing on every finite transitive graph bounded by a universal constant times degree times diameter squared?


\entrylabel{Scope}

The2018 spectral-embedding paper still states this as an open question. Relaxation bounds of this order do not automatically prove the mixing bound.


\subsection*{P0202-A001 | Diameter-squared upper bound for transitive mixing explicit structural calculation}

\entrylabel{Argument or result}

The cycle illustrates the missing spectral information. On the N-cycle,\allowbreak{} lazy eigenvalues are1-sin\textasciicircum{}2(pi*k/\allowbreak{}N). Since sin(pi*k/\allowbreak{}N)>=\allowbreak{}2*k/\allowbreak{}N for k<=\allowbreak{}N/\allowbreak{}2,\allowbreak{} the chi-square bound from a point is at most2*sum\_\allowbreak{}(k>=\allowbreak{}1) exp(-8*t*k\textasciicircum{}2/\allowbreak{}N\textasciicircum{}2). At t=\allowbreak{}C*N\textasciicircum{}2 this sum is as small as desired for a universal C,\allowbreak{} proving the proposed diameter-squared order in this case without a log N penalty.


\entrylabel{Second check}

Use symmetry k versus N-k and(1-x)\textasciicircum{}(2t)<=\allowbreak{}exp(-2tx). The Gaussian series converges and tends to0 as C increases. The cycle diameter is floor(N/\allowbreak{}2) and its degree is2.


\entrylabel{Remaining obstruction}

For a general transitive graph,\allowbreak{} the gap alone bounds all N-1 modes by the same exponential and introduces log N. A uniform bound on low-frequency multiplicities or spectral mass strong enough to remove that factor is missing.


Verification:\allowbreak{} author checked limited argument | Independent review:\allowbreak{} none | Full target solved:\allowbreak{} False


\clearpage

\section*{P0203  Mixing exponent of uniform binary row additions}

\entrylabel{Primary source}

Levin-Peres:\allowbreak{} Markov Chains and Mixing Times,\allowbreak{} second edition;\allowbreak{} Chapter26 Question9,\allowbreak{} PDFp377;\allowbreak{} Peres current page Question2


\url{https://pages.uoregon.edu/dlevin/MARKOV/mcmt2e.pdf}


Source statement date:\allowbreak{} 2017. Archived access:\allowbreak{} 2026-09-06T21:\allowbreak{}43:\allowbreak{}30Z.


Historical:\allowbreak{} source\_\allowbreak{}says\_\allowbreak{}open | Current:\allowbreak{} unverified | Result:\allowbreak{} attempted\_\allowbreak{}unsolved


\entrylabel{Question}

Determine the mixing order for the walk on GL\_\allowbreak{}N(2) that adds one uniformly chosen distinct row to another.


\entrylabel{Scope}

The same question on the Peres author page is counted once. The book phrase saying order2\textasciicircum{}N nodes is a typo;\allowbreak{} the group size is of order2\textasciicircum{}(N\textasciicircum{}2). The asymptotic question concerns N tending to infinity.


\subsection*{P0203-A001 | Mixing exponent of uniform binary row additions explicit structural calculation}

\entrylabel{Argument or result}

The entropy/\allowbreak{}support lower bound has an exact leading constant. After t steps at most[N*(N-1)]\textasciicircum{}t states can be reached. Since |GL\_\allowbreak{}N(2)|=\allowbreak{}2\textasciicircum{}(N\textasciicircum{}2)*product\_\allowbreak{}(j=\allowbreak{}1)\textasciicircum{}N(1-2\textasciicircum{}-j),\allowbreak{} and the product tends to a positive constant,\allowbreak{} total variation is at least1-[N*(N-1)]\textasciicircum{}t/\allowbreak{}|GL\_\allowbreak{}N(2)|. Thus tmix(1/\allowbreak{}4)>=\allowbreak{}(N\textasciicircum{}2*log2-O(1))/\allowbreak{}log[N*(N-1)],\allowbreak{} asymptotically at least(log2/\allowbreak{}2)*N\textasciicircum{}2/\allowbreak{}log N.


\entrylabel{Second check}

The group-size formula chooses successive linearly independent columns:\allowbreak{} product\_\allowbreak{}(k=\allowbreak{}0)\textasciicircum{}(N-1)(2\textasciicircum{}N-2\textasciicircum{}k). The infinite product is positive because sum2\textasciicircum{}-j converges;\allowbreak{} log(1-x) is bounded below by-2x for0<=\allowbreak{}x<=\allowbreak{}1/\allowbreak{}2 after treating the first factor separately.


\entrylabel{Remaining obstruction}

Counting update words supplies no upper bound because many words cancel or represent the same matrix. The known O(N\textasciicircum{}3) upper bound leaves a logarithmic and polynomial gap.


Verification:\allowbreak{} author checked limited argument | Independent review:\allowbreak{} none | Full target solved:\allowbreak{} False


\clearpage

\section*{P0204  Mixing on large monotone subsets of the hypercube}

\entrylabel{Primary source}

Yuval Peres:\allowbreak{} related open Markov-chain problems;\allowbreak{} Related open problems,\allowbreak{} Question1


\url{https://www.yuval-peres-books.com/markov-chains-and-mixing-times/}


Source statement date:\allowbreak{} author\_\allowbreak{}page\_\allowbreak{}accessed2026-09-06. Archived access:\allowbreak{} 2026-09-06T21:\allowbreak{}45:\allowbreak{}37Z.


Historical:\allowbreak{} source\_\allowbreak{}says\_\allowbreak{}open | Current:\allowbreak{} unverified | Result:\allowbreak{} attempted\_\allowbreak{}unsolved


\entrylabel{Question}

For an increasing subset V of the binary N-cube with at least half its vertices,\allowbreak{} is lazy induced-subgraph random-walk mixing always O(N*log N)?


\entrylabel{Scope}

The author page gives an O(N\textasciicircum{}3) theorem of Ding-Mossel and conjectures N*log N. This is the induced graph walk,\allowbreak{} whose stationary law is degree weighted,\allowbreak{} not automatically uniform.


\subsection*{P0204-A001 | Mixing on large monotone subsets of the hypercube explicit structural calculation}

\entrylabel{Argument or result}

Every such induced graph is connected and has diameter at most2*N:\allowbreak{} from any x in V one may flip zero coordinates upwards to the all-one vertex while remaining in V. Every vertex has degree at least the number of its zero coordinates,\allowbreak{} but the degrees need not be constant. For V equal to the cube with only the all-zero vertex removed,\allowbreak{} a unit-vector vertex has degree N-1 while a vertex of Hamming weight at least2 has degree N (for N>=\allowbreak{}3). Thus uniform stationarity is already false in an almost-full monotone set.


\entrylabel{Second check}

For an undirected graph,\allowbreak{} detailed balance holds with pi(x)=\allowbreak{}deg(x)/\allowbreak{}sum\_\allowbreak{}y deg(y). On an edge between a unit vector and a weight-two vertex,\allowbreak{} a uniform law would have unequal transition flow1/\allowbreak{}[2*(N-1)] versus1/\allowbreak{}(2*N). The up-path connectivity argument is deterministic and uses only monotonicity.


\entrylabel{Remaining obstruction}

Applying full-cube product Fourier estimates without correcting the stationary law is invalid. Diameter and volume alone give much weaker general bounds;\allowbreak{} the monotone geometry must control bottlenecks quantitatively.


Verification:\allowbreak{} author checked limited argument | Independent review:\allowbreak{} none | Full target solved:\allowbreak{} False


\clearpage

\section*{P0205  Metropolis relaxation for geometric-time Cayley targets}

\entrylabel{Primary source}

Aldous open problems:\allowbreak{} Metropolis on Cayley graphs;\allowbreak{} Questions(i)-(iii),\allowbreak{} first posted March2009


\url{https://www.stat.berkeley.edu/~aldous/Research/OP/cayley.html}


Source statement date:\allowbreak{} 2009-03. Archived access:\allowbreak{} 2026-09-06T22:\allowbreak{}05:\allowbreak{}06Z.


Historical:\allowbreak{} source\_\allowbreak{}says\_\allowbreak{}open | Current:\allowbreak{} unverified | Result:\allowbreak{} partial\_\allowbreak{}counterexample\_\allowbreak{}author\_\allowbreak{}checked


\entrylabel{Question}

Control the relaxation time of Metropolis random walk targeting the distribution of a Cayley walk at an independent geometric time;\allowbreak{} in particular,\allowbreak{} test monotonic decrease in the geometric parameter.


\entrylabel{Scope}

The source defines p in[0,\allowbreak{}1] but writes tau(infinity) in its proposed bounds;\allowbreak{} that reference parameter is undefined and is retained as a source ambiguity. The monotonicity question itself is precise. Its negative answer below is a partial answer to the grouped request for general bounds.


\subsection*{P0205-A001 | Metropolis relaxation for geometric-time Cayley targets explicit structural calculation}

\entrylabel{Argument or result}

Monotone decrease is false already on the two-element Cayley group,\allowbreak{} with relaxation defined by1/\allowbreak{}(1-lambda\_\allowbreak{}2). For the nonlazy base walk that flips states,\allowbreak{} mu\_\allowbreak{}p=\allowbreak{}(1/\allowbreak{}(2-p),\allowbreak{}(1-p)/\allowbreak{}(2-p)),\allowbreak{} and its Metropolis matrix is[[p,\allowbreak{}1-p],\allowbreak{}[1,\allowbreak{}0]]. Its second eigenvalue is p-1,\allowbreak{} hence tau(p)=\allowbreak{}1/\allowbreak{}(2-p),\allowbreak{} which is strictly increasing. The same conclusion survives a lazy base walk:\allowbreak{} the base matrix has all entries1/\allowbreak{}2,\allowbreak{} mu\_\allowbreak{}p=\allowbreak{}((1+\allowbreak{}p)/\allowbreak{}2,\allowbreak{}(1-p)/\allowbreak{}2),\allowbreak{} the Metropolis off-diagonal rates are(1-p)/\allowbreak{}(2*(1+\allowbreak{}p)) and1/\allowbreak{}2,\allowbreak{} and tau(p)=\allowbreak{}1+\allowbreak{}p.


\entrylabel{Second check}

For the flip walk,\allowbreak{} sum the geometric probabilities over even versus odd times. For the lazy base,\allowbreak{} X\_\allowbreak{}0 is the identity and X\_\allowbreak{}k is uniform for every k>=\allowbreak{}1. The nontrivial eigenvalue of a two-state stochastic matrix is1 minus the sum of its off-diagonal probabilities,\allowbreak{} giving both formulas directly. If laziness is imposed after forming the nonlazy Metropolis kernel,\allowbreak{} tau becomes2/\allowbreak{}(2-p),\allowbreak{} still increasing.


\entrylabel{Remaining obstruction}

The literal monotonicity direction is refuted under the standard spectral-gap relaxation convention,\allowbreak{} but this does not answer the other universal-bound questions or resolve the source undefined tau(infinity). If relaxation meant inverse absolute gap for a nonlazy chain,\allowbreak{} that convention differs;\allowbreak{} the lazy example has nonnegative spectrum and removes that ambiguity. Prior-art status remains unverified.


Verification:\allowbreak{} author checked limited argument | Independent review:\allowbreak{} none | Full target solved:\allowbreak{} False


\clearpage

\section*{P0208  Extremal arc-transitive walk parameters away from cycles}

\entrylabel{Primary source}

Aldous-Fill:\allowbreak{} Reversible Markov Chains and Random Walks on Graphs,\allowbreak{} 2002 recompiled2014;\allowbreak{} Open Problem7.27,\allowbreak{} PDFp259;\allowbreak{} definitionp254


\url{https://www.stat.berkeley.edu/~aldous/RWG/book.pdf}


Source statement date:\allowbreak{} 1994;\allowbreak{}2002manuscript;\allowbreak{}2014recompilation. Archived access:\allowbreak{} 2026-09-06T22:\allowbreak{}03:\allowbreak{}09Z.


Historical:\allowbreak{} source\_\allowbreak{}says\_\allowbreak{}open | Current:\allowbreak{} unverified | Result:\allowbreak{} candidate\_\allowbreak{}full\_\allowbreak{}solution


\entrylabel{Question}

On arc-transitive n-vertex graphs other than n-cycles,\allowbreak{} are maximal commute time O(n log n),\allowbreak{} relaxation time O(n),\allowbreak{} and the maximal-commute to twice stationary-average-hitting ratio asymptotic to1?


\entrylabel{Scope}

The three linked bounds are one question. The source definition has a variable-assignment typo;\allowbreak{} the intended standard ordered-edge transitivity is used,\allowbreak{} as in the surrounding tori examples. Wreath graphs themselves are established prior art.


\entrylabel{Dated literature check}

Focused primary searches2026-09-06 found source question and related established results,\allowbreak{} but no verified later disposition. Absence from search is not proof of current openness or novelty.


\subsection*{P0208-A001 | Extremal arc-transitive walk parameters away from cycles explicit counterexample}

\entrylabel{Argument or result}

For odd m>=\allowbreak{}5 the established wreath graph C\_\allowbreak{}m[edgeless K2] is simple,\allowbreak{} connected,\allowbreak{}4-regular and arc-transitive on n=\allowbreak{}2m vertices. Its base projection is cycle walk. Exact hitting means are j*(m-j)+\allowbreak{}m between different fibers and2m between twins. Thus tau\_\allowbreak{}star=\allowbreak{}2*floor(m\textasciicircum{}2/\allowbreak{}4)+\allowbreak{}2m,\allowbreak{} tau0=\allowbreak{}(m\textasciicircum{}2-1)/\allowbreak{}6+\allowbreak{}m,\allowbreak{} and tau2=\allowbreak{}1/\allowbreak{}(1-cos(2*pi/\allowbreak{}m)). The first two scale quadratically in n,\allowbreak{} and tau\_\allowbreak{}star/\allowbreak{}(2*tau0) tends to3/\allowbreak{}2,\allowbreak{} refuting all three proposed bounds.


\entrylabel{Second check}

The invariant fiber-constant and fiber-antisymmetric subspaces give the full transition spectrum:\allowbreak{} cycle eigenvalues plus m zeros. Its fundamental trace agrees with the direct stationary hitting-time sum. Neighbor hitting equals2m-1=\allowbreak{}n-1 as required by the source lemma.


\entrylabel{Remaining obstruction}

Verify prior disposition of this1994 question:\allowbreak{} the graph construction is classical. The source typo and rate-one convention are explicit. No mathematical gap found by the author,\allowbreak{} but independent review and novelty remain pending.


Verification:\allowbreak{} author checked complete candidate | Independent review:\allowbreak{} none | Full target solved:\allowbreak{} False


\entrylabel{Full proof packet}

proofs/\allowbreak{}P0208\_\allowbreak{}wreath\_\allowbreak{}graph\_\allowbreak{}counterexample\_\allowbreak{}candidate.md


\clearpage

\section*{P0209  Uniform stationary normal approximation from gap and bounded observable}

\entrylabel{Primary source}

Aldous-Fill:\allowbreak{} Reversible Markov Chains and Random Walks on Graphs,\allowbreak{} 2002 recompiled2014;\allowbreak{} Open Problem4.30,\allowbreak{} PDFp137;\allowbreak{} continuous ratesp26


\url{https://www.stat.berkeley.edu/~aldous/RWG/book.pdf}


Source statement date:\allowbreak{} 1994;\allowbreak{}2002manuscript;\allowbreak{}2014recompilation. Archived access:\allowbreak{} 2026-09-06T22:\allowbreak{}03:\allowbreak{}09Z.


Historical:\allowbreak{} source\_\allowbreak{}says\_\allowbreak{}open | Current:\allowbreak{} unverified | Result:\allowbreak{} candidate\_\allowbreak{}full\_\allowbreak{}solution


\entrylabel{Question}

For stationary finite reversible continuous-time chains and Epi g=\allowbreak{}0,\allowbreak{}Epi g\textasciicircum{}2=\allowbreak{}1,\allowbreak{} does a universal normal-approximation error bound psi(||g||infinity,\allowbreak{}t/\allowbreak{}tau2) tend to0 at fixed first argument?


\entrylabel{Scope}

The source uses asymptotic variance normalization,\allowbreak{} permits arbitrary continuous transition rates,\allowbreak{} and starts in stationarity. The counterexample fixes ||g||infinity exactly2 and also survives uniformizing rates to at most1.


\entrylabel{Dated literature check}

Focused primary searches2026-09-06 found source question and related established results,\allowbreak{} but no verified later disposition. Absence from search is not proof of current openness or novelty.


\subsection*{P0209-A001 | Uniform stationary normal approximation from gap and bounded observable explicit counterexample}

\entrylabel{Argument or result}

Let epsilon=\allowbreak{}s\textasciicircum{}-2,\allowbreak{}R=\allowbreak{}s\textasciicircum{}4,\allowbreak{} Y have rates0to1=\allowbreak{}epsilon,\allowbreak{}1to0=\allowbreak{}1-epsilon,\allowbreak{} and independently let fair sign Z flip at R/\allowbreak{}2. Gap=\allowbreak{}1. Set A=\allowbreak{}2*epsilon+\allowbreak{}sqrt((1-epsilon)*(1-4*epsilon)),\allowbreak{}B=\allowbreak{}(2-A)/\allowbreak{}(1-epsilon),\allowbreak{}g=\allowbreak{}A*Z+\allowbreak{}B*(Y-epsilon). Then mean0,\allowbreak{}variance1,\allowbreak{}supnorm2 exactly and sigma\textasciicircum{}2=\allowbreak{}2*A\textasciicircum{}2/\allowbreak{}R+\allowbreak{}2*B\textasciicircum{}2*epsilon*(1-epsilon). By time s the chance ever Y=\allowbreak{}1 is at most epsilon*(s+\allowbreak{}1);\allowbreak{} off that event the slow integral is negligible after sigma*sqrt(s) normalization. The fast normalized variance is at most A\textasciicircum{}2/\allowbreak{}[R*B\textasciicircum{}2*epsilon*(1-epsilon)] tending0. Thus the normalized integral tends in probability to0,\allowbreak{} with normal Kolmogorov distance liminf at least1/\allowbreak{}2,\allowbreak{} although s/\allowbreak{}tau2 tends to infinity.


\entrylabel{Second check}

The product generator has eigenvalues0,\allowbreak{}-1,\allowbreak{}-R,\allowbreak{}-(R+\allowbreak{}1). Exact covariance integration verifies sigma and confirms finite-time variance divided by sigma\textasciicircum{}2*s tends to1 despite distributional collapse:\allowbreak{} rare events destroy uniform integrability. Scaling Q to Q/\allowbreak{}R and time to R*s leaves the standardized law and time/\allowbreak{}gap ratio unchanged,\allowbreak{} and bounds outgoing rates by1.


\entrylabel{Remaining obstruction}

No gap identified in the written quantifier and normalization checks. Current literature and independent review remain unverified;\allowbreak{} no original full-solution promotion.


Verification:\allowbreak{} author checked complete candidate | Independent review:\allowbreak{} none | Full target solved:\allowbreak{} False


\entrylabel{Full proof packet}

proofs/\allowbreak{}P0209\_\allowbreak{}uniform\_\allowbreak{}clt\_\allowbreak{}counterexample\_\allowbreak{}candidate.md


\clearpage

\section*{P0210  Martingales detectable by computationally practical stopping rules}

\entrylabel{Primary source}

Aldous open problems:\allowbreak{} Martingale for practical purposes;\allowbreak{} Problem and polynomial-constraint formulation


\url{https://www.stat.berkeley.edu/~aldous/Research/OP/fields.html}


Source statement date:\allowbreak{} undated author page;\allowbreak{}accessed2026-09-06. Archived access:\allowbreak{} 2026-09-06T22:\allowbreak{}04:\allowbreak{}12Z.


Historical:\allowbreak{} source\_\allowbreak{}says\_\allowbreak{}open | Current:\allowbreak{} unverified | Result:\allowbreak{} attempted\_\allowbreak{}unsolved


\entrylabel{Question}

Give a useful probabilistic definition of martingale for practical purposes,\allowbreak{} or a tractable class defined by polynomially many stopping-time constraints.


\entrylabel{Scope}

The two formulations are grouped as one conceptual question. A complexity model and input representation are not specified by the source;\allowbreak{} no financial claim or backtest is involved.


\entrylabel{Dated literature check}

Author-maintained page and relevant linked primary sources checked2026-09-06. Focused later-literature searches did not establish a complete resolution;\allowbreak{} current status remains unverified.


\subsection*{P0210-A001 | Martingales detectable by computationally practical stopping rules structural calculation}

\entrylabel{Argument or result}

Any prescribed polynomial-size family of stopping tests can miss bounded conditional drift on an exponentially large history space. Reveal m fair bits by time m,\allowbreak{} set X\_\allowbreak{}t=\allowbreak{}0 through time m and X\_\allowbreak{}(m+\allowbreak{}1)=\allowbreak{}f(Y). For predetermined stopping rules T\_\allowbreak{}j,\allowbreak{} each optional-mean constraint is the linear equation E[f(Y)*1\_\allowbreak{}\{T\_\allowbreak{}j=\allowbreak{}m+\allowbreak{}1\}]=\allowbreak{}0. If the number of tests plus the mean-zero condition is less than2\textasciicircum{}m,\allowbreak{} a nonzero f exists in their common nullspace. Rescale it to supnorm1. Then the process passes every supplied test but E[X\_\allowbreak{}(m+\allowbreak{}1)|F\_\allowbreak{}m]=\allowbreak{}f(Y) is nonzero.


\entrylabel{Second check}

The witness T*=\allowbreak{}m+\allowbreak{}1 on\{f(Y)>0\} and m otherwise is a stopping time and has E X\_\allowbreak{}T*=\allowbreak{}E f\_\allowbreak{}+\allowbreak{}>0. Predetermined test rules are functions of the revealed-bit history;\allowbreak{} their final-time continuation indicators do not depend on the unknown terminal value. The construction does not claim f or T* is computationally efficient.


\entrylabel{Remaining obstruction}

A polynomial number of linear tests alone cannot reproduce the full martingale property. Selecting a useful computational subclass requires an explicit representation,\allowbreak{} complexity bound and error tolerance;\allowbreak{} the nullspace construction does not supply that theory.


Verification:\allowbreak{} author checked limited argument | Independent review:\allowbreak{} none | Full target solved:\allowbreak{} False


\clearpage

\section*{P0211  Relaxation of linear-polymer coagulation and fragmentation}

\entrylabel{Primary source}

Aldous open problems:\allowbreak{} Coagulation-fragmentation mixing;\allowbreak{} Problem;\allowbreak{} linked coag\_\allowbreak{}frag\_\allowbreak{}mix.pdf pp1-4


\url{https://www.stat.berkeley.edu/~aldous/Research/OP/coag_frag.html}


Source statement date:\allowbreak{} 2012-08;\allowbreak{}updated2018-02. Archived access:\allowbreak{} 2026-09-06T22:\allowbreak{}04:\allowbreak{}42Z.


Historical:\allowbreak{} source\_\allowbreak{}says\_\allowbreak{}open | Current:\allowbreak{} unverified | Result:\allowbreak{} attempted\_\allowbreak{}unsolved


\entrylabel{Question}

Determine the order of the relaxation time for the reversible polymer coagulation-fragmentation chain,\allowbreak{} especially fragmentation rate beta\_\allowbreak{}n comparable to1/\allowbreak{}n.


\entrylabel{Scope}

Clusters merge at rate2/\allowbreak{}n per pair;\allowbreak{} each internal link splits at rate beta. Fixed-beta and small-beta regimes are grouped. The source proves bounds between asymptotic1/\allowbreak{}sqrt(beta\textasciicircum{}2+\allowbreak{}4beta) and1/\allowbreak{}beta.


\entrylabel{Dated literature check}

Author-maintained page and relevant linked primary sources checked2026-09-06. Focused later-literature searches did not establish a complete resolution;\allowbreak{} current status remains unverified.


\subsection*{P0211-A001 | Relaxation of linear-polymer coagulation and fragmentation structural calculation}

\entrylabel{Argument or result}

The cluster-count coordinate K has exact rates K->K+\allowbreak{}1 at beta*(n-K) and K->K-1 at K*(K-1)/\allowbreak{}n. Its deterministic equilibrium density mu solves mu\textasciicircum{}2=\allowbreak{}beta*(1-mu),\allowbreak{} and linearizing the density drift gives restoring rate beta+\allowbreak{}2*mu=\allowbreak{}sqrt(beta\textasciicircum{}2+\allowbreak{}4*beta). At beta=\allowbreak{}b/\allowbreak{}n this is order n\textasciicircum{}-1/\allowbreak{}2. The source edge coupling has discrepancy killing rate at least beta+\allowbreak{}(K-1)/\allowbreak{}n,\allowbreak{} so it also suggests this scale once K is near sqrt(b*n).


\entrylabel{Second check}

For n=\allowbreak{}3 the partition states3,\allowbreak{}21,\allowbreak{}111 have rates3->21=\allowbreak{}2beta,\allowbreak{}21->3=\allowbreak{}2/\allowbreak{}3,\allowbreak{}21->111=\allowbreak{}beta,\allowbreak{}111->21=\allowbreak{}2. Direct characteristic-polynomial calculation gives gap[3beta+\allowbreak{}8/\allowbreak{}3-sqrt(beta\textasciicircum{}2+\allowbreak{}16/\allowbreak{}9)]/\allowbreak{}2. Thus finite-n small-beta behavior cannot simply be substituted for the joint n,\allowbreak{}beta\_\allowbreak{}n limit.


\entrylabel{Remaining obstruction}

The required upgrade is a uniform exponential bound on the coupling survival functional E exp[-integral(beta+\allowbreak{}(K\_\allowbreak{}u-1)/\allowbreak{}n)du],\allowbreak{} including atypically small K. Replacing K by its stationary mean inside the exponential has the wrong Jensen direction and is not a proof of the desired upper bound.


Verification:\allowbreak{} author checked limited argument | Independent review:\allowbreak{} none | Full target solved:\allowbreak{} False


\clearpage

\section*{P0212  Low-density genealogical limit of branching coalescing walks}

\entrylabel{Primary source}

Aldous open problems:\allowbreak{} Low-density coalescing branching random walk;\allowbreak{} Problem A,\allowbreak{}2012;\allowbreak{} updateMarch2018


\url{https://www.stat.berkeley.edu/~aldous/Research/OP/cbrw.html}


Source statement date:\allowbreak{} 2012;\allowbreak{}updated2018-03. Archived access:\allowbreak{} 2026-09-06T22:\allowbreak{}04:\allowbreak{}49Z.


Historical:\allowbreak{} source\_\allowbreak{}says\_\allowbreak{}open | Current:\allowbreak{} unverified | Result:\allowbreak{} attempted\_\allowbreak{}unsolved


\entrylabel{Question}

Prove the proposed low-density,\allowbreak{} time-rescaled lineage limit for transient branching-coalescing walk and the constrained-Ising analogue,\allowbreak{} after discarding short-lived branching jitter.


\entrylabel{Scope}

The two explicitly linked microscopic models are grouped;\allowbreak{} the undefined dynamic-graph research direction is not a second admitted problem. Spatial positions are forgotten but splitting and merging histories remain.


\entrylabel{Dated literature check}

Author-maintained page and relevant linked primary sources checked2026-09-06. Focused later-literature searches did not establish a complete resolution;\allowbreak{} current status remains unverified.


\subsection*{P0212-A001 | Low-density genealogical limit of branching coalescing walks structural calculation}

\entrylabel{Argument or result}

Immediate parent-daughter returns renormalize the branching clock. For two independent simple walks started at neighboring sites in Z\textasciicircum{}3,\allowbreak{} their relative walk is a time change of simple walk. The probability they never meet is q=\allowbreak{}1-G(e,\allowbreak{}0)/\allowbreak{}G(0,\allowbreak{}0)=\allowbreak{}1/\allowbreak{}G(0,\allowbreak{}0),\allowbreak{} using G(0,\allowbreak{}0)=\allowbreak{}1+\allowbreak{}G(e,\allowbreak{}0) by the first-step identity and lattice symmetry. Thus a rare branch at rate beta produces a macroscopically surviving split at leading rate beta*q,\allowbreak{} rather than beta.


\entrylabel{Second check}

The equality of the two Green-function expressions follows from translation and reflection symmetry of all six neighbors. It concerns an isolated pair;\allowbreak{} the chance of interference by other particles has not been set to zero without justification.


\entrylabel{Remaining obstruction}

A thinning heuristic suggests time scaling1/\allowbreak{}(beta*q),\allowbreak{} but a proof needs control of third-particle encounters,\allowbreak{} dependent return episodes and joint Palm genealogies. Stationary balance of split and merge counts does not imply independent Poisson clocks.


Verification:\allowbreak{} author checked limited argument | Independent review:\allowbreak{} none | Full target solved:\allowbreak{} False


\clearpage

\section*{P0213  Constant-factor online scheduling with subadditive batch durations}

\entrylabel{Primary source}

Aldous open problems:\allowbreak{} Bacon and eggs batch scheduling;\allowbreak{} Problem,\allowbreak{}postedMarch2009


\url{https://www.stat.berkeley.edu/~aldous/Research/OP/bacon.html}


Source statement date:\allowbreak{} 2009-03. Archived access:\allowbreak{} 2026-09-06T22:\allowbreak{}04:\allowbreak{}58Z.


Historical:\allowbreak{} source\_\allowbreak{}says\_\allowbreak{}open | Current:\allowbreak{} unverified | Result:\allowbreak{} attempted\_\allowbreak{}unsolved


\entrylabel{Question}

Find an explicit online batch-scheduling algorithm with a universal ratio of long-run weighted waiting cost to the optimal online cost,\allowbreak{} over monotone strictly subadditive service-time functions.


\entrylabel{Scope}

Arrivals form a rate1 Poisson process with uniform types. The source assumes finite optimal cost. A finite deterministic queue example below is only an obstruction to a proposed greedy step,\allowbreak{} not a steady-state counterexample.


\entrylabel{Dated literature check}

Author-maintained page and relevant linked primary sources checked2026-09-06. Focused later-literature searches did not establish a complete resolution;\allowbreak{} current status remains unverified.


\subsection*{P0213-A001 | Constant-factor online scheduling with subadditive batch durations structural calculation}

\entrylabel{Argument or result}

Greedy batching can be arbitrarily poor for a fixed two-job queue when holding costs differ. Let S(a)=\allowbreak{}1,\allowbreak{}S(b)=\allowbreak{}L,\allowbreak{}S(\{a,\allowbreak{}b\})=\allowbreak{}L+\allowbreak{}epsilon with0<epsilon<1 and L>1. Give a holding rate M and b rate1. Processing both immediately costs(M+\allowbreak{}1)*(L+\allowbreak{}epsilon). Processing a first and then b costs M+\allowbreak{}1+\allowbreak{}L. For M much larger than L their ratio approaches L+\allowbreak{}epsilon,\allowbreak{} although service times are strictly monotone and subadditive.


\entrylabel{Second check}

The individual completion times are L+\allowbreak{}epsilon for both in the batch,\allowbreak{} and1,\allowbreak{}L+\allowbreak{}1 in the sequential schedule. This compares exactly the same initial jobs without incoming work. For the actual stationary model,\allowbreak{} Campbell reasoning also gives the elementary lower bound W\_\allowbreak{}*>=\allowbreak{}integral c(a)*S(\{a\}) da,\allowbreak{} since each task must spend at least its singleton service time in its completing batch.


\entrylabel{Remaining obstruction}

Embedding the bad queue event in a stable Poisson-arrival system while controlling its frequency and other waiting costs remains necessary. The two-job calculation does not disprove existence of a different constant-factor algorithm.


Verification:\allowbreak{} author checked limited argument | Independent review:\allowbreak{} none | Full target solved:\allowbreak{} False


\clearpage

\section*{P0214  Fluctuations of shortest routes in Poisson proximity networks}

\entrylabel{Primary source}

Aldous open problems:\allowbreak{} Shortest routes in random proximity networks;\allowbreak{} Main question;\allowbreak{}2019 simulation study discussed


\url{https://www.stat.berkeley.edu/~aldous/Research/OP/random_proximity.html}


Source statement date:\allowbreak{} 2019 study discussed;\allowbreak{}page accessed2026-09-06. Archived access:\allowbreak{} 2026-09-06T22:\allowbreak{}05:\allowbreak{}16Z.


Historical:\allowbreak{} source\_\allowbreak{}says\_\allowbreak{}open | Current:\allowbreak{} unverified | Result:\allowbreak{} attempted\_\allowbreak{}unsolved


\entrylabel{Question}

Determine rigorous orders of route-length variance and transversal deviation for shortest paths in random proximity networks on planar Poisson points.


\entrylabel{Scope}

Variance and transversal exponents are grouped as the source fluctuation family. A law of large numbers is established and is not counted as a solution.


\entrylabel{Dated literature check}

Author-maintained page and relevant linked primary sources checked2026-09-06. Focused later-literature searches did not establish a complete resolution;\allowbreak{} current status remains unverified.


\subsection*{P0214-A001 | Fluctuations of shortest routes in Poisson proximity networks structural calculation}

\entrylabel{Argument or result}

A deterministic route of length D connecting points at distance r and reaching perpendicular deviation T must satisfy D>=\allowbreak{}sqrt(r\textasciicircum{}2+\allowbreak{}4*T\textasciicircum{}2). To see this,\allowbreak{} split at a point of maximal deviation and minimize the sum of its distances to the two endpoints over the longitudinal coordinate;\allowbreak{} the minimum occurs above the midpoint. Hence T<=\allowbreak{}sqrt(D\textasciicircum{}2-r\textasciicircum{}2)/\allowbreak{}2.


\entrylabel{Second check}

When D/\allowbreak{}r tends to a stretch constant mu>1,\allowbreak{} the bound gives only T/\allowbreak{}r<=\allowbreak{}sqrt(mu\textasciicircum{}2-1)/\allowbreak{}2 asymptotically. It does not force sublinear wandering. This explicitly explains why the existing route-length shape theorem does not already identify a transversal exponent.


\entrylabel{Remaining obstruction}

A KPZ-type prediction needs estimates for centered route excess and correlations of competing near-geodesics. The deterministic Euclidean inequality leaves an order-r allowance whenever the limiting stretch exceeds1.


Verification:\allowbreak{} author checked limited argument | Independent review:\allowbreak{} none | Full target solved:\allowbreak{} False


\clearpage

\section*{P0215  Growth exponent for common subcladograms of random trees}

\entrylabel{Primary source}

Aldous open problems:\allowbreak{} Largest common random substructures;\allowbreak{} Original2003 Section2 Example3 PDFp5;\allowbreak{} author updates


\url{https://www.stat.berkeley.edu/~aldous/Research/OP/substructures.html}


Source statement date:\allowbreak{} 2003-05;\allowbreak{}updated2022. Archived access:\allowbreak{} 2026-09-06T22:\allowbreak{}05:\allowbreak{}42Z.


Historical:\allowbreak{} source\_\allowbreak{}says\_\allowbreak{}open | Current:\allowbreak{} unverified | Result:\allowbreak{} attempted\_\allowbreak{}unsolved


\entrylabel{Question}

Establish existence and values of the power-law growth exponents for the largest common induced leaf-labeled subtree of two independent uniform cladograms or coalescent trees.


\entrylabel{Scope}

The two source models form one admitted family. The2024 Budzinski-Senizergues paper improves the uniform-model upper exponent below1/\allowbreak{}2;\allowbreak{} that established bound is not a new solution or proof of a limiting exponent.


\entrylabel{Dated literature check}

Author-maintained page and relevant linked primary sources checked2026-09-06. Focused later-literature searches did not establish a complete resolution;\allowbreak{} current status remains unverified.


\subsection*{P0215-A001 | Growth exponent for common subcladograms of random trees structural calculation}

\entrylabel{Argument or result}

For the uniform model,\allowbreak{} sampling consistency gives an exact first-moment count:\allowbreak{} the expected number of agreeing k-subsets is binom(n,\allowbreak{}k)/\allowbreak{}(2k-3)!!. Thus P(C\_\allowbreak{}n>=\allowbreak{}k) is at most this ratio. Stirling gives log of the ratio=\allowbreak{}k*[log n-2log k+\allowbreak{}2-log2]+\allowbreak{}o(k) when k is of order sqrt n. Consequently the elementary bound stops at order sqrt n,\allowbreak{} with threshold constant e/\allowbreak{}sqrt2.


\entrylabel{Second check}

The agreement probability is the collision probability of two uniform induced k-cladograms,\allowbreak{} exactly1/\allowbreak{}(2k-3)!!. Existence of a common larger subset implies existence of one of size k. The coalescent model has a nonuniform induced law,\allowbreak{} so its collision probability must instead be sum\_\allowbreak{}T p\_\allowbreak{}T\textasciicircum{}2;\allowbreak{} the uniform formula is not transferred to it.


\entrylabel{Remaining obstruction}

First moments overcount many strongly overlapping agreeing subsets. They cannot reach the stronger known exponent bound or prove convergence of log E C\_\allowbreak{}n/\allowbreak{}log n. The source centroid recursion constructs particular matches and is not automatically optimal.


Verification:\allowbreak{} author checked limited argument | Independent review:\allowbreak{} none | Full target solved:\allowbreak{} False


\clearpage

\section*{P0216  Limit constant for common random two-dimensional partial orders}

\entrylabel{Primary source}

Aldous open problems:\allowbreak{} Largest common random substructures;\allowbreak{} Original2003 Section2 Example4 PDFpp5-6


\url{https://www.stat.berkeley.edu/~aldous/Research/OP/substructures.html}


Source statement date:\allowbreak{} 2003-05. Archived access:\allowbreak{} 2026-09-06T22:\allowbreak{}05:\allowbreak{}42Z.


Historical:\allowbreak{} source\_\allowbreak{}says\_\allowbreak{}open | Current:\allowbreak{} unverified | Result:\allowbreak{} attempted\_\allowbreak{}unsolved


\entrylabel{Question}

Prove that the mean size of the largest common induced suborder of two independent uniform two-dimensional partial orders is asymptotic to c*n\textasciicircum{}(1/\allowbreak{}3),\allowbreak{} and determine c.


\entrylabel{Scope}

The source already gives order n\textasciicircum{}(1/\allowbreak{}3) upper and lower arguments. The calculation here makes a constant in its grid lower construction explicit;\allowbreak{} it is a rederivation of an established method,\allowbreak{} not resolution of the limit question.


\entrylabel{Dated literature check}

Author-maintained page and relevant linked primary sources checked2026-09-06. Focused later-literature searches did not establish a complete resolution;\allowbreak{} current status remains unverified.


\subsection*{P0216-A001 | Limit constant for common random two-dimensional partial orders structural calculation}

\entrylabel{Argument or result}

Partition each square into m-by-m cells,\allowbreak{} m=\allowbreak{}floor(n\textasciicircum{}(1/\allowbreak{}3)). Keep a label when its two points occupy the same cell,\allowbreak{} and retain it only if no other kept label shares that row or column. The retained labels have identical comparisons in both posets. For a fixed label the match probability is1/\allowbreak{}m\textasciicircum{}2;\allowbreak{} conditional on its cell,\allowbreak{} another label is a conflicting match with probability(2m-1)/\allowbreak{}m\textasciicircum{}4. Therefore the expected retained number is n/\allowbreak{}m\textasciicircum{}2*[1-(2m-1)/\allowbreak{}m\textasciicircum{}4]\textasciicircum{}(n-1),\allowbreak{} asymptotic to exp(-2)*n\textasciicircum{}(1/\allowbreak{}3).


\entrylabel{Second check}

Two retained labels have distinct row and column indices,\allowbreak{} so each of their coordinate orderings is determined by cell indices in both independent samples. Independence over labels gives the exact no-conflict probability. The bound is liminf E C\_\allowbreak{}n/\allowbreak{}n\textasciicircum{}(1/\allowbreak{}3)>=\allowbreak{}exp(-2).


\entrylabel{Remaining obstruction}

The isolated-cell matching discards many usable labels and supplies no matching upper constant. Optimizing a local grid construction does not establish existence of the global optimal limit.


Verification:\allowbreak{} author checked limited argument | Independent review:\allowbreak{} none | Full target solved:\allowbreak{} False


\clearpage

\section*{P0217  Percolation of merging planar empires}

\entrylabel{Primary source}

Aldous open problems:\allowbreak{} Percolation and empires;\allowbreak{} Two PROBLEM paragraphs;\allowbreak{}2009 and2017 updates


\url{https://www.stat.berkeley.edu/~aldous/Research/OP/empires.html}


Source statement date:\allowbreak{} 2001 discussion;\allowbreak{}2009-10 writeup;\allowbreak{}2017-06 update. Archived access:\allowbreak{} 2026-09-06T22:\allowbreak{}05:\allowbreak{}50Z.


Historical:\allowbreak{} source\_\allowbreak{}says\_\allowbreak{}open | Current:\allowbreak{} unverified | Result:\allowbreak{} attempted\_\allowbreak{}unsolved


\entrylabel{Question}

Characterize merging-rate conditions for finite-time infinite empires,\allowbreak{} in particular whether rate1 per adjacent pair percolates.


\entrylabel{Scope}

General rate criteria and the distinguished constant-rate case are grouped. The source already proves some sufficient conditions in its2010 paper;\allowbreak{} the main constant-rate case is not inferred from bond percolation.


\entrylabel{Dated literature check}

Author-maintained page and relevant linked primary sources checked2026-09-06. Focused later-literature searches did not establish a complete resolution;\allowbreak{} current status remains unverified.


\subsection*{P0217-A001 | Percolation of merging planar empires structural calculation}

\entrylabel{Argument or result}

The natural area-moment Lyapunov calculation works for strongly suppressed rates. In a finite periodic partition,\allowbreak{} F=\allowbreak{}sum\_\allowbreak{}A area(A)\textasciicircum{}2 changes by2ab when areas a,\allowbreak{}b merge. If r(A,\allowbreak{}B)<=\allowbreak{}1/\allowbreak{}N(A,\allowbreak{}B),\allowbreak{} where N is at least each endpoint region adjacency degree,\allowbreak{} then the generator obeys L F<=\allowbreak{}sum\_\allowbreak{}edges(a\textasciicircum{}2+\allowbreak{}b\textasciicircum{}2)/\allowbreak{}N(A,\allowbreak{}B)<=\allowbreak{}F. Hence E F\_\allowbreak{}t<=\allowbreak{}e\textasciicircum{}t F\_\allowbreak{}0. Under an infinite-volume limit preserving this bound and translation invariance,\allowbreak{} the mean area at a typical point stays finite,\allowbreak{} excluding an infinite empire.


\entrylabel{Second check}

Use2ab<=\allowbreak{}a\textasciicircum{}2+\allowbreak{}b\textasciicircum{}2 and sum over neighbors of each region of1/\allowbreak{}N<=\allowbreak{}1. This independently explains the source suppressed-rate example. For rate1 the neighbor sum becomes the region degree,\allowbreak{} which has no such uniform control.


\entrylabel{Remaining obstruction}

For rate1,\allowbreak{} knowing each original boundary survives with probability exp(-t) does not give product bounds for contours:\allowbreak{} multiple segments separating the same two current empires share one merger clock. The missing joint survival estimate prevents the naive Peierls argument.


Verification:\allowbreak{} author checked limited argument | Independent review:\allowbreak{} none | Full target solved:\allowbreak{} False


\clearpage

\section*{P0218  Growth exponents in the spatial city-formation model}

\entrylabel{Primary source}

Aldous:\allowbreak{} A spatial model of city growth and formation;\allowbreak{} Section3 equations4-9 and Section3.1;\allowbreak{} author open page


\url{https://arxiv.org/pdf/1209.5120}


Source statement date:\allowbreak{} 2012-09-23. Archived access:\allowbreak{} 2026-09-06T22:\allowbreak{}16:\allowbreak{}15Z.


Historical:\allowbreak{} source\_\allowbreak{}says\_\allowbreak{}open | Current:\allowbreak{} unverified | Result:\allowbreak{} attempted\_\allowbreak{}unsolved


\entrylabel{Question}

Prove the proposed balanced/\allowbreak{}unbalanced growth regimes and city-count exponents in the Aldous-Huang spatial immigration model.


\entrylabel{Scope}

All parameter regimes,\allowbreak{} including alpha1,\allowbreak{} are grouped. A city of size N has influence c0*N\textasciicircum{}alpha*r\textasciicircum{}-beta;\allowbreak{} an arrival founds a city with probability1/\allowbreak{}(1+\allowbreak{}maximum influence). The source heuristic is not a theorem.


\entrylabel{Dated literature check}

Author-maintained page and relevant linked primary sources checked2026-09-06. Focused later-literature searches did not establish a complete resolution;\allowbreak{} current status remains unverified.


\subsection*{P0218-A001 | Growth exponents in the spatial city-formation model structural calculation}

\entrylabel{Argument or result}

A rigorous one-sided bound supports the balanced exponent when0<alpha<1 and beta>2alpha. Write d=\allowbreak{}beta/\allowbreak{}2-alpha>0. A union bound over influence disks gives area\{I(y,\allowbreak{}t)>u\}<=\allowbreak{}pi*c0\textasciicircum{}(2/\allowbreak{}beta)*u\textasciicircum{}(-2/\allowbreak{}beta)*M\_\allowbreak{}t\textasciicircum{}(1-2alpha/\allowbreak{}beta)*t\textasciicircum{}(2alpha/\allowbreak{}beta),\allowbreak{} using concavity of the population power. Thus the conditional founding probability is at least1/\allowbreak{}[2*(1+\allowbreak{}K*t\textasciicircum{}alpha*M\_\allowbreak{}t\textasciicircum{}d)],\allowbreak{} K=\allowbreak{}(2pi)\textasciicircum{}(beta/\allowbreak{}2)*c0. Put theta=\allowbreak{}(1-alpha)/\allowbreak{}(1+\allowbreak{}d) and m=\allowbreak{}floor(kappa*t\textasciicircum{}theta). Until M reaches m,\allowbreak{} each of the last floor(t/\allowbreak{}2) arrivals has founding chance at least p=\allowbreak{}1/\allowbreak{}[2*(1+\allowbreak{}K*t\textasciicircum{}alpha*m\textasciicircum{}d)]. Choose kappa small so this Bernoulli count has mean at least2m. A stopped coupling and Chernoff then give P(M\_\allowbreak{}t<m)<=\allowbreak{}exp(-c*m) for large t.


\entrylabel{Second check}

The disk estimate is valid even at the square boundary because clipping only reduces area. Sum\_\allowbreak{}i N\_\allowbreak{}i\textasciicircum{}(2alpha/\allowbreak{}beta)<=\allowbreak{}M\textasciicircum{}(1-2alpha/\allowbreak{}beta)*t\textasciicircum{}(2alpha/\allowbreak{}beta). On the event M\_\allowbreak{}t<m the threshold has never been reached,\allowbreak{} so domination by independent Bernoulli(p) increments throughout the half-interval is legitimate. The exponent identity theta*(1+\allowbreak{}d)=\allowbreak{}1-alpha makes the mean comparable to m.


\entrylabel{Remaining obstruction}

This proves a lower city-count growth bound in the balanced regime,\allowbreak{} not the matching upper bound,\allowbreak{} comparable populations,\allowbreak{} unbalanced regime,\allowbreak{} or alpha1 logarithmic law. The upper bound requires control of spatial clustering and influence overlap that the union bound does not supply.


Verification:\allowbreak{} author checked limited argument | Independent review:\allowbreak{} none | Full target solved:\allowbreak{} False


\clearpage

\section*{P0219  High-density oriented percolation flow asymptotics}

\entrylabel{Primary source}

Aldous open problems:\allowbreak{} Discrete Hammersley process and oriented percolation flows;\allowbreak{} Conjecture,\allowbreak{}formulated2004


\url{https://www.stat.berkeley.edu/~aldous/Research/OP/hammersley_flow.html}


Source statement date:\allowbreak{} 2004 formulation;\allowbreak{}undated author page accessed2026-09-06. Archived access:\allowbreak{} 2026-09-06T22:\allowbreak{}10:\allowbreak{}49Z.


Historical:\allowbreak{} source\_\allowbreak{}says\_\allowbreak{}open | Current:\allowbreak{} unverified | Result:\allowbreak{} attempted\_\allowbreak{}unsolved


\entrylabel{Question}

For the oriented square-lattice maximal-flow constant v(p),\allowbreak{} prove1-v(p) is asymptotic to sqrt(2*(1-p)) as p tends to1.


\entrylabel{Scope}

The p-down-to-critical regime is not counted as an invented second problem. The author heuristic contains two normalization errors,\allowbreak{} recorded explicitly below;\allowbreak{} correcting them still supports the displayed conjecture.


\entrylabel{Dated literature check}

Author-maintained page and relevant linked primary sources checked2026-09-06. Focused later-literature searches did not establish a complete resolution;\allowbreak{} current status remains unverified.


\subsection*{P0219-A001 | High-density oriented percolation flow asymptotics structural calculation}

\entrylabel{Argument or result}

For a continuum Poisson process of intensity delta in[0,\allowbreak{}n]\textasciicircum{}2,\allowbreak{} rescaling coordinates by sqrt(delta) gives intensity1 in[0,\allowbreak{}n*sqrt(delta)]\textasciicircum{}2. Hence its minimum increasing-chain cover has the law U*(n*sqrt(delta),\allowbreak{}1),\allowbreak{} not U*(n/\allowbreak{}sqrt(delta),\allowbreak{}1). By the poset chain-cover/\allowbreak{}antichain relation it equals the longest decreasing subsequence,\allowbreak{} whose leading Poissonized Hammersley asymptotic is2m for a unit-intensity m-square. With delta=\allowbreak{}2*(1-p),\allowbreak{} dividing2*n*sqrt(delta) by2n recovers sqrt(2*(1-p)).


\entrylabel{Second check}

The expected number of Poisson points,\allowbreak{} delta*n\textasciicircum{}2,\allowbreak{} is invariant under the corrected scaling and would fail under the printed inverse scaling. Reflection converts decreasing to increasing subsequences,\allowbreak{} preserving the known leading constant2. These are checks of the continuum surrogate only.


\entrylabel{Remaining obstruction}

One still needs a uniform comparison of the lattice closed-edge path-cover problem with the continuum antichain problem as n tends to infinity and p tends to1. The two corrected constants do not justify interchanging those limits.


Verification:\allowbreak{} author checked limited argument | Independent review:\allowbreak{} none | Full target solved:\allowbreak{} False


\clearpage

\section*{P0220  Explicit stationary law of the leftmost Hammersley drift-jump particles}

\entrylabel{Primary source}

Aldous open problems:\allowbreak{} One-dimensional drift-jump particle process;\allowbreak{} Problem,\allowbreak{}postedFebruary2018;\allowbreak{}1993notes Section4.1


\url{https://www.stat.berkeley.edu/~aldous/Research/OP/left_Hamm.html}


Source statement date:\allowbreak{} 2018-02. Archived access:\allowbreak{} 2026-09-06T22:\allowbreak{}10:\allowbreak{}42Z.


Historical:\allowbreak{} source\_\allowbreak{}says\_\allowbreak{}open | Current:\allowbreak{} unverified | Result:\allowbreak{} attempted\_\allowbreak{}unsolved


\entrylabel{Question}

Give a useful explicit joint stationary law for the ordered particles which drift at speed x and jump left at rate1 Poisson space-time points.


\entrylabel{Scope}

The1993 notes already contain a patience-sorting representation and small-k calculations. The1999 paper Section5.2 uses the stationary process for limit theorems. Those established representations are not counted as newly solved open work.


\entrylabel{Dated literature check}

Author-maintained page and relevant linked primary sources checked2026-09-06. Focused later-literature searches did not establish a complete resolution;\allowbreak{} current status remains unverified.


\subsection*{P0220-A001 | Explicit stationary law of the leftmost Hammersley drift-jump particles structural calculation}

\entrylabel{Argument or result}

The established rank representation can be expanded into a convergent combinatorial density series. Let B=\allowbreak{}(1=\allowbreak{}b1<...<bk) be the limiting first-row labels of the insertion tableau,\allowbreak{} and let q\_\allowbreak{}B be its probability. The first b=\allowbreak{}bk labels form a uniform permutation,\allowbreak{} so q\_\allowbreak{}B=\allowbreak{}(1/\allowbreak{}b!)*sum over partitions lambda of b with first row length k of f\textasciicircum{}lambda times the number of standard tableaux of shape lambda with first row B. Independent unit-rate Poisson spacings then give density exp(-xk)*sum\_\allowbreak{}B q\_\allowbreak{}B*product\_\allowbreak{}(i=\allowbreak{}2..k)[(x\_\allowbreak{}i-x\_\allowbreak{}(i-1))\textasciicircum{}(b\_\allowbreak{}i-b\_\allowbreak{}(i-1)-1)/\allowbreak{}(b\_\allowbreak{}i-b\_\allowbreak{}(i-1)-1)!].


\entrylabel{Second check}

Restriction of an insertion tableau to labels at most b is the tableau of the corresponding restricted subword. The last specified first-row entry b pins the first k labels,\allowbreak{} irrespective of later larger labels. Conditional on B,\allowbreak{} the Poisson arrival times have independent gamma increments,\allowbreak{} so every term is a normalized conditional density and nonnegative summation is valid. For k1 this gives Exp(1);\allowbreak{} known Poissonized LIS formulas give P(X2>x)=\allowbreak{}exp(-x)*I0(2sqrt(x)) and E X2=\allowbreak{}e.


\entrylabel{Remaining obstruction}

This repackages a known representation into explicit but rapidly growing finite-tableau coefficients and an infinite sum. It does not establish a new tractable closed joint law or its novelty,\allowbreak{} and should not be promoted as a full solution to the source request.


Verification:\allowbreak{} author checked limited argument | Independent review:\allowbreak{} none | Full target solved:\allowbreak{} False


\clearpage

\section*{P0221  Limit of the optimal online random spanning-tree cost}

\entrylabel{Primary source}

Aldous-Angel-Berestycki:\allowbreak{} Online random weight minimal spanning trees,\allowbreak{} draft2008;\allowbreak{} Conjecture PDFp2;\allowbreak{}Theorem3 PDFp4;\allowbreak{} author open page


\url{https://www.stat.berkeley.edu/~aldous/Research/OP/Beres2008.pdf}


Source statement date:\allowbreak{} 2008-06-16;\allowbreak{}authorpage2018. Archived access:\allowbreak{} 2026-09-06T22:\allowbreak{}39:\allowbreak{}04Z.


Historical:\allowbreak{} source\_\allowbreak{}says\_\allowbreak{}open | Current:\allowbreak{} unverified | Result:\allowbreak{} attempted\_\allowbreak{}unsolved


\entrylabel{Question}

Prove existence and determine the limiting optimal mean cost of an online spanning tree of the complete randomly weighted graph.


\entrylabel{Scope}

Combinatorial random-order and Poisson offer models are grouped,\allowbreak{} but their finite-n costs differ. The draft explicitly lacks a complete asymptotic equivalence proof. Numerical estimates near1.9 are not certified limits.


\entrylabel{Dated literature check}

Author-maintained page and relevant linked primary sources checked2026-09-06. Focused later-literature searches did not establish a complete resolution;\allowbreak{} current status remains unverified.


\subsection*{P0221-A001 | Limit of the optimal online random spanning-tree cost structural calculation}

\entrylabel{Argument or result}

The Poisson dynamic-programming recursion has exact small-state solutions. For two current components of sizes x,\allowbreak{}y,\allowbreak{} F=\allowbreak{}(xy/\allowbreak{}2)\emph{F\textasciicircum{}2 and the positive solution is2/\allowbreak{}(xy). For three singleton vertices each merged state has cost1,\allowbreak{} giving F=\allowbreak{}(3/\allowbreak{}2)}(F-1)\textasciicircum{}2. The admissible root exceeding the child cost is(4+\allowbreak{}sqrt7)/\allowbreak{}3;\allowbreak{} the other root is below1 and violates strict monotonicity.


\entrylabel{Second check}

The first recursion follows by integrating the accepted-offer surplus over weights below the continuation-cost difference. For n2 the Poisson cost is2 while the combinatorial model must buy its sole Uniform(0,\allowbreak{}1) edge and has cost1/\allowbreak{}2,\allowbreak{} demonstrating why finite-n model identification is invalid.


\entrylabel{Remaining obstruction}

The recursion depends on the entire component-size partition and on n. Standard coagulation hydrodynamic limits assume a prescribed kernel,\allowbreak{} whereas these acceptance kernels are themselves defined by unknown value functions. A compactness and uniqueness argument for their joint limit is missing.


Verification:\allowbreak{} author checked limited argument | Independent review:\allowbreak{} none | Full target solved:\allowbreak{} False


\clearpage

\section*{P0222  Near-unit stretch versus length in spatial networks}

\entrylabel{Primary source}

Aldous-Lando:\allowbreak{} The stretch-length tradeoff in geometric networks;\allowbreak{} Abstract and dense-network limit question;\allowbreak{} author open page


\url{https://arxiv.org/pdf/1404.2653}


Source statement date:\allowbreak{} 2015. Archived access:\allowbreak{} 2026-09-06T22:\allowbreak{}48:\allowbreak{}41Z.


Historical:\allowbreak{} source\_\allowbreak{}says\_\allowbreak{}open | Current:\allowbreak{} unverified | Result:\allowbreak{} attempted\_\allowbreak{}unsolved


\entrylabel{Question}

Establish existence and value of the exponent governing minimum network length per unit area as the allowed all-pairs stretch s decreases to1.


\entrylabel{Scope}

The Poisson average-case problem is the probability target;\allowbreak{} the source worst-case analogue is grouped as a comparison. This is different from the mean-route-length constraint for SIRSNs.


\entrylabel{Dated literature check}

Primary source and available author updates read2026-09-06;\allowbreak{} focused later searches recorded in source registry. No current-open certification or novelty claim.


\subsection*{P0222-A001 | Near-unit stretch versus length in spatial networks structural attempt}

\entrylabel{Argument or result}

Every path of stretch at most1+\allowbreak{}epsilon between endpoints at distance r stays in an ellipse whose maximum transverse half-width is(r/\allowbreak{}2)*sqrt(2epsilon+\allowbreak{}epsilon\textasciicircum{}2). Indeed a point of deviation T forces the total path length to be at least sqrt(r\textasciicircum{}2+\allowbreak{}4T\textasciicircum{}2). Thus as epsilon decreases,\allowbreak{} each local pair needs a usable road within a tube of angular width order sqrt(epsilon). This is an explicit geometric constraint for a possible direction-packing lower bound.


\entrylabel{Second check}

The ellipse inequality follows by minimizing the sum of distances from a path point to its two endpoints over the longitudinal coordinate. It is shared with the deterministic lemma in P0214,\allowbreak{} but the sourced optimization problem here has a different target.


\entrylabel{Remaining obstruction}

Counting constrained pairs is insufficient because one road segment can serve many pairs. A stochastic-geometric bound on this reuse,\allowbreak{} or a matching construction with controlled length,\allowbreak{} is needed to obtain a sharp exponent.


Verification:\allowbreak{} author checked limited argument | Independent review:\allowbreak{} none | Full target solved:\allowbreak{} False


\clearpage

\section*{P0223  Sharp relaxation exponent for local tree rotations and triangulation flips}

\entrylabel{Primary source}

Aldous:\allowbreak{} Mixing times for branch rotation and triangulation walk;\allowbreak{} PROBLEM,\allowbreak{}PDFp2;\allowbreak{} historyp3


\url{https://www.stat.berkeley.edu/~aldous/206-RWG/RWGpapers/clad-mix.pdf}


Source statement date:\allowbreak{} circa2000;\allowbreak{}authorpageupdate2018. Archived access:\allowbreak{} 2026-09-06T22:\allowbreak{}48:\allowbreak{}36Z.


Historical:\allowbreak{} source\_\allowbreak{}says\_\allowbreak{}open | Current:\allowbreak{} unverified | Result:\allowbreak{} attempted\_\allowbreak{}unsolved


\entrylabel{Question}

Prove an order n\textasciicircum{}(3/\allowbreak{}2) relaxation-time upper bound for the branch-rotation cladogram chain and the related convex-polygon triangulation flip chain.


\entrylabel{Scope}

The two chains are conservatively grouped as the source rotation family. The2026 Alev-Frishberg-Sarantis-Tetali theorem improves the triangulation bound to n\textasciicircum{}2 times polylogarithms and explicitly leaves the conjectured exponent unresolved.


\entrylabel{Dated literature check}

Primary source and available author updates read2026-09-06;\allowbreak{} focused later searches recorded in source registry. No current-open certification or novelty claim.


\subsection*{P0223-A001 | Sharp relaxation exponent for local tree rotations and triangulation flips structural attempt}

\entrylabel{Argument or result}

A canonical-path construction that first rotates every tree to one fixed comb has exponential congestion,\allowbreak{} even if each route is short. Let N be the number of states and use uniform stationary mass1/\allowbreak{}N. Routing the all-pairs demand pi(x)*pi(y) through that comb sends total order1 flow through a vertex whose total stationary transition capacity is1/\allowbreak{}N. Some incident edge therefore has load-to-capacity ratio at least order N. For triangulations N is a Catalan number of order4\textasciicircum{}n/\allowbreak{}n\textasciicircum{}(3/\allowbreak{}2).


\entrylabel{Second check}

The total capacity through the chosen vertex is sum\_\allowbreak{}y pi(v)P(v,\allowbreak{}y)=\allowbreak{}1/\allowbreak{}N. Pairs not starting or ending there have total demand(1-1/\allowbreak{}N)\textasciicircum{}2 and must cross it in the proposed construction. Summing their edge loads gives the claimed bottleneck independently of path length.


\entrylabel{Remaining obstruction}

A proof needs routes or transport flows distributed over many intermediate shapes. The comb normal form is useful for irreducibility but cannot furnish the conjectured polynomial relaxation bound.


Verification:\allowbreak{} author checked limited argument | Independent review:\allowbreak{} none | Full target solved:\allowbreak{} False


\clearpage

\section*{P0224  Feasible length and stretch statistics of scale-invariant networks}

\entrylabel{Primary source}

Aldous:\allowbreak{} Scale-invariant random spatial networks;\allowbreak{} Open Problems29 and30(i),\allowbreak{}PDFp48


\url{https://arxiv.org/pdf/1204.0817}


Source statement date:\allowbreak{} 2012preprint;\allowbreak{}2014journal. Archived access:\allowbreak{} 2026-09-06T22:\allowbreak{}15:\allowbreak{}59Z.


Historical:\allowbreak{} source\_\allowbreak{}says\_\allowbreak{}open | Current:\allowbreak{} unverified | Result:\allowbreak{} attempted\_\allowbreak{}unsolved


\entrylabel{Question}

Describe the feasible triples(Delta,\allowbreak{}ell,\allowbreak{}p(1)) for SIRSNs and obtain sharp bounds for the minimum sampled-network length ell at fixed mean route stretch Delta.


\entrylabel{Scope}

The feasible-region question and its lower-boundary function are one optimization family. Source30(ii),\allowbreak{} attainment,\allowbreak{} is separately admitted because it concerns closure and existence.


\entrylabel{Dated literature check}

Primary source and available author updates read2026-09-06;\allowbreak{} focused later searches recorded in source registry. No current-open certification or novelty claim.


\subsection*{P0224-A001 | Feasible length and stretch statistics of scale-invariant networks structural attempt}

\entrylabel{Argument or result}

The feasible set is convex under the source axioms. Choose one of two entire SIRSN laws once using an independent Bernoulli variable,\allowbreak{} and use that law for all routes. Every realization retains compatibility and all required distributional invariances. Each of Delta,\allowbreak{}ell,\allowbreak{}p(1) is an expectation or monotone limit of expectations,\allowbreak{} hence is the same convex combination of the corresponding statistics. Taking near-minimizers shows ell\_\allowbreak{}*(lambda*Delta1+\allowbreak{}(1-lambda)*Delta2)<=\allowbreak{}lambda*ell\_\allowbreak{}*(Delta1)+\allowbreak{}(1-lambda)*ell\_\allowbreak{}*(Delta2).


\entrylabel{Second check}

The random choice is global,\allowbreak{} so routes from the two networks are never spliced together. No ergodicity assumption appears in the cited definition;\allowbreak{} imposing one would invalidate this mixture argument. Finite expected statistics remain finite for a two-law mixture.


\entrylabel{Remaining obstruction}

Convexity does not locate the feasible region or give its sharp boundary. It also does not guarantee an extremal law exists;\allowbreak{} a convex set need not be closed.


Verification:\allowbreak{} author checked limited argument | Independent review:\allowbreak{} none | Full target solved:\allowbreak{} False


\clearpage

\section*{P0225  Attainment of the optimal SIRSN length-stretch tradeoff}

\entrylabel{Primary source}

Aldous:\allowbreak{} Scale-invariant random spatial networks;\allowbreak{} Open Problem30(ii),\allowbreak{}PDFp48


\url{https://arxiv.org/pdf/1204.0817}


Source statement date:\allowbreak{} 2012preprint;\allowbreak{}2014journal. Archived access:\allowbreak{} 2026-09-06T22:\allowbreak{}15:\allowbreak{}59Z.


Historical:\allowbreak{} source\_\allowbreak{}says\_\allowbreak{}open | Current:\allowbreak{} unverified | Result:\allowbreak{} attempted\_\allowbreak{}unsolved


\entrylabel{Question}

Does a SIRSN attain the infimum sampled-network length at a prescribed mean route stretch,\allowbreak{} and what structure can an optimizer have?


\entrylabel{Scope}

This is the existence part of the source optimization question,\allowbreak{} not a separately counted parameter choice.


\entrylabel{Dated literature check}

Primary source and available author updates read2026-09-06;\allowbreak{} focused later searches recorded in source registry. No current-open certification or novelty claim.


\subsection*{P0225-A001 | Attainment of the optimal SIRSN length-stretch tradeoff structural attempt}

\entrylabel{Argument or result}

Uniform expected length in every compact set supplies tightness for the unlabelled length measures:\allowbreak{} Markov bounds on an exhausting sequence of compact sets control local masses,\allowbreak{} and bounded-mass sets of Radon measures are vaguely relatively compact. However the prescribed feasible-route class is not closed under even uniform convergence with bounded lengths. Inscribed polygonal approximations to a semicircle are finite unions of straight segments with uniformly bounded length,\allowbreak{} while their limit is a smooth arc and is not a jagged route made of straight segments.


\entrylabel{Second check}

For the semicircle,\allowbreak{} no nontrivial subarc is a line segment;\allowbreak{} countably many straight segments lying in the arc can therefore cover only countably many points. The source itself discusses a possible extension from jagged to rectifiable routes.


\entrylabel{Remaining obstruction}

A compactness proof for length measures alone loses route labels,\allowbreak{} compatibility and the allowed-curve condition. One must choose a richer topology and prove closure or justify a relaxed model before claiming attainment.


Verification:\allowbreak{} author checked limited argument | Independent review:\allowbreak{} none | Full target solved:\allowbreak{} False


\clearpage

\section*{P0226  Local finiteness of scale-invariant network traffic measures}

\entrylabel{Primary source}

Aldous:\allowbreak{} Scale-invariant random spatial networks;\allowbreak{} Open Problem31,\allowbreak{}PDFp49


\url{https://arxiv.org/pdf/1204.0817}


Source statement date:\allowbreak{} 2012preprint;\allowbreak{}2014journal. Archived access:\allowbreak{} 2026-09-06T22:\allowbreak{}15:\allowbreak{}59Z.


Historical:\allowbreak{} source\_\allowbreak{}says\_\allowbreak{}open | Current:\allowbreak{} unverified | Result:\allowbreak{} attempted\_\allowbreak{}unsolved


\entrylabel{Question}

For source-destination displacement density|z|\textasciicircum{}-beta with2<beta<4,\allowbreak{} prove suitable local finiteness of induced route traffic on the major-road set E(infinity,\allowbreak{}1).


\entrylabel{Scope}

The target is traffic restricted to major roads. Finiteness of the expected traffic on the full dense route network is a different,\allowbreak{} stronger statement and fails by the calculation below.


\entrylabel{Dated literature check}

Primary source and available author updates read2026-09-06;\allowbreak{} focused later searches recorded in source registry. No current-open certification or novelty claim.


\subsection*{P0226-A001 | Local finiteness of scale-invariant network traffic measures structural attempt}

\entrylabel{Argument or result}

For a bounded area A,\allowbreak{} translation invariance and Tonelli give expected full-network traffic through A equal to area(A)*integral\_\allowbreak{}R2 |z|\textasciicircum{}-beta*E len R(0,\allowbreak{}z) dz. Scale invariance makes E len R(0,\allowbreak{}z)=\allowbreak{}Delta*|z|,\allowbreak{} yielding area(A)*2*pi*Delta*integral\_\allowbreak{}0\textasciicircum{}infinity r\textasciicircum{}(2-beta) dr. This integral diverges for every beta:\allowbreak{} at infinity for beta<=\allowbreak{}3,\allowbreak{} at zero for beta>=\allowbreak{}3.


\entrylabel{Second check}

Translate each route by its source and integrate the length of its intersection with A over all source locations. Each infinitesimal route segment contributes area(A). This mass-transport identity uses nonnegative integrands,\allowbreak{} so it remains valid when the integral is infinite.


\entrylabel{Remaining obstruction}

The divergent full-network expectation neither disproves almost-sure local finiteness nor answers the restricted-major-road question. A proof must exploit the major-road restriction and separately control long routes and short routes which take large excursions.


Verification:\allowbreak{} author checked limited argument | Independent review:\allowbreak{} none | Full target solved:\allowbreak{} False


\clearpage

\section*{P0227  Implications among SIRSN finiteness and geodesic properties}

\entrylabel{Primary source}

Aldous:\allowbreak{} Scale-invariant random spatial networks;\allowbreak{} Open Problem32(i)-(iv),\allowbreak{}PDFp49;\allowbreak{}definitionspp10-11,\allowbreak{}44-45


\url{https://arxiv.org/pdf/1204.0817}


Source statement date:\allowbreak{} 2012preprint;\allowbreak{}2014journal. Archived access:\allowbreak{} 2026-09-06T22:\allowbreak{}15:\allowbreak{}59Z.


Historical:\allowbreak{} source\_\allowbreak{}says\_\allowbreak{}open | Current:\allowbreak{} unverified | Result:\allowbreak{} attempted\_\allowbreak{}unsolved


\entrylabel{Question}

Determine the four proposed reverse implications between weak and full SIRSN finiteness,\allowbreak{} unique rays,\allowbreak{} pairwise coalescence,\allowbreak{} uniform coalescence and endpoint route continuity.


\entrylabel{Scope}

All four linked implications are grouped as one source family. Equation16 is finite ell;\allowbreak{}20 is finite p(1);\allowbreak{}49 is pairwise ray coalescence;\allowbreak{}50 is uniform small-neighborhood-to-far-target coalescence;\allowbreak{}51 is route coincidence away from nearby endpoints.


\entrylabel{Dated literature check}

Primary source and available author updates read2026-09-06;\allowbreak{} focused later searches recorded in source registry. No current-open certification or novelty claim.


\subsection*{P0227-A001 | Implications among SIRSN finiteness and geodesic properties structural attempt}

\entrylabel{Argument or result}

The available finite-intensity estimate does not yield the first implication:\allowbreak{} p(lambda,\allowbreak{}1)<=\allowbreak{}intensity S(lambda)=\allowbreak{}ell*sqrt(lambda),\allowbreak{} which diverges with the sampling density. For the third implication,\allowbreak{} finite coalescence radii at each finite sample do not bound their increasing supremum;\allowbreak{} a concrete scalar analogue is R\_\allowbreak{}k=\allowbreak{}k,\allowbreak{} finite at every k and divergent in the limit. For the second,\allowbreak{} stationary families of parallel rays show that uniqueness of a ray at each origin does not alone imply coalescence. For the fourth,\allowbreak{} continuity at each fixed destination gives no uniform estimate as the destination recedes to infinity.


\entrylabel{Second check}

The scaling identity for sampled intensity follows by rescaling a rate1 Poisson process by lambda\textasciicircum{}-1/\allowbreak{}2. The parallel-ray construction is deliberately only a counterexample to the isolated logical inference;\allowbreak{} it is not claimed to satisfy the full connected-route SIRSN axioms.


\entrylabel{Remaining obstruction}

Actual counterexamples must preserve compatibility,\allowbreak{} all scale invariances and finite required moments. Alternatively a proof needs a uniform estimate absent from the displayed bounds. No member of the four-part family is settled here.


Verification:\allowbreak{} author checked limited argument | Independent review:\allowbreak{} none | Full target solved:\allowbreak{} False


\clearpage

\section*{P0228  Uniqueness of the unbounded major-road component}

\entrylabel{Primary source}

Aldous:\allowbreak{} Scale-invariant random spatial networks;\allowbreak{} Open Problem33,\allowbreak{}PDFp49


\url{https://arxiv.org/pdf/1204.0817}


Source statement date:\allowbreak{} 2012preprint;\allowbreak{}2014journal. Archived access:\allowbreak{} 2026-09-06T22:\allowbreak{}15:\allowbreak{}59Z.


Historical:\allowbreak{} source\_\allowbreak{}says\_\allowbreak{}open | Current:\allowbreak{} unverified | Result:\allowbreak{} attempted\_\allowbreak{}unsolved


\entrylabel{Question}

Must the major-road set E(infinity,\allowbreak{}1) of every SIRSN have only one unbounded connected component almost surely?


\entrylabel{Scope}

Disconnected finite major-road pieces are allowed by definition. The question concerns unbounded components only.


\entrylabel{Dated literature check}

Primary source and available author updates read2026-09-06;\allowbreak{} focused later searches recorded in source registry. No current-open certification or novelty claim.


\subsection*{P0228-A001 | Uniqueness of the unbounded major-road component structural attempt}

\entrylabel{Argument or result}

A precise sufficient reduction is available:\allowbreak{} if every unbounded major-road component contains a tail of a singly infinite geodesic and all such geodesics coalesce,\allowbreak{} then uniqueness follows. Two purported unbounded components would contain two tails,\allowbreak{} whose common eventual tail lies in both,\allowbreak{} forcing the components to coincide.


\entrylabel{Second check}

The conclusion uses connected-component maximality,\allowbreak{} not merely that the full network is connected. Finite local total length does not by itself supply the needed geodesic tail in every unbounded component,\allowbreak{} because the route union is not assumed to be a locally finite combinatorial graph.


\entrylabel{Remaining obstruction}

Neither the universal ray-coalescence property nor the component-tail property is proved from the source axioms. A Burton-Keane style argument would require a valid local-modification principle,\allowbreak{} which arbitrary SIRSN laws need not have.


Verification:\allowbreak{} author checked limited argument | Independent review:\allowbreak{} none | Full target solved:\allowbreak{} False


\clearpage

\section*{P0229  Integrable supremum of routes to a bounded region}

\entrylabel{Primary source}

Aldous:\allowbreak{} Scale-invariant random spatial networks;\allowbreak{} Open Problem34,\allowbreak{}PDFp50


\url{https://arxiv.org/pdf/1204.0817}


Source statement date:\allowbreak{} 2012preprint;\allowbreak{}2014journal. Archived access:\allowbreak{} 2026-09-06T22:\allowbreak{}15:\allowbreak{}59Z.


Historical:\allowbreak{} source\_\allowbreak{}says\_\allowbreak{}open | Current:\allowbreak{} unverified | Result:\allowbreak{} attempted\_\allowbreak{}unsolved


\entrylabel{Question}

Under what additional assumptions does the expected supremum of route lengths from0 to an infinite independent uniform sample in the unit disk remain finite?


\entrylabel{Scope}

Pointwise E len R(0,\allowbreak{}x)=\allowbreak{}Delta*|x| does not control a random supremum. The sufficient increment condition below is stronger than the base axioms and is not claimed to follow from them.


\entrylabel{Dated literature check}

Primary source and available author updates read2026-09-06;\allowbreak{} focused later searches recorded in source registry. No current-open certification or novelty claim.


\subsection*{P0229-A001 | Integrable supremum of routes to a bounded region structural attempt}

\entrylabel{Argument or result}

Let L(x)=\allowbreak{}len R(0,\allowbreak{}x). Suppose on a containing square E|L(x)-L(y)|\textasciicircum{}p<=\allowbreak{}C*|x-y|\textasciicircum{}p for a p>2,\allowbreak{} with a jointly measurable version and finite mean at one anchor. On the dyadic grid of mesh2\textasciicircum{}-k there are order2\textasciicircum{}(2k) parent-child or neighbor pairs. The expected maximum increment is at most Cprime*2\textasciicircum{}(-k*(1-2/\allowbreak{}p)). This is summable,\allowbreak{} so telescoping along grid ancestors constructs a continuous modification with finite expected supremum.


\entrylabel{Second check}

Use max\_\allowbreak{}j |Z\_\allowbreak{}j|<=\allowbreak{}(sum\_\allowbreak{}j |Z\_\allowbreak{}j|\textasciicircum{}p)\textasciicircum{}(1/\allowbreak{}p) and Jensen,\allowbreak{} then sum the resulting geometric series. The increment bound implies stochastic continuity,\allowbreak{} identifying the modification at each deterministic point;\allowbreak{} Fubini identifies it at each independent uniform sample point,\allowbreak{} and countability handles the entire sampled sequence.


\entrylabel{Remaining obstruction}

The source only gives first moments of individual route lengths,\allowbreak{} not p-moment increments of L. Its compatibility condition does not make Euclidean route length an automatic metric with a triangle inequality,\allowbreak{} so a metric-chaining shortcut is unjustified.


Verification:\allowbreak{} author checked limited argument | Independent review:\allowbreak{} none | Full target solved:\allowbreak{} False


\clearpage

\section*{P0230  Asymptotic total length of finite SIRSN samples}

\entrylabel{Primary source}

Aldous:\allowbreak{} Scale-invariant random spatial networks;\allowbreak{} Open Problem35,\allowbreak{}PDFp50


\url{https://arxiv.org/pdf/1204.0817}


Source statement date:\allowbreak{} 2012preprint;\allowbreak{}2014journal. Archived access:\allowbreak{} 2026-09-06T22:\allowbreak{}15:\allowbreak{}59Z.


Historical:\allowbreak{} source\_\allowbreak{}says\_\allowbreak{}open | Current:\allowbreak{} unverified | Result:\allowbreak{} attempted\_\allowbreak{}unsolved


\entrylabel{Question}

For k uniform points in a square of area k,\allowbreak{} prove that the mean total length of their spanning route network is asymptotic to ell*k.


\entrylabel{Scope}

Length includes route pieces outside the sampling square. This global excursion contribution is the main source obstruction.


\entrylabel{Dated literature check}

Primary source and available author updates read2026-09-06;\allowbreak{} focused later searches recorded in source registry. No current-open certification or novelty claim.


\subsection*{P0230-A001 | Asymptotic total length of finite SIRSN samples structural attempt}

\entrylabel{Argument or result}

A local truncation argument gives the lower bound liminf E length/\allowbreak{}k>=\allowbreak{}ell. Tile the bulk of the large square by unit cells. For fixed R and M,\allowbreak{} retain in each cell only the length,\allowbreak{} truncated at M,\allowbreak{} of routes between sampled endpoints within distance R of that cell. Its local binomial point configuration converges in total variation to the corresponding rate1 Poisson configuration. The retained length is bounded by M,\allowbreak{} so its expectation converges. Summing disjoint unit-cell contributions and then taking k to infinity gives the Poisson local truncated expectation as a lower bound. Let M and then R increase:\allowbreak{} monotone convergence recovers the full sampled-network intensity ell.


\entrylabel{Second check}

Only a bounded endpoint region is used before taking k to infinity;\allowbreak{} its binomial count tends to a Poisson count and conditional positions have the same uniform law. Truncation avoids any unproved uniform-integrability step. Boundary cells are an asymptotically negligible fraction for fixed R.


\entrylabel{Remaining obstruction}

This proves no upper bound on total length. Rare paths can leave the sampling square by a large distance,\allowbreak{} and their exterior contribution is absent from the local lower-bound argument. Controlling those excursions remains necessary for equality.


Verification:\allowbreak{} author checked limited argument | Independent review:\allowbreak{} none | Full target solved:\allowbreak{} False


\clearpage

\section*{P0231  A zero-one law for Poisson-coloring frontier dimension}

\entrylabel{Primary source}

Basdevant-Blanc-Curien-Singh:\allowbreak{} Fractal properties of the frontier in Poissonian coloring,\allowbreak{}2023;\allowbreak{} Discussion before Open Question1.2,\allowbreak{}PDFp3


\url{https://alea.impa.br/articles/v20/20-59.pdf}


Source statement date:\allowbreak{} 2023-11. Archived access:\allowbreak{} 2026-09-06T22:\allowbreak{}39:\allowbreak{}35Z.


Historical:\allowbreak{} source\_\allowbreak{}says\_\allowbreak{}open | Current:\allowbreak{} unverified | Result:\allowbreak{} attempted\_\allowbreak{}unsolved


\entrylabel{Question}

Is the Hausdorff dimension of the limiting Poisson-coloring frontier almost surely deterministic?


\entrylabel{Scope}

The strict inequalities d-1<dimension<d are already proved in Theorem1.1. This entry concerns the explicitly remaining zero-one question,\allowbreak{} not that resolved conjecture.


\entrylabel{Dated literature check}

Primary source and available author updates read2026-09-06;\allowbreak{} focused later searches recorded in source registry. No current-open certification or novelty claim.


\subsection*{P0231-A001 | A zero-one law for Poisson-coloring frontier dimension structural attempt}

\entrylabel{Argument or result}

The dimension event is not visibly a tail event of the independent arriving points. Changing an early point can change the color of infinitely many later descendants,\allowbreak{} so Kolmogorov zero-one cannot be invoked without a finite-modification invariance proof. The potential route is a regeneration principle producing,\allowbreak{} inside every surviving interface neighborhood,\allowbreak{} copies with the full essential supremum of the dimension law;\allowbreak{} countable-union stability of dimension would then force the maximal value.


\entrylabel{Second check}

The mere facts that the frontier is closed and has dimension strictly between1 and2 do not imply determinism:\allowbreak{} a Bernoulli mixture of two fixed compact fractal curves with different dimensions satisfies those facts and has random dimension. This is a logical control,\allowbreak{} not a Poisson-coloring counterexample.


\entrylabel{Remaining obstruction}

Independent arrivals do not make recursively inherited colors independent across subregions. A regeneration or local-universality statement strong enough to remove early-history effects is missing.


Verification:\allowbreak{} author checked limited argument | Independent review:\allowbreak{} none | Full target solved:\allowbreak{} False


\clearpage

\section*{P0232  Continuous curves in the Poisson-coloring frontier}

\entrylabel{Primary source}

Basdevant-Blanc-Curien-Singh:\allowbreak{} Fractal properties of the frontier in Poissonian coloring,\allowbreak{}2023;\allowbreak{} Open Question1.2,\allowbreak{}PDFpp3-4


\url{https://alea.impa.br/articles/v20/20-59.pdf}


Source statement date:\allowbreak{} 2023-11. Archived access:\allowbreak{} 2026-09-06T22:\allowbreak{}39:\allowbreak{}35Z.


Historical:\allowbreak{} source\_\allowbreak{}says\_\allowbreak{}open | Current:\allowbreak{} unverified | Result:\allowbreak{} attempted\_\allowbreak{}unsolved


\entrylabel{Question}

Is the planar limiting frontier a countable union of continuous curves?


\entrylabel{Scope}

A curve means an image of a continuous map from[0,\allowbreak{}1]. It is not required to be rectifiable. Point and segment variants are grouped.


\entrylabel{Dated literature check}

Primary source and available author updates read2026-09-06;\allowbreak{} focused later searches recorded in source registry. No current-open certification or novelty claim.


\subsection*{P0232-A001 | Continuous curves in the Poisson-coloring frontier structural attempt}

\entrylabel{Argument or result}

The established dimension lower bound already rules out a countable union of rectifiable curves:\allowbreak{} every rectifiable curve has Hausdorff dimension at most1,\allowbreak{} and the dimension of a countable union is the supremum of the component dimensions. Thus any affirmative construction must use nonrectifiable curves. Compact connectedness alone also cannot supply path connectedness,\allowbreak{} as shown by the topologist sine-curve closure.


\entrylabel{Second check}

The sine-curve closure of\{(x,\allowbreak{}sin(1/\allowbreak{}x)):\allowbreak{}0<x<=\allowbreak{}1\} together with the limiting vertical segment is compact and connected but not path connected. This does not by itself refute a countable-curve covering;\allowbreak{} it specifically blocks the inference from connectedness to a single continuous parametrization.


\entrylabel{Remaining obstruction}

A scale-uniform crossing/\allowbreak{}tightness estimate is needed to construct continuous interface paths. The source lacks the correlation inequality used to establish such estimates in standard planar models.


Verification:\allowbreak{} author checked limited argument | Independent review:\allowbreak{} none | Full target solved:\allowbreak{} False


\clearpage

\section*{P0233  Absence of double points on Poisson-coloring interface curves}

\entrylabel{Primary source}

Basdevant-Blanc-Curien-Singh:\allowbreak{} Fractal properties of the frontier in Poissonian coloring,\allowbreak{}2023;\allowbreak{} Open Question1.3,\allowbreak{}PDFp4


\url{https://alea.impa.br/articles/v20/20-59.pdf}


Source statement date:\allowbreak{} 2023-11. Archived access:\allowbreak{} 2026-09-06T22:\allowbreak{}39:\allowbreak{}35Z.


Historical:\allowbreak{} source\_\allowbreak{}says\_\allowbreak{}open | Current:\allowbreak{} unverified | Result:\allowbreak{} attempted\_\allowbreak{}unsolved


\entrylabel{Question}

If the limiting frontier is a union of curves,\allowbreak{} are its component curves almost surely simple?


\entrylabel{Scope}

This conditional question is distinct from existence of curves. The segment-model single-curve consequence is grouped.


\entrylabel{Dated literature check}

Primary source and available author updates read2026-09-06;\allowbreak{} focused later searches recorded in source registry. No current-open certification or novelty claim.


\subsection*{P0233-A001 | Absence of double points on Poisson-coloring interface curves structural attempt}

\entrylabel{Argument or result}

A quantitative four-arm estimate would settle the proposed double-point obstruction. In a bounded region cover by order epsilon\textasciicircum{}-2 boxes. If the probability that a box supports four alternating color arms from scale epsilon to a fixed macroscopic radius is at most C*epsilon\textasciicircum{}(2+\allowbreak{}eta),\allowbreak{} a union bound makes the chance of any such box at most Cprime*epsilon\textasciicircum{}eta. Taking epsilon down to0 and then countably many macroscopic radii excludes points with four nontrivial alternating arms.


\entrylabel{Second check}

A double point with four arms reaching a positive radius must lie in one of the grid boxes and induce a slightly enlarged box arm event. Changing fixed geometric constants does not alter the positive exponent margin.


\entrylabel{Remaining obstruction}

No four-arm exponent exceeding2 is established for this dependent coloring process. Excluding four-arm points would still have to be connected carefully to the exact parametrization and component statement,\allowbreak{} rather than silently equated with every form of self-touching.


Verification:\allowbreak{} author checked limited argument | Independent review:\allowbreak{} none | Full target solved:\allowbreak{} False


\clearpage

\section*{P0234  A safe margin around extended initial color seeds}

\entrylabel{Primary source}

Basdevant-Blanc-Curien-Singh:\allowbreak{} Fractal properties of the frontier in Poissonian coloring,\allowbreak{}2023;\allowbreak{} Open Question1.4,\allowbreak{}PDFp4


\url{https://alea.impa.br/articles/v20/20-59.pdf}


Source statement date:\allowbreak{} 2023-11. Archived access:\allowbreak{} 2026-09-06T22:\allowbreak{}39:\allowbreak{}35Z.


Historical:\allowbreak{} source\_\allowbreak{}says\_\allowbreak{}open | Current:\allowbreak{} unverified | Result:\allowbreak{} attempted\_\allowbreak{}unsolved


\entrylabel{Question}

If the initial red seed is a segment or ball disjoint from the blue seed,\allowbreak{} does the limiting blue region almost surely avoid the entire initial red seed?


\entrylabel{Scope}

Pointwise protection of each fixed initial point is established in related arguments and does not imply simultaneous protection of a continuum. Disjoint initial colors are part of the meaningful scope.


\entrylabel{Dated literature check}

Primary source and available author updates read2026-09-06;\allowbreak{} focused later searches recorded in source registry. No current-open certification or novelty claim.


\subsection*{P0234-A001 | A safe margin around extended initial color seeds structural attempt}

\entrylabel{Argument or result}

A sufficient uniform estimate can be stated precisely. Let the compact seed K have covering number at most C*delta\textasciicircum{}-d0. If every x in K satisfies P(dist(x,\allowbreak{}B\_\allowbreak{}infinity)<=\allowbreak{}2delta)<=\allowbreak{}Cprime*delta\textasciicircum{}alpha uniformly,\allowbreak{} with alpha>d0,\allowbreak{} a delta-net and union bound give P(K intersects B\_\allowbreak{}infinity)<=\allowbreak{}Cdouble*delta\textasciicircum{}(alpha-d0),\allowbreak{} hence zero. Compactness converts zero separation into intersection for the closed limiting blue set.


\entrylabel{Second check}

Pointwise zero intersection probability is inadequate:\allowbreak{} for a uniformly random U in a segment,\allowbreak{} every fixed x differs from U almost surely,\allowbreak{} but the random singleton\{U\} intersects the segment surely. The net argument explicitly identifies the missing uniform rate.


\entrylabel{Remaining obstruction}

The available shielding results give positive protection probability or fixed-point statements,\allowbreak{} not a uniform small-distance exponent exceeding the dimension of the seed. The full safe-margin question remains unresolved here.


Verification:\allowbreak{} author checked limited argument | Independent review:\allowbreak{} none | Full target solved:\allowbreak{} False


\clearpage

\section*{P0235  A Feller topology for limits of reversible Markov chains}

\entrylabel{Primary source}

Aldous:\allowbreak{} Limits for processes over general networks,\allowbreak{}2018;\allowbreak{} Conjecture,\allowbreak{}slide5;\allowbreak{}author current compact page


\url{https://www.stat.berkeley.edu/~aldous/Research/OP/talk_towsner_2018.pdf}


Source statement date:\allowbreak{} 2018. Archived access:\allowbreak{} 2026-09-06T22:\allowbreak{}39:\allowbreak{}17Z.


Historical:\allowbreak{} source\_\allowbreak{}says\_\allowbreak{}open | Current:\allowbreak{} unverified | Result:\allowbreak{} attempted\_\allowbreak{}unsolved


\entrylabel{Question}

Can the measurable limit process of finite reversible chains be given a natural complete separable state metric with the Feller property?


\entrylabel{Scope}

Towsner2015 already supplies the measure-theoretic existence and refined uniqueness results. This entry concerns only the remaining topological upgrade.


\entrylabel{Dated literature check}

Primary source and available author updates read2026-09-06;\allowbreak{} focused later searches recorded in source registry. No current-open certification or novelty claim.


\subsection*{P0235-A001 | A Feller topology for limits of reversible Markov chains structural attempt}

\entrylabel{Argument or result}

The suggested unweighted heat-kernel metric can be infinite in a basic admissible example. For Brownian motion on the two-torus,\allowbreak{} the squared row-distance integrated against exp(-t)dt equals sum\_\allowbreak{}(k in Z\textasciicircum{}2) |1-exp(2pi*i*k.h)|\textasciicircum{}2/\allowbreak{}[1+\allowbreak{}8pi\textasciicircum{}2*|k|\textasciicircum{}2]. At h=\allowbreak{}(1/\allowbreak{}2,\allowbreak{}0),\allowbreak{} every odd k1 has numerator4. Each dyadic annulus contains order N\textasciicircum{}2 such terms of order N\textasciicircum{}-2,\allowbreak{} so the series diverges. A repair is to use w(t)=\allowbreak{}exp(-t)/\allowbreak{}(1+\allowbreak{}G(2t)),\allowbreak{} where G is the integrated diagonal heat trace:\allowbreak{} Fubini gives finite weighted row norm for almost every state.


\entrylabel{Second check}

Parseval computes the row norm exactly;\allowbreak{} integrating the nonnegative Fourier series is justified by Tonelli. For the repaired weight,\allowbreak{} the average row squared norm is integral w(t)*G(2t)dt<=\allowbreak{}1. This produces an almost-everywhere finite Hilbert-space pseudometric after identifying indistinguishable rows.


\entrylabel{Remaining obstruction}

Brownian motion already has a good usual topology,\allowbreak{} so the divergence refutes only the illustrative formula,\allowbreak{} not the existence conjecture. Completing a heat-kernel embedding can add states;\allowbreak{} proving valid transition kernels and the Feller property on that completion is still missing.


Verification:\allowbreak{} author checked limited argument | Independent review:\allowbreak{} none | Full target solved:\allowbreak{} False


\clearpage

\section*{P0236  Infinite mean route length at each large fixed Euclidean separation}

\entrylabel{Primary source}

Aldous:\allowbreak{} Route Lengths in Invariant Spatial Tree Networks;\allowbreak{} Introduction conjecture,\allowbreak{}PDFp2;\allowbreak{}Section3.3


\url{https://arxiv.org/pdf/2103.00669}


Source statement date:\allowbreak{} 2021-03. Archived access:\allowbreak{} 2026-09-06T22:\allowbreak{}38:\allowbreak{}07Z.


Historical:\allowbreak{} source\_\allowbreak{}says\_\allowbreak{}open | Current:\allowbreak{} unverified | Result:\allowbreak{} attempted\_\allowbreak{}unsolved


\entrylabel{Question}

Is there a universal r0 such that every invariant tree network on planar rate1 Poisson points has infinite mean within-tree distance at almost every fixed separation r>=\allowbreak{}r0?


\entrylabel{Scope}

Theorem2 proves the averaged statement for pairs at distance at most r. The stronger fixed-separation conjecture and the authors possible thin-annulus counterconstruction remain distinct from that theorem.


\entrylabel{Dated literature check}

Primary source and available author updates read2026-09-06;\allowbreak{} focused later searches recorded in source registry. No current-open certification or novelty claim.


\subsection*{P0236-A001 | Infinite mean route length at each large fixed Euclidean separation structural attempt}

\entrylabel{Argument or result}

Divergence of the averaged mean cannot alone imply pointwise infinite means. A scalar control is rho(s)=\allowbreak{}1+\allowbreak{}1/\allowbreak{}|s-a| for s different from a,\allowbreak{} with any finite value assigned at a. Then rho is finite at every s,\allowbreak{} but integral\_\allowbreak{}0\textasciicircum{}r s*rho(s)ds is infinite for every r>a. Thus even all sufficiently large cumulative pair averages can diverge while no fixed-radius value is infinite.


\entrylabel{Second check}

The planar radial averaging density is proportional to s,\allowbreak{} so a singularity at a positive a survives the weighting. This example is only a function-level obstruction;\allowbreak{} it is not asserted to arise from any invariant tree network.


\entrylabel{Remaining obstruction}

One needs regularity or geometric constraints on the actual Palm mean function rho,\allowbreak{} or a completed invariant-tree counterconstruction. The source thin-annulus idea still requires cycle removal and connectivity without destroying the desired finite conditional means.


Verification:\allowbreak{} author checked limited argument | Independent review:\allowbreak{} none | Full target solved:\allowbreak{} False


\clearpage

\section*{P0242  Optimal stationary stopping under Markov-chain lumping}

\entrylabel{Primary source}

Aldous-Fill:\allowbreak{} Reversible Markov Chains and Random Walks on Graphs,\allowbreak{} 2002 recompiled2014;\allowbreak{} Open Problem4.45,\allowbreak{}PDF149;\allowbreak{} definition4.39,\allowbreak{}p148


\url{https://www.stat.berkeley.edu/~aldous/RWG/book.pdf}


Source statement date:\allowbreak{} 1994-10-11. Archived access:\allowbreak{} 2026-09-06T22:\allowbreak{}03:\allowbreak{}09Z.


Historical:\allowbreak{} source\_\allowbreak{}says\_\allowbreak{}open | Current:\allowbreak{} unverified | Result:\allowbreak{} attempted\_\allowbreak{}unsolved


\entrylabel{Question}

Determine the sharp universal comparison between the worst optimal expected stationary stopping time of a finite reversible chain and that of its stationary-flow lumping.


\entrylabel{Scope}

This is stationary-flow lumping by a map,\allowbreak{} not the chain watched on a subset. The source asks both a universal K and K=\allowbreak{}1;\allowbreak{} later lifting literature supplies comparison results whose exact constant and stopping convention still need checking.


\entrylabel{Dated literature check}

Focused primary-source searches on2026-09-06;\allowbreak{} no current-open or novelty certification. See source\_\allowbreak{}question\_\allowbreak{}registry\_\allowbreak{}markov\_\allowbreak{}002.json.


\subsection*{P0242-A001 | Optimal stationary stopping under Markov-chain lumping structural calculation}

\entrylabel{Argument or result}

For a genuinely lumpable partition,\allowbreak{} the K=\allowbreak{}1 inequality follows directly. Start the fine chain in any state of a chosen coarse block and use a fine-chain stopping time attaining stationarity. Its projection has the coarse-chain path law. Any extra fine information can be represented by auxiliary conditional randomness in the coarse filtration,\allowbreak{} so this yields a randomized coarse stopping rule with stationary terminal law and the same expected duration. Taking maxima proves the claim in this special case.


\entrylabel{Second check}

Equivalently disintegrate the stopped fine path given each successively observed coarse path;\allowbreak{} the conditional continuation probabilities give an adapted randomized coarse stopping rule. The conclusion requires strong lumpability,\allowbreak{} because otherwise the projected fine path is not the Markov chain defined by stationary-flow matching.


\entrylabel{Remaining obstruction}

For arbitrary partitions,\allowbreak{} equality of one-step stationary flows does not identify multi-step path laws. The stopping-time projection argument therefore fails at its first premise. A full comparison needs a different potential-theoretic argument and the later lifting result must be reconciled.


Verification:\allowbreak{} author checked limited argument | Independent review:\allowbreak{} none | Full target solved:\allowbreak{} False


\clearpage

\section*{P0243  Sharp constant-one comparisons of reversible mixing and stopping parameters}

\entrylabel{Primary source}

Aldous-Fill:\allowbreak{} Reversible Markov Chains and Random Walks on Graphs,\allowbreak{} 2002 recompiled2014;\allowbreak{} Open Problem4.48(a,\allowbreak{}b),\allowbreak{}PDF154;\allowbreak{} definitions119-120


\url{https://www.stat.berkeley.edu/~aldous/RWG/book.pdf}


Source statement date:\allowbreak{} 1994-10-11. Archived access:\allowbreak{} 2026-09-06T22:\allowbreak{}03:\allowbreak{}09Z.


Historical:\allowbreak{} source\_\allowbreak{}says\_\allowbreak{}open | Current:\allowbreak{} unverified | Result:\allowbreak{} attempted\_\allowbreak{}unsolved


\entrylabel{Question}

For every finite reversible chain,\allowbreak{} decide whether the pairwise total-variation e\textasciicircum{}-1 threshold tau1 is at most Kemeny mean hitting time tau0,\allowbreak{} and whether relaxation time tau2 is at most the worst optimal expected stationary stopping time.


\entrylabel{Scope}

The two displayed constant-one questions are grouped. The stopping half has a short complete eigenfunction derivation below,\allowbreak{} but this is not a full solution of the grouped problem or a novelty claim.


\entrylabel{Dated literature check}

Focused primary-source searches on2026-09-06;\allowbreak{} no current-open or novelty certification. See source\_\allowbreak{}question\_\allowbreak{}registry\_\allowbreak{}markov\_\allowbreak{}002.json.


\subsection*{P0243-A001 | Sharp constant-one comparisons of reversible mixing and stopping parameters structural calculation}

\entrylabel{Argument or result}

The stopping half follows from Dynkin optional sampling. Let Lf=\allowbreak{}-lambda*f with lambda=\allowbreak{}1/\allowbreak{}tau2,\allowbreak{} mean\_\allowbreak{}pi f=\allowbreak{}0,\allowbreak{} and choose x with f(x)=\allowbreak{}M=\allowbreak{}max f>0. For any integrable randomized stopping time U from x with X\_\allowbreak{}U distributed as pi,\allowbreak{} E f(X\_\allowbreak{}U)-f(x)=\allowbreak{}-lambda*E integral\_\allowbreak{}0\textasciicircum{}U f(X\_\allowbreak{}s)ds. Thus M=\allowbreak{}lambda*E integral f<=\allowbreak{}lambda*M*E U,\allowbreak{} giving E U>=\allowbreak{}tau2. The discrete-time proof replaces the integral by sum\_\allowbreak{}(s<U)f(X\_\allowbreak{}s).


\entrylabel{Second check}

Apply the identity first at U wedge n. The bounded function and integrability of U justify dominated convergence of the integral and of f(X\_\allowbreak{}(U wedge n)). This avoids the false assertion that an exponential eigenfunction martingale is automatically uniformly integrable at U.


\entrylabel{Remaining obstruction}

The bound on tau1 by tau0 remains unresolved here. Worst-start total variation cannot be bounded solely by an unweighted spectral trace without controlling eigenfunction amplitudes. The source convention uses the continuized threshold and pairwise total variation,\allowbreak{} not the usual distance-to-pi threshold.


Verification:\allowbreak{} author checked limited argument | Independent review:\allowbreak{} none | Full target solved:\allowbreak{} False


\clearpage

\section*{P0244  Sharp extremal walk parameters among regular graphs}

\entrylabel{Primary source}

Aldous-Fill:\allowbreak{} Reversible Markov Chains and Random Walks on Graphs,\allowbreak{} 2002 recompiled2014;\allowbreak{} Open Problem6.14,\allowbreak{}PDF217


\url{https://www.stat.berkeley.edu/~aldous/RWG/book.pdf}


Source statement date:\allowbreak{} 1994-10-31. Archived access:\allowbreak{} 2026-09-06T22:\allowbreak{}03:\allowbreak{}09Z.


Historical:\allowbreak{} source\_\allowbreak{}says\_\allowbreak{}open | Current:\allowbreak{} unverified | Result:\allowbreak{} attempted\_\allowbreak{}unsolved


\entrylabel{Question}

Determine the sharp asymptotic maxima of hitting,\allowbreak{} commute,\allowbreak{} cover,\allowbreak{} average-hitting and relaxation times over connected regular n-vertex graphs,\allowbreak{} as listed in the source necklace-graph conjecture.


\entrylabel{Scope}

All seven linked extremal formulas are grouped once. Abdi-Ghorbani2024 confirms the relaxation conjecture for additional near-maximum-diameter classes and describes the remaining general conjecture;\allowbreak{} this does not establish the whole family.


\entrylabel{Dated literature check}

Focused primary-source searches on2026-09-06;\allowbreak{} no current-open or novelty certification. See source\_\allowbreak{}question\_\allowbreak{}registry\_\allowbreak{}markov\_\allowbreak{}002.json.


\subsection*{P0244-A001 | Sharp extremal walk parameters among regular graphs structural calculation}

\entrylabel{Argument or result}

The cycle is an exact baseline but not the proposed extremizer:\allowbreak{} H(i,\allowbreak{}j)=\allowbreak{}r(n-r),\allowbreak{} maximal commute=\allowbreak{}2*floor(n\textasciicircum{}2/\allowbreak{}4),\allowbreak{} tau0=\allowbreak{}(n\textasciicircum{}2-1)/\allowbreak{}6,\allowbreak{} and tau2=\allowbreak{}1/\allowbreak{}(1-cos(2*pi/\allowbreak{}n)). In particular the cycle relaxation constant is1/\allowbreak{}(2*pi\textasciicircum{}2),\allowbreak{} one third of the proposed regular-graph extremal3/\allowbreak{}(2*pi\textasciicircum{}2). The source Feige cover-and-return bound implies tau2<=\allowbreak{}tau0<=\allowbreak{}tau*/\allowbreak{}2<=\allowbreak{}13*n\textasciicircum{}2/\allowbreak{}12,\allowbreak{} supplying a rigorous but nonsharp universal quadratic scale.


\entrylabel{Second check}

The cycle identities follow independently from the difference equation for hitting times and the Fourier spectrum. The spectral ordering tau2<=\allowbreak{}tau0 and stationary average of commute times tau0<=\allowbreak{}tau*/\allowbreak{}2 justify the chain of inequalities without confusing cover time with commute time.


\entrylabel{Remaining obstruction}

Optimizing diameter alone cannot prove the extremal spectral claim:\allowbreak{} the cited2024 work explicitly separates large diameter from minimal algebraic connectivity. Sharp resistance-profile control and all of the cover-function constants remain missing.


Verification:\allowbreak{} author checked limited argument | Independent review:\allowbreak{} none | Full target solved:\allowbreak{} False


\clearpage

\section*{P0245  Exact minimum stationary cover time for reversible chains}

\entrylabel{Primary source}

Aldous-Fill:\allowbreak{} Reversible Markov Chains and Random Walks on Graphs,\allowbreak{} 2002 recompiled2014;\allowbreak{} Open Problem6.30,\allowbreak{}PDF230


\url{https://www.stat.berkeley.edu/~aldous/RWG/book.pdf}


Source statement date:\allowbreak{} 1994-10-31. Archived access:\allowbreak{} 2026-09-06T22:\allowbreak{}03:\allowbreak{}09Z.


Historical:\allowbreak{} source\_\allowbreak{}says\_\allowbreak{}open | Current:\allowbreak{} unverified | Result:\allowbreak{} attempted\_\allowbreak{}unsolved


\entrylabel{Question}

Prove or refute E\_\allowbreak{}pi C>=\allowbreak{}(n-1)H\_\allowbreak{}(n-1) for every irreducible reversible discrete-time chain on n states.


\entrylabel{Scope}

The comparison is the nonlazy complete graph. Feige asymptotic lower bounds for unweighted graphs do not establish this exact weighted-chain statement.


\entrylabel{Dated literature check}

Focused primary-source searches on2026-09-06;\allowbreak{} no current-open or novelty certification. See source\_\allowbreak{}question\_\allowbreak{}registry\_\allowbreak{}markov\_\allowbreak{}002.json.


\subsection*{P0245-A001 | Exact minimum stationary cover time for reversible chains structural calculation}

\entrylabel{Argument or result}

The assertion holds for two states. Write P(0,\allowbreak{}1)=\allowbreak{}a and P(1,\allowbreak{}0)=\allowbreak{}b,\allowbreak{}0<a,\allowbreak{}b<=\allowbreak{}1. Stationarity gives pi=\allowbreak{}(b,\allowbreak{}a)/\allowbreak{}(a+\allowbreak{}b),\allowbreak{} and E\_\allowbreak{}pi C=\allowbreak{}b/\allowbreak{}[a(a+\allowbreak{}b)]+\allowbreak{}a/\allowbreak{}[b(a+\allowbreak{}b)]=\allowbreak{}(a\textasciicircum{}2+\allowbreak{}b\textasciicircum{}2)/\allowbreak{}[ab(a+\allowbreak{}b)]>=\allowbreak{}1,\allowbreak{} since a\textasciicircum{}2+\allowbreak{}b\textasciicircum{}2>=\allowbreak{}2ab and a+\allowbreak{}b<=\allowbreak{}2. Equality requires a=\allowbreak{}b=\allowbreak{}1.


\entrylabel{Second check}

The initial state is already visited at time0,\allowbreak{} so the two-state cover time is exactly the first hit of the other state. This fixes the common off-by-one and laziness errors. A deterministic directed cycle would violate the general bound for large n but is nonreversible and is not a counterexample.


\entrylabel{Remaining obstruction}

A many-state coupon-collector argument needs to control the distribution of the next unvisited target;\allowbreak{} its selection is biased by the history. Individual stationary hitting-time inequalities cannot simply be summed in random discovery order.


Verification:\allowbreak{} author checked limited argument | Independent review:\allowbreak{} none | Full target solved:\allowbreak{} False


\clearpage

\section*{P0246  Universal leading cover-time upper bound under transitivity}

\entrylabel{Primary source}

Aldous-Fill:\allowbreak{} Reversible Markov Chains and Random Walks on Graphs,\allowbreak{} 2002 recompiled2014;\allowbreak{} Open Problem7.6,\allowbreak{}PDF243;\allowbreak{} Corollary7.5


\url{https://www.stat.berkeley.edu/~aldous/RWG/book.pdf}


Source statement date:\allowbreak{} 1994-01-31. Archived access:\allowbreak{} 2026-09-06T22:\allowbreak{}03:\allowbreak{}09Z.


Historical:\allowbreak{} source\_\allowbreak{}says\_\allowbreak{}open | Current:\allowbreak{} unverified | Result:\allowbreak{} attempted\_\allowbreak{}unsolved


\entrylabel{Question}

For every growing sequence of symmetric reversible n-state chains,\allowbreak{} is E C<=\allowbreak{}(1+\allowbreak{}o(1))*tau0*log n without a gap-to-hitting assumption?


\entrylabel{Scope}

The source already proves the bound when tau2/\allowbreak{}tau0 tends to0. That established regime is only a partial result,\allowbreak{} not a new solution.


\entrylabel{Dated literature check}

Focused primary-source searches on2026-09-06;\allowbreak{} no current-open or novelty certification. See source\_\allowbreak{}question\_\allowbreak{}registry\_\allowbreak{}markov\_\allowbreak{}002.json.


\subsection*{P0246-A001 | Universal leading cover-time upper bound under transitivity structural calculation}

\entrylabel{Argument or result}

For continuous-time rate-one random walk on the cycle,\allowbreak{} tau0=\allowbreak{}(n\textasciicircum{}2-1)/\allowbreak{}6 while E C=\allowbreak{}n(n-1)/\allowbreak{}2. To see the latter,\allowbreak{} when exactly k consecutive vertices have been visited the walker is at an endpoint,\allowbreak{} and the mean time to enlarge the range is k,\allowbreak{} by gambler-ruin exit-time calculation on a path of length k+\allowbreak{}1. Summing k=\allowbreak{}1,\allowbreak{}...,\allowbreak{}n-1 gives E C. Hence E C/\allowbreak{}(tau0*log n) is asymptotic to3/\allowbreak{}log n and tends to0 even though tau2/\allowbreak{}tau0 stays positive.


\entrylabel{Second check}

Poissonizing the jump chain preserves the mean duration because the holding means are one. The cycle explicitly shows that failure of the source sufficient condition does not imply failure of the desired upper bound.


\entrylabel{Remaining obstruction}

One needs a uniform dichotomy or bound between chains with small tau2/\allowbreak{}tau0 and cycle-like slower geometries. The general Matthews bound still loses a factor2 through maximal versus average hitting time.


Verification:\allowbreak{} author checked limited argument | Independent review:\allowbreak{} none | Full target solved:\allowbreak{} False


\clearpage

\section*{P0247  Asymptotic uniformity of the final discovered vertex}

\entrylabel{Primary source}

Aldous-Fill:\allowbreak{} Reversible Markov Chains and Random Walks on Graphs,\allowbreak{} 2002 recompiled2014;\allowbreak{} Open Problem7.10,\allowbreak{}PDF246


\url{https://www.stat.berkeley.edu/~aldous/RWG/book.pdf}


Source statement date:\allowbreak{} 1994-01-31. Archived access:\allowbreak{} 2026-09-06T22:\allowbreak{}03:\allowbreak{}09Z.


Historical:\allowbreak{} source\_\allowbreak{}says\_\allowbreak{}open | Current:\allowbreak{} unverified | Result:\allowbreak{} attempted\_\allowbreak{}unsolved


\entrylabel{Question}

For simple random walk from a fixed vertex on growing finite vertex-transitive graphs,\allowbreak{} does the law of the last new vertex approach the uniform law in total variation?


\entrylabel{Scope}

Uniform stationary initialization makes the marginal exactly uniform by symmetry,\allowbreak{} but that is weaker than the fixed-start question.


\entrylabel{Dated literature check}

Focused primary-source searches on2026-09-06;\allowbreak{} no current-open or novelty certification. See source\_\allowbreak{}question\_\allowbreak{}registry\_\allowbreak{}markov\_\allowbreak{}002.json.


\subsection*{P0247-A001 | Asymptotic uniformity of the final discovered vertex structural calculation}

\entrylabel{Argument or result}

Let K(x,\allowbreak{}y) be the probability that y is last when the walk starts at x. Automorphism equivariance and transitivity imply every row sums to1 and every column has the same sum,\allowbreak{} hence K is doubly stochastic and pi*K=\allowbreak{}pi. This proves exact uniformity after a uniform start. However K can be far from the constant matrix:\allowbreak{} on the complete graph,\allowbreak{} K(x,\allowbreak{}x)=\allowbreak{}0 and K(x,\allowbreak{}y)=\allowbreak{}1/\allowbreak{}(n-1) for y different from x,\allowbreak{} giving total-variation distance exactly1/\allowbreak{}n.


\entrylabel{Second check}

The complete-graph value follows by permutation symmetry of all vertices other than the start. For a generic doubly stochastic equivariant kernel,\allowbreak{} the identity kernel is a counterexample to uniform row convergence,\allowbreak{} so double stochasticity alone cannot solve the walk-specific question.


\entrylabel{Remaining obstruction}

A useful proof must exploit discovery dynamics or show decorrelation from the initial vertex. The symmetry calculation does not do this,\allowbreak{} and randomizing the initial vertex would erase the actual question.


Verification:\allowbreak{} author checked limited argument | Independent review:\allowbreak{} none | Full target solved:\allowbreak{} False


\clearpage

\section*{P0250  Kalai absorption bound from a directed cut parameter}

\entrylabel{Primary source}

Aldous-Fill:\allowbreak{} Reversible Markov Chains and Random Walks on Graphs,\allowbreak{} 2002 recompiled2014;\allowbreak{} Open Problem9.27,\allowbreak{}PDF332


\url{https://www.stat.berkeley.edu/~aldous/RWG/book.pdf}


Source statement date:\allowbreak{} 1995-04-21. Archived access:\allowbreak{} 2026-09-06T22:\allowbreak{}03:\allowbreak{}09Z.


Historical:\allowbreak{} source\_\allowbreak{}says\_\allowbreak{}open | Current:\allowbreak{} unverified | Result:\allowbreak{} attempted\_\allowbreak{}unsolved


\entrylabel{Question}

For a strictly downward chain on states1,\allowbreak{}...,\allowbreak{}n absorbed at1,\allowbreak{} let kappa=\allowbreak{}max\_\allowbreak{}(A excludes1)|A|*|Ac|/\allowbreak{}[n*sum\_\allowbreak{}(i in A,\allowbreak{}j outside A)P(i,\allowbreak{}j)]. Is max\_\allowbreak{}i E\_\allowbreak{}i T1 bounded by a polynomial in kappa*log n?


\entrylabel{Scope}

The chain is acyclic before absorption and is not reversible. Applying reversible Cheeger inequalities to its raw transition matrix is invalid.


\entrylabel{Dated literature check}

Focused primary-source searches on2026-09-06;\allowbreak{} no current-open or novelty certification. See source\_\allowbreak{}question\_\allowbreak{}registry\_\allowbreak{}markov\_\allowbreak{}002.json.


\subsection*{P0250-A001 | Kalai absorption bound from a directed cut parameter structural calculation}

\entrylabel{Argument or result}

For the chain choosing its next state uniformly from1,\allowbreak{}...,\allowbreak{}i-1,\allowbreak{} the exact recursion h\_\allowbreak{}i=\allowbreak{}1+\allowbreak{}(1/\allowbreak{}(i-1))*sum\_\allowbreak{}(j<i)h\_\allowbreak{}j gives h\_\allowbreak{}i=\allowbreak{}H\_\allowbreak{}(i-1),\allowbreak{} so logarithmic absorption is possible. For the deterministic predecessor chain,\allowbreak{} h\_\allowbreak{}n=\allowbreak{}n-1 and the cut A=\allowbreak{}\{floor(n/\allowbreak{}2)+\allowbreak{}1,\allowbreak{}...,\allowbreak{}n\} has total outflow1,\allowbreak{} forcing kappa>=\allowbreak{}floor(n/\allowbreak{}2)*ceil(n/\allowbreak{}2)/\allowbreak{}n,\allowbreak{} of order n. Thus this slow example is detected by the cut parameter.


\entrylabel{Second check}

The harmonic recursion follows by subtracting consecutive averaged recurrences:\allowbreak{} h\_\allowbreak{}(i+\allowbreak{}1)-h\_\allowbreak{}i=\allowbreak{}1/\allowbreak{}i. The predecessor cut has exactly one crossing transition,\allowbreak{} so no stationary-measure normalization is being silently inserted.


\entrylabel{Remaining obstruction}

The cut condition controls average outward probability over a set,\allowbreak{} while the law of a chain conditioned on a long descent can concentrate on atypical vertices within that set. A level-set or occupation-flow argument must bridge that mismatch.


Verification:\allowbreak{} author checked limited argument | Independent review:\allowbreak{} none | Full target solved:\allowbreak{} False


\clearpage

\section*{P0251  Mixing of the random-pair binary crossover chain}

\entrylabel{Primary source}

Aldous-Fill:\allowbreak{} Reversible Markov Chains and Random Walks on Graphs,\allowbreak{} 2002 recompiled2014;\allowbreak{} Open Problem12.5,\allowbreak{}PDF400-401


\url{https://www.stat.berkeley.edu/~aldous/RWG/book.pdf}


Source statement date:\allowbreak{} 1999-10-11. Archived access:\allowbreak{} 2026-09-06T22:\allowbreak{}03:\allowbreak{}09Z.


Historical:\allowbreak{} source\_\allowbreak{}says\_\allowbreak{}open | Current:\allowbreak{} unverified | Result:\allowbreak{} attempted\_\allowbreak{}unsolved


\entrylabel{Question}

For K even binary words of length L,\allowbreak{} repeatedly choose a uniform pairing and independently swap each paired coordinate with probability1/\allowbreak{}2. Prove the proposed O(log K*log L) mixing bound within each fixed column-sum class.


\entrylabel{Scope}

The source uses a family of words and a fresh uniform permutation to place pairs;\allowbreak{} shared pairings couple the columns. Replace logarithms by log(K+\allowbreak{}1),\allowbreak{}log(L+\allowbreak{}1) when stating small-parameter bounds. General evolutionary algorithms are different chains.


\entrylabel{Dated literature check}

Focused primary-source searches on2026-09-06;\allowbreak{} no current-open or novelty certification. See source\_\allowbreak{}question\_\allowbreak{}registry\_\allowbreak{}markov\_\allowbreak{}002.json.


\subsection*{P0251-A001 | Mixing of the random-pair binary crossover chain structural calculation}

\entrylabel{Argument or result}

The naive Hamming coupling in the source has exact one-column drift E(W\_\allowbreak{}next|W=\allowbreak{}w)=\allowbreak{}w-w\textasciicircum{}2/\allowbreak{}[2(K-1)]. At the last possible disagreement w=\allowbreak{}2,\allowbreak{} the chance of removing it is1/\allowbreak{}(K-1),\allowbreak{} so this coupling itself requires order K further steps from that state. It therefore cannot prove logarithmic K mixing simply by improving its large-w drift estimate.


\entrylabel{Second check}

When w=\allowbreak{}2 there is one opposite mismatch of each sign. They coalesce only if paired,\allowbreak{} probability1/\allowbreak{}(K-1),\allowbreak{} and then the prescribed opposite-swap coupling agrees forever. Waiting is geometric. This is a limitation of the coupling,\allowbreak{} not a mixing-time lower bound.


\entrylabel{Remaining obstruction}

One needs a different coupling,\allowbreak{} entropy estimate or representation argument which does not insist on coordinatewise agreement at every generation. Marginal one-column mixing also does not automatically control joint mixing because pairings are shared across all columns.


Verification:\allowbreak{} author checked limited argument | Independent review:\allowbreak{} none | Full target solved:\allowbreak{} False


\clearpage

\section*{P0252  Sharp harmonic-mean characterization of meeting times}

\entrylabel{Primary source}

Aldous-Fill:\allowbreak{} Reversible Markov Chains and Random Walks on Graphs,\allowbreak{} 2002 recompiled2014;\allowbreak{} Open Problem14.8,\allowbreak{}PDF463


\url{https://www.stat.berkeley.edu/~aldous/RWG/book.pdf}


Source statement date:\allowbreak{} 1994-03-10. Archived access:\allowbreak{} 2026-09-06T22:\allowbreak{}03:\allowbreak{}09Z.


Historical:\allowbreak{} source\_\allowbreak{}says\_\allowbreak{}open | Current:\allowbreak{} unverified | Result:\allowbreak{} attempted\_\allowbreak{}unsolved


\entrylabel{Question}

Is the maximal expected meeting time of two independent copies of any finite reversible continuous-time chain comparable to J=\allowbreak{}[sum\_\allowbreak{}i pi\_\allowbreak{}i/\allowbreak{}max(E\_\allowbreak{}pi T\_\allowbreak{}i,\allowbreak{}tau1)]\textasciicircum{}-1 by universal constants?


\entrylabel{Scope}

The source proves the upper comparison;\allowbreak{} the question is the converse. Its tau1 is the pairwise total-variation e\textasciicircum{}-1 threshold.


\entrylabel{Dated literature check}

Focused primary-source searches on2026-09-06;\allowbreak{} no current-open or novelty certification. See source\_\allowbreak{}question\_\allowbreak{}registry\_\allowbreak{}markov\_\allowbreak{}002.json.


\subsection*{P0252-A001 | Sharp harmonic-mean characterization of meeting times structural calculation}

\entrylabel{Argument or result}

For a two-state chain with rates a<=\allowbreak{}b,\allowbreak{} tau1=\allowbreak{}tau2=\allowbreak{}1/\allowbreak{}(a+\allowbreak{}b),\allowbreak{} and the only distinct pair meets after an exponential time of mean tau2. Direct stationary hitting calculations give J/\allowbreak{}tau2=\allowbreak{}b*(a+\allowbreak{}b)/\allowbreak{}(a\textasciicircum{}2+\allowbreak{}b\textasciicircum{}2)=\allowbreak{}(1+\allowbreak{}r)/\allowbreak{}(1+\allowbreak{}r\textasciicircum{}2),\allowbreak{} r=\allowbreak{}a/\allowbreak{}b. This ratio ranges up to(1+\allowbreak{}sqrt2)/\allowbreak{}2,\allowbreak{} attained at r=\allowbreak{}sqrt2-1. Thus the converse cannot hold with constant1 even for two states,\allowbreak{} although a universal larger constant is consistent with this control.


\entrylabel{Second check}

E\_\allowbreak{}pi T0=\allowbreak{}a/\allowbreak{}[b(a+\allowbreak{}b)] and E\_\allowbreak{}pi T1=\allowbreak{}b/\allowbreak{}[a(a+\allowbreak{}b)]. Substituting the two maxima into the harmonic mean gives the formula;\allowbreak{} differentiating(1+\allowbreak{}r)/\allowbreak{}(1+\allowbreak{}r\textasciicircum{}2) gives its exact maximum. Meeting occurs at the first jump of either distinct coordinate,\allowbreak{} total rate a+\allowbreak{}b.


\entrylabel{Remaining obstruction}

The two-state calculation does not bound J/\allowbreak{}max E M on arbitrary chains. Rare spatial bottlenecks can affect stationary hitting times and two-particle meeting differently.


Verification:\allowbreak{} author checked limited argument | Independent review:\allowbreak{} none | Full target solved:\allowbreak{} False


\clearpage

\section*{P0253  Maximal stationary antivoter variance}

\entrylabel{Primary source}

Aldous-Fill:\allowbreak{} Reversible Markov Chains and Random Walks on Graphs,\allowbreak{} 2002 recompiled2014;\allowbreak{} Unnumbered question before Corollary14.24,\allowbreak{}PDF478


\url{https://www.stat.berkeley.edu/~aldous/RWG/book.pdf}


Source statement date:\allowbreak{} 1994-03-10. Archived access:\allowbreak{} 2026-09-06T22:\allowbreak{}03:\allowbreak{}09Z.


Historical:\allowbreak{} source\_\allowbreak{}says\_\allowbreak{}open | Current:\allowbreak{} unverified | Result:\allowbreak{} attempted\_\allowbreak{}unsolved


\entrylabel{Question}

Among admissible regular n-vertex graphs,\allowbreak{} is the stationary variance of the sum of antivoter spins maximized by the complete graph?


\entrylabel{Scope}

Use the source irreducibility convention for nonbipartite noncycle graphs. The separate request for sufficient normal-limit conditions was answered in later Stein-method literature and is not counted here.


\entrylabel{Dated literature check}

Focused primary-source searches on2026-09-06;\allowbreak{} no current-open or novelty certification. See source\_\allowbreak{}question\_\allowbreak{}registry\_\allowbreak{}markov\_\allowbreak{}002.json.


\subsection*{P0253-A001 | Maximal stationary antivoter variance structural calculation}

\entrylabel{Argument or result}

On the complete graph all distinct-pair correlations equal c. At stationarity,\allowbreak{} the generator applied to eta\_\allowbreak{}i*eta\_\allowbreak{}j gives c=\allowbreak{}-(1+\allowbreak{}(n-2)c)/\allowbreak{}(n-1),\allowbreak{} hence c=\allowbreak{}-1/\allowbreak{}(2n-3). Therefore Var(S)=\allowbreak{}n+\allowbreak{}n(n-1)c=\allowbreak{}n(n-2)/\allowbreak{}(2n-3). The general source inequality Var(S)<n/\allowbreak{}2 has the same leading term but leaves the exact gap n/\allowbreak{}[2(2n-3)].


\entrylabel{Second check}

The correlation equation can also be obtained by conditioning on which member of the pair updates and whether it samples its partner. Each route gives the same c. In particular merely proving Var(S)<=\allowbreak{}n/\allowbreak{}2 cannot resolve the proposed finite-n optimizer.


\entrylabel{Remaining obstruction}

A sharp bound on the return-parity expression for the meeting chain,\allowbreak{} or on averaged spin correlations,\allowbreak{} is missing. Pairwise correlations need not all be bounded by the complete-graph value.


Verification:\allowbreak{} author checked limited argument | Independent review:\allowbreak{} none | Full target solved:\allowbreak{} False


\clearpage

\section*{P0254  Mixing of gap-jump exclusion on a discrete circle}

\entrylabel{Primary source}

Aldous-Fill:\allowbreak{} Reversible Markov Chains and Random Walks on Graphs,\allowbreak{} 2002 recompiled2014;\allowbreak{} Example14.32 conjectures,\allowbreak{}PDF485-486


\url{https://www.stat.berkeley.edu/~aldous/RWG/book.pdf}


Source statement date:\allowbreak{} 1994-03-10. Archived access:\allowbreak{} 2026-09-06T22:\allowbreak{}03:\allowbreak{}09Z.


Historical:\allowbreak{} source\_\allowbreak{}says\_\allowbreak{}open | Current:\allowbreak{} unverified | Result:\allowbreak{} attempted\_\allowbreak{}unsolved


\entrylabel{Question}

For m particles on K circular sites,\allowbreak{} choose a vacant site and a direction by the source per-step rule,\allowbreak{} then move the first encountered particle to the site. Prove the conjectured tau2=\allowbreak{}Theta(m\textasciicircum{}3) and tau1=\allowbreak{}O(m\textasciicircum{}3 log m) uniformly in K>m.


\entrylabel{Scope}

The step first chooses a site uniformly from all K,\allowbreak{} doing nothing if occupied. The continuum-circle limit keeps m fixed before K grows;\allowbreak{} both spectral and mixing targets are grouped.


\entrylabel{Dated literature check}

Focused primary-source searches on2026-09-06;\allowbreak{} no current-open or novelty certification. See source\_\allowbreak{}question\_\allowbreak{}registry\_\allowbreak{}markov\_\allowbreak{}002.json.


\subsection*{P0254-A001 | Mixing of gap-jump exclusion on a discrete circle structural calculation}

\entrylabel{Argument or result}

For m=\allowbreak{}1 the spectrum is exactly accessible. From occupied position x,\allowbreak{} any chosen different site y becomes the new position regardless of direction,\allowbreak{} and choosing x leaves it there. Thus every step produces the uniform law on K sites:\allowbreak{} P(x,\allowbreak{}y)=\allowbreak{}1/\allowbreak{}K. The nonconstant eigenvalues are0,\allowbreak{} tau2=\allowbreak{}1 and the discrete mixing time is one. This confirms K-independence in the simplest sector but contains no information about the m\textasciicircum{}3 growth.


\entrylabel{Second check}

The direction contributes no factor1/\allowbreak{}2 because there is only one particle and either direction encounters that same particle. Treating the rule as nearest-neighbor exclusion would incorrectly predict a K\textasciicircum{}2 relaxation time.


\entrylabel{Remaining obstruction}

With several particles the moved particle is chosen by the entire gap geometry. The obvious common-site common-direction coupling can retain a last disagreement for a long time,\allowbreak{} and no uniform sharp coalescence estimate has been obtained.


Verification:\allowbreak{} author checked limited argument | Independent review:\allowbreak{} none | Full target solved:\allowbreak{} False


\clearpage

\section*{P0255  Variation of nearest-unvisited walk length under a random start}

\entrylabel{Primary source}

Aldous:\allowbreak{} The Nearest Unvisited Vertex Walk on Random Graphs,\allowbreak{}internal date26August2021;\allowbreak{} Conjecture5 and stronger Open Problem4,\allowbreak{}PDF8


\url{https://arxiv.org/pdf/1912.13175}


Source statement date:\allowbreak{} 2021-08-26. Archived access:\allowbreak{} 2026-09-06T23:\allowbreak{}13:\allowbreak{}46Z.


Historical:\allowbreak{} source\_\allowbreak{}says\_\allowbreak{}open | Current:\allowbreak{} unverified | Result:\allowbreak{} attempted\_\allowbreak{}unsolved


\entrylabel{Question}

For every finite positive-length graph G and a uniform initial vertex V,\allowbreak{} is sd(L\_\allowbreak{}NUV(G,\allowbreak{}V))/\allowbreak{}E L\_\allowbreak{}NUV(G,\allowbreak{}V) bounded by an absolute constant?


\entrylabel{Scope}

The deterministic max/\allowbreak{}min question is recorded as a stronger companion within this one probability entry. The arXiv wrapper says March2021 while the internal manuscript date is August26,\allowbreak{}2021;\allowbreak{} both are preserved.


\entrylabel{Dated literature check}

Focused primary-source searches on2026-09-06;\allowbreak{} no current-open or novelty certification. See source\_\allowbreak{}question\_\allowbreak{}registry\_\allowbreak{}markov\_\allowbreak{}002.json.


\subsection*{P0255-A001 | Variation of nearest-unvisited walk length under a random start structural calculation}

\entrylabel{Argument or result}

Let m and M be the minimum and maximum start-dependent lengths,\allowbreak{} R=\allowbreak{}M/\allowbreak{}m. The elementary bound Var(X)<=\allowbreak{}(M-E X)*(E X-m) implies sd(X)/\allowbreak{}E X<=\allowbreak{}(R-1)/\allowbreak{}(2*sqrt(R)) for any X in[m,\allowbreak{}M]. Classical nearest-neighbor-tour comparison gives R=\allowbreak{}O(log n),\allowbreak{} since every walk length is at least MST length and its length is at most O(log n) times a TSP tour,\allowbreak{} at most twice MST length. Hence a rigorous O(sqrt(log n)) coefficient-of-variation bound follows,\allowbreak{} short of an absolute constant.


\entrylabel{Second check}

The variance inequality follows by averaging(X-m)*(M-X)>=\allowbreak{}0. Maximizing(M-u)*(u-m)/\allowbreak{}u\textasciicircum{}2 over u in[m,\allowbreak{}M] gives u=\allowbreak{}2Mm/\allowbreak{}(M+\allowbreak{}m) and the displayed exact ratio. This is sharper than simply bounding standard deviation by the range.


\entrylabel{Remaining obstruction}

The deterministic logarithmic approximation examples do not automatically give a large random-start coefficient of variation. One must also control how many starting vertices exhibit the long route.


Verification:\allowbreak{} author checked limited argument | Independent review:\allowbreak{} none | Full target solved:\allowbreak{} False


\clearpage

\section*{P0256  Weak concentration of nearest-unvisited walks under exponential edge noise}

\entrylabel{Primary source}

Aldous:\allowbreak{} The Nearest Unvisited Vertex Walk on Random Graphs,\allowbreak{}internal date26August2021;\allowbreak{} Section2.5,\allowbreak{}PDF8-9


\url{https://arxiv.org/pdf/1912.13175}


Source statement date:\allowbreak{} 2021-08-26. Archived access:\allowbreak{} 2026-09-06T23:\allowbreak{}13:\allowbreak{}46Z.


Historical:\allowbreak{} source\_\allowbreak{}says\_\allowbreak{}open | Current:\allowbreak{} unverified | Result:\allowbreak{} attempted\_\allowbreak{}unsolved


\entrylabel{Question}

Find and prove the proposed analogue of FPP weak concentration for NUV walk length when deterministic edge means are replaced by independent exponential lengths and no individual edge contributes a macroscopic share.


\entrylabel{Scope}

The source asks for an analogue and does not fix a complete sufficient-condition theorem. Any proposed maximal-edge condition must be stated and proved,\allowbreak{} not attributed to the source as an established exact conjecture.


\entrylabel{Dated literature check}

Focused primary-source searches on2026-09-06;\allowbreak{} no current-open or novelty certification. See source\_\allowbreak{}question\_\allowbreak{}registry\_\allowbreak{}markov\_\allowbreak{}002.json.


\subsection*{P0256-A001 | Weak concentration of nearest-unvisited walks under exponential edge noise structural calculation}

\entrylabel{Argument or result}

A standard bounded-difference proof faces a real discontinuity. On the line with vertices A=\allowbreak{}0,\allowbreak{}B=\allowbreak{}1,\allowbreak{}D=\allowbreak{}3,\allowbreak{}C=\allowbreak{}-1-epsilon,\allowbreak{} start at A. For epsilon>0 small,\allowbreak{} the NUV order is A,\allowbreak{}B,\allowbreak{}D,\allowbreak{}C and length7+\allowbreak{}epsilon. With C=\allowbreak{}-1+\allowbreak{}epsilon,\allowbreak{} the order is A,\allowbreak{}C,\allowbreak{}B,\allowbreak{}D and length5-2epsilon. An arbitrarily small point perturbation changes the route length by nearly2. Thus NUV length lacks the elementary geometric Lipschitz property available for optimized MST or TSP costs.


\entrylabel{Second check}

Both examples use true Euclidean shortest-path distances and strict choices for positive epsilon. The discontinuity comes from a tie limit,\allowbreak{} not a violation of the triangle inequality. It is a deterministic methodological obstruction,\allowbreak{} not a random-exponential concentration counterexample.


\entrylabel{Remaining obstruction}

One needs probabilistic control of near-tie cascades or stabilization,\allowbreak{} rather than a uniform deterministic one-coordinate Lipschitz estimate. The no-macroscopic-edge hypothesis alone has not been shown sufficient here.


Verification:\allowbreak{} author checked limited argument | Independent review:\allowbreak{} none | Full target solved:\allowbreak{} False


\clearpage

\section*{P0257  Law of large numbers and constants for NUV lengths}

\entrylabel{Primary source}

Aldous:\allowbreak{} The Nearest Unvisited Vertex Walk on Random Graphs,\allowbreak{}internal date26August2021;\allowbreak{} Equation4,\allowbreak{}PDF9;\allowbreak{} Section5,\allowbreak{}PDF20;\allowbreak{} companion talk2022p23


\url{https://arxiv.org/pdf/1912.13175}


Source statement date:\allowbreak{} 2021-08-26. Archived access:\allowbreak{} 2026-09-06T23:\allowbreak{}13:\allowbreak{}46Z.


Historical:\allowbreak{} source\_\allowbreak{}says\_\allowbreak{}open | Current:\allowbreak{} unverified | Result:\allowbreak{} attempted\_\allowbreak{}unsolved


\entrylabel{Question}

Prove existence and identify the linear-scale NUV length constants in the source Euclidean,\allowbreak{} iid-weighted-grid and complete-graph mean-field models.


\entrylabel{Scope}

The three model versions are conservatively grouped as one first-order limit programme. Tightness and order bounds proved in the source are not counted as solutions.


\entrylabel{Dated literature check}

Focused primary-source searches on2026-09-06;\allowbreak{} no current-open or novelty certification. See source\_\allowbreak{}question\_\allowbreak{}registry\_\allowbreak{}markov\_\allowbreak{}002.json.


\subsection*{P0257-A001 | Law of large numbers and constants for NUV lengths structural calculation}

\entrylabel{Argument or result}

An elementary geometric lower bound comes from nearest-neighbor distances d\_\allowbreak{}v:\allowbreak{} the ordered discovery walk has a departure from every vertex except its final one,\allowbreak{} and every such segment is at least d\_\allowbreak{}v. Hence L\_\allowbreak{}NUV>=\allowbreak{}sum\_\allowbreak{}v d\_\allowbreak{}v-max\_\allowbreak{}v d\_\allowbreak{}v. For unit-intensity planar Poisson points,\allowbreak{} Palm P(d\_\allowbreak{}v>r)=\allowbreak{}exp(-pi*r\textasciicircum{}2),\allowbreak{} so E d\_\allowbreak{}v=\allowbreak{}1/\allowbreak{}2. In the growing-square binomial version,\allowbreak{} the bulk expectation tends to1/\allowbreak{}2 and E max d\_\allowbreak{}v=\allowbreak{}O(sqrt(log n)),\allowbreak{} yielding liminf E L\_\allowbreak{}NUV/\allowbreak{}n>=\allowbreak{}1/\allowbreak{}2. In the complete graph with independent mean-n exponential edge lengths,\allowbreak{} the analogous bound gives liminf E L\_\allowbreak{}NUV/\allowbreak{}n>=\allowbreak{}1.


\entrylabel{Second check}

The mean-field minimum at each vertex has mean n/\allowbreak{}(n-1),\allowbreak{} and a union bound gives expected maximum O(log n);\allowbreak{} independence between vertices is unnecessary for these expectation bounds. The Euclidean boundary affects only a vanishing bulk fraction for the lower bound,\allowbreak{} and the maximum estimate follows from empty-ball bounds.


\entrylabel{Remaining obstruction}

These are lower constants only and may be weaker than established optimized-tour bounds. No subadditivity or stabilization argument proves convergence of the greedy length,\allowbreak{} and the source finite-size simulation constants are not theorem values.


Verification:\allowbreak{} author checked limited argument | Independent review:\allowbreak{} none | Full target solved:\allowbreak{} False


\clearpage

\section*{P0258  Variance scale and Gaussian fluctuations of NUV lengths}

\entrylabel{Primary source}

Aldous:\allowbreak{} The Nearest Unvisited Vertex Walk on Random Graphs,\allowbreak{}internal date26August2021;\allowbreak{} Equation4 and discussion,\allowbreak{}PDF9-10 and21


\url{https://arxiv.org/pdf/1912.13175}


Source statement date:\allowbreak{} 2021-08-26. Archived access:\allowbreak{} 2026-09-06T23:\allowbreak{}13:\allowbreak{}46Z.


Historical:\allowbreak{} source\_\allowbreak{}says\_\allowbreak{}open | Current:\allowbreak{} unverified | Result:\allowbreak{} attempted\_\allowbreak{}unsolved


\entrylabel{Question}

Determine the variance order and limiting centered distribution of NUV lengths in the source random spatial and mean-field models,\allowbreak{} including the proposed linear variance and Gaussian limit.


\entrylabel{Scope}

All variance and central-limit refinements are grouped. The source itself notes that small simulations cast doubt on a universal linear lower bound.


\entrylabel{Dated literature check}

Focused primary-source searches on2026-09-06;\allowbreak{} no current-open or novelty certification. See source\_\allowbreak{}question\_\allowbreak{}registry\_\allowbreak{}markov\_\allowbreak{}002.json.


\subsection*{P0258-A001 | Variance scale and Gaussian fluctuations of NUV lengths structural calculation}

\entrylabel{Argument or result}

The variance decomposes exactly as Var(L)=\allowbreak{}Var(E[L|G])+\allowbreak{}E Var(L|G),\allowbreak{} separating graph disorder from random-start disorder. This implies a lower bound from either term but gives no additivity across mandatory local edges. Even if two vertices force an edge into every NUV walk,\allowbreak{} its random contribution can covary negatively with the remainder;\allowbreak{} the scalar example U+\allowbreak{}(1-U)=\allowbreak{}1 shows why mandatory participation alone cannot prove a variance lower bound.


\entrylabel{Second check}

The law of total variance is applied with the entire weighted graph fixed and the starting vertex random. Conditional independence of local edge contributions is not asserted:\allowbreak{} the chosen trajectory depends on the full configuration.


\entrylabel{Remaining obstruction}

A lower bound needs orthogonal coordinate projections or a conditional variance decomposition controlling cancellations. A Gaussian limit further needs a Lindeberg or stabilization estimate for the actual greedy functional;\allowbreak{} neither has been proved.


Verification:\allowbreak{} author checked limited argument | Independent review:\allowbreak{} none | Full target solved:\allowbreak{} False


\clearpage

\section*{P0259  Never-visited density for the planar Poisson greedy walk}

\entrylabel{Primary source}

Aldous:\allowbreak{} The Nearest Unvisited Vertex Walk on Random Graphs,\allowbreak{}internal date26August2021;\allowbreak{} Section5 local weak convergence discussion,\allowbreak{}PDF20


\url{https://arxiv.org/pdf/1912.13175}


Source statement date:\allowbreak{} 2021-08-26. Archived access:\allowbreak{} 2026-09-06T23:\allowbreak{}13:\allowbreak{}46Z.


Historical:\allowbreak{} source\_\allowbreak{}says\_\allowbreak{}open | Current:\allowbreak{} unverified | Result:\allowbreak{} attempted\_\allowbreak{}unsolved


\entrylabel{Question}

For the nearest-unvisited greedy walk on a planar Poisson point process,\allowbreak{} estimate the number of never-visited points in a radius-r ball as r tends to infinity.


\entrylabel{Scope}

The source attributes this to a longstanding folklore problem and cites Bordenave-Foss-Last. A finite-graph NUV cover is automatic and does not answer the infinite-plane question.


\entrylabel{Dated literature check}

Focused primary-source searches on2026-09-06;\allowbreak{} no current-open or novelty certification. See source\_\allowbreak{}question\_\allowbreak{}registry\_\allowbreak{}markov\_\allowbreak{}002.json.


\subsection*{P0259-A001 | Never-visited density for the planar Poisson greedy walk structural calculation}

\entrylabel{Argument or result}

Local finiteness implies a simple but weak path fact:\allowbreak{} if the infinite greedy walk is well-defined for all steps,\allowbreak{} then it eventually leaves every bounded ball,\allowbreak{} since it never revisits a chosen point and a bounded ball contains finitely many Poisson points. This does not imply it visited all those finitely many points before leaving. A deterministic locally finite configuration can have a path going to infinity while leaving an entire bounded subset untouched.


\entrylabel{Second check}

The distinction is between eventual departure and exhaustive visitation. The trivial upper bound on never-visited points is the Poisson count,\allowbreak{} asymptotic pi*r\textasciicircum{}2 almost surely,\allowbreak{} but no matching order or limiting density follows from self-avoidance.


\entrylabel{Remaining obstruction}

The finite-window path law is sensitive to points outside the window. Identifying a translation-invariant limit of disjoint doubly infinite greedy paths,\allowbreak{} or quantitative exploration of missed points,\allowbreak{} remains the main obstruction.


Verification:\allowbreak{} author checked limited argument | Independent review:\allowbreak{} none | Full target solved:\allowbreak{} False


\clearpage

\section*{P0266  Critical variational SPDE maximal second-moment bound}

\entrylabel{Primary source}

Agresti-Veraar:\allowbreak{} Nonlinear SPDEs and maximal regularity,\allowbreak{}published NoDEA2025;\allowbreak{} Problem4,\allowbreak{}journal page130/\allowbreak{}150;\allowbreak{}PDF131;\allowbreak{}arXivPDF92


\url{https://repository.tudelft.nl/file/File_783e6459-5d3b-4957-bcd3-64da5b8d9025}


Source statement date:\allowbreak{} 2025-09-01. Archived access:\allowbreak{} 2026-09-06T23:\allowbreak{}44:\allowbreak{}20Z.


Historical:\allowbreak{} source\_\allowbreak{}says\_\allowbreak{}open | Current:\allowbreak{} unverified | Result:\allowbreak{} candidate\_\allowbreak{}full\_\allowbreak{}solution


\entrylabel{Question}

Does Theorem6.4 imply E sup\_\allowbreak{}\{t<=\allowbreak{}T\}||u(t)||\_\allowbreak{}H\textasciicircum{}2 <=\allowbreak{} C\_\allowbreak{}T(1+\allowbreak{}E||u0||\_\allowbreak{}H\textasciicircum{}2),\allowbreak{} including the source critical Allen-Cahn model gamma\textasciicircum{}2=\allowbreak{}2?


\entrylabel{Scope}

Published source explicitly includes critical quadratic scalar noise. Current-open status and originality remain unverified;\allowbreak{} candidate counterexample is unpromoted.


\subsection*{P0266-A001 | Spatially constant reduction;\allowbreak{} reciprocal BES4;\allowbreak{} finite-horizon harmonic stopping and Gaussian integral}

\entrylabel{Argument or result}

For deterministic constant initial condition a>0,\allowbreak{} the critical model reduces to U\_\allowbreak{}t=\allowbreak{}e\textasciicircum{}t/\allowbreak{}(sqrt(2)R\_\allowbreak{}\{(e\textasciicircum{}(2t)-1)/\allowbreak{}2\}) with R a BES4 process. An exact finite-horizon hitting identity gives P(sup(2R\textasciicircum{}2)\textasciicircum{}(-1)>=\allowbreak{}x)\textasciitilde{}c/\allowbreak{}x with c>0,\allowbreak{} hence E sup U\_\allowbreak{}t\textasciicircum{}2=\allowbreak{}infinity for every T>0. The Gelfand-triple and torus assumptions are checked in the full proof.


\entrylabel{Second check}

Second route checks the entire abstract theorem in the scalar Gelfand triple and separately checks the torus coercivity and critical interpolation estimates. Gaussian expectation is derived by a positive Laplace integral. Author checks only,\allowbreak{} not independent review.


\entrylabel{Remaining obstruction}

No gap found in author checks. Later-literature reconciliation and independent proof review remain necessary before promotion.


Verification:\allowbreak{} author checked candidate | Independent review:\allowbreak{} none | Full target solved:\allowbreak{} False


\entrylabel{Full proof packet}

proofs/\allowbreak{}P0266\_\allowbreak{}P0267\_\allowbreak{}critical\_\allowbreak{}allen\_\allowbreak{}cahn\_\allowbreak{}counterexamples\_\allowbreak{}candidate.md


\clearpage

\section*{P0267  Endpoint energy-space continuity for critical variational SPDEs}

\entrylabel{Primary source}

Agresti-Veraar:\allowbreak{} Nonlinear SPDEs and maximal regularity,\allowbreak{}published NoDEA2025;\allowbreak{} Problem5,\allowbreak{}journal page130/\allowbreak{}150;\allowbreak{}PDF131;\allowbreak{}arXivPDF92


\url{https://repository.tudelft.nl/file/File_783e6459-5d3b-4957-bcd3-64da5b8d9025}


Source statement date:\allowbreak{} 2025-09-01. Archived access:\allowbreak{} 2026-09-06T23:\allowbreak{}44:\allowbreak{}20Z.


Historical:\allowbreak{} source\_\allowbreak{}says\_\allowbreak{}open | Current:\allowbreak{} unverified | Result:\allowbreak{} candidate\_\allowbreak{}full\_\allowbreak{}solution


\entrylabel{Question}

Under Theorem6.4,\allowbreak{} does L2 convergence of initial conditions imply L2(Omega;\allowbreak{}L2(0,\allowbreak{}T;\allowbreak{}V) intersect C([0,\allowbreak{}T];\allowbreak{}H)) convergence of solutions?


\entrylabel{Scope}

Published source explicitly includes critical quadratic scalar noise. Current-open status and originality remain unverified;\allowbreak{} candidate counterexample is unpromoted.


\subsection*{P0267-A001 | Spatially constant reduction;\allowbreak{} reciprocal BES4;\allowbreak{} finite-horizon harmonic stopping and Gaussian integral}

\entrylabel{Argument or result}

Take the same critical Allen-Cahn equation with initial constants1/\allowbreak{}n converging to0 and a common Brownian driver. The zero-initial solution is0,\allowbreak{} whereas each positive-initial solution has infinite expected squared supremum on every positive time interval. Thus the claimed L2 convergence fails. These are different source properties with a common counterexample,\allowbreak{} counted once per numbered question.


\entrylabel{Second check}

Second route checks the entire abstract theorem in the scalar Gelfand triple and separately checks the torus coercivity and critical interpolation estimates. Gaussian expectation is derived by a positive Laplace integral. Author checks only,\allowbreak{} not independent review.


\entrylabel{Remaining obstruction}

No gap found in author checks. Later-literature reconciliation and independent proof review remain necessary before promotion.


Verification:\allowbreak{} author checked candidate | Independent review:\allowbreak{} none | Full target solved:\allowbreak{} False


\entrylabel{Full proof packet}

proofs/\allowbreak{}P0266\_\allowbreak{}P0267\_\allowbreak{}critical\_\allowbreak{}allen\_\allowbreak{}cahn\_\allowbreak{}counterexamples\_\allowbreak{}candidate.md


\clearpage

\section*{P0268  Abstract stochastic maximal regularity with non-small multiplicative noise}

\entrylabel{Primary source}

Agresti-Veraar:\allowbreak{} Nonlinear SPDEs and Maximal Regularity,\allowbreak{}2025 extended survey;\allowbreak{} Problem1,\allowbreak{}PDF91


\url{https://arxiv.org/pdf/2501.18561}


Source statement date:\allowbreak{} 2025-01-30. Archived access:\allowbreak{} 2026-09-06T23:\allowbreak{}40:\allowbreak{}06Z.


Historical:\allowbreak{} source\_\allowbreak{}says\_\allowbreak{}open | Current:\allowbreak{} unverified | Result:\allowbreak{} unsolved


\entrylabel{Question}

Find abstract operator conditions yielding stochastic maximal Lp regularity for (A,\allowbreak{}B),\allowbreak{} beyond the small-B perturbation theorem.


\entrylabel{Scope}

The source asks for a satisfactory structural criterion,\allowbreak{} rather than any sufficient condition:\allowbreak{} the small-norm sufficient condition is already established.


\subsection*{P0268-A001 | Explicit model calculation and hypothesis obstruction}

\entrylabel{Argument or result}

For simultaneously diagonal A and B on a Hilbert space,\allowbreak{} the homogeneous coordinate equation dX=\allowbreak{}-lambda*X dt+\allowbreak{}b*sqrt(lambda)\emph{X dW gives E|X\_\allowbreak{}t|\textasciicircum{}p=\allowbreak{}|X\_\allowbreak{}0|\textasciicircum{}p exp\{p}((p-1)*b\textasciicircum{}2/\allowbreak{}2-1)*lambda*t\}. Thus positivity of A alone cannot supply a p-independent stochastic parabolicity condition. In the p=\allowbreak{}2 coercive regime b\textasciicircum{}2<2,\allowbreak{} applying Ito to ||A\textasciicircum{}(1/\allowbreak{}2)u||\textasciicircum{}2 yields control of E integral ||Au||\textasciicircum{}2 for additive input in D(A\textasciicircum{}(1/\allowbreak{}2)),\allowbreak{} after absorbing its cross term by Young inequality.


\entrylabel{Second check}

The moment identity follows both from the explicit stochastic exponential and from Ito applied to |X|\textasciicircum{}p. The coercive calculation uses the weighted norm A\textasciicircum{}(1/\allowbreak{}2)u,\allowbreak{} so it controls the full domain D(A),\allowbreak{} not merely the energy domain.


\entrylabel{Remaining obstruction}

A criterion handling noncommuting A,\allowbreak{}B in Banach spaces must replace scalar coercivity and preserve the correct p-dependence. No such general criterion is proved here. Brzezniak-Veraar2012 already studies this p-dependence;\allowbreak{} the diagonal calculation is a known-model rederivation.


Verification:\allowbreak{} author checked partial | Independent review:\allowbreak{} none | Full target solved:\allowbreak{} False


\clearpage

\section*{P0269  Necessary and sufficient operator criterion for stochastic maximal regularity}

\entrylabel{Primary source}

Agresti-Veraar:\allowbreak{} Nonlinear SPDEs and Maximal Regularity,\allowbreak{}2025 extended survey;\allowbreak{} Problem2,\allowbreak{}PDF92


\url{https://arxiv.org/pdf/2501.18561}


Source statement date:\allowbreak{} 2025-01-30. Archived access:\allowbreak{} 2026-09-06T23:\allowbreak{}40:\allowbreak{}06Z.


Historical:\allowbreak{} source\_\allowbreak{}says\_\allowbreak{}open | Current:\allowbreak{} unverified | Result:\allowbreak{} unsolved


\entrylabel{Question}

For a sectorial operator A of angle below pi/\allowbreak{}2 on Lq,\allowbreak{} q>2,\allowbreak{} characterize stochastic maximal Lp regularity of (A,\allowbreak{}0).


\entrylabel{Scope}

A bounded H-infinity calculus is an established sufficient condition. Repeating that condition is not a solution of the necessary-and-sufficient problem.


\subsection*{P0269-A001 | Explicit model calculation and hypothesis obstruction}

\entrylabel{Argument or result}

In the positive diagonal Hilbert model u\_\allowbreak{}k(t)=\allowbreak{}integral\_\allowbreak{}0\textasciicircum{}t exp[-lambda\_\allowbreak{}k*(t-s)]g\_\allowbreak{}k(s)dW\_\allowbreak{}s,\allowbreak{} Ito isometry and Tonelli give E integral\_\allowbreak{}0\textasciicircum{}T sum\_\allowbreak{}k lambda\_\allowbreak{}k\textasciicircum{}2 |u\_\allowbreak{}k(t)|\textasciicircum{}2 dt <=\allowbreak{} (1/\allowbreak{}2) E integral\_\allowbreak{}0\textasciicircum{}T sum\_\allowbreak{}k lambda\_\allowbreak{}k |g\_\allowbreak{}k(s)|\textasciicircum{}2 ds. The exact kernel integral is lambda\_\allowbreak{}k*(1-exp[-2lambda\_\allowbreak{}k*(T-s)])/\allowbreak{}2. This isolates the quadratic square-function estimate behind the property.


\entrylabel{Second check}

This proof works first for finitely many coordinates and then passes by monotone convergence. Correlations caused by using the same Brownian motion do not enter a Hilbert norm squared,\allowbreak{} which is the sum of coordinate squares.


\entrylabel{Remaining obstruction}

For q>2,\allowbreak{} the spatial Lq norm lies outside the quadratic sum and the Hilbert Fubini computation no longer identifies the necessary operator square-function property. Establishing an intrinsic equivalence,\allowbreak{} rather than a sufficient functional-calculus hypothesis,\allowbreak{} remains open here.


Verification:\allowbreak{} author checked partial | Independent review:\allowbreak{} none | Full target solved:\allowbreak{} False


\clearpage

\section*{P0270  Critical time regularization at Hilbert exponent two}

\entrylabel{Primary source}

Agresti-Veraar:\allowbreak{} Nonlinear SPDEs and Maximal Regularity,\allowbreak{}2025 extended survey;\allowbreak{} Problem6,\allowbreak{}PDF92;\allowbreak{}Theorem5.7


\url{https://arxiv.org/pdf/2501.18561}


Source statement date:\allowbreak{} 2025-01-30. Archived access:\allowbreak{} 2026-09-06T23:\allowbreak{}40:\allowbreak{}06Z.


Historical:\allowbreak{} source\_\allowbreak{}says\_\allowbreak{}open | Current:\allowbreak{} unverified | Result:\allowbreak{} unsolved


\entrylabel{Question}

Does the time regularization theorem5.7 remain valid for p=\allowbreak{}2 with Hilbert X0 when a nonlinear criticality inequality is attained?


\entrylabel{Scope}

The question concerns the complete hypotheses and conclusion of theorem5.7;\allowbreak{} it is not the unrelated assertion that arbitrary additive noise gains time regularity in the energy domain.


\subsection*{P0270-A001 | Explicit model calculation and hypothesis obstruction}

\entrylabel{Argument or result}

A diagnostic linear calculation shows why a naive time-difference argument is insufficient. For du\_\allowbreak{}k+\allowbreak{}lambda\_\allowbreak{}k*u\_\allowbreak{}k dt=\allowbreak{}x\_\allowbreak{}k dW,\allowbreak{} a time increment includes fresh noise integral\_\allowbreak{}t\textasciicircum{}(t+\allowbreak{}h) exp[-lambda\_\allowbreak{}k*(t+\allowbreak{}h-s)]x\_\allowbreak{}k dW\_\allowbreak{}s. Its A\textasciicircum{}(1/\allowbreak{}2)-weighted second moment is (1/\allowbreak{}2)sum |x\_\allowbreak{}k|\textasciicircum{}2*(1-exp[-2lambda\_\allowbreak{}k*h]). Integrating h\textasciicircum{}(-1-2alpha) over h gives a term comparable to sum lambda\_\allowbreak{}k\textasciicircum{}(2alpha)|x\_\allowbreak{}k|\textasciicircum{}2. Thus even a constant-in-time coefficient needs additional spatial regularity to give W\textasciicircum{}(alpha,\allowbreak{}2) time regularity in D(A\textasciicircum{}(1/\allowbreak{}2)).


\entrylabel{Second check}

The change of variables z=\allowbreak{}lambda\_\allowbreak{}k*h gives the lambda\_\allowbreak{}k\textasciicircum{}(2alpha) factor,\allowbreak{} and the integral near zero is finite exactly when alpha<1/\allowbreak{}2. Agresti-Veraar2026,\allowbreak{} arXiv2509.07803,\allowbreak{} proves this precise obstruction for the linear problem.


\entrylabel{Remaining obstruction}

This known obstruction is not a counterexample to the full nonlinear bootstrap theorem,\allowbreak{} whose diffusion mapping assumptions may supply the missing spatial input. A proof must use those assumptions at the critical endpoint;\allowbreak{} interpolation with a positive exponent margin is unavailable.


Verification:\allowbreak{} author checked partial | Independent review:\allowbreak{} none | Full target solved:\allowbreak{} False


\clearpage

\section*{P0271  Global stochastic reaction diffusion under a triangular reaction structure}

\entrylabel{Primary source}

Agresti-Veraar:\allowbreak{} Nonlinear SPDEs and Maximal Regularity,\allowbreak{}2025 extended survey;\allowbreak{} Problem7,\allowbreak{}PDF93;\allowbreak{}system8.14


\url{https://arxiv.org/pdf/2501.18561}


Source statement date:\allowbreak{} 2025-01-30. Archived access:\allowbreak{} 2026-09-06T23:\allowbreak{}40:\allowbreak{}06Z.


Historical:\allowbreak{} source\_\allowbreak{}says\_\allowbreak{}open | Current:\allowbreak{} unverified | Result:\allowbreak{} unsolved


\entrylabel{Question}

Under the six positivity,\allowbreak{} polynomial-growth,\allowbreak{} mass-control and nonnegative lower-triangular reaction conditions in Problem7,\allowbreak{} prove global well-posedness with non-small regular transport noise. Include the source question of allowable superlinear reaction noise.


\entrylabel{Scope}

All linked growth extensions are grouped with the principal triangular-system question. The exact six conditions and stochastic parabolicity are retained by reference to the primary source.


\subsection*{P0271-A001 | Explicit model calculation and hypothesis obstruction}

\entrylabel{Argument or result}

In the spatially constant reduction with globally Lipschitz reaction noise,\allowbreak{} total mass S=\allowbreak{}sum\_\allowbreak{}i U\_\allowbreak{}i is nonnegative and its drift is at most K0+\allowbreak{}K1*S. On the positive orthant,\allowbreak{} |U|<=\allowbreak{}S,\allowbreak{} so the total diffusion coefficient satisfies ||sum\_\allowbreak{}i g\_\allowbreak{}i(U)||<=\allowbreak{}C*(1+\allowbreak{}S). Ito applied to S\textasciicircum{}2,\allowbreak{} stopping at levelN,\allowbreak{} gives E S\_\allowbreak{}(t wedge tauN)\textasciicircum{}2 <=\allowbreak{} C\_\allowbreak{}T*(1+\allowbreak{}S0\textasciicircum{}2). The standard stopped maximal estimate gives a uniform bound on E sup\_\allowbreak{}(s<=\allowbreak{}T) S\_\allowbreak{}(s wedge tauN)\textasciicircum{}2,\allowbreak{} hence the finite-dimensional system does not explode.


\entrylabel{Second check}

The diffusion bound follows from global Lipschitz continuity of each g\_\allowbreak{}i and positivity,\allowbreak{} rather than from cancellation between its noises. Markov inequality at the exit time then makes P(tauN<=\allowbreak{}T) tend to zero.


\entrylabel{Remaining obstruction}

For spatially varying solutions,\allowbreak{} an L1 mass bound is not an L-infinity or high-Lp bound. Different species have different transport fields,\allowbreak{} so their sum does not satisfy a closed scalar advection-diffusion equation. The deterministic triangular duality argument has not been transferred in this generality. The2026 entropy-dissipation paper treats specially chosen regularizing noise,\allowbreak{} not every source-admissible coefficient.


Verification:\allowbreak{} author checked partial | Independent review:\allowbreak{} none | Full target solved:\allowbreak{} False


\clearpage

\section*{P0272  Global stochastic reaction diffusion with quadratic reactions and mass control}

\entrylabel{Primary source}

Agresti-Veraar:\allowbreak{} Nonlinear SPDEs and Maximal Regularity,\allowbreak{}2025 extended survey;\allowbreak{} Problem8,\allowbreak{}PDF93;\allowbreak{}system8.14


\url{https://arxiv.org/pdf/2501.18561}


Source statement date:\allowbreak{} 2025-01-30. Archived access:\allowbreak{} 2026-09-06T23:\allowbreak{}40:\allowbreak{}06Z.


Historical:\allowbreak{} source\_\allowbreak{}says\_\allowbreak{}open | Current:\allowbreak{} unverified | Result:\allowbreak{} unsolved


\entrylabel{Question}

Under the positivity,\allowbreak{} Lipschitz reaction-noise,\allowbreak{} mass-control and quadratic-growth hypotheses of Problem8,\allowbreak{} prove global well-posedness in arbitrary spatial dimension with non-small regular transport noise.


\entrylabel{Scope}

The superlinear-noise extension is grouped once. This differs from Problem7 because a quadratic mass-controlled system need not have its lower-triangular reaction structure.


\subsection*{P0272-A001 | Explicit model calculation and hypothesis obstruction}

\entrylabel{Argument or result}

For the displayed Lotka-Volterra reaction f\_\allowbreak{}i=\allowbreak{}-tau\_\allowbreak{}i*u\_\allowbreak{}i+\allowbreak{}u\_\allowbreak{}i*sum\_\allowbreak{}j c\_\allowbreak{}ij*u\_\allowbreak{}j with c\_\allowbreak{}ij=\allowbreak{}-c\_\allowbreak{}ji and tau\_\allowbreak{}i>=\allowbreak{}0,\allowbreak{} summing gives sum\_\allowbreak{}i f\_\allowbreak{}i=\allowbreak{}-sum\_\allowbreak{}i tau\_\allowbreak{}i*u\_\allowbreak{}i<=\allowbreak{}0:\allowbreak{} every quadratic cross term cancels in a pair. However the L2 energy reaction is sum\_\allowbreak{}(i,\allowbreak{}j)c\_\allowbreak{}ij*u\_\allowbreak{}i\textasciicircum{}2*u\_\allowbreak{}j,\allowbreak{} whose pair contribution c\_\allowbreak{}ij*u\_\allowbreak{}i*u\_\allowbreak{}j*(u\_\allowbreak{}i-u\_\allowbreak{}j) has no fixed sign. For two species c12=\allowbreak{}1,\allowbreak{} u1=\allowbreak{}2,\allowbreak{}u2=\allowbreak{}1 gives a positive cubic contribution2.


\entrylabel{Second check}

This exact two-species computation proves that mass dissipation does not imply the usual L2 coercivity. It blocks a direct invocation of the coercive SPDE theorem while respecting the nonnegative state space.


\entrylabel{Remaining obstruction}

One needs duality and low-regularity parabolic estimates to transfer mass control into higher integrability. The2025 stochastic-flow De Giorgi-Nash-Moser result supplies a new scalar tool with sufficiently smooth transport fields,\allowbreak{} but no complete implication for the source system and its merely Holder transport coefficients is established here.


Verification:\allowbreak{} author checked partial | Independent review:\allowbreak{} none | Full target solved:\allowbreak{} False


\clearpage

\section*{P0273  Local stochastic Navier Stokes theory with no slip and non-small transport noise}

\entrylabel{Primary source}

Agresti-Veraar:\allowbreak{} Nonlinear SPDEs and Maximal Regularity,\allowbreak{}2025 extended survey;\allowbreak{} Problem9,\allowbreak{}PDF94;\allowbreak{}Theorem8.26


\url{https://arxiv.org/pdf/2501.18561}


Source statement date:\allowbreak{} 2025-01-30. Archived access:\allowbreak{} 2026-09-06T23:\allowbreak{}40:\allowbreak{}06Z.


Historical:\allowbreak{} source\_\allowbreak{}says\_\allowbreak{}open | Current:\allowbreak{} unverified | Result:\allowbreak{} unsolved


\entrylabel{Question}

Establish local well-posedness for three-dimensional Navier-Stokes equations on a smooth bounded domain with no-slip boundary conditions and non-small transport noise satisfying the natural stochastic parabolicity condition.


\entrylabel{Scope}

The small-noise case is already covered by a perturbation argument. The task is local existence,\allowbreak{} not deterministic global regularity.


\subsection*{P0273-A001 | Explicit model calculation and hypothesis obstruction}

\entrylabel{Argument or result}

For a constant transport vector b in an unbounded or periodic domain,\allowbreak{} the change of variables x going to x+\allowbreak{}b*W\_\allowbreak{}t removes Stratonovich transport noise and leaves a deterministic equation with the corresponding Ito drift correction. On a bounded domain this map moves the boundary;\allowbreak{} even if u initially vanishes on the fixed boundary,\allowbreak{} its translated values generally do not. Thus the simplest transport-removal proof does not preserve the required operator domain.


\entrylabel{Second check}

Take a smooth scalar test function vanishing exactly at a flat boundary and a vector b with nonzero normal component. Its translated boundary value is generally nonzero for W\_\allowbreak{}t different from zero. This directly demonstrates the failure before invoking any nonlinear estimate.


\entrylabel{Remaining obstruction}

A tangential stochastic flow might preserve the boundary but introduces random coefficients and its interaction with the Leray projection and Stokes no-slip domain. The required stochastic maximal regularity for the non-small noise remains unproved here.


Verification:\allowbreak{} author checked partial | Independent review:\allowbreak{} none | Full target solved:\allowbreak{} False


\clearpage

\section*{P0274  Optimal critical spatial smoothness for stochastic Navier Stokes transport noise}

\entrylabel{Primary source}

Agresti-Veraar:\allowbreak{} Nonlinear SPDEs and Maximal Regularity,\allowbreak{}2025 extended survey;\allowbreak{} Problem10,\allowbreak{}PDF94


\url{https://arxiv.org/pdf/2501.18561}


Source statement date:\allowbreak{} 2025-01-30. Archived access:\allowbreak{} 2026-09-06T23:\allowbreak{}40:\allowbreak{}06Z.


Historical:\allowbreak{} source\_\allowbreak{}says\_\allowbreak{}open | Current:\allowbreak{} unverified | Result:\allowbreak{} unsolved


\entrylabel{Question}

For smooth nonzero transport coefficients on the three-torus or R3,\allowbreak{} establish local well-posedness in a scaling-critical space of smoothness at most -1/\allowbreak{}2 and determine the optimal threshold.


\entrylabel{Scope}

The question is for the source general coefficient class. A favorable constant-noise model is a partial case,\allowbreak{} not its resolution.


\subsection*{P0274-A001 | Explicit model calculation and hypothesis obstruction}

\entrylabel{Argument or result}

For constant divergence-free transport fields on a torus or R3,\allowbreak{} write the equation in Stratonovich form with the matching second-order Ito correction. Let Z\_\allowbreak{}t=\allowbreak{}sum b\_\allowbreak{}n W\_\allowbreak{}t\textasciicircum{}n. If v is a deterministic Navier-Stokes solution,\allowbreak{} u(t,\allowbreak{}x)=\allowbreak{}v(t,\allowbreak{}x+\allowbreak{}Z\_\allowbreak{}t) solves the corresponding stochastic equation by Ito-Wentzell. Translation is an isometry on Lq and homogeneous Besov norms. Hence in this special commuting case a deterministic local theory transfers without any loss of its spatial smoothness threshold.


\entrylabel{Second check}

The quadratic-variation drift is (1/\allowbreak{}2)sum (b\_\allowbreak{}n dot grad)\textasciicircum{}2u. Stating this correction is essential:\allowbreak{} deleting Ito transport without correcting viscosity would solve a different equation. The Leray projection commutes with constant translations.


\entrylabel{Remaining obstruction}

For spatially varying coefficients the flow is not a translation,\allowbreak{} the transformed Laplacian has random variable coefficients,\allowbreak{} and the pressure projection no longer simply commutes. No general endpoint estimate or optimality proof is obtained. The2026 high-probability global result uses chosen noise and smoother initial spaces,\allowbreak{} so its abstract does not settle this question.


Verification:\allowbreak{} author checked partial | Independent review:\allowbreak{} none | Full target solved:\allowbreak{} False


\clearpage

\section*{P0275  Supremum endpoint Serrin criterion for stochastic Navier Stokes}

\entrylabel{Primary source}

Agresti-Veraar:\allowbreak{} Nonlinear SPDEs and Maximal Regularity,\allowbreak{}2025 extended survey;\allowbreak{} Problem11,\allowbreak{}PDF94


\url{https://arxiv.org/pdf/2501.18561}


Source statement date:\allowbreak{} 2025-01-30. Archived access:\allowbreak{} 2026-09-06T23:\allowbreak{}40:\allowbreak{}06Z.


Historical:\allowbreak{} source\_\allowbreak{}says\_\allowbreak{}open | Current:\allowbreak{} unverified | Result:\allowbreak{} unsolved


\entrylabel{Question}

For the source maximal critical-space stochastic Navier-Stokes solution (u,\allowbreak{}sigma),\allowbreak{} prove that finite sigma with sup\_\allowbreak{}(t<sigma)||u(t)||\_\allowbreak{}(B\textasciicircum{}(3/\allowbreak{}q-1)\_\allowbreak{}(q,\allowbreak{}p)) finite has probability zero.


\entrylabel{Scope}

The published theorem requires a stronger limiting condition. A bounded critical norm is not automatically a strong limit in that space.


\subsection*{P0275-A001 | Explicit model calculation and hypothesis obstruction}

\entrylabel{Argument or result}

The obstruction to a purely topological continuation argument is already visible in R3:\allowbreak{} choose a nonzero smooth compactly supported divergence-free vector field phi and put v\_\allowbreak{}n(x)=\allowbreak{}n*phi(n*x). Then ||v\_\allowbreak{}n||\_\allowbreak{}L3=\allowbreak{}||phi||\_\allowbreak{}L3 for every n,\allowbreak{} while v\_\allowbreak{}n tends to zero almost everywhere and has no strongly convergent L3 subsequence. If it had a strong subsequence,\allowbreak{} an almost-everywhere convergent subsubsequence would force its limit to zero,\allowbreak{} contradicting the constant positive norm.


\entrylabel{Second check}

The scaling is exact because n\textasciicircum{}3 from the cubic amplitude cancels n\textasciicircum{}-3 from the volume change. This is a sequence in the critical space,\allowbreak{} not a claimed trajectory of Navier-Stokes.


\entrylabel{Remaining obstruction}

Energy or compactness information supplied by the equation must rule out such concentration near a hypothetical blow-up time. Boundedness of the Banach-space norm alone cannot replace the source limit condition. No stochastic backward-uniqueness or concentration-compactness argument is completed.


Verification:\allowbreak{} author checked partial | Independent review:\allowbreak{} none | Full target solved:\allowbreak{} False


\clearpage

\section*{P0276  Product condition and cutoff for unweighted Abelian Cayley graphs}

\entrylabel{Primary source}

Salez:\allowbreak{} Modern aspects of Markov chains,\allowbreak{} entropy curvature and cutoff,\allowbreak{}2025;\allowbreak{} Problem2.4,\allowbreak{}PDF26-27


\url{https://arxiv.org/pdf/2508.21055}


Source statement date:\allowbreak{} 2025-08-28. Archived access:\allowbreak{} 2026-09-06T22:\allowbreak{}03:\allowbreak{}23Z.


Historical:\allowbreak{} source\_\allowbreak{}says\_\allowbreak{}open | Current:\allowbreak{} unverified | Result:\allowbreak{} candidate\_\allowbreak{}full\_\allowbreak{}solution


\entrylabel{Question}

Is total-variation cutoff equivalent to the spectral product condition for simple random walks on finite Abelian Cayley graphs?


\entrylabel{Scope}

Literal unrestricted uniform-generator statement;\allowbreak{} explicit unbounded-degree unweighted Abelian Cayley counterexample candidate. No claim about bounded-degree or random-generator variants;\allowbreak{} novelty and independent review unverified.


\subsection*{P0276-A001 | Explicit model calculation and hypothesis obstruction}

\entrylabel{Argument or result}

The weighted reset construction gives a useful obstruction that must not be mislabelled. Let Q\_\allowbreak{}d be continuous-time rate-one simple walk on the d-cube and add independent resets to uniform at rate theta\_\allowbreak{}d. Because the uniform projection commutes with its generator,\allowbreak{} the total-variation distance is exactly exp(-theta\_\allowbreak{}d*t) times the cube distance. Taking theta\_\allowbreak{}d*t\_\allowbreak{}d to1,\allowbreak{} where t\_\allowbreak{}d=\allowbreak{}(d/\allowbreak{}4)*log d,\allowbreak{} produces a gradual exp(-s) profile before the cube mixing scale although the inverse-gap-to-mixing ratio tends to zero. This destroys cutoff for the weighted walk.


\entrylabel{Second check}

The semigroup identity is e\textasciicircum{}(t(L+\allowbreak{}theta(Pi-I)))=\allowbreak{}Pi+\allowbreak{}e\textasciicircum{}(-theta*t)*(e\textasciicircum{}(tL)-Pi). The cube first-level eigenvalue is2/\allowbreak{}d,\allowbreak{} so adding theta of order1/\allowbreak{}(d log d) leaves relaxation of order d. The example thus really satisfies the product condition.


\entrylabel{Remaining obstruction}

To settle the stated uniform-generator problem one must either transfer the reset mechanism to an unweighted Abelian Cayley graph and verify the full state-space mixing profile,\allowbreak{} or prove a structural obstruction to that transfer. A possible auxiliary-factor construction has been identified for further checking;\allowbreak{} no full result is claimed at intake. The2026 local-product theorem uses a stronger constant and is not a resolution by its abstract.


Verification:\allowbreak{} author checked partial | Independent review:\allowbreak{} none | Full target solved:\allowbreak{} False


\subsection*{P0276-A002 | Unweighted Abelian auxiliary-factor realization of a rare uniform reset;\allowbreak{} exact Fourier spectrum}

\entrylabel{Argument or result}

The Cayley generator set on H\_\allowbreak{}d times Z\_\allowbreak{}m consists of all basis-coordinate increments paired with a nonzero auxiliary increment,\allowbreak{} and all nonzero cube increments paired with zero. Choosing m\textasciitilde{}2\textasciicircum{}d log(d)/\allowbreak{}4 gives full-state TV profile e\textasciicircum{}-s for0<s<1 and0 for s>1 on scale T\textasciitilde{}d log(d)/\allowbreak{}4. The exact full character spectrum gives gap\textasciitilde{}2/\allowbreak{}d. Hence the product condition holds while t\_\allowbreak{}mix(3/\allowbreak{}4)/\allowbreak{}t\_\allowbreak{}mix(1/\allowbreak{}4) tends to log(4/\allowbreak{}3),\allowbreak{} so cutoff fails.


\entrylabel{Second check}

Exact projected semigroup and conditional auxiliary coupling are checked separately;\allowbreak{} a complete character spectrum rules out a hidden slow auxiliary mode. The cube profile is proved by an exact chi-square formula and a magnetization Chebyshev argument.


\entrylabel{Remaining obstruction}

No gap found in author checks,\allowbreak{} but exact source scope,\allowbreak{} independent review and prior-art reconciliation remain necessary. The degree is unbounded and the generator set is specially chosen.


Verification:\allowbreak{} author checked candidate | Independent review:\allowbreak{} none | Full target solved:\allowbreak{} False


\entrylabel{Full proof packet}

proofs/\allowbreak{}P0276\_\allowbreak{}unweighted\_\allowbreak{}abelian\_\allowbreak{}cutoff\_\allowbreak{}counterexample\_\allowbreak{}candidate.md


\clearpage

\section*{P0277  Cutoff for conservative exclusion in dimension at least two}

\entrylabel{Primary source}

Salez:\allowbreak{} Modern aspects of Markov chains,\allowbreak{} entropy curvature and cutoff,\allowbreak{}2025;\allowbreak{} Problem2.5,\allowbreak{}PDF28


\url{https://arxiv.org/pdf/2508.21055}


Source statement date:\allowbreak{} 2025-08-28. Archived access:\allowbreak{} 2026-09-06T22:\allowbreak{}03:\allowbreak{}23Z.


Historical:\allowbreak{} source\_\allowbreak{}says\_\allowbreak{}open | Current:\allowbreak{} unverified | Result:\allowbreak{} unsolved


\entrylabel{Question}

Establish total-variation cutoff for the symmetric conservative exclusion process on growing lattices or tori of dimension at least two,\allowbreak{} in the many-particle regimes intended by the source.


\entrylabel{Scope}

The particle-growth regime must be stated:\allowbreak{} one particle or one hole reduces to ordinary torus walk and does not have cutoff. Results with reservoirs are nonconservative and do not answer this problem.


\subsection*{P0277-A001 | Explicit model calculation and hypothesis obstruction}

\entrylabel{Argument or result}

For a single-particle eigenfunction phi of the underlying graph Laplacian with mean zero,\allowbreak{} F(eta)=\allowbreak{}sum\_\allowbreak{}x eta\_\allowbreak{}x*phi(x) is an exclusion eigenfunction with the same eigenvalue. On a uniform m-particle configuration among N sites,\allowbreak{} direct sampling-without-replacement gives Var\_\allowbreak{}pi F =\allowbreak{} m*(N-m)/\allowbreak{}(N*(N-1))*sum\_\allowbreak{}x phi(x)\textasciicircum{}2. The initial mean decays exactly exponentially. This yields the natural lower-bound scale by comparing the initial eigenfunction amplitude with its equilibrium standard deviation.


\entrylabel{Second check}

The variance formula follows from E eta\_\allowbreak{}x=\allowbreak{}m/\allowbreak{}N and Cov(eta\_\allowbreak{}x,\allowbreak{}eta\_\allowbreak{}y)=\allowbreak{}-m*(N-m)/\allowbreak{}(N\textasciicircum{}2*(N-1)) for x different from y,\allowbreak{} using sum phi=\allowbreak{}0. All times must use a single rate convention;\allowbreak{} multiplying the generator by N changes the location but not cutoff.


\entrylabel{Remaining obstruction}

This spectral observable only proves a lower obstruction to mixing;\allowbreak{} it gives no matching total-variation upper bound. The one-dimensional height monotonicity argument fails on higher-dimensional tori. Salez2022 reservoir cutoff and Tran2023 Glauber-exclusion are different processes.


Verification:\allowbreak{} author checked partial | Independent review:\allowbreak{} none | Full target solved:\allowbreak{} False


\clearpage

\section*{P0278  Cutoff for unit-rate zero range on higher dimensional lattices}

\entrylabel{Primary source}

Salez:\allowbreak{} Modern aspects of Markov chains,\allowbreak{} entropy curvature and cutoff,\allowbreak{}2025;\allowbreak{} Problem2.6,\allowbreak{}PDF29


\url{https://arxiv.org/pdf/2508.21055}


Source statement date:\allowbreak{} 2025-08-28. Archived access:\allowbreak{} 2026-09-06T22:\allowbreak{}03:\allowbreak{}23Z.


Historical:\allowbreak{} source\_\allowbreak{}says\_\allowbreak{}open | Current:\allowbreak{} unverified | Result:\allowbreak{} unsolved


\entrylabel{Question}

Establish total-variation cutoff for the unit-rate conservative zero-range process on growing lattices or tori in dimension at least two,\allowbreak{} with an explicit particle-growth regime.


\entrylabel{Scope}

The source question does not specify density. Fixed one-particle chains must be excluded from any universal affirmative formulation. Unit rate means each occupied site emits at rate one,\allowbreak{} regardless of its occupancy.


\subsection*{P0278-A001 | Explicit model calculation and hypothesis obstruction}

\entrylabel{Argument or result}

On a symmetric regular graph the stationary measure is uniform over weak compositions of m into N parts. Stars-and-bars gives P\_\allowbreak{}pi(eta\_\allowbreak{}1=\allowbreak{}0)=\allowbreak{}(N-1)/\allowbreak{}(m+\allowbreak{}N-1),\allowbreak{} so the expected number of occupied sites is N*m/\allowbreak{}(m+\allowbreak{}N-1). Starting with all m particles at one site,\allowbreak{} the first departure clock has rate one under the per-site convention,\allowbreak{} whereas a fictitious independent-particle model would have rate m. Hence coupling to independent walkers without a time correction misses the slow dissolution of a pile.


\entrylabel{Second check}

The count for eta\_\allowbreak{}1=\allowbreak{}0 is binomial(m+\allowbreak{}N-2,\allowbreak{}N-2),\allowbreak{} divided by binomial(m+\allowbreak{}N-1,\allowbreak{}N-1). Detailed balance holds because moving one particle between occupied source and target has the reverse rate one and both configurations have equal stationary weight.


\entrylabel{Remaining obstruction}

The stationary occupancy identity does not quantify the concentration of pile-dissolution time or the subsequent spatial relaxation on a lattice. Mean-field cutoff results do not automatically transfer through comparison of spectral gaps.


Verification:\allowbreak{} author checked partial | Independent review:\allowbreak{} none | Full target solved:\allowbreak{} False


\clearpage

\section*{P0279  Cutoff for bounded-degree Ising throughout the rapid-mixing phase}

\entrylabel{Primary source}

Salez:\allowbreak{} Modern aspects of Markov chains,\allowbreak{} entropy curvature and cutoff,\allowbreak{}2025;\allowbreak{} Problem2.8,\allowbreak{}PDF31


\url{https://arxiv.org/pdf/2508.21055}


Source statement date:\allowbreak{} 2025-08-28. Archived access:\allowbreak{} 2026-09-06T22:\allowbreak{}03:\allowbreak{}23Z.


Historical:\allowbreak{} source\_\allowbreak{}says\_\allowbreak{}open | Current:\allowbreak{} unverified | Result:\allowbreak{} unsolved


\entrylabel{Question}

Prove cutoff for the Ising Gibbs sampler on bounded-degree graphs throughout the rapid-mixing phase specified by the source,\allowbreak{} beyond the high-temperature perturbative regime.


\entrylabel{Scope}

This is not the already proved lattice strong-spatial-mixing theorem or an assertion about arbitrary polynomial mixing at a critical point. The uniform structural regime and parameters need to be fixed.


\subsection*{P0279-A001 | Explicit model calculation and hypothesis obstruction}

\entrylabel{Argument or result}

At zero interaction with rate-one independent refresh clocks,\allowbreak{} the spins at time t from all plus are independent with mean e\textasciicircum{}-t. The squared chi-square distance to uniform is exactly (1+\allowbreak{}e\textasciicircum{}-2t)\textasciicircum{}N-1. This gives a total-variation upper bound tending to zero at t=\allowbreak{}(log N)/\allowbreak{}2+\allowbreak{}c as c tends to infinity. The total magnetization has mean N*e\textasciicircum{}-t and variance at most N;\allowbreak{} comparing this with its equilibrium mean0 and varianceN gives distance tending to one at t=\allowbreak{}(log N)/\allowbreak{}2-c as c tends to infinity. Thus the product model has cutoff with a bounded window.


\entrylabel{Second check}

The chi-square formula comes from multiplying the one-spin density second moments. For the lower bound,\allowbreak{} threshold at half the initial magnetization mean and apply Chebyshev under both laws. The per-site update convention is explicit.


\entrylabel{Remaining obstruction}

In the interacting model,\allowbreak{} a bounded entropy-decay constant proves logarithmic-order mixing but does not by itself identify a narrower window. The missing step is concentration of information or an equivalent local-product estimate uniformly over the full uniqueness regime,\allowbreak{} not another order-of-magnitude bound.


Verification:\allowbreak{} author checked partial | Independent review:\allowbreak{} none | Full target solved:\allowbreak{} False


\clearpage

\section*{P0280  Cutoff for bounded-degree hardcore throughout the rapid-mixing phase}

\entrylabel{Primary source}

Salez:\allowbreak{} Modern aspects of Markov chains,\allowbreak{} entropy curvature and cutoff,\allowbreak{}2025;\allowbreak{} Problem2.9,\allowbreak{}PDF32


\url{https://arxiv.org/pdf/2508.21055}


Source statement date:\allowbreak{} 2025-08-28. Archived access:\allowbreak{} 2026-09-06T22:\allowbreak{}03:\allowbreak{}23Z.


Historical:\allowbreak{} source\_\allowbreak{}says\_\allowbreak{}open | Current:\allowbreak{} unverified | Result:\allowbreak{} unsolved


\entrylabel{Question}

Prove cutoff for the hardcore Gibbs sampler on bounded-degree graphs throughout the source rapid-mixing regime,\allowbreak{} beyond the small-fugacity information-percolation range.


\entrylabel{Scope}

The question concerns fixed fugacity in a uniform structural phase. Critical or size-dependent fugacity limits may require different statements.


\subsection*{P0280-A001 | Explicit model calculation and hypothesis obstruction}

\entrylabel{Argument or result}

For an edgeless graph and fixed fugacity zeta>0,\allowbreak{} the stationary coordinates are independent Bernoulli(p),\allowbreak{} p=\allowbreak{}zeta/\allowbreak{}(1+\allowbreak{}zeta). Under rate-one refreshes and initial all ones,\allowbreak{} q\_\allowbreak{}t=\allowbreak{}p+\allowbreak{}(1-p)*e\textasciicircum{}-t. A one-coordinate chi-square calculation gives (q\_\allowbreak{}t-p)\textasciicircum{}2/\allowbreak{}[p*(1-p)] and hence full chi-square distance [1+\allowbreak{}((1-p)/\allowbreak{}p)*e\textasciicircum{}-2t]\textasciicircum{}N-1. An occupancy-sum Chebyshev lower bound matches the scale (log N)/\allowbreak{}2,\allowbreak{} yielding the familiar product cutoff.


\entrylabel{Second check}

The expression changes if one starts at all zeros,\allowbreak{} replacing (1-p)/\allowbreak{}p by p/\allowbreak{}(1-p);\allowbreak{} for fixed p both give the same leading time and bounded window. On a graph with edges the all-one configuration is inadmissible,\allowbreak{} so this is explicitly an edgeless baseline.


\entrylabel{Remaining obstruction}

Exclusion of adjacent particles creates long dependent update histories. Spectral independence and O(N log N) discrete-time mixing do not alone prove a vanishing relative window. No matching information-percolation or local-variance estimate in the full phase is established here.


Verification:\allowbreak{} author checked partial | Independent review:\allowbreak{} none | Full target solved:\allowbreak{} False


\clearpage

\section*{P0281  Log concavity of finite particle TASEP transition probabilities}

\entrylabel{Primary source}

AIM All roads to KPZ universality:\allowbreak{} workshop open problems,\allowbreak{}March19 2025;\allowbreak{} Section3,\allowbreak{}PDF2;\allowbreak{}follow-up Q2,\allowbreak{}PDF3


\url{https://lpetrov.cc/LaTeX/2025-03-19-AIM-KPZ/2025-03-19-AIM-KPZ.pdf}


Source statement date:\allowbreak{} 2025-03-19. Archived access:\allowbreak{} 2026-09-06T23:\allowbreak{}41:\allowbreak{}31Z.


Historical:\allowbreak{} source\_\allowbreak{}says\_\allowbreak{}open | Current:\allowbreak{} unverified | Result:\allowbreak{} unsolved


\entrylabel{Question}

For continuous-time TASEP with n particles started from the step configuration,\allowbreak{} determine whether its fixed-time transition mass function on ordered particle configurations is log concave.


\entrylabel{Scope}

The lattice midpoint assertion is considered when the midpoint is integral and is a legal ordered configuration. A real-variable extension or a stronger rounded-midpoint notion would require a separate definition.


\subsection*{P0281-A001 | Exact identities,\allowbreak{} finite controls and structural obstruction}

\entrylabel{Argument or result}

For one particle the displacement has Poisson mass p\_\allowbreak{}k=\allowbreak{}e\textasciicircum{}-t*t\textasciicircum{}k/\allowbreak{}k!,\allowbreak{} so p\_\allowbreak{}k\textasciicircum{}2/\allowbreak{}(p\_\allowbreak{}(k-1)*p\_\allowbreak{}(k+\allowbreak{}1))=\allowbreak{}(k+\allowbreak{}1)/\allowbreak{}k for k>=\allowbreak{}1. Thus the scalar nearest-neighbor log-concavity inequality holds strictly. Independent freely moving particles would give a product of these masses and preserve this inequality on every integral coordinate midpoint. Actual TASEP jump histories have blocking constraints and do not form that independent product.


\entrylabel{Second check}

The ratio follows by cancellation of adjacent factorials;\allowbreak{} the same conclusion follows from discrete concavity of -log(k!). The independent-particle calculation is explicitly a comparison model,\allowbreak{} not a TASEP formula.


\entrylabel{Remaining obstruction}

Blocking changes holding rates according to the configuration. The determinantal transition formula contains signed terms,\allowbreak{} and log concavity is not preserved under signed summation. No inequality controlling these cancellations for all n and configurations is obtained. The workshop follow-up reports no resolution;\allowbreak{} a targeted September2026 search found no matching theorem,\allowbreak{} which is not a current-openness certificate.


Verification:\allowbreak{} author checked partial | Independent review:\allowbreak{} none | Full target solved:\allowbreak{} False


\clearpage

\section*{P0282  Chen log concavity conjecture for fixed size random permutation LIS}

\entrylabel{Primary source}

AIM all roads to KPZ workshop follow-up report,\allowbreak{}2025;\allowbreak{} Log-concavity Q1,\allowbreak{}PDF3;\allowbreak{}BKM Conjecture1


\url{https://aimath.org/pastworkshops/roadtokpzrep.pdf}


Source statement date:\allowbreak{} 2025\_\allowbreak{}workshop\_\allowbreak{}followup\_\allowbreak{}undated. Archived access:\allowbreak{} 2026-09-07T00:\allowbreak{}14:\allowbreak{}14Z.


Historical:\allowbreak{} source\_\allowbreak{}says\_\allowbreak{}open | Current:\allowbreak{} unverified | Result:\allowbreak{} unsolved


\entrylabel{Question}

For a uniform permutation of exactly n elements,\allowbreak{} is the probability mass function of the longest increasing subsequence length log concave for every n?


\entrylabel{Scope}

Poissonized Plancherel log concavity is an established theorem and is excluded from this target. Involution variants are not split off as extra entries.


\subsection*{P0282-A001 | Exact identities,\allowbreak{} finite controls and structural obstruction}

\entrylabel{Argument or result}

The RSK formula gives n!*P(L\_\allowbreak{}n=\allowbreak{}k)=\allowbreak{}sum\_\allowbreak{}(lambda partitions n,\allowbreak{}lambda\_\allowbreak{}1=\allowbreak{}k) f\_\allowbreak{}lambda\textasciicircum{}2,\allowbreak{} where f\_\allowbreak{}lambda=\allowbreak{}n!/\allowbreak{}product of hook lengths. Exact integer evaluation for n=\allowbreak{}1 through15 verifies all adjacent log-concavity inequalities and total mass n!. For n=\allowbreak{}4 the counts are (1,\allowbreak{}13,\allowbreak{}9,\allowbreak{}1). The proof question is whether a shape-preserving injection can compare the neighboring sums while keeping the same n.


\entrylabel{Second check}

The independent total-mass check sum f\_\allowbreak{}lambda\textasciicircum{}2=\allowbreak{}n! detects partition or hook errors. Mixing alone cannot justify de-Poissonization:\allowbreak{} half each of conditional masses (.49,\allowbreak{}.02,\allowbreak{}.49) and (.01,\allowbreak{}.98,\allowbreak{}.01) gives the log-concave mass (.25,\allowbreak{}.50,\allowbreak{}.25),\allowbreak{} although the first conditional law is not log concave. This is a logical obstruction,\allowbreak{} not a counterexample to Chen.


\entrylabel{Remaining obstruction}

Neither the known Poissonized theorem nor log-concavity of the limiting Tracy-Widom density controls the coefficientwise inequality at every fixed n. BKM arXiv2412.15116 retains Conjecture1 and explicitly distinguishes its theorem. The finite checks do not supply the missing injection or coefficientwise inequality.


Verification:\allowbreak{} author checked partial | Independent review:\allowbreak{} none | Full target solved:\allowbreak{} False


\clearpage

\section*{P0283  Exact binary longest common subsequence time constant}

\entrylabel{Primary source}

AIM All roads to KPZ universality:\allowbreak{} workshop open problems,\allowbreak{}March19 2025;\allowbreak{} Section6,\allowbreak{}PDF3


\url{https://lpetrov.cc/LaTeX/2025-03-19-AIM-KPZ/2025-03-19-AIM-KPZ.pdf}


Source statement date:\allowbreak{} 2025-03-19. Archived access:\allowbreak{} 2026-09-06T23:\allowbreak{}41:\allowbreak{}31Z.


Historical:\allowbreak{} source\_\allowbreak{}says\_\allowbreak{}open | Current:\allowbreak{} unverified | Result:\allowbreak{} unsolved


\entrylabel{Question}

Determine the exact almost-sure linear growth constant for the longest common subsequence of two independent strings of iid uniform binary letters of equal length.


\entrylabel{Scope}

This target is the mean or law-of-large-numbers constant. The distinct fluctuation-law question is P0284. Other alphabet sizes are grouped rather than counted as parameter variants.


\subsection*{P0283-A001 | Exact identities,\allowbreak{} finite controls and structural obstruction}

\entrylabel{Argument or result}

Concatenating optimal alignments from independent blocks gives E L\_\allowbreak{}(m+\allowbreak{}n)>=\allowbreak{}E L\_\allowbreak{}m+\allowbreak{}E L\_\allowbreak{}n,\allowbreak{} hence the usual superadditive existence argument and lower bounds gamma>=\allowbreak{}E L\_\allowbreak{}m/\allowbreak{}m. At block size2,\allowbreak{} four equal pairs among the16 binary-word pairs contribute length2,\allowbreak{} two complementary constant pairs contribute0,\allowbreak{} and the other ten contribute1. Thus E L\_\allowbreak{}2=\allowbreak{}9/\allowbreak{}8 and gamma>=\allowbreak{}9/\allowbreak{}16. For q uniform letters,\allowbreak{} match indicators eta\_\allowbreak{}ij=\allowbreak{}1\{X\_\allowbreak{}i=\allowbreak{}Y\_\allowbreak{}j\} are pairwise independent,\allowbreak{} but the four indicators around a rectangle are all1 with probability q\textasciicircum{}-3,\allowbreak{} not q\textasciicircum{}-4.


\entrylabel{Second check}

The block2 count is independently checked by the dynamic-programming recurrence on all16 pairs. The rectangle event requires four letters to have one common value,\allowbreak{} giving q choices out of q\textasciicircum{}4 outcomes. Sharing just one endpoint imposes two independent equality constraints,\allowbreak{} explaining pairwise independence.


\entrylabel{Remaining obstruction}

The exact time constant for independent Bernoulli last-passage weights cannot be substituted because the rectangle dependence changes the optimization environment. The block lower bound is far from an exact constant;\allowbreak{} no matching upper construction or Busemann law for the dependent environment is obtained.


Verification:\allowbreak{} author checked partial | Independent review:\allowbreak{} none | Full target solved:\allowbreak{} False


\clearpage

\section*{P0284  Gaussian fluctuations for uniform binary longest common subsequences}

\entrylabel{Primary source}

AIM all roads to KPZ workshop follow-up report,\allowbreak{}2025;\allowbreak{} Longest common subsequence,\allowbreak{}PDF4


\url{https://aimath.org/pastworkshops/roadtokpzrep.pdf}


Source statement date:\allowbreak{} 2025\_\allowbreak{}workshop\_\allowbreak{}followup\_\allowbreak{}undated. Archived access:\allowbreak{} 2026-09-07T00:\allowbreak{}14:\allowbreak{}14Z.


Historical:\allowbreak{} source\_\allowbreak{}says\_\allowbreak{}open | Current:\allowbreak{} unverified | Result:\allowbreak{} unsolved


\entrylabel{Question}

For two independent uniform binary words,\allowbreak{} prove or disprove asymptotic normality of the centered longest common subsequence length normalized by its standard deviation,\allowbreak{} and establish the needed fluctuation scale.


\entrylabel{Scope}

Nonuniform-letter CLTs and conditional theorems assuming a variance lower bound do not settle the uniform-letter question. All fluctuation-scale subquestions are grouped with this distributional target.


\subsection*{P0284-A001 | Exact identities,\allowbreak{} finite controls and structural obstruction}

\entrylabel{Argument or result}

Changing one of the2n input letters changes L\_\allowbreak{}n by at most1. Efron-Stein therefore gives Var L\_\allowbreak{}n <=\allowbreak{}(1/\allowbreak{}2)*sum\_\allowbreak{}(j=\allowbreak{}1)\textasciicircum{}(2n) E[(L\_\allowbreak{}n-L\_\allowbreak{}n\textasciicircum{}(j))\textasciicircum{}2]<=\allowbreak{}n/\allowbreak{}2 in the binary model:\allowbreak{} a resampled fair bit differs with probability1/\allowbreak{}2. Bounded differences also gives P(|L\_\allowbreak{}n-E L\_\allowbreak{}n|>=\allowbreak{}t)<=\allowbreak{}2 exp(-t\textasciicircum{}2/\allowbreak{}n). These exact upper bounds rule out fluctuations substantially larger than sqrt(n),\allowbreak{} while allowing smaller scales.


\entrylabel{Second check}

The Lipschitz bound follows by deleting at most one matched position from an optimal subsequence and applying the same argument in reverse. The factor n/\allowbreak{}2 uses the resampling mismatch probability;\allowbreak{} the cruder distribution-free bound is n. The known Houdré-Islak2023 CLT explicitly assumes a variance lower bound,\allowbreak{} rather than deriving it for every uniform alphabet. The archived2023 primary paper,\allowbreak{} equation1.4 onPDF3,\allowbreak{} attributes the exact general bound n*(1-sum p\_\allowbreak{}k\textasciicircum{}2) to Steele;\allowbreak{} this campaign calculation is a rederivation of that known bound.


\entrylabel{Remaining obstruction}

Neither concentration nor pairwise independence of matching indicators proves a positive linear variance lower bound. Independent Bernoulli matching models can have different limiting laws. No uniform-model lower bound or direct normal approximation avoiding it is proved here.


Verification:\allowbreak{} author checked partial | Independent review:\allowbreak{} none | Full target solved:\allowbreak{} False


\clearpage

\section*{P0285  Finitary ordering permutation of a KPZ line ensemble vertical slice}

\entrylabel{Primary source}

AIM All roads to KPZ universality:\allowbreak{} workshop open problems,\allowbreak{}March19 2025;\allowbreak{} Section7,\allowbreak{}PDF3;\allowbreak{}follow-up PDF4-5


\url{https://lpetrov.cc/LaTeX/2025-03-19-AIM-KPZ/2025-03-19-AIM-KPZ.pdf}


Source statement date:\allowbreak{} 2025-03-19. Archived access:\allowbreak{} 2026-09-06T23:\allowbreak{}41:\allowbreak{}31Z.


Historical:\allowbreak{} source\_\allowbreak{}says\_\allowbreak{}open | Current:\allowbreak{} unverified | Result:\allowbreak{} unsolved


\entrylabel{Question}

For a fixed vertical slice of the KPZ line ensemble in the source normalization,\allowbreak{} is the permutation that orders its levels well defined and finitary? If it is not finitary,\allowbreak{} quantify its displacement.


\entrylabel{Scope}

The source distinguishes a uniform Gibbs interaction and a level-dependent interaction. The question is retained as one target;\allowbreak{} conclusions for one normalization cannot be transferred to the other by ignoring deterministic level shifts.


\subsection*{P0285-A001 | Exact identities,\allowbreak{} finite controls and structural obstruction}

\entrylabel{Argument or result}

A sufficient criterion is summability over all pairs of the probabilities of an inversion relative to the reference level order. Borel-Cantelli then gives only finitely many inversions,\allowbreak{} so an existing ranking permutation differs from the identity on finitely many labels. Merely making high-level inversion probabilities tend to zero is insufficient. Independently swap the two labels in block \{2j-1,\allowbreak{}2j\} with probability1/\allowbreak{}j for j>=\allowbreak{}2. This defines a permutation with displacement at most1 and marginal displacement probability tending to zero,\allowbreak{} yet it has infinitely many swaps almost surely.


\entrylabel{Second check}

The independent second Borel-Cantelli lemma applies because sum1/\allowbreak{}j diverges. Each realized transposition is confined to a disjoint two-label block,\allowbreak{} so the example is a valid permutation rather than an ordering of an unbounded collection without a first element.


\entrylabel{Remaining obstruction}

The toy permutation is not a KPZ counterexample. The missing source-specific estimate is a summable bound on deep-level inversions,\allowbreak{} together with existence of the infinite ranking. Bounded displacement alone cannot answer the finitary part,\allowbreak{} as the explicit example shows.


Verification:\allowbreak{} author checked partial | Independent review:\allowbreak{} none | Full target solved:\allowbreak{} False


\clearpage

\section*{P0286  Stationary two color ASEP covariance decay on a large ring}

\entrylabel{Primary source}

AIM All roads to KPZ universality:\allowbreak{} workshop open problems,\allowbreak{}March19 2025;\allowbreak{} Section5,\allowbreak{}PDF3


\url{https://lpetrov.cc/LaTeX/2025-03-19-AIM-KPZ/2025-03-19-AIM-KPZ.pdf}


Source statement date:\allowbreak{} 2025-03-19. Archived access:\allowbreak{} 2026-09-06T23:\allowbreak{}41:\allowbreak{}31Z.


Historical:\allowbreak{} source\_\allowbreak{}says\_\allowbreak{}open | Current:\allowbreak{} unverified | Result:\allowbreak{} unsolved


\entrylabel{Question}

For stationary two-color ASEP on a growing ring with fixed color densities,\allowbreak{} determine the spatial decay and scaling of occupation covariances,\allowbreak{} including the near-half-density regime discussed in the source.


\entrylabel{Scope}

Finite-ring conservation corrections must be separated from an infinite-volume or mesoscopic decay assertion. Color densities alone do not make the asymmetric multi-species stationary law uniform.


\subsection*{P0286-A001 | Exact identities,\allowbreak{} finite controls and structural obstruction}

\entrylabel{Argument or result}

With fixed color counts N\_\allowbreak{}i,\allowbreak{} the exact sum rule is sum\_\allowbreak{}x Cov(eta\_\allowbreak{}i(0),\allowbreak{}eta\_\allowbreak{}j(x))=\allowbreak{}Cov(eta\_\allowbreak{}i(0),\allowbreak{}N\_\allowbreak{}j)=\allowbreak{}0. At zero separation the same-color covariance is rho\_\allowbreak{}i*(1-rho\_\allowbreak{}i),\allowbreak{} and the different-color covariance is -rho\_\allowbreak{}i*rho\_\allowbreak{}j. In the symmetric uniform-arrangement baseline,\allowbreak{} exchangeability gives off-site values -rho\_\allowbreak{}i*(1-rho\_\allowbreak{}i)/\allowbreak{}(N-1) for the same color and +\allowbreak{}rho\_\allowbreak{}i*rho\_\allowbreak{}j/\allowbreak{}(N-1) for different colors.


\entrylabel{Second check}

The off-site formulas follow both from sampling without replacement and from the sum rule plus exchangeability. They sum exactly to cancel the zero-site covariance. A decay claim uniform in the ring size must account for this order1/\allowbreak{}N correction.


\entrylabel{Remaining obstruction}

For biased multi-color dynamics the stationary correlations depend on the ordering structure and cannot be derived from the conservation identities alone. These identities give signed integral constraints,\allowbreak{} not the requested asymptotic decay kernel. No analysis of the relevant stationary horizon or matrix-product correlations is completed.


Verification:\allowbreak{} author checked partial | Independent review:\allowbreak{} none | Full target solved:\allowbreak{} False


\clearpage

\section*{P0287  Invariant measures for full line TASEP with killing at one site}

\entrylabel{Primary source}

AIM all roads to KPZ workshop follow-up report,\allowbreak{}2025;\allowbreak{} ASEP on graphs or trees,\allowbreak{}S2,\allowbreak{}PDF5


\url{https://aimath.org/pastworkshops/roadtokpzrep.pdf}


Source statement date:\allowbreak{} 2025\_\allowbreak{}workshop\_\allowbreak{}followup\_\allowbreak{}undated. Archived access:\allowbreak{} 2026-09-07T00:\allowbreak{}14:\allowbreak{}14Z.


Historical:\allowbreak{} source\_\allowbreak{}says\_\allowbreak{}open | Current:\allowbreak{} unverified | Result:\allowbreak{} unsolved


\entrylabel{Question}

Characterize invariant measures for unit-rate full-line TASEP with particles at site0 removed at a fixed positive rate kappa,\allowbreak{} as proposed in workshop systemS2.


\entrylabel{Scope}

The empty invariant law is trivial and does not resolve classification of nontrivial stationary flows. Boundary or tail densities must be included when specifying a particular nontrivial invariant law.


\subsection*{P0287-A001 | Exact identities,\allowbreak{} finite controls and structural obstruction}

\entrylabel{Argument or result}

Writing J\_\allowbreak{}i=\allowbreak{}E[eta\_\allowbreak{}i*(1-eta\_\allowbreak{}(i+\allowbreak{}1))],\allowbreak{} stationarity of eta\_\allowbreak{}i gives J\_\allowbreak{}(i-1)-J\_\allowbreak{}i=\allowbreak{}kappa*1\_\allowbreak{}(i=\allowbreak{}0)*E eta\_\allowbreak{}0. There is also an exact obstruction to any independent Bernoulli invariant law whose site densities r\_\allowbreak{}i all lie strictly between0 and1. For adjacent sites,\allowbreak{} subtract the derivative of r\_\allowbreak{}i*r\_\allowbreak{}(i+\allowbreak{}1) from the generator derivative of E[eta\_\allowbreak{}i*eta\_\allowbreak{}(i+\allowbreak{}1)]. At a product law the covariance derivative equals J\_\allowbreak{}i*(r\_\allowbreak{}(i+\allowbreak{}1)-r\_\allowbreak{}i),\allowbreak{} with killing terms cancelling. Stationarity therefore forces all adjacent densities equal. The one-site balance at0 then forces kappa*r\_\allowbreak{}0=\allowbreak{}0,\allowbreak{} a contradiction.


\entrylabel{Second check}

The internal jump i to i+\allowbreak{}1 leaves eta\_\allowbreak{}i*eta\_\allowbreak{}(i+\allowbreak{}1) equal to0 both before and after,\allowbreak{} whereas it changes the two marginal means by -J\_\allowbreak{}i and +\allowbreak{}J\_\allowbreak{}i. Incoming and outgoing exterior jumps factor under the product law,\allowbreak{} giving the displayed covariance derivative. Independent killing at0 contributes the same product term in both expressions and cancels.


\entrylabel{Remaining obstruction}

Thus even a piecewise Bernoulli product with strictly interior densities cannot be the desired invariant law. The source matrix-product approach needs nontrivial correlations and a telescoping identity at the killing defect. Neither this obstruction nor the current-balance equation supplies that identity or classifies the correlated stationary measures.


Verification:\allowbreak{} author checked partial | Independent review:\allowbreak{} none | Full target solved:\allowbreak{} False


\clearpage

\section*{P0288  Rank one polynomial asymptotics for lambda transient Markov chains}

\entrylabel{Primary source}

Pollett:\allowbreak{} Quasi stationary distributions for queueing and other models,\allowbreak{}2022;\allowbreak{} Section2 Conjecture(ii),\allowbreak{}PDF2;\allowbreak{}discussionPDF3


\url{https://arxiv.org/pdf/2202.05428}


Source statement date:\allowbreak{} 2022-02-14. Archived access:\allowbreak{} 2026-09-07T00:\allowbreak{}46:\allowbreak{}35Z.


Historical:\allowbreak{} source\_\allowbreak{}says\_\allowbreak{}open | Current:\allowbreak{} unverified | Result:\allowbreak{} candidate\_\allowbreak{}full\_\allowbreak{}solution


\entrylabel{Question}

For a lambda-transient irreducible countable transient class with almost-sure absorption,\allowbreak{} must every finite strictly positive limit of t\textasciicircum{}a*exp(lambda*t)*p\_\allowbreak{}t(i,\allowbreak{}j) factor as A*x\_\allowbreak{}i*m\_\allowbreak{}j,\allowbreak{} with invariant or subinvariant positive factors?


\entrylabel{Scope}

The question is the factorization assertion(ii). Neighboring existence,\allowbreak{} common-exponent and survival claims are not multiplied into extra entries or all declared resolved.


\subsection*{P0288-A001 | Exact rank-one resolvent perturbation and spectral-edge inversion}

\entrylabel{Argument or result}

Nearest-neighbor continuous-time walk on Z,\allowbreak{} killed at rate kappa everywhere and additional ratev at0,\allowbreak{} has t\textasciicircum{}(3/\allowbreak{}2)*e\textasciicircum{}(kappa*t)*p\_\allowbreak{}t(i,\allowbreak{}j) tending to [(|i|+\allowbreak{}2/\allowbreak{}v)(|j|+\allowbreak{}2/\allowbreak{}v)+\allowbreak{}ij]/\allowbreak{}(4*sqrt(pi)). The class is lambda-transient with lambda=\allowbreak{}kappa. Atv2 the0,\allowbreak{}1 minor equals1/\allowbreak{}(16pi)>0,\allowbreak{} ruling out rank-one factorization.


\entrylabel{Second check}

The odd sector is unaffected by the defect. Fourier Gaussian scaling gives the independent difference limit C(i,\allowbreak{}j)-C(i,\allowbreak{}-j)=\allowbreak{}ij/\allowbreak{}(2sqrt(pi)),\allowbreak{} incompatible with a positive symmetric reflection-invariant rank-one kernel. The explicit Stieltjes density justifies pointwise inversion rather than relying on an invalid bare Tauberian inference.


\entrylabel{Remaining obstruction}

The original proof packet received a scoped independent source and spectral review,\allowbreak{} pinned before the later attribution addition. Foley-McDonald published a discrete two-end construction from which the parent reviewer derives the same general continuous-time obstruction. Current exact-publication status and originality remain unverified;\allowbreak{} no original-solution promotion. The separate survival assertion is not settled.


Verification:\allowbreak{} complete author checked candidate | Independent review:\allowbreak{} scoped parent review of original packet completed no novelty certificate | Full target solved:\allowbreak{} False


\subsection*{P0288-A002 | Read parent independent review and verify its Poissonization limit and minor}

\entrylabel{Argument or result}

Published Foley-McDonald two-end discrete kernel has a rank-two asymptotic. Parent-derived rate-one Poissonization gives the same general continuous-time factorization obstruction;\allowbreak{} exact generators are not identified. Full argument is appended with explicit attribution.


\entrylabel{Second check}

Poisson concentration plus the global a\_\allowbreak{}n<=\allowbreak{}C(1+\allowbreak{}n)\textasciicircum{}(-3/\allowbreak{}2) bound proves the expectation limit;\allowbreak{} determinant13*41-23\textasciicircum{}2=\allowbreak{}4,\allowbreak{} so diagonal weights give positive limiting minor mu\textasciicircum{}(-3)/\allowbreak{}pi.


\entrylabel{Remaining obstruction}

No novelty certificate. Whether the exact continuous-time corollary or delta-killing realization was previously published remains unverified. This update adds attribution after the independently reviewed original proof hash.


Verification:\allowbreak{} author checked attribution and reduction | Independent review:\allowbreak{} parent contributed reduction original packet only scoped review | Full target solved:\allowbreak{} False


\entrylabel{Full proof packet}

proofs/\allowbreak{}P0288\_\allowbreak{}pollett\_\allowbreak{}rank\_\allowbreak{}two\_\allowbreak{}transition\_\allowbreak{}counterexample\_\allowbreak{}candidate.md


\clearpage

\section*{P0289  Ross modulation monotonicity and a criterion for Cox workload ordering}

\entrylabel{Primary source}

Leskela:\allowbreak{} Ross second conjecture and supermodular stochastic ordering,\allowbreak{}2022;\allowbreak{} Problem1,\allowbreak{}journal214,\allowbreak{}PDF3


\url{https://acris.aalto.fi/ws/portalfiles/portal/89616571/Ross_s_second_conjecture_and_supermodular_stochastic_ordering.pdf}


Source statement date:\allowbreak{} 2022-05-01. Archived access:\allowbreak{} 2026-09-07T00:\allowbreak{}46:\allowbreak{}33Z.


Historical:\allowbreak{} source\_\allowbreak{}says\_\allowbreak{}open | Current:\allowbreak{} unverified | Result:\allowbreak{} candidate\_\allowbreak{}partial\_\allowbreak{}counterexample


\entrylabel{Question}

Is the mean stationary workload in a stable Cox/\allowbreak{}GI/\allowbreak{}1 queue always nonincreasing in the modulation speed of its stationary ergodic intensity? If not,\allowbreak{} characterize when the order holds.


\entrylabel{Scope}

The yes-or-no monotonicity clause and the source request for a necessary-and-sufficient condition are grouped. A counterexample to the first clause does not finish the characterization request.


\subsection*{P0289-A001 | Explicit model calculation,\allowbreak{} comparison identities and remaining hypothesis obstruction}

\entrylabel{Argument or result}

A complete candidate refutation of universal monotonicity is in the proof packet. Use deterministic unit service,\allowbreak{} epsilon=\allowbreak{}1/\allowbreak{}10000,\allowbreak{} and a1024-state clockwise CTMC with first Fourier eigenvalue -tan(pi/\allowbreak{}1024)+\allowbreak{}i. The intensity is epsilon*[2+\allowbreak{}cos(2pi*X/\allowbreak{}1024+\allowbreak{}pi/\allowbreak{}2048)],\allowbreak{} with distinct positive state values. A workload-minus-parallel-shot monotonicity argument gives the uniform expansion w(c)=\allowbreak{}epsilon+\allowbreak{}2epsilon\textasciicircum{}2+\allowbreak{}(epsilon\textasciicircum{}2/\allowbreak{}2)*integral\_\allowbreak{}0\textasciicircum{}1(1-u)*exp[-c*u*tan(pi/\allowbreak{}1024)]*cos(c*u)du+\allowbreak{}R\_\allowbreak{}c,\allowbreak{} where0<=\allowbreak{}R\_\allowbreak{}c<=\allowbreak{}27epsilon\textasciicircum{}3/\allowbreak{}[2(1-3epsilon)]. It follows exactly that w(3pi)-w(2pi)>epsilon\textasciicircum{}2*2491/\allowbreak{}1612800>0.


\entrylabel{Second check}

The second-order kernel1-u is independently recovered by integrating the interaction workload of two arrivals over the age of the younger one. The proof bounds the remainder through domination by a stable homogeneous M/\allowbreak{}D/\allowbreak{}1 queue,\allowbreak{} rather than numerically fitting a light-traffic series. The rational lower-bound arithmetic is recorded separately.


\entrylabel{Remaining obstruction}

Universal monotonicity has a full counterexample candidate,\allowbreak{} but the necessary-and-sufficient condition requested in the same source problem remains open here. Earlier Ross-type literature,\allowbreak{} including Miyoshi-Rolski2004 and Heyman,\allowbreak{} must be compared in full before novelty is claimed. No independent review is claimed.


Verification:\allowbreak{} author checked partial | Independent review:\allowbreak{} none | Full target solved:\allowbreak{} False


\entrylabel{Full proof packet}

proofs/\allowbreak{}P0289\_\allowbreak{}ross\_\allowbreak{}modulation\_\allowbreak{}counterexample\_\allowbreak{}candidate.md


\clearpage

\section*{P0290  Ross workload ordering for an image of a doubly monotone Markov chain}

\entrylabel{Primary source}

Leskela:\allowbreak{} Ross second conjecture and supermodular stochastic ordering,\allowbreak{}2022;\allowbreak{} Problem2,\allowbreak{}journal214,\allowbreak{}PDF3


\url{https://acris.aalto.fi/ws/portalfiles/portal/89616571/Ross_s_second_conjecture_and_supermodular_stochastic_ordering.pdf}


Source statement date:\allowbreak{} 2022-05-01. Archived access:\allowbreak{} 2026-09-07T00:\allowbreak{}46:\allowbreak{}33Z.


Historical:\allowbreak{} source\_\allowbreak{}says\_\allowbreak{}open | Current:\allowbreak{} unverified | Result:\allowbreak{} unsolved


\entrylabel{Question}

Does Ross workload monotonicity hold when the intensity is an arbitrary deterministic,\allowbreak{} possibly Boolean,\allowbreak{} function of a stationary doubly stochastically monotone CTMC?


\entrylabel{Scope}

An arbitrary image need not preserve the order on the original state space. Failure of a sufficient supermodular-order criterion is not itself failure of the workload assertion.


\subsection*{P0290-A001 | Explicit model calculation,\allowbreak{} comparison identities and remaining hypothesis obstruction}

\entrylabel{Argument or result}

Consider the reversible birth-death chain on states1,\allowbreak{}2,\allowbreak{}3 with symmetric edge rates1 on1-2 and2 on2-3. It is doubly stochastically monotone. Its slow nonzero eigenvalue is3-sqrt(3),\allowbreak{} with eigenvector (1,\allowbreak{}sqrt(3)-2,\allowbreak{}1-sqrt(3)). The last two components have the same sign,\allowbreak{} so P\_\allowbreak{}t(2,\allowbreak{}3)>1/\allowbreak{}3 for sufficiently large t. Label intensities f(1)=\allowbreak{}1,\allowbreak{}f(2)=\allowbreak{}0,\allowbreak{}f(3)=\allowbreak{}2. For g=\allowbreak{}1\{f>=\allowbreak{}1\},\allowbreak{}h=\allowbreak{}1\{f>=\allowbreak{}2\},\allowbreak{} covariance at lag t is -(P\_\allowbreak{}t(2,\allowbreak{}3)-1/\allowbreak{}3)/\allowbreak{}3<0. Hence the image process does not satisfy the supermodular time-slowing order,\allowbreak{} whose infinite-lag limit would require this covariance nonnegative.


\entrylabel{Second check}

The eigenvector sums to zero and direct multiplication by the symmetric generator gives eigenvalue -(3-sqrt3). The remaining nonzero mode decays at rate3+\allowbreak{}sqrt3,\allowbreak{} so it cannot cancel the eventual positive transition overshoot. The covariance formula uses g=\allowbreak{}1-1\{X=\allowbreak{}2\} and h=\allowbreak{}1\{X=\allowbreak{}3\}.


\entrylabel{Remaining obstruction}

This blocks transferring the established sufficient supermodular-order argument through arbitrary f. It does not produce a queue-workload reversal. In fact same-observable covariance under a reversible environment is a nonnegative mixture of exponentials,\allowbreak{} so the simplest light-traffic coefficient remains monotone. Higher-order queue effects or a different comparison method are needed.


Verification:\allowbreak{} author checked partial | Independent review:\allowbreak{} none | Full target solved:\allowbreak{} False


\clearpage

\section*{P0291  A nonmonotone Markov intensity with supermodular time ordering}

\entrylabel{Primary source}

Leskela:\allowbreak{} Ross second conjecture and supermodular stochastic ordering,\allowbreak{}2022;\allowbreak{} Problem3,\allowbreak{}journal215,\allowbreak{}PDF4


\url{https://acris.aalto.fi/ws/portalfiles/portal/89616571/Ross_s_second_conjecture_and_supermodular_stochastic_ordering.pdf}


Source statement date:\allowbreak{} 2022-05-01. Archived access:\allowbreak{} 2026-09-07T00:\allowbreak{}46:\allowbreak{}33Z.


Historical:\allowbreak{} source\_\allowbreak{}says\_\allowbreak{}open | Current:\allowbreak{} unverified | Result:\allowbreak{} unsolved


\entrylabel{Question}

Does a stationary ergodic CTMC on ordered intensity values exist that is not stochastically monotone but whose speed-scaled process decreases in supermodular stochastic order?


\entrylabel{Scope}

This is an existence question about the stronger process order,\allowbreak{} distinct from workload ordering. Discrete-time damped variants are not counted separately.


\subsection*{P0291-A001 | Explicit model calculation,\allowbreak{} comparison identities and remaining hypothesis obstruction}

\entrylabel{Argument or result}

A necessary condition follows by comparing a finite lag with the infinite-lag independent limit:\allowbreak{} Cov(g(X0),\allowbreak{}h(Xt))>=\allowbreak{}0 for every bounded increasing g,\allowbreak{}h and every t>=\allowbreak{}0. A simple nonmonotone candidate fails decisively. For the stationary clockwise three-cycle with rate-one jumps and state values0,\allowbreak{}1,\allowbreak{}2,\allowbreak{} the centered state observable has covariance (2/\allowbreak{}3)*exp(-3t/\allowbreak{}2)*cos(sqrt3*t/\allowbreak{}2). At t=\allowbreak{}2pi/\allowbreak{}sqrt3 this is negative,\allowbreak{} so that chain cannot be the desired example.


\entrylabel{Second check}

The two nonconstant eigenvalues of the three-cycle generator are -3/\allowbreak{}2 plus or minus i*sqrt3/\allowbreak{}2. The centered observable has variance2/\allowbreak{}3 and equal conjugate Fourier weights;\allowbreak{} this independently yields the displayed covariance and sign.


\entrylabel{Remaining obstruction}

Positive pairwise association would still be weaker than the required all-dimensional supermodular order. No nonmonotone chain satisfying the full order,\allowbreak{} or proof that none exists,\allowbreak{} has been found. Reset chains do satisfy the order but are already stochastically monotone and therefore do not answer the source question.


Verification:\allowbreak{} author checked partial | Independent review:\allowbreak{} none | Full target solved:\allowbreak{} False


\clearpage

\section*{P0292  Characterize CTMC speed monotonicity in supermodular order}

\entrylabel{Primary source}

Leskela:\allowbreak{} Ross second conjecture and supermodular stochastic ordering,\allowbreak{}2022;\allowbreak{} Problem4,\allowbreak{}journal215,\allowbreak{}PDF4


\url{https://acris.aalto.fi/ws/portalfiles/portal/89616571/Ross_s_second_conjecture_and_supermodular_stochastic_ordering.pdf}


Source statement date:\allowbreak{} 2022-05-01. Archived access:\allowbreak{} 2026-09-07T00:\allowbreak{}46:\allowbreak{}33Z.


Historical:\allowbreak{} source\_\allowbreak{}says\_\allowbreak{}open | Current:\allowbreak{} unverified | Result:\allowbreak{} unsolved


\entrylabel{Question}

Give a simple necessary-and-sufficient condition on a stationary ergodic CTMC for its speed-scaled process to decrease in supermodular stochastic order.


\entrylabel{Scope}

The existing doubly stochastic monotonicity condition is sufficient. Repeating it is not the requested characterization.


\subsection*{P0292-A001 | Explicit model calculation,\allowbreak{} comparison identities and remaining hypothesis obstruction}

\entrylabel{Argument or result}

Finite-dimensional supermodular order cannot be reduced to pairwise correlation alone. Let three fair bits have the uniform even-parity law on000,\allowbreak{}011,\allowbreak{}101,\allowbreak{}110. Every pair is independent,\allowbreak{} exactly as for three iid fair bits,\allowbreak{} but the supermodular function phi(x)=\allowbreak{}x1*x2*x3 has expectations0 and1/\allowbreak{}8 respectively. The odd-parity law has the same pairwise distributions and expectation1/\allowbreak{}4. Thus identical pairwise data permits opposite comparisons for a legitimate supermodular test.


\entrylabel{Second check}

On the Boolean cube,\allowbreak{} each mixed second difference of the product function equals the remaining coordinate and is nonnegative,\allowbreak{} so phi is supermodular. Enumerating the four parity configurations verifies every pair marginal and the three distinct product expectations.


\entrylabel{Remaining obstruction}

These are static diagnostic laws,\allowbreak{} not a CTMC counterexample. A necessary-and-sufficient Markov generator condition must control the entire sequence of transition kernels at all lags and all supermodular test functions. Covariance inequalities and a spectral-gap calculation leave those higher-dimensional constraints unresolved.


Verification:\allowbreak{} author checked partial | Independent review:\allowbreak{} none | Full target solved:\allowbreak{} False


\clearpage

\section*{P0293  Order one over N accuracy for least workload mean field queues}

\entrylabel{Primary source}

Hellemans-Van Houdt:\allowbreak{} Open problems in refined mean field approximations,\allowbreak{}2022;\allowbreak{} Numbered open problem1,\allowbreak{}PDF3;\allowbreak{}LL(d)modelPDF2


\url{https://win.uantwerpen.be/~vanhoudt/papers/2022/Questa_100.pdf}


Source statement date:\allowbreak{} 2022-02-02. Archived access:\allowbreak{} 2026-09-07T00:\allowbreak{}46:\allowbreak{}29Z.


Historical:\allowbreak{} source\_\allowbreak{}says\_\allowbreak{}open | Current:\allowbreak{} unverified | Result:\allowbreak{} unsolved


\entrylabel{Question}

Establish order1/\allowbreak{}N stationary mean-workload accuracy for the continuous-state least-loaded-of-d queueing model in the source and identify assumptions supporting a broader continuous-state theorem.


\entrylabel{Scope}

The target uses the source nontrivial interacting workload model. Trivial independent continuous-state examples do not resolve the intended extension;\allowbreak{} load and sampling convention must be fixed.


\subsection*{P0293-A001 | Explicit model calculation,\allowbreak{} comparison identities and remaining hypothesis obstruction}

\entrylabel{Argument or result}

Continuous state alone cannot guarantee1/\allowbreak{}N weak accuracy for arbitrary observables. If X\_\allowbreak{}i are iid standard normal and their empirical mean is M\_\allowbreak{}N,\allowbreak{} then E|M\_\allowbreak{}N|=\allowbreak{}sqrt(2/\allowbreak{}(pi*N)),\allowbreak{} whereas the limiting observable is0. The nonsmooth absolute value loses a power of sqrtN. In contrast,\allowbreak{} a twice differentiable test has its centered first Taylor term vanish and its expected second term of order1/\allowbreak{}N under a covariance bound. This locates the regularity input a weak-error proof needs.


\entrylabel{Second check}

The Gaussian example is exact because M\_\allowbreak{}N has variance1/\allowbreak{}N. In the queue model,\allowbreak{} the observable mean workload is linear,\allowbreak{} but the routing mechanism and the empty-workload boundary create nonsmooth dynamics;\allowbreak{} smoothness of the observable does not automatically make its Poisson solution smooth.


\entrylabel{Remaining obstruction}

One needs uniform stationary moment bounds,\allowbreak{} quantitative propagation of chaos and a sufficiently regular Poisson equation or semigroup derivative estimate for the unbounded workload state. The source simulations and the discrete-state refined theorem do not provide those estimates. No matching general continuous-workload proof is supplied.


Verification:\allowbreak{} author checked partial | Independent review:\allowbreak{} none | Full target solved:\allowbreak{} False


\clearpage

\section*{P0294  Compute a first order bias correction for continuous workload mean field models}

\entrylabel{Primary source}

Hellemans-Van Houdt:\allowbreak{} Open problems in refined mean field approximations,\allowbreak{}2022;\allowbreak{} Numbered open problem2,\allowbreak{}PDF3


\url{https://win.uantwerpen.be/~vanhoudt/papers/2022/Questa_100.pdf}


Source statement date:\allowbreak{} 2022-02-02. Archived access:\allowbreak{} 2026-09-07T00:\allowbreak{}46:\allowbreak{}29Z.


Historical:\allowbreak{} source\_\allowbreak{}says\_\allowbreak{}open | Current:\allowbreak{} unverified | Result:\allowbreak{} unsolved


\entrylabel{Question}

Develop an effective first-order stationary bias correction for continuous-state mean-field models such as least-loaded-of-d queues,\allowbreak{} with a justified improved error order.


\entrylabel{Scope}

This asks for a computable correction after the source accuracy question,\allowbreak{} not merely a fitted constant from simulations.


\subsection*{P0294-A001 | Explicit model calculation,\allowbreak{} comparison identities and remaining hypothesis obstruction}

\entrylabel{Argument or result}

An exactly solvable continuous-state control model shows how covariance enters the correction. Let dX\_\allowbreak{}i=\allowbreak{}(-a*X\_\allowbreak{}i+\allowbreak{}b*M\_\allowbreak{}N)dt+\allowbreak{}sigma*dB\_\allowbreak{}i with a>b>=\allowbreak{}0 and independent Brownian motions. The empirical mean is an OU process of rate a-b and variance sigma\textasciicircum{}2/\allowbreak{}[2(a-b)N]. Orthogonal deviations have marginal variance sigma\textasciicircum{}2*(1-1/\allowbreak{}N)/\allowbreak{}(2a). Therefore E X\_\allowbreak{}i\textasciicircum{}2=\allowbreak{}sigma\textasciicircum{}2/\allowbreak{}(2a)+\allowbreak{}(sigma\textasciicircum{}2/\allowbreak{}(2N))*(1/\allowbreak{}(a-b)-1/\allowbreak{}a),\allowbreak{} an exact first-order correction with no further error.


\entrylabel{Second check}

The orthogonal projection onto the constant vector and its complement diagonalizes the linear drift. Their independent Brownian projections give the two variances,\allowbreak{} and summing them reproduces the one-coordinate stationary variance directly from the Lyapunov equation.


\entrylabel{Remaining obstruction}

This classical Gaussian calculation is a baseline,\allowbreak{} not a solution for the least-workload PDMP. Its interaction involves routing minima and a reflecting or sticky empty boundary;\allowbreak{} the appropriate infinite-dimensional covariance and response equations must be derived and solved with certified truncation and tail error. The source empirical approximation lambda\textasciicircum{}2/\allowbreak{}[2(1-lambda)] is not asserted exact.


Verification:\allowbreak{} author checked partial | Independent review:\allowbreak{} none | Full target solved:\allowbreak{} False


\clearpage

\section*{P0295  Fluid limit with dependent service and patience marks}

\entrylabel{Primary source}

Moyal-Perry:\allowbreak{} Many server limits with dependent service and patience times,\allowbreak{}2022;\allowbreak{} Section2 FWLLN,\allowbreak{}PDF2;\allowbreak{}discussionPDF3


\url{https://users.iems.northwestern.edu/~perry/QueuesOpen.pdf}


Source statement date:\allowbreak{} 2022-02-11. Archived access:\allowbreak{} 2026-09-07T00:\allowbreak{}46:\allowbreak{}38Z.


Historical:\allowbreak{} source\_\allowbreak{}says\_\allowbreak{}open | Current:\allowbreak{} unverified | Result:\allowbreak{} unsolved


\entrylabel{Question}

Prove the functional fluid law of large numbers for the source GI/\allowbreak{}Gdep/\allowbreak{}n+\allowbreak{}Gdep queue when each customer has an iid jointly distributed service-patience pair,\allowbreak{} including uniqueness of the limiting integral equation.


\entrylabel{Scope}

The joint mark law is essential. Exponential marginals do not restore the independent-service martingale representation. Appropriate initial conditions and regularity remain part of the target.


\subsection*{P0295-A001 | Explicit model calculation,\allowbreak{} comparison identities and remaining hypothesis obstruction}

\entrylabel{Argument or result}

An explicit density shows the missing state information. Let patience D and E be independent mean-one exponentials and service X=\allowbreak{}D+\allowbreak{}E,\allowbreak{} with joint density f(x,\allowbreak{}y)=\allowbreak{}exp(-x) on x>y>0. A customer entering service after waiting w has X conditional on D>w distributed as w+\allowbreak{}Gamma(2,\allowbreak{}1). At service age a,\allowbreak{} its hazard is0 if a<w and (a-w)/\allowbreak{}(1+\allowbreak{}a-w) if a>w. Thus two customers of the same current service age can have different completion hazards because their entry waiting times differ.


\entrylabel{Second check}

Memorylessness gives D conditional on D>w equal to w plus a fresh exponential independent of E. The Gamma(2,\allowbreak{}1) survival function is (1+\allowbreak{}t)*exp(-t),\allowbreak{} yielding the stated hazard. The joint density integrates to1 and has precisely this conditional law.


\entrylabel{Remaining obstruction}

The unmarked service-age measure alone cannot close the dynamics. An augmented measure retaining entry-wait or residual-service marks can supply a martingale representation,\allowbreak{} but its tightness,\allowbreak{} boundary flux and uniqueness of the limit remain to be proved. A2025 unified fluid model lead concerns a directly constructed model and has not been verified as the requested stochastic limit theorem.


Verification:\allowbreak{} author checked partial | Independent review:\allowbreak{} none | Full target solved:\allowbreak{} False


\clearpage

\section*{P0296  Interchange stationary and many server limits with dependent marks}

\entrylabel{Primary source}

Moyal-Perry:\allowbreak{} Many server limits with dependent service and patience times,\allowbreak{}2022;\allowbreak{} Section2 stationary WLLN equation2,\allowbreak{}PDF2;\allowbreak{}discussionPDF3


\url{https://users.iems.northwestern.edu/~perry/QueuesOpen.pdf}


Source statement date:\allowbreak{} 2022-02-11. Archived access:\allowbreak{} 2026-09-07T00:\allowbreak{}46:\allowbreak{}38Z.


Historical:\allowbreak{} source\_\allowbreak{}says\_\allowbreak{}open | Current:\allowbreak{} unverified | Result:\allowbreak{} unsolved


\entrylabel{Question}

Under the source hypotheses and a unique appropriate stationary fluid state,\allowbreak{} prove that the stationary total customer count divided by n converges to that state in the dependent-service/\allowbreak{}patience queue.


\entrylabel{Scope}

The displayed overloaded fixed-point formula requires an admissible nonnegative waiting threshold. An underloaded regime in which that threshold does not exist is not silently forced into the formula.


\subsection*{P0296-A001 | Explicit model calculation,\allowbreak{} comparison identities and remaining hypothesis obstruction}

\entrylabel{Argument or result}

For the density f(x,\allowbreak{}y)=\allowbreak{}exp(-x) on x>y>0,\allowbreak{} E[X*1\{D>w\}]=\allowbreak{}(w+\allowbreak{}2)*exp(-w). The source capacity equation is therefore lambda*(w+\allowbreak{}2)*exp(-w)=\allowbreak{}1;\allowbreak{} a positive root exists when2lambda>1. More generally,\allowbreak{} uniqueness of a deterministic fixed point alone cannot justify stationary interchange. The planar drift b(x,\allowbreak{}y)=\allowbreak{}((1-r\textasciicircum{}2)x-y,\allowbreak{}(1-r\textasciicircum{}2)y+\allowbreak{}x) has only the equilibrium0 but has an attracting unit-circle orbit. Adding independent noise of size1/\allowbreak{}sqrtN gives invariant density proportional to exp[N*(r\textasciicircum{}2-r\textasciicircum{}4/\allowbreak{}2)],\allowbreak{} which concentrates on the circle as N grows,\allowbreak{} not at the unique equilibrium.


\entrylabel{Second check}

The drift matrix at a point is (1-r\textasciicircum{}2)I+\allowbreak{}J,\allowbreak{} whose determinant is (1-r\textasciicircum{}2)\textasciicircum{}2+\allowbreak{}1,\allowbreak{} so no nonzero fixed point exists. The rotational drift is divergence-free and orthogonal to the radial gradient,\allowbreak{} verifying the invariant density by the Fokker-Planck equation. Its exponent is maximized exactly at r=\allowbreak{}1.


\entrylabel{Remaining obstruction}

The diffusion is a logical counterexample to an insufficient interchange argument,\allowbreak{} not to this queue. A queue proof still needs stationary tightness and global attraction or a characterization of all invariant fluid laws,\allowbreak{} beyond uniqueness of a stationary point. Neither a finite-time FWLLN nor the scalar capacity equation alone supplies that conclusion.


Verification:\allowbreak{} author checked partial | Independent review:\allowbreak{} none | Full target solved:\allowbreak{} False


\clearpage

\section*{P0297  Exact Hausdorff dimension of Liouville quantum gravity}

\entrylabel{Primary source}

Rhodes-Vargas:\allowbreak{} Two decades of probabilistic approach to Liouville conformal field theory,\allowbreak{}2025;\allowbreak{} Section3,\allowbreak{}PDF4


\url{https://arxiv.org/pdf/2509.21053}


Source statement date:\allowbreak{} 2025-09-25. Archived access:\allowbreak{} 2026-09-07T00:\allowbreak{}40:\allowbreak{}35Z.


Historical:\allowbreak{} source\_\allowbreak{}says\_\allowbreak{}open | Current:\allowbreak{} unverified | Result:\allowbreak{} unsolved


\entrylabel{Question}

Determine the almost sure Hausdorff dimension d\_\allowbreak{}gamma of the gamma-LQG metric for general gamma in (0,\allowbreak{}2),\allowbreak{} beyond the known value4 at gamma=\allowbreak{}sqrt(8/\allowbreak{}3).


\entrylabel{Scope}

The parameter family is one question. The formal exponential of the GFF is not an ordinary smooth Riemannian metric.


\subsection*{P0297-A001 | Explicit calculation or obstruction with source-specific scope}

\entrylabel{Argument or result}

Constant shifts of the field multiply quantum area by exp(gamma*c) and metric distances by exp(xi*c),\allowbreak{} where xi is the metric scaling parameter. If a metric-volume exponent d is identified consistently with these scalings,\allowbreak{} the necessary relation is xi*d=\allowbreak{}gamma. This reproduces d=\allowbreak{}4 when xi=\allowbreak{}gamma/\allowbreak{}4 at the Brownian-map parameter. It does not calculate xi for other parameters.


\entrylabel{Second check}

The two scaling laws give the same volume rescaling only if exp(xi*d*c)=\allowbreak{}exp(gamma*c) for every real c. Merely inserting the smooth-metric exponent gamma/\allowbreak{}2 would force d=\allowbreak{}2 and contradict the known Brownian value.


\entrylabel{Remaining obstruction}

The missing quantity is the nontrivial first-passage metric normalization exponent,\allowbreak{} together with the theorem relating metric-ball volumes to it. Gaussian multiplicative chaos moments determine quantum area multifractality but do not supply the metric exponent. No new dimension formula is established.


Verification:\allowbreak{} author checked partial | Independent review:\allowbreak{} none | Full target solved:\allowbreak{} False


\clearpage

\section*{P0298  General Toda three-point structure constants}

\entrylabel{Primary source}

Rhodes-Vargas:\allowbreak{} Two decades of probabilistic approach to Liouville conformal field theory,\allowbreak{}2025;\allowbreak{} Section9,\allowbreak{}Path integrals for other CFTs,\allowbreak{}PDF18


\url{https://arxiv.org/pdf/2509.21053}


Source statement date:\allowbreak{} 2025-09-25. Archived access:\allowbreak{} 2026-09-07T00:\allowbreak{}40:\allowbreak{}35Z.


Historical:\allowbreak{} source\_\allowbreak{}says\_\allowbreak{}open | Current:\allowbreak{} unverified | Result:\allowbreak{} unsolved


\entrylabel{Question}

Determine the general three-point correlation function of probabilistically constructed Toda conformal field theory beyond the special Fateev-Litvinov parameter family.


\entrylabel{Scope}

The known sl3 special-family formula arXiv2208.12085 is excluded from the target. Lie-algebra and charge parameter copies are grouped.


\subsection*{P0298-A001 | Explicit calculation or obstruction with source-specific scope}

\entrylabel{Argument or result}

The constant-mode integral can be evaluated conditionally on the nonconstant field. Write the independent simple-root coordinates as y\_\allowbreak{}i and let the linear charge term be sum a\_\allowbreak{}i*y\_\allowbreak{}i. If a\_\allowbreak{}i>0 and M\_\allowbreak{}i>0,\allowbreak{} substituting u\_\allowbreak{}i=\allowbreak{}mu\_\allowbreak{}i*M\_\allowbreak{}i*exp(gamma*y\_\allowbreak{}i) gives integral exp(sum a\_\allowbreak{}i*y\_\allowbreak{}i-sum mu\_\allowbreak{}i*M\_\allowbreak{}i*exp(gamma*y\_\allowbreak{}i))dy =\allowbreak{} gamma\textasciicircum{}(-r)*product[(mu\_\allowbreak{}i*M\_\allowbreak{}i)\textasciicircum{}(-a\_\allowbreak{}i/\allowbreak{}gamma)*Gamma(a\_\allowbreak{}i/\allowbreak{}gamma)]. A fixed root-coordinate determinant multiplies this expression when integrating in an orthonormal basis.


\entrylabel{Second check}

The Jacobian is gamma\textasciicircum{}(-r)*product du\_\allowbreak{}i/\allowbreak{}u\_\allowbreak{}i;\allowbreak{} the remaining integrals are independent Gamma integrals. This is an exact finite-dimensional integration conditional on the chaos masses,\allowbreak{} not an independence assertion about those masses.


\entrylabel{Remaining obstruction}

The remaining expectation is a joint negative moment of correlated GMC integrals carrying the three insertions. Their joint law does not factor because the root fields share the vector GFF. Existing semi-degenerate identities do not determine this general joint moment. No general structure-constant formula is obtained.


Verification:\allowbreak{} author checked partial | Independent review:\allowbreak{} none | Full target solved:\allowbreak{} False


\clearpage

\section*{P0299  Probabilistic conformal bootstrap for Toda theory}

\entrylabel{Primary source}

Rhodes-Vargas:\allowbreak{} Two decades of probabilistic approach to Liouville conformal field theory,\allowbreak{}2025;\allowbreak{} Section9,\allowbreak{}Path integrals for other CFTs,\allowbreak{}PDF18


\url{https://arxiv.org/pdf/2509.21053}


Source statement date:\allowbreak{} 2025-09-25. Archived access:\allowbreak{} 2026-09-07T00:\allowbreak{}40:\allowbreak{}35Z.


Historical:\allowbreak{} source\_\allowbreak{}says\_\allowbreak{}open | Current:\allowbreak{} unverified | Result:\allowbreak{} unsolved


\entrylabel{Question}

Establish conformal bootstrap identities for the probabilistic Toda CFT,\allowbreak{} including the spectral decomposition and sewing of correlation functions.


\entrylabel{Scope}

This is the correlation-factorization and completeness theorem,\allowbreak{} distinct from merely finding a three-point formula. Existing path-integral constructions and Ward identities are not counted as its solution.


\subsection*{P0299-A001 | Explicit calculation or obstruction with source-specific scope}

\entrylabel{Argument or result}

An annulus semigroup with a self-adjoint generator would yield an abstract spectral resolution,\allowbreak{} but the spectral theorem alone does not identify the multiplicity or the W-algebra descendant states required in a bootstrap integral. A diagnostic is multiplication by x on L2([0,\allowbreak{}1]) versus multiplication by x on L2([0,\allowbreak{}1]) direct-summed with itself:\allowbreak{} both have spectrum [0,\allowbreak{}1],\allowbreak{} identical semigroup norm exp(0)=\allowbreak{}1 for positive times,\allowbreak{} and different spectral multiplicity.


\entrylabel{Second check}

The second operator has two orthogonal spectral channels over every interval of positive measure. Equality of spectral sets and operator norms therefore cannot imply a complete one-channel expansion.


\entrylabel{Remaining obstruction}

A proof needs the actual Toda scattering transform,\allowbreak{} its Plancherel measure and multiplicities,\allowbreak{} identification of descendants,\allowbreak{} and justified sewing integrals. Neither an abstract existence theorem for the path integral nor a conjectural structure constant resolves these completeness and integrability requirements.


Verification:\allowbreak{} author checked partial | Independent review:\allowbreak{} none | Full target solved:\allowbreak{} False


\clearpage

\section*{P0300  Conformal bootstrap for the noncompact H3 WZW model}

\entrylabel{Primary source}

Rhodes-Vargas:\allowbreak{} Two decades of probabilistic approach to Liouville conformal field theory,\allowbreak{}2025;\allowbreak{} Section9,\allowbreak{}Path integrals for other CFTs,\allowbreak{}PDF18


\url{https://arxiv.org/pdf/2509.21053}


Source statement date:\allowbreak{} 2025-09-25. Archived access:\allowbreak{} 2026-09-07T00:\allowbreak{}40:\allowbreak{}35Z.


Historical:\allowbreak{} source\_\allowbreak{}says\_\allowbreak{}open | Current:\allowbreak{} unverified | Result:\allowbreak{} unsolved


\entrylabel{Question}

Implement a rigorous conformal bootstrap for the noncompact H3=\allowbreak{}SL(2,\allowbreak{}C)/\allowbreak{}SU(2) WZW model constructed from its probabilistic path integral.


\entrylabel{Scope}

The compact-group VOA bootstrap and the H3 path-integral existence theorem are established starting points. Other noncompact groups are not separately counted.


\subsection*{P0300-A001 | Explicit calculation or obstruction with source-specific scope}

\entrylabel{Argument or result}

A compact-target trace argument fails already for the zero-mode heat semigroup on H3. Homogeneity makes its heat-kernel diagonal p\_\allowbreak{}t(x,\allowbreak{}x) constant and strictly positive for fixed t>0. Since H3 has infinite Riemannian volume,\allowbreak{} integral p\_\allowbreak{}t(x,\allowbreak{}x)dx is infinite,\allowbreak{} so the heat semigroup is not trace class.


\entrylabel{Second check}

The diagonal constancy follows by applying any transitive isometry to the Laplacian;\allowbreak{} integrating the same positive value over increasing finite-volume sets diverges. This is a baseline obstruction,\allowbreak{} not a claim about every regularized WZW amplitude.


\entrylabel{Remaining obstruction}

The noncompact model requires a continuous spectral measure and suitable localized or insertion-weighted amplitudes. A discrete trace expansion borrowed from compact targets is unjustified. The required affine-symmetry scattering and sewing analysis remains unsolved here.


Verification:\allowbreak{} author checked partial | Independent review:\allowbreak{} none | Full target solved:\allowbreak{} False


\clearpage

\section*{P0301  Full probabilistic noncompact imaginary Liouville theory}

\entrylabel{Primary source}

Rhodes-Vargas:\allowbreak{} Two decades of probabilistic approach to Liouville conformal field theory,\allowbreak{}2025;\allowbreak{} Section9,\allowbreak{}PDF18;\allowbreak{}Usciati et al2505.09390,\allowbreak{}PDF3-4


\url{https://arxiv.org/pdf/2509.21053}


Source statement date:\allowbreak{} 2025-09-25. Archived access:\allowbreak{} 2026-09-07T00:\allowbreak{}40:\allowbreak{}35Z.


Historical:\allowbreak{} source\_\allowbreak{}says\_\allowbreak{}open | Current:\allowbreak{} unverified | Result:\allowbreak{} unsolved


\entrylabel{Question}

Complete the probabilistic construction and identification of noncompact imaginary Liouville theory in the general non-neutral regime described by the survey.


\entrylabel{Scope}

Usciati-Guillarmou-Rhodes-Santachiara give a proposed construction and exact circle results;\allowbreak{} their sphere construction depends on conjectured imaginary-chaos Laplace asymptotics. Chatterjee2504.02348 proves timelike results under charge neutrality. These partial results must not be erased or counted as new.


\subsection*{P0301-A001 | Explicit calculation or obstruction with source-specific scope}

\entrylabel{Argument or result}

The circle contour calculation shows why simply integrating the constant mode over a real period imposes neutrality. For an integer screening expansion,\allowbreak{} the integral of exp(i*k*beta*c) over a full period vanishes unless k=\allowbreak{}0. A noncompact complex Hankel-type contour can avoid that selection rule,\allowbreak{} but requires a large-argument estimate for the complex GMC Laplace transform along its two infinite legs. Finite exponential moments only ensure finiteness for each fixed argument,\allowbreak{} not decay along those legs.


\entrylabel{Second check}

For a deterministic complex random variable M=\allowbreak{}-1,\allowbreak{} every exponential moment is finite but E exp(-z*M)=\allowbreak{}exp(z) grows on the positive real ray. Thus existence of all moments cannot provide the needed contour decay. This diagnostic does not assert that the actual imaginary GMC has this behavior.


\entrylabel{Remaining obstruction}

The missing sphere Laplace estimate and identification with the imaginary DOZZ expression remain unresolved in the read source. Current status requires later papers beyond the read2025 proposal and April2026 charge-neutral result;\allowbreak{} it is not certified open solely from the survey.


Verification:\allowbreak{} author checked partial | Independent review:\allowbreak{} none | Full target solved:\allowbreak{} False


\clearpage

\section*{P0302  One critically loaded queue under general k-limited polling inputs}

\entrylabel{Primary source}

Boon and Winands:\allowbreak{} Open problems for critically loaded k-limited polling systems (2022);\allowbreak{} Conjecture1 and Problem1,\allowbreak{}journal282,\allowbreak{}PDF3


\url{https://pure.uva.nl/ws/files/120928955/Open_problems_for_critically_loaded_k_limited_polling_systems.pdf}


Source statement date:\allowbreak{} 2022-03-29. Archived access:\allowbreak{} 2026-09-07T01:\allowbreak{}25:\allowbreak{}41Z.


Historical:\allowbreak{} source\_\allowbreak{}says\_\allowbreak{}open | Current:\allowbreak{} unverified | Result:\allowbreak{} unsolved


\entrylabel{Question}

Prove the general finite-variance renewal-input heavy-traffic joint limit in Boon-Winands Conjecture1:\allowbreak{} one critically loaded queue has the stated scaled exponential limit,\allowbreak{} the other converges to the stated vacation queue,\allowbreak{} and the components become independent.


\entrylabel{Scope}

The exponential and phase-type special cases are already established. The source regime lambda1/\allowbreak{}k1<lambda2/\allowbreak{}k2 and scaling with fixed interarrival coefficient of variation are essential.


\subsection*{P0302-A001 | Explicit calculation or obstruction with source-specific scope}

\entrylabel{Argument or result}

The source switch-over variance absorption has a useful exact batch check. When the critical queue serves exactly k2 jobs per cycle,\allowbreak{} its batch time is sum of k2 iid B2 variables plus the total switch-over S. The batch variance divided by k2 is Var(B2)+\allowbreak{}Var(S)/\allowbreak{}k2. Assigning the same S/\allowbreak{}k2 to every effective job instead would give marginal variance Var(B2)+\allowbreak{}Var(S)/\allowbreak{}k2\textasciicircum{}2 but introduce correlations between those jobs. Ignoring those correlations would produce the wrong diffusion coefficient.


\entrylabel{Second check}

The variance of the sum of the correlated effective jobs contains k2*(k2-1) covariance terms Var(S)/\allowbreak{}k2\textasciicircum{}2 and exactly restores Var(S). The mean effective service is E B2+\allowbreak{}E S/\allowbreak{}k2,\allowbreak{} agreeing with the source.


\entrylabel{Remaining obstruction}

This verifies the proposed effective mean and variance but not the joint stationary limit. General renewal residual lives,\allowbreak{} the stable-queue vacation environment,\allowbreak{} asymptotic independence and stationary interchange still need a proof. The source simulation correlation near zero is not a proof of independence.


Verification:\allowbreak{} author checked partial | Independent review:\allowbreak{} none | Full target solved:\allowbreak{} False


\clearpage

\section*{P0303  Simultaneous criticality in k-limited polling}

\entrylabel{Primary source}

Boon and Winands:\allowbreak{} Open problems for critically loaded k-limited polling systems (2022);\allowbreak{} Problem2,\allowbreak{}journal283,\allowbreak{}PDF4


\url{https://pure.uva.nl/ws/files/120928955/Open_problems_for_critically_loaded_k_limited_polling_systems.pdf}


Source statement date:\allowbreak{} 2022-03-29. Archived access:\allowbreak{} 2026-09-07T01:\allowbreak{}25:\allowbreak{}41Z.


Historical:\allowbreak{} source\_\allowbreak{}says\_\allowbreak{}open | Current:\allowbreak{} unverified | Result:\allowbreak{} unsolved


\entrylabel{Question}

Determine the heavy-traffic joint queue-length distribution when both queues in a k-limited polling system become critically loaded simultaneously.


\entrylabel{Scope}

The two-queue source problem is one entry;\allowbreak{} extensions to more queues and parameter choices are grouped. The one-critical-queue theorem does not apply at equality of the load ratios.


\subsection*{P0303-A001 | Explicit calculation or obstruction with source-specific scope}

\entrylabel{Argument or result}

The no-switch-over comparison has an exact simplification:\allowbreak{} if all jobs require deterministic time2 and arrivals are the superposition of independent Poisson streams,\allowbreak{} total workload is the workload of a work-conserving M/\allowbreak{}D/\allowbreak{}1 queue,\allowbreak{} independently of the polling order. Introducing deterministic switch-over times destroys work conservation even when the sum of a service time and a switch time remains2. Empty visits then consume capacity,\allowbreak{} and their boundary frequencies enter the two-dimensional reflected limit.


\entrylabel{Second check}

Total unfinished service decreases at unit rate whenever any work is present in the first system. In the second it is constant during each switch,\allowbreak{} including switches caused by an empty queue. Thus equal nominal time per saturated visit does not imply equal sample-path workload.


\entrylabel{Remaining obstruction}

The missing boundary occupation and reflection structure prevent identifying the limiting joint law from a scalar workload formula. The source suggests a two-dimensional transform boundary-value problem. No solution of that problem or valid interchange of its heavy-traffic limit is supplied.


Verification:\allowbreak{} author checked partial | Independent review:\allowbreak{} none | Full target solved:\allowbreak{} False


\clearpage

\section*{P0305  Hausdorff dimensions of bulk LQG geodesic star sets}

\entrylabel{Primary source}

Bhatia Gwynne Harper:\allowbreak{} Dimension lower bounds in random geometry via Lipschitz functions (2026);\allowbreak{} Introduction PDF2;\allowbreak{}discussionPDF6;\allowbreak{}Theorem1.2 baseline


\url{https://arxiv.org/pdf/2606.10496}


Source statement date:\allowbreak{} 2026-06-09. Archived access:\allowbreak{} 2026-09-07T01:\allowbreak{}25:\allowbreak{}42Z.


Historical:\allowbreak{} source\_\allowbreak{}says\_\allowbreak{}open | Current:\allowbreak{} unverified | Result:\allowbreak{} unsolved


\entrylabel{Question}

Determine the Hausdorff dimensions of bulk k-star sets for general gamma-LQG;\allowbreak{} in particular prove or refute the conjectured metric dimension2 of the three-star set.


\entrylabel{Scope}

The different k values are grouped into a single source family. Brownian-map dimensions5-k for k<=\allowbreak{}5 and the2026 lower bound2 for three-stars are known results,\allowbreak{} not new solutions.


\subsection*{P0305-A001 | Explicit calculation or obstruction with source-specific scope}

\entrylabel{Argument or result}

A difference-of-distance map f(x)=\allowbreak{}(D(x,\allowbreak{}a)-D(x,\allowbreak{}b),\allowbreak{}D(x,\allowbreak{}a)-D(x,\allowbreak{}c)) is2-Lipschitz in each coordinate. The source uses non-constancy of such maps to get dimension lower bounds. This mechanism supplies no upper bound:\allowbreak{} the coordinate projection from a square to a line is Lipschitz and has a one-dimensional image,\allowbreak{} while its non-constancy set is the entire two-dimensional square.


\entrylabel{Second check}

The triangle inequality bounds each distance difference by2D(x,\allowbreak{}y). For f(x,\allowbreak{}y)=\allowbreak{}x on [0,\allowbreak{}1]\textasciicircum{}2,\allowbreak{} every relative neighborhood contains different x values,\allowbreak{} so its non-constancy set is the whole square despite a one-dimensional image.


\entrylabel{Remaining obstruction}

An upper bound needs quantitative control of fibers or a covering estimate for points supporting multiple noncoalesced geodesics. Lipschitz mapping and topological confluence alone do not provide it. No matching upper bound for general gamma is proved.


Verification:\allowbreak{} author checked partial | Independent review:\allowbreak{} none | Full target solved:\allowbreak{} False


\clearpage

\section*{P0306  One-dimensional boundary two-star set in LQG}

\entrylabel{Primary source}

Bhatia Gwynne Harper:\allowbreak{} Dimension lower bounds in random geometry via Lipschitz functions (2026);\allowbreak{} Conjecture1.4,\allowbreak{}PDF7;\allowbreak{}Theorem1.3 lower boundPDF6


\url{https://arxiv.org/pdf/2606.10496}


Source statement date:\allowbreak{} 2026-06-09. Archived access:\allowbreak{} 2026-09-07T01:\allowbreak{}25:\allowbreak{}42Z.


Historical:\allowbreak{} source\_\allowbreak{}says\_\allowbreak{}open | Current:\allowbreak{} unverified | Result:\allowbreak{} unsolved


\entrylabel{Question}

For the gamma-LQG metric associated with a Neumann GFF on the upper half plane,\allowbreak{} is the metric Hausdorff dimension of the boundary two-star set exactly1?


\entrylabel{Scope}

The boundary intersection uses the ambient LQG metric restricted to the boundary,\allowbreak{} not an intrinsic boundary path metric. The established lower bound1 is not the requested equality.


\subsection*{P0306-A001 | Explicit calculation or obstruction with source-specific scope}

\entrylabel{Argument or result}

Topological one-dimensionality of the boundary cannot prove this metric upper bound. On [0,\allowbreak{}1],\allowbreak{} the snowflake metric D\_\allowbreak{}alpha(x,\allowbreak{}y)=\allowbreak{}|x-y|\textasciicircum{}alpha for0<alpha<1 induces the usual topology but has Hausdorff dimension1/\allowbreak{}alpha. A cover with Euclidean diameters r\_\allowbreak{}j has D\_\allowbreak{}alpha diameters r\_\allowbreak{}j\textasciicircum{}alpha,\allowbreak{} making the s-dimensional metric Hausdorff sum the alpha*s-dimensional Euclidean sum.


\entrylabel{Second check}

The equality of the covering sums proves dim\_\allowbreak{}Dalpha([0,\allowbreak{}1])=\allowbreak{}1/\allowbreak{}alpha directly. The snowflake is only a diagnostic for the restriction metric;\allowbreak{} it is not asserted to be an LQG surface or a geodesic metric.


\entrylabel{Remaining obstruction}

One needs a covering argument specifically for the random boundary two-star subset in the restricted LQG metric. The source lower bound and the fact that the boundary is a topological line do not establish the conjectured upper bound.


Verification:\allowbreak{} author checked partial | Independent review:\allowbreak{} none | Full target solved:\allowbreak{} False


\clearpage

\section*{P0307  Exact dimension of the LQG metric net}

\entrylabel{Primary source}

Bhatia Gwynne Harper:\allowbreak{} Dimension lower bounds in random geometry via Lipschitz functions (2026);\allowbreak{} IntroductionPDF3;\allowbreak{}discussionPDF6 afterTheorem1.3


\url{https://arxiv.org/pdf/2606.10496}


Source statement date:\allowbreak{} 2026-06-09. Archived access:\allowbreak{} 2026-09-07T01:\allowbreak{}25:\allowbreak{}42Z.


Historical:\allowbreak{} source\_\allowbreak{}says\_\allowbreak{}open | Current:\allowbreak{} unverified | Result:\allowbreak{} unsolved


\entrylabel{Question}

Determine the Hausdorff dimension of the LQG metric net;\allowbreak{} in particular establish the expected metric dimension3 in the Brownian-map case.


\entrylabel{Scope}

One metric-net dimension family is counted. Dimensions of geodesic stars and of intersections of two nets are different targets. The2026 source lower bounds are not claimed as new.


\subsection*{P0307-A001 | Explicit calculation or obstruction with source-specific scope}

\entrylabel{Argument or result}

A fixed metric sphere and the union of spheres over all radii have different dimension behavior. In the Euclidean plane,\allowbreak{} each circle about the origin has Hausdorff dimension1 while their union over radii in[1,\allowbreak{}2] is an annulus of dimension2. Thus even a sharp dimension theorem for each outer boundary cannot be passed to the metric net by countable stability:\allowbreak{} the union is uncountable.


\entrylabel{Second check}

The annulus contains a nonempty open set,\allowbreak{} proving its dimension2;\allowbreak{} each positive-radius circle is a smooth one-dimensional manifold. Countable unions would preserve the supremum of dimensions,\allowbreak{} so the radii parameter cannot be silently replaced by a countable dense set.


\entrylabel{Remaining obstruction}

For LQG one needs joint control of the outer boundaries as radius varies,\allowbreak{} including the jumps created by swallowed components. The source non-constancy representation supplies a lower bound,\allowbreak{} but not the required upper covering or exact law. The dimension3 statement remains unproved here.


Verification:\allowbreak{} author checked partial | Independent review:\allowbreak{} none | Full target solved:\allowbreak{} False


\clearpage

\section*{P0308  Canonical and grand-canonical equivalence for long-range Riesz gases}

\entrylabel{Primary source}

Lewin:\allowbreak{} Coulomb and Riesz gases,\allowbreak{} the known and the unknown,\allowbreak{}2022;\allowbreak{} SectionIII,\allowbreak{}PDF27-28


\url{https://arxiv.org/pdf/2202.09240}


Source statement date:\allowbreak{} 2022-04-30. Archived access:\allowbreak{} 2026-09-07T00:\allowbreak{}34:\allowbreak{}23Z.


Historical:\allowbreak{} source\_\allowbreak{}says\_\allowbreak{}open | Current:\allowbreak{} unverified | Result:\allowbreak{} unsolved


\entrylabel{Question}

Prove the conjectured equivalence of canonical and grand-canonical thermodynamic free energies for neutral Riesz gases throughout the remaining positive long-range exponents0<s<d,\allowbreak{} beyond the known s>=\allowbreak{}d-2 regime.


\entrylabel{Scope}

All remaining exponents are grouped. Nonpositive exponents can make the grand-canonical energy unbounded below and are outside this formulation. The source discusses a possible route when s>=\allowbreak{}d-4.


\subsection*{P0308-A001 | Source-specific calculation,\allowbreak{} conditional implication or exact obstruction}

\entrylabel{Argument or result}

The grand partition function contains its exactly neutral canonical summand,\allowbreak{} giving one free-energy inequality immediately. Concentration of the particle number around neutrality does not give the opposite inequality. For example,\allowbreak{} a law supported equally at N-floor(sqrt(N)) and N+\allowbreak{}floor(sqrt(N)) has particle number divided by N tending to1 but assigns probability zero to exactly N.


\entrylabel{Second check}

The relative deviation is at most1/\allowbreak{}sqrt(N);\allowbreak{} nevertheless its central mass is identically zero. Thus a law-of-large-numbers neutrality argument supplies no logarithmic lower bound on the exact neutral sector.


\entrylabel{Remaining obstruction}

A proof needs uniform energy and entropy control when inserting or removing the subextensive excess charge,\allowbreak{} or a local partition-function comparison. Long-range background energy makes that repair nonlocal. No reverse free-energy inequality is proved here.


Verification:\allowbreak{} author checked partial | Independent review:\allowbreak{} none | Full target solved:\allowbreak{} False


\clearpage

\section*{P0309  Positive-temperature crystalline Gibbs phases for three-dimensional Riesz gases}

\entrylabel{Primary source}

Lewin:\allowbreak{} Coulomb and Riesz gases,\allowbreak{} the known and the unknown,\allowbreak{}2022;\allowbreak{} SectionVI,\allowbreak{}PDF59-61,\allowbreak{}66-67


\url{https://arxiv.org/pdf/2202.09240}


Source statement date:\allowbreak{} 2022-04-30. Archived access:\allowbreak{} 2026-09-07T00:\allowbreak{}34:\allowbreak{}23Z.


Historical:\allowbreak{} source\_\allowbreak{}says\_\allowbreak{}open | Current:\allowbreak{} unverified | Result:\allowbreak{} unsolved


\entrylabel{Question}

Prove existence of a crystalline Gibbs phase at sufficiently low positive temperature for the three-dimensional Coulomb or Riesz gas in the regimes predicted in the review.


\entrylabel{Scope}

This is the positive-temperature phase-existence question;\allowbreak{} the lattice-energy minimization problem at zero temperature and individual numerical transition constants are not separate entries.


\subsection*{P0309-A001 | Source-specific calculation,\allowbreak{} conditional implication or exact obstruction}

\entrylabel{Argument or result}

The harmonic phonon calculation around an assumed crystal has infrared displacement variance proportional to integral over |k|<1 of |k|\textasciicircum{}-2 dk. In dimension d this is a constant times integral from0 to1 of r\textasciicircum{}(d-3)dr,\allowbreak{} finite in d=\allowbreak{}3 and logarithmically divergent in d=\allowbreak{}2. Therefore the elementary infrared objection to positional order is absent in three dimensions.


\entrylabel{Second check}

The radial Jacobian r\textasciicircum{}(d-1) produces the exponent d-3. This only checks the Gaussian approximation conditional on a lattice;\allowbreak{} it neither selects a lattice nor constructs an infinite-volume Gibbs measure.


\entrylabel{Remaining obstruction}

One must control defects,\allowbreak{} anharmonic fluctuations and alternative configurations uniformly in volume and show nontrivial positional order. Finite variance within an assumed crystal cannot prove that the full Gibbs law is crystalline. Numerical BCC/\allowbreak{}FCC phase diagrams are not a proof.


Verification:\allowbreak{} author checked partial | Independent review:\allowbreak{} none | Full target solved:\allowbreak{} False


\clearpage

\section*{P0310  Uniqueness for the one-dimensional short-range Riesz Gibbs gas}

\entrylabel{Primary source}

Lewin:\allowbreak{} Coulomb and Riesz gases,\allowbreak{} the known and the unknown,\allowbreak{}2022;\allowbreak{} SectionVI,\allowbreak{}PDF61;\allowbreak{}Boursier Section1.3,\allowbreak{}PDF5


\url{https://arxiv.org/pdf/2202.09240}


Source statement date:\allowbreak{} 2022-04-30. Archived access:\allowbreak{} 2026-09-07T00:\allowbreak{}34:\allowbreak{}23Z.


Historical:\allowbreak{} source\_\allowbreak{}says\_\allowbreak{}open | Current:\allowbreak{} unverified | Result:\allowbreak{} unsolved


\entrylabel{Question}

Establish the expected uniqueness of the one-dimensional short-range Riesz Gibbs state in the unresolved exponent range1<s<=\allowbreak{}2.


\entrylabel{Scope}

The s>2 uniqueness regime is already known. Boursier proves a unique thermodynamic limit of the circular finite-volume laws for s>1,\allowbreak{} which is not by itself uniqueness among all DLR Gibbs states.


\subsection*{P0310-A001 | Source-specific calculation,\allowbreak{} conditional implication or exact obstruction}

\entrylabel{Argument or result}

A lattice interface calculation separates summability of a pair potential from summability of its boundary influence. For potential r\textasciicircum{}-s,\allowbreak{} its sum over r is finite when s>1,\allowbreak{} whereas the total weight of pairs crossing one cut is sum over r of r*r\textasciicircum{}-s,\allowbreak{} which diverges when s<=\allowbreak{}2. Hence the usual finite-boundary-energy argument cannot simply use integrability of the potential.


\entrylabel{Second check}

At distance r there are exactly r integer pairs crossing the cut between0 and1,\allowbreak{} so the interface series is sum r\textasciicircum{}(1-s). This is an exact lattice diagnostic,\allowbreak{} not an assertion of lattice constraints on the continuum gas.


\entrylabel{Remaining obstruction}

A uniqueness proof must exploit convex repulsion or finer cancellation beyond finite interface cost,\allowbreak{} and control all boundary conditions. Uniqueness of a particular periodic approximation does not exclude other infinite-volume specifications.


Verification:\allowbreak{} author checked partial | Independent review:\allowbreak{} none | Full target solved:\allowbreak{} False


\clearpage

\section*{P0311  Optimal truncated pair-correlation decay in the Riesz fluid phase}

\entrylabel{Primary source}

Lewin:\allowbreak{} Coulomb and Riesz gases,\allowbreak{} the known and the unknown,\allowbreak{}2022;\allowbreak{} SectionVI,\allowbreak{}PDF63,\allowbreak{}67,\allowbreak{}69


\url{https://arxiv.org/pdf/2202.09240}


Source statement date:\allowbreak{} 2022-04-30. Archived access:\allowbreak{} 2026-09-07T00:\allowbreak{}34:\allowbreak{}23Z.


Historical:\allowbreak{} source\_\allowbreak{}says\_\allowbreak{}open | Current:\allowbreak{} unverified | Result:\allowbreak{} unsolved


\entrylabel{Question}

Prove the predicted fluid-phase truncated pair-correlation tail of order |x|\textasciicircum{}(-(2d-s)) for long-range Riesz gases,\allowbreak{} with the exceptional Coulomb-type exponents and phase restrictions treated as in the source.


\entrylabel{Scope}

One tail family is counted. Correlation of ordered gaps in the one-dimensional circular gas is a different observable and does not automatically prove this point-density asymptotic. Current openness is not certified for every parameter.


\subsection*{P0311-A001 | Source-specific calculation,\allowbreak{} conditional implication or exact obstruction}

\entrylabel{Argument or result}

The formal random-phase closure gives structure factor S(k)=\allowbreak{}1/\allowbreak{}(1+\allowbreak{}beta*rho*Vhat(k)). Since the Riesz Fourier transform scales as |k|\textasciicircum{}(s-d),\allowbreak{} this predicts S(k) of order |k|\textasciicircum{}(d-s) near0. Fourier homogeneity then predicts a real-space correction with degree -(2d-s).


\entrylabel{Second check}

Under k->a*k,\allowbreak{} the inverse Fourier transform of a homogeneous degree alpha symbol has degree -(d+\allowbreak{}alpha);\allowbreak{} substituting alpha=\allowbreak{}d-s gives -(2d-s). Polynomial/\allowbreak{}even-integer symbols can instead contribute distributions at the origin,\allowbreak{} explaining why exceptional exponents cannot be ignored.


\entrylabel{Remaining obstruction}

The closure replaces higher correlations by a Gaussian response and has no error estimate here. Establishing its leading nonanalytic coefficient,\allowbreak{} excluding crystalline oscillations and passing to a pointwise tail remain unresolved. The scaling calculation is a heuristic conditional implication.


Verification:\allowbreak{} author checked partial | Independent review:\allowbreak{} none | Full target solved:\allowbreak{} False


\clearpage

\section*{P0312  Infinite-volume BBGKY equation for the general sine-beta process}

\entrylabel{Primary source}

Lewin:\allowbreak{} Coulomb and Riesz gases,\allowbreak{} the known and the unknown,\allowbreak{}2022;\allowbreak{} Equation147 and discussion,\allowbreak{}PDF63


\url{https://arxiv.org/pdf/2202.09240}


Source statement date:\allowbreak{} 2022-04-30. Archived access:\allowbreak{} 2026-09-07T00:\allowbreak{}34:\allowbreak{}23Z.


Historical:\allowbreak{} source\_\allowbreak{}says\_\allowbreak{}open | Current:\allowbreak{} unverified | Result:\allowbreak{} unsolved


\entrylabel{Question}

Justify the renormalized infinite-volume BBGKY equation described in the review for the sine-beta process at general beta.


\entrylabel{Scope}

The source specifically distinguishes known DLR equations from the still-unjustified BBGKY passage. Later stochastic-dynamics work may overlap;\allowbreak{} current status remains unverified until its hypotheses are compared.


\subsection*{P0312-A001 | Source-specific calculation,\allowbreak{} conditional implication or exact obstruction}

\entrylabel{Argument or result}

Finite-volume integration by parts produces the force sum over other particles of1/\allowbreak{}(x-y). At unit density,\allowbreak{} the absolute far-field integral is integral from1 toR of dy/\allowbreak{}y=\allowbreak{}log R. Thus a bound obtained by taking absolute values before the infinite-volume limit cannot be uniform. A symmetric principal value may cancel the mean term,\allowbreak{} but then fluctuations and truncated correlations must be controlled.


\entrylabel{Second check}

The two-sided integral of |1/\allowbreak{}y| doubles the logarithmic divergence;\allowbreak{} the symmetric integral of1/\allowbreak{}y itself is zero. This exhibits exactly the cancellation destroyed by an absolute-value estimate.


\entrylabel{Remaining obstruction}

The missing step is uniform integrability of the renormalized force and justification of differentiation under the infinite conditional law. Existence of a DLR density does not alone justify that interchange. No general BBGKY identity is established here.


Verification:\allowbreak{} author checked partial | Independent review:\allowbreak{} none | Full target solved:\allowbreak{} False


\clearpage

\section*{P0313  Translation invariance in low-dimensional long-range Riesz gases}

\entrylabel{Primary source}

Lewin:\allowbreak{} Coulomb and Riesz gases,\allowbreak{} the known and the unknown,\allowbreak{}2022;\allowbreak{} SectionVI,\allowbreak{}PDF63


\url{https://arxiv.org/pdf/2202.09240}


Source statement date:\allowbreak{} 2022-04-30. Archived access:\allowbreak{} 2026-09-07T00:\allowbreak{}34:\allowbreak{}23Z.


Historical:\allowbreak{} source\_\allowbreak{}says\_\allowbreak{}open | Current:\allowbreak{} unverified | Result:\allowbreak{} unsolved


\entrylabel{Question}

Prove the expected absence of translation-symmetry breaking for one-dimensional Riesz gases0<=\allowbreak{}s<1 and two-dimensional Riesz gases1<=\allowbreak{}s<2 in the regimes discussed in the review.


\entrylabel{Scope}

The dimensions and exponent ranges are grouped as the source translation-invariance question. Leble2024 proves translation invariance for local bulk limits of the two-dimensional Coulomb gas s=\allowbreak{}0,\allowbreak{} a separate established case.


\subsection*{P0313-A001 | Source-specific calculation,\allowbreak{} conditional implication or exact obstruction}

\entrylabel{Argument or result}

A naive Mermin-Wagner deformation estimate based on absolute Hessians is divergent. For V(r)=\allowbreak{}r\textasciicircum{}-s,\allowbreak{} the Hessian has size of order r\textasciicircum{}-s-2. Multiplication by a displacement difference bounded crudely by a constant times r gives the second-order energy integral of order integral r\textasciicircum{}(d-1)*r\textasciicircum{}2*r\textasciicircum{}-s-2 dr =\allowbreak{} integral r\textasciicircum{}(d-s-1)dr.


\entrylabel{Second check}

The tail integral diverges whenever s<=\allowbreak{}d,\allowbreak{} covering the long-range regimes. For the logarithmic case the corresponding Hessian r\textasciicircum{}-2 gives the same limiting power at s=\allowbreak{}0.


\entrylabel{Remaining obstruction}

One needs a deformation adapted to neutrality or screening,\allowbreak{} with cancellation of long-range contributions and control of rare configurations. Absolute Hessian integrability cannot establish the stated symmetry theorem. Existing periodic-limit uniqueness is not automatically a theorem for every Gibbs state.


Verification:\allowbreak{} author checked partial | Independent review:\allowbreak{} none | Full target solved:\allowbreak{} False


\clearpage

\section*{P0314  Uniqueness of the general sine-beta DLR Gibbs measure}

\entrylabel{Primary source}

Dumaz and Malvy:\allowbreak{} Long-Range Correlation of Sine beta Point Process (2026);\allowbreak{} Introduction,\allowbreak{}PDF2


\url{https://arxiv.org/pdf/2603.15289}


Source statement date:\allowbreak{} 2026-03-16. Archived access:\allowbreak{} 2026-09-07T01:\allowbreak{}33:\allowbreak{}42Z.


Historical:\allowbreak{} source\_\allowbreak{}says\_\allowbreak{}open | Current:\allowbreak{} unverified | Result:\allowbreak{} unsolved


\entrylabel{Question}

Prove or refute uniqueness of the sine-beta Gibbs measure under its Dobrushin-Lanford-Ruelle characterization for general positive beta.


\entrylabel{Scope}

The March2026 primary source explicitly calls uniqueness for the DLR characterization open. A stationary free-energy characterization and special determinantal cases do not by themselves settle all DLR measures.


\subsection*{P0314-A001 | Source-specific calculation,\allowbreak{} conditional implication or exact obstruction}

\entrylabel{Argument or result}

A tempting variational route is to prove that the specific relative entropy of two candidate states vanishes. Vanishing entropy per volume alone cannot imply equality of states. On n binary sites take P\_\allowbreak{}n=\allowbreak{}delta\_\allowbreak{}(0,\allowbreak{}...,\allowbreak{}0) and Q\_\allowbreak{}n=\allowbreak{}(delta\_\allowbreak{}(0,\allowbreak{}...,\allowbreak{}0)+\allowbreak{}delta\_\allowbreak{}(1,\allowbreak{}...,\allowbreak{}1))/\allowbreak{}2. Then D(P\_\allowbreak{}n||Q\_\allowbreak{}n)=\allowbreak{}log2 for every n,\allowbreak{} hence D/\allowbreak{}n tends to0,\allowbreak{} while the limiting cylinder event X\_\allowbreak{}1=\allowbreak{}1 has probabilities0 and1/\allowbreak{}2.


\entrylabel{Second check}

The density dP\_\allowbreak{}n/\allowbreak{}dQ\_\allowbreak{}n is2 on its only supported configuration,\allowbreak{} so the relative entropy is exactly log2. The two consistent infinite-volume measures are visibly different.


\entrylabel{Remaining obstruction}

A successful variational argument needs a uniqueness mechanism beyond zero specific entropy,\allowbreak{} plus a theorem placing every DLR state in the admissible stationary finite-energy class. The diagnostic does not construct a second sine-beta Gibbs measure.


Verification:\allowbreak{} author checked partial | Independent review:\allowbreak{} none | Full target solved:\allowbreak{} False


\clearpage

\section*{P0315  Forrester-Haldane pointwise pair-correlation asymptotics for sine beta}

\entrylabel{Primary source}

Dumaz and Malvy:\allowbreak{} Long-Range Correlation of Sine beta Point Process (2026);\allowbreak{} Introduction,\allowbreak{}PDF2-3;\allowbreak{}Lewin equations149-150,\allowbreak{}PDF64


\url{https://arxiv.org/pdf/2603.15289}


Source statement date:\allowbreak{} 2026-03-16. Archived access:\allowbreak{} 2026-09-07T01:\allowbreak{}33:\allowbreak{}42Z.


Historical:\allowbreak{} source\_\allowbreak{}says\_\allowbreak{}open | Current:\allowbreak{} unverified | Result:\allowbreak{} unsolved


\entrylabel{Question}

Prove the predicted oscillatory large-distance expansion of the sine-beta truncated two-point correlation function for general beta,\allowbreak{} beyond the established special values and weaker averaged decay estimates.


\entrylabel{Scope}

All harmonics and beta regimes are one asymptotic family. Lewin uses unit density;\allowbreak{} Dumaz-Malvy use density1/\allowbreak{}(2*pi),\allowbreak{} so distances must be rescaled before comparing coefficients.


\subsection*{P0315-A001 | Source-specific calculation,\allowbreak{} conditional implication or exact obstruction}

\entrylabel{Argument or result}

At beta=\allowbreak{}2 and unit density the exact formula is rho\_\allowbreak{}T(r)=\allowbreak{}-sin\textasciicircum{}2(pi*r)/\allowbreak{}(pi\textasciicircum{}2*r\textasciicircum{}2)=\allowbreak{}-1/\allowbreak{}(2*pi\textasciicircum{}2*r\textasciicircum{}2)+\allowbreak{}cos(2*pi*r)/\allowbreak{}(2*pi\textasciicircum{}2*r\textasciicircum{}2). This verifies the collision of the nonoscillatory r\textasciicircum{}-2 and oscillatory r\textasciicircum{}-4/\allowbreak{}beta powers at beta=\allowbreak{}2. It also shows why a naive nonzero pointwise ratio to a single power cannot hold:\allowbreak{} the expression vanishes at every nonzero integer.


\entrylabel{Second check}

The double-angle identity gives the two terms exactly. Rescaling from unit density to density1/\allowbreak{}(2*pi) replaces pi*r by r/\allowbreak{}2 and the prefactor by1/\allowbreak{}pi\textasciicircum{}2 in the source normalization.


\entrylabel{Remaining obstruction}

For general beta the March2026 theorem bounds a locally averaged correlation by a nonsharp power;\allowbreak{} it does not give the exact oscillatory coefficients or pointwise remainder. The desired expansion needs additive error control that remains meaningful at zeros. No generic asymptotic is proved.


Verification:\allowbreak{} author checked partial | Independent review:\allowbreak{} none | Full target solved:\allowbreak{} False


\clearpage

\section*{P0316  Two phase transitions in negative-exponent one-dimensional Riesz gases}

\entrylabel{Primary source}

Dereudre and Digneaux:\allowbreak{} Rigidity of one dimensional point processes via optimal transport;\allowbreak{} Introduction and Figure1,\allowbreak{}PDF2-3;\allowbreak{}Lewin PDF68


\url{https://arxiv.org/pdf/2510.17257}


Source statement date:\allowbreak{} 2025-10-20. Archived access:\allowbreak{} 2026-09-07T01:\allowbreak{}33:\allowbreak{}44Z.


Historical:\allowbreak{} source\_\allowbreak{}says\_\allowbreak{}open | Current:\allowbreak{} unverified | Result:\allowbreak{} unsolved


\entrylabel{Question}

Prove or refute the predicted temperature-dependent phase diagram,\allowbreak{} including the suggested two transitions,\allowbreak{} for one-dimensional Riesz gases with -1<s<0.


\entrylabel{Scope}

The2025 source explicitly leaves this interval open. Its rigidity and hyperuniformity result for -2<s<=\allowbreak{}-1 is established prior work and outside this unresolved interval.


\subsection*{P0316-A001 | Source-specific calculation,\allowbreak{} conditional implication or exact obstruction}

\entrylabel{Argument or result}

For V(r)=\allowbreak{}-r\textasciicircum{}p with p=\allowbreak{}-s in(0,\allowbreak{}1),\allowbreak{} the positive second derivative is p*(1-p)*r\textasciicircum{}(p-2). The harmonic lattice stiffness therefore has small-frequency order |k|\textasciicircum{}(1-p)=\allowbreak{}|k|\textasciicircum{}(1+\allowbreak{}s),\allowbreak{} obtained from the sum of (1-cos(k*r))*r\textasciicircum{}(-2-s). Its Gaussian displacement variance involves integral\_\allowbreak{}0\textasciicircum{}1 k\textasciicircum{}(-1-s)dk,\allowbreak{} finite for every s<0 at any fixed temperature.


\entrylabel{Second check}

The substitution u=\allowbreak{}k*r in the continuum comparison gives the stiffness exponent1+\allowbreak{}s. The integral is finite because -1-s>-1. Thus the harmonic approximation alone predicts finite roughness throughout the entire negative interval and cannot locate a temperature transition.


\entrylabel{Remaining obstruction}

Defects and changes of ordering must be controlled beyond the Gaussian expansion. The finite harmonic variance is neither a proof of a crystal nor a refutation of the proposed fluid and intermediate phases. No transition temperatures are established.


Verification:\allowbreak{} author checked partial | Independent review:\allowbreak{} none | Full target solved:\allowbreak{} False


\clearpage

\section*{P0317  Hyperuniformity of general long-range Riesz Gibbs gases}

\entrylabel{Primary source}

Lewin:\allowbreak{} Coulomb and Riesz gases,\allowbreak{} the known and the unknown,\allowbreak{}2022;\allowbreak{} SectionVI.C,\allowbreak{}PDF69


\url{https://arxiv.org/pdf/2202.09240}


Source statement date:\allowbreak{} 2022-04-30. Archived access:\allowbreak{} 2026-09-07T00:\allowbreak{}34:\allowbreak{}23Z.


Historical:\allowbreak{} source\_\allowbreak{}says\_\allowbreak{}open | Current:\allowbreak{} unverified | Result:\allowbreak{} unsolved


\entrylabel{Question}

Prove the expected hyperuniformity of the long-range Riesz gas in the general regimes where the review predicts it,\allowbreak{} beyond presently established special cases.


\entrylabel{Scope}

The question concerns subvolume number variance in growing windows. Fixed total particle number in a finite canonical ensemble is not itself hyperuniformity of an infinite-volume limit;\allowbreak{} newer one-dimensional and Coulomb results are partial established cases.


\subsection*{P0317-A001 | Source-specific calculation,\allowbreak{} conditional implication or exact obstruction}

\entrylabel{Argument or result}

Consider N independent uniform points in a box of volume N. The canonical truncated pair density is -1/\allowbreak{}N within the box,\allowbreak{} whose integral over the full box is -1. In any fixed bounded window A,\allowbreak{} the count has variance |A|*(1-|A|/\allowbreak{}N),\allowbreak{} tending to |A|;\allowbreak{} the local point-process limit is a unit-intensity Poisson process and is not hyperuniform.


\entrylabel{Second check}

The count is Binomial(N,\allowbreak{}|A|/\allowbreak{}N),\allowbreak{} so its variance formula is exact. Pointwise the truncated density tends to0 while its integral stays-1,\allowbreak{} demonstrating loss of the sum-rule mass at spatial infinity.


\entrylabel{Remaining obstruction}

The canonical neutrality sum rule must be supplemented by uniform correlation-tail or screening estimates before passage to an infinite-volume variance law. Those estimates are not provided here for the full Riesz family.


Verification:\allowbreak{} author checked partial | Independent review:\allowbreak{} none | Full target solved:\allowbreak{} False


\clearpage

\section*{P0318  Sharp number-variance growth exponents for long-range Riesz gases}

\entrylabel{Primary source}

Lewin:\allowbreak{} Coulomb and Riesz gases,\allowbreak{} the known and the unknown,\allowbreak{}2022;\allowbreak{} SectionVI.C,\allowbreak{}PDF69


\url{https://arxiv.org/pdf/2202.09240}


Source statement date:\allowbreak{} 2022-04-30. Archived access:\allowbreak{} 2026-09-07T00:\allowbreak{}34:\allowbreak{}23Z.


Historical:\allowbreak{} source\_\allowbreak{}says\_\allowbreak{}open | Current:\allowbreak{} unverified | Result:\allowbreak{} unsolved


\entrylabel{Question}

Establish the predicted three regimes for the variance of the number of particles in a ball of radius R:\allowbreak{} order R\textasciicircum{}s for s>d-1,\allowbreak{} order R\textasciicircum{}(d-1)*log R at s=\allowbreak{}d-1,\allowbreak{} and order R\textasciicircum{}(d-1) for s<d-1.


\entrylabel{Scope}

The three regimes are one source family. The question asks for sharp growth beyond hyperuniformity alone. Existing one-dimensional gap fluctuation and selected Coulomb results must be treated as known subcases.


\subsection*{P0318-A001 | Source-specific calculation,\allowbreak{} conditional implication or exact obstruction}

\entrylabel{Argument or result}

Conditional on a stationary unit-density process with an integrable truncated pair density h satisfying integral h=\allowbreak{}-1,\allowbreak{} Var N(B\_\allowbreak{}R)=\allowbreak{}-integral h(z)*(|B\_\allowbreak{}R|-|B\_\allowbreak{}R intersect(B\_\allowbreak{}R-z)|)dz. For a ball,\allowbreak{} the intersection deficit is of order R\textasciicircum{}(d-1)*min(|z|,\allowbreak{}R). If h has a noncanceling negative tail of order |z|\textasciicircum{}(-(2d-s)),\allowbreak{} the contribution from1<|z|<R is of order R\textasciicircum{}(d-1)*integral\_\allowbreak{}1\textasciicircum{}R r\textasciicircum{}(s-d)dr,\allowbreak{} yielding exactly the three predicted powers.


\entrylabel{Second check}

The variance identity follows by expanding the second factorial moment and using the sum rule. The tail beyond R contributes order R\textasciicircum{}s and is subleading in the surface regime;\allowbreak{} at equality the integral is logarithmic.


\entrylabel{Remaining obstruction}

This is a conditional calculation,\allowbreak{} not a proof that a Gibbs state satisfies the required sum rule,\allowbreak{} signed tail,\allowbreak{} uniform bounds or positive surface coefficient. Oscillatory tails may require cancellation rather than absolute comparisons. These unresolved hypotheses carry the probabilistic content.


Verification:\allowbreak{} author checked partial | Independent review:\allowbreak{} none | Full target solved:\allowbreak{} False


\clearpage

\section*{P0320  Number rigidity in the remaining Riesz exponent band}

\entrylabel{Primary source}

Thoma:\allowbreak{} Non-rigidity Properties of the Coulomb Gas;\allowbreak{} Introduction,\allowbreak{}PDF2


\url{https://arxiv.org/pdf/2303.11486}


Source statement date:\allowbreak{} 2023-03-20. Archived access:\allowbreak{} 2026-09-07T01:\allowbreak{}35:\allowbreak{}29Z.


Historical:\allowbreak{} source\_\allowbreak{}says\_\allowbreak{}open | Current:\allowbreak{} unverified | Result:\allowbreak{} unsolved


\entrylabel{Question}

Determine number rigidity or non-rigidity for Riesz gases in the remaining admissible band s in(d-2,\allowbreak{}d-1],\allowbreak{} which Thoma identifies as open after disproving the older general threshold conjecture.


\entrylabel{Scope}

The earlier threshold claim is retained separately as an excluded historical conjecture P0319. The interval is one question,\allowbreak{} restricted to the Riesz models for which an infinite-volume gas is defined. Later status remains unverified.


\subsection*{P0320-A001 | Source-specific calculation,\allowbreak{} conditional implication or exact obstruction}

\entrylabel{Argument or result}

A symmetric displacement comparison cancels the first-order force:\allowbreak{} V(x+\allowbreak{}h)+\allowbreak{}V(x-h)-2V(x) is controlled by |h|\textasciicircum{}2 times a local Hessian bound. For V(x)=\allowbreak{}|x|\textasciicircum{}-s,\allowbreak{} the far-field Hessian integral has radial order integral\_\allowbreak{}1\textasciicircum{}infinity r\textasciicircum{}(d-s-3)dr,\allowbreak{} finite exactly when s>d-2. This explains why a second-order transport argument might reach the stated band despite a nonintegrable first derivative.


\entrylabel{Second check}

Taylor expansion with integral remainder yields the second difference;\allowbreak{} convergence requires d-s-3<-1,\allowbreak{} equivalent to s>d-2. The equality endpoint is logarithmically divergent and must not be included by this argument.


\entrylabel{Remaining obstruction}

Integrability of the deterministic tail is only one ingredient. Constructing a number-changing comparison also requires entropy,\allowbreak{} label matching,\allowbreak{} local singularity control and uniform conditional probability bounds. No non-rigidity theorem for the open band is proved here.


Verification:\allowbreak{} author checked partial | Independent review:\allowbreak{} none | Full target solved:\allowbreak{} False


\clearpage

\section*{P0321  Existence of a hexatic phase for two-dimensional Riesz gases}

\entrylabel{Primary source}

Lewin:\allowbreak{} Coulomb and Riesz gases,\allowbreak{} the known and the unknown,\allowbreak{}2022;\allowbreak{} SectionVI.B,\allowbreak{}PDF66-67


\url{https://arxiv.org/pdf/2202.09240}


Source statement date:\allowbreak{} 2022-04-30. Archived access:\allowbreak{} 2026-09-07T00:\allowbreak{}34:\allowbreak{}23Z.


Historical:\allowbreak{} source\_\allowbreak{}says\_\allowbreak{}open | Current:\allowbreak{} unverified | Result:\allowbreak{} unsolved


\entrylabel{Question}

Prove or refute the intermediate hexatic phase predicted in two-dimensional Riesz-gas phase diagrams,\allowbreak{} distinguished by orientational order and absence of crystalline positional order.


\entrylabel{Scope}

This is a finite-temperature order-type question distinct from existence of a three-dimensional crystal. Numerical phase boundaries are motivation,\allowbreak{} not rigorous evidence of existence.


\subsection*{P0321-A001 | Source-specific calculation,\allowbreak{} conditional implication or exact obstruction}

\entrylabel{Argument or result}

A rotationally invariant Gibbs mixture can have zero mean orientational order while retaining long-range two-point order. As a diagnostic let theta\_\allowbreak{}x=\allowbreak{}Theta+\allowbreak{}eta\_\allowbreak{}x,\allowbreak{} where Theta is uniform modulo2*pi and the eta\_\allowbreak{}x are iid. Then E exp(6*i*theta\_\allowbreak{}x)=\allowbreak{}0,\allowbreak{} but for x distinct from y,\allowbreak{} E exp(6*i*(theta\_\allowbreak{}x-theta\_\allowbreak{}y))=\allowbreak{}|E exp(6*i*eta\_\allowbreak{}x)|\textasciicircum{}2,\allowbreak{} which can be strictly positive.


\entrylabel{Second check}

The global angle cancels in the difference and independence factors the two local terms into conjugates. Taking all eta\_\allowbreak{}x=\allowbreak{}0 gives two-point order1 despite the zero one-point mean.


\entrylabel{Remaining obstruction}

Consequently a vanishing rotation-averaged one-point order parameter cannot rule out the proposed phase. A proof must establish the proper bond-orientational correlations together with positional decorrelation for the actual Gibbs law,\allowbreak{} with defects and thermodynamic limits controlled.


Verification:\allowbreak{} author checked partial | Independent review:\allowbreak{} none | Full target solved:\allowbreak{} False


\clearpage

\section*{P0322  Critical gap-correlation decay in the one-dimensional s=\allowbreak{}1 Riesz gas}

\entrylabel{Primary source}

Boursier:\allowbreak{} Decay of correlations and thermodynamic limit for circular Riesz gas;\allowbreak{} Section1.2,\allowbreak{}PDF4


\url{https://arxiv.org/pdf/2209.00396}


Source statement date:\allowbreak{} 2023-07-19. Archived access:\allowbreak{} 2026-09-07T01:\allowbreak{}33:\allowbreak{}41Z.


Historical:\allowbreak{} source\_\allowbreak{}says\_\allowbreak{}open | Current:\allowbreak{} unverified | Result:\allowbreak{} unsolved


\entrylabel{Question}

Determine the sharp decay of ordered-gap correlations in the critical one-dimensional Riesz gas at s=\allowbreak{}1,\allowbreak{} where the source conjectures a distance\textasciicircum{}-1 times logarithmic correction.


\entrylabel{Scope}

The source treats s=\allowbreak{}1 as a distinct marginal regime omitted from its theorems for s<1 and s>1. This entry does not count separate values of an unknown logarithmic exponent.


\subsection*{P0322-A001 | Source-specific calculation,\allowbreak{} conditional implication or exact obstruction}

\entrylabel{Argument or result}

The harmonic heuristic at the marginal exponent leads to a logarithmically divergent Fourier symbol. The identity sum\_\allowbreak{}(n>=\allowbreak{}1) cos(n*theta)/\allowbreak{}n=\allowbreak{}-log(2*sin(theta/\allowbreak{}2)) for0<theta<2*pi behaves as log(1/\allowbreak{}theta) near0. Inverting such a symbol suggests a slowly varying factor1/\allowbreak{}log(1/\allowbreak{}theta),\allowbreak{} rather than an ordinary power-law inverse.


\entrylabel{Second check}

The identity follows by taking the real part of sum z\textasciicircum{}n/\allowbreak{}n=\allowbreak{}-log(1-z) and an Abel limit. Since 2*sin(theta/\allowbreak{}2) is asymptotic to theta,\allowbreak{} the logarithmic divergence is explicit.


\entrylabel{Remaining obstruction}

The actual random Hessian must first be compared with this translation-invariant harmonic operator,\allowbreak{} and a precise inverse-Fourier estimate with correct gap observable is required. The calculation does not identify the logarithmic exponent or prove a critical decay bound.


Verification:\allowbreak{} author checked partial | Independent review:\allowbreak{} none | Full target solved:\allowbreak{} False


\clearpage

\section*{P0323  DLR characterization of the circular Riesz thermodynamic limit}

\entrylabel{Primary source}

Boursier:\allowbreak{} Decay of correlations and thermodynamic limit for circular Riesz gas;\allowbreak{} Section1.3,\allowbreak{}PDF5


\url{https://arxiv.org/pdf/2209.00396}


Source statement date:\allowbreak{} 2023-07-19. Archived access:\allowbreak{} 2026-09-07T01:\allowbreak{}33:\allowbreak{}41Z.


Historical:\allowbreak{} source\_\allowbreak{}says\_\allowbreak{}open | Current:\allowbreak{} unverified | Result:\allowbreak{} unsolved


\entrylabel{Question}

Prove that the thermodynamic limit of the circular one-dimensional Riesz gas constructed in Boursier satisfies the anticipated infinite-volume DLR equations.


\entrylabel{Scope}

The paper proves convergence of the finite circular laws for0<s<1 and s>1 and explicitly says the DLR characterization remains open. This is existence of the specification identity for that limit,\allowbreak{} distinct from uniqueness among all Gibbs states.


\subsection*{P0323-A001 | Source-specific calculation,\allowbreak{} conditional implication or exact obstruction}

\entrylabel{Argument or result}

Weak convergence of joint laws does not automatically preserve a conditional specification that depends increasingly sensitively on the conditioning variable. Let Y be uniform on[0,\allowbreak{}1] and X\_\allowbreak{}n=\allowbreak{}1 if the fractional part of n*Y is below1/\allowbreak{}2. Every X\_\allowbreak{}n is a measurable deterministic function of Y,\allowbreak{} yet (Y,\allowbreak{}X\_\allowbreak{}n) converges weakly to a uniform Y independent of a fair Bernoulli variable.


\entrylabel{Second check}

For a continuous test function of Y,\allowbreak{} sums over the n alternating half-cells are Riemann sums giving half the integral for either value of X\_\allowbreak{}n. This proves the stated joint convergence;\allowbreak{} the finite conditional probabilities take only0 or1 whereas the limiting conditional probability is1/\allowbreak{}2.


\entrylabel{Remaining obstruction}

The Riesz problem therefore needs continuity or uniform integrability of the renormalized exterior potential and conditional partition functions,\allowbreak{} beyond local weak convergence. The example diagnoses that missing interchange;\allowbreak{} it is not a Riesz counterexample.


Verification:\allowbreak{} author checked partial | Independent review:\allowbreak{} none | Full target solved:\allowbreak{} False


\clearpage

\section*{P0324  Analyticity of the one-dimensional Riesz free energy}

\entrylabel{Primary source}

Boursier:\allowbreak{} Decay of correlations and thermodynamic limit for circular Riesz gas;\allowbreak{} Section1.3,\allowbreak{}PDF5


\url{https://arxiv.org/pdf/2209.00396}


Source statement date:\allowbreak{} 2023-07-19. Archived access:\allowbreak{} 2026-09-07T01:\allowbreak{}33:\allowbreak{}41Z.


Historical:\allowbreak{} source\_\allowbreak{}says\_\allowbreak{}open | Current:\allowbreak{} unverified | Result:\allowbreak{} unsolved


\entrylabel{Question}

Prove the expected real analyticity in inverse temperature beta of the thermodynamic free energy of the one-dimensional long-range Riesz gas0<s<1 defined in the source.


\entrylabel{Scope}

The source discusses deriving twice differentiability from its correlation and convergence results and expects full analyticity. Neither a possible C2 consequence nor convergence of periodic laws is counted as analyticity.


\subsection*{P0324-A001 | Source-specific calculation,\allowbreak{} conditional implication or exact obstruction}

\entrylabel{Argument or result}

Finite-volume derivatives of log Z are energy cumulants:\allowbreak{} the second is a variance and the kth derivative is the kth centered cumulant up to sign. Analyticity in a beta neighborhood requires control at all orders,\allowbreak{} typically |cumulant\_\allowbreak{}k|/\allowbreak{}volume bounded by C\textasciicircum{}k*k!,\allowbreak{} not merely a two-point decay estimate. Even real-analytic finite-volume functions can converge to a nonanalytic limit:\allowbreak{} f\_\allowbreak{}n(beta)=\allowbreak{}sqrt((beta-beta0)\textasciicircum{}2+\allowbreak{}1/\allowbreak{}n) tends to |beta-beta0| for any fixed positive beta0.


\entrylabel{Second check}

Each f\_\allowbreak{}n has a positive radicand on the real axis and is real analytic there,\allowbreak{} while the limit has different one-sided derivatives at beta0. Its complex singularities approach beta0 at distance1/\allowbreak{}sqrt(n),\allowbreak{} showing the missing uniform analytic radius.


\entrylabel{Remaining obstruction}

Uniform higher-order cumulant estimates or a convergent expansion are needed throughout the desired temperature range. The existing two-point control does not by itself provide them. No full analyticity theorem is established here.


Verification:\allowbreak{} author checked partial | Independent review:\allowbreak{} none | Full target solved:\allowbreak{} False


\clearpage

\section*{P0325  Uniqueness of the microscopic bulk limit of the two-dimensional Coulomb gas}

\entrylabel{Primary source}

Leble:\allowbreak{} Les log-gas et leurs amis,\allowbreak{} HDR June2026;\allowbreak{} Section2.3,\allowbreak{}PDF25


\url{https://tleble.perso.math.cnrs.fr/assets/papers/HDR_Leble.pdf}


Source statement date:\allowbreak{} 2026-06-10. Archived access:\allowbreak{} 2026-09-07T01:\allowbreak{}48:\allowbreak{}40Z.


Historical:\allowbreak{} source\_\allowbreak{}says\_\allowbreak{}open | Current:\allowbreak{} unverified | Result:\allowbreak{} unsolved


\entrylabel{Question}

For beta different from2,\allowbreak{} does the law of the bulk point configuration of the two-dimensional one-component plasma converge as particle number tends to infinity,\allowbreak{} rather than only along subsequences?


\entrylabel{Scope}

The source finite-volume laws and bulk observation points are those in Section2.3. The Ginibre beta=\allowbreak{}2 case is known. This concerns convergence of the designated finite-volume approximations,\allowbreak{} not uniqueness of every DLR state.


\subsection*{P0325-A001 | Explicit model calculation or source-specific limiting obstruction}

\entrylabel{Argument or result}

The new theorem gives translation invariance,\allowbreak{} number rigidity and DLR identities for every subsequential limit. These properties need not identify a measure:\allowbreak{} two distinct measures can satisfy the same invariance and integrability constraints. Even at the elementary sequence level,\allowbreak{} alternating between two different stationary point-process laws gives a tight sequence with those shared properties but no limit.


\entrylabel{Second check}

As an explicit invariance-only diagnostic,\allowbreak{} alternate a uniform translate of Z\textasciicircum{}2 and a uniform translate of the rectangular lattice with spacings2 and1/\allowbreak{}2. Both have unit intensity and are stationary;\allowbreak{} a ball of radius0.4 has different second factorial moments under the two laws. This does not claim either law satisfies the Coulomb DLR specification.


\entrylabel{Remaining obstruction}

The missing step is a characterization that uniquely selects the Coulomb bulk law,\allowbreak{} or a quantitative comparison of all finite-volume boundary approximations. The shared qualitative properties do not prove uniqueness of the subsequential limits.


Verification:\allowbreak{} author checked partial | Independent review:\allowbreak{} none | Full target solved:\allowbreak{} False


\clearpage

\section*{P0326  High-temperature uniqueness and variational characterization of the two-dimensional plasma}

\entrylabel{Primary source}

Leble:\allowbreak{} Les log-gas et leurs amis,\allowbreak{} HDR June2026;\allowbreak{} Section2.4,\allowbreak{}PDF25


\url{https://tleble.perso.math.cnrs.fr/assets/papers/HDR_Leble.pdf}


Source statement date:\allowbreak{} 2026-06-10. Archived access:\allowbreak{} 2026-09-07T01:\allowbreak{}48:\allowbreak{}40Z.


Historical:\allowbreak{} source\_\allowbreak{}says\_\allowbreak{}open | Current:\allowbreak{} unverified | Result:\allowbreak{} unsolved


\entrylabel{Question}

For sufficiently small positive beta,\allowbreak{} prove that the two-dimensional one-component plasma has a unique DLR solution and that this is also the unique free-energy minimizer.


\entrylabel{Scope}

The two descriptions form the source liquid-phase uniqueness objective. Exponential decay of correlations is a separate observable recorded in P0327. Uniqueness of a bulk approximation alone would not establish this stronger assertion.


\subsection*{P0326-A001 | Explicit model calculation or source-specific limiting obstruction}

\entrylabel{Argument or result}

The usual absolute-integrability high-temperature cluster estimate does not apply directly to the logarithmic pair interaction. In two dimensions its force has magnitude1/\allowbreak{}r,\allowbreak{} whose integral over1<|x|<R is2*pi*(R-1). Multiplication by an arbitrarily small fixed beta still diverges as R tends to infinity.


\entrylabel{Second check}

The polar Jacobian r cancels the factor1/\allowbreak{}r exactly. Thus beta-smallness cannot turn this absolute first-force bound into a finite contraction constant at infinite volume.


\entrylabel{Remaining obstruction}

An argument must incorporate the neutralizing background and screening before estimating influences. A direct short-range Dobrushin estimate fails to exploit those cancellations. No uniqueness theorem or equivalence of all DLR states with variational minimizers is proved.


Verification:\allowbreak{} author checked partial | Independent review:\allowbreak{} none | Full target solved:\allowbreak{} False


\clearpage

\section*{P0327  Exponential clustering in the high-temperature two-dimensional plasma}

\entrylabel{Primary source}

Leble:\allowbreak{} Les log-gas et leurs amis,\allowbreak{} HDR June2026;\allowbreak{} Section2.4,\allowbreak{}PDF25


\url{https://tleble.perso.math.cnrs.fr/assets/papers/HDR_Leble.pdf}


Source statement date:\allowbreak{} 2026-06-10. Archived access:\allowbreak{} 2026-09-07T01:\allowbreak{}48:\allowbreak{}40Z.


Historical:\allowbreak{} source\_\allowbreak{}says\_\allowbreak{}open | Current:\allowbreak{} unverified | Result:\allowbreak{} unsolved


\entrylabel{Question}

Prove exponential decay of the truncated two-point correlation function of the two-dimensional one-component plasma for all sufficiently small positive beta.


\entrylabel{Scope}

The beta=\allowbreak{}2 Gaussian decay is established. The generic non-Coulomb power-law prediction in P0311 excludes this Coulomb mechanism;\allowbreak{} this is not another exponent copy of that question.


\subsection*{P0327-A001 | Explicit model calculation or source-specific limiting obstruction}

\entrylabel{Argument or result}

Linear Debye screening gives the formal structure factor S(k)=\allowbreak{}|k|\textasciicircum{}2/\allowbreak{}(|k|\textasciicircum{}2+\allowbreak{}m\textasciicircum{}2),\allowbreak{} with m\textasciicircum{}2 proportional to beta times the background density. Its truncated-correlation Fourier transform is proportional to -m\textasciicircum{}2/\allowbreak{}(|k|\textasciicircum{}2+\allowbreak{}m\textasciicircum{}2). In two dimensions the inverse transform is a multiple of the modified Bessel function K0(m*r),\allowbreak{} whose large-distance behavior is proportional to exp(-m*r)/\allowbreak{}sqrt(m*r).


\entrylabel{Second check}

The resolvent of -Delta+\allowbreak{}m\textasciicircum{}2 has Fourier multiplier1/\allowbreak{}(|k|\textasciicircum{}2+\allowbreak{}m\textasciicircum{}2);\allowbreak{} its heat-kernel integral yields the positive radial K0 kernel. Thus screening formally creates an exponential length m\textasciicircum{}-1 rather than a Riesz power tail.


\entrylabel{Remaining obstruction}

The linear-response closure is not proved for the interacting Gibbs law. A rigorous result needs a uniform control of nonlinear correlations and error terms at distances beyond the screening length. The calculation identifies a mechanism,\allowbreak{} not a clustering theorem.


Verification:\allowbreak{} author checked partial | Independent review:\allowbreak{} none | Full target solved:\allowbreak{} False


\clearpage

\section*{P0328  Gravitational allocation for the two-dimensional one-component plasma}

\entrylabel{Primary source}

Leble:\allowbreak{} Les log-gas et leurs amis,\allowbreak{} HDR June2026;\allowbreak{} Section3.3 perspective,\allowbreak{}PDF30


\url{https://tleble.perso.math.cnrs.fr/assets/papers/HDR_Leble.pdf}


Source statement date:\allowbreak{} 2026-06-10. Archived access:\allowbreak{} 2026-09-07T01:\allowbreak{}48:\allowbreak{}40Z.


Historical:\allowbreak{} source\_\allowbreak{}says\_\allowbreak{}open | Current:\allowbreak{} unverified | Result:\allowbreak{} unsolved


\entrylabel{Question}

Construct and control the gravitational allocation from Lebesgue measure to the infinite two-dimensional Coulomb gas,\allowbreak{} in the sense of the proposed extension of the Nazarov-Sodin-Volberg construction;\allowbreak{} the source says this remains open even at beta=\allowbreak{}2.


\entrylabel{Scope}

Existence of some finite-cost stationary transport coupling is already known. The target is the specified gradient-flow allocation,\allowbreak{} including meaningful cells and sufficient cell-size control.


\subsection*{P0328-A001 | Explicit model calculation or source-specific limiting obstruction}

\entrylabel{Argument or result}

Gauss-law heuristics relate unit source strength to unit cell area,\allowbreak{} provided the flow is defined,\allowbreak{} flux through the cell boundary vanishes and trajectories are captured. Those hypotheses do not give a second-moment bound on allocation distance. A random cell-radius bound R that is finite almost surely can still have infinite second moment:\allowbreak{} choose P(R>r)=\allowbreak{}1/\allowbreak{}r for r>=\allowbreak{}1,\allowbreak{} giving E R\textasciicircum{}2=\allowbreak{}2*integral r*P(R>r)dr=\allowbreak{}infinity.


\entrylabel{Second check}

The displayed survival function defines a proper probability law with R finite almost surely. Its second-moment tail integral diverges linearly,\allowbreak{} so boundedness of individual random cells alone cannot justify finite expected quadratic transport cost.


\entrylabel{Remaining obstruction}

The missing work includes renormalizing the infinite force field,\allowbreak{} ruling out escaping trajectories and quantifying cell tails. The already established abstract transport coupling is not necessarily generated by the gravitational field.


Verification:\allowbreak{} author checked partial | Independent review:\allowbreak{} none | Full target solved:\allowbreak{} False


\clearpage

\section*{P0329  Jancovici-Lebowitz-Manificat discrepancy deviations for general beta}

\entrylabel{Primary source}

Leble:\allowbreak{} Les log-gas et leurs amis,\allowbreak{} HDR June2026;\allowbreak{} Section2.2,\allowbreak{}PDF21;\allowbreak{}Serfaty Section10.3,\allowbreak{}PDF224


\url{https://tleble.perso.math.cnrs.fr/assets/papers/HDR_Leble.pdf}


Source statement date:\allowbreak{} 2026-06-10. Archived access:\allowbreak{} 2026-09-07T01:\allowbreak{}48:\allowbreak{}40Z.


Historical:\allowbreak{} source\_\allowbreak{}says\_\allowbreak{}open | Current:\allowbreak{} unverified | Result:\allowbreak{} unsolved


\entrylabel{Question}

Prove the predicted charge-discrepancy deviation exponents for the two-dimensional Coulomb gas at general beta,\allowbreak{} especially the unresolved small-deviation regime alpha>1/\allowbreak{}2 close to1/\allowbreak{}2.


\entrylabel{Scope}

For a microscopic-radius R disk,\allowbreak{} the predicted powers are phi(alpha)=\allowbreak{}2*alpha-1 for1/\allowbreak{}2<alpha<=\allowbreak{}1,\allowbreak{}3*alpha-2 for1<=\allowbreak{}alpha<=\allowbreak{}2,\allowbreak{}and2*alpha foralpha>=\allowbreak{}2. All regimes are one family. Beta=\allowbreak{}2 and some larger-deviation regimes are known;\allowbreak{} no unit leading constant is asserted.


\subsection*{P0329-A001 | Explicit model calculation or source-specific limiting obstruction}

\entrylabel{Argument or result}

If the discrepancy were Gaussian with variance c*R,\allowbreak{} a threshold R\textasciicircum{}alpha would cost logarithmic probability of order -R\textasciicircum{}(2*alpha-1)/\allowbreak{}(2*c). This exactly reproduces the first predicted power. It cannot produce the second and third powers:\allowbreak{} using the same Gaussian ansatz at alpha=\allowbreak{}3 would predict exponent5,\allowbreak{} whereas the large-excess prediction is6.


\entrylabel{Second check}

The normal tail logarithm at threshold u standard deviations is -u\textasciicircum{}2/\allowbreak{}2 up to lower-order logarithms. Substituting u=\allowbreak{}R\textasciicircum{}alpha/\allowbreak{}sqrt(c*R) yields2*alpha-1. The piecewise source exponents agree at alpha=\allowbreak{}1 and2,\allowbreak{} but have different slopes.


\entrylabel{Remaining obstruction}

A proof must control moderate deviations far beyond an ordinary CLT and account for nonlinear electrostatic cost of large charge rearrangements. The Gaussian calculation explains only one branch conditionally;\allowbreak{} it proves no rare-event asymptotic for the plasma.


Verification:\allowbreak{} author checked partial | Independent review:\allowbreak{} none | Full target solved:\allowbreak{} False


\clearpage

\section*{P0330  Higher-order free-energy expansion of the two-dimensional Coulomb gas}

\entrylabel{Primary source}

Serfaty:\allowbreak{} Lectures on Coulomb and Riesz gases,\allowbreak{}2024;\allowbreak{} Section1.3,\allowbreak{}PDF21;\allowbreak{}Section9.3,\allowbreak{}PDF199-201


\url{https://arxiv.org/pdf/2407.21194}


Source statement date:\allowbreak{} 2024-07-30. Archived access:\allowbreak{} 2026-09-07T01:\allowbreak{}48:\allowbreak{}41Z.


Historical:\allowbreak{} source\_\allowbreak{}says\_\allowbreak{}open | Current:\allowbreak{} unverified | Result:\allowbreak{} unsolved


\entrylabel{Question}

Establish the predicted free-energy expansion beyond the thermodynamic order-N term for general beta in the two-dimensional logarithmic gas,\allowbreak{} including the proposed boundary,\allowbreak{} logarithmic and order-one contributions.


\entrylabel{Scope}

The family of coefficients is one question. Determinantal special cases and the order-N theorem are already established. The preliminary book has visible normalization inconsistencies in the schematic formula on PDF200;\allowbreak{} exact coefficients must be matched to its proved formula9.3.5 before use.


\subsection*{P0330-A001 | Explicit model calculation or source-specific limiting obstruction}

\entrylabel{Argument or result}

An exact beta=\allowbreak{}2 baseline is the planar density proportional to product\_\allowbreak{}(i<j)|z\_\allowbreak{}i-z\_\allowbreak{}j|\textasciicircum{}2 * exp(-N*sum|z\_\allowbreak{}i|\textasciicircum{}2),\allowbreak{} with labelled-particle Lebesgue normalization. Orthogonality of complex monomials gives Z\_\allowbreak{}N=\allowbreak{}N!*pi\textasciicircum{}N*product\_\allowbreak{}(k=\allowbreak{}0)\textasciicircum{}(N-1)k!\emph{N\textasciicircum{}(-N}(N+\allowbreak{}1)/\allowbreak{}2). Stirling and the Barnes-G expansion yield a term (N/\allowbreak{}2)*log N and a logarithmic term(5/\allowbreak{}12)*log N. They therefore check normalization and show why an unqualified(1/\allowbreak{}4)*N*log N in a schematic beta-dependent formula cannot be copied literally.


\entrylabel{Second check}

The monomial squared norm is pi*k!*N\textasciicircum{}(-k-1);\allowbreak{} expanding the Vandermonde determinant and integrating gives the N! product. The log N coefficient combines1/\allowbreak{}2 from log N! and-1/\allowbreak{}12 from log G(N+\allowbreak{}1),\allowbreak{} giving5/\allowbreak{}12. The N log N coefficient is1-1/\allowbreak{}2=\allowbreak{}1/\allowbreak{}2.


\entrylabel{Remaining obstruction}

These exact identities are known determinantal calculations,\allowbreak{} not a solution at general beta. Beyond beta=\allowbreak{}2,\allowbreak{} one needs uniform boundary free-energy corrections and an order-one determinant/\allowbreak{}field-theory identification with controlled remainder. No generic higher-order expansion is proved.


Verification:\allowbreak{} author checked partial | Independent review:\allowbreak{} none | Full target solved:\allowbreak{} False


\clearpage

\section*{P0331  Sharpness of the Debye scale in Coulomb local laws}

\entrylabel{Primary source}

Serfaty:\allowbreak{} Lectures on Coulomb and Riesz gases,\allowbreak{}2024;\allowbreak{} Equations8.0.1-8.0.2,\allowbreak{}PDF163;\allowbreak{}conjecturePDF164


\url{https://arxiv.org/pdf/2407.21194}


Source statement date:\allowbreak{} 2024-07-30. Archived access:\allowbreak{} 2026-09-07T01:\allowbreak{}48:\allowbreak{}41Z.


Historical:\allowbreak{} source\_\allowbreak{}says\_\allowbreak{}open | Current:\allowbreak{} unverified | Result:\allowbreak{} unsolved


\entrylabel{Question}

Prove the conjectured sharpness of the temperature-dependent minimal Debye scale in the Coulomb local-energy laws,\allowbreak{} including the proposed Poissonian behavior below that scale.


\entrylabel{Scope}

The source scale is rho\_\allowbreak{}beta=\allowbreak{}C*max(1,\allowbreak{}sqrt(chi(beta)/\allowbreak{}beta)) in microscopic coordinates;\allowbreak{} chi(beta) is the source dimension-dependent correction. In d>=\allowbreak{}3 it is1. All temperature and dimension cases are one scale-sharpness question.


\subsection*{P0331-A001 | Explicit model calculation or source-specific limiting obstruction}

\entrylabel{Argument or result}

The linear screened structure factor S(k)=\allowbreak{}k\textasciicircum{}2/\allowbreak{}(k\textasciicircum{}2+\allowbreak{}m\textasciicircum{}2),\allowbreak{} m of order sqrt(beta),\allowbreak{} is close to1 for frequencies k much larger than m and small for k much smaller than m. Thus it predicts Poisson-sized fluctuations for wavelengths much shorter than beta\textasciicircum{}-1/\allowbreak{}2 and suppressed fluctuations above that length. This is consistent with the d>=\allowbreak{}3 scale in the theorem.


\entrylabel{Second check}

The exact algebra gives S(k)=\allowbreak{}1+\allowbreak{}O(m\textasciicircum{}2/\allowbreak{}k\textasciicircum{}2) in the first regime and S(k)=\allowbreak{}O(k\textasciicircum{}2/\allowbreak{}m\textasciicircum{}2) in the second. The crossover solves k\textasciicircum{}2=\allowbreak{}m\textasciicircum{}2,\allowbreak{} so its reciprocal length is proportional to beta\textasciicircum{}-1/\allowbreak{}2.


\entrylabel{Remaining obstruction}

This is only a linear-response diagnostic. Sharp failure of the nonlinear energy estimate8.0.1 below rho\_\allowbreak{}beta is stronger than a structure-factor heuristic and needs a lower bound or a Poisson limit in a specified joint temperature-volume regime. The two-dimensional correction is not justified by this calculation.


Verification:\allowbreak{} author checked partial | Independent review:\allowbreak{} none | Full target solved:\allowbreak{} False


\clearpage

\section*{P0332  Regularity threshold for bounded Coulomb-gas linear-statistic fluctuations}

\entrylabel{Primary source}

Serfaty:\allowbreak{} Lectures on Coulomb and Riesz gases,\allowbreak{}2024;\allowbreak{} Section10.3,\allowbreak{}PDF223


\url{https://arxiv.org/pdf/2407.21194}


Source statement date:\allowbreak{} 2024-07-30. Archived access:\allowbreak{} 2026-09-07T01:\allowbreak{}48:\allowbreak{}41Z.


Historical:\allowbreak{} source\_\allowbreak{}says\_\allowbreak{}open | Current:\allowbreak{} unverified | Result:\allowbreak{} unsolved


\entrylabel{Question}

Determine the sharp regularity threshold on test functions for bounded fluctuations of linear statistics of Coulomb gases,\allowbreak{} beyond the established smooth-function CLTs and selected indicator results.


\entrylabel{Scope}

The source poses a general regularity question. Different dimensions and smoothness scales are grouped. A Sobolev prediction from the limiting Gaussian field is not automatically a uniform finite-particle variance theorem.


\subsection*{P0332-A001 | Explicit model calculation or source-specific limiting obstruction}

\entrylabel{Argument or result}

In the two-dimensional bulk Gaussian-field approximation,\allowbreak{} the limiting variance is proportional to the Dirichlet energy integral|grad f|\textasciicircum{}2. Mollify an indicator of a smooth bounded region over a boundary layer of thickness epsilon. Its gradient has magnitude of order epsilon\textasciicircum{}-1 over area of order perimeter*epsilon,\allowbreak{} giving energy of order perimeter/\allowbreak{}epsilon. This demonstrates why merely approximating an indicator by smooth functions does not preserve a uniform variance bound.


\entrylabel{Second check}

In signed-distance coordinates the leading integral separates into boundary length times epsilon\textasciicircum{}-1 times the squared derivative integral of the fixed transition profile. Thus the divergence is a direct scaling calculation.


\entrylabel{Remaining obstruction}

One needs uniform estimates in the mollification scale,\allowbreak{} identification of the maximal admissible function space and treatment of boundary terms. The smooth CLT cannot be applied with an increasingly rough test function without such uniformity.


Verification:\allowbreak{} author checked partial | Independent review:\allowbreak{} none | Full target solved:\allowbreak{} False


\clearpage

\section*{P0333  Multivariate fluctuation limits for empirical transport with Euclidean cost}

\entrylabel{Primary source}

del Barrio et al:\allowbreak{} Distributional Limit Theory for Optimal Transport,\allowbreak{}July2026revision;\allowbreak{} Problem1,\allowbreak{}PDF29;\allowbreak{}Theorem3.1 baselinePDF10


\url{https://arxiv.org/pdf/2505.19104}


Source statement date:\allowbreak{} 2026-07-13. Archived access:\allowbreak{} 2026-09-07T01:\allowbreak{}49:\allowbreak{}06Z.


Historical:\allowbreak{} source\_\allowbreak{}says\_\allowbreak{}open | Current:\allowbreak{} unverified | Result:\allowbreak{} unsolved


\entrylabel{Question}

Find broad conditions in dimension d>1 under which sqrt(n)*(T1(P\_\allowbreak{}n,\allowbreak{}Q)-E T1(P\_\allowbreak{}n,\allowbreak{}Q)) converges weakly,\allowbreak{} and identify when its limit is Gaussian.


\entrylabel{Scope}

T1 uses Euclidean distance,\allowbreak{} not its square. Known discrete-support and one-dimensional limits are excluded from new-solution claims. Trivial point-mass targets do not answer the intended general criterion.


\subsection*{P0333-A001 | Explicit model calculation or source-specific limiting obstruction}

\entrylabel{Argument or result}

If Q is a point mass at y,\allowbreak{} T1(P\_\allowbreak{}n,\allowbreak{}Q) is exactly the sample mean of ||X\_\allowbreak{}i-y||,\allowbreak{} so the classical CLT applies under a finite second moment. This verifies one known degenerate baseline. For general Q the Efron-Stein linearization would reduce the question to convergence in L2 of empirical optimal dual potentials to a single population potential;\allowbreak{} Euclidean cost can have nonunique transport structures.


\entrylabel{Second check}

For Q=\allowbreak{}delta\_\allowbreak{}y the coupling is forced and the variance is Var(||X-y||). Conversely,\allowbreak{} on a two-point space with P=\allowbreak{}Q=\allowbreak{}(1/\allowbreak{}2,\allowbreak{}1/\allowbreak{}2),\allowbreak{} the empirical cost is the interpoint distance times|P\_\allowbreak{}n(\{x\})-1/\allowbreak{}2|,\allowbreak{} so its sqrt(n) limit is a centered absolute normal after expectation centering,\allowbreak{} not Gaussian.


\entrylabel{Remaining obstruction}

The known special cases show both behaviors but provide no general multivariate condition controlling uniqueness and convergence of the relevant dual potential. Nonuniqueness and nondifferentiability cannot be dismissed solely from finite moments.


Verification:\allowbreak{} author checked partial | Independent review:\allowbreak{} none | Full target solved:\allowbreak{} False


\clearpage

\section*{P0334  Necessary and sufficient conditions for an Lp quantile-process limit}

\entrylabel{Primary source}

del Barrio et al:\allowbreak{} Distributional Limit Theory for Optimal Transport,\allowbreak{}July2026revision;\allowbreak{} Problem2,\allowbreak{}PDF29;\allowbreak{}Sections2.2-2.4,\allowbreak{}PDF7-9


\url{https://arxiv.org/pdf/2505.19104}


Source statement date:\allowbreak{} 2026-07-13. Archived access:\allowbreak{} 2026-09-07T01:\allowbreak{}49:\allowbreak{}06Z.


Historical:\allowbreak{} source\_\allowbreak{}says\_\allowbreak{}open | Current:\allowbreak{} unverified | Result:\allowbreak{} unsolved


\entrylabel{Question}

Find necessary and sufficient conditions for sqrt(n)*(F\_\allowbreak{}n\textasciicircum{}(-1)-F\textasciicircum{}(-1)) to converge weakly as an Lp(0,\allowbreak{}1)-valued random element,\allowbreak{} removing unnecessary endpoint monotonicity assumptions.


\entrylabel{Scope}

The p=\allowbreak{}1 criterion is established in Beare-Kaji2502.01254. The remaining p>1 family is one question. The paper gives sufficient conditions involving J\_\allowbreak{}p(P) and endpoint monotonicity.


\subsection*{P0334-A001 | Explicit model calculation or source-specific limiting obstruction}

\entrylabel{Argument or result}

For a positive density f the only possible local Gaussian limit is B(t)/\allowbreak{}f(F\textasciicircum{}(-1)(t)),\allowbreak{} where B is a Brownian bridge. Its expected Lp norm to the pth power equals E|Z|\textasciicircum{}p times J\_\allowbreak{}p(P),\allowbreak{} with J\_\allowbreak{}p(P)=\allowbreak{}integral\_\allowbreak{}0\textasciicircum{}1 [t*(1-t)]\textasciicircum{}(p/\allowbreak{}2)/\allowbreak{}f(F\textasciicircum{}(-1)(t))\textasciicircum{}p dt and Z standard normal. This makes the necessary tail-integrability quantity explicit.


\entrylabel{Second check}

At each t,\allowbreak{} B(t) has the law sqrt(t*(1-t))*Z;\allowbreak{} Tonelli gives the integral identity. For a Banach-valued Gaussian limit,\allowbreak{} its norm has finite moments,\allowbreak{} so finiteness of J\_\allowbreak{}p is necessary in this Gaussian-density setting.


\entrylabel{Remaining obstruction}

Local finite-dimensional convergence and finiteness of the candidate limit norm do not establish tightness of the empirical quantile tails. The missing uniform endpoint control without monotonicity is precisely the unresolved sufficiency direction. No complete criterion is proved.


Verification:\allowbreak{} author checked partial | Independent review:\allowbreak{} none | Full target solved:\allowbreak{} False


\clearpage

\section*{P0335  Nondegenerate centered fluctuations of null quadratic empirical transport}

\entrylabel{Primary source}

del Barrio et al:\allowbreak{} Distributional Limit Theory for Optimal Transport,\allowbreak{}July2026revision;\allowbreak{} Problem3 and equation42,\allowbreak{}PDF29


\url{https://arxiv.org/pdf/2505.19104}


Source statement date:\allowbreak{} 2026-07-13. Archived access:\allowbreak{} 2026-09-07T01:\allowbreak{}49:\allowbreak{}06Z.


Historical:\allowbreak{} source\_\allowbreak{}says\_\allowbreak{}open | Current:\allowbreak{} unverified | Result:\allowbreak{} unsolved


\entrylabel{Question}

Determine whether n*(T2(P\_\allowbreak{}n,\allowbreak{}P)-E T2(P\_\allowbreak{}n,\allowbreak{}P)) has a nondegenerate limit for the uniform measure on the unit square,\allowbreak{} identify that limit,\allowbreak{} and determine the correct analogous behavior in other dimensions.


\entrylabel{Scope}

T2 is squared Euclidean transport cost. All dimension extensions are grouped. The usual sqrt(n) theorem gives zero in this null case;\allowbreak{} known one-dimensional limits are starting points.


\subsection*{P0335-A001 | Explicit model calculation or source-specific limiting obstruction}

\entrylabel{Argument or result}

The linearized matching-energy model gives n times the centered cost as a quadratic Fourier sum of terms (|Z\_\allowbreak{}k|\textasciicircum{}2-1)/\allowbreak{}|k|\textasciicircum{}2. For independent normalized Gaussian modes,\allowbreak{} the variance is proportional to sum\_\allowbreak{}(0<|k|<=\allowbreak{}K)|k|\textasciicircum{}-4,\allowbreak{} which converges in d<4,\allowbreak{} grows as log K in d=\allowbreak{}4,\allowbreak{} and as K\textasciicircum{}(d-4) in d>4. This makes a dimension-independent n normalization doubtful.


\entrylabel{Second check}

The shell at radius r contains order r\textasciicircum{}(d-1) lattice modes,\allowbreak{} so comparison with integral\_\allowbreak{}1\textasciicircum{}K r\textasciicircum{}(d-5)dr gives the three regimes. This calculation is exact for the specified independent Gaussian surrogate,\allowbreak{} not for the empirical matching cost.


\entrylabel{Remaining obstruction}

One must control the nonlinear optimal matching remainder at fluctuation scale and identify dependence among empirical Fourier modes as the cutoff grows. The Gaussian surrogate neither proves the square conjecture nor refutes a proposed empirical limit in higher dimension.


Verification:\allowbreak{} author checked partial | Independent review:\allowbreak{} none | Full target solved:\allowbreak{} False


\clearpage

\section*{P0336  Selected limit of empirical Euclidean transport between separated rectangles}

\entrylabel{Primary source}

del Barrio et al:\allowbreak{} Distributional Limit Theory for Optimal Transport,\allowbreak{}July2026revision;\allowbreak{} Problem4,\allowbreak{}PDF29-30


\url{https://arxiv.org/pdf/2505.19104}


Source statement date:\allowbreak{} 2026-07-13. Archived access:\allowbreak{} 2026-09-07T01:\allowbreak{}49:\allowbreak{}06Z.


Historical:\allowbreak{} source\_\allowbreak{}says\_\allowbreak{}open | Current:\allowbreak{} unverified | Result:\allowbreak{} unsolved


\entrylabel{Question}

For independent uniform samples from [0,\allowbreak{}1]\textasciicircum{}2 and [2,\allowbreak{}3]x[0,\allowbreak{}1],\allowbreak{} find the limiting empirical optimal transport plan for Euclidean distance as sample size tends to infinity.


\entrylabel{Scope}

The empirical plan is almost surely a map,\allowbreak{} but the population problem has many optimal couplings. The task is to identify the plan selected by empirical optimization,\allowbreak{} not merely its limiting cost2.


\subsection*{P0336-A001 | Explicit model calculation or source-specific limiting obstruction}

\entrylabel{Argument or result}

Write a horizontal displacement d=\allowbreak{}y1-x1 in[1,\allowbreak{}3] and vertical displacement h=\allowbreak{}y2-x2. The excess Euclidean cost over horizontal displacement is sqrt(d\textasciicircum{}2+\allowbreak{}h\textasciicircum{}2)-d=\allowbreak{}h\textasciicircum{}2/\allowbreak{}(sqrt(d\textasciicircum{}2+\allowbreak{}h\textasciicircum{}2)+\allowbreak{}d). Since |h|<=\allowbreak{}1,\allowbreak{} this lies between h\textasciicircum{}2/\allowbreak{}(sqrt(10)+\allowbreak{}3) and h\textasciicircum{}2/\allowbreak{}2. Therefore any sequence of couplings with cost tending to2 must have mean squared vertical displacement tending to0.


\entrylabel{Second check}

The mean horizontal displacement under any coupling of the two fixed uniform marginals is2. The displayed denominator bounds are uniform on d in[1,\allowbreak{}3],\allowbreak{}|h|<=\allowbreak{}1,\allowbreak{} so the conclusion follows by integrating the positive excess. For empirical marginals the mean horizontal displacement converges to2 by the law of large numbers,\allowbreak{} giving the same conclusion along empirical optimal plans whose costs converge to2.


\entrylabel{Remaining obstruction}

Any subsequential limit is thus supported on equal vertical coordinates,\allowbreak{} as expected,\allowbreak{} but the horizontal coupling remains undetermined. Every such population coupling costs2. Resolving the lower-order empirical selection mechanism remains open here.


Verification:\allowbreak{} author checked partial | Independent review:\allowbreak{} none | Full target solved:\allowbreak{} False


\clearpage

\section*{P0337  Consistent estimation of the Sinkhorn null-limit spectrum}

\entrylabel{Primary source}

del Barrio et al:\allowbreak{} Distributional Limit Theory for Optimal Transport,\allowbreak{}July2026revision;\allowbreak{} Theorem4.5 and following paragraph,\allowbreak{}PDF20-21


\url{https://arxiv.org/pdf/2505.19104}


Source statement date:\allowbreak{} 2026-07-13. Archived access:\allowbreak{} 2026-09-07T01:\allowbreak{}49:\allowbreak{}06Z.


Historical:\allowbreak{} source\_\allowbreak{}says\_\allowbreak{}open | Current:\allowbreak{} unverified | Result:\allowbreak{} unsolved


\entrylabel{Question}

Construct consistent estimators of the nonnegative spectral sequence determining the weighted chi-square null limit of the empirical Sinkhorn divergence.


\entrylabel{Scope}

The source writes the limit as epsilon/\allowbreak{}2 times sum lambda\_\allowbreak{}i\textasciicircum{}2*N\_\allowbreak{}i\textasciicircum{}2,\allowbreak{} with sum lambda\_\allowbreak{}i\textasciicircum{}2 finite. The target includes enough spectral control to approximate the full limiting law,\allowbreak{} not merely a fixed finite list.


\subsection*{P0337-A001 | Explicit model calculation or source-specific limiting obstruction}

\entrylabel{Argument or result}

If a finite-rank approximation of the relevant positive operator were consistently estimated in operator norm,\allowbreak{} Weyl bounds would give convergence of every fixed ordered eigenvalue. This does not control the weighted chi-square tail. As a diagnostic,\allowbreak{} diagonal matrices with n eigenvalues all equal to1/\allowbreak{}n have operator norm tending to0 and every fixed eigenvalue tending to0,\allowbreak{} while their trace stays1.


\entrylabel{Second check}

For independent standard normals Z\_\allowbreak{}i,\allowbreak{} the expectation of sum\_\allowbreak{}(i=\allowbreak{}1)\textasciicircum{}n Z\_\allowbreak{}i\textasciicircum{}2/\allowbreak{}n is1 and its variance is2/\allowbreak{}n,\allowbreak{} so the quadratic form tends to1 despite the vanishing individual eigenvalues. This is exact and shows that coordinatewise spectral consistency is insufficient for the full law.


\entrylabel{Remaining obstruction}

The required estimator must control the trace-class weights lambda\_\allowbreak{}i\textasciicircum{}2,\allowbreak{} or an equivalent uniform tail,\allowbreak{} while estimating the population-dependent linearized Schrodinger operator. The fixed-rank perturbation argument supplies no such uniform tail bound.


Verification:\allowbreak{} author checked partial | Independent review:\allowbreak{} none | Full target solved:\allowbreak{} False


\clearpage

\section*{P0338  Minimax rates for conditional optimal-transport maps}

\entrylabel{Primary source}

del Barrio et al:\allowbreak{} Distributional Limit Theory for Optimal Transport,\allowbreak{}July2026revision;\allowbreak{} Section6.2 concluding paragraph,\allowbreak{}PDF29


\url{https://arxiv.org/pdf/2505.19104}


Source statement date:\allowbreak{} 2026-07-13. Archived access:\allowbreak{} 2026-09-07T01:\allowbreak{}49:\allowbreak{}06Z.


Historical:\allowbreak{} source\_\allowbreak{}says\_\allowbreak{}open | Current:\allowbreak{} unverified | Result:\allowbreak{} unsolved


\entrylabel{Question}

Determine sharp minimax rates for estimating conditional optimal-transport maps in the regression setting referenced by the source.


\entrylabel{Scope}

The source notes available nonsharp mixing-based rates. A meaningful minimax statement must specify covariate smoothness,\allowbreak{} response regularity,\allowbreak{} loss and reference measure;\allowbreak{} the entry records that unresolved model-class dependence rather than inventing a universal exponent.


\subsection*{P0338-A001 | Explicit model calculation or source-specific limiting obstruction}

\entrylabel{Argument or result}

In a location-family submodel Y=\allowbreak{}m(X)+\allowbreak{}Z with known noise law and a reference matching Z,\allowbreak{} the conditional transport map is T\_\allowbreak{}x(u)=\allowbreak{}m(x)+\allowbreak{}u. Thus estimating the map under integrated squared error is at least as hard as estimating the regression function m:\allowbreak{} the map difference is exactly the shift-estimation error for every reference point u.


\entrylabel{Second check}

Integrating |T\_\allowbreak{}hat\_\allowbreak{}x(u)-T\_\allowbreak{}x(u)|\textasciicircum{}2 in the shift submodel over a unit-mass reference measure gives |m\_\allowbreak{}hat(x)-m(x)|\textasciicircum{}2. For an arbitrary map estimator,\allowbreak{} averaging T\_\allowbreak{}hat\_\allowbreak{}x(u)-u over the reference law produces an estimator of m(x),\allowbreak{} whose squared error is bounded by the map loss by Jensen. Thus a uniformly faster map estimator would imply an equally faster regression estimator on that submodel.


\entrylabel{Remaining obstruction}

This reduction supplies a lower-bound route but no sharp rate for the general conditional map. One must combine covariate localization with the response-dimensional transport difficulty and prove a matching estimator under precisely specified regularity assumptions.


Verification:\allowbreak{} author checked partial | Independent review:\allowbreak{} none | Full target solved:\allowbreak{} False


\clearpage

\section*{P0339  Sample complexity of Gromov-Wasserstein plans and dual potentials}

\entrylabel{Primary source}

del Barrio et al:\allowbreak{} Distributional Limit Theory for Optimal Transport,\allowbreak{}July2026revision;\allowbreak{} Section5.2 concluding paragraph,\allowbreak{}PDF27


\url{https://arxiv.org/pdf/2505.19104}


Source statement date:\allowbreak{} 2026-07-13. Archived access:\allowbreak{} 2026-09-07T01:\allowbreak{}49:\allowbreak{}06Z.


Historical:\allowbreak{} source\_\allowbreak{}says\_\allowbreak{}open | Current:\allowbreak{} unverified | Result:\allowbreak{} unsolved


\entrylabel{Question}

Determine sample-complexity bounds for Gromov-Wasserstein optimal plans and associated dual potentials,\allowbreak{} beyond known bounds for the optimal cost.


\entrylabel{Scope}

The source explicitly says this complexity is unknown. Plans and potentials are grouped as the source optimizer-estimation question. Nonuniqueness and metric-space isometries require an identified loss or a set-valued target.


\subsection*{P0339-A001 | Explicit model calculation or source-specific limiting obstruction}

\entrylabel{Argument or result}

Cost convergence cannot determine a selected optimizer when symmetries create multiple optima. For a uniform two-point metric space with distance1,\allowbreak{} both the identity coupling and the swap coupling have zero Gromov-Wasserstein distortion between the space and itself. Their supports are disjoint. Alternating those two plans preserves exactly zero cost and has no limit as a labelled coupling.


\entrylabel{Second check}

Each coupling pairs points by an isometry,\allowbreak{} so the difference of pairwise distances is identically zero. The two atomic couplings assign masses1/\allowbreak{}2 to disjoint pairs and have total variation distance1.


\entrylabel{Remaining obstruction}

A rate theorem needs either a quotient by symmetries,\allowbreak{} distance to the full optimizer set,\allowbreak{} or an explicit unique selection with stability assumptions. The known cost rate alone cannot give a plan or potential rate in a fixed labelled norm.


Verification:\allowbreak{} author checked partial | Independent review:\allowbreak{} none | Full target solved:\allowbreak{} False


\clearpage

\section*{P0340  Growth of the minimal drawn length of critical beta-splitting trees}

\entrylabel{Primary source}

Aldous-Janson:\allowbreak{} Critical beta-splitting random tree II,\allowbreak{} overview and open problems;\allowbreak{} Open Problem14,\allowbreak{}PDF54-55


\url{https://arxiv.org/pdf/2303.02529}


Source statement date:\allowbreak{} 2025-04-17. Archived access:\allowbreak{} 2026-09-07T02:\allowbreak{}02:\allowbreak{}03Z.


Historical:\allowbreak{} source\_\allowbreak{}says\_\allowbreak{}open | Current:\allowbreak{} unverified | Result:\allowbreak{} candidate\_\allowbreak{}full\_\allowbreak{}solution


\entrylabel{Question}

Determine the expected total vertical length in the minimal integer-height drawing of DTCS(n);\allowbreak{} the source specifically proposes growth on the scale n*(log n)\textasciicircum{}4.


\entrylabel{Scope}

The source drawing recursion is authoritative. This is expected drawn length,\allowbreak{} not ordinary total edge count or the continuous-time total length.


\subsection*{P0340-A001 | Exact tree identity,\allowbreak{} source-specific recurrence or explicit limiting obstruction}

\entrylabel{Argument or result}

For every full binary tree,\allowbreak{} drawn length equals the sum of all internal subtree heights plus the root height,\allowbreak{} and is at most the sum of all leaf depths. Hence its expectation is O(n*log\textasciicircum{}2 n),\allowbreak{} already refuting the proposed rate. A stronger proof gives E drawn length/\allowbreak{}n -> sum\_\allowbreak{}(m>=\allowbreak{}2) a(m)*E height(T\_\allowbreak{}m)/\allowbreak{}m,\allowbreak{} a finite constant at least2. A geometric halving bound proves E height(T\_\allowbreak{}m)=\allowbreak{}O(log\textasciicircum{}2 m),\allowbreak{} and a uniform dyadic occupation bound controls the series tail.


\entrylabel{Second check}

Independent induction and telescoping give the same exact tree inequality. Exact checks passed for6918 ordered binary trees and23 rational height/\allowbreak{}occupation systems;\allowbreak{} five deliberately false identities were rejected. The initial float-contaminated checker failure is retained.


\entrylabel{Remaining obstruction}

The source scale is contradicted and a linear expectation limit is derived,\allowbreak{} but there is no independent review or originality certificate. No LLN,\allowbreak{} variance,\allowbreak{} CLT or closed form for the series constant is claimed.


Verification:\allowbreak{} author checked candidate | Independent review:\allowbreak{} none | Full target solved:\allowbreak{} False


\entrylabel{Full proof packet}

proofs/\allowbreak{}P0340\_\allowbreak{}beta\_\allowbreak{}tree\_\allowbreak{}drawn\_\allowbreak{}length\_\allowbreak{}counterexample\_\allowbreak{}candidate.md


\clearpage

\section*{P0341  Uniform occupation estimates and the polynomial-test ansatz}

\entrylabel{Primary source}

Aldous-Janson:\allowbreak{} Critical beta-splitting random tree II,\allowbreak{} overview and open problems;\allowbreak{} Open Problem8 and Ansatz20,\allowbreak{}PDF39-40


\url{https://arxiv.org/pdf/2303.02529}


Source statement date:\allowbreak{} 2025-04-17. Archived access:\allowbreak{} 2026-09-07T02:\allowbreak{}02:\allowbreak{}03Z.


Historical:\allowbreak{} source\_\allowbreak{}says\_\allowbreak{}open | Current:\allowbreak{} unverified | Result:\allowbreak{} candidate\_\allowbreak{}partial\_\allowbreak{}counterexample


\entrylabel{Question}

Find quantitative bounds on the tagged harmonic-descent occupation probabilities and determine the validity of the displayed weighted-convergence ansatz for all nonnegative polynomially bounded test sequences.


\entrylabel{Scope}

Brandenberger-Chin-Mossel2505.24821 supplies fixed-state and growing-state error bounds and a decaying-test special case. Those established advances are credited;\allowbreak{} the full source ansatz permits bounded sparse functions.


\subsection*{P0341-A001 | Exact tree identity,\allowbreak{} source-specific recurrence or explicit limiting obstruction}

\entrylabel{Argument or result}

The bounded test f(j)=\allowbreak{}1 when j is a power of2 and0 otherwise has finite limiting sum S because a(2\textasciicircum{}r)=\allowbreak{}O(r*2\textasciicircum{}-r). Along n=\allowbreak{}2\textasciicircum{}K,\allowbreak{} the certain root visit contributes1 in addition to every fixed initial set of terms tending to S. Thus the liminf is at least S+\allowbreak{}1,\allowbreak{} contradicting the proposed limit S. The test g(j)=\allowbreak{}j*f(j) separately refutes the divergent-series clause:\allowbreak{} the actual sum is at least2\textasciicircum{}K whereas the proposed truncated sum is of order K\textasciicircum{}2.


\entrylabel{Second check}

Both constructions use the exact root identity a(n,\allowbreak{}n)=\allowbreak{}1,\allowbreak{} nonnegativity and the established pointwise formula. No computed rate is needed. The stronger dyadic occupation bound in P0340 proves validity for suitable uniformly decreasing tails and explains why it does not cover these sparse root-sensitive tests.


\entrylabel{Remaining obstruction}

The unrestricted ansatz is false. A comprehensive sharp two-parameter error theory and the best admissible class of tests remain unresolved here;\allowbreak{} therefore the broader numbered problem is only partly answered.


Verification:\allowbreak{} author checked partial | Independent review:\allowbreak{} none | Full target solved:\allowbreak{} False


\entrylabel{Full proof packet}

proofs/\allowbreak{}P0341\_\allowbreak{}beta\_\allowbreak{}tree\_\allowbreak{}sparse\_\allowbreak{}occupation\_\allowbreak{}counterexample\_\allowbreak{}candidate.md


\clearpage

\section*{P0342  Transferring a stationary fringe-process CLT to finite beta-splitting trees}

\entrylabel{Primary source}

Aldous-Janson:\allowbreak{} Critical beta-splitting random tree II,\allowbreak{} overview and open problems;\allowbreak{} Open Problem12,\allowbreak{}PDF44-45


\url{https://arxiv.org/pdf/2303.02529}


Source statement date:\allowbreak{} 2025-04-17. Archived access:\allowbreak{} 2026-09-07T02:\allowbreak{}02:\allowbreak{}03Z.


Historical:\allowbreak{} source\_\allowbreak{}says\_\allowbreak{}open | Current:\allowbreak{} unverified | Result:\allowbreak{} candidate\_\allowbreak{}partial\_\allowbreak{}counterexample


\entrylabel{Question}

Determine whether a CLT for a functional along the stationary infinite fringe process necessarily gives the same CLT for its sum over the leaves of DTCS(n).


\entrylabel{Scope}

The source gives no quantitative locality condition in the numbered question. The explicit example below has a degenerate Gaussian limit;\allowbreak{} a strictly positive variance or fixed-neighborhood interpretation is retained as an unresolved scope restriction.


\subsection*{P0342-A001 | Exact tree identity,\allowbreak{} source-specific recurrence or explicit limiting obstruction}

\entrylabel{Argument or result}

Define a bounded continuous functional to be0 on infinite trees and n\textasciicircum{}-1/\allowbreak{}4*(B\_\allowbreak{}n-p\_\allowbreak{}n) on finite n-leaf trees,\allowbreak{} where B\_\allowbreak{}n indicates that the smaller root child has at most floor(sqrt n) leaves and p\_\allowbreak{}n is its exact probability. The original degree-two root and n are unchanged by leaf rerooting. Thus every infinite fringe sum is0,\allowbreak{} while the finite sum divided by sqrt n is n\textasciicircum{}1/\allowbreak{}4*(B\_\allowbreak{}n-p\_\allowbreak{}n). The split law gives p\_\allowbreak{}n ->1/\allowbreak{}2,\allowbreak{} so these finite normalized sums are not tight.


\entrylabel{Second check}

The functional is continuous at infinite trees because its magnitude is bounded by n\textasciicircum{}-1/\allowbreak{}4 and local convergence to an infinite bounded-degree tree forces n ->infinity. Both finite and infinite means are0. Sixteen exact root-split probability identities were checked;\allowbreak{} the full non-tightness proof is analytic.


\entrylabel{Remaining obstruction}

This disproves transfer if zero-variance Gaussian limits are admitted. For a class containing a nondegenerate fringe CLT functional and closed under this perturbation,\allowbreak{} a universal transfer rule also fails by a tight-difference argument. A concrete positive-variance or quantitatively local transfer theorem remains unsolved.


Verification:\allowbreak{} author checked partial | Independent review:\allowbreak{} none | Full target solved:\allowbreak{} False


\entrylabel{Full proof packet}

proofs/\allowbreak{}P0342\_\allowbreak{}beta\_\allowbreak{}tree\_\allowbreak{}fringe\_\allowbreak{}transfer\_\allowbreak{}scope\_\allowbreak{}counterexample.md


\clearpage

\section*{P0343  Correlations of critical-tree leaves at mesoscopic separation}

\entrylabel{Primary source}

Aldous-Janson:\allowbreak{} Critical beta-splitting random tree II,\allowbreak{} overview and open problems;\allowbreak{} Open Problem1,\allowbreak{}PDF22


\url{https://arxiv.org/pdf/2303.02529}


Source statement date:\allowbreak{} 2025-04-17. Archived access:\allowbreak{} 2026-09-07T02:\allowbreak{}02:\allowbreak{}03Z.


Historical:\allowbreak{} source\_\allowbreak{}says\_\allowbreak{}open | Current:\allowbreak{} unverified | Result:\allowbreak{} unsolved


\entrylabel{Question}

For two leaves separated by r\_\allowbreak{}n positions in the cyclic interval representation,\allowbreak{} prove that their continuous-height correlation tends to1-alpha when log(r\_\allowbreak{}n)/\allowbreak{}log(n) tends to alpha in[0,\allowbreak{}1].


\entrylabel{Scope}

The macroscopic alpha=\allowbreak{}1 case is already established in the corrected published Aldous-Pittel result. All separation exponents form one question,\allowbreak{} not separate variants.


\subsection*{P0343-A001 | Exact tree identity,\allowbreak{} source-specific recurrence or explicit limiting obstruction}

\entrylabel{Argument or result}

The proposed common-ancestry approximation shares a Brownian path up to time s=\allowbreak{}1-alpha and then uses two independent continuations. For X=\allowbreak{}B\_\allowbreak{}s+\allowbreak{}B1\_\allowbreak{}(1-s) and Y=\allowbreak{}B\_\allowbreak{}s+\allowbreak{}B2\_\allowbreak{}(1-s),\allowbreak{} the variances are1 and covariance is s,\allowbreak{} so the predicted correlation is exactly1-alpha. This isolates the common time that a two-tagged-clade limit must identify.


\entrylabel{Second check}

The covariance follows by expanding the product:\allowbreak{} the independent centered continuation terms vanish,\allowbreak{} leaving E B\_\allowbreak{}s\textasciicircum{}2=\allowbreak{}s. The endpoint alpha=\allowbreak{}0 gives full correlation and alpha=\allowbreak{}1 zero correlation,\allowbreak{} consistent with the heuristic limits.


\entrylabel{Remaining obstruction}

One needs a two-leaf process approximation valid for moving separations and sufficient uniform second-moment control to pass from weak limits to covariances. A one-leaf CLT alone proves neither requirement.


Verification:\allowbreak{} author checked partial | Independent review:\allowbreak{} none | Full target solved:\allowbreak{} False


\clearpage

\section*{P0344  Heights of prescribed interval-labelled leaves in the critical tree}

\entrylabel{Primary source}

Aldous-Janson:\allowbreak{} Critical beta-splitting random tree II,\allowbreak{} overview and open problems;\allowbreak{} Open Problem2,\allowbreak{}PDF22-23


\url{https://arxiv.org/pdf/2303.02529}


Source statement date:\allowbreak{} 2025-04-17. Archived access:\allowbreak{} 2026-09-07T02:\allowbreak{}02:\allowbreak{}03Z.


Historical:\allowbreak{} source\_\allowbreak{}says\_\allowbreak{}open | Current:\allowbreak{} unverified | Result:\allowbreak{} unsolved


\entrylabel{Question}

Analyze the height distribution of a prescribed leaf i(n) in the original interval-splitting DTCS(n) and CTCS(n) models.


\entrylabel{Scope}

Uniformly randomized leaf labels change the question. Boundary leaves and moving interior labels are one prescribed-location family.


\subsection*{P0344-A001 | Exact tree identity,\allowbreak{} source-specific recurrence or explicit limiting obstruction}

\entrylabel{Argument or result}

For the leftmost leaf the next size has the un-biased split law q(m,\allowbreak{}i),\allowbreak{} whereas a uniformly selected leaf uses size bias. Conditional on choosing the smaller left child,\allowbreak{} the leading weights are proportional to1/\allowbreak{}i. Therefore for0<u<1,\allowbreak{} the probability of I<=\allowbreak{}m\textasciicircum{}u,\allowbreak{} conditional on I<=\allowbreak{}m/\allowbreak{}2,\allowbreak{} tends to u,\allowbreak{} because the ratio of the relevant harmonic sums tends to u. This yields the proposed logarithmic contraction log I/\allowbreak{}log m ->Uniform(0,\allowbreak{}1). Taking another logarithm suggests a renewal drift of1/\allowbreak{}2 per split and height around2 log log n.


\entrylabel{Second check}

The factor1/\allowbreak{}2 comes from left-right symmetry;\allowbreak{} the mean of minus log U is1. Both facts are exact in the limiting surrogate. This explains why importing the uniform-leaf log\textasciicircum{}2 n scale into the boundary problem is invalid.


\entrylabel{Remaining obstruction}

The state-dependent renewal approximation must be controlled through the terminal small-size regime,\allowbreak{} and moving interior locations require an additional position coordinate. No full law,\allowbreak{} CLT or continuous-time analogue is proved here.


Verification:\allowbreak{} author checked partial | Independent review:\allowbreak{} none | Full target solved:\allowbreak{} False


\clearpage

\section*{P0345  The leading constant for maximum discrete critical-tree height}

\entrylabel{Primary source}

Aldous-Janson:\allowbreak{} Critical beta-splitting random tree II,\allowbreak{} overview and open problems;\allowbreak{} Open Problem4,\allowbreak{}PDF28


\url{https://arxiv.org/pdf/2303.02529}


Source statement date:\allowbreak{} 2025-04-17. Archived access:\allowbreak{} 2026-09-07T02:\allowbreak{}02:\allowbreak{}03Z.


Historical:\allowbreak{} source\_\allowbreak{}says\_\allowbreak{}open | Current:\allowbreak{} unverified | Result:\allowbreak{} unsolved


\entrylabel{Question}

Prove that the maximum graph height of DTCS(n),\allowbreak{} divided by(log n)\textasciicircum{}2,\allowbreak{} converges in probability and identify its constant.


\entrylabel{Scope}

The continuous-time maximum question was answered in the August2026 preprint;\allowbreak{} it is a different observable and is excluded separately. All conjectured constants for the discrete law are one problem.


\subsection*{P0345-A001 | Exact tree identity,\allowbreak{} source-specific recurrence or explicit limiting obstruction}

\entrylabel{Argument or result}

A tagged chain halves its size with conditional probability at least1/\allowbreak{}(2H\_\allowbreak{}(n-1));\allowbreak{} coupling to a negative-binomial waiting time and taking a union bound over leaves gives E maximum height=\allowbreak{}O(log\textasciicircum{}2 n). The uniform random-leaf CLT gives the lower bound maximum height/\allowbreak{}(log n)\textasciicircum{}2 >=\allowbreak{}3/\allowbreak{}pi\textasciicircum{}2-o(1) in probability. Thus the order can be obtained without identifying the extremal constant.


\entrylabel{Second check}

The upper bound is proved in detail in P0340 using a geometric exponential moment,\allowbreak{} not by integrating a probability-only asymptotic. The lower bound uses that the maximum dominates every sampled leaf and that the leaf fluctuations are of smaller order(log n)\textasciicircum{}1.5.


\entrylabel{Remaining obstruction}

The upper estimate is too loose to determine a common limit. Dependence among candidate deep branches and the tail behavior beyond a typical-leaf CLT remain the obstruction.


Verification:\allowbreak{} author checked partial | Independent review:\allowbreak{} none | Full target solved:\allowbreak{} False


\clearpage

\section*{P0346  Distinct and repeated clade shapes in critical beta-splitting trees}

\entrylabel{Primary source}

Aldous-Janson:\allowbreak{} Critical beta-splitting random tree II,\allowbreak{} overview and open problems;\allowbreak{} Open Problem7,\allowbreak{}PDF39


\url{https://arxiv.org/pdf/2303.02529}


Source statement date:\allowbreak{} 2025-04-17. Archived access:\allowbreak{} 2026-09-07T02:\allowbreak{}02:\allowbreak{}03Z.


Historical:\allowbreak{} source\_\allowbreak{}says\_\allowbreak{}open | Current:\allowbreak{} unverified | Result:\allowbreak{} unsolved


\entrylabel{Question}

Determine the distributions and growth of the number of distinct clade shapes,\allowbreak{} the largest repeated clade and the smallest missing clade shape in DTCS(n).


\entrylabel{Scope}

The three linked occupancy observables are grouped as one source question. The size-one clade and fixed-shape CLTs are known baselines.


\subsection*{P0346-A001 | Exact tree identity,\allowbreak{} source-specific recurrence or explicit limiting obstruction}

\entrylabel{Argument or result}

For any fixed clade shape chi of size m with positive model probability,\allowbreak{} its expected count is asymptotically n*p(chi)/\allowbreak{}m. The later fixed-clade variance bound is O(n),\allowbreak{} so Chebyshev gives probability O(1/\allowbreak{}n) that its count is zero. For any fixed finite list of distinct positive-probability shapes,\allowbreak{} all are therefore present with probability tending to1. It follows that the number of distinct shapes tends to infinity in probability,\allowbreak{} while it is always at most n-1 for internal clades.


\entrylabel{Second check}

For a count N with E N\textasciitilde{}c*n and Var N=\allowbreak{}O(n),\allowbreak{} P(N=\allowbreak{}0)<=\allowbreak{}Var N/\allowbreak{}(E N)\textasciicircum{}2=\allowbreak{}O(1/\allowbreak{}n). A finite union bound then proves the stated divergence without assuming independence between shapes.


\entrylabel{Remaining obstruction}

This supplies no growth rate:\allowbreak{} the constants and count estimates are not uniform over growing shapes. Largest repetitions and smallest missing shapes need joint tail information across an exponentially growing family.


Verification:\allowbreak{} author checked partial | Independent review:\allowbreak{} none | Full target solved:\allowbreak{} False


\clearpage

\section*{P0347  Distributions of balance indices for critical beta-splitting trees}

\entrylabel{Primary source}

Aldous-Janson:\allowbreak{} Critical beta-splitting random tree II,\allowbreak{} overview and open problems;\allowbreak{} Open Problem10,\allowbreak{}PDF41-43


\url{https://arxiv.org/pdf/2303.02529}


Source statement date:\allowbreak{} 2025-04-17. Archived access:\allowbreak{} 2026-09-07T02:\allowbreak{}02:\allowbreak{}03Z.


Historical:\allowbreak{} source\_\allowbreak{}says\_\allowbreak{}open | Current:\allowbreak{} unverified | Result:\allowbreak{} candidate\_\allowbreak{}partial\_\allowbreak{}counterexample


\entrylabel{Question}

Analyze the distributions of the source balance indices under DTCS(n),\allowbreak{} beyond the preliminary expectation heuristics.


\entrylabel{Scope}

Indices and parameter variants are grouped. Some source one-split constants are incorrect;\allowbreak{} correcting them is partial progress,\allowbreak{} not a solution of the distributional problem.


\subsection*{P0347-A001 | Exact tree identity,\allowbreak{} source-specific recurrence or explicit limiting obstruction}

\entrylabel{Argument or result}

For a split I distributed by q(m,\allowbreak{}i),\allowbreak{} the exact quadratic toll is E(m-2I)\textasciicircum{}2=\allowbreak{}m\textasciicircum{}2-2m*(m-1)/\allowbreak{}H\_\allowbreak{}(m-1),\allowbreak{} hence asymptotic to m\textasciicircum{}2 rather than the source m\textasciicircum{}2/\allowbreak{}4. Also E min(I,\allowbreak{}m-I)<=\allowbreak{}m/\allowbreak{}H\_\allowbreak{}(m-1),\allowbreak{} so E|m-2I|=\allowbreak{}m-O(m/\allowbreak{}log m),\allowbreak{} rather than m/\allowbreak{}2. If S is the external path length and C the Colless index,\allowbreak{} S-C is the sum of2*min(child sizes). Therefore E(S-C)<=\allowbreak{}2n*sum\_\allowbreak{}m a(n,\allowbreak{}m)/\allowbreak{}H\_\allowbreak{}(m-1)=\allowbreak{}2n*E continuous leaf depth=\allowbreak{}O(n log n). Since E S\textasciitilde{}(3/\allowbreak{}pi\textasciicircum{}2)n log\textasciicircum{}2 n,\allowbreak{} E C has that same leading constant,\allowbreak{} twice the source heuristic coefficient.


\entrylabel{Second check}

The quadratic identity follows by E[I(m-I)]=\allowbreak{}m*(m-1)/\allowbreak{}(2H\_\allowbreak{}(m-1)). The Colless bound uses only the exact clade-count identity and the established mean continuous leaf-depth theorem. It does not rely on the false unrestricted occupation ansatz.


\entrylabel{Remaining obstruction}

The mean Colless correction and one-split quadratic correction do not identify variances,\allowbreak{} normalized laws or joint distributions of the balance indices. Those remain unsolved here.


Verification:\allowbreak{} author checked partial | Independent review:\allowbreak{} none | Full target solved:\allowbreak{} False


\clearpage

\section*{P0348  The drawn width profile of a critical beta-splitting cladogram}

\entrylabel{Primary source}

Aldous-Janson:\allowbreak{} Critical beta-splitting random tree II,\allowbreak{} overview and open problems;\allowbreak{} Open Problem15,\allowbreak{}PDF55


\url{https://arxiv.org/pdf/2303.02529}


Source statement date:\allowbreak{} 2025-04-17. Archived access:\allowbreak{} 2026-09-07T02:\allowbreak{}02:\allowbreak{}03Z.


Historical:\allowbreak{} source\_\allowbreak{}says\_\allowbreak{}open | Current:\allowbreak{} unverified | Result:\allowbreak{} unsolved


\entrylabel{Question}

Analyze the process counting vertical drawing segments crossing each integer height interval in the minimal cladogram representation of DTCS(n).


\entrylabel{Scope}

The process is measured upward from leaves in the source minimal drawing. Ordinary breadth by distance from the original root is a different quantity.


\subsection*{P0348-A001 | Exact tree identity,\allowbreak{} source-specific recurrence or explicit limiting obstruction}

\entrylabel{Argument or result}

For every finite tree,\allowbreak{} W(h)=\allowbreak{}number of internal clades whose own graph height exceeds h,\allowbreak{} plus1 if the whole-tree height exceeds h. Thus E W\_\allowbreak{}n(h)/\allowbreak{}n converges,\allowbreak{} for each fixed nonnegative integer h,\allowbreak{} to w(h)=\allowbreak{}sum\_\allowbreak{}(m>=\allowbreak{}2) a(m)*P(height(T\_\allowbreak{}m)>h)/\allowbreak{}m. The dyadic occupation bound in P0340 controls the tails uniformly,\allowbreak{} justifying the passage from fixed sizes to the full sum. Summing the limiting mean profile over h gives the linear expected-length constant from P0340.


\entrylabel{Second check}

The identity follows by cutting every edge at level h:\allowbreak{} the surviving portion above that level is a full binary tree with one more boundary branch than internal nodes,\allowbreak{} unless it is empty. It gives W(0)=\allowbreak{}n and sum\_\allowbreak{}h W(h)=\allowbreak{}drawn length;\allowbreak{} balanced and comb trees verify both relations.


\entrylabel{Remaining obstruction}

Only the limiting mean profile and its summability are obtained. Concentration,\allowbreak{} covariance,\allowbreak{} a random profile scaling limit and height regimes growing with n remain unproved.


Verification:\allowbreak{} author checked partial | Independent review:\allowbreak{} none | Full target solved:\allowbreak{} False


\clearpage

\section*{P0349  Complex powers of clade sizes in the critical tree}

\entrylabel{Primary source}

Aldous-Janson:\allowbreak{} Critical beta-splitting random tree II,\allowbreak{} overview and open problems;\allowbreak{} Open Problem16,\allowbreak{}PDF56


\url{https://arxiv.org/pdf/2303.02529}


Source statement date:\allowbreak{} 2025-04-17. Archived access:\allowbreak{} 2026-09-07T02:\allowbreak{}02:\allowbreak{}03Z.


Historical:\allowbreak{} source\_\allowbreak{}says\_\allowbreak{}open | Current:\allowbreak{} unverified | Result:\allowbreak{} unsolved


\entrylabel{Question}

Give a detailed asymptotic analysis of the sum of complex powers of all internal clade sizes in DTCS(n).


\entrylabel{Scope}

All powers form one family. The p=\allowbreak{}1 expectation and polynomially decaying weighted-sum cases are established baselines;\allowbreak{} they are not new open solutions.


\subsection*{P0349-A001 | Exact tree identity,\allowbreak{} source-specific recurrence or explicit limiting obstruction}

\entrylabel{Argument or result}

For real p<1,\allowbreak{} the normalized expectation is sum\_\allowbreak{}m a(n,\allowbreak{}m)*m\textasciicircum{}(p-1). A dyadic block contains at most O(log m) expected tagged visits,\allowbreak{} uniformly in n,\allowbreak{} so its contribution is O(log m*m\textasciicircum{}(p-1));\allowbreak{} this is summable over dyadic scales. Hence the limit equals sum\_\allowbreak{}m a(m)*m\textasciicircum{}(p-1),\allowbreak{} consistent with the later decaying-test theorem. At p=\allowbreak{}1 the known expectation grows as n*log\textasciicircum{}2 n. For p>1,\allowbreak{} the same block estimate yields an upper bound O(n\textasciicircum{}p log n),\allowbreak{} while the root alone supplies n\textasciicircum{}p.


\entrylabel{Second check}

The dyadic sum is a convergent geometric series when p<1 and is dominated by its largest block when p>1. For complex alpha,\allowbreak{} absolute tail control uses Re(alpha);\allowbreak{} this is a bound on expectations,\allowbreak{} not a complex distributional limit.


\entrylabel{Remaining obstruction}

The intermediate logarithmic factor,\allowbreak{} exact constants,\allowbreak{} variances and normalized laws in the nondecaying and complex regimes remain unresolved. The known p<1 limit is only a baseline check.


Verification:\allowbreak{} author checked partial | Independent review:\allowbreak{} none | Full target solved:\allowbreak{} False


\clearpage

\section*{P0350  A branchpoint point-process scaling limit of the critical fringe tree}

\entrylabel{Primary source}

Aldous-Janson:\allowbreak{} Critical beta-splitting random tree II,\allowbreak{} overview and open problems;\allowbreak{} Open Problem17,\allowbreak{}PDF57


\url{https://arxiv.org/pdf/2303.02529}


Source statement date:\allowbreak{} 2025-04-17. Archived access:\allowbreak{} 2026-09-07T02:\allowbreak{}02:\allowbreak{}03Z.


Historical:\allowbreak{} source\_\allowbreak{}says\_\allowbreak{}open | Current:\allowbreak{} unverified | Result:\allowbreak{} unsolved


\entrylabel{Question}

Study the rescaled branchpoint-position process suggested by the infinite discrete fringe tree of the critical beta-splitting model.


\entrylabel{Scope}

The source proposes a scaling-limit program without specifying a proved normalization. The task is recorded at that scope,\allowbreak{} rather than inventing an established limit law.


\subsection*{P0350-A001 | Exact tree identity,\allowbreak{} source-specific recurrence or explicit limiting obstruction}

\entrylabel{Argument or result}

The first upward jump from a distinguished leaf to an ancestor of size j has exact probability6/\allowbreak{}[pi\textasciicircum{}2*(j-1)\textasciicircum{}2]. Consequently the first-ancestor size has tail asymptotic6/\allowbreak{}(pi\textasciicircum{}2*m) and infinite mean. This immediately rules out using a finite-mean renewal heuristic for the raw ancestor sizes.


\entrylabel{Second check}

The tail follows from integral comparison for sum\_\allowbreak{}(j>m)(j-1)\textasciicircum{}-2;\allowbreak{} summing the resulting1/\allowbreak{}m tail diverges. These are exact consequences of the source reverse harmonic-descent transition rule.


\entrylabel{Remaining obstruction}

Later upward transitions depend on the current clade size,\allowbreak{} so they are not independent copies of the first jump. A point-process limit requires a scale and joint control of these dependent jumps and their spatial positions.


Verification:\allowbreak{} author checked partial | Independent review:\allowbreak{} none | Full target solved:\allowbreak{} False


\clearpage

\section*{P0351  Connections between the critical fragmentation and beta(2,\allowbreak{}1) coalescent}

\entrylabel{Primary source}

Aldous-Janson:\allowbreak{} Critical beta-splitting random tree II,\allowbreak{} overview and open problems;\allowbreak{} Open Problem19,\allowbreak{}PDF58


\url{https://arxiv.org/pdf/2303.02529}


Source statement date:\allowbreak{} 2025-04-17. Archived access:\allowbreak{} 2026-09-07T02:\allowbreak{}02:\allowbreak{}03Z.


Historical:\allowbreak{} source\_\allowbreak{}says\_\allowbreak{}open | Current:\allowbreak{} unverified | Result:\allowbreak{} unsolved


\entrylabel{Question}

Identify precise probabilistic connections between CTCS(infinity) and the beta(2,\allowbreak{}1) coalescent beyond their already known common tagged-subordinator and leading-CLT relations.


\entrylabel{Scope}

The equality in law of selected frequency processes is already stated in the source. The open objective concerns stronger structural connections,\allowbreak{} not reproving that marginal identity.


\subsection*{P0351-A001 | Exact tree identity,\allowbreak{} source-specific recurrence or explicit limiting obstruction}

\entrylabel{Argument or result}

A literal time reversal of the finite binary fragmentation cannot equal the usual beta(2,\allowbreak{}1) coalescent:\allowbreak{} on three labelled blocks the latter has a positive rate for merging all three at once. Its Lambda-coalescent triple-merger rate is integral x*2x dx=\allowbreak{}2/\allowbreak{}3. The finite fragmentation has only binary splits,\allowbreak{} so reversing any of its jump paths produces only binary mergers.


\entrylabel{Second check}

The merger-rate calculation uses Lambda(dx)=\allowbreak{}2x dx on(0,\allowbreak{}1). A jump from three blocks to one changes the block count by2,\allowbreak{} whereas every reversed binary split changes it by1. Thus this proposed strongest pathwise identification fails already for three labels.


\entrylabel{Remaining obstruction}

This counterexample to naive reversal does not rule out a size-biased duality,\allowbreak{} a transformation of time or labels,\allowbreak{} or another coupling. The shared subordinator alone does not specify such a full-process relation.


Verification:\allowbreak{} author checked partial | Independent review:\allowbreak{} none | Full target solved:\allowbreak{} False


\clearpage

\section*{P0355  Endogeny of nonminimal discounted-tree fixed points}

\entrylabel{Primary source}

Aldous and Bandyopadhyay:\allowbreak{} A survey of max-type recursive distributional equations;\allowbreak{} Conjecture34;\allowbreak{}PDF26-30


\url{https://www.stat.berkeley.edu/~aldous/Papers/me107.pdf}


Source statement date:\allowbreak{} 2005. Archived access:\allowbreak{} 2026-09-07T01:\allowbreak{}59:\allowbreak{}45Z.


Historical:\allowbreak{} source\_\allowbreak{}says\_\allowbreak{}open | Current:\allowbreak{} unverified | Result:\allowbreak{} candidate\_\allowbreak{}partial\_\allowbreak{}counterexample


\entrylabel{Question}

Determine whether every positive-index invariant RTP from Corollary33 and condition39 is nonendogenous.


\entrylabel{Scope}

The counterexample is restricted to literal zero-mark/\allowbreak{}unit-weight boundary assumptions. The intended strictly discounted nondegenerate regime remains unsolved.


\entrylabel{Dated literature check}

Bounded primary-source and focused later-literature search on2026-09-07;\allowbreak{} later geometric endogeny and discounted boundedness advances credited. No full current-openness certificate.


\subsection*{P0355-A001 | Exact recurrence,\allowbreak{} conditional-law calculation,\allowbreak{} or source-specific limiting obstruction}

\entrylabel{Argument or result}

Set N=\allowbreak{}2,\allowbreak{} eta=\allowbreak{}0,\allowbreak{} xi1=\allowbreak{}xi2=\allowbreak{}1 and Y=\allowbreak{}1 deterministically. Every ray sum is0,\allowbreak{} the homogeneous equation fixes Y,\allowbreak{} condition39 holds,\allowbreak{} and every iterate from aY has law delta\_\allowbreak{}a. The all-a tree process is invariant and endogenous because its root is constant. Thus the literal conjecture needs an additional assumption.


\entrylabel{Second check}

The conditional root law is delta\_\allowbreak{}a and its only self-coupling is diagonal,\allowbreak{} independently confirming endogeny. Positive constant marks would make the ray sum infinite;\allowbreak{} strict constant discounts would destroy the chosen positive homogeneous fixed point.


\entrylabel{Remaining obstruction}

This is a degenerate scope correction,\allowbreak{} not a solution of the substantive nondegenerate version or a novelty claim.


Verification:\allowbreak{} author checked partial | Independent review:\allowbreak{} none | Full target solved:\allowbreak{} False


\clearpage

\section*{P0356  A general boundedness criterion for discounted tree sums}

\entrylabel{Primary source}

Aldous and Bandyopadhyay:\allowbreak{} A survey of max-type recursive distributional equations;\allowbreak{} Open Problem31;\allowbreak{}PDF26-28


\url{https://www.stat.berkeley.edu/~aldous/Papers/me107.pdf}


Source statement date:\allowbreak{} 2005. Archived access:\allowbreak{} 2026-09-07T01:\allowbreak{}59:\allowbreak{}45Z.


Historical:\allowbreak{} source\_\allowbreak{}says\_\allowbreak{}open | Current:\allowbreak{} unverified | Result:\allowbreak{} unsolved


\entrylabel{Question}

Characterize when the supremum over rays of the sum of positive marks multiplied by ancestral edge weights is almost surely finite.


\entrylabel{Scope}

Includes random discounts and dependent mark/\allowbreak{}offspring vectors. Aidekon-Hu-Shi2025 and Zhang2026 solve substantial regimes,\allowbreak{} which are credited rather than counted as new results.


\entrylabel{Dated literature check}

Bounded primary-source and focused later-literature search on2026-09-07;\allowbreak{} later geometric endogeny and discounted boundedness advances credited. No full current-openness certificate.


\subsection*{P0356-A001 | Exact recurrence,\allowbreak{} conditional-law calculation,\allowbreak{} or source-specific limiting obstruction}

\entrylabel{Argument or result}

On a deterministic m-ary tree with discount c=\allowbreak{}m\textasciicircum{}-H,\allowbreak{} take positive marks with tail x\textasciicircum{}-1/\allowbreak{}H*(log x)\textasciicircum{}-r for large x. Integration of the moment tail shows E eta\textasciicircum{}(1/\allowbreak{}H) is finite exactly when r>1. Zhang2026 Theorem1 then gives boundedness exactly for r>1 in this example. Both r<=\allowbreak{}1 and r>1 have the same tail exponent zeta=\allowbreak{}1/\allowbreak{}H=\allowbreak{}tstar,\allowbreak{} so a criterion using that exponent alone cannot resolve the critical boundary.


\entrylabel{Second check}

For alpha=\allowbreak{}1/\allowbreak{}H,\allowbreak{} E eta\textasciicircum{}alpha has tail integral alpha*integral x\textasciicircum{}(alpha-1)P(eta>x)dx,\allowbreak{} which reduces to integral dx/\allowbreak{}[x*(log x)\textasciicircum{}r]. Its convergence threshold is r=\allowbreak{}1. The cited theorem applies because deterministic offspring have all moments and no extinction.


\entrylabel{Remaining obstruction}

This exact critical example is a consequence of a published preprint theorem. A necessary and sufficient criterion for general random discounts and mark/\allowbreak{}offspring dependence is not obtained.


Verification:\allowbreak{} author checked partial | Independent review:\allowbreak{} none | Full target solved:\allowbreak{} False


\clearpage

\section*{P0357  Rays with persistent extremal branching-random-walk local times}

\entrylabel{Primary source}

Aidekon-Hu-Shi:\allowbreak{} Boundedness of discounted tree sums;\allowbreak{} Open question1(i),\allowbreak{}(ii);\allowbreak{}PDF2-5


\url{https://arxiv.org/pdf/2409.01048}


Source statement date:\allowbreak{} 2025-06-09. Archived access:\allowbreak{} 2026-09-07T02:\allowbreak{}01:\allowbreak{}54Z.


Historical:\allowbreak{} source\_\allowbreak{}says\_\allowbreak{}open | Current:\allowbreak{} unverified | Result:\allowbreak{} unsolved


\entrylabel{Question}

Determine whether the supremum over rays of the normalized liminf of level local times is positive and finite,\allowbreak{} with normalization log(n)/\allowbreak{}n under Theorem1.2 and1/\allowbreak{}n\textasciicircum{}2 under Theorem1.6(ii),\allowbreak{} and identify its value.


\entrylabel{Scope}

The two regimes of this single numbered local-time question are grouped. The source already determines the corresponding levelwise maxima and certain limsups.


\entrylabel{Dated literature check}

Bounded primary-source and focused later-literature search on2026-09-07;\allowbreak{} later geometric endogeny and discounted boundedness advances credited. No full current-openness certificate.


\subsection*{P0357-A001 | Exact recurrence,\allowbreak{} conditional-law calculation,\allowbreak{} or source-specific limiting obstruction}

\entrylabel{Argument or result}

In the boundary regime Theorem1.6(ii) immediately bounds the requested supremum of liminfs between0 and tstar/\allowbreak{}(2*theta). Interchanging the optimizing ray with liminf cannot justify positivity. On the compact space of binary sequences,\allowbreak{} f\_\allowbreak{}n(x)=\allowbreak{}1\{x\_\allowbreak{}n=\allowbreak{}1 and x\_\allowbreak{}(n+\allowbreak{}1)=\allowbreak{}0\} has sup\_\allowbreak{}x f\_\allowbreak{}n=\allowbreak{}1 for every n,\allowbreak{} but every x has liminf\_\allowbreak{}n f\_\allowbreak{}n(x)=\allowbreak{}0 because the pattern10 cannot occur at every sufficiently large position.


\entrylabel{Second check}

The deterministic compact-space example uses continuous cylinder indicators,\allowbreak{} so compactness and continuity alone do not fix the quantifier exchange. The upper bound follows from liminf f\_\allowbreak{}n(x)<=\allowbreak{}liminf sup\_\allowbreak{}x f\_\allowbreak{}n(x) for every fixed ray.


\entrylabel{Remaining obstruction}

A single infinite ray must realize high occupation at all sufficiently large levels. Separate favorable rays for each level and existing limsup constructions do not provide that nested compatibility.


Verification:\allowbreak{} author checked partial | Independent review:\allowbreak{} none | Full target solved:\allowbreak{} False


\clearpage

\section*{P0358  The extremal ray-height constant in a heavy-tailed branching random walk}

\entrylabel{Primary source}

Aidekon-Hu-Shi:\allowbreak{} Boundedness of discounted tree sums;\allowbreak{} Open question2;\allowbreak{}Theorem1.2;\allowbreak{}PDF2-5


\url{https://arxiv.org/pdf/2409.01048}


Source statement date:\allowbreak{} 2025-06-09. Archived access:\allowbreak{} 2026-09-07T02:\allowbreak{}01:\allowbreak{}54Z.


Historical:\allowbreak{} source\_\allowbreak{}says\_\allowbreak{}open | Current:\allowbreak{} unverified | Result:\allowbreak{} unsolved


\entrylabel{Question}

Identify the positive finite constant a1 in inf\_\allowbreak{}ray limsup\_\allowbreak{}n V(ray\_\allowbreak{}n)/\allowbreak{}sqrt(n*log n)=\allowbreak{}a1 under Theorem1.2.


\entrylabel{Scope}

This is the value of the source constant,\allowbreak{} not the already established existence theorem. Its full tilted-tail and logarithmic moment assumptions are retained by reference.


\entrylabel{Dated literature check}

Bounded primary-source and focused later-literature search on2026-09-07;\allowbreak{} later geometric endogeny and discounted boundedness advances credited. No full current-openness certificate.


\subsection*{P0358-A001 | Exact recurrence,\allowbreak{} conditional-law calculation,\allowbreak{} or source-specific limiting obstruction}

\entrylabel{Argument or result}

Rescaling every displacement by c>0 preserves the theorem class with tstar replaced by tstar/\allowbreak{}c:\allowbreak{} the tilted moment measure transforms by x->cx,\allowbreak{} its positive mean is multiplied by c,\allowbreak{} and the tail lower bound retains the allowed polynomial form after changing constants or exponents. Pathwise the variational quantity is multiplied by c. Therefore any proposed distribution-free numerical value of a1 is impossible;\allowbreak{} a valid formula must transform homogeneously under spatial scaling.


\entrylabel{Second check}

The equality inf\_\allowbreak{}ray limsup c*V\_\allowbreak{}n/\allowbreak{}sqrt(n log n)=\allowbreak{}c*a1 follows directly for c>0. The invariant product tstar*a1 is a possible dimensionless quantity,\allowbreak{} but scaling does not prove it is universal.


\entrylabel{Remaining obstruction}

The upper and lower constructions depend on rare negative drops and reproduction. Their sharp matching coefficient is not extracted here;\allowbreak{} dimensional consistency only restricts possible formulas.


Verification:\allowbreak{} author checked partial | Independent review:\allowbreak{} none | Full target solved:\allowbreak{} False


\clearpage

\section*{P0359  The multifractal spectrum of extremal ray local times}

\entrylabel{Primary source}

Aidekon-Hu-Shi:\allowbreak{} Boundedness of discounted tree sums;\allowbreak{} Open question3;\allowbreak{}PDF4-5


\url{https://arxiv.org/pdf/2409.01048}


Source statement date:\allowbreak{} 2025-06-09. Archived access:\allowbreak{} 2026-09-07T02:\allowbreak{}01:\allowbreak{}54Z.


Historical:\allowbreak{} source\_\allowbreak{}says\_\allowbreak{}open | Current:\allowbreak{} unverified | Result:\allowbreak{} unsolved


\entrylabel{Question}

Determine the Hausdorff dimension of rays with limsup,\allowbreak{} or liminf,\allowbreak{} N\_\allowbreak{}ray\textasciicircum{}n/\allowbreak{}n\textasciicircum{}2=\allowbreak{}a under Theorem1.6(ii).


\entrylabel{Scope}

All a and the related limsup/\allowbreak{}liminf spectra form one source question. Dimension uses the usual generation ultrametric on the tree boundary.


\entrylabel{Dated literature check}

Bounded primary-source and focused later-literature search on2026-09-07;\allowbreak{} later geometric endogeny and discounted boundedness advances credited. No full current-openness certificate.


\subsection*{P0359-A001 | Exact recurrence,\allowbreak{} conditional-law calculation,\allowbreak{} or source-specific limiting obstruction}

\entrylabel{Argument or result}

Write A=\allowbreak{}tstar/\allowbreak{}(2*theta). The source levelwise maximum theorem implies that for every ray limsup N\_\allowbreak{}ray\textasciicircum{}n/\allowbreak{}n\textasciicircum{}2<=\allowbreak{}A,\allowbreak{} on one common probability-one event. Hence the level sets in the question are empty for a>A. It is insufficient to count high-local-time vertices at isolated levels:\allowbreak{} Hausdorff covers use generation depths,\allowbreak{} whereas the local-time level n is a spatial coordinate and its visits may occur over widely separated generations.


\entrylabel{Second check}

The uniform bound follows by dominating every ray local time by the same supremum at each level. No uncountable intersection is needed because the maximum bound already holds simultaneously for all rays.


\entrylabel{Remaining obstruction}

For0<a<=\allowbreak{}A,\allowbreak{} a sharp upper cover estimate and a lower-bound measure supported on compatible recurrent rays are both missing. No positive-dimensional spectrum is claimed.


Verification:\allowbreak{} author checked partial | Independent review:\allowbreak{} none | Full target solved:\allowbreak{} False


\clearpage

\section*{P0360  Does frozen-percolation endogeny imply unrestricted almost sure uniqueness?}

\entrylabel{Primary source}

Rath-Swart-Szoke:\allowbreak{} A phase transition between endogeny and nonendogeny;\allowbreak{} Question1 and discussion;\allowbreak{}PDF7-9


\url{https://arxiv.org/pdf/2103.14408}


Source statement date:\allowbreak{} 2022-10-27. Archived access:\allowbreak{} 2026-09-07T02:\allowbreak{}02:\allowbreak{}00Z.


Historical:\allowbreak{} source\_\allowbreak{}says\_\allowbreak{}open | Current:\allowbreak{} unverified | Result:\allowbreak{} unsolved


\entrylabel{Question}

When the burning-time RTP is endogenous,\allowbreak{} must all solutions to the frozen-percolation equation driven by the same innovations agree,\allowbreak{} even without stationarity,\allowbreak{} adaptation or subtree independence?


\entrylabel{Scope}

Includes the source conjecture for geometric freezing sets theta<1/\allowbreak{}2;\allowbreak{} these are not counted as separate parameter questions.


\entrylabel{Dated literature check}

Bounded primary-source and focused later-literature search on2026-09-07;\allowbreak{} later geometric endogeny and discounted boundedness advances credited. No full current-openness certificate.


\subsection*{P0360-A001 | Exact recurrence,\allowbreak{} conditional-law calculation,\allowbreak{} or source-specific limiting obstruction}

\entrylabel{Argument or result}

Let A be the antitone map on proposed frozen sets defined by the percolation equation. If F is any fixed point then repeated antitonicity gives A\textasciicircum{}(2k)(empty)<=\allowbreak{}F<=\allowbreak{}A\textasciicircum{}(2k+\allowbreak{}1)(empty),\allowbreak{} after ordering the alternating iterates by inclusion. Thus equality of the even and odd limiting envelopes would prove uniqueness without distributional assumptions. The source observes that for theta>=\allowbreak{}1/\allowbreak{}2 the second iterate is empty,\allowbreak{} so these envelopes can remain the trivial two-cycle even in an endogenous regime.


\entrylabel{Second check}

The sandwich is proved inductively:\allowbreak{} applying an order-reversing map reverses both inequalities. Endogeny only concerns the specified RTP law,\allowbreak{} so it cannot itself identify a different fixed point outside that law.


\entrylabel{Remaining obstruction}

For theta<1/\allowbreak{}2 the required collapse of the envelopes is unproved here. For1/\allowbreak{}2<=\allowbreak{}theta<=\allowbreak{}theta\_\allowbreak{}star the elementary envelope method is specifically too coarse,\allowbreak{} despite known endogeny.


Verification:\allowbreak{} author checked partial | Independent review:\allowbreak{} none | Full target solved:\allowbreak{} False


\clearpage

\section*{P0361  Uniqueness of the nondiagonal scale-invariant frozen-percolation coupling}

\entrylabel{Primary source}

Rath-Swart-Szoke:\allowbreak{} A phase transition between endogeny and nonendogeny;\allowbreak{} Conjecture15;\allowbreak{}PDF10-11


\url{https://arxiv.org/pdf/2103.14408}


Source statement date:\allowbreak{} 2022-10-27. Archived access:\allowbreak{} 2026-09-07T02:\allowbreak{}02:\allowbreak{}00Z.


Historical:\allowbreak{} source\_\allowbreak{}says\_\allowbreak{}open | Current:\allowbreak{} unverified | Result:\allowbreak{} unsolved


\entrylabel{Question}

For theta\_\allowbreak{}star<theta<=\allowbreak{}1,\allowbreak{} prove that the diagonal and minimal bivariate fixed points are the only solutions in the source space M\_\allowbreak{}theta\textasciicircum{}(2).


\entrylabel{Scope}

The theorem for theta=\allowbreak{}1 is established. Symmetry,\allowbreak{} the exact marginal rho\_\allowbreak{}Xi,\allowbreak{} the union-tail bound and invariance under every geometric scaling are mandatory.


\entrylabel{Dated literature check}

Bounded primary-source and focused later-literature search on2026-09-07;\allowbreak{} later geometric endogeny and discounted boundedness advances credited. No full current-openness certificate.


\subsection*{P0361-A001 | Exact recurrence,\allowbreak{} conditional-law calculation,\allowbreak{} or source-specific limiting obstruction}

\entrylabel{Argument or result}

A tempting proposed continuum of fixed points is the convex mixture nu\_\allowbreak{}p=\allowbreak{}p*nu\_\allowbreak{}diag+\allowbreak{}(1-p)*nu\_\allowbreak{}min. It has the required marginals and scaling,\allowbreak{} but the branching update is nonlinear:\allowbreak{} with two independent child pairs,\allowbreak{} T2(nu\_\allowbreak{}p)=\allowbreak{}p\textasciicircum{}2*T2(nu\_\allowbreak{}diag)+\allowbreak{}(1-p)\textasciicircum{}2*T2(nu\_\allowbreak{}min)+\allowbreak{}2p(1-p)*C,\allowbreak{} where C is the symmetrized cross-input output law. Since the endpoints are fixed,\allowbreak{} the mixture is fixed only if C=\allowbreak{}(nu\_\allowbreak{}diag+\allowbreak{}nu\_\allowbreak{}min)/\allowbreak{}2 for0<p<1.


\entrylabel{Second check}

Expanding the four choices of child-input laws gives the identity exactly. Convexity of the admissible class does not imply convexity of the fixed-point set.


\entrylabel{Remaining obstruction}

No identity forces the cross-input law to equal the endpoint average. Controlling the nonlinear recurrence and ruling out further scale-invariant solutions remain unresolved.


Verification:\allowbreak{} author checked partial | Independent review:\allowbreak{} none | Full target solved:\allowbreak{} False


\clearpage

\section*{P0362  Constructing frozen percolation on the PWIT as a finite-graph limit}

\entrylabel{Primary source}

Frozen percolation on the binary tree is nonendogenous;\allowbreak{} Problem1;\allowbreak{}PDF12-14


\url{https://arxiv.org/pdf/1910.09213}


Source statement date:\allowbreak{} 2020-12-03. Archived access:\allowbreak{} 2026-09-07T02:\allowbreak{}40:\allowbreak{}11Z.


Historical:\allowbreak{} source\_\allowbreak{}says\_\allowbreak{}open | Current:\allowbreak{} unverified | Result:\allowbreak{} unsolved


\entrylabel{Question}

Construct frozen percolation on the Poisson weighted infinite tree and prove local weak convergence from the complete-graph lightning and deterministic-threshold models cited in Problem1.


\entrylabel{Scope}

The finite-cluster size distributions and mean-field coagulation limits in the cited literature are established,\allowbreak{} not a full marked-process local limit.


\entrylabel{Dated literature check}

Bounded primary-source and focused later-literature search on2026-09-07;\allowbreak{} later geometric endogeny and discounted boundedness advances credited. No full current-openness certificate.


\subsection*{P0362-A001 | Exact recurrence,\allowbreak{} conditional-law calculation,\allowbreak{} or source-specific limiting obstruction}

\entrylabel{Argument or result}

Before incorporating freezing,\allowbreak{} a vertex in K\_\allowbreak{}n has Binomial(n-1,\allowbreak{}t/\allowbreak{}n) incident edges activated by rescaled time t,\allowbreak{} converging to Poisson(t). Joint exploration of any fixed bounded finite neighborhood has collision probability tending to0,\allowbreak{} giving the usual PWIT local structure. Freezing is not a local continuous function of this structure:\allowbreak{} a cluster can be frozen by a mark arbitrarily far outside the explored neighborhood.


\entrylabel{Second check}

The binomial probability generating function (1+\allowbreak{}(t/\allowbreak{}n)*(z-1))\textasciicircum{}(n-1) converges to exp(t*(z-1)). This validates the local input limit but says nothing about the infinite-cluster stopping rule.


\entrylabel{Remaining obstruction}

A tight marked-process construction and control of influence from arbitrarily large components are required. Convergence of component masses alone does not identify all local edge histories.


Verification:\allowbreak{} author checked partial | Independent review:\allowbreak{} none | Full target solved:\allowbreak{} False


\clearpage

\section*{P0363  A finite-graph approximation to frozen percolation on the three-regular tree}

\entrylabel{Primary source}

Frozen percolation on the binary tree is nonendogenous;\allowbreak{} Problem2;\allowbreak{}PDF15


\url{https://arxiv.org/pdf/1910.09213}


Source statement date:\allowbreak{} 2020-12-03. Archived access:\allowbreak{} 2026-09-07T02:\allowbreak{}40:\allowbreak{}11Z.


Historical:\allowbreak{} source\_\allowbreak{}says\_\allowbreak{}open | Current:\allowbreak{} unverified | Result:\allowbreak{} unsolved


\entrylabel{Question}

Show that the canonical frozen-percolation process on the three-regular tree is the local weak limit of a suitable sequence of finite graphs.


\entrylabel{Scope}

This asks for finite graphs,\allowbreak{} not the known growing-cluster-threshold approximation on an infinite binary tree.


\entrylabel{Dated literature check}

Bounded primary-source and focused later-literature search on2026-09-07;\allowbreak{} later geometric endogeny and discounted boundedness advances credited. No full current-openness certificate.


\subsection*{P0363-A001 | Exact recurrence,\allowbreak{} conditional-law calculation,\allowbreak{} or source-specific limiting obstruction}

\entrylabel{Argument or result}

For a fixed finite ball B and bounded time t,\allowbreak{} the chance of any lightning in B is at most |B|*lambda\_\allowbreak{}n*t,\allowbreak{} which vanishes under the source scaling lambda\_\allowbreak{}n->0. Yet nontrivial freezing must survive in the limiting process. Therefore all limiting freezing in a fixed ball must be transmitted from components extending outside every fixed ball,\allowbreak{} so simply coupling local clocks and discarding distant vertices cannot prove the result.


\entrylabel{Second check}

The bound follows from1-exp(-|B|*lambda\_\allowbreak{}n*t)<=\allowbreak{}|B|*lambda\_\allowbreak{}n*t. The simultaneous condition n*lambda\_\allowbreak{}n->infinity allows giant components to freeze quickly while fixed neighborhoods see no direct marks.


\entrylabel{Remaining obstruction}

One needs quantitative control on mesoscopic component growth and transmitted freeze times. The local graph limit and vanishing local lightning do not determine the required boundary randomness.


Verification:\allowbreak{} author checked partial | Independent review:\allowbreak{} none | Full target solved:\allowbreak{} False


\clearpage

\section*{P0364  Uniqueness in law from the unoriented freezing rule and tree symmetry}

\entrylabel{Primary source}

Frozen percolation on the binary tree is nonendogenous;\allowbreak{} Question3 and Theorem1;\allowbreak{}PDF3,\allowbreak{}15


\url{https://arxiv.org/pdf/1910.09213}


Source statement date:\allowbreak{} 2020-12-03. Archived access:\allowbreak{} 2026-09-07T02:\allowbreak{}40:\allowbreak{}11Z.


Historical:\allowbreak{} source\_\allowbreak{}says\_\allowbreak{}open | Current:\allowbreak{} unverified | Result:\allowbreak{} unsolved


\entrylabel{Question}

Do the freezing rule and invariance under automorphisms of the three-regular tree uniquely determine the joint law of activation times and frozen edges?


\entrylabel{Scope}

The known characterization adds oriented subtree-independence. That additional hypothesis is exactly what this question seeks to remove.


\entrylabel{Dated literature check}

Bounded primary-source and focused later-literature search on2026-09-07;\allowbreak{} later geometric endogeny and discounted boundedness advances credited. No full current-openness certificate.


\subsection*{P0364-A001 | Exact recurrence,\allowbreak{} conditional-law calculation,\allowbreak{} or source-specific limiting obstruction}

\entrylabel{Argument or result}

Nonendogeny gives an edge event Z with positive conditional variance given all activation times U. Starting from the canonical joint law,\allowbreak{} the density1+\allowbreak{}epsilon*(Z-E[Z|U]),\allowbreak{} with0<epsilon<1,\allowbreak{} preserves the full U law and all almost sure freezing rules,\allowbreak{} but changes E Z by epsilon*E Var(Z|U)>0. This constructs another law satisfying the rule alone.


\entrylabel{Second check}

The density is positive and has conditional expectation1 given U. Its change in the selected edge mean is E[Z*(Z-E[Z|U])]=\allowbreak{}E Var(Z|U). The source already notes the resulting failure of rule-only uniqueness;\allowbreak{} this is a derivation of that known baseline.


\entrylabel{Remaining obstruction}

Selecting a distinguished edge breaks full automorphism invariance. No invariant tilt or valid invariant averaging construction is produced,\allowbreak{} so the actual rule-plus-symmetry question remains unsolved.


Verification:\allowbreak{} author checked partial | Independent review:\allowbreak{} none | Full target solved:\allowbreak{} False


\clearpage

\section*{P0365  Classifying endogenous laws among all frozen-percolation RDE solutions}

\entrylabel{Primary source}

Frozen percolation on the binary tree is nonendogenous;\allowbreak{} Question4;\allowbreak{}PDF15 and Proposition37,\allowbreak{}PDF29-30


\url{https://arxiv.org/pdf/1910.09213}


Source statement date:\allowbreak{} 2020-12-03. Archived access:\allowbreak{} 2026-09-07T02:\allowbreak{}40:\allowbreak{}11Z.


Historical:\allowbreak{} source\_\allowbreak{}says\_\allowbreak{}open | Current:\allowbreak{} unverified | Result:\allowbreak{} unsolved


\entrylabel{Question}

Determine which of the complete family of marginal solutions to the marked-binary-tree frozen-percolation RDE have endogenous RTPs.


\entrylabel{Scope}

The later geometric-freezing-set phase transition supplies a substantial subclass. General closed freezing sets and their RDE laws remain the classification target.


\entrylabel{Dated literature check}

Bounded primary-source and focused later-literature search on2026-09-07;\allowbreak{} later geometric endogeny and discounted boundedness advances credited. No full current-openness certificate.


\subsection*{P0365-A001 | Exact recurrence,\allowbreak{} conditional-law calculation,\allowbreak{} or source-specific limiting obstruction}

\entrylabel{Argument or result}

The scalar marginal equation is integral\_\allowbreak{}[0,\allowbreak{}t] s*mu(ds)=\allowbreak{}F(t)\textasciicircum{}2. It classifies one-dimensional laws but contains no conditional coupling information. Finite freezing sets produce endogenous processes by chronological induction. The geometric sets approaching0 include nonendogenous laws above theta\_\allowbreak{}star,\allowbreak{} despite each finite truncation having a deterministic chronological construction. Thus finiteness or approximation by finite freezing sets cannot alone establish endogeny of the limit.


\entrylabel{Second check}

Substituting the special law F(t)=\allowbreak{}t/\allowbreak{}2 yields t\textasciicircum{}2/\allowbreak{}4 on both sides of the scalar equation,\allowbreak{} while its RTP is known nonendogenous. A solved scalar RDE is therefore not an endogeny test.


\entrylabel{Remaining obstruction}

A condition on the geometry of a general freezing set that controls its bivariate fixed points is missing. No classification beyond the credited finite and geometric cases is obtained.


Verification:\allowbreak{} author checked partial | Independent review:\allowbreak{} none | Full target solved:\allowbreak{} False


\clearpage

\section*{P0366  The extra randomness in nonendogenous frozen percolation}

\entrylabel{Primary source}

Frozen percolation on the binary tree is nonendogenous;\allowbreak{} Question5;\allowbreak{}PDF15


\url{https://arxiv.org/pdf/1910.09213}


Source statement date:\allowbreak{} 2020-12-03. Archived access:\allowbreak{} 2026-09-07T02:\allowbreak{}40:\allowbreak{}11Z.


Historical:\allowbreak{} source\_\allowbreak{}says\_\allowbreak{}open | Current:\allowbreak{} unverified | Result:\allowbreak{} unsolved


\entrylabel{Question}

Give an explicit characterization of the randomness beyond all activation times needed to construct the binary-tree burning-time process.


\entrylabel{Scope}

An abstract regular-conditional-distribution sampler is not the requested structural characterization;\allowbreak{} the source also records triviality of the burning-process tail sigma-field.


\entrylabel{Dated literature check}

Bounded primary-source and focused later-literature search on2026-09-07;\allowbreak{} later geometric endogeny and discounted boundedness advances credited. No full current-openness certificate.


\subsection*{P0366-A001 | Exact recurrence,\allowbreak{} conditional-law calculation,\allowbreak{} or source-specific limiting obstruction}

\entrylabel{Argument or result}

Because the tree is countable and the state space is standard Borel,\allowbreak{} the whole process can formally be sampled from its conditional law using one independent Uniform(0,\allowbreak{}1),\allowbreak{} or countably many uniforms obtained from disjoint digit subsequences. This supplies an existence representation but simply encodes the unknown conditional kernel. Triviality of the process tail sigma-field rules out identifying the missing randomness with a nontrivial variable in that tail.


\entrylabel{Second check}

A measurable sampler applied to the conditional law preserves all countable-coordinate joint distributions. The sampler does not compute that law or explain its relation to cluster geometry,\allowbreak{} so it is not claimed as a solution.


\entrylabel{Remaining obstruction}

The desired canonical noise space,\allowbreak{} coupling to the activation field,\allowbreak{} and geometric interpretation of its independent information are not identified.


Verification:\allowbreak{} author checked partial | Independent review:\allowbreak{} none | Full target solved:\allowbreak{} False


\clearpage

\section*{P0367  Support of the conditional burning-time law}

\entrylabel{Primary source}

Frozen percolation on the binary tree is nonendogenous;\allowbreak{} Question6;\allowbreak{}PDF15


\url{https://arxiv.org/pdf/1910.09213}


Source statement date:\allowbreak{} 2020-12-03. Archived access:\allowbreak{} 2026-09-07T02:\allowbreak{}40:\allowbreak{}11Z.


Historical:\allowbreak{} source\_\allowbreak{}says\_\allowbreak{}open | Current:\allowbreak{} unverified | Result:\allowbreak{} unsolved


\entrylabel{Question}

Is the conditional law of the root burning time given all activation times supported on the entire interval allowed above the ordinary percolation time Xprime,\allowbreak{} including the never-burning state?


\entrylabel{Scope}

Intervals are interpreted inside the source state space [0,\allowbreak{}1] union \{infinity\};\allowbreak{} no false counterexample is based on literal real numbers between1 and infinity.


\entrylabel{Dated literature check}

Bounded primary-source and focused later-literature search on2026-09-07;\allowbreak{} later geometric endogeny and discounted boundedness advances credited. No full current-openness certificate.


\subsection*{P0367-A001 | Exact recurrence,\allowbreak{} conditional-law calculation,\allowbreak{} or source-specific limiting obstruction}

\entrylabel{Argument or result}

For a Borel set A write q\_\allowbreak{}A=\allowbreak{}P(X\_\allowbreak{}root in A|U). A conditionally independent replica gives E q\_\allowbreak{}A\textasciicircum{}2=\allowbreak{}P(X\_\allowbreak{}root in A,\allowbreak{}Xprime\_\allowbreak{}root in A),\allowbreak{} so the explicit bivariate fixed point determines the second moment of q\_\allowbreak{}A. Cauchy-Schwarz yields P(q\_\allowbreak{}A>0)>=\allowbreak{}P(X\_\allowbreak{}root in A)\textasciicircum{}2/\allowbreak{}P(X\_\allowbreak{}root in A,\allowbreak{}Xprime\_\allowbreak{}root in A). This gives a quantitative positive-probability support statement whenever the marginal mass is positive.


\entrylabel{Second check}

Apply (E q\_\allowbreak{}A)\textasciicircum{}2<=\allowbreak{}(E q\_\allowbreak{}A\textasciicircum{}2)*P(q\_\allowbreak{}A>0). The ordinary-open percolation time supplies the necessary lower support bound pathwise,\allowbreak{} but the displayed moment inequality does not imply conditional positivity almost surely.


\entrylabel{Remaining obstruction}

Full support requires ruling out random conditional gaps for every rational interval above the random lower endpoint. Second moments and nonendogeny alone do not exclude such gaps.


Verification:\allowbreak{} author checked partial | Independent review:\allowbreak{} none | Full target solved:\allowbreak{} False


\clearpage

\section*{P0368  All bivariate frozen-percolation fixed points and their attraction domains}

\entrylabel{Primary source}

Frozen percolation on the binary tree is nonendogenous;\allowbreak{} Problem7;\allowbreak{}PDF16


\url{https://arxiv.org/pdf/1910.09213}


Source statement date:\allowbreak{} 2020-12-03. Archived access:\allowbreak{} 2026-09-07T02:\allowbreak{}40:\allowbreak{}11Z.


Historical:\allowbreak{} source\_\allowbreak{}says\_\allowbreak{}open | Current:\allowbreak{} unverified | Result:\allowbreak{} unsolved


\entrylabel{Question}

Classify every fixed point and domain of attraction of d(mu\_\allowbreak{}h)/\allowbreak{}dh=\allowbreak{}T2(mu\_\allowbreak{}h)-mu\_\allowbreak{}h for frozen percolation on the oriented binary tree.


\entrylabel{Scope}

The two scale-invariant fixed points are known;\allowbreak{} this question permits nonscale-invariant couplings and asks for their dynamics.


\entrylabel{Dated literature check}

Bounded primary-source and focused later-literature search on2026-09-07;\allowbreak{} later geometric endogeny and discounted boundedness advances credited. No full current-openness certificate.


\subsection*{P0368-A001 | Exact recurrence,\allowbreak{} conditional-law calculation,\allowbreak{} or source-specific limiting obstruction}

\entrylabel{Argument or result}

Starting a law at a convex mixture of the two known fixed points does not make the evolution stationary. The two-child update has a cross-input term of weight2p(1-p),\allowbreak{} so along the mixture line the derivative is2p(1-p)*(C-(nu\_\allowbreak{}diag+\allowbreak{}nu\_\allowbreak{}min)/\allowbreak{}2). Thus an apparent line of equilibria in a linear interpolation plot can be spurious. The source suggests sharing lightning until a fixed time and separating it afterward as a possible source of other equilibria,\allowbreak{} which would require a consistent infinite-process construction.


\entrylabel{Second check}

The derivative formula follows by expanding the independent child labels and subtracting the original mixture. It is an exact signed-measure identity,\allowbreak{} not a numerical simulation.


\entrylabel{Remaining obstruction}

No invariant family containing the proposed time-split couplings has been constructed,\allowbreak{} and no global Lyapunov function or complete basin decomposition is obtained.


Verification:\allowbreak{} author checked partial | Independent review:\allowbreak{} none | Full target solved:\allowbreak{} False


\clearpage

\section*{P0369  Nonendogeny from instability of the diagonal coupling}

\entrylabel{Primary source}

Frozen percolation on the binary tree is nonendogenous;\allowbreak{} Question8;\allowbreak{}PDF16


\url{https://arxiv.org/pdf/1910.09213}


Source statement date:\allowbreak{} 2020-12-03. Archived access:\allowbreak{} 2026-09-07T02:\allowbreak{}40:\allowbreak{}11Z.


Historical:\allowbreak{} source\_\allowbreak{}says\_\allowbreak{}open | Current:\allowbreak{} unverified | Result:\allowbreak{} unsolved


\entrylabel{Question}

Find general conditions under which instability of the trivial bivariate fixed point proves that an RTP is nonendogenous.


\entrylabel{Scope}

The meaning of instability and the topology must be specified. The attempt below treats a closed scalar family of couplings only.


\entrylabel{Dated literature check}

Bounded primary-source and focused later-literature search on2026-09-07;\allowbreak{} later geometric endogeny and discounted boundedness advances credited. No full current-openness certificate.


\subsection*{P0369-A001 | Exact recurrence,\allowbreak{} conditional-law calculation,\allowbreak{} or source-specific limiting obstruction}

\entrylabel{Argument or result}

Suppose a continuous invariant family of admissible self-couplings is parameterized by disagreement d in[0,\allowbreak{}D],\allowbreak{} with diagonal d=\allowbreak{}0 and T2 acting by a continuous f with f(0)=\allowbreak{}0. If f(d)>d for all sufficiently small positive d,\allowbreak{} then because f(D)<=\allowbreak{}D the intermediate value theorem yields a nonzero fixed point. Under the standard bivariate uniqueness criterion this proves nonendogeny. A derivative fprime(0)>1 is one sufficient way to obtain the strict inequality.


\entrylabel{Second check}

Choose delta>0 with f(delta)-delta>0 and compare with f(D)-D<=\allowbreak{}0. A root in[delta,\allowbreak{}D] is a nondiagonal fixed coupling. The invariant-family assumption is essential:\allowbreak{} a chosen discrepancy statistic need not close under the bivariate update.


\entrylabel{Remaining obstruction}

For general state spaces,\allowbreak{} no scalar closure or compact trapping family is available,\allowbreak{} and linearized growth can leave the chosen family. This conditional criterion does not answer the fully general question.


Verification:\allowbreak{} author checked partial | Independent review:\allowbreak{} none | Full target solved:\allowbreak{} False


\clearpage

\section*{P0370  Endogeny of frozen percolation on higher-degree trees and the PWIT}

\entrylabel{Primary source}

Frozen percolation on the binary tree is nonendogenous;\allowbreak{} Question9;\allowbreak{}PDF16;\allowbreak{}later frozen\_\allowbreak{}transition\_\allowbreak{}2021 PDF3


\url{https://arxiv.org/pdf/1910.09213}


Source statement date:\allowbreak{} 2020-12-03. Archived access:\allowbreak{} 2026-09-07T02:\allowbreak{}40:\allowbreak{}11Z.


Historical:\allowbreak{} source\_\allowbreak{}says\_\allowbreak{}open | Current:\allowbreak{} unverified | Result:\allowbreak{} unsolved


\entrylabel{Question}

Determine endogeny for the PWIT and for d-regular trees beyond the solved degree-three case.


\entrylabel{Scope}

All degrees and the related infinite-degree limiting tree are grouped as the source asks. The degree-three theorem is excluded from the open target.


\entrylabel{Dated literature check}

Bounded primary-source and focused later-literature search on2026-09-07;\allowbreak{} later geometric endogeny and discounted boundedness advances credited. No full current-openness certificate.


\subsection*{P0370-A001 | Exact recurrence,\allowbreak{} conditional-law calculation,\allowbreak{} or source-specific limiting obstruction}

\entrylabel{Argument or result}

The proposed transfer through a near-critical binary scaling limit cannot rely merely on convergence in distribution:\allowbreak{} deterministic measurable functions of richer input fields can converge jointly to a variable with extra randomness. A concrete example is U uniform[0,\allowbreak{}1] and X\_\allowbreak{}n=\allowbreak{}1\{fractional\_\allowbreak{}part(nU)<1/\allowbreak{}2\}. Each X\_\allowbreak{}n is U-measurable,\allowbreak{} yet (U,\allowbreak{}X\_\allowbreak{}n) converges weakly to an independent pair (U,\allowbreak{}Bernoulli(1/\allowbreak{}2)),\allowbreak{} as follows by averaging any continuous test function over the n equal subintervals.


\entrylabel{Second check}

On each subinterval the first half and second half give asymptotically equal Riemann sums. The limit conditional law is nondegenerate although all finite-n conditional laws are point masses.


\entrylabel{Remaining obstruction}

A transfer of nonendogeny from the binary scaling limit needs a quantitative bivariate discrepancy that survives the approximation,\allowbreak{} not just a marginal marked-tree limit.


Verification:\allowbreak{} author checked partial | Independent review:\allowbreak{} none | Full target solved:\allowbreak{} False


\clearpage

\section*{P0371  Existence of infinite-cluster frozen percolation on higher-dimensional lattices}

\entrylabel{Primary source}

Frozen percolation on the binary tree is nonendogenous;\allowbreak{} Question10;\allowbreak{}PDF17


\url{https://arxiv.org/pdf/1910.09213}


Source statement date:\allowbreak{} 2020-12-03. Archived access:\allowbreak{} 2026-09-07T02:\allowbreak{}40:\allowbreak{}11Z.


Historical:\allowbreak{} source\_\allowbreak{}says\_\allowbreak{}open | Current:\allowbreak{} unverified | Result:\allowbreak{} unsolved


\entrylabel{Question}

For which dimensions d>=\allowbreak{}3 does nearest-neighbor Z\textasciicircum{}d admit a frozen-percolation process satisfying the source infinite-cluster freezing rule?


\entrylabel{Scope}

This is freezing at infinite size,\allowbreak{} not constant-rate cluster freezing or finite-threshold models. All dimensions are one question.


\entrylabel{Dated literature check}

Bounded primary-source and focused later-literature search on2026-09-07;\allowbreak{} later geometric endogeny and discounted boundedness advances credited. No full current-openness certificate.


\subsection*{P0371-A001 | Exact recurrence,\allowbreak{} conditional-law calculation,\allowbreak{} or source-specific limiting obstruction}

\entrylabel{Argument or result}

A compactness limit of finite-threshold processes does not automatically preserve the rule,\allowbreak{} because percolation is not a continuous local event. On Z\textasciicircum{}d,\allowbreak{} configurations with only edges inside[-n,\allowbreak{}n]\textasciicircum{}d open converge locally to the all-open configuration. Every approximating root cluster is finite,\allowbreak{} while the limiting root cluster is infinite. Thus an assertion that no finite approximation froze cannot be passed through a local limit without extra control.


\entrylabel{Second check}

For every fixed ball,\allowbreak{} all its edges are eventually open;\allowbreak{} this proves local convergence directly. Each finite box has finitely many vertices,\allowbreak{} so none of its open clusters is infinite.


\entrylabel{Remaining obstruction}

One must rule out discontinuous creation of infinite clusters at limiting times while enforcing the freezing equation. The cited high-dimensional self-destructive-percolation results do not themselves construct the required dynamics.


Verification:\allowbreak{} author checked partial | Independent review:\allowbreak{} none | Full target solved:\allowbreak{} False


\clearpage

\section*{P0372  Spherical shape and uniform bulk density of ARW aggregates}

\entrylabel{Primary source}

Levine and Silvestri:\allowbreak{} Universality conjectures for activated random walk;\allowbreak{} Conjectures 1-2;\allowbreak{} PDF3-4


\url{https://arxiv.org/pdf/2306.01698}


Source statement date:\allowbreak{} 2023-06-14. Archived access:\allowbreak{} 2026-09-07T05:\allowbreak{}09:\allowbreak{}51Z.


Historical:\allowbreak{} source\_\allowbreak{}says\_\allowbreak{}open | Current:\allowbreak{} unverified | Result:\allowbreak{} unsolved


\entrylabel{Question}

For fixed sleep rate lambda>0 on Z\textasciicircum{}d,\allowbreak{} d>=\allowbreak{}2,\allowbreak{} prove the point-source visited-set ball limit and constant-density weak scaling limit of the final sleepers,\allowbreak{} with one common density.


\entrylabel{Scope}

The one-dimensional ball/\allowbreak{}density results are credited to arw\_\allowbreak{}density\_\allowbreak{}2024,\allowbreak{} Theorem1.1 and Corollary8.9. Visited sites and final occupied sites are different observables,\allowbreak{} grouped here.


\entrylabel{Dated literature check}

Focused searches on 2026-09-07 for ARW density,\allowbreak{} cutoff,\allowbreak{} local limits,\allowbreak{} hyperuniformity and incompressibility;\allowbreak{} four later primary papers read at stated theorem scope. No certificate of present openness.


\subsection*{P0372-A001 | Conservation,\allowbreak{} explicit coupling reduction,\allowbreak{} exact law comparison or spectral calculation}

\entrylabel{Argument or result}

In a finite killed domain V,\allowbreak{} let J(x) count departures from x during stabilization of n particles at 0,\allowbreak{} and let h be the final sleeper indicator. Conservation and the fresh uniform jump choice give $\mathbb E h(x)=n1_{x=0}+\Delta\mathbb E J(x)$,\allowbreak{} where $\Delta f(x)=(2d)^{-1}\sum_{y\sim x}f(y)-f(x)$ with zero extension outside V. This follows by stopping after K departures,\allowbreak{} taking expectations of arrivals,\allowbreak{} and then passing to the finite-domain absorption limit. It identifies the candidate continuum Poisson equation,\allowbreak{} but not its unknown density or free boundary.


\entrylabel{Second check}

The finite process has finite expected absorption time:\allowbreak{} from every reachable state there is a finite positive-probability sequence moving all particles to sinks,\allowbreak{} and the state space for fixed n,\allowbreak{}V is finite. Summing the identity over V gives final mass =\allowbreak{} initial mass minus expected boundary flux,\allowbreak{} independently checking the sign.


\entrylabel{Remaining obstruction}

A limit of the expected odometer is insufficient:\allowbreak{} concentration,\allowbreak{} a uniform bulk sleeper density,\allowbreak{} and inner/\allowbreak{}outer control of the visited region are all still required.


Verification:\allowbreak{} author checked partial | Independent review:\allowbreak{} none | Full target solved:\allowbreak{} False


\clearpage

\section*{P0373  Full microscopic limit of the ARW point-source aggregate}

\entrylabel{Primary source}

Levine and Silvestri:\allowbreak{} Universality conjectures for activated random walk;\allowbreak{} Conjectures 3-5;\allowbreak{} PDF4-5


\url{https://arxiv.org/pdf/2306.01698}


Source statement date:\allowbreak{} 2023-06-14. Archived access:\allowbreak{} 2026-09-07T05:\allowbreak{}09:\allowbreak{}51Z.


Historical:\allowbreak{} source\_\allowbreak{}says\_\allowbreak{}open | Current:\allowbreak{} unverified | Result:\allowbreak{} unsolved


\entrylabel{Question}

Prove convergence of every finite-window law of the point-source stabilization,\allowbreak{} with translation-invariant limit having the aggregate density almost surely.


\entrylabel{Scope}

The three tightly linked existence/\allowbreak{}symmetry/\allowbreak{}density claims are counted once. In d=\allowbreak{}1,\allowbreak{} one-site convergence and shift invariance of every subsequential limit are known from arw\_\allowbreak{}local\_\allowbreak{}density\_\allowbreak{}2026,\allowbreak{} Theorem1.1 and Proposition1.2.


\entrylabel{Dated literature check}

Focused searches on 2026-09-07 for ARW density,\allowbreak{} cutoff,\allowbreak{} local limits,\allowbreak{} hyperuniformity and incompressibility;\allowbreak{} four later primary papers read at stated theorem scope. No certificate of present openness.


\subsection*{P0373-A001 | Conservation,\allowbreak{} explicit coupling reduction,\allowbreak{} exact law comparison or spectral calculation}

\entrylabel{Argument or result}

Compactness gives subsequential limits,\allowbreak{} but the new one-site theorem does not identify them. At density 1/\allowbreak{}2,\allowbreak{} compare iid Bernoulli occupation with a checkerboard having a uniform random phase on Z. Both are shift-invariant,\allowbreak{} ergodic under the full translation action,\allowbreak{} and have deterministic density 1/\allowbreak{}2. Both have the same one-site laws. Their adjacent covariance is respectively 0 and -1/\allowbreak{}4. Thus even the one-site theorem plus ergodicity would leave a real joint-distribution gap.


\entrylabel{Second check}

The shifted checkerboard has adjacent product identically zero and single-site mean 1/\allowbreak{}2. Translation swaps its two phases,\allowbreak{} so its only translation-invariant events have probability zero or one. These are auxiliary measures,\allowbreak{} not claimed ARW limits.


\entrylabel{Remaining obstruction}

Need uniqueness of all cylinder probabilities,\allowbreak{} or a coupling that suppresses differences between all subsequences. One-site flattening cannot provide it alone.


Verification:\allowbreak{} author checked partial | Independent review:\allowbreak{} none | Full target solved:\allowbreak{} False


\clearpage

\section*{P0374  Quadrature-domain scaling limit from extended ARW sources}

\entrylabel{Primary source}

Levine and Silvestri:\allowbreak{} Universality conjectures for activated random walk;\allowbreak{} Conjecture 6;\allowbreak{} PDF6-7


\url{https://arxiv.org/pdf/2306.01698}


Source statement date:\allowbreak{} 2023-06-14. Archived access:\allowbreak{} 2026-09-07T05:\allowbreak{}09:\allowbreak{}51Z.


Historical:\allowbreak{} source\_\allowbreak{}says\_\allowbreak{}open | Current:\allowbreak{} unverified | Result:\allowbreak{} unsolved


\entrylabel{Question}

For bounded open A with null-volume boundary,\allowbreak{} prove that stabilizing one particle at each lattice point of A has the source quadrature-domain scaling limit at density zeta\_\allowbreak{}a.


\entrylabel{Scope}

The full geometric and integrability conventions of Conjecture6 apply. This is a multiple-source obstacle problem,\allowbreak{} not another point-source radius question.


\entrylabel{Dated literature check}

Focused searches on 2026-09-07 for ARW density,\allowbreak{} cutoff,\allowbreak{} local limits,\allowbreak{} hyperuniformity and incompressibility;\allowbreak{} four later primary papers read at stated theorem scope. No certificate of present openness.


\subsection*{P0374-A001 | Conservation,\allowbreak{} explicit coupling reduction,\allowbreak{} exact law comparison or spectral calculation}

\entrylabel{Argument or result}

Let h be harmonic on a finite domain containing the evolution,\allowbreak{} and stop the process before exit and after at most K jumps. The sum of h over all particle locations is a martingale:\allowbreak{} sleeping and waking do not change locations,\allowbreak{} and each jump has zero conditional h-increment. For superharmonic h the corresponding sum is a supermartingale. Consequently,\allowbreak{} any tight scaling limit for which these stopped quantities are uniformly integrable must preserve harmonic moments and satisfy the superharmonic quadrature inequality. Testing h=\allowbreak{}1 and h=\allowbreak{}-1 forces |A|=\allowbreak{}zeta\_\allowbreak{}a*|Astar|;\allowbreak{} linear harmonic tests force equality of centroids.


\entrylabel{Second check}

The generator calculation is a sum of active-particle discrete Laplacians of h. The jump truncation and spatial stopping make optional stopping legitimate. Removing those stops is a separate step,\allowbreak{} not implicit in the displayed identity.


\entrylabel{Remaining obstruction}

No proof of tightness,\allowbreak{} uniform integrability,\allowbreak{} constant terminal density,\allowbreak{} or identification of the limiting support is obtained. Necessary moment constraints do not establish the conjectured limit.


Verification:\allowbreak{} author checked partial | Independent review:\allowbreak{} none | Full target solved:\allowbreak{} False


\clearpage

\section*{P0375  Independence of ARW bulk microstructure from the initial source shape}

\entrylabel{Primary source}

Levine and Silvestri:\allowbreak{} Universality conjectures for activated random walk;\allowbreak{} Conjecture 7;\allowbreak{} PDF7


\url{https://arxiv.org/pdf/2306.01698}


Source statement date:\allowbreak{} 2023-06-14. Archived access:\allowbreak{} 2026-09-07T05:\allowbreak{}09:\allowbreak{}51Z.


Historical:\allowbreak{} source\_\allowbreak{}says\_\allowbreak{}open | Current:\allowbreak{} unverified | Result:\allowbreak{} unsolved


\entrylabel{Question}

For an interior point of the limiting aggregate from any allowed extended source A,\allowbreak{} prove that its microscopic law is the same as the point-source law alpha.


\entrylabel{Scope}

Interior-point condition and the existence of the relevant macroscopic support are retained. The source-shape universality claim is distinct from existence of alpha.


\entrylabel{Dated literature check}

Focused searches on 2026-09-07 for ARW density,\allowbreak{} cutoff,\allowbreak{} local limits,\allowbreak{} hyperuniformity and incompressibility;\allowbreak{} four later primary papers read at stated theorem scope. No certificate of present openness.


\subsection*{P0375-A001 | Conservation,\allowbreak{} explicit coupling reduction,\allowbreak{} exact law comparison or spectral calculation}

\entrylabel{Argument or result}

A useful sufficient condition is a common local coding from instruction stacks. Suppose the output on a finite window W can be determined by stacks within a random radius R,\allowbreak{} with the same coding rule for both source geometries after their fronts pass,\allowbreak{} and P(R>r)<=\allowbreak{}C*r\textasciicircum{}(-beta). Coupling the stacks then bounds total variation of the two W-laws by 2*C*r\textasciicircum{}(-beta) whenever both available bulk neighborhoods have radius r. With r proportional to inverse mesh,\allowbreak{} this gives a quantitative microscopic universality rate.


\entrylabel{Second check}

On the intersection of the two successful coding events the outputs are identical;\allowbreak{} the union bound supplies the factor two. The argument remains valid without independence of the coding radii.


\entrylabel{Remaining obstruction}

A common finite-radius coding rule is unproved at criticality and may itself be too strong. This is an explicit coupling reduction,\allowbreak{} not an assertion that ARW has a finite influence radius.


Verification:\allowbreak{} author checked partial | Independent review:\allowbreak{} none | Full target solved:\allowbreak{} False


\clearpage

\section*{P0376  Universal critical limit obtained by sprinkling subcritical ARW states}

\entrylabel{Primary source}

Levine and Silvestri:\allowbreak{} Universality conjectures for activated random walk;\allowbreak{} Conjecture 9;\allowbreak{} PDF7


\url{https://arxiv.org/pdf/2306.01698}


Source statement date:\allowbreak{} 2023-06-14. Archived access:\allowbreak{} 2026-09-07T05:\allowbreak{}09:\allowbreak{}51Z.


Historical:\allowbreak{} source\_\allowbreak{}says\_\allowbreak{}open | Current:\allowbreak{} unverified | Result:\allowbreak{} unsolved


\entrylabel{Question}

As Poisson sprinkling raises any stationary ergodic initial state to the critical density from below,\allowbreak{} prove a common limiting stabilized law with almost-sure critical density.


\entrylabel{Scope}

The initial state,\allowbreak{} sprinkling,\allowbreak{} and local-limit order are exactly those of Conjecture9. Average density and almost-sure density are kept distinct.


\entrylabel{Dated literature check}

Focused searches on 2026-09-07 for ARW density,\allowbreak{} cutoff,\allowbreak{} local limits,\allowbreak{} hyperuniformity and incompressibility;\allowbreak{} four later primary papers read at stated theorem scope. No certificate of present openness.


\subsection*{P0376-A001 | Conservation,\allowbreak{} explicit coupling reduction,\allowbreak{} exact law comparison or spectral calculation}

\entrylabel{Argument or result}

For a stabilizable stationary system,\allowbreak{} mass transport preserves the mean particle density:\allowbreak{} send one unit from every initial particle to its final sleeping site. Thus any local subsequential limit of the bounded occupation indicators has one-site mean zeta\_\allowbreak{}c. This still does not imply almost-sure density. To see the precise gap,\allowbreak{} fill consecutive blocks of length m with iid Bernoulli(rho) values,\allowbreak{} constant within a block,\allowbreak{} and randomize the block origin uniformly. Each resulting process is stationary ergodic with density rho. As m tends to infinity its local law converges to the mixture rho*delta\_\allowbreak{}all\_\allowbreak{}ones+\allowbreak{}(1-rho)*delta\_\allowbreak{}all\_\allowbreak{}zeros,\allowbreak{} whose density is random.


\entrylabel{Second check}

A fixed window crosses a block boundary with probability at most its diameter/\allowbreak{}m. Away from a boundary it sees a single Bernoulli variable. Ergodicity of each finite-block process follows by iid block ergodicity and the cyclic action on the random phase.


\entrylabel{Remaining obstruction}

Need control of correlations or ergodic components that is uniform near criticality,\allowbreak{} together with uniqueness and independence from the initial law. Conservation of mean does not supply any of these.


Verification:\allowbreak{} author checked partial | Independent review:\allowbreak{} none | Full target solved:\allowbreak{} False


\clearpage

\section*{P0377  Universal stationary ARW density in higher-dimensional growing domains}

\entrylabel{Primary source}

Levine and Silvestri:\allowbreak{} Universality conjectures for activated random walk;\allowbreak{} Conjectures 11-12;\allowbreak{} PDF8-9


\url{https://arxiv.org/pdf/2306.01698}


Source statement date:\allowbreak{} 2023-06-14. Archived access:\allowbreak{} 2026-09-07T05:\allowbreak{}09:\allowbreak{}51Z.


Historical:\allowbreak{} source\_\allowbreak{}says\_\allowbreak{}open | Current:\allowbreak{} unverified | Result:\allowbreak{} unsolved


\entrylabel{Question}

For d>=\allowbreak{}2 and vanishing boundary-to-volume exhaustions,\allowbreak{} prove concentration of the wired stationary density and equality with the fixed-energy and aggregate critical densities.


\entrylabel{Scope}

Concentration,\allowbreak{} exhaustion independence,\allowbreak{} and equality of the related critical densities are one grouped density question. The interval/\allowbreak{}cycle result in d=\allowbreak{}1 is known.


\entrylabel{Dated literature check}

Focused searches on 2026-09-07 for ARW density,\allowbreak{} cutoff,\allowbreak{} local limits,\allowbreak{} hyperuniformity and incompressibility;\allowbreak{} four later primary papers read at stated theorem scope. No certificate of present openness.


\subsection*{P0377-A001 | Conservation,\allowbreak{} explicit coupling reduction,\allowbreak{} exact law comparison or spectral calculation}

\entrylabel{Argument or result}

The exact stationary sampler starts one active particle at every site. Write D\_\allowbreak{}V for the total number eventually killed at the boundary. Then S\_\allowbreak{}V=\allowbreak{}|V|-D\_\allowbreak{}V pathwise. If J(x) counts jumps at x,\allowbreak{} the departure identity gives E D\_\allowbreak{}V=\allowbreak{}sum\_\allowbreak{}\{x in V\} q\_\allowbreak{}V(x)*E J(x),\allowbreak{} where q\_\allowbreak{}V(x) is the one-jump exit probability. This shows why a vanishing number of boundary sites cannot by itself imply negligible dissipation:\allowbreak{} expected boundary odometers can grow enough to carry a volume-order flux.


\entrylabel{Second check}

Count exiting jumps directly for the pathwise identity,\allowbreak{} or sum the discrete Poisson equation for the expected final occupation;\allowbreak{} both yield the same boundary-flux formula. No assumption of independent final sleepers is used.


\entrylabel{Remaining obstruction}

Need a concentration estimate for D\_\allowbreak{}V and a limit for its volume-normalized expectation that is independent of the exhaustion. Surface-to-volume geometry alone does not control repeated boundary visits.


Verification:\allowbreak{} author checked partial | Independent review:\allowbreak{} none | Full target solved:\allowbreak{} False


\clearpage

\section*{P0378  Exhaustion-independent infinite-volume wired ARW measure}

\entrylabel{Primary source}

Levine and Silvestri:\allowbreak{} Universality conjectures for activated random walk;\allowbreak{} Conjecture 13;\allowbreak{} PDF9


\url{https://arxiv.org/pdf/2306.01698}


Source statement date:\allowbreak{} 2023-06-14. Archived access:\allowbreak{} 2026-09-07T05:\allowbreak{}09:\allowbreak{}51Z.


Historical:\allowbreak{} source\_\allowbreak{}says\_\allowbreak{}open | Current:\allowbreak{} unverified | Result:\allowbreak{} unsolved


\entrylabel{Question}

Prove local convergence of wired stationary measures along every exhaustion of Z\textasciicircum{}d,\allowbreak{} with the same limit.


\entrylabel{Scope}

Compactness proves existence of subsequences only. The January2026 point-source paper still identifies this as a joint-law question.


\entrylabel{Dated literature check}

Focused searches on 2026-09-07 for ARW density,\allowbreak{} cutoff,\allowbreak{} local limits,\allowbreak{} hyperuniformity and incompressibility;\allowbreak{} four later primary papers read at stated theorem scope. No certificate of present openness.


\subsection*{P0378-A001 | Conservation,\allowbreak{} explicit coupling reduction,\allowbreak{} exact law comparison or spectral calculation}

\entrylabel{Argument or result}

Finite-volume stationarity gives pi\_\allowbreak{}V*K\_\allowbreak{}v,\allowbreak{}V=\allowbreak{}pi\_\allowbreak{}V for addition at v. To pass this identity through a local limit,\allowbreak{} truncate the avalanche at a radius r. If its influence on an observed window leaves that ball with probability at most epsilon\_\allowbreak{}r,\allowbreak{} uniformly in outer V,\allowbreak{} then for each bounded local test f the discrepancy between the full and truncated expectations is at most 2*||f||\_\allowbreak{}infinity*epsilon\_\allowbreak{}r. The truncated operator is local and continuous,\allowbreak{} so subsequential limits inherit addition invariance if epsilon\_\allowbreak{}r tends to zero.


\entrylabel{Second check}

Use the same stacks up to the exit from the radius-r ball. On the no-exit event the two evaluations agree,\allowbreak{} giving the bound directly. This requires a uniform probability estimate;\allowbreak{} finite-volume almost-sure stabilization alone is not enough.


\entrylabel{Remaining obstruction}

Uniform localization is missing and may fail for critical avalanches. Even inherited addition invariance would still need a uniqueness theorem to identify all subsequential limits.


Verification:\allowbreak{} author checked partial | Independent review:\allowbreak{} none | Full target solved:\allowbreak{} False


\clearpage

\section*{P0379  Equality of the three infinite-volume critical ARW laws}

\entrylabel{Primary source}

Levine and Silvestri:\allowbreak{} Universality conjectures for activated random walk;\allowbreak{} Question 14;\allowbreak{} PDF9;\allowbreak{} Section9


\url{https://arxiv.org/pdf/2306.01698}


Source statement date:\allowbreak{} 2023-06-14. Archived access:\allowbreak{} 2026-09-07T05:\allowbreak{}09:\allowbreak{}51Z.


Historical:\allowbreak{} source\_\allowbreak{}says\_\allowbreak{}open | Current:\allowbreak{} unverified | Result:\allowbreak{} unsolved


\entrylabel{Question}

Determine whether the wired,\allowbreak{} subcritical-sprinkling,\allowbreak{} and point-source microscopic limits pi,\allowbreak{} mu,\allowbreak{} and alpha coincide.


\entrylabel{Scope}

The question includes the required existence of its limits by reference to the preceding entries. Matching the mean densities or a finite covariance table is not equality of laws.


\entrylabel{Dated literature check}

Focused searches on 2026-09-07 for ARW density,\allowbreak{} cutoff,\allowbreak{} local limits,\allowbreak{} hyperuniformity and incompressibility;\allowbreak{} four later primary papers read at stated theorem scope. No certificate of present openness.


\subsection*{P0379-A001 | Conservation,\allowbreak{} explicit coupling reduction,\allowbreak{} exact law comparison or spectral calculation}

\entrylabel{Argument or result}

Even agreement of every one- and two-site marginal need not identify a measure. On Z,\allowbreak{} partition into triples with a uniformly random phase and use iid uniform even-parity bit triples (000,\allowbreak{}011,\allowbreak{}101,\allowbreak{}110). Every pair of distinct bits is independent Bernoulli(1/\allowbreak{}2),\allowbreak{} including pairs across triples. The measure is translation invariant. For three consecutive sites,\allowbreak{} however,\allowbreak{} the all-ones probability is (1/\allowbreak{}3)*0+\allowbreak{}(2/\allowbreak{}3)*(1/\allowbreak{}8)=\allowbreak{}1/\allowbreak{}12,\allowbreak{} whereas iid Bernoulli(1/\allowbreak{}2) gives 1/\allowbreak{}8. Thus the pair-correlation experiments cannot distinguish all possible candidate laws.


\entrylabel{Second check}

For an aligned triple,\allowbreak{} 111 is forbidden. In each nonaligned phase the triple crosses a block boundary;\allowbreak{} two bits from one block are independent and the third is independent of them,\allowbreak{} giving 1/\allowbreak{}8. Averaging the three phases gives 1/\allowbreak{}12.


\entrylabel{Remaining obstruction}

Need a coupling or a characterization using all cylinder events. This auxiliary counterexample refutes only the proposed inference from pair data,\allowbreak{} not equality of the actual ARW limits.


Verification:\allowbreak{} author checked partial | Independent review:\allowbreak{} none | Full target solved:\allowbreak{} False


\clearpage

\section*{P0380  Equality of bulk one-site density and wired global ARW density}

\entrylabel{Primary source}

Levine and Silvestri:\allowbreak{} Universality conjectures for activated random walk;\allowbreak{} Conjecture 15;\allowbreak{} PDF9


\url{https://arxiv.org/pdf/2306.01698}


Source statement date:\allowbreak{} 2023-06-14. Archived access:\allowbreak{} 2026-09-07T05:\allowbreak{}09:\allowbreak{}51Z.


Historical:\allowbreak{} source\_\allowbreak{}says\_\allowbreak{}open | Current:\allowbreak{} unverified | Result:\allowbreak{} unsolved


\entrylabel{Question}

Show that the one-site occupation probability under the infinite-volume wired measure equals the limiting finite-volume stationary density.


\entrylabel{Scope}

This is the boundary-versus-bulk assertion for wired measures. The 2026 point-source one-site theorem does not directly prove it.


\entrylabel{Dated literature check}

Focused searches on 2026-09-07 for ARW density,\allowbreak{} cutoff,\allowbreak{} local limits,\allowbreak{} hyperuniformity and incompressibility;\allowbreak{} four later primary papers read at stated theorem scope. No certificate of present openness.


\subsection*{P0380-A001 | Conservation,\allowbreak{} explicit coupling reduction,\allowbreak{} exact law comparison or spectral calculation}

\entrylabel{Argument or result}

Let rho\_\allowbreak{}V(x) be the wired single-site mean. A sufficient estimate is sup\_\allowbreak{}\{dist(x,\allowbreak{}boundary V)>=\allowbreak{}r\}|rho\_\allowbreak{}V(x)-rho\_\allowbreak{}infinity|<=\allowbreak{}epsilon\_\allowbreak{}r,\allowbreak{} uniformly over V,\allowbreak{} with epsilon\_\allowbreak{}r tending to zero. Splitting the mean over the depth-r boundary layer and its complement gives |E S\_\allowbreak{}V/\allowbreak{}|V|-rho\_\allowbreak{}infinity|<=\allowbreak{}epsilon\_\allowbreak{}r+\allowbreak{}|boundary\_\allowbreak{}layer\_\allowbreak{}r(V)|/\allowbreak{}|V|. For each fixed r,\allowbreak{} the last ratio vanishes along lattice Folner exhaustions;\allowbreak{} then r can grow after V.


\entrylabel{Second check}

Occupation indicators lie between zero and one,\allowbreak{} so the entire boundary contribution is bounded by its cardinality. A fixed-depth boundary layer has size at most a constant depending on d,\allowbreak{}r times |boundary V|.


\entrylabel{Remaining obstruction}

Ordinary local convergence at a fixed origin does not give the required uniform interior estimate. A growing central defect can have zero volume fraction but determine every fixed-window limit.


Verification:\allowbreak{} author checked partial | Independent review:\allowbreak{} none | Full target solved:\allowbreak{} False


\clearpage

\section*{P0381  Higher-dimensional hockey-stick density evolution for wired ARW}

\entrylabel{Primary source}

Levine and Silvestri:\allowbreak{} Universality conjectures for activated random walk;\allowbreak{} Conjecture 17;\allowbreak{} PDF9-10


\url{https://arxiv.org/pdf/2306.01698}


Source statement date:\allowbreak{} 2023-06-14. Archived access:\allowbreak{} 2026-09-07T05:\allowbreak{}09:\allowbreak{}51Z.


Historical:\allowbreak{} source\_\allowbreak{}says\_\allowbreak{}open | Current:\allowbreak{} unverified | Result:\allowbreak{} unsolved


\entrylabel{Question}

On d-dimensional boxes with d>=\allowbreak{}2 and uniform driving from empty,\allowbreak{} prove that after floor(t*L\textasciicircum{}d) additions the density converges to min(t,\allowbreak{}zeta\_\allowbreak{}c).


\entrylabel{Scope}

The full interval theorem is proved in arw\_\allowbreak{}hockey\_\allowbreak{}2024,\allowbreak{} Theorem1. All times t are one question.


\entrylabel{Dated literature check}

Focused searches on 2026-09-07 for ARW density,\allowbreak{} cutoff,\allowbreak{} local limits,\allowbreak{} hyperuniformity and incompressibility;\allowbreak{} four later primary papers read at stated theorem scope. No certificate of present openness.


\subsection*{P0381-A001 | Conservation,\allowbreak{} explicit coupling reduction,\allowbreak{} exact law comparison or spectral calculation}

\entrylabel{Argument or result}

Writing N=\allowbreak{}L\textasciicircum{}d,\allowbreak{} S\_\allowbreak{}k for retained particles,\allowbreak{} and E\_\allowbreak{}k for cumulative losses,\allowbreak{} the exact identity E\_\allowbreak{}k=\allowbreak{}k-S\_\allowbreak{}k converts the conjecture to E\_\allowbreak{}floor(tN)/\allowbreak{}N -> (t-zeta\_\allowbreak{}c)\_\allowbreak{}+\allowbreak{}. At stationarity one addition has mean loss one,\allowbreak{} because 0=\allowbreak{}E(S\_\allowbreak{}\{k+\allowbreak{}1\}-S\_\allowbreak{}k)=\allowbreak{}1-E(loss\_\allowbreak{}\{k+\allowbreak{}1\}). That balance does not say when stationarity is reached. The necessary transient estimate is a bound on the total excess loss before zeta\_\allowbreak{}c*N and on the total retention after it.


\entrylabel{Second check}

The cumulative loss identity is pathwise,\allowbreak{} so a convergence theorem for either variable is equivalent to one for the other,\allowbreak{} including its constants. The stationary balance is legitimate on a finite state space.


\entrylabel{Remaining obstruction}

No bound makes the transient loss negligible below criticality and matches it to the linear excess above criticality in d>=\allowbreak{}2. Stationary mean balance alone cannot establish the hockey-stick curve.


Verification:\allowbreak{} author checked partial | Independent review:\allowbreak{} none | Full target solved:\allowbreak{} False


\clearpage

\section*{P0382  Critical-density total-variation cutoff for wired ARW in higher dimensions}

\entrylabel{Primary source}

Levine and Silvestri:\allowbreak{} Universality conjectures for activated random walk;\allowbreak{} Conjecture 19;\allowbreak{} PDF11


\url{https://arxiv.org/pdf/2306.01698}


Source statement date:\allowbreak{} 2023-06-14. Archived access:\allowbreak{} 2026-09-07T05:\allowbreak{}09:\allowbreak{}51Z.


Historical:\allowbreak{} source\_\allowbreak{}says\_\allowbreak{}open | Current:\allowbreak{} unverified | Result:\allowbreak{} unsolved


\entrylabel{Question}

Prove total-variation cutoff at zeta\_\allowbreak{}s times the volume for uniformly driven wired ARW on lattice balls in d>=\allowbreak{}2.


\entrylabel{Scope}

The interval theorem is established by arw\_\allowbreak{}cutoff\_\allowbreak{}2025,\allowbreak{} Theorem1.1. Separation cutoff at volume is a different metric and is not substituted.


\entrylabel{Dated literature check}

Focused searches on 2026-09-07 for ARW density,\allowbreak{} cutoff,\allowbreak{} local limits,\allowbreak{} hyperuniformity and incompressibility;\allowbreak{} four later primary papers read at stated theorem scope. No certificate of present openness.


\subsection*{P0382-A001 | Conservation,\allowbreak{} explicit coupling reduction,\allowbreak{} exact law comparison or spectral calculation}

\entrylabel{Argument or result}

Assuming stationary-density concentration at rho,\allowbreak{} the lower bound at every t<(rho-epsilon)N follows immediately and sharply from the particle-count test. Starting empty,\allowbreak{} S\_\allowbreak{}t<=\allowbreak{}t deterministically. For A=\allowbreak{}\{S>(rho-epsilon/\allowbreak{}2)N\},\allowbreak{} the time-t law gives probability zero whereas stationarity gives probability tending to one. Therefore the total-variation distance tends to one. This reduces the unresolved direction to forgetting the configuration just above rho*N.


\entrylabel{Second check}

The TV lower bound is the difference of probabilities of the single event A,\allowbreak{} with the intermediate threshold strictly between t/\allowbreak{}N and rho. No spectral approximation or mixing definition change occurs.


\entrylabel{Remaining obstruction}

Need the matching upper bound after (rho+\allowbreak{}epsilon)N additions,\allowbreak{} including spatial activation coverage or a successful coupling. A small relaxation time and the density lower bound do not by themselves prove this upper bound.


Verification:\allowbreak{} author checked partial | Independent review:\allowbreak{} none | Full target solved:\allowbreak{} False


\clearpage

\section*{P0383  Exponential incompressibility of the infinite-volume wired ARW state}

\entrylabel{Primary source}

Levine and Silvestri:\allowbreak{} Universality conjectures for activated random walk;\allowbreak{} Conjecture 20;\allowbreak{} PDF11


\url{https://arxiv.org/pdf/2306.01698}


Source statement date:\allowbreak{} 2023-06-14. Archived access:\allowbreak{} 2026-09-07T05:\allowbreak{}09:\allowbreak{}51Z.


Historical:\allowbreak{} source\_\allowbreak{}says\_\allowbreak{}open | Current:\allowbreak{} unverified | Result:\allowbreak{} unsolved


\entrylabel{Question}

For every zeta>zeta\_\allowbreak{}s,\allowbreak{} prove a volume-exponential upper tail for the number of sleepers in an L-sided cube under pi.


\entrylabel{Scope}

Conditional on existence of pi and the intended mean density. Pairwise negative covariance is not assumed to imply negative association.


\entrylabel{Dated literature check}

Focused searches on 2026-09-07 for ARW density,\allowbreak{} cutoff,\allowbreak{} local limits,\allowbreak{} hyperuniformity and incompressibility;\allowbreak{} four later primary papers read at stated theorem scope. No certificate of present openness.


\subsection*{P0383-A001 | Conservation,\allowbreak{} explicit coupling reduction,\allowbreak{} exact law comparison or spectral calculation}

\entrylabel{Argument or result}

A route through negative association would suffice:\allowbreak{} for N=\allowbreak{}L\textasciicircum{}d indicators with common mean rho and negative association,\allowbreak{} $\mathbb E e^{\theta S}\le(1-\rho+\rho e^\theta)^N$ for theta>0. Chernoff optimization at $e^\theta=\zeta(1-\rho)/(\rho(1-\zeta))$ yields $\mathbb P(S\ge\zeta N)\le e^{-N D(\zeta\Vert\rho)}$ for rho<zeta<1. Here D is Bernoulli relative entropy. At zeta=\allowbreak{}1 the bound is rho\textasciicircum{}N,\allowbreak{} obtained by a limit.


\entrylabel{Second check}

Differentiate the logarithm of the Chernoff bound to check the optimizer;\allowbreak{} strict convexity of Bernoulli entropy makes D positive for zeta>rho. The factorization uses association for products of increasing functions on disjoint coordinate sets,\allowbreak{} a stronger property than pairwise correlation signs.


\entrylabel{Remaining obstruction}

Negative association,\allowbreak{} or a comparably strong exponential-moment estimate,\allowbreak{} is unproved here for pi. A covariance table cannot justify this step.


Verification:\allowbreak{} author checked partial | Independent review:\allowbreak{} none | Full target solved:\allowbreak{} False


\clearpage

\section*{P0384  Sharp first slow-stabilization threshold of the ARW free chain}

\entrylabel{Primary source}

Levine and Silvestri:\allowbreak{} Universality conjectures for activated random walk;\allowbreak{} Conjecture 21 and Question22;\allowbreak{} PDF12


\url{https://arxiv.org/pdf/2306.01698}


Source statement date:\allowbreak{} 2023-06-14. Archived access:\allowbreak{} 2026-09-07T05:\allowbreak{}09:\allowbreak{}51Z.


Historical:\allowbreak{} source\_\allowbreak{}says\_\allowbreak{}open | Current:\allowbreak{} unverified | Result:\allowbreak{} unsolved


\entrylabel{Question}

For the free addition chain on growing tori,\allowbreak{} prove concentration of the first superlinear stabilization time at critical density,\allowbreak{} and determine the proposed linear-to-exponential jump-count transition.


\entrylabel{Scope}

The cycle fixed-energy polynomial/\allowbreak{}exponential threshold is known from arw\_\allowbreak{}density\_\allowbreak{}2024. It concerns a different initialization and does not automatically prove this first-passage statement for a chain of successive stable states.


\entrylabel{Dated literature check}

Focused searches on 2026-09-07 for ARW density,\allowbreak{} cutoff,\allowbreak{} local limits,\allowbreak{} hyperuniformity and incompressibility;\allowbreak{} four later primary papers read at stated theorem scope. No certificate of present openness.


\subsection*{P0384-A001 | Conservation,\allowbreak{} explicit coupling reduction,\allowbreak{} exact law comparison or spectral calculation}

\entrylabel{Argument or result}

For a fixed threshold f(N),\allowbreak{} define tau\_\allowbreak{}f=\allowbreak{}min\{k:\allowbreak{}U\_\allowbreak{}k>=\allowbreak{}f(N)\}. An explicit sufficient pair of bounds is sup\_\allowbreak{}\{k<=\allowbreak{}(rho-epsilon)N\} P(U\_\allowbreak{}k>=\allowbreak{}f(N))=\allowbreak{}o(1/\allowbreak{}N),\allowbreak{} and P(U\_\allowbreak{}floor((rho+\allowbreak{}epsilon)N)>=\allowbreak{}f(N))=\allowbreak{}1-o(1). The union bound gives the lower bound on tau\_\allowbreak{}f,\allowbreak{} and the single late-time event gives the upper bound. Merely proving each early probability is o(1) cannot be enough:\allowbreak{} independent rare spikes with probability 1/\allowbreak{}N per step have a nonvanishing probability of a first crossing among N steps.


\entrylabel{Second check}

The rare-spike example has no-crossing probability (1-1/\allowbreak{}N)\textasciicircum{}floor(aN) -> exp(-a). It is an auxiliary obstruction to exchanging pointwise high probability with a stopping-time bound,\allowbreak{} not an ARW construction.


\entrylabel{Remaining obstruction}

Uniform early-time tails at the needed scale,\allowbreak{} the appropriate f,\allowbreak{} and a late-time slow-stabilization estimate for the free-chain input distribution are missing.


Verification:\allowbreak{} author checked partial | Independent review:\allowbreak{} none | Full target solved:\allowbreak{} False


\clearpage

\section*{P0385  Asymptotic equivalence of supercritical free and wake ARW states}

\entrylabel{Primary source}

Levine and Silvestri:\allowbreak{} Universality conjectures for activated random walk;\allowbreak{} Conjecture 23;\allowbreak{} PDF12


\url{https://arxiv.org/pdf/2306.01698}


Source statement date:\allowbreak{} 2023-06-14. Archived access:\allowbreak{} 2026-09-07T05:\allowbreak{}09:\allowbreak{}51Z.


Historical:\allowbreak{} source\_\allowbreak{}says\_\allowbreak{}open | Current:\allowbreak{} unverified | Result:\allowbreak{} unsolved


\entrylabel{Question}

At a fixed density zeta in (zeta\_\allowbreak{}c,\allowbreak{}1),\allowbreak{} prove that the free-chain law and the stationary wake-chain law on growing tori converge in total variation.


\entrylabel{Scope}

The same particle count is essential. This question concerns the law before an extra full wake,\allowbreak{} not just mixing of the wake kernel.


\entrylabel{Dated literature check}

Focused searches on 2026-09-07 for ARW density,\allowbreak{} cutoff,\allowbreak{} local limits,\allowbreak{} hyperuniformity and incompressibility;\allowbreak{} four later primary papers read at stated theorem scope. No certificate of present openness.


\subsection*{P0385-A001 | Conservation,\allowbreak{} explicit coupling reduction,\allowbreak{} exact law comparison or spectral calculation}

\entrylabel{Argument or result}

Let K be the wake-and-stabilize kernel on k-particle configurations,\allowbreak{} and delta=\allowbreak{}sup\_\allowbreak{}\{x,\allowbreak{}y\}TV(K(x,\allowbreak{}.),\allowbreak{}K(y,\allowbreak{}.)). If nu is stationary,\allowbreak{} convexity gives TV(mu*K,\allowbreak{}nu)<=\allowbreak{}delta for every input law mu. Consequently TV(mu,\allowbreak{}nu)<=\allowbreak{}TV(mu,\allowbreak{}mu*K)+\allowbreak{}delta. Even a proof that delta tends to zero only settles the target once the free-chain law mu is shown nearly invariant under the additional full wake.


\entrylabel{Second check}

Write nu=\allowbreak{}nu*K and integrate TV(K(x,\allowbreak{}.),\allowbreak{}K(y,\allowbreak{}.)) against mu(dx)nu(dy) for the first inequality;\allowbreak{} apply the triangle inequality for the second. This does not silently replace mu with its one-step image.


\entrylabel{Remaining obstruction}

The approximate wake invariance of the free-chain law and a vanishing Dobrushin coefficient are both unproved here. The reduction identifies the extra distributional obligation.


Verification:\allowbreak{} author checked partial | Independent review:\allowbreak{} none | Full target solved:\allowbreak{} False


\clearpage

\section*{P0386  Uniform loss of initial-state information in dense ARW stabilization}

\entrylabel{Primary source}

Levine and Silvestri:\allowbreak{} Universality conjectures for activated random walk;\allowbreak{} Conjecture 24;\allowbreak{} PDF12-13


\url{https://arxiv.org/pdf/2306.01698}


Source statement date:\allowbreak{} 2023-06-14. Archived access:\allowbreak{} 2026-09-07T05:\allowbreak{}09:\allowbreak{}51Z.


Historical:\allowbreak{} source\_\allowbreak{}says\_\allowbreak{}open | Current:\allowbreak{} unverified | Result:\allowbreak{} unsolved


\entrylabel{Question}

For all active initial configurations on a torus with the same particle count between zeta*L\textasciicircum{}d and L\textasciicircum{}d,\allowbreak{} zeta>zeta\_\allowbreak{}c,\allowbreak{} show their stabilized laws become uniformly close in total variation.


\entrylabel{Scope}

The upper bound on particle count is necessary for stabilization and is made explicit. The survey credits a small-sleep-rate regime;\allowbreak{} it is excluded from any novelty claim.


\entrylabel{Dated literature check}

Focused searches on 2026-09-07 for ARW density,\allowbreak{} cutoff,\allowbreak{} local limits,\allowbreak{} hyperuniformity and incompressibility;\allowbreak{} four later primary papers read at stated theorem scope. No certificate of present openness.


\subsection*{P0386-A001 | Conservation,\allowbreak{} explicit coupling reduction,\allowbreak{} exact law comparison or spectral calculation}

\entrylabel{Argument or result}

A particlewise coupling with k labelled particles would prove the result if every pair can be matched to share its future path and sleep instructions before it becomes permanently inactive,\allowbreak{} with individual failure probability o(1/\allowbreak{}k). Indeed,\allowbreak{} the union bound gives failure probability at most k*epsilon\_\allowbreak{}k,\allowbreak{} and successful matching yields identical unlabelled terminal configurations. Pairwise success with probability merely 1-o(1) is too weak when k is of volume order.


\entrylabel{Second check}

For independent failures of probability 1/\allowbreak{}k,\allowbreak{} the chance of at least one failure tends to 1-e\textasciicircum{}-1,\allowbreak{} showing the exact scale needed by this route. No independence is required for the sufficient union bound.


\entrylabel{Remaining obstruction}

Long global stabilization time does not imply every labelled particle remains available for matching. Constructing a valid interacting coupling with uniform o(1/\allowbreak{}k) failure remains the missing step.


Verification:\allowbreak{} author checked partial | Independent review:\allowbreak{} none | Full target solved:\allowbreak{} False


\clearpage

\section*{P0387  Slow-to-one-step mixing transition of the ARW wake chain}

\entrylabel{Primary source}

Levine and Silvestri:\allowbreak{} Universality conjectures for activated random walk;\allowbreak{} Conjecture 25;\allowbreak{} PDF13


\url{https://arxiv.org/pdf/2306.01698}


Source statement date:\allowbreak{} 2023-06-14. Archived access:\allowbreak{} 2026-09-07T05:\allowbreak{}09:\allowbreak{}51Z.


Historical:\allowbreak{} source\_\allowbreak{}says\_\allowbreak{}open | Current:\allowbreak{} unverified | Result:\allowbreak{} unsolved


\entrylabel{Question}

Determine the proposed diffusive lower bound below critical density and one-step asymptotic mixing above it for the wake chain on tori.


\entrylabel{Scope}

The empty configuration is a trivial exception to the literal lower bound. The substantive question concerns positive fixed density;\allowbreak{} mixing tolerance and worst-start versus individual-start quantifiers need specification.


\entrylabel{Dated literature check}

Focused searches on 2026-09-07 for ARW density,\allowbreak{} cutoff,\allowbreak{} local limits,\allowbreak{} hyperuniformity and incompressibility;\allowbreak{} four later primary papers read at stated theorem scope. No certificate of present openness.


\subsection*{P0387-A001 | Conservation,\allowbreak{} explicit coupling reduction,\allowbreak{} exact law comparison or spectral calculation}

\entrylabel{Argument or result}

The exactly solvable one-particle sector provides a dilute check. On the cycle of length L,\allowbreak{} waking a particle runs rate-one simple random walk until an independent Exp(lambda) sleep time. For $\theta=2\pi/L$,\allowbreak{} its first Fourier multiplier is $r_L=\lambda/(\lambda+1-\cos\theta)$. After k wakes,\allowbreak{} a point start has cosine expectation $r_L^k$,\allowbreak{} whereas stationarity is uniform with expectation zero. Hence TV is at least $r_L^k/2$,\allowbreak{} and reaching TV<=\allowbreak{}1/\allowbreak{}4 requires $k\ge\log2/(-\log r_L)\sim\lambda\log2\,L^2/(2\pi^2)$.


\entrylabel{Second check}

The multiplier is the integral of lambda*exp(-lambda*t)*exp(-(1-cos(theta))*t) over positive t. Equivalently it is the resolvent eigenvalue of the walk generator. For zero particles the state space is a singleton and mixing takes zero steps.


\entrylabel{Remaining obstruction}

One particle has density tending to zero,\allowbreak{} so this does not prove the positive-fixed-density lower bound or the supercritical upper bound. Neither is claimed solved.


Verification:\allowbreak{} author checked partial | Independent review:\allowbreak{} none | Full target solved:\allowbreak{} False


\clearpage

\section*{P0388  Hyperuniformity of the infinite-volume wired ARW state}

\entrylabel{Primary source}

Levine and Silvestri:\allowbreak{} Universality conjectures for activated random walk;\allowbreak{} Conjecture 26;\allowbreak{} PDF13-14


\url{https://arxiv.org/pdf/2306.01698}


Source statement date:\allowbreak{} 2023-06-14. Archived access:\allowbreak{} 2026-09-07T05:\allowbreak{}09:\allowbreak{}51Z.


Historical:\allowbreak{} source\_\allowbreak{}says\_\allowbreak{}open | Current:\allowbreak{} unverified | Result:\allowbreak{} unsolved


\entrylabel{Question}

Prove a power saving below volume order for sleeper-number variance in large cubes under the wired infinite-volume measure.


\entrylabel{Scope}

Existence and translation invariance of the measure remain prerequisites. The finite complete-graph result is a different model.


\entrylabel{Dated literature check}

Focused searches on 2026-09-07 for ARW density,\allowbreak{} cutoff,\allowbreak{} local limits,\allowbreak{} hyperuniformity and incompressibility;\allowbreak{} four later primary papers read at stated theorem scope. No certificate of present openness.


\subsection*{P0388-A001 | Conservation,\allowbreak{} explicit coupling reduction,\allowbreak{} exact law comparison or spectral calculation}

\entrylabel{Argument or result}

For a stationary occupation field with covariance c(z),\allowbreak{} $\mathrm{Var}(S_B)=\sum_z c(z)|B\cap(B-z)|$. If $\sum_z c(z)=0$ and $\sum_z |z|_1|c(z)|<\infty$,\allowbreak{} a cube of side L satisfies $\mathrm{Var}(S_B)\le C_d L^{d-1}\sum_z|z|_1|c(z)|$. Subtract |B| times the zero sum,\allowbreak{} then bound the lost overlap by $C_d L^{d-1}\min(|z|_1,L)$. This would give the desired exponent d-1,\allowbreak{} including bounded variance when d=\allowbreak{}1.


\entrylabel{Second check}

Expand the variance as a double sum of covariances and count ordered pairs with displacement z. For a cube the overlap equals the product of (L-|z\_\allowbreak{}i|)\_\allowbreak{}+\allowbreak{},\allowbreak{} verifying the geometric bound directly.


\entrylabel{Remaining obstruction}

Neither the infinite-volume covariance cancellation nor the weighted absolute summability is proved for ARW. Finite-torus conservation gives a zero sum at each size but does not justify interchanging that sum with the infinite-volume limit.


Verification:\allowbreak{} author checked partial | Independent review:\allowbreak{} none | Full target solved:\allowbreak{} False


\clearpage

\section*{P0389  Onset of mesoscopic hyperuniformity in the ARW free chain}

\entrylabel{Primary source}

Levine and Silvestri:\allowbreak{} Universality conjectures for activated random walk;\allowbreak{} Conjecture 27;\allowbreak{} PDF15


\url{https://arxiv.org/pdf/2306.01698}


Source statement date:\allowbreak{} 2023-06-14. Archived access:\allowbreak{} 2026-09-07T05:\allowbreak{}09:\allowbreak{}51Z.


Historical:\allowbreak{} source\_\allowbreak{}says\_\allowbreak{}open | Current:\allowbreak{} unverified | Result:\allowbreak{} candidate\_\allowbreak{}partial\_\allowbreak{}counterexample


\entrylabel{Question}

Determine the subcritical linear-variance and critical/\allowbreak{}supercritical subvolume-variance regimes for particle counts in windows of the free-chain torus configuration.


\entrylabel{Scope}

The literal source permits every box with size-uniform constants. That lower bound fails for boxes filling almost the whole torus. The intended mesoscopic or bounded-volume-fraction regime remains open.


\entrylabel{Dated literature check}

Focused searches on 2026-09-07 for ARW density,\allowbreak{} cutoff,\allowbreak{} local limits,\allowbreak{} hyperuniformity and incompressibility;\allowbreak{} four later primary papers read at stated theorem scope. No certificate of present openness.


\subsection*{P0389-A001 | Conservation,\allowbreak{} explicit coupling reduction,\allowbreak{} exact law comparison or spectral calculation}

\entrylabel{Argument or result}

At step k the free chain has exactly k sleeping particles on N=\allowbreak{}L\textasciicircum{}d sites. Thus Var(S\_\allowbreak{}B)=\allowbreak{}Var(S\_\allowbreak{}Bc) and Var(S\_\allowbreak{}B)<=\allowbreak{}min(|B|,\allowbreak{}N-|B|)\textasciicircum{}2/\allowbreak{}4. Taking the proper box B=\allowbreak{}[0,\allowbreak{}L-2] times the full other d-1 coordinate intervals leaves only L\textasciicircum{}(d-1) sites outside B. In d=\allowbreak{}1,\allowbreak{} B has L-1 sites and Var(S\_\allowbreak{}B)<=\allowbreak{}1/\allowbreak{}4,\allowbreak{} contradicting a size-uniform lower bound c*(L-1) at any fixed positive subcritical density. The contradiction does not depend on a simulation or the full-torus endpoint.


\entrylabel{Second check}

On a cycle,\allowbreak{} S\_\allowbreak{}B=\allowbreak{}k-I\_\allowbreak{}x for the omitted site x. Therefore Var(S\_\allowbreak{}B)=\allowbreak{}Var(I\_\allowbreak{}x)<=\allowbreak{}1/\allowbreak{}4 exactly. Translation invariance from uniform driving even gives Var(I\_\allowbreak{}x)=\allowbreak{}(k/\allowbreak{}L)*(1-k/\allowbreak{}L).


\entrylabel{Remaining obstruction}

This is a scope counterexample to the unrestricted window quantifier. It does not settle variance growth for windows small compared with the torus,\allowbreak{} nor the critical exponent.


Verification:\allowbreak{} author checked partial | Independent review:\allowbreak{} none | Full target solved:\allowbreak{} False


\clearpage

\section*{P0390  Hyperuniformity transition in stationary ARW wake states}

\entrylabel{Primary source}

Levine and Silvestri:\allowbreak{} Universality conjectures for activated random walk;\allowbreak{} Question28;\allowbreak{} PDF15-16


\url{https://arxiv.org/pdf/2306.01698}


Source statement date:\allowbreak{} 2023-06-14. Archived access:\allowbreak{} 2026-09-07T05:\allowbreak{}09:\allowbreak{}51Z.


Historical:\allowbreak{} source\_\allowbreak{}says\_\allowbreak{}open | Current:\allowbreak{} unverified | Result:\allowbreak{} candidate\_\allowbreak{}partial\_\allowbreak{}counterexample


\entrylabel{Question}

Determine whether stationary wake-chain configurations have linear window variance below criticality and a power saving at and above criticality.


\entrylabel{Scope}

As with the free chain,\allowbreak{} a fixed global particle count requires a restriction on observation windows for a positive uniform lower bound. The stationary and transient free ensembles are distinct source questions.


\entrylabel{Dated literature check}

Focused searches on 2026-09-07 for ARW density,\allowbreak{} cutoff,\allowbreak{} local limits,\allowbreak{} hyperuniformity and incompressibility;\allowbreak{} four later primary papers read at stated theorem scope. No certificate of present openness.


\subsection*{P0390-A001 | Conservation,\allowbreak{} explicit coupling reduction,\allowbreak{} exact law comparison or spectral calculation}

\entrylabel{Argument or result}

Every wake-chain state has the same total k. On a cycle,\allowbreak{} the complement-of-one-site window therefore has variance rho\_\allowbreak{}N*(1-rho\_\allowbreak{}N)<=\allowbreak{}1/\allowbreak{}4 under the unique stationary law,\allowbreak{} where rho\_\allowbreak{}N=\allowbreak{}k/\allowbreak{}N;\allowbreak{} the whole torus has variance zero. More generally,\allowbreak{} translation invariance forces sum\_\allowbreak{}y Cov(I\_\allowbreak{}x,\allowbreak{}I\_\allowbreak{}y)=\allowbreak{}Cov(I\_\allowbreak{}x,\allowbreak{}k)=\allowbreak{}0,\allowbreak{} but this finite-volume cancellation does not imply mesoscopic hyperuniformity. For comparison,\allowbreak{} uniform k-subsets have Var(S\_\allowbreak{}B)=\allowbreak{}|B|*rho\_\allowbreak{}N*(1-rho\_\allowbreak{}N)*(N-|B|)/\allowbreak{}(N-1),\allowbreak{} which is volume order for |B|=\allowbreak{}o(N).


\entrylabel{Second check}

Under the uniform k-subset comparison law,\allowbreak{} E I\_\allowbreak{}x=\allowbreak{}k/\allowbreak{}N and E I\_\allowbreak{}x I\_\allowbreak{}y=\allowbreak{}k*(k-1)/\allowbreak{}(N*(N-1));\allowbreak{} substitution yields the displayed variance. This comparison law is not asserted to be the ARW stationary distribution.


\entrylabel{Remaining obstruction}

The literal all-window lower bound is false;\allowbreak{} the meaningful small-window transition is unresolved. Exact global covariance cancellation alone cannot prove a subvolume variance bound.


Verification:\allowbreak{} author checked partial | Independent review:\allowbreak{} none | Full target solved:\allowbreak{} False


\clearpage

\section*{P0391  Diffusion limit for an arbitrary symmetric M/\allowbreak{}G/\allowbreak{}1 discipline}

\entrylabel{Primary source}

Zwart:\allowbreak{} Conjectures on Symmetric Queues in Heavy Traffic;\allowbreak{} Problem 1;\allowbreak{} PDF1-2


\url{https://arxiv.org/pdf/2202.02366}


Source statement date:\allowbreak{} 2022-02-08. Archived access:\allowbreak{} 2026-09-07T02:\allowbreak{}50:\allowbreak{}54Z.


Historical:\allowbreak{} source\_\allowbreak{}says\_\allowbreak{}open | Current:\allowbreak{} unverified | Result:\allowbreak{} unsolved


\entrylabel{Question}

For every work-conserving symmetric discipline gamma,\allowbreak{} prove convergence of Q\_\allowbreak{}gamma(r\textasciicircum{}2 t)/\allowbreak{}r to the common reflected Brownian limit when the service second moment is finite and r(1-rho\_\allowbreak{}r) tends to beta.


\entrylabel{Scope}

Use empty initial state,\allowbreak{} as in the source discussion of universal marginal laws. PS and LCFS are known special cases,\allowbreak{} not new solutions.


\entrylabel{Dated literature check}

Focused 2026-09-07 search for symmetric limits,\allowbreak{} correlated redundancy,\allowbreak{} slowdown and noisy scheduling;\allowbreak{} matched later slowdown theorem and noisy-SJF limits. No current-openness certificate.


\subsection*{P0391-A001 | Transform asymptotics,\allowbreak{} exact special case,\allowbreak{} coupling reduction or identifiability check}

\entrylabel{Argument or result}

Put m=\allowbreak{}E B,\allowbreak{} m2=\allowbreak{}E B\textasciicircum{}2 and h=\allowbreak{}m2/\allowbreak{}(2m). The workload diffusion has drift -beta and variance rate m2/\allowbreak{}m. If uniform-on-compacts state-space collapse Q/\allowbreak{}r-W/\allowbreak{}(rh) were established,\allowbreak{} its queue-length image would have drift -beta/\allowbreak{}h and variance rate m2/\allowbreak{}(m*h\textasciicircum{}2). For beta>0 the stationary exponential rate of this candidate is 2*(beta/\allowbreak{}h)/\allowbreak{}(m2/\allowbreak{}(m*h\textasciicircum{}2))=\allowbreak{}beta,\allowbreak{} exactly the rate obtained from P(Q>=\allowbreak{}k)=\allowbreak{}rho\_\allowbreak{}r\textasciicircum{}k. This fixes the constants and reduces a proof to a precise collapse estimate.


\entrylabel{Second check}

For exponential service of rate mu,\allowbreak{} h=\allowbreak{}1/\allowbreak{}mu and the proposed variance rate is 2*mu,\allowbreak{} matching the birth-death generator directly. The stationary calculation is a consistency test,\allowbreak{} not an interchange-of-limits proof.


\entrylabel{Remaining obstruction}

Uniform state-space collapse for arbitrary gamma remains unproved. Equality of fixed-time marginals does not establish process convergence or tightness.


Verification:\allowbreak{} author checked partial | Independent review:\allowbreak{} none | Full target solved:\allowbreak{} False


\clearpage

\section*{P0392  Infinite-variance scaling limits of symmetric queue lengths}

\entrylabel{Primary source}

Zwart:\allowbreak{} Conjectures on Symmetric Queues in Heavy Traffic;\allowbreak{} Problem 2;\allowbreak{} PDF2-3


\url{https://arxiv.org/pdf/2202.02366}


Source statement date:\allowbreak{} 2022-02-08. Archived access:\allowbreak{} 2026-09-07T02:\allowbreak{}50:\allowbreak{}54Z.


Historical:\allowbreak{} source\_\allowbreak{}says\_\allowbreak{}open | Current:\allowbreak{} unverified | Result:\allowbreak{} unsolved


\entrylabel{Question}

For a regularly varying service tail L(x)x\textasciicircum{}(-alpha),\allowbreak{} 1<alpha<2,\allowbreak{} and c\_\allowbreak{}r*Fbar(c\_\allowbreak{}r)=\allowbreak{}1/\allowbreak{}r,\allowbreak{} identify and prove the process limit of Q\_\allowbreak{}gamma(c\_\allowbreak{}r*t)/\allowbreak{}r for general symmetric gamma.


\entrylabel{Scope}

The LCFS limit is known;\allowbreak{} the source identifies PS tightness and the general-discipline limit as unresolved. The following stationary calculation assumes beta>0.


\entrylabel{Dated literature check}

Focused 2026-09-07 search for symmetric limits,\allowbreak{} correlated redundancy,\allowbreak{} slowdown and noisy scheduling;\allowbreak{} matched later slowdown theorem and noisy-SJF limits. No current-openness certificate.


\subsection*{P0392-A001 | Transform asymptotics,\allowbreak{} exact special case,\allowbreak{} coupling reduction or identifiability check}

\entrylabel{Argument or result}

The Pollaczek-Khinchine transform is E exp(-sW)=\allowbreak{}(1-rho)*s/\allowbreak{}[s-lambda*(1-E exp(-sB))]. Regular variation gives E exp(-sB)=\allowbreak{}1-ms+\allowbreak{}a*s\textasciicircum{}alpha*L(1/\allowbreak{}s)+\allowbreak{}o(s\textasciicircum{}alpha L(1/\allowbreak{}s)),\allowbreak{} where a=\allowbreak{}-Gamma(1-alpha)>0. Substituting s=\allowbreak{}u/\allowbreak{}c\_\allowbreak{}r and r(1-rho\_\allowbreak{}r)->beta yields E exp(-uW/\allowbreak{}c\_\allowbreak{}r)->beta/\allowbreak{}[beta+\allowbreak{}(a/\allowbreak{}m)*u\textasciicircum{}(alpha-1)]. In contrast Q/\allowbreak{}r has stationary exponential limit of rate beta. Consequently a positive linear state-space collapse W/\allowbreak{}c\_\allowbreak{}r=\allowbreak{}c*Q/\allowbreak{}r cannot hold in this regime.


\entrylabel{Second check}

The normalization c\_\allowbreak{}r\textasciicircum{}(1-alpha)*L(c\_\allowbreak{}r)=\allowbreak{}1/\allowbreak{}r makes the transform coefficient a/\allowbreak{}(m*beta). The limiting workload transform has infinite derivative at zero,\allowbreak{} whereas every scaled exponential has finite mean;\allowbreak{} the two laws cannot coincide.


\entrylabel{Remaining obstruction}

This rules out only linear collapse and supplies a necessary stationary check. It neither identifies the temporal law nor proves tightness;\allowbreak{} even a nonlinear relation is not excluded by this marginal calculation.


Verification:\allowbreak{} author checked partial | Independent review:\allowbreak{} none | Full target solved:\allowbreak{} False


\clearpage

\section*{P0393  Uniform busy-cycle maximum bounds for symmetric queues}

\entrylabel{Primary source}

Zwart:\allowbreak{} Conjectures on Symmetric Queues in Heavy Traffic;\allowbreak{} Discussion,\allowbreak{} final paragraph before references;\allowbreak{} PDF3


\url{https://arxiv.org/pdf/2202.02366}


Source statement date:\allowbreak{} 2022-02-08. Archived access:\allowbreak{} 2026-09-07T02:\allowbreak{}50:\allowbreak{}54Z.


Historical:\allowbreak{} source\_\allowbreak{}says\_\allowbreak{}open | Current:\allowbreak{} unverified | Result:\allowbreak{} unsolved


\entrylabel{Question}

Obtain busy-cycle queue-maximum tail bounds for arbitrary symmetric disciplines strong enough to support the regenerative tightness route described in the source.


\entrylabel{Scope}

This is the separately identified excursion-tail question. It is not counted again for each service distribution or traffic parameter.


\entrylabel{Dated literature check}

Focused 2026-09-07 search for symmetric limits,\allowbreak{} correlated redundancy,\allowbreak{} slowdown and noisy scheduling;\allowbreak{} matched later slowdown theorem and noisy-SJF limits. No current-openness certificate.


\subsection*{P0393-A001 | Transform asymptotics,\allowbreak{} exact special case,\allowbreak{} coupling reduction or identifiability check}

\entrylabel{Argument or result}

With exponential service of rate mu,\allowbreak{} the queue length under any symmetric gamma is a birth-death chain:\allowbreak{} the total departure hazard is mu times sum\_\allowbreak{}i gamma(n,\allowbreak{}i)=\allowbreak{}mu. For a busy period starting at one,\allowbreak{} the probability that its maximum reaches K is (q-1)/\allowbreak{}(q\textasciicircum{}K-1),\allowbreak{} q=\allowbreak{}mu/\allowbreak{}lambda,\allowbreak{} with critical value 1/\allowbreak{}K. If r(1-rho\_\allowbreak{}r)->beta>0 and K=\allowbreak{}floor(x*r),\allowbreak{} r times this probability tends to beta/\allowbreak{}(exp(beta*x)-1). Thus a usable general tail estimate must be compatible with this exact critical excursion scale.


\entrylabel{Second check}

Solve p\_\allowbreak{}j=\allowbreak{}(lambda*p\_\allowbreak{}(j+\allowbreak{}1)+\allowbreak{}mu*p\_\allowbreak{}(j-1))/\allowbreak{}(lambda+\allowbreak{}mu),\allowbreak{} p\_\allowbreak{}0=\allowbreak{}0,\allowbreak{}p\_\allowbreak{}K=\allowbreak{}1,\allowbreak{} and evaluate p\_\allowbreak{}1. Independently,\allowbreak{} the stopped martingale q\textasciicircum{}Q gives the same expression.


\entrylabel{Remaining obstruction}

General service distributions destroy this one-dimensional Markov calculation. A uniform bound must control residual-service composition within an excursion;\allowbreak{} fixed-time insensitivity does not determine the maximum.


Verification:\allowbreak{} author checked partial | Independent review:\allowbreak{} none | Full target solved:\allowbreak{} False


\clearpage

\section*{P0394  Efficient input estimation from periodically observed queue workloads}

\entrylabel{Primary source}

Ravner:\allowbreak{} Queue input estimation from discrete workload observations;\allowbreak{} Problem statement and final discussion;\allowbreak{} PDF2-3


\url{https://arxiv.org/pdf/2202.10750}


Source statement date:\allowbreak{} 2022-02-23. Archived access:\allowbreak{} 2026-09-07T02:\allowbreak{}50:\allowbreak{}55Z.


Historical:\allowbreak{} source\_\allowbreak{}says\_\allowbreak{}open | Current:\allowbreak{} unverified | Result:\allowbreak{} unsolved


\entrylabel{Question}

Construct statistically efficient estimators of the finite-dimensional compound-Poisson input parameters from discrete M/\allowbreak{}G/\allowbreak{}1 workload observations,\allowbreak{} retaining their temporal dependence.


\entrylabel{Scope}

Consistency and transform-based methods already exist. The target is attainable asymptotic efficiency,\allowbreak{} not simply another consistent moment estimator.


\entrylabel{Dated literature check}

Focused 2026-09-07 search for symmetric limits,\allowbreak{} correlated redundancy,\allowbreak{} slowdown and noisy scheduling;\allowbreak{} matched later slowdown theorem and noisy-SJF limits. No current-openness certificate.


\subsection*{P0394-A001 | Transform asymptotics,\allowbreak{} exact special case,\allowbreak{} coupling reduction or identifiability check}

\entrylabel{Argument or result}

For a stable M/\allowbreak{}M/\allowbreak{}1 benchmark let a=\allowbreak{}P(W>0)=\allowbreak{}lambda/\allowbreak{}mu and b=\allowbreak{}E W=\allowbreak{}a/\allowbreak{}(mu-lambda). Inversion gives mu=\allowbreak{}a/\allowbreak{}[b*(1-a)] and lambda=\allowbreak{}a\textasciicircum{}2/\allowbreak{}[b*(1-a)]. Empirical busy fraction and mean workload therefore yield a consistent estimator on a fixed sampling grid by ergodicity. Its asymptotic covariance,\allowbreak{} when the joint CLT holds,\allowbreak{} is J*Sigma*J\textasciicircum{}T,\allowbreak{} where J is the Jacobian of this inversion and Sigma=\allowbreak{}sum over integer lags of Cov((1\{W\_\allowbreak{}0>0\},\allowbreak{}W\_\allowbreak{}0),\allowbreak{}(1\{W\_\allowbreak{}l>0\},\allowbreak{}W\_\allowbreak{}l)). The off-zero lags are part of the information calculation.


\entrylabel{Second check}

Substitution of the inverted parameters returns both a and b exactly. At lambda=\allowbreak{}1,\allowbreak{}mu=\allowbreak{}2 one gets a=\allowbreak{}1/\allowbreak{}2,\allowbreak{}b=\allowbreak{}1/\allowbreak{}2 and recovers both rates. Treating Sigma as a one-observation variance would discard the dependence at every finite grid spacing.


\entrylabel{Remaining obstruction}

The estimator is a benchmark only. Neither an efficiency bound for general input families nor an estimator attaining it has been proved here.


Verification:\allowbreak{} author checked partial | Independent review:\allowbreak{} none | Full target solved:\allowbreak{} False


\clearpage

\section*{P0395  Network decomposition with several bottlenecks and internal flow covariance}

\entrylabel{Primary source}

Whitt-You:\allowbreak{} New decomposition approximations for queueing networks;\allowbreak{} Sections2.1.1-2.1.2;\allowbreak{} PDF2


\url{https://www.columbia.edu/~ww2040/view_RQNA_format_020222.pdf}


Source statement date:\allowbreak{} 2022-02-02. Archived access:\allowbreak{} 2026-09-07T02:\allowbreak{}50:\allowbreak{}58Z.


Historical:\allowbreak{} source\_\allowbreak{}says\_\allowbreak{}open | Current:\allowbreak{} unverified | Result:\allowbreak{} unsolved


\entrylabel{Question}

Extend the robust queueing network analyzer to several simultaneously critical nodes while estimating and propagating the internal cross-flow covariance needed in non-tree networks.


\entrylabel{Scope}

The multiple-bottleneck and covariance-fitting subquestions are grouped because the proposed limit calculation needs both.


\entrylabel{Dated literature check}

Focused 2026-09-07 search for symmetric limits,\allowbreak{} correlated redundancy,\allowbreak{} slowdown and noisy scheduling;\allowbreak{} matched later slowdown theorem and noisy-SJF limits. No current-openness certificate.


\subsection*{P0395-A001 | Transform asymptotics,\allowbreak{} exact special case,\allowbreak{} coupling reduction or identifiability check}

\entrylabel{Argument or result}

Separate rate and IDC functions cannot determine the required joint input law. Let N1,\allowbreak{}N2 each have rate lambda Poisson marginals. If independent,\allowbreak{} Var(N1(t)+\allowbreak{}N2(t))=\allowbreak{}2*lambda*t;\allowbreak{} if N1=\allowbreak{}N2,\allowbreak{} the variance is 4*lambda*t. Both inputs have identical individual IDCs equal to one,\allowbreak{} but the total IDC is respectively one and two. Thus any multi-bottleneck diffusion approximation based solely on the marginal IDC data has a missing covariance parameter before reflection is even applied.


\entrylabel{Second check}

Expand Var(A+\allowbreak{}B)=\allowbreak{}Var A+\allowbreak{}Var B+\allowbreak{}2 Cov(A,\allowbreak{}B). The independent and common-input examples make the cross term respectively zero and lambda*t. They are auxiliary arrival models,\allowbreak{} not claimed counterexamples to an analyzer that retains cross covariances.


\entrylabel{Remaining obstruction}

Need a consistent joint covariance model and its propagation through multidimensional reflected Brownian flow limits. This information obstruction does not provide those estimates or an approximation error bound.


Verification:\allowbreak{} author checked partial | Independent review:\allowbreak{} none | Full target solved:\allowbreak{} False


\clearpage

\section*{P0396  IDC-based decomposition for networks of multiserver queues}

\entrylabel{Primary source}

Whitt-You:\allowbreak{} New decomposition approximations for queueing networks;\allowbreak{} Section2.2(i),\allowbreak{} Section3.2(i);\allowbreak{} PDF2-3


\url{https://www.columbia.edu/~ww2040/view_RQNA_format_020222.pdf}


Source statement date:\allowbreak{} 2022-02-02. Archived access:\allowbreak{} 2026-09-07T02:\allowbreak{}50:\allowbreak{}58Z.


Historical:\allowbreak{} source\_\allowbreak{}says\_\allowbreak{}open | Current:\allowbreak{} unverified | Result:\allowbreak{} unsolved


\entrylabel{Question}

Extend the single-server IDC-based robust queueing analysis to multiple servers at each network node,\allowbreak{} with a justified performance approximation.


\entrylabel{Scope}

The source distinguishes this from the existing multiserver QNA and generic robust-queueing models.


\entrylabel{Dated literature check}

Focused 2026-09-07 search for symmetric limits,\allowbreak{} correlated redundancy,\allowbreak{} slowdown and noisy scheduling;\allowbreak{} matched later slowdown theorem and noisy-SJF limits. No current-openness certificate.


\subsection*{P0396-A001 | Transform asymptotics,\allowbreak{} exact special case,\allowbreak{} coupling reduction or identifiability check}

\entrylabel{Argument or result}

Even without future arrivals,\allowbreak{} aggregate workload cannot close a two-server state description. Residual workloads (2,\allowbreak{}0) and (1,\allowbreak{}1) both sum to two,\allowbreak{} but their immediate workload drain rates are one and two. Their times to empty are two and one. A scalar workload approximation therefore needs an additional busy-server or residual-work descriptor,\allowbreak{} or a proved collapse regime that removes this ambiguity.


\entrylabel{Second check}

At time t in [0,\allowbreak{}1],\allowbreak{} the two total workloads are 2-t and 2-2t. This exact deterministic calculation isolates a state-information gap before any stochastic approximation.


\entrylabel{Remaining obstruction}

Need an augmented state and a tractable covariance propagation rule,\allowbreak{} together with an asymptotic justification. A scalar workload replacement alone cannot be universally exact.


Verification:\allowbreak{} author checked partial | Independent review:\allowbreak{} none | Full target solved:\allowbreak{} False


\clearpage

\section*{P0397  Multiclass and non-Markovian-routing robust queueing decomposition}

\entrylabel{Primary source}

Whitt-You:\allowbreak{} New decomposition approximations for queueing networks;\allowbreak{} Section2.2(ii),\allowbreak{} Section3.2(ii);\allowbreak{} PDF2-3


\url{https://www.columbia.edu/~ww2040/view_RQNA_format_020222.pdf}


Source statement date:\allowbreak{} 2022-02-02. Archived access:\allowbreak{} 2026-09-07T02:\allowbreak{}50:\allowbreak{}58Z.


Historical:\allowbreak{} source\_\allowbreak{}says\_\allowbreak{}open | Current:\allowbreak{} unverified | Result:\allowbreak{} unsolved


\entrylabel{Question}

Extend the IDC network approximation to class-dependent or more general routing without losing the dependencies that determine downstream congestion.


\entrylabel{Scope}

The source describes aggregation into approximate Markovian routing as an existing method,\allowbreak{} not an exact solution.


\entrylabel{Dated literature check}

Focused 2026-09-07 search for symmetric limits,\allowbreak{} correlated redundancy,\allowbreak{} slowdown and noisy scheduling;\allowbreak{} matched later slowdown theorem and noisy-SJF limits. No current-openness certificate.


\subsection*{P0397-A001 | Transform asymptotics,\allowbreak{} exact special case,\allowbreak{} coupling reduction or identifiability check}

\entrylabel{Argument or result}

Consider iid job marks taking two values with equal probability. A job has first-node service requirement B=\allowbreak{}1 or B=\allowbreak{}3. In model A only the B=\allowbreak{}3 jobs are routed downstream;\allowbreak{} in model B only B=\allowbreak{}1 jobs are routed downstream. Both models have the same marginal service distribution and the same routing probability 1/\allowbreak{}2. But E[B*R] is 3/\allowbreak{}2 in A and 1/\allowbreak{}2 in B,\allowbreak{} and Cov(B,\allowbreak{}R) is respectively 1/\allowbreak{}2 and -1/\allowbreak{}2. Marginal service and routing data thus do not identify the marked flow dependence entering a downstream approximation.


\entrylabel{Second check}

E B=\allowbreak{}2 and E R=\allowbreak{}1/\allowbreak{}2 in both examples,\allowbreak{} so subtracting their product one gives the two covariance signs. The examples use legal iid per-job marks;\allowbreak{} no independence between B and R was assumed.


\entrylabel{Remaining obstruction}

Need a marked or class-resolved closure and bounds for its error. The scalar marginal primitives alone cannot recover this cross moment.


Verification:\allowbreak{} author checked partial | Independent review:\allowbreak{} none | Full target solved:\allowbreak{} False


\clearpage

\section*{P0398  Robust network approximations under time-varying arrivals}

\entrylabel{Primary source}

Whitt-You:\allowbreak{} New decomposition approximations for queueing networks;\allowbreak{} Section2.2(iii),\allowbreak{} Section3.2(iii);\allowbreak{} PDF2-3


\url{https://www.columbia.edu/~ww2040/view_RQNA_format_020222.pdf}


Source statement date:\allowbreak{} 2022-02-02. Archived access:\allowbreak{} 2026-09-07T02:\allowbreak{}50:\allowbreak{}58Z.


Historical:\allowbreak{} source\_\allowbreak{}says\_\allowbreak{}open | Current:\allowbreak{} unverified | Result:\allowbreak{} unsolved


\entrylabel{Question}

Develop and justify a time-varying extension of the robust queueing network analyzer,\allowbreak{} accounting for propagation delays and changing congestion.


\entrylabel{Scope}

Single-queue time-varying robust approximations are already cited in the source. Network behavior is the target.


\entrylabel{Dated literature check}

Focused 2026-09-07 search for symmetric limits,\allowbreak{} correlated redundancy,\allowbreak{} slowdown and noisy scheduling;\allowbreak{} matched later slowdown theorem and noisy-SJF limits. No current-openness certificate.


\subsection*{P0398-A001 | Transform asymptotics,\allowbreak{} exact special case,\allowbreak{} coupling reduction or identifiability check}

\entrylabel{Argument or result}

An exact infinite-server benchmark shows the necessary temporal response. With exponential service rate mu and arrival intensity lambda(t)=\allowbreak{}lambda0+\allowbreak{}epsilon*cos(omega*t),\allowbreak{} the periodic mean occupancy is lambda0/\allowbreak{}mu+\allowbreak{}epsilon*(mu*cos(omega*t)+\allowbreak{}omega*sin(omega*t))/\allowbreak{}(mu\textasciicircum{}2+\allowbreak{}omega\textasciicircum{}2). It solves mprime=\allowbreak{}lambda(t)-mu*m. The instantaneous stationary substitution lambda(t)/\allowbreak{}mu has neither the correct amplitude nor the phase delay when omega is comparable to mu.


\entrylabel{Second check}

Differentiate the displayed expression and substitute into the mean equation. Equivalently integrate the survival kernel m(t)=\allowbreak{}integral from minus infinity to t of lambda(s)*exp(-mu*(t-s)) ds. These independently give the same amplitude epsilon/\allowbreak{}sqrt(mu\textasciicircum{}2+\allowbreak{}omega\textasciicircum{}2).


\entrylabel{Remaining obstruction}

A finite-server interacting network requires an additional queueing and routing response,\allowbreak{} plus a controlled approximation regime. The linear benchmark constrains such a model but does not solve it.


Verification:\allowbreak{} author checked partial | Independent review:\allowbreak{} none | Full target solved:\allowbreak{} False


\clearpage

\section*{P0399  Identifiable and efficient probe inference in tandem queueing networks}

\entrylabel{Primary source}

Antunes-Jacinto-Pacheco:\allowbreak{} Statistical inference in queueing networks with probing information;\allowbreak{} Problem statement;\allowbreak{} PDF2


\url{https://rdpc.uevora.pt/bitstream/10174/32450/3/paper_Q65.pdf}


Source statement date:\allowbreak{} 2022-02-25. Archived access:\allowbreak{} 2026-09-07T02:\allowbreak{}51:\allowbreak{}02Z.


Historical:\allowbreak{} source\_\allowbreak{}says\_\allowbreak{}open | Current:\allowbreak{} unverified | Result:\allowbreak{} unsolved


\entrylabel{Question}

Design a minimally perturbing probing strategy and inference method for the unknown arrival streams and packet-size distributions in the source tandem FIFO network.


\entrylabel{Scope}

The full per-stream target needs identifiability assumptions. The source allows n queues without an explicit n>=\allowbreak{}2 restriction;\allowbreak{} the following obstruction concerns n=\allowbreak{}1.


\entrylabel{Dated literature check}

Focused 2026-09-07 search for symmetric limits,\allowbreak{} correlated redundancy,\allowbreak{} slowdown and noisy scheduling;\allowbreak{} matched later slowdown theorem and noisy-SJF limits. No current-openness certificate.


\subsection*{P0399-A001 | Transform asymptotics,\allowbreak{} exact special case,\allowbreak{} coupling reduction or identifiability check}

\entrylabel{Argument or result}

For n=\allowbreak{}1,\allowbreak{} streams 0 and 1 have the same route. If both are independent Poisson processes with the same positive packet-size law G,\allowbreak{} their superposition is a marked Poisson process of rate lambda0+\allowbreak{}lambda1 and mark law G. Every probe departure is a function of that unlabelled input and the chosen probes. Therefore all positive pairs (lambda0,\allowbreak{}lambda1) with a common sum give exactly the same probe-data distribution,\allowbreak{} including under adaptive strategies whose randomization is independent of the hidden stream labels.


\entrylabel{Second check}

Couple the common unlabelled Poisson input and probe randomization,\allowbreak{} then assign stream labels by independent thinning with the appropriate probabilities. All queue and probe paths agree pathwise while the separate rates differ.


\entrylabel{Remaining obstruction}

Separate per-stream rates cannot be recovered in this allowed one-node case. Identifiable multi-node models,\allowbreak{} optimal probe design,\allowbreak{} asymptotic normality and efficiency remain open targets. This is a scope obstruction,\allowbreak{} with no novelty claim.


Verification:\allowbreak{} author checked partial | Independent review:\allowbreak{} none | Full target solved:\allowbreak{} False


\clearpage

\section*{P0400  Response-time laws with correlated redundancy replicas}

\entrylabel{Primary source}

Gardner:\allowbreak{} Correlation in redundancy systems;\allowbreak{} Q1;\allowbreak{} PDF1,\allowbreak{}3


\url{https://kgardner.people.amherst.edu/papers/gardner_questa_100views.pdf}


Source statement date:\allowbreak{} 2022-05-05. Archived access:\allowbreak{} 2026-09-07T02:\allowbreak{}54:\allowbreak{}03Z.


Historical:\allowbreak{} source\_\allowbreak{}says\_\allowbreak{}open | Current:\allowbreak{} unverified | Result:\allowbreak{} unsolved


\entrylabel{Question}

Derive useful exact or justified approximate response-time laws for redundancy-d queues with correlated replica service requirements and specified scheduling rules.


\entrylabel{Scope}

Known product-form results for independent exponential replicas and numerical FCFS methods are credited. Replica starting times need not coincide in a loaded system.


\entrylabel{Dated literature check}

Focused 2026-09-07 search for symmetric limits,\allowbreak{} correlated redundancy,\allowbreak{} slowdown and noisy scheduling;\allowbreak{} matched later slowdown theorem and noisy-SJF limits. No current-openness certificate.


\subsection*{P0400-A001 | Transform asymptotics,\allowbreak{} exact special case,\allowbreak{} coupling reduction or identifiability check}

\entrylabel{Argument or result}

For an isolated job in the S-and-X model,\allowbreak{} all d copies start together and completion time is X*min(S1,\allowbreak{}...,\allowbreak{}Sd),\allowbreak{} where X is independent of the iid slowdowns. Consequently E T=\allowbreak{}E X times integral\_\allowbreak{}0\textasciicircum{}infinity P(S>t)\textasciicircum{}d dt. This gives a required zero-load boundary value for any response approximation. With unit-rate exponential S the value is E X/\allowbreak{}d;\allowbreak{} with deterministic S=\allowbreak{}1 it is E X for every d.


\entrylabel{Second check}

Use the tail-integral identity for a nonnegative minimum,\allowbreak{} or condition on X. The deterministic and exponential checks display why equal one-copy mean service time cannot determine the redundancy gain.


\entrylabel{Remaining obstruction}

The isolated-job result is exact but classical. Positive-load response depends on staggered starts,\allowbreak{} cancellations and inter-queue dependence;\allowbreak{} no formula or error bound for that regime is established.


Verification:\allowbreak{} author checked partial | Independent review:\allowbreak{} none | Full target solved:\allowbreak{} False


\clearpage

\section*{P0401  Optimizing replication degree under correlated service requirements}

\entrylabel{Primary source}

Gardner:\allowbreak{} Correlation in redundancy systems;\allowbreak{} Q2;\allowbreak{} PDF1,\allowbreak{}3


\url{https://kgardner.people.amherst.edu/papers/gardner_questa_100views.pdf}


Source statement date:\allowbreak{} 2022-05-05. Archived access:\allowbreak{} 2026-09-07T02:\allowbreak{}54:\allowbreak{}03Z.


Historical:\allowbreak{} source\_\allowbreak{}says\_\allowbreak{}open | Current:\allowbreak{} unverified | Result:\allowbreak{} unsolved


\entrylabel{Question}

Determine the replication degree d minimizing steady-state mean response time for correlated-replica redundancy queues.


\entrylabel{Scope}

All degrees are one optimization question. Independent-replica monotonicity results do not resolve general correlation.


\entrylabel{Dated literature check}

Focused 2026-09-07 search for symmetric limits,\allowbreak{} correlated redundancy,\allowbreak{} slowdown and noisy scheduling;\allowbreak{} matched later slowdown theorem and noisy-SJF limits. No current-openness certificate.


\subsection*{P0401-A001 | Transform asymptotics,\allowbreak{} exact special case,\allowbreak{} coupling reduction or identifiability check}

\entrylabel{Argument or result}

For a simultaneously started isolated job,\allowbreak{} total server effort is d*X*min\_\allowbreak{}i S\_\allowbreak{}i. Unit-rate iid exponential slowdowns give mean effort E X for every d,\allowbreak{} while identical deterministic slowdowns give d*E X. The corresponding completion times are E X/\allowbreak{}d and E X. Thus the same replica count can buy latency without extra mean work in one model,\allowbreak{} or multiply mean work without reducing latency in another. A degree optimizer must account for both the joint slowdown law and load.


\entrylabel{Second check}

The minimum of d unit-rate exponentials has rate d,\allowbreak{} so multiplying its mean by d cancels exactly. For deterministic slowdowns every copy runs the whole job,\allowbreak{} giving the second formula pathwise.


\entrylabel{Remaining obstruction}

This does not determine the optimum at positive load. Queue-state-dependent staggering can change both effort and latency,\allowbreak{} and an optimization theorem needs a controlled response function.


Verification:\allowbreak{} author checked partial | Independent review:\allowbreak{} none | Full target solved:\allowbreak{} False


\clearpage

\section*{P0402  Stability beyond processor sharing with correlated redundancy copies}

\entrylabel{Primary source}

Gardner:\allowbreak{} Correlation in redundancy systems;\allowbreak{} Q3;\allowbreak{} PDF1,\allowbreak{}3


\url{https://kgardner.people.amherst.edu/papers/gardner_questa_100views.pdf}


Source statement date:\allowbreak{} 2022-05-05. Archived access:\allowbreak{} 2026-09-07T02:\allowbreak{}54:\allowbreak{}03Z.


Historical:\allowbreak{} source\_\allowbreak{}says\_\allowbreak{}open | Current:\allowbreak{} unverified | Result:\allowbreak{} unsolved


\entrylabel{Question}

Characterize the stability region of correlated-replica redundancy-d systems under FCFS and other scheduling disciplines beyond the known PS results.


\entrylabel{Scope}

The following fully replicated endpoint d=\allowbreak{}k is a solvable check outside the survey focus d much smaller than k;\allowbreak{} it does not resolve the general regime.


\entrylabel{Dated literature check}

Focused 2026-09-07 search for symmetric limits,\allowbreak{} correlated redundancy,\allowbreak{} slowdown and noisy scheduling;\allowbreak{} matched later slowdown theorem and noisy-SJF limits. No current-openness certificate.


\subsection*{P0402-A001 | Transform asymptotics,\allowbreak{} exact special case,\allowbreak{} coupling reduction or identifiability check}

\entrylabel{Argument or result}

When d=\allowbreak{}k and every queue is FCFS,\allowbreak{} all queues have the same job order and begin the next job simultaneously after a cancellation. The system is exactly an M/\allowbreak{}G/\allowbreak{}1 queue of rate k*lambda and service duration M=\allowbreak{}min(B1,\allowbreak{}...,\allowbreak{}Bk),\allowbreak{} where the vector B is iid across jobs but may be correlated within each job. Its strict stability condition is k*lambda*E M<1. In the S-and-X case M=\allowbreak{}X*min\_\allowbreak{}i S\_\allowbreak{}i. Identical copies yield lambda<1/\allowbreak{}(k*E B);\allowbreak{} independent unit exponentials yield lambda<1.


\entrylabel{Second check}

Induct over completion epochs:\allowbreak{} all copies of the head job depart together,\allowbreak{} leaving the same ordered job list at every server. This proves the one-server reduction rather than merely comparing average work.


\entrylabel{Remaining obstruction}

For d<k,\allowbreak{} different compatibility sets break synchronization. The exact endpoint cannot justify a stability criterion based only on d times the isolated minimum service duration.


Verification:\allowbreak{} author checked partial | Independent review:\allowbreak{} none | Full target solved:\allowbreak{} False


\clearpage

\section*{P0403  Optimal scheduling for correlated redundancy systems}

\entrylabel{Primary source}

Gardner:\allowbreak{} Correlation in redundancy systems;\allowbreak{} Q4;\allowbreak{} PDF1,\allowbreak{}3


\url{https://kgardner.people.amherst.edu/papers/gardner_questa_100views.pdf}


Source statement date:\allowbreak{} 2022-05-05. Archived access:\allowbreak{} 2026-09-07T02:\allowbreak{}54:\allowbreak{}03Z.


Historical:\allowbreak{} source\_\allowbreak{}says\_\allowbreak{}open | Current:\allowbreak{} unverified | Result:\allowbreak{} unsolved


\entrylabel{Question}

Identify scheduling rules minimizing mean response time,\allowbreak{} including how the optimal rule changes with the replica correlation and service law.


\entrylabel{Scope}

The source credits least-redundant-first optimality in special nested exponential models. The following deterministic comparison is only a finite-workload diagnostic.


\entrylabel{Dated literature check}

Focused 2026-09-07 search for symmetric limits,\allowbreak{} correlated redundancy,\allowbreak{} slowdown and noisy scheduling;\allowbreak{} matched later slowdown theorem and noisy-SJF limits. No current-openness certificate.


\subsection*{P0403-A001 | Transform asymptotics,\allowbreak{} exact special case,\allowbreak{} coupling reduction or identifiability check}

\entrylabel{Argument or result}

Take two servers and two simultaneously present jobs with identical replica durations a<=\allowbreak{}b across servers. Replicating the shorter job on both servers first gives completion times a and a+\allowbreak{}b,\allowbreak{} with mean a+\allowbreak{}b/\allowbreak{}2. Assigning one job to each server immediately gives completion times a and b,\allowbreak{} with mean (a+\allowbreak{}b)/\allowbreak{}2. The second policy improves the mean by a/\allowbreak{}2. Therefore a proof based on always assigning every free server to the currently shortest job cannot be valid across the correlated-copy model class.


\entrylabel{Second check}

List the two completion epochs directly;\allowbreak{} at a=\allowbreak{}b=\allowbreak{}1 the means are 3/\allowbreak{}2 versus 1. This is a two-job scheduling example,\allowbreak{} not a stationary optimality theorem.


\entrylabel{Remaining obstruction}

Need a dynamic rule and a stochastic comparison or control proof for the full arriving system. Finite-workload preference need not identify a steady-state optimizer.


Verification:\allowbreak{} author checked partial | Independent review:\allowbreak{} none | Full target solved:\allowbreak{} False


\clearpage

\section*{P0404  Tractable concurrent-replica policies with broad stability guarantees}

\entrylabel{Primary source}

Gardner:\allowbreak{} Correlation in redundancy systems;\allowbreak{} Q5;\allowbreak{} PDF2-3


\url{https://kgardner.people.amherst.edu/papers/gardner_questa_100views.pdf}


Source statement date:\allowbreak{} 2022-05-05. Archived access:\allowbreak{} 2026-09-07T02:\allowbreak{}54:\allowbreak{}03Z.


Historical:\allowbreak{} source\_\allowbreak{}says\_\allowbreak{}open | Current:\allowbreak{} unverified | Result:\allowbreak{} unsolved


\entrylabel{Question}

Design analytically tractable redundancy policies that retain useful simultaneous service and achieve a provably maximal stability region under a specified correlated-replica model.


\entrylabel{Scope}

RIQ and delta-probe maximal-stability results with only one served replica are already known. Baseline stability is weaker than a maximal-region theorem.


\entrylabel{Dated literature check}

Focused 2026-09-07 search for symmetric limits,\allowbreak{} correlated redundancy,\allowbreak{} slowdown and noisy scheduling;\allowbreak{} matched later slowdown theorem and noisy-SJF limits. No current-openness certificate.


\subsection*{P0404-A001 | Transform asymptotics,\allowbreak{} exact special case,\allowbreak{} coupling reduction or identifiability check}

\entrylabel{Argument or result}

A conservative construction uses independent primary routing and serves extra replicas only on otherwise idle servers,\allowbreak{} with instantaneous preemption of an extra whenever a primary needs that server. Couple each primary to the same job and primary service requirement in the no-replication system. An extra completion can only remove work earlier,\allowbreak{} and an extra never consumes primary capacity,\allowbreak{} so each primary workload is no larger than its baseline counterpart. Under iid uniform primary routing,\allowbreak{} the ordinary M/\allowbreak{}G/\allowbreak{}1 load condition lambda*E B\_\allowbreak{}primary<1 supplies a baseline stability benchmark.


\entrylabel{Second check}

At primary arrivals add the same work in both coupled systems;\allowbreak{} between arrivals the proposed primary workload has at least the baseline service opportunity;\allowbreak{} cancellation subtracts nonnegative work. The comparison requires zero preemption/\allowbreak{}cancellation overhead and no hidden primary slowdown.


\entrylabel{Remaining obstruction}

This comparison does not prove maximality or quantify response gains,\allowbreak{} and is not claimed as a new policy. For general service laws redundancy can alter achievable capacity;\allowbreak{} a universal baseline argument is insufficient.


Verification:\allowbreak{} author checked partial | Independent review:\allowbreak{} none | Full target solved:\allowbreak{} False


\clearpage

\section*{P0405  Characterizing optimality of MAX-H in multiclass queues}

\entrylabel{Primary source}

Aalto and Scully:\allowbreak{} Minimizing the mean slowdown in the M/\allowbreak{}G/\allowbreak{}1 queue;\allowbreak{} Section4 after Theorem3;\allowbreak{} PDF15


\url{https://aaltodoc.aalto.fi/server/api/core/bitstreams/a9d987aa-1806-45f5-a7be-7bce2d8cc8f0/content}


Source statement date:\allowbreak{} 2023-08-17. Archived access:\allowbreak{} 2026-09-07T05:\allowbreak{}39:\allowbreak{}06Z.


Historical:\allowbreak{} source\_\allowbreak{}says\_\allowbreak{}open | Current:\allowbreak{} unverified | Result:\allowbreak{} unsolved


\entrylabel{Question}

Find a necessary and sufficient characterization,\allowbreak{} or a sharper usable sufficient condition,\allowbreak{} for MAX-H optimality in the multiclass M/\allowbreak{}G/\allowbreak{}1 generalized holding-cost problem.


\entrylabel{Scope}

Theorem3 gives a sufficient monotonicity condition only. The concluding prose loosely says necessary and sufficient;\allowbreak{} the theorem and PDF15 explicitly explain that necessity is unresolved. Single-class FCFS is already resolved.


\entrylabel{Dated literature check}

Focused 2026-09-07 search for symmetric limits,\allowbreak{} correlated redundancy,\allowbreak{} slowdown and noisy scheduling;\allowbreak{} matched later slowdown theorem and noisy-SJF limits. No current-openness certificate.


\subsection*{P0405-A001 | Transform asymptotics,\allowbreak{} exact special case,\allowbreak{} coupling reduction or identifiability check}

\entrylabel{Argument or result}

For one class and attained service a,\allowbreak{} the index is the supremum over b>a of the Fbar-weighted average of h\_\allowbreak{}c(y)=\allowbreak{}c(y)h(y) on [a,\allowbreak{}b]. Its value G(a) therefore satisfies H(a)<=\allowbreak{}G(a)<=\allowbreak{}U(a),\allowbreak{} where H is the weighted average over [a,\allowbreak{}infinity) and U(a)=\allowbreak{}sup\_\allowbreak{}\{y>=\allowbreak{}a\}h\_\allowbreak{}c(y). A sufficient statewise certificate for the MAX-H chosen job i is H\_\allowbreak{}i(a\_\allowbreak{}i)>=\allowbreak{}max\_\allowbreak{}\{j!=\allowbreak{}i\} U\_\allowbreak{}j(a\_\allowbreak{}j). Then G\_\allowbreak{}i>=\allowbreak{}H\_\allowbreak{}i>=\allowbreak{}U\_\allowbreak{}j>=\allowbreak{}G\_\allowbreak{}j,\allowbreak{} so that choice is Gittins-optimal. This supplies an explicit conservative certificate computable from service and cost laws.


\entrylabel{Second check}

The upper bound follows because an average cannot exceed its integrand supremum;\allowbreak{} the lower bound follows by sending the stopping threshold b to infinity under the source integrability assumptions. The ordering chain is valid for ties as well.


\entrylabel{Remaining obstruction}

The certificate is often too strong or vacuous when U is infinite,\allowbreak{} and checking it on every reachable state is itself hard. It does not give the missing necessity characterization or improve the paper uniformly.


Verification:\allowbreak{} author checked partial | Independent review:\allowbreak{} none | Full target solved:\allowbreak{} False


\clearpage

\section*{P0406  Finite-system performance of size-based scheduling with impatience and noisy sizes}

\entrylabel{Primary source}

Ibrahim:\allowbreak{} Personalized scheduling in service systems;\allowbreak{} Problem statement and discussion;\allowbreak{} PDF2


\url{https://discovery.ucl.ac.uk/10147615/1/problem_Ibrahim.pdf}


Source statement date:\allowbreak{} 2022-04-01. Archived access:\allowbreak{} 2026-09-07T02:\allowbreak{}50:\allowbreak{}59Z.


Historical:\allowbreak{} source\_\allowbreak{}says\_\allowbreak{}open | Current:\allowbreak{} unverified | Result:\allowbreak{} unsolved


\entrylabel{Question}

Quantify the performance of size-based scheduling in multiserver queues with abandonment and noisy service-time estimates beyond the established overloaded many-server throughput limit.


\entrylabel{Scope}

Dong-Ibrahim later proves noisy-SJF asymptotic throughput results under its calibration assumptions. Those results are credited;\allowbreak{} finite-server performance and error control remain the retained part of the broad source question.


\entrylabel{Dated literature check}

Focused 2026-09-07 search for symmetric limits,\allowbreak{} correlated redundancy,\allowbreak{} slowdown and noisy scheduling;\allowbreak{} matched later slowdown theorem and noisy-SJF limits. No current-openness certificate.


\subsection*{P0406-A001 | Transform asymptotics,\allowbreak{} exact special case,\allowbreak{} coupling reduction or identifiability check}

\entrylabel{Argument or result}

The relevant ranking variable need not be the raw prediction. Let two prediction classes be z1<z2,\allowbreak{} but their true service times be deterministically 10 and 1. With one free server,\allowbreak{} no further arrivals and infinite patience for this local check,\allowbreak{} raw-prediction order has completion times 10,\allowbreak{}11 and mean 10.5;\allowbreak{} ordering by the conditional mean service times gives 1,\allowbreak{}11 and mean 6. Thus calibration monotonicity E[B|Z=\allowbreak{}z] increasing in z is a substantive assumption,\allowbreak{} not a consequence of having a prediction. Posterior service and patience information must enter a general performance model.


\entrylabel{Second check}

The completion sums are 21 and 12,\allowbreak{} respectively. The example is deliberately outside the later paper calibration hypothesis,\allowbreak{} so it cannot refute that theorem.


\entrylabel{Remaining obstruction}

Need a finite-server state description and quantitative waiting/\allowbreak{}abandonment analysis under stated prediction and patience laws. This local misranking example provides neither an error bound nor a general optimizer.


Verification:\allowbreak{} author checked partial | Independent review:\allowbreak{} none | Full target solved:\allowbreak{} False


\clearpage

\section*{P0408  Recovering an ideal Voronoi tessellation from vertices or its dual}

\entrylabel{Primary source}

D Achille Curien Enriquez Lyons Unel:\allowbreak{} Ideal Poisson Voronoi tessellations on hyperbolic spaces;\allowbreak{} Question7.2;\allowbreak{} PDF48


\url{https://arxiv.org/pdf/2303.16831}


Source statement date:\allowbreak{} see\_\allowbreak{}archived\_\allowbreak{}source. Archived access:\allowbreak{} 2026-09-07T05:\allowbreak{}50:\allowbreak{}05Z.


Historical:\allowbreak{} source\_\allowbreak{}says\_\allowbreak{}open | Current:\allowbreak{} unverified | Result:\allowbreak{} unsolved


\entrylabel{Question}

Determine whether the maps from a hyperbolic IPVT to its vertex set and to its ideal Delaunay tessellation are almost surely injective.


\entrylabel{Scope}

The two inverse-map questions are grouped. Theorem3.11 already reconstructs the ideal nuclei from richer tessellation information.


\entrylabel{Dated literature check}

Focused 2026-09-07 primary literature check;\allowbreak{} 2026 convergence/\allowbreak{}face-density and 2024 indistinguishability results matched. Present openness and novelty not certified.


\subsection*{P0408-A001 | Exact integral manipulation,\allowbreak{} Palm calculus,\allowbreak{} mass transport,\allowbreak{} limiting obstruction or conditional construction}

\entrylabel{Argument or result}

A sufficient reconstruction route is quantitative local decoding. Suppose for every rational test ball and every accuracy epsilon there is a measurable decoder from the observed vertex set in radius R whose answer about the tessellation in that ball fails with probability epsilon\_\allowbreak{}R tending to zero. Choose R\_\allowbreak{}k with epsilon\_\allowbreak{}Rk<=\allowbreak{}2\textasciicircum{}(-k). Borel-Cantelli then recovers all countably many test answers almost surely,\allowbreak{} yielding the complete random closed tessellation. This turns injectivity into a concrete conditional-uncertainty estimate.


\entrylabel{Second check}

The first Borel-Cantelli lemma requires no independence. A countable family of closed-set hit tests determines a random closed-set law and its realized closed sets under the standard Fell Borel structure.


\entrylabel{Remaining obstruction}

No such decoder or vanishing-error estimate is constructed. Vertex coordinates alone do not supply adjacency automatically,\allowbreak{} and the dual-map question needs its own decoding argument.


Verification:\allowbreak{} author checked partial | Independent review:\allowbreak{} none | Full target solved:\allowbreak{} False


\clearpage

\section*{P0409  Joint law of neighboring ideal Voronoi vertices}

\entrylabel{Primary source}

D Achille Curien Enriquez Lyons Unel:\allowbreak{} Ideal Poisson Voronoi tessellations on hyperbolic spaces;\allowbreak{} Question7.3;\allowbreak{} PDF48;\allowbreak{} Proposition4.4


\url{https://arxiv.org/pdf/2303.16831}


Source statement date:\allowbreak{} see\_\allowbreak{}archived\_\allowbreak{}source. Archived access:\allowbreak{} 2026-09-07T05:\allowbreak{}50:\allowbreak{}05Z.


Historical:\allowbreak{} source\_\allowbreak{}says\_\allowbreak{}open | Current:\allowbreak{} unverified | Result:\allowbreak{} unsolved


\entrylabel{Question}

Find the joint distribution of all neighboring vertices viewed from a typical vertex of the hyperbolic IPVT.


\entrylabel{Scope}

Typical means the vertex Palm law,\allowbreak{} not the cell containing a fixed point. The source gives Gamma(d+\allowbreak{}1) separation at a typical vertex.


\entrylabel{Dated literature check}

Focused 2026-09-07 primary literature check;\allowbreak{} 2026 convergence/\allowbreak{}face-density and 2024 indistinguishability results matched. Present openness and novelty not certified.


\subsection*{P0409-A001 | Exact integral manipulation,\allowbreak{} Palm calculus,\allowbreak{} mass transport,\allowbreak{} limiting obstruction or conditional construction}

\entrylabel{Argument or result}

The vertex graph has constant degree d+\allowbreak{}1. Under its stationary Palm law,\allowbreak{} rerooting at a uniformly selected neighbor preserves the marked-vertex marginal by mass transport:\allowbreak{} send f of the receiving marked vertex along every directed edge with weight 1/\allowbreak{}(d+\allowbreak{}1). Outgoing expectation is the neighbor average and incoming mass is f at the root. Thus the separation mark of a uniformly chosen neighbor also has the source Gamma(d+\allowbreak{}1) marginal. Any proposed joint neighbor formula must reproduce this law,\allowbreak{} not the Exp(1) law for a volume-typical point.


\entrylabel{Second check}

Constant degree makes the incoming weights sum to one. Without constant degree the stationary neighbor-rooted law would instead be degree-biased. The mass transport applies to the full marked configuration.


\entrylabel{Remaining obstruction}

Marginal preservation does not determine dependence between neighboring separation marks,\allowbreak{} lengths or angles. A joint Palm/\allowbreak{}Mecke calculation remains missing.


Verification:\allowbreak{} author checked partial | Independent review:\allowbreak{} none | Full target solved:\allowbreak{} False


\clearpage

\section*{P0410  Hole probabilities for lower-dimensional ideal Voronoi faces}

\entrylabel{Primary source}

D Achille Curien Enriquez Lyons Unel:\allowbreak{} Ideal Poisson Voronoi tessellations on hyperbolic spaces;\allowbreak{} Question7.4,\allowbreak{} hole-probability portion;\allowbreak{} PDF48


\url{https://arxiv.org/pdf/2303.16831}


Source statement date:\allowbreak{} see\_\allowbreak{}archived\_\allowbreak{}source. Archived access:\allowbreak{} 2026-09-07T05:\allowbreak{}50:\allowbreak{}05Z.


Historical:\allowbreak{} source\_\allowbreak{}says\_\allowbreak{}open | Current:\allowbreak{} unverified | Result:\allowbreak{} unsolved


\entrylabel{Question}

Compute the probabilities that observation balls miss k-faces of the hyperbolic IPVT for 0<=\allowbreak{}k<=\allowbreak{}d-2.


\entrylabel{Scope}

The face-volume-density portion of the source question is resolved by ideal\_\allowbreak{}face\_\allowbreak{}volumes\_\allowbreak{}2026,\allowbreak{} Corollary2.1. All dimensions of the remaining hole question are grouped.


\entrylabel{Dated literature check}

Focused 2026-09-07 primary literature check;\allowbreak{} 2026 convergence/\allowbreak{}face-density and 2024 indistinguishability results matched. Present openness and novelty not certified.


\subsection*{P0410-A001 | Exact integral manipulation,\allowbreak{} Palm calculus,\allowbreak{} mass transport,\allowbreak{} limiting obstruction or conditional construction}

\entrylabel{Argument or result}

For vertices,\allowbreak{} let N\_\allowbreak{}B count those in a ball B and let alpha\textasciicircum{}(2) be the second factorial moment measure. Bonferroni gives 1-E N\_\allowbreak{}B <=\allowbreak{} P(N\_\allowbreak{}B=\allowbreak{}0) <=\allowbreak{} 1-E N\_\allowbreak{}B+\allowbreak{}(1/\allowbreak{}2)*E[N\_\allowbreak{}B(N\_\allowbreak{}B-1)],\allowbreak{} with the lower bound truncated at zero if needed. The known intensity determines E N\_\allowbreak{}B=\allowbreak{}I\_\allowbreak{}d,\allowbreak{}0*vol(B),\allowbreak{} while the next term is alpha\textasciicircum{}(2)(B times B). Hence the new mean-density formulas identify only the first term of a small-ball expansion;\allowbreak{} a second-order joint-vertex calculation is already required for a quantitative remainder.


\entrylabel{Second check}

For each nonnegative integer n,\allowbreak{} 1-n<=\allowbreak{}1\{n=\allowbreak{}0\}<=\allowbreak{}1-n+\allowbreak{}n*(n-1)/\allowbreak{}2. Taking expectations proves the bounds without assuming Poisson vertices.


\entrylabel{Remaining obstruction}

The factorial moment measures and a controlled inclusion-exclusion remainder are unavailable here. Treating the vertex process as Poisson merely because the nuclei are Poisson would be unjustified.


Verification:\allowbreak{} author checked partial | Independent review:\allowbreak{} none | Full target solved:\allowbreak{} False


\clearpage

\section*{P0411  Boundary-hole monotonicity and coefficient positivity for ideal cells}

\entrylabel{Primary source}

D Achille Curien Enriquez Lyons Unel:\allowbreak{} Ideal Poisson Voronoi tessellations on hyperbolic spaces;\allowbreak{} Conjecture7.5;\allowbreak{} PDF48;\allowbreak{} Proposition5.1 PDF35


\url{https://arxiv.org/pdf/2303.16831}


Source statement date:\allowbreak{} 2025-03-19 (manuscript date;\allowbreak{} arXiv v3 uploaded 2025-06-10). Archived access:\allowbreak{} 2026-09-07T05:\allowbreak{}50:\allowbreak{}05Z.


Historical:\allowbreak{} source\_\allowbreak{}says\_\allowbreak{}open | Current:\allowbreak{} unverified | Result:\allowbreak{} candidate\_\allowbreak{}partial\_\allowbreak{}counterexample


\entrylabel{Question}

Establish the proposed normalized-hole monotonicity,\allowbreak{} odd-dimensional positive rational polynomial form,\allowbreak{} and even-dimensional positive power-series form,\allowbreak{} with the related decreasing-density behavior.


\entrylabel{Scope}

The source already gives d=\allowbreak{}3 explicitly and observes odd-dimensional rationality. The conjunction is counted once. A private partial packet proves monotonicity and odd-dimensional coefficients and gives a negative sixth Maclaurin coefficient in d=\allowbreak{}2 for the literal displayed distance variable. Remaining even-dimensional density/\allowbreak{}median claims are not proved;\allowbreak{} no novelty certification.


\entrylabel{Dated literature check}

Focused 2026-09-07 primary literature check;\allowbreak{} 2026 convergence/\allowbreak{}face-density and 2024 indistinguishability results matched. Present openness and novelty not certified.


\subsection*{P0411-A001 | Exact integral manipulation,\allowbreak{} Palm calculus,\allowbreak{} mass transport,\allowbreak{} limiting obstruction or conditional construction}

\entrylabel{Argument or result}

Put n=\allowbreak{}d-1 and J=\allowbreak{}1/\allowbreak{}H,\allowbreak{} where H is the boundary-hole probability. The exact source integral,\allowbreak{} after x=\allowbreak{}e\textasciicircum{}(-u)z and q=\allowbreak{}e\textasciicircum{}(2u),\allowbreak{} gives e\textasciicircum{}(-nu)J=\allowbreak{}c\_\allowbreak{}d\textasciicircum{}(-1) integral [sqrt((q+\allowbreak{}1)\textasciicircum{}2/\allowbreak{}4+\allowbreak{}|z|\textasciicircum{}2)-(q-1)/\allowbreak{}2]\textasciicircum{}(-2n) dz. The bracket strictly decreases with q for z nonzero,\allowbreak{} so e\textasciicircum{}(nu)H strictly decreases. When n=\allowbreak{}2m,\allowbreak{} a separate radial substitution gives J=\allowbreak{}e\textasciicircum{}(2mu)*p\_\allowbreak{}m(e\textasciicircum{}(2u)),\allowbreak{} where p\_\allowbreak{}m has degree m and strictly positive rational coefficients,\allowbreak{} written explicitly as integer-argument beta integrals in the proof packet. Subsequent exact calculation:\allowbreak{} in dimension two the required f(u)=\allowbreak{}exp(-u)/\allowbreak{}H\_\allowbreak{}2(u) has coefficient [u\textasciicircum{}6]f=\allowbreak{}61/\allowbreak{}240-4/\allowbreak{}(5*pi)<0,\allowbreak{} since pi<22/\allowbreak{}7. This is a candidate counterexample to the positive Maclaurin-coefficient clause in the displayed distance variable.


\entrylabel{Second check}

Two independent exact coefficient calculations agree for m=\allowbreak{}1 through20,\allowbreak{} with every coefficient positive and p\_\allowbreak{}m(1)=\allowbreak{}1. Twenty high-precision comparisons with the original angular integral agree below 1e-45 relative error. The known d=\allowbreak{}3 formula is recovered. These checks supplement,\allowbreak{} not replace,\allowbreak{} the written argument. Sixth coefficient checked by exact reciprocal-cosh recurrence,\allowbreak{} direct high-precision differentiation,\allowbreak{} and Cauchy extraction from the original angular integral;\allowbreak{} agreement within 5.5e-62.


\entrylabel{Remaining obstruction}

No independent review or novelty certification. The literal even-dimensional power-series clause has a candidate counterexample;\allowbreak{} a different intended variable would require source clarification. Even-dimensional decreasing-density/\allowbreak{}median claims beyond the source-known d=\allowbreak{}2 case are not proved here. The grouped entry remains partial.


Verification:\allowbreak{} author checked partial | Independent review:\allowbreak{} none | Full target solved:\allowbreak{} False


\clearpage

\section*{P0412  A general transfer theorem between horospherical and hyperbolic typical values}

\entrylabel{Primary source}

D Achille Curien Enriquez Lyons Unel:\allowbreak{} Ideal Poisson Voronoi tessellations on hyperbolic spaces;\allowbreak{} Question7.6;\allowbreak{} PDF48;\allowbreak{} Section5.3


\url{https://arxiv.org/pdf/2303.16831}


Source statement date:\allowbreak{} see\_\allowbreak{}archived\_\allowbreak{}source. Archived access:\allowbreak{} 2026-09-07T05:\allowbreak{}50:\allowbreak{}05Z.


Historical:\allowbreak{} source\_\allowbreak{}says\_\allowbreak{}open | Current:\allowbreak{} unverified | Result:\allowbreak{} unsolved


\entrylabel{Question}

Explain by a general theorem the matching typical geometric statistics of the stationary horospherical ideal-cell model and the hyperbolic IPVT.


\entrylabel{Scope}

The source records several matching examples. Equality for those functionals does not by itself identify the two Palm constructions.


\entrylabel{Dated literature check}

Focused 2026-09-07 primary literature check;\allowbreak{} 2026 convergence/\allowbreak{}face-density and 2024 indistinguishability results matched. Present openness and novelty not certified.


\subsection*{P0412-A001 | Exact integral manipulation,\allowbreak{} Palm calculus,\allowbreak{} mass transport,\allowbreak{} limiting obstruction or conditional construction}

\entrylabel{Argument or result}

A change of Palm sampling usually introduces a weight. If random measure tau is w times sigma,\allowbreak{} then the Campbell identity gives E\_\allowbreak{}tau\textasciicircum{}0 f=\allowbreak{}E\_\allowbreak{}sigma\textasciicircum{}0(w*f)/\allowbreak{}E\_\allowbreak{}sigma\textasciicircum{}0 w,\allowbreak{} provided the intensities are finite and positive. As a diagnostic,\allowbreak{} equal-probability marks of weights one and two become probabilities 1/\allowbreak{}3 and 2/\allowbreak{}3 under weighted sampling. A transfer theorem must therefore identify why the actual geometric weight cancels for the invariant functionals under consideration.


\entrylabel{Second check}

Integrate f against tau=\allowbreak{}w*sigma and apply the Campbell formula on a unit-volume window;\allowbreak{} division by the new intensity gives the displayed identity. The finite two-mark example checks the normalization.


\entrylabel{Remaining obstruction}

The actual horospherical-to-hyperbolic Palm weight,\allowbreak{} and the class of functionals for which it cancels,\allowbreak{} are not identified. Matching a few moments does not remove this obligation.


Verification:\allowbreak{} author checked partial | Independent review:\allowbreak{} none | Full target solved:\allowbreak{} False


\clearpage

\section*{P0413  Low-intensity variance of typical hyperbolic Voronoi cell volume}

\entrylabel{Primary source}

D Achille Curien Enriquez Lyons Unel:\allowbreak{} Ideal Poisson Voronoi tessellations on hyperbolic spaces;\allowbreak{} Question7.7;\allowbreak{} PDF48


\url{https://arxiv.org/pdf/2303.16831}


Source statement date:\allowbreak{} see\_\allowbreak{}archived\_\allowbreak{}source. Archived access:\allowbreak{} 2026-09-07T05:\allowbreak{}50:\allowbreak{}05Z.


Historical:\allowbreak{} source\_\allowbreak{}says\_\allowbreak{}open | Current:\allowbreak{} unverified | Result:\allowbreak{} unsolved


\entrylabel{Question}

Determine the finite limit,\allowbreak{} if any,\allowbreak{} of lambda times the typical-cell volume variance at actual nucleus intensity lambda,\allowbreak{} and its proposed coupling representation.


\entrylabel{Scope}

Use lambda for the actual intensity in this entry. The source uses an auxiliary reparametrization in parts of its proofs;\allowbreak{} no reparametrized intensity is substituted in the Palm identity.


\entrylabel{Dated literature check}

Focused 2026-09-07 primary literature check;\allowbreak{} 2026 convergence/\allowbreak{}face-density and 2024 indistinguishability results matched. Present openness and novelty not certified.


\subsection*{P0413-A001 | Exact integral manipulation,\allowbreak{} Palm calculus,\allowbreak{} mass transport,\allowbreak{} limiting obstruction or conditional construction}

\entrylabel{Argument or result}

Let V be typical-cell volume and V0 the zero-cell volume. Volume bias gives E V0=\allowbreak{}E V\textasciicircum{}2/\allowbreak{}E V,\allowbreak{} while mass conservation gives E V=\allowbreak{}1/\allowbreak{}lambda. Therefore E V0-E V=\allowbreak{}lambda*Var(V) exactly whenever these moments are finite. If a finite-intensity coupling produces a nonnegative excess D\_\allowbreak{}lambda converging in law to the proposed ideal-cell excess D,\allowbreak{} then Fatou gives only E D<=\allowbreak{}liminf E D\_\allowbreak{}lambda. A uniform bound on E D\_\allowbreak{}lambda\textasciicircum{}(1+\allowbreak{}epsilon) would supply the missing uniform integrability and hence the desired equality of limits.


\entrylabel{Second check}

Substituting E V=\allowbreak{}1/\allowbreak{}lambda into the variance decomposition cancels the squared mean term exactly. Nonnegative variables equal to n with probability 1/\allowbreak{}n and zero otherwise show why weak convergence alone cannot pass the expected excess to the limit.


\entrylabel{Remaining obstruction}

Neither a suitable finite-intensity excess coupling nor uniform integrability is proved. The algebraic variance identity is known Palm calculus and is not claimed as a solution.


Verification:\allowbreak{} author checked partial | Independent review:\allowbreak{} none | Full target solved:\allowbreak{} False


\clearpage

\section*{P0414  Hausdorff dimension of exceptional rays in the planar ideal Voronoi tree}

\entrylabel{Primary source}

D Achille Curien Enriquez Lyons Unel:\allowbreak{} Ideal Poisson Voronoi tessellations on hyperbolic spaces;\allowbreak{} Question7.9;\allowbreak{} PDF49


\url{https://arxiv.org/pdf/2303.16831}


Source statement date:\allowbreak{} see\_\allowbreak{}archived\_\allowbreak{}source. Archived access:\allowbreak{} 2026-09-07T05:\allowbreak{}50:\allowbreak{}05Z.


Historical:\allowbreak{} source\_\allowbreak{}says\_\allowbreak{}open | Current:\allowbreak{} unverified | Result:\allowbreak{} unsolved


\entrylabel{Question}

Determine the almost-sure Hausdorff dimension of rays whose separation-mark averages fail the value three obtained along cell-boundary rays.


\entrylabel{Scope}

The ray space is uncountable. Ergodicity along each ordinary boundary ray is not a simultaneous theorem about every ray.


\entrylabel{Dated literature check}

Focused 2026-09-07 primary literature check;\allowbreak{} 2026 convergence/\allowbreak{}face-density and 2024 indistinguishability results matched. Present openness and novelty not certified.


\subsection*{P0414-A001 | Exact integral manipulation,\allowbreak{} Palm calculus,\allowbreak{} mass transport,\allowbreak{} limiting obstruction or conditional construction}

\entrylabel{Argument or result}

A first-moment route would use a uniform large-deviation bound for all rooted length-n paths. If each such path has probability at most exp(-n*I) of a prescribed average deviation,\allowbreak{} there are 3*2\textasciicircum{}(n-1) paths in the 3-regular tree,\allowbreak{} so the expected number of bad prefixes is at most (3/\allowbreak{}2)*exp(n*(log(2)-I)). In a visual metric where a depth-n cylinder has diameter theta\textasciicircum{}n,\allowbreak{} this implies the limsup bad-ray set has dimension at most max(0,\allowbreak{}(log(2)-I)/\allowbreak{}log(1/\allowbreak{}theta));\allowbreak{} if I>log2 it is empty by Borel-Cantelli.


\entrylabel{Second check}

Sum the s-powers of cylinder diameters and use Tonelli. The expected total is summable whenever log2-I+\allowbreak{}s*log(theta)<0. Independence between distinct paths is unnecessary for this upper bound.


\entrylabel{Remaining obstruction}

No uniform deviation rate along arbitrary rays is proved,\allowbreak{} and a matching lower bound would require a correlated branching construction. One-vertex Gamma marginals are insufficient.


Verification:\allowbreak{} author checked partial | Independent review:\allowbreak{} none | Full target solved:\allowbreak{} False


\clearpage

\section*{P0415  Canonical intensity and typical-cell laws for unbounded hyperbolic cells}

\entrylabel{Primary source}

Hyperbolic meets stochastic geometry,\allowbreak{} Oberwolfach57/\allowbreak{}2025;\allowbreak{} Random tessellations overview;\allowbreak{} PDF16


\url{https://homeweb.unifr.ch/kellerha/pub/OWR-2025.pdf}


Source statement date:\allowbreak{} see\_\allowbreak{}archived\_\allowbreak{}source. Archived access:\allowbreak{} 2026-09-07T02:\allowbreak{}49:\allowbreak{}39Z.


Historical:\allowbreak{} source\_\allowbreak{}says\_\allowbreak{}open | Current:\allowbreak{} unverified | Result:\allowbreak{} unsolved


\entrylabel{Question}

Develop natural replacements for cell intensity and the typical-cell distribution when a stationary hyperbolic tessellation contains unbounded cells.


\entrylabel{Scope}

The source explains that an equivariant finite center cannot be selected for infinite-volume cells by the usual mass-transport construction.


\entrylabel{Dated literature check}

Focused 2026-09-07 primary literature check;\allowbreak{} 2026 convergence/\allowbreak{}face-density and 2024 indistinguishability results matched. Present openness and novelty not certified.


\subsection*{P0415-A001 | Exact integral manipulation,\allowbreak{} Palm calculus,\allowbreak{} mass transport,\allowbreak{} limiting obstruction or conditional construction}

\entrylabel{Argument or result}

The bounded-cell Palm relation already exposes the obstruction to simply retaining the old normalization. If finite cell intensity gamma and a probability-valued typical-cell law satisfied gamma*E[volume(C\_\allowbreak{}typ)]=\allowbreak{}1,\allowbreak{} assigning positive typical probability to infinite-volume cells would make the left side infinite for every gamma>0. Thus a putative extension must change the sampling measure,\allowbreak{} allow a sigma-finite object,\allowbreak{} or restrict the observables. A boundary- or horosphere-based sampling scheme is a possible route,\allowbreak{} but needs its own invariance statement.


\entrylabel{Second check}

The identity is obtained by assigning each location to its cell and integrating its unit mass via Campbell. The divergence conclusion follows from nonnegative integration and is not repaired by merely choosing another finite center.


\entrylabel{Remaining obstruction}

No canonical replacement measure or uniqueness principle is constructed. The contradiction applies to preserving the usual finite-intensity volume Palm identity,\allowbreak{} not to every conceivable generalized notion.


Verification:\allowbreak{} author checked partial | Independent review:\allowbreak{} none | Full target solved:\allowbreak{} False


\clearpage

\section*{P0416  Amenability of ideal Bernoulli Voronoi relations on nonexact groups}

\entrylabel{Primary source}

Jardon Sanchez Mellick Poulin Wrobel:\allowbreak{} Exactness and topology of invariant random equivalence relations;\allowbreak{} Question1;\allowbreak{} PDF18;\allowbreak{} Theorem4.5


\url{https://antoinegpoulin.github.io/pdf/exac.pdf}


Source statement date:\allowbreak{} see\_\allowbreak{}archived\_\allowbreak{}source. Archived access:\allowbreak{} 2026-09-07T05:\allowbreak{}50:\allowbreak{}12Z.


Historical:\allowbreak{} source\_\allowbreak{}says\_\allowbreak{}open | Current:\allowbreak{} unverified | Result:\allowbreak{} unsolved


\entrylabel{Question}

Determine whether ideal Bernoulli Voronoi equivalence relations are amenable for arbitrary countable groups,\allowbreak{} beyond the theorem for exact groups.


\entrylabel{Scope}

Limits of arbitrary finite invariant equivalence relations can be nonamenable on nonexact groups. That theorem is not automatically a Voronoi counterexample.


\entrylabel{Dated literature check}

Focused 2026-09-07 primary literature check;\allowbreak{} 2026 convergence/\allowbreak{}face-density and 2024 indistinguishability results matched. Present openness and novelty not certified.


\subsection*{P0416-A001 | Exact integral manipulation,\allowbreak{} Palm calculus,\allowbreak{} mass transport,\allowbreak{} limiting obstruction or conditional construction}

\entrylabel{Argument or result}

Uniform probability measures on finite approximating cells do not automatically pass to invariant probability kernels on the limiting classes. The elementary sequence uniform[-n,\allowbreak{}n] on Z has mass 1/\allowbreak{}(2n+\allowbreak{}1) at each included point;\allowbreak{} every coordinate tends to zero,\allowbreak{} so its pointwise limit is the zero measure,\allowbreak{} not a probability measure. For the Voronoi limit,\allowbreak{} a proof must retain an appropriate Reiter or hyperfinite approximation with a uniform error bound,\allowbreak{} rather than taking an unqualified limit of the finite-cell averages.


\entrylabel{Second check}

Sum each finite kernel to obtain one,\allowbreak{} but sum its pointwise limit to obtain zero. This directly demonstrates failure of the needed exchange of infinite summation and limit.


\entrylabel{Remaining obstruction}

No Voronoi-specific uniform amenability estimate is proved. Conversely,\allowbreak{} the known arbitrary finite-relation counterexamples need not arise from low-density Voronoi cells.


Verification:\allowbreak{} author checked partial | Independent review:\allowbreak{} none | Full target solved:\allowbreak{} False


\clearpage

\section*{P0417  Regular exponential volume growth without full graph IPVT convergence}

\entrylabel{Primary source}

D Achille and Khezeli:\allowbreak{} Convergence towards Ideal Poisson Voronoi tessellations with a focus on Diestel Leader graphs;\allowbreak{} Question7.1;\allowbreak{} PDF49;\allowbreak{} Theorem1.3


\url{https://arxiv.org/pdf/2606.30517}


Source statement date:\allowbreak{} 2026-07-20. Archived access:\allowbreak{} 2026-09-07T05:\allowbreak{}52:\allowbreak{}03Z.


Historical:\allowbreak{} source\_\allowbreak{}says\_\allowbreak{}open | Current:\allowbreak{} unverified | Result:\allowbreak{} unsolved


\entrylabel{Question}

Find or rule out a transitive graph with log-regularly-varying growth for which the boundary-measure convergence required in Theorem1.3 fails.


\entrylabel{Scope}

The intensity phase condition in Theorem1.3 is retained. Ordinary full convergence across all intensity phases is a different issue.


\entrylabel{Dated literature check}

Focused 2026-09-07 primary literature check;\allowbreak{} 2026 convergence/\allowbreak{}face-density and 2024 indistinguishability results matched. Present openness and novelty not certified.


\subsection*{P0417-A001 | Exact integral manipulation,\allowbreak{} Palm calculus,\allowbreak{} mass transport,\allowbreak{} limiting obstruction or conditional construction}

\entrylabel{Argument or result}

Write nu\_\allowbreak{}n for the uniform ball measure in the horocompactification and sigma\_\allowbreak{}n for the uniform sphere measure. If F(n-1)/\allowbreak{}F(n)->1/\allowbreak{}p,\allowbreak{} then nu\_\allowbreak{}n is asymptotically (1-1/\allowbreak{}p)*sigma\_\allowbreak{}n+\allowbreak{}(1/\allowbreak{}p)*nu\_\allowbreak{}(n-1). Alternating sphere limits alpha,\allowbreak{}beta would therefore produce ball subsequential limits (p*alpha+\allowbreak{}beta)/\allowbreak{}(p+\allowbreak{}1) and (alpha+\allowbreak{}p*beta)/\allowbreak{}(p+\allowbreak{}1),\allowbreak{} which differ whenever alpha differs from beta. This isolates the required mechanism:\allowbreak{} stable volume ratios can coexist algebraically with oscillating directional mass.


\entrylabel{Second check}

Solve the two linear equations a=\allowbreak{}(1-1/\allowbreak{}p)*alpha+\allowbreak{}b/\allowbreak{}p and b=\allowbreak{}(1-1/\allowbreak{}p)*beta+\allowbreak{}a/\allowbreak{}p. Direct substitution gives the stated two-cycle.


\entrylabel{Remaining obstruction}

The missing step is realizing such directional oscillation in a transitive or Cayley graph with the prescribed volume growth. An abstract sequence of sphere measures is not such a graph.


Verification:\allowbreak{} author checked partial | Independent review:\allowbreak{} none | Full target solved:\allowbreak{} False


\clearpage

\section*{P0418  Non-Poisson low-intensity ideal Voronoi diagrams}

\entrylabel{Primary source}

D Achille and Khezeli:\allowbreak{} Convergence towards Ideal Poisson Voronoi tessellations with a focus on Diestel Leader graphs;\allowbreak{} Question7.2;\allowbreak{} PDF49


\url{https://arxiv.org/pdf/2606.30517}


Source statement date:\allowbreak{} 2026-07-20. Archived access:\allowbreak{} 2026-09-07T05:\allowbreak{}52:\allowbreak{}03Z.


Historical:\allowbreak{} source\_\allowbreak{}says\_\allowbreak{}open | Current:\allowbreak{} unverified | Result:\allowbreak{} unsolved


\entrylabel{Question}

Construct a nontrivial low-intensity ideal diagram with a law different from the IPVT using a non-Poisson nuclei process.


\entrylabel{Scope}

The source suggests the alpha-hyperbolic ensemble. Merely making the finite-intensity point process non-Poisson need not alter the limiting diagram.


\entrylabel{Dated literature check}

Focused 2026-09-07 primary literature check;\allowbreak{} 2026 convergence/\allowbreak{}face-density and 2024 indistinguishability results matched. Present openness and novelty not certified.


\subsection*{P0418-A001 | Exact integral manipulation,\allowbreak{} Palm calculus,\allowbreak{} mass transport,\allowbreak{} limiting obstruction or conditional construction}

\entrylabel{Argument or result}

A simple global Cox mixture cannot achieve the target when the underlying Poisson model has a unique intensity-independent limit. Let Z be any positive finite random variable and,\allowbreak{} conditionally on Z,\allowbreak{} take a Poisson nuclei process of intensity lambda*Z. For every bounded continuous diagram test f,\allowbreak{} conditional convergence at each fixed Z=\allowbreak{}z gives E[f(V\_\allowbreak{}lambda)|Z=\allowbreak{}z]->E[f(V\_\allowbreak{}IPVT)]. Dominated convergence then yields the same unconditional IPVT law. Randomizing the global intensity changes finite-scale counts but not this unique limit.


\entrylabel{Second check}

The conditional intensity tends to zero for every 0<z<infinity,\allowbreak{} and boundedness of f supplies domination. An atom at Z=\allowbreak{}0 would add the trivial empty-nuclei case,\allowbreak{} which is excluded by the stated assumptions.


\entrylabel{Remaining obstruction}

Need spatial dependence that survives the boundary-and-delay limit. The Cox comparison does not analyze the determinantal candidate or construct a distinct nontrivial limit.


Verification:\allowbreak{} author checked partial | Independent review:\allowbreak{} none | Full target solved:\allowbreak{} False


\clearpage

\section*{P0419  Equal-growth L1 products beyond the divergent tilted-volume condition}

\entrylabel{Primary source}

D Achille and Khezeli:\allowbreak{} Convergence towards Ideal Poisson Voronoi tessellations with a focus on Diestel Leader graphs;\allowbreak{} Question7.3;\allowbreak{} PDF49;\allowbreak{} Proposition5.14(i)


\url{https://arxiv.org/pdf/2606.30517}


Source statement date:\allowbreak{} 2026-07-20. Archived access:\allowbreak{} 2026-09-07T05:\allowbreak{}52:\allowbreak{}03Z.


Historical:\allowbreak{} source\_\allowbreak{}says\_\allowbreak{}open | Current:\allowbreak{} unverified | Result:\allowbreak{} unsolved


\entrylabel{Question}

Determine IPVT convergence for equal-growth product factors when the tilted-volume integrals required in Proposition5.14(i) need not diverge.


\entrylabel{Scope}

The source condition is integral exp(-b*s)dF(s)=\allowbreak{}infinity for each factor,\allowbreak{} with their common growth exponent b. The following is a radial comparison model only.


\entrylabel{Dated literature check}

Focused 2026-09-07 primary literature check;\allowbreak{} 2026 convergence/\allowbreak{}face-density and 2024 indistinguishability results matched. Present openness and novelty not certified.


\subsection*{P0419-A001 | Exact integral manipulation,\allowbreak{} Palm calculus,\allowbreak{} mass transport,\allowbreak{} limiting obstruction or conditional construction}

\entrylabel{Argument or result}

Consider two equal radial densities proportional at large r to exp(b*r)*(1+\allowbreak{}r)\textasciicircum{}(-a),\allowbreak{} a>1. On a large L1 sphere,\allowbreak{} the split at radius s has weight proportional to (1+\allowbreak{}s)\textasciicircum{}(-a)*(1+\allowbreak{}R-s)\textasciicircum{}(-a). The contribution from epsilon*R<=\allowbreak{}s<=\allowbreak{}(1-epsilon)*R is O(R\textasciicircum{}(1-2a)),\allowbreak{} whereas either fixed endpoint region contributes order R\textasciicircum{}(-a). Their ratio tends to zero because a>1. Thus equal exponential growth alone need not force both factor radii to diverge;\allowbreak{} endpoint concentration can replace the tensor product of two boundary laws.


\entrylabel{Second check}

In the middle interval both denominators are bounded below by constants times R\textasciicircum{}a and its length is at most R. Near an endpoint,\allowbreak{} one denominator remains bounded while the other has order R\textasciicircum{}a. These independent upper and lower estimates establish the comparison.


\entrylabel{Remaining obstruction}

The model is not asserted to be a homogeneous geometric counterexample. A general theorem needs precise radial regularity,\allowbreak{} the endpoint mixture weights,\allowbreak{} and nondegeneracy of the resulting ideal diagram.


Verification:\allowbreak{} author checked partial | Independent review:\allowbreak{} none | Full target solved:\allowbreak{} False


\clearpage

\section*{P0420  Low-intensity Voronoi limits for products with the L2 metric}

\entrylabel{Primary source}

D Achille and Khezeli:\allowbreak{} Convergence towards Ideal Poisson Voronoi tessellations with a focus on Diestel Leader graphs;\allowbreak{} Question7.4;\allowbreak{} PDF49


\url{https://arxiv.org/pdf/2606.30517}


Source statement date:\allowbreak{} 2026-07-20. Archived access:\allowbreak{} 2026-09-07T05:\allowbreak{}52:\allowbreak{}03Z.


Historical:\allowbreak{} source\_\allowbreak{}says\_\allowbreak{}open | Current:\allowbreak{} unverified | Result:\allowbreak{} unsolved


\entrylabel{Question}

Give conditions and identify the limiting boundary measure for products with the L2 metric,\allowbreak{} including comparison with an appropriately weighted L1 model.


\entrylabel{Scope}

The source conjectures a relation under mild conditions. A formal saddle-point calculation alone is not a convergence theorem.


\entrylabel{Dated literature check}

Focused 2026-09-07 primary literature check;\allowbreak{} 2026 convergence/\allowbreak{}face-density and 2024 indistinguishability results matched. Present openness and novelty not certified.


\subsection*{P0420-A001 | Exact integral manipulation,\allowbreak{} Palm calculus,\allowbreak{} mass transport,\allowbreak{} limiting obstruction or conditional construction}

\entrylabel{Argument or result}

If radial shell densities have leading exponents b*r1 and bprime*r2 with both positive,\allowbreak{} their product on r1\textasciicircum{}2+\allowbreak{}r2\textasciicircum{}2=\allowbreak{}R\textasciicircum{}2 is governed by maximizing b*r1+\allowbreak{}bprime*r2. Cauchy-Schwarz gives the unique maximum B*R at (r1,\allowbreak{}r2)=\allowbreak{}R*(b,\allowbreak{}bprime)/\allowbreak{}B,\allowbreak{} B=\allowbreak{}sqrt(b\textasciicircum{}2+\allowbreak{}bprime\textasciicircum{}2). Taylor expansion of the product distance then suggests the limiting horofunction (b*h1+\allowbreak{}bprime*h2)/\allowbreak{}B. The second derivative in r1 at the maximizer is -B\textasciicircum{}3/\allowbreak{}(bprime\textasciicircum{}2*R),\allowbreak{} so the radial split fluctuates on the sqrt(R) scale under regular Laplace asymptotics.


\entrylabel{Second check}

The maximum follows independently from a Lagrange multiplier. Differentiate b*s+\allowbreak{}bprime*sqrt(R\textasciicircum{}2-s\textasciicircum{}2) twice to verify the curvature. Zero growth exponents require separate endpoint analysis.


\entrylabel{Remaining obstruction}

Need uniform asymptotics for shell measures and angular laws,\allowbreak{} control of the saddle-point tails,\allowbreak{} and convergence of the boundary measures. Exponential rates alone do not justify these steps.


Verification:\allowbreak{} author checked partial | Independent review:\allowbreak{} none | Full target solved:\allowbreak{} False


\clearpage

\section*{P0421  Convergence of edge-measured product-graph Voronoi diagrams}

\entrylabel{Primary source}

D Achille and Khezeli:\allowbreak{} Convergence towards Ideal Poisson Voronoi tessellations with a focus on Diestel Leader graphs;\allowbreak{} Question7.5;\allowbreak{} PDF49;\allowbreak{} Theorem4.12


\url{https://arxiv.org/pdf/2606.30517}


Source statement date:\allowbreak{} 2026-07-20. Archived access:\allowbreak{} 2026-09-07T05:\allowbreak{}52:\allowbreak{}03Z.


Historical:\allowbreak{} source\_\allowbreak{}says\_\allowbreak{}open | Current:\allowbreak{} unverified | Result:\allowbreak{} unsolved


\entrylabel{Question}

Characterize when weighted edge-measured product graphs satisfy the full low-intensity convergence criterion of Theorem4.12.


\entrylabel{Scope}

The theorem involves boundary behavior of edge endpoints;\allowbreak{} a vertex measure alone can lose this information.


\entrylabel{Dated literature check}

Focused 2026-09-07 primary literature check;\allowbreak{} 2026 convergence/\allowbreak{}face-density and 2024 indistinguishability results matched. Present openness and novelty not certified.


\subsection*{P0421-A001 | Exact integral manipulation,\allowbreak{} Palm calculus,\allowbreak{} mass transport,\allowbreak{} limiting obstruction or conditional construction}

\entrylabel{Argument or result}

For a finite weighted edge set,\allowbreak{} choosing an edge proportionally to its weight and then a uniform endpoint samples vertices with probability proportional to weighted degree. Thus a uniform-vertex limit need not equal the endpoint limit when degrees or weights vary. Even knowledge of both endpoint marginals does not identify their joint boundary law:\allowbreak{} an equal mixture of pairs (0,\allowbreak{}0),\allowbreak{}(1,\allowbreak{}1) and an equal mixture of (0,\allowbreak{}1),\allowbreak{}(1,\allowbreak{}0) have the same marginals but disjoint supports. Product-edge convergence must retain the endpoint-pair measure and horizontal/\allowbreak{}vertical edge weights.


\entrylabel{Second check}

Each incident edge contributes half its sampling weight to a vertex,\allowbreak{} proving the degree bias. The two binary pair laws both have Bernoulli(1/\allowbreak{}2) marginals,\allowbreak{} checking the independent loss of correlation information.


\entrylabel{Remaining obstruction}

No convergence estimate for the actual product-graph endpoint measure is proved. The pair-law example is an information diagnostic,\allowbreak{} not a constructed graph counterexample.


Verification:\allowbreak{} author checked partial | Independent review:\allowbreak{} none | Full target solved:\allowbreak{} False


\clearpage

\section*{P0422  More than two distinguishable types of group IPVT cells}

\entrylabel{Primary source}

D Achille and Khezeli:\allowbreak{} Convergence towards Ideal Poisson Voronoi tessellations with a focus on Diestel Leader graphs;\allowbreak{} Question7.6;\allowbreak{} PDF49;\allowbreak{} Proposition6.7


\url{https://arxiv.org/pdf/2606.30517}


Source statement date:\allowbreak{} 2026-07-20. Archived access:\allowbreak{} 2026-09-07T05:\allowbreak{}52:\allowbreak{}03Z.


Historical:\allowbreak{} source\_\allowbreak{}says\_\allowbreak{}open | Current:\allowbreak{} unverified | Result:\allowbreak{} unsolved


\entrylabel{Question}

Construct or rule out group IPVTs with a finite number of distinguishability types greater than two,\allowbreak{} or infinitely many types.


\entrylabel{Scope}

The directed Diestel-Leader/\allowbreak{}group-action example already has two types. The symmetry group is part of the statement:\allowbreak{} the undirected automorphism group can merge those types.


\entrylabel{Dated literature check}

Focused 2026-09-07 primary literature check;\allowbreak{} 2026 convergence/\allowbreak{}face-density and 2024 indistinguishability results matched. Present openness and novelty not certified.


\subsection*{P0422-A001 | Exact integral manipulation,\allowbreak{} Palm calculus,\allowbreak{} mass transport,\allowbreak{} limiting obstruction or conditional construction}

\entrylabel{Argument or result}

A sufficient construction route is an equivariant recoverable nucleus mark. Suppose the limiting boundary decomposes into invariant measurable sets A1,\allowbreak{}...,\allowbreak{}Ak of positive intensity,\allowbreak{} nuclei in every Ai produce nonempty cells,\allowbreak{} and the associated boundary label can be reconstructed measurably from each cell and its diagram. Then selecting cells whose label is Ai is an equivariant rule,\allowbreak{} so these classes are distinguishable. Independent restrictions of the boundary Poisson process preserve all k label populations,\allowbreak{} but geometric visibility of those nuclei remains an additional requirement.


\entrylabel{Second check}

Equivariance of the recovered label implies invariance of each selection rule. Positive intensity alone is not enough:\allowbreak{} nuclei with empty cells would not provide a cell type,\allowbreak{} which is why nonemptiness was explicitly assumed.


\entrylabel{Remaining obstruction}

No group with the required boundary decomposition and visibility properties is constructed. The criterion gives at least k distinguishable classes,\allowbreak{} not exactly k;\allowbreak{} within-class indistinguishability still requires proof.


Verification:\allowbreak{} author checked partial | Independent review:\allowbreak{} none | Full target solved:\allowbreak{} False


\clearpage

\section*{P0425  Optimal parallel allocation with unknown nonexponential job sizes}

\entrylabel{Primary source}

Harchol-Balter:\allowbreak{} Open problems in queueing theory inspired by datacenter computing;\allowbreak{} Section3.2;\allowbreak{} PDF8


\url{https://www.cs.cmu.edu/~harchol/Papers/QUESTA21.pdf}


Source statement date:\allowbreak{} 2021-01-27. Archived access:\allowbreak{} 2026-09-07T02:\allowbreak{}54:\allowbreak{}43Z.


Historical:\allowbreak{} source\_\allowbreak{}says\_\allowbreak{}open | Current:\allowbreak{} unverified | Result:\allowbreak{} unsolved


\entrylabel{Question}

For Poisson arrivals and a common concave speedup curve,\allowbreak{} characterize the mean-response-time optimal allocation when inherent sizes are unknown and nonexponential.


\entrylabel{Scope}

The exponential EQUI theorem is already known;\allowbreak{} the scheduler observes attained service and the size law. Multiple servers and sublinear parallel speedup are essential to the remaining question.


\entrylabel{Dated literature check}

Focused later-primary-source checks on 2026-09-07;\allowbreak{} offline concave scheduling,\allowbreak{} general FCFS threshold representation,\allowbreak{} SEK-SMOD and fork-join asymptotics credited. No present-openness certification.


\subsection*{P0425-A001 | Controlled generators,\allowbreak{} exact finite-state diagnostic,\allowbreak{} coupling,\allowbreak{} renewal reward or capacity comparison}

\entrylabel{Argument or result}

For an absolutely continuous size law with hazard h,\allowbreak{} a job at attained service a allocated b servers completes at instantaneous rate h(a)*s(b),\allowbreak{} while its age advances at rate s(b). A formal controlled generator therefore has both an age-drift term s(b)*partial\_\allowbreak{}a V and a completion term h(a)*s(b)*(V\_\allowbreak{}after\_\allowbreak{}departure-V). Maximizing the departure rate alone drops the learning and continuation terms. The common exponential case removes age learning but does not justify that omission for a general size law.


\entrylabel{Second check}

Conditional survival over work da is 1-h(a)*da+\allowbreak{}o(da);\allowbreak{} substituting da=\allowbreak{}s(b)*dt derives both generator terms from the same clock. A nonexponential two-point law instead requires jump boundary conditions,\allowbreak{} so the density formula is explicitly limited.


\entrylabel{Remaining obstruction}

No solution of the infinite-state average-cost control equation or verification theorem is supplied.


Verification:\allowbreak{} author checked partial | Independent review:\allowbreak{} none | Full target solved:\allowbreak{} False


\clearpage

\section*{P0426  Optimal known-size parallel allocation with Poisson arrivals}

\entrylabel{Primary source}

Harchol-Balter:\allowbreak{} Open problems in queueing theory inspired by datacenter computing;\allowbreak{} Section3.3;\allowbreak{} PDF10


\url{https://www.cs.cmu.edu/~harchol/Papers/QUESTA21.pdf}


Source statement date:\allowbreak{} 2021-01-27. Archived access:\allowbreak{} 2026-09-07T02:\allowbreak{}54:\allowbreak{}43Z.


Historical:\allowbreak{} source\_\allowbreak{}says\_\allowbreak{}open | Current:\allowbreak{} unverified | Result:\allowbreak{} unsolved


\entrylabel{Question}

Determine optimal mean-response-time allocation for known-size parallel jobs arriving over time under a common concave speedup curve.


\entrylabel{Scope}

The all-jobs-at-zero case has later general concave algorithms;\allowbreak{} it is excluded separately. Those offline results do not solve the Poisson-arrival problem.


\entrylabel{Dated literature check}

Focused later-primary-source checks on 2026-09-07;\allowbreak{} offline concave scheduling,\allowbreak{} general FCFS threshold representation,\allowbreak{} SEK-SMOD and fork-join asymptotics credited. No present-openness certification.


\subsection*{P0426-A001 | Controlled generators,\allowbreak{} exact finite-state diagnostic,\allowbreak{} coupling,\allowbreak{} renewal reward or capacity comparison}

\entrylabel{Argument or result}

With remaining-work vector x,\allowbreak{} arrival rate lambda,\allowbreak{} and relative value V,\allowbreak{} the formal average-cost equation is g=\allowbreak{}m+\allowbreak{}lambda*E[V(x appended with S)-V(x)]-max\_\allowbreak{}\{sum b\_\allowbreak{}i<=\allowbreak{}n\} sum\_\allowbreak{}i s(b\_\allowbreak{}i)*partial\_\allowbreak{}i V(x). Interior allocations must satisfy s\_\allowbreak{}prime(b\_\allowbreak{}i)*partial\_\allowbreak{}i V=\allowbreak{}nu. Thus the offline allocation can persist only if the stochastic value gradients preserve the required ratios;\allowbreak{} arrivals explicitly couple all coordinates through the appended-state term.


\entrylabel{Second check}

Between arrivals the residual of job i decreases at rate s(b\_\allowbreak{}i),\allowbreak{} giving the minus sign. For linear speedup the maximization concentrates resources on the largest value gradient;\allowbreak{} this is consistent with a priority rule,\allowbreak{} not proof of which priority.


\entrylabel{Remaining obstruction}

The value gradients,\allowbreak{} boundary conditions at completion,\allowbreak{} recurrence assumptions and optimality verification remain unknown.


Verification:\allowbreak{} author checked partial | Independent review:\allowbreak{} none | Full target solved:\allowbreak{} False


\clearpage

\section*{P0427  Optimal allocation between two exponential parallelism classes}

\entrylabel{Primary source}

Harchol-Balter:\allowbreak{} Open problems in queueing theory inspired by datacenter computing;\allowbreak{} Section3.4;\allowbreak{} PDF10-11


\url{https://www.cs.cmu.edu/~harchol/Papers/QUESTA21.pdf}


Source statement date:\allowbreak{} 2021-01-27. Archived access:\allowbreak{} 2026-09-07T02:\allowbreak{}54:\allowbreak{}43Z.


Historical:\allowbreak{} source\_\allowbreak{}says\_\allowbreak{}open | Current:\allowbreak{} unverified | Result:\allowbreak{} unsolved


\entrylabel{Question}

Characterize optimal allocation when two classes have different ordered concave speedup curves and unknown exponential inherent sizes.


\entrylabel{Scope}

The source already shows that instantaneous maximum-throughput allocation need not minimize delay. This attempt supplies a small diagnostic,\allowbreak{} not a new counterexample claim.


\entrylabel{Dated literature check}

Focused later-primary-source checks on 2026-09-07;\allowbreak{} offline concave scheduling,\allowbreak{} general FCFS threshold representation,\allowbreak{} SEK-SMOD and fork-join asymptotics credited. No present-openness certification.


\subsection*{P0427-A001 | Controlled generators,\allowbreak{} exact finite-state diagnostic,\allowbreak{} coupling,\allowbreak{} renewal reward or capacity comparison}

\entrylabel{Argument or result}

Even with two jobs and no arrivals,\allowbreak{} continuation values change the allocation. Take four divisible servers and independent Exp(1) inherent job sizes. Use normalized concave speedups s1(b)=\allowbreak{}min(b,\allowbreak{}sqrt(b)) and s2(b)=\allowbreak{}min(b,\allowbreak{}b\textasciicircum{}(3/\allowbreak{}4)),\allowbreak{} including linear fractional-server behavior below one server. Allocate b to class one and 4-b to class two until the first completion,\allowbreak{} then all four to the survivor. Expected sum of completion times is V(b)=\allowbreak{}[2+\allowbreak{}s1(b)/\allowbreak{}(2*sqrt(2))+\allowbreak{}s2(4-b)/\allowbreak{}2]/\allowbreak{}[s1(b)+\allowbreak{}s2(4-b)]. The throughput maximizer is b=\allowbreak{}1:\allowbreak{} total rate increases below one and decreases above it. Yet V(81/\allowbreak{}64)<V(1). The slower departure rate improves the continuation state.


\entrylabel{Second check}

The first-completion time has mean 1/\allowbreak{}(r1+\allowbreak{}r2) and incurs holding cost two;\allowbreak{} survivor means are 1/\allowbreak{}(2*sqrt(2)) and 1/\allowbreak{}2. Independently,\allowbreak{} the sign of the right derivative V\_\allowbreak{}prime(1) is the sign of 3\textasciicircum{}(3/\allowbreak{}4)*(3/\allowbreak{}(2*sqrt(2))+\allowbreak{}1/\allowbreak{}2)-4,\allowbreak{} which is negative by the rational bounds 3\textasciicircum{}(3/\allowbreak{}4)<23/\allowbreak{}10 and 1/\allowbreak{}(2*sqrt(2))<9/\allowbreak{}25.


\entrylabel{Remaining obstruction}

The no-arrival diagnostic does not determine a stationary optimal policy with Poisson arrivals,\allowbreak{} nor prove a universal allocation formula.


Verification:\allowbreak{} author checked partial | Independent review:\allowbreak{} none | Full target solved:\allowbreak{} False


\clearpage

\section*{P0428  Mean-response optimization for stochastically arriving DAG jobs}

\entrylabel{Primary source}

Harchol-Balter:\allowbreak{} Open problems in queueing theory inspired by datacenter computing;\allowbreak{} Section4;\allowbreak{} PDF13-14


\url{https://www.cs.cmu.edu/~harchol/Papers/QUESTA21.pdf}


Source statement date:\allowbreak{} 2021-01-27. Archived access:\allowbreak{} 2026-09-07T02:\allowbreak{}54:\allowbreak{}43Z.


Historical:\allowbreak{} source\_\allowbreak{}says\_\allowbreak{}open | Current:\allowbreak{} unverified | Result:\allowbreak{} unsolved


\entrylabel{Question}

Find optimal scheduling or tractable mean job-response laws for DAG jobs with random arrivals and task durations,\allowbreak{} retaining task precedence.


\entrylabel{Scope}

Different graph shapes and information models must be specified. This attempt uses known finite DAGs and independent exponential task times as a concrete submodel.


\entrylabel{Dated literature check}

Focused later-primary-source checks on 2026-09-07;\allowbreak{} offline concave scheduling,\allowbreak{} general FCFS threshold representation,\allowbreak{} SEK-SMOD and fork-join asymptotics credited. No present-openness certification.


\subsection*{P0428-A001 | Controlled generators,\allowbreak{} exact finite-state diagnostic,\allowbreak{} coupling,\allowbreak{} renewal reward or capacity comparison}

\entrylabel{Argument or result}

Total work does not identify a job duration. Two independent Exp(mu) tasks in parallel have mean completion time 3/\allowbreak{}(2*mu),\allowbreak{} while the same two tasks in a serial chain have mean 2/\allowbreak{}mu;\allowbreak{} both have mean total work 2/\allowbreak{}mu. For general DAGs a Markov state must retain the completed downward-closed vertex set and currently served eligible tasks. Lower bounds are the expected critical-path duration and expected total work divided by server count,\allowbreak{} but neither retains inter-job interference.


\entrylabel{Second check}

For two exponential tasks,\allowbreak{} E[max]=\allowbreak{}E[X]+\allowbreak{}E[Y]-E[min]=\allowbreak{}2/\allowbreak{}mu-1/\allowbreak{}(2*mu). A chain requires the sum. These two exact laws independently isolate the precedence information lost by a work-only state.


\entrylabel{Remaining obstruction}

No optimal stochastic DAG allocation or stationary response formula is derived. Worst-case offline guarantees are not substituted for the stochastic objective.


Verification:\allowbreak{} author checked partial | Independent review:\allowbreak{} none | Full target solved:\allowbreak{} False


\clearpage

\section*{P0429  Largest growing fork-join batch compatible with asymptotic independence}

\entrylabel{Primary source}

Harchol-Balter:\allowbreak{} Open problems in queueing theory inspired by datacenter computing;\allowbreak{} Section5;\allowbreak{} PDF16


\url{https://www.cs.cmu.edu/~harchol/Papers/QUESTA21.pdf}


Source statement date:\allowbreak{} 2021-01-27. Archived access:\allowbreak{} 2026-09-07T02:\allowbreak{}54:\allowbreak{}43Z.


Historical:\allowbreak{} source\_\allowbreak{}says\_\allowbreak{}open | Current:\allowbreak{} unverified | Result:\allowbreak{} candidate\_\allowbreak{}partial\_\allowbreak{}counterexample


\entrylabel{Question}

Determine how quickly the number k(n) of tasks per uniformly routed job can grow while the corresponding k(n) stationary queues remain asymptotically independent.


\entrylabel{Scope}

The source credits k=\allowbreak{}o(n\textasciicircum{}(1/\allowbreak{}4));\allowbreak{} fixed-dimensional convergence alone does not answer the growing-dimensional question. Private author-only packet extends the credited coupling under exponential service. Later primary-source leads do not establish novelty.


\entrylabel{Dated literature check}

Focused later-primary-source checks on 2026-09-07;\allowbreak{} offline concave scheduling,\allowbreak{} general FCFS threshold representation,\allowbreak{} SEK-SMOD and fork-join asymptotics credited. No present-openness certification.


\subsection*{P0429-A001 | Controlled generators,\allowbreak{} exact finite-state diagnostic,\allowbreak{} coupling,\allowbreak{} renewal reward or capacity comparison}

\entrylabel{Argument or result}

A finite-horizon input coupling suggests a cubic collision scale. Fix k tagged servers. Jobs arrive at rate Lambda=\allowbreak{}lambda*n/\allowbreak{}k and choose a uniform k-subset. If N is the number of tagged servers hit by an arrival,\allowbreak{} E[N(N-1)]=\allowbreak{}k*(k-1)*k*(k-1)/\allowbreak{}(n*(n-1)). Remove arrivals that hit at least two tagged queues and restore their lost marginal rates by independent Poisson streams. A union bound on removed events and added events is at most (3/\allowbreak{}2)*Lambda*T*E[N(N-1)],\allowbreak{} hence at most 3*lambda*T*k*(k-1)\textasciicircum{}2/\allowbreak{}(2*(n-1)). The resulting streams are independent with the correct marginal rates. With independent task marks and queues initially empty,\allowbreak{} this also bounds total variation of their workload paths on [0,\allowbreak{}T]. Extending the credited source coupling to stationary starts gives TV<=\allowbreak{}2*k*r(T)+\allowbreak{}3*lambda*T*k*(k-1)\textasciicircum{}2/\allowbreak{}(2*(n-1)),\allowbreak{} where r(T) is the stationary marginal emptying tail. With independent exponential service of rate mu>lambda,\allowbreak{} a stopped generator calculation gives r(T)<=\allowbreak{}(1+\allowbreak{}sqrt(lambda/\allowbreak{}mu))*exp(-(sqrt(mu)-sqrt(lambda))\textasciicircum{}2*T). Taking T=\allowbreak{}2*log(n)/\allowbreak{}(sqrt(mu)-sqrt(lambda))\textasciicircum{}2 gives sufficient k=\allowbreak{}o((n/\allowbreak{}log(n))\textasciicircum{}(1/\allowbreak{}3)). A conditional polynomial tail r(T)<=\allowbreak{}C*T\textasciicircum{}(-p) gives sufficient exponent p/\allowbreak{}(3*p+\allowbreak{}1).


\entrylabel{Second check}

The hypergeometric factorial moment follows by summing ordered pairs of tagged-server indicators. On N>=\allowbreak{}2,\allowbreak{} both 1<=\allowbreak{}N(N-1)/\allowbreak{}2 and N<=\allowbreak{}N(N-1),\allowbreak{} giving the factor 3/\allowbreak{}2. Small subset enumeration checks the moment exactly. Separate exact checks cover 77 hypergeometric/\allowbreak{}restoration cases,\allowbreak{} 81 generator identities,\allowbreak{} nine geometric exponential moments and five polynomial balancing identities. The complete packet explicitly checks stationary coupling,\allowbreak{} marginal rate restoration,\allowbreak{} PASTA and heavy-traffic limits;\allowbreak{} these are author checks,\allowbreak{} not independent review.


\entrylabel{Remaining obstruction}

Partial exponential-service sufficient regime only. No optimal threshold,\allowbreak{} boundary k of order n\textasciicircum{}(1/\allowbreak{}3),\allowbreak{} general-service improvement,\allowbreak{} correlated-mark result,\allowbreak{} uniform heavy-traffic result or novelty certification. The basic source coupling is credited.


Verification:\allowbreak{} author checked partial | Independent review:\allowbreak{} none | Full target solved:\allowbreak{} False


\clearpage

\section*{P0430  Response laws for limited fork-join shortest-queue dispatching}

\entrylabel{Primary source}

Harchol-Balter:\allowbreak{} Open problems in queueing theory inspired by datacenter computing;\allowbreak{} Section5;\allowbreak{} PDF16


\url{https://www.cs.cmu.edu/~harchol/Papers/QUESTA21.pdf}


Source statement date:\allowbreak{} 2021-01-27. Archived access:\allowbreak{} 2026-09-07T02:\allowbreak{}54:\allowbreak{}43Z.


Historical:\allowbreak{} source\_\allowbreak{}says\_\allowbreak{}open | Current:\allowbreak{} unverified | Result:\allowbreak{} unsolved


\entrylabel{Question}

Analyze job response time when k distinct tasks are sent to the k shortest queues or the k least-workload queues.


\entrylabel{Scope}

Random-routing independence results do not transfer automatically to state-dependent routing;\allowbreak{} queue-length and workload policies use different information.


\entrylabel{Dated literature check}

Focused later-primary-source checks on 2026-09-07;\allowbreak{} offline concave scheduling,\allowbreak{} general FCFS threshold representation,\allowbreak{} SEK-SMOD and fork-join asymptotics credited. No present-openness certification.


\subsection*{P0430-A001 | Controlled generators,\allowbreak{} exact finite-state diagnostic,\allowbreak{} coupling,\allowbreak{} renewal reward or capacity comparison}

\entrylabel{Argument or result}

For queue-length vector q and a simultaneous unit increment on selected set A of size k,\allowbreak{} the quadratic potential changes by 2*sum\_\allowbreak{}\{i in A\}q\_\allowbreak{}i+\allowbreak{}k. Selecting the k shortest queues minimizes this one-step drift over all distinct-server choices. Uniform selection instead gives mean increment 2*k*sum\_\allowbreak{}i q\_\allowbreak{}i/\allowbreak{}n+\allowbreak{}k. This establishes a local balancing property,\allowbreak{} but the response of the job is the maximum of its task sojourn times,\allowbreak{} which also depends on residual service and the joint selected-state law.


\entrylabel{Second check}

Expand (q\_\allowbreak{}i+\allowbreak{}1)\textasciicircum{}2-q\_\allowbreak{}i\textasciicircum{}2 separately at each selected coordinate. A pairwise exchange replacing a selected larger queue by an unselected smaller queue decreases the drift,\allowbreak{} proving the minimizer.


\entrylabel{Remaining obstruction}

No stationary maximum-sojourn law or optimality for mean job response follows from the quadratic drift inequality.


Verification:\allowbreak{} author checked partial | Independent review:\allowbreak{} none | Full target solved:\allowbreak{} False


\clearpage

\section*{P0431  Straggler replication in queued limited fork-join systems}

\entrylabel{Primary source}

Harchol-Balter:\allowbreak{} Open problems in queueing theory inspired by datacenter computing;\allowbreak{} Section5;\allowbreak{} PDF16-17


\url{https://www.cs.cmu.edu/~harchol/Papers/QUESTA21.pdf}


Source statement date:\allowbreak{} 2021-01-27. Archived access:\allowbreak{} 2026-09-07T02:\allowbreak{}54:\allowbreak{}43Z.


Historical:\allowbreak{} source\_\allowbreak{}says\_\allowbreak{}open | Current:\allowbreak{} unverified | Result:\allowbreak{} unsolved


\entrylabel{Question}

Analyze and optimize selective straggler replication when a fork-join job completes only after every task completes and replicas also encounter queues.


\entrylabel{Scope}

This adds a job-level maximum and delayed replication to the simpler correlated single-task redundancy questions P0400-P0404.


\entrylabel{Dated literature check}

Focused later-primary-source checks on 2026-09-07;\allowbreak{} offline concave scheduling,\allowbreak{} general FCFS threshold representation,\allowbreak{} SEK-SMOD and fork-join asymptotics credited. No present-openness certification.


\subsection*{P0431-A001 | Controlled generators,\allowbreak{} exact finite-state diagnostic,\allowbreak{} coupling,\allowbreak{} renewal reward or capacity comparison}

\entrylabel{Argument or result}

In the no-queue,\allowbreak{} independent exponential diagnostic,\allowbreak{} two simultaneous replicas each at rate mu complete after Exp(2*mu),\allowbreak{} while expected total server effort is 2/\allowbreak{}(2*mu)=\allowbreak{}1/\allowbreak{}mu,\allowbreak{} equal to one original. With a deterministic overhead c followed by Exp(mu),\allowbreak{} cancellation at first completion instead consumes mean effort 2*c+\allowbreak{}1/\allowbreak{}mu versus c+\allowbreak{}1/\allowbreak{}mu without replication. Thus a service-time shift can reverse the apparent free speedup before any waiting is considered. In a k-task job this replication helps only if the replicated task controls the maximum.


\entrylabel{Second check}

The minimum of independent exponential variables has the sum of their rates. Both replicas execute the common deterministic startup interval before either can complete,\allowbreak{} proving the extra effort c.


\entrylabel{Remaining obstruction}

The diagnostic does not quantify queue feedback,\allowbreak{} correlation,\allowbreak{} delayed starts or the stationary fork-join response distribution.


Verification:\allowbreak{} author checked partial | Independent review:\allowbreak{} none | Full target solved:\allowbreak{} False


\clearpage

\section*{P0432  Optimal known-size multiserver scheduling beyond SRPT}

\entrylabel{Primary source}

Harchol-Balter:\allowbreak{} Open problems in queueing theory inspired by datacenter computing;\allowbreak{} Section7.1.1;\allowbreak{} PDF21;\allowbreak{} later SEK-SMOD paper Theorem4.1


\url{https://www.cs.cmu.edu/~harchol/Papers/QUESTA21.pdf}


Source statement date:\allowbreak{} 2021-01-27. Archived access:\allowbreak{} 2026-09-07T02:\allowbreak{}54:\allowbreak{}43Z.


Historical:\allowbreak{} source\_\allowbreak{}says\_\allowbreak{}open | Current:\allowbreak{} unverified | Result:\allowbreak{} unsolved


\entrylabel{Question}

Characterize nonasymptotic optimal mean-response scheduling and sharp performance bounds in the known-size M/\allowbreak{}G/\allowbreak{}k queue.


\entrylabel{Scope}

The later SEK-SMOD result already outperforms SRPT-k in its stated scope. Existence of any improvement over SRPT-k is no longer admitted as open.


\entrylabel{Dated literature check}

Focused later-primary-source checks on 2026-09-07;\allowbreak{} offline concave scheduling,\allowbreak{} general FCFS threshold representation,\allowbreak{} SEK-SMOD and fork-join asymptotics credited. No present-openness certification.


\subsection*{P0432-A001 | Controlled generators,\allowbreak{} exact finite-state diagnostic,\allowbreak{} coupling,\allowbreak{} renewal reward or capacity comparison}

\entrylabel{Argument or result}

A first-arrival decomposition explains why no-arrival optimality is insufficient. For two policies coupled until a selected state,\allowbreak{} assume their cost difference through recoupling is integrable and let D(t) be its conditional expectation given that the next arrival occurs at time t,\allowbreak{} including their specified continuation rules. Under a Poisson rate lambda,\allowbreak{} the expected difference is integral\_\allowbreak{}0\textasciicircum{}infinity lambda*exp(-lambda*t)*D(t) dt. The sign of D on an intermediate interval can outweigh both early-arrival harm and a favorable no-arrival endpoint. A global comparison must also bound the cost until recoupling;\allowbreak{} comparing only the first completed job misses that term.


\entrylabel{Second check}

Condition on the exponential next-arrival time to derive the integral. Conditional continuations still include all later arrivals;\allowbreak{} replacing them by a no-further-arrivals trajectory would invalidate the identity.


\entrylabel{Remaining obstruction}

No new policy comparison or sharp universal lower bound is proved. The later source resolves the improvement-existence subquestion,\allowbreak{} not full optimal control.


Verification:\allowbreak{} author checked partial | Independent review:\allowbreak{} none | Full target solved:\allowbreak{} False


\clearpage

\section*{P0433  Optimal dispatching to SRPT queues outside heavy traffic}

\entrylabel{Primary source}

Harchol-Balter:\allowbreak{} Open problems in queueing theory inspired by datacenter computing;\allowbreak{} Section7.1.1;\allowbreak{} PDF21-22


\url{https://www.cs.cmu.edu/~harchol/Papers/QUESTA21.pdf}


Source statement date:\allowbreak{} 2021-01-27. Archived access:\allowbreak{} 2026-09-07T02:\allowbreak{}54:\allowbreak{}43Z.


Historical:\allowbreak{} source\_\allowbreak{}says\_\allowbreak{}open | Current:\allowbreak{} unverified | Result:\allowbreak{} unsolved


\entrylabel{Question}

Determine a mean-response-optimal dispatcher to separate SRPT servers at nonasymptotic loads.


\entrylabel{Scope}

Guardrails heavy-traffic optimality is credited. This attempt distinguishes current workload from the incremental delay imposed on existing jobs.


\entrylabel{Dated literature check}

Focused later-primary-source checks on 2026-09-07;\allowbreak{} offline concave scheduling,\allowbreak{} general FCFS threshold representation,\allowbreak{} SEK-SMOD and fork-join asymptotics credited. No present-openness certification.


\subsection*{P0433-A001 | Controlled generators,\allowbreak{} exact finite-state diagnostic,\allowbreak{} coupling,\allowbreak{} renewal reward or capacity comparison}

\entrylabel{Argument or result}

Consider two unit-speed SRPT queues with no future arrivals:\allowbreak{} queue one holds one job of residual size 100;\allowbreak{} queue two holds ten jobs of residual size 0.1 each. A new job of size 0.05 preempts either queue and finishes in 0.05. Sending it to queue one delays one existing completion by 0.05;\allowbreak{} sending it to queue two delays ten existing completions by 0.05 each. Thus the total incremental completion cost is 0.10 versus 0.55,\allowbreak{} although queue two has much less workload. Least workload is not the correct one-step objective for this size-aware setting.


\entrylabel{Second check}

The newcomer is shorter than every existing residual,\allowbreak{} so every existing completion at the chosen queue shifts by exactly 0.05. Adding its own completion time gives the two displayed costs.


\entrylabel{Remaining obstruction}

Future arrivals and the stationary distribution of dispatcher states are omitted. The example is a diagnostic against a proposed greedy criterion,\allowbreak{} not the optimal dispatcher.


Verification:\allowbreak{} author checked partial | Independent review:\allowbreak{} none | Full target solved:\allowbreak{} False


\clearpage

\section*{P0434  Sharp approximation ratio for simple alternatives to Gittins}

\entrylabel{Primary source}

Harchol-Balter:\allowbreak{} Open problems in queueing theory inspired by datacenter computing;\allowbreak{} Section7.1.2;\allowbreak{} PDF22-23


\url{https://www.cs.cmu.edu/~harchol/Papers/QUESTA21.pdf}


Source statement date:\allowbreak{} 2021-01-27. Archived access:\allowbreak{} 2026-09-07T02:\allowbreak{}54:\allowbreak{}43Z.


Historical:\allowbreak{} source\_\allowbreak{}says\_\allowbreak{}open | Current:\allowbreak{} unverified | Result:\allowbreak{} unsolved


\entrylabel{Question}

Improve or identify the sharp worst-case mean-response approximation factor for M-SERPT or other simple age-based policies in M/\allowbreak{}G/\allowbreak{}1.


\entrylabel{Scope}

The source credits a factor-five guarantee. Heavy-traffic factor-two results and the conjectured sharp factor two are distinct from a uniform all-load guarantee.


\entrylabel{Dated literature check}

Focused later-primary-source checks on 2026-09-07;\allowbreak{} offline concave scheduling,\allowbreak{} general FCFS threshold representation,\allowbreak{} SEK-SMOD and fork-join asymptotics credited. No present-openness certification.


\subsection*{P0434-A001 | Controlled generators,\allowbreak{} exact finite-state diagnostic,\allowbreak{} coupling,\allowbreak{} renewal reward or capacity comparison}

\entrylabel{Argument or result}

A two-job calculation isolates the value of probing. Job A has size 0.01 with probability 0.9 and 100 with probability 0.1;\allowbreak{} job B has known size 2. SERPT initially selects B because E[A]=\allowbreak{}10.009. Completing B before A gives expected sum of completion times 14.009. Serving A for 0.01,\allowbreak{} then serving B if A survives,\allowbreak{} gives 0.9*(0.01+\allowbreak{}2.01)+\allowbreak{}0.1*(102+\allowbreak{}2.01)=\allowbreak{}12.219. The Gittins stopping ratio at probe duration 0.01 is 0.01/\allowbreak{}0.9,\allowbreak{} smaller than B rank 2. Any approximation proof based only on expected residual size loses this probing advantage.


\entrylabel{Second check}

In the long-A branch,\allowbreak{} total A completion time is 102 after the intervening two units for B;\allowbreak{} omitting those two units would give a false larger improvement. Both strategies are enumerated by the two possible A sizes.


\entrylabel{Remaining obstruction}

This finite instance does not establish a stationary approximation factor,\allowbreak{} nor distinguish the full worst-case ratios of SERPT and M-SERPT.


Verification:\allowbreak{} author checked partial | Independent review:\allowbreak{} none | Full target solved:\allowbreak{} False


\clearpage

\section*{P0435  Optimal nonasymptotic scheduling with unknown sizes in M/\allowbreak{}G/\allowbreak{}k}

\entrylabel{Primary source}

Harchol-Balter:\allowbreak{} Open problems in queueing theory inspired by datacenter computing;\allowbreak{} Section7.1.2;\allowbreak{} PDF23


\url{https://www.cs.cmu.edu/~harchol/Papers/QUESTA21.pdf}


Source statement date:\allowbreak{} 2021-01-27. Archived access:\allowbreak{} 2026-09-07T02:\allowbreak{}54:\allowbreak{}43Z.


Historical:\allowbreak{} source\_\allowbreak{}says\_\allowbreak{}open | Current:\allowbreak{} unverified | Result:\allowbreak{} unsolved


\entrylabel{Question}

Characterize optimal nonanticipating mean-response policies with multiple servers when sizes are unknown but their law and attained service are observed.


\entrylabel{Scope}

Gittins single-server optimality and multiserver heavy-traffic results are known. No universal finite-load optimality is assumed.


\entrylabel{Dated literature check}

Focused later-primary-source checks on 2026-09-07;\allowbreak{} offline concave scheduling,\allowbreak{} general FCFS threshold representation,\allowbreak{} SEK-SMOD and fork-join asymptotics credited. No present-openness certification.


\subsection*{P0435-A001 | Controlled generators,\allowbreak{} exact finite-state diagnostic,\allowbreak{} coupling,\allowbreak{} renewal reward or capacity comparison}

\entrylabel{Argument or result}

For exponential sizes with a common rate mu,\allowbreak{} every nonidling policy has the same queue-count death rate mu*min(m,\allowbreak{}k),\allowbreak{} so scheduling is immaterial for that special law. With nonexponential sizes the departure intensity becomes sum\_\allowbreak{}i h(a\_\allowbreak{}i) over served jobs,\allowbreak{} and the age vector has policy-dependent drift. Thus queue count ceases to close:\allowbreak{} two states with equal m but different service ages can have different next-departure laws. A proposed extension of the exponential argument must retain this age information and its effect on future capacity.


\entrylabel{Second check}

Independent exponential residuals survive conditioning unchanged,\allowbreak{} deriving the special-case death rate. For a deterministic size,\allowbreak{} an age just below completion and age zero have distinct completion times despite the same queue count,\allowbreak{} checking failure of closure.


\entrylabel{Remaining obstruction}

No optimal age-vector control is constructed. The reduction supplies a known special case and an obstruction to extending its proof.


Verification:\allowbreak{} author checked partial | Independent review:\allowbreak{} none | Full target solved:\allowbreak{} False


\clearpage

\section*{P0436  Optimal finite-deadline success scheduling}

\entrylabel{Primary source}

Harchol-Balter:\allowbreak{} Open problems in queueing theory inspired by datacenter computing;\allowbreak{} Section7.2;\allowbreak{} PDF24-25


\url{https://www.cs.cmu.edu/~harchol/Papers/QUESTA21.pdf}


Source statement date:\allowbreak{} 2021-01-27. Archived access:\allowbreak{} 2026-09-07T02:\allowbreak{}54:\allowbreak{}43Z.


Historical:\allowbreak{} source\_\allowbreak{}says\_\allowbreak{}open | Current:\allowbreak{} unverified | Result:\allowbreak{} unsolved


\entrylabel{Question}

Minimize the fraction of known-size Poisson arrivals missing a common relative deadline D,\allowbreak{} with an explicit convention that dropped jobs count as failures.


\entrylabel{Scope}

This is a finite-threshold objective;\allowbreak{} large-time tail asymptotics alone do not solve it. The source discusses preemptive and dropping policies.


\entrylabel{Dated literature check}

Focused later-primary-source checks on 2026-09-07;\allowbreak{} offline concave scheduling,\allowbreak{} general FCFS threshold representation,\allowbreak{} SEK-SMOD and fork-join asymptotics credited. No present-openness certification.


\subsection*{P0436-A001 | Controlled generators,\allowbreak{} exact finite-state diagnostic,\allowbreak{} coupling,\allowbreak{} renewal reward or capacity comparison}

\entrylabel{Argument or result}

A deadline-success value must carry remaining slack as well as remaining work. For a job arriving at a with size x,\allowbreak{} slack at t is a+\allowbreak{}D-t-x\_\allowbreak{}remaining. Two equal-work jobs can have opposite feasibility depending on arrival age. Moreover an admission rule using only current feasibility ignores the opportunity cost of accepted work:\allowbreak{} a size close to D may be feasible now but block many later short jobs. In a finite-horizon control equation the immediate reward is a unit success at completion before the deadline,\allowbreak{} rather than the running holding cost used for mean response.


\entrylabel{Second check}

Check two residual-size 0.2 jobs with slacks 0.1 and -0.1:\allowbreak{} the latter cannot succeed even with immediate exclusive service. This separates a feasibility state from a workload-only state without assuming an optimal policy.


\entrylabel{Remaining obstruction}

The stationary success-maximizing policy,\allowbreak{} admissible preemption/\allowbreak{}drop rules and its verification remain unresolved.


Verification:\allowbreak{} author checked partial | Independent review:\allowbreak{} none | Full target solved:\allowbreak{} False


\clearpage

\section*{P0437  Dropping probability and optimality gap of DTLH}

\entrylabel{Primary source}

Harchol-Balter:\allowbreak{} Open problems in queueing theory inspired by datacenter computing;\allowbreak{} Section7.2;\allowbreak{} PDF24


\url{https://www.cs.cmu.edu/~harchol/Papers/QUESTA21.pdf}


Source statement date:\allowbreak{} 2021-01-27. Archived access:\allowbreak{} 2026-09-07T02:\allowbreak{}54:\allowbreak{}43Z.


Historical:\allowbreak{} source\_\allowbreak{}says\_\allowbreak{}open | Current:\allowbreak{} unverified | Result:\allowbreak{} unsolved


\entrylabel{Question}

Determine the stationary loss probability and performance gap of Drop-The-Larger-Hopeless in the source model.


\entrylabel{Scope}

A precise implementation must specify which eligible larger waiting job is removed. Loss probabilities refer to all arrivals,\allowbreak{} including removed old jobs.


\entrylabel{Dated literature check}

Focused later-primary-source checks on 2026-09-07;\allowbreak{} offline concave scheduling,\allowbreak{} general FCFS threshold representation,\allowbreak{} SEK-SMOD and fork-join asymptotics credited. No present-openness certification.


\subsection*{P0437-A001 | Controlled generators,\allowbreak{} exact finite-state diagnostic,\allowbreak{} coupling,\allowbreak{} renewal reward or capacity comparison}

\entrylabel{Argument or result}

Total workload alone does not determine the next admission outcome. Let the in-service residual be 0.1,\allowbreak{} the waiting workload be 0.8,\allowbreak{} D=\allowbreak{}1,\allowbreak{} and the new job have size 0.5. If the waiting line is one size-0.8 job,\allowbreak{} removing it makes the newcomer feasible. If the waiting line is two size-0.4 jobs,\allowbreak{} neither is larger than the newcomer,\allowbreak{} so the newcomer is dropped. Both states have total workload 0.9. A Markov analysis needs the ordered waiting-size/\allowbreak{}slack state,\allowbreak{} not just a reflected scalar workload.


\entrylabel{Second check}

Before replacement the prospective completion delay is 1.4. Removing 0.8 leaves delay 0.6. In the two-job state no waiting size exceeds 0.5. The in-service job is explicitly excluded from the removal candidates.


\entrylabel{Remaining obstruction}

No stationary ordered-state law or approximation ratio is obtained;\allowbreak{} the example diagnoses insufficient state information.


Verification:\allowbreak{} author checked partial | Independent review:\allowbreak{} none | Full target solved:\allowbreak{} False


\clearpage

\section*{P0438  Feasibility of class-specific response quantiles}

\entrylabel{Primary source}

Harchol-Balter:\allowbreak{} Open problems in queueing theory inspired by datacenter computing;\allowbreak{} Section7.3;\allowbreak{} PDF25-26


\url{https://www.cs.cmu.edu/~harchol/Papers/QUESTA21.pdf}


Source statement date:\allowbreak{} 2021-01-27. Archived access:\allowbreak{} 2026-09-07T02:\allowbreak{}54:\allowbreak{}43Z.


Historical:\allowbreak{} source\_\allowbreak{}says\_\allowbreak{}open | Current:\allowbreak{} unverified | Result:\allowbreak{} unsolved


\entrylabel{Question}

Characterize scheduling feasibility for several classes with different response deadlines and allowed miss fractions.


\entrylabel{Scope}

The source gives distinct class SLOs. A class-blind aggregate quantile cannot substitute for all class-specific constraints.


\entrylabel{Dated literature check}

Focused later-primary-source checks on 2026-09-07;\allowbreak{} offline concave scheduling,\allowbreak{} general FCFS threshold representation,\allowbreak{} SEK-SMOD and fork-join asymptotics credited. No present-openness certification.


\subsection*{P0438-A001 | Controlled generators,\allowbreak{} exact finite-state diagnostic,\allowbreak{} coupling,\allowbreak{} renewal reward or capacity comparison}

\entrylabel{Argument or result}

For class i with nonparallel intrinsic service S\_\allowbreak{}i and deadline D\_\allowbreak{}i,\allowbreak{} every policy obeys P(T\_\allowbreak{}i>D\_\allowbreak{}i)>=\allowbreak{}P(S\_\allowbreak{}i>D\_\allowbreak{}i). Also an aggregate miss rate is a weighted average sum\_\allowbreak{}i p\_\allowbreak{}i*m\_\allowbreak{}i and does not bound each m\_\allowbreak{}i:\allowbreak{} if a rare class has probability 0.005 and all its jobs miss while the rest succeed,\allowbreak{} aggregate misses are 0.005 but that class has miss rate one. Thus aggregate feasibility and even workload stability leave separate unavoidable class constraints.


\entrylabel{Second check}

Response time contains service time pathwise. The rare-class example follows by total probability and does not rely on simulated queue data.


\entrylabel{Remaining obstruction}

These necessary conditions are not sufficient;\allowbreak{} no policy is proved to attain the feasible vector of quantiles.


Verification:\allowbreak{} author checked partial | Independent review:\allowbreak{} none | Full target solved:\allowbreak{} False


\clearpage

\section*{P0439  Server provisioning for a finite response-time SLO}

\entrylabel{Primary source}

Harchol-Balter:\allowbreak{} Open problems in queueing theory inspired by datacenter computing;\allowbreak{} Section7.3;\allowbreak{} PDF26


\url{https://www.cs.cmu.edu/~harchol/Papers/QUESTA21.pdf}


Source statement date:\allowbreak{} 2021-01-27. Archived access:\allowbreak{} 2026-09-07T02:\allowbreak{}54:\allowbreak{}43Z.


Historical:\allowbreak{} source\_\allowbreak{}says\_\allowbreak{}open | Current:\allowbreak{} unverified | Result:\allowbreak{} unsolved


\entrylabel{Question}

Find the minimum server count and a feasible scheduling policy that meet a prescribed response-time quantile under a specified stochastic workload.


\entrylabel{Scope}

The exponential FCFS model is a known benchmark only. General service laws and policy choice are part of the source question.


\entrylabel{Dated literature check}

Focused later-primary-source checks on 2026-09-07;\allowbreak{} offline concave scheduling,\allowbreak{} general FCFS threshold representation,\allowbreak{} SEK-SMOD and fork-join asymptotics credited. No present-openness certification.


\subsection*{P0439-A001 | Controlled generators,\allowbreak{} exact finite-state diagnostic,\allowbreak{} coupling,\allowbreak{} renewal reward or capacity comparison}

\entrylabel{Argument or result}

For the stable M/\allowbreak{}M/\allowbreak{}k FCFS benchmark with total arrival rate Lambda and service rate mu per server,\allowbreak{} let C be Erlang-C and a=\allowbreak{}k*mu-Lambda. Waiting time is zero with probability 1-C and otherwise Exp(a),\allowbreak{} independent of the tagged Exp(mu) service. For a!=\allowbreak{}mu this gives P(T>t)=\allowbreak{}(1-C)*exp(-mu*t)+\allowbreak{}C*[mu*exp(-a*t)-a*exp(-mu*t)]/\allowbreak{}(mu-a). At a=\allowbreak{}mu the bracket is exp(-mu*t)*(1+\allowbreak{}mu*t). This gives a reproducible benchmark for each k;\allowbreak{} simply using the waiting-time tail as the response-time tail would omit service.


\entrylabel{Second check}

Convolve the two exponential distributions,\allowbreak{} and take the continuous limit a->mu to obtain the Erlang-two tail. At t=\allowbreak{}0 the expression equals one,\allowbreak{} and at C=\allowbreak{}0 it reduces to the service tail.


\entrylabel{Remaining obstruction}

The benchmark neither proves the minimum k over all policies nor transfers to M/\allowbreak{}G/\allowbreak{}k. No new general staffing rule is established.


Verification:\allowbreak{} author checked partial | Independent review:\allowbreak{} none | Full target solved:\allowbreak{} False


\clearpage

\section*{P0440  Waiting-time tails with deterministic multiserver setup}

\entrylabel{Primary source}

Williams Wang Harchol-Balter:\allowbreak{} An Upper Bound on the M M k Queue With Deterministic Setup Times;\allowbreak{} Section8;\allowbreak{} PDF39


\url{https://arxiv.org/pdf/2512.16854}


Source statement date:\allowbreak{} 2025-12-18. Archived access:\allowbreak{} 2026-09-07T06:\allowbreak{}25:\allowbreak{}33Z.


Historical:\allowbreak{} source\_\allowbreak{}says\_\allowbreak{}open | Current:\allowbreak{} unverified | Result:\allowbreak{} unsolved


\entrylabel{Question}

Obtain higher-moment or tail bounds for the M/\allowbreak{}M/\allowbreak{}k system with deterministic simultaneous setup times and the paper specified cancellation policy.


\entrylabel{Scope}

The source already supplies multiplicatively matching mean bounds in its parameter regime;\allowbreak{} that mean result is not recounted as open.


\entrylabel{Dated literature check}

Focused later-primary-source checks on 2026-09-07;\allowbreak{} offline concave scheduling,\allowbreak{} general FCFS threshold representation,\allowbreak{} SEK-SMOD and fork-join asymptotics credited. No present-openness certification.


\subsection*{P0440-A001 | Controlled generators,\allowbreak{} exact finite-state diagnostic,\allowbreak{} coupling,\allowbreak{} renewal reward or capacity comparison}

\entrylabel{Argument or result}

For a state with m jobs and z busy servers,\allowbreak{} the count part of the generator applied to exp(theta*m) is [Lambda*(exp(theta)-1)+\allowbreak{}mu*z*(exp(-theta)-1)]*exp(theta*m). It is negative when z is sufficiently large,\allowbreak{} but positive when every server is in setup. Thus a one-step exponential drift estimate cannot hold uniformly on large-m states using m alone. A useful Lyapunov function must include residual setup times,\allowbreak{} or the proof must use a time interval long enough for setup completion while controlling cancellations.


\entrylabel{Second check}

Each arrival multiplies the test function by exp(theta);\allowbreak{} each completion multiplies it by exp(-theta). Deterministic setup progress does not change m immediately,\allowbreak{} which is precisely the missing generator term in a count-only test.


\entrylabel{Remaining obstruction}

No tail exponent or higher-moment bound with verified constants is obtained. Mean bounds alone do not certify the tail.


Verification:\allowbreak{} author checked partial | Independent review:\allowbreak{} none | Full target solved:\allowbreak{} False


\clearpage

\section*{P0441  Deterministic DelayedOff waiting-power tradeoff}

\entrylabel{Primary source}

Williams Wang Harchol-Balter:\allowbreak{} An Upper Bound on the M M k Queue With Deterministic Setup Times;\allowbreak{} Section8;\allowbreak{} PDF39-40


\url{https://arxiv.org/pdf/2512.16854}


Source statement date:\allowbreak{} 2025-12-18. Archived access:\allowbreak{} 2026-09-07T06:\allowbreak{}25:\allowbreak{}33Z.


Historical:\allowbreak{} source\_\allowbreak{}says\_\allowbreak{}open | Current:\allowbreak{} unverified | Result:\allowbreak{} unsolved


\entrylabel{Question}

Analyze waiting and power for deterministic setup when idle servers remain on for a grace interval gamma before switching off.


\entrylabel{Scope}

The source distinguishes these policies from its instantaneous-off model and credits earlier exponential-setup studies.


\entrylabel{Dated literature check}

Focused later-primary-source checks on 2026-09-07;\allowbreak{} offline concave scheduling,\allowbreak{} general FCFS threshold representation,\allowbreak{} SEK-SMOD and fork-join asymptotics credited. No present-openness certification.


\subsection*{P0441-A001 | Controlled generators,\allowbreak{} exact finite-state diagnostic,\allowbreak{} coupling,\allowbreak{} renewal reward or capacity comparison}

\entrylabel{Argument or result}

A single idle-gap diagnostic is exact:\allowbreak{} for next-arrival gap G\textasciitilde{}Exp(Lambda),\allowbreak{} the chance the server becomes cold is exp(-Lambda*gamma),\allowbreak{} and expected on-idle duration is E[min(G,\allowbreak{}gamma)]=\allowbreak{}(1-exp(-Lambda*gamma))/\allowbreak{}Lambda. If a cold arrival triggers deterministic setup beta,\allowbreak{} that arrival incurs beta*exp(-Lambda*gamma) expected setup delay. The local marginal tradeoff follows by differentiation,\allowbreak{} but a full renewal cycle also includes later arrivals during setup,\allowbreak{} active service and possible setup cancellation.


\entrylabel{Second check}

Integrate P(G>t)=\allowbreak{}exp(-Lambda*t) over [0,\allowbreak{}gamma]. Separately condition on G>gamma for the cold-arrival delay. Both formulas have correct gamma=\allowbreak{}0 and gamma->infinity limits.


\entrylabel{Remaining obstruction}

The idle-gap calculation is not a stationary many-server mean or a globally optimal grace period.


Verification:\allowbreak{} author checked partial | Independent review:\allowbreak{} none | Full target solved:\allowbreak{} False


\clearpage

\section*{P0442  Optimal waiting-power Pareto frontier with deterministic setup}

\entrylabel{Primary source}

Williams Wang Harchol-Balter:\allowbreak{} An Upper Bound on the M M k Queue With Deterministic Setup Times;\allowbreak{} Section8;\allowbreak{} PDF40


\url{https://arxiv.org/pdf/2512.16854}


Source statement date:\allowbreak{} 2025-12-18. Archived access:\allowbreak{} 2026-09-07T06:\allowbreak{}25:\allowbreak{}33Z.


Historical:\allowbreak{} source\_\allowbreak{}says\_\allowbreak{}open | Current:\allowbreak{} unverified | Result:\allowbreak{} unsolved


\entrylabel{Question}

Characterize the attainable and optimal tradeoff between mean waiting time and average power over admissible multiserver on-off policies.


\entrylabel{Scope}

This ranges over all policies,\allowbreak{} whereas P0441 analyzes one specified family. Always-on is the source-known minimum-wait endpoint.


\entrylabel{Dated literature check}

Focused later-primary-source checks on 2026-09-07;\allowbreak{} offline concave scheduling,\allowbreak{} general FCFS threshold representation,\allowbreak{} SEK-SMOD and fork-join asymptotics credited. No present-openness certification.


\subsection*{P0442-A001 | Controlled generators,\allowbreak{} exact finite-state diagnostic,\allowbreak{} coupling,\allowbreak{} renewal reward or capacity comparison}

\entrylabel{Argument or result}

Let a stationary occupation measure record active,\allowbreak{} idle-on,\allowbreak{} setting-up and off servers. Average power is a linear functional of that measure,\allowbreak{} while mean waiting equals E[Q]/\allowbreak{}Lambda by Little law when finite. Minimizing E[Q]/\allowbreak{}Lambda+\allowbreak{}eta*E[power] for eta>=\allowbreak{}0 produces supported frontier points. If policies share a common regenerative empty/\allowbreak{}reset state and have finite cycle moments,\allowbreak{} randomizing between them at those epochs produces convex mixtures of cycle reward rates. The mixing weight is cycle-length biased;\allowbreak{} averaging policy labels without that correction gives the wrong frontier.


\entrylabel{Second check}

If a cycle uses policy A with probability p,\allowbreak{} its reward rate is [p*E[R\_\allowbreak{}A]+\allowbreak{}(1-p)*E[R\_\allowbreak{}B]]/\allowbreak{}[p*E[L\_\allowbreak{}A]+\allowbreak{}(1-p)*E[L\_\allowbreak{}B]],\allowbreak{} directly from renewal reward. The coefficient of A rate is p*E[L\_\allowbreak{}A] divided by the same denominator.


\entrylabel{Remaining obstruction}

No feasible occupation-measure characterization,\allowbreak{} optimal control solution,\allowbreak{} or proof that DelayedOff spans the frontier is supplied.


Verification:\allowbreak{} author checked partial | Independent review:\allowbreak{} none | Full target solved:\allowbreak{} False


\clearpage

\section*{P0443  Closed-form performance bounds for general setup distributions}

\entrylabel{Primary source}

Williams Wang Harchol-Balter:\allowbreak{} An Upper Bound on the M M k Queue With Deterministic Setup Times;\allowbreak{} Section8;\allowbreak{} PDF40


\url{https://arxiv.org/pdf/2512.16854}


Source statement date:\allowbreak{} 2025-12-18. Archived access:\allowbreak{} 2026-09-07T06:\allowbreak{}25:\allowbreak{}33Z.


Historical:\allowbreak{} source\_\allowbreak{}says\_\allowbreak{}open | Current:\allowbreak{} unverified | Result:\allowbreak{} unsolved


\entrylabel{Question}

Extend closed-form multiserver setup bounds to exponential,\allowbreak{} phase-type or two-point setup times under explicitly stated restart/\allowbreak{}cancellation semantics.


\entrylabel{Scope}

Exact numerical matrix formulas already exist in some models;\allowbreak{} a simple parameter-explicit bound and its error are the remaining target.


\entrylabel{Dated literature check}

Focused later-primary-source checks on 2026-09-07;\allowbreak{} offline concave scheduling,\allowbreak{} general FCFS threshold representation,\allowbreak{} SEK-SMOD and fork-join asymptotics credited. No present-openness certification.


\subsection*{P0443-A001 | Controlled generators,\allowbreak{} exact finite-state diagnostic,\allowbreak{} coupling,\allowbreak{} renewal reward or capacity comparison}

\entrylabel{Argument or result}

The mean setup time alone loses a basic renewal quantity. If setup durations B are sampled independently with finite second moment and a time observer samples a stationary setup interval,\allowbreak{} the residual mean is E[B\textasciicircum{}2]/\allowbreak{}(2*E[B]). Deterministic B=\allowbreak{}beta gives beta/\allowbreak{}2,\allowbreak{} while Exp(mean beta) gives beta. But an arrival that initiates a setup samples the original B,\allowbreak{} not this length-biased residual. Therefore neither substituting the same mean nor blindly inserting the residual-life formula yields a correct multiserver approximation.


\entrylabel{Second check}

Renewal reward integrates the residual triangular area B\textasciicircum{}2/\allowbreak{}2 in each interval and divides by E[B]. The initiation observer is an independent endpoint observer,\allowbreak{} distinguishing its sampling law.


\entrylabel{Remaining obstruction}

The inspection-paradox calculation is a diagnostic only. Actual queue-coupled setup intervals and cancellation can violate its independent renewal assumptions.


Verification:\allowbreak{} author checked partial | Independent review:\allowbreak{} none | Full target solved:\allowbreak{} False


\clearpage

\section*{P0444  Computational complexity of the multiserver-job stability threshold}

\entrylabel{Primary source}

The Multiserver Job Stochastic Recurrence Equation for Cloud Computing Performance Evaluation;\allowbreak{} Section5;\allowbreak{} PDF11


\url{https://arxiv.org/pdf/2601.20653}


Source statement date:\allowbreak{} 2026-01-28. Archived access:\allowbreak{} 2026-09-07T06:\allowbreak{}25:\allowbreak{}16Z.


Historical:\allowbreak{} source\_\allowbreak{}says\_\allowbreak{}open | Current:\allowbreak{} unverified | Result:\allowbreak{} unsolved


\entrylabel{Question}

Determine the complexity and certified approximation accuracy of the saturated pile growth rate gamma that defines the FCFS multiserver-job stability threshold.


\entrylabel{Scope}

The source already establishes the threshold representation lambda\_\allowbreak{}c=\allowbreak{}1/\allowbreak{}gamma. A stopped simulation with small successive change is not automatically a certified estimate.


\entrylabel{Dated literature check}

Focused later-primary-source checks on 2026-09-07;\allowbreak{} offline concave scheduling,\allowbreak{} general FCFS threshold representation,\allowbreak{} SEK-SMOD and fork-join asymptotics credited. No present-openness certification.


\subsection*{P0444-A001 | Controlled generators,\allowbreak{} exact finite-state diagnostic,\allowbreak{} coupling,\allowbreak{} renewal reward or capacity comparison}

\entrylabel{Argument or result}

For the special case where every job needs b servers,\allowbreak{} put c=\allowbreak{}floor(s/\allowbreak{}b). The saturated system is a c-machine list schedule. Its maximum completion time after n jobs is at least total service/\allowbreak{}c and at most total service/\allowbreak{}c+\allowbreak{}(1-1/\allowbreak{}c)*max\_\allowbreak{}\{j<=\allowbreak{}n\}S\_\allowbreak{}j. For iid finite-mean service,\allowbreak{} total service/\allowbreak{}n->E[S] and max S\_\allowbreak{}j/\allowbreak{}n->0 almost surely,\allowbreak{} so gamma=\allowbreak{}E[S]/\allowbreak{}c. Heterogeneous demands break this scalar packing reduction. In the general estimator,\allowbreak{} a small difference between consecutive sample means does not bound the distance from their limiting value.


\entrylabel{Second check}

The list-scheduling upper bound is obtained by looking just before the job that finishes last starts:\allowbreak{} all c machine loads are at least its start time,\allowbreak{} and its own service accounts for the remaining gap. Divide by n and use the finite-mean iid maximal-term law.


\entrylabel{Remaining obstruction}

No general finite-sample certificate or computational complexity bound is derived;\allowbreak{} the equal-demand case is a credited classical scheduling reduction.


Verification:\allowbreak{} author checked partial | Independent review:\allowbreak{} none | Full target solved:\allowbreak{} False


\clearpage

\section*{P0445  Monotone workload recurrences beyond FCFS multiserver jobs}

\entrylabel{Primary source}

The Multiserver Job Stochastic Recurrence Equation for Cloud Computing Performance Evaluation;\allowbreak{} Section10;\allowbreak{} PDF22


\url{https://arxiv.org/pdf/2601.20653}


Source statement date:\allowbreak{} 2026-01-28. Archived access:\allowbreak{} 2026-09-07T06:\allowbreak{}25:\allowbreak{}16Z.


Historical:\allowbreak{} source\_\allowbreak{}says\_\allowbreak{}open | Current:\allowbreak{} unverified | Result:\allowbreak{} unsolved


\entrylabel{Question}

Identify scheduling disciplines beyond FCFS whose multiserver-job workload maps preserve the monotonicity needed by the source framework.


\entrylabel{Scope}

The paper proves monotonicity for its FCFS map. A different queue discipline needs its own state and order,\allowbreak{} not merely the same name for workload.


\entrylabel{Dated literature check}

Focused later-primary-source checks on 2026-09-07;\allowbreak{} offline concave scheduling,\allowbreak{} general FCFS threshold representation,\allowbreak{} SEK-SMOD and fork-join asymptotics credited. No present-openness certification.


\subsection*{P0445-A001 | Controlled generators,\allowbreak{} exact finite-state diagnostic,\allowbreak{} coupling,\allowbreak{} renewal reward or capacity comparison}

\entrylabel{Argument or result}

Compositions of coordinatewise isotone primitive maps remain isotone:\allowbreak{} if x<=\allowbreak{}y,\allowbreak{} first F(x)<=\allowbreak{}F(y),\allowbreak{} then G(F(x))<=\allowbreak{}G(F(y)). Consequently,\allowbreak{} a policy using an exogenously fixed sequence of the same monotone primitives inherits the comparison. An adaptive scheduler instead selects F\_\allowbreak{}\{a(x)\} at state x and F\_\allowbreak{}\{a(y)\} at y;\allowbreak{} isotonicity of each separate F\_\allowbreak{}a gives no comparison between those two values. In one dimension,\allowbreak{} F0(x)=\allowbreak{}x and F1(x)=\allowbreak{}x+\allowbreak{}10 are both increasing,\allowbreak{} but selecting F1 below threshold one and F0 above it makes the chosen map decrease at the threshold.


\entrylabel{Second check}

The scalar pair x=\allowbreak{}0.9,\allowbreak{}y=\allowbreak{}1.1 yields outputs 10.9 and 1.1. This disproves the proposed generic inference from branchwise isotonicity,\allowbreak{} not monotonicity of any named queue discipline.


\entrylabel{Remaining obstruction}

No new non-FCFS discipline is proved monotone. A queue-specific comparison requires controlling policy switches and the admissible arrival order.


Verification:\allowbreak{} author checked partial | Independent review:\allowbreak{} none | Full target solved:\allowbreak{} False


\clearpage

\section*{P0446  Perfect multiserver-job samples with a stability-matching upper bound}

\entrylabel{Primary source}

The Multiserver Job Stochastic Recurrence Equation for Cloud Computing Performance Evaluation;\allowbreak{} Section10;\allowbreak{} PDF22


\url{https://arxiv.org/pdf/2601.20653}


Source statement date:\allowbreak{} 2026-01-28. Archived access:\allowbreak{} 2026-09-07T06:\allowbreak{}25:\allowbreak{}16Z.


Historical:\allowbreak{} source\_\allowbreak{}says\_\allowbreak{}open | Current:\allowbreak{} unverified | Result:\allowbreak{} unsolved


\entrylabel{Question}

Construct computable stationary upper bounds with the same stability region as the FCFS multiserver-job queue,\allowbreak{} enabling perfect rather than sub-perfect samples.


\entrylabel{Scope}

The source already obtains monotone lower samples that eventually couple;\allowbreak{} detecting that coupling is a separate problem.


\entrylabel{Dated literature check}

Focused later-primary-source checks on 2026-09-07;\allowbreak{} offline concave scheduling,\allowbreak{} general FCFS threshold representation,\allowbreak{} SEK-SMOD and fork-join asymptotics credited. No present-openness certification.


\subsection*{P0446-A001 | Controlled generators,\allowbreak{} exact finite-state diagnostic,\allowbreak{} coupling,\allowbreak{} renewal reward or capacity comparison}

\entrylabel{Argument or result}

Serializing every job yields a simple potential dominating workload but may destroy the needed stability region. With two servers,\allowbreak{} single-server-demand iid service of mean one,\allowbreak{} and arrival rate 1.5,\allowbreak{} the actual M/\allowbreak{}G/\allowbreak{}2 has load 0.75,\allowbreak{} while an all-jobs-serialized M/\allowbreak{}G/\allowbreak{}1 has load 1.5 and no finite stationary workload. Thus this upper-bound construction cannot support perfect sampling throughout the true stable region. Nor does observing a lower sequence remain constant over a finite number of lookback extensions certify that it will never jump later.


\entrylabel{Second check}

The respective work capacities are two and one units per unit time. A deterministic nondecreasing sequence that stays zero for N steps and then becomes one independently demonstrates the logical failure of finite plateau certification.


\entrylabel{Remaining obstruction}

No stability-matching dominating chain or detectable coalescence criterion is constructed.


Verification:\allowbreak{} author checked partial | Independent review:\allowbreak{} none | Full target solved:\allowbreak{} False


\clearpage

\section*{P0447  Finite-size response characterization for classical fork-join queues}

\entrylabel{Primary source}

Harchol-Balter:\allowbreak{} Open problems in queueing theory inspired by datacenter computing;\allowbreak{} Section5;\allowbreak{} PDF15;\allowbreak{} later fork-join extreme-value theorem


\url{https://www.cs.cmu.edu/~harchol/Papers/QUESTA21.pdf}


Source statement date:\allowbreak{} 2021-01-27. Archived access:\allowbreak{} 2026-09-07T02:\allowbreak{}54:\allowbreak{}43Z.


Historical:\allowbreak{} source\_\allowbreak{}says\_\allowbreak{}open | Current:\allowbreak{} unverified | Result:\allowbreak{} unsolved


\entrylabel{Question}

Obtain a tractable sharp finite-n response characterization for classical fork-join queues beyond the two-server special case.


\entrylabel{Scope}

Later normal-limit results identify a logarithmic centering constant in their light-tailed regime. This entry retains only the finite-size joint-law question,\allowbreak{} not the older unresolved-constant claim.


\entrylabel{Dated literature check}

Focused later-primary-source checks on 2026-09-07;\allowbreak{} offline concave scheduling,\allowbreak{} general FCFS threshold representation,\allowbreak{} SEK-SMOD and fork-join asymptotics credited. No present-openness certification.


\subsection*{P0447-A001 | Controlled generators,\allowbreak{} exact finite-state diagnostic,\allowbreak{} coupling,\allowbreak{} renewal reward or capacity comparison}

\entrylabel{Argument or result}

In a stable exponential-service fork-join model with common Poisson arrivals,\allowbreak{} independent task services alone give E[T\_\allowbreak{}job]>=\allowbreak{}H\_\allowbreak{}n/\allowbreak{}mu,\allowbreak{} where H\_\allowbreak{}n is the nth harmonic sum,\allowbreak{} because waiting times are nonnegative. An independence approximation applied to stationary task sojourns would give H\_\allowbreak{}n/\allowbreak{}(mu-Lambda),\allowbreak{} but shared arrivals make those sojourns dependent. Computing each marginal M/\allowbreak{}M/\allowbreak{}1 law therefore does not identify the maximum. The gap between these two scales also warns against replacing a finite-n maximum by its asymptotic centering without a remainder bound.


\entrylabel{Second check}

For iid Exp(mu) service requirements,\allowbreak{} integrate 1-(1-exp(-mu*t))\textasciicircum{}n to obtain H\_\allowbreak{}n/\allowbreak{}mu. The actual lower bound is pathwise since each task sojourn exceeds its own service.


\entrylabel{Remaining obstruction}

No finite-n joint transform or sharp error bound is obtained. The independent-sojourn expression is explicitly an approximation unless a separate association comparison is invoked.


Verification:\allowbreak{} author checked partial | Independent review:\allowbreak{} none | Full target solved:\allowbreak{} False


\clearpage

\section*{P0450  Conditional generation limits for the Erlang weighted tree}

\entrylabel{Primary source}

Moharrami,\allowbreak{} Subramanian,\allowbreak{} Liu and Sundaresan:\allowbreak{} The Erlang Weighted Tree,\allowbreak{} A New Branching Process (v5,\allowbreak{} 27 May 2023);\allowbreak{} Section4,\allowbreak{} question1;\allowbreak{} PDF15-16


\url{https://arxiv.org/pdf/2002.03993}


Source statement date:\allowbreak{} 2023-05-27. Archived access:\allowbreak{} 2026-09-07T07:\allowbreak{}23:\allowbreak{}15Z.


Historical:\allowbreak{} source\_\allowbreak{}says\_\allowbreak{}open | Current:\allowbreak{} unverified | Result:\allowbreak{} unsolved


\entrylabel{Question}

Determine survival-conditioned generation laws and their proper normalization for critical and subcritical Erlang weighted trees,\allowbreak{} including conditioning on later survival.


\entrylabel{Scope}

The source groups these conditional-limit versions. Its displayed notation mixes a counting measure with its type space;\allowbreak{} the target must distinguish total population from the finite point-measure law. Critical unscaled tightness cannot be assumed.


\entrylabel{Dated literature check}

2026-09-07 bounded primary searches for EWT,\allowbreak{} La-Kabkab and bilateral preference degree results;\allowbreak{} 2025 connectivity theorem credited. No current-open or novelty certification.


\subsection*{P0450-A001 | Exact conditioning,\allowbreak{} spectral integral estimate,\allowbreak{} deterministic graph comparison or order-statistic asymptotics}

\entrylabel{Argument or result}

An exact classical branching diagnostic shows why the critical normalization matters. For geometric offspring with pgf f(s)=\allowbreak{}1/\allowbreak{}(2-s),\allowbreak{} iteration gives f\_\allowbreak{}l(s)=\allowbreak{}(l-(l-1)*s)/\allowbreak{}(l+\allowbreak{}1-l*s). Thus survival has probability 1/\allowbreak{}(l+\allowbreak{}1),\allowbreak{} and the conditional positive population is geometric with success 1/\allowbreak{}(l+\allowbreak{}1). Its mass on \{1,\allowbreak{}...,\allowbreak{}M\} tends to zero for each fixed M,\allowbreak{} while Z\_\allowbreak{}l/\allowbreak{}l conditioned on survival converges to Exp(1). This is a single-type diagnostic,\allowbreak{} not the EWT. For EWT a possible route uses the positive right eigenfunction to form a size-biased spine,\allowbreak{} but convergence under survival conditioning requires a separate comparison to that bias.


\entrylabel{Second check}

Direct substitution verifies f(f\_\allowbreak{}l(s))=\allowbreak{}f\_\allowbreak{}(l+\allowbreak{}1)(s);\allowbreak{} subtract f\_\allowbreak{}l(0),\allowbreak{} divide by survival probability and expand the geometric series. The conditional tail at floor(l*x) is (l/\allowbreak{}(l+\allowbreak{}1))\textasciicircum{}floor(l*x),\allowbreak{} tending to exp(-x).


\entrylabel{Remaining obstruction}

No EWT Yaglom or Q-process theorem is proved. A size-biased spine law is not automatically the law conditioned on nonextinction.


Verification:\allowbreak{} author checked partial | Independent review:\allowbreak{} none | Full target solved:\allowbreak{} False


\clearpage

\section*{P0451  Unimodularity and reversibility of the Erlang branching spine}

\entrylabel{Primary source}

Moharrami,\allowbreak{} Subramanian,\allowbreak{} Liu and Sundaresan:\allowbreak{} The Erlang Weighted Tree,\allowbreak{} A New Branching Process (v5,\allowbreak{} 27 May 2023);\allowbreak{} Section4,\allowbreak{} question2;\allowbreak{} PDF16


\url{https://arxiv.org/pdf/2002.03993}


Source statement date:\allowbreak{} 2023-05-27. Archived access:\allowbreak{} 2026-09-07T07:\allowbreak{}23:\allowbreak{}15Z.


Historical:\allowbreak{} source\_\allowbreak{}says\_\allowbreak{}open | Current:\allowbreak{} unverified | Result:\allowbreak{} unsolved


\entrylabel{Question}

Characterize the relationship between unimodularity of the marked branching tree and reversibility of its source-defined continuous-state spine kernel.


\entrylabel{Scope}

The reversibility of this particular kernel is already established in the source. The open target is the general relationship and correct change of root,\allowbreak{} not another proof of detailed balance.


\entrylabel{Dated literature check}

2026-09-07 bounded primary searches for EWT,\allowbreak{} La-Kabkab and bilateral preference degree results;\allowbreak{} 2025 connectivity theorem credited. No current-open or novelty certification.


\subsection*{P0451-A001 | Exact conditioning,\allowbreak{} spectral integral estimate,\allowbreak{} deterministic graph comparison or order-statistic asymptotics}

\entrylabel{Argument or result}

For the displayed kernel p(x,\allowbreak{}y)=\allowbreak{}min(x,\allowbreak{}y)*g2(y)*f0(y)/\allowbreak{}(beta0*f0(x)),\allowbreak{} the measure with density pi(x)=\allowbreak{}g2(x)*f0(x)\textasciicircum{}2 has symmetric pair flow pi(x)p(x,\allowbreak{}y)=\allowbreak{}min(x,\allowbreak{}y)*g2(x)*g2(y)*f0(x)*f0(y)/\allowbreak{}beta0. This recovers the credited reversibility by inspection. The obstacle to identifying it with ordinary unimodular rerooting is the reproductive f0 bias:\allowbreak{} unimodularity is a mass-transport identity for the root law,\allowbreak{} while the spine favors descendants through the right eigenfunction. One must construct and normalize a two-sided marked spine with the appropriate Radon-Nikodym factor before claiming an equivalence.


\entrylabel{Second check}

Cancel exactly one factor of f0(x),\allowbreak{} then exchange x and y. For an ordinary unimodular graph,\allowbreak{} stationary simple-random-walk rerooting instead uses degree bias,\allowbreak{} exposing the missing identification.


\entrylabel{Remaining obstruction}

No general equivalence,\allowbreak{} converse,\allowbreak{} or root-bias formula beyond this explicit kernel calculation is proved.


Verification:\allowbreak{} author checked partial | Independent review:\allowbreak{} none | Full target solved:\allowbreak{} False


\clearpage

\section*{P0452  Giant-component density from Erlang weighted-tree survival}

\entrylabel{Primary source}

Moharrami,\allowbreak{} Subramanian,\allowbreak{} Liu and Sundaresan:\allowbreak{} The Erlang Weighted Tree,\allowbreak{} A New Branching Process (v5,\allowbreak{} 27 May 2023);\allowbreak{} Section4,\allowbreak{} question3;\allowbreak{} PDF16;\allowbreak{} Corollary2.2


\url{https://arxiv.org/pdf/2002.03993}


Source statement date:\allowbreak{} 2023-05-27. Archived access:\allowbreak{} 2026-09-07T07:\allowbreak{}23:\allowbreak{}15Z.


Historical:\allowbreak{} source\_\allowbreak{}says\_\allowbreak{}open | Current:\allowbreak{} unverified | Result:\allowbreak{} unsolved


\entrylabel{Question}

Prove the lower bound matching the known upper bound:\allowbreak{} the largest La-Kabkab component has asymptotic vertex fraction equal to EWT survival probability.


\entrylabel{Scope}

The later logarithmic connectivity theorem concerns growing preference budgets and does not supply the fixed-distribution giant lower bound.


\entrylabel{Dated literature check}

2026-09-07 bounded primary searches for EWT,\allowbreak{} La-Kabkab and bilateral preference degree results;\allowbreak{} 2025 connectivity theorem credited. No current-open or novelty certification.


\subsection*{P0452-A001 | Exact conditioning,\allowbreak{} spectral integral estimate,\allowbreak{} deterministic graph comparison or order-statistic asymptotics}

\entrylabel{Argument or result}

The local-limit argument controls the fraction of vertices in components above a fixed size;\allowbreak{} it does not put those vertices in one component. To break that inference,\allowbreak{} compare one cycle of size n with a disjoint union of cycles of length floor(sqrt(n)),\allowbreak{} adjusting a negligible remainder. Both local weak limits are the bi-infinite path,\allowbreak{} which is infinite with probability one. The largest component fractions instead tend to one and zero. Thus the missing step for La-Kabkab is a genuine uniqueness/\allowbreak{}coalescence or sprinkling argument controlling dependencies among mutually preferred edges;\allowbreak{} local convergence alone cannot provide it.


\entrylabel{Second check}

For any fixed radius r,\allowbreak{} all vertices except a vanishing remainder in both examples have the same radius-r rooted path neighborhood once cycle lengths exceed 2r+\allowbreak{}1. The largest block sizes are n and O(sqrt(n)).


\entrylabel{Remaining obstruction}

This is a counterexample to a proposed proof implication on general graph sequences,\allowbreak{} not a counterexample within the La-Kabkab model.


Verification:\allowbreak{} author checked partial | Independent review:\allowbreak{} none | Full target solved:\allowbreak{} False


\clearpage

\section*{P0453  Weaker degree-tail hypotheses for Erlang weighted-tree growth}

\entrylabel{Primary source}

Moharrami,\allowbreak{} Subramanian,\allowbreak{} Liu and Sundaresan:\allowbreak{} The Erlang Weighted Tree,\allowbreak{} A New Branching Process (v5,\allowbreak{} 27 May 2023);\allowbreak{} Section4,\allowbreak{} question4;\allowbreak{} PDF16;\allowbreak{} Section5.4


\url{https://arxiv.org/pdf/2002.03993}


Source statement date:\allowbreak{} 2023-05-27. Archived access:\allowbreak{} 2026-09-07T07:\allowbreak{}23:\allowbreak{}15Z.


Historical:\allowbreak{} source\_\allowbreak{}says\_\allowbreak{}open | Current:\allowbreak{} unverified | Result:\allowbreak{} unsolved


\entrylabel{Question}

Relax the exponential-moment hypothesis in the EWT growth,\allowbreak{} spectral and extinction conclusions,\allowbreak{} with the conditions for each conclusion stated separately.


\entrylabel{Scope}

The question groups related conclusions. A compact mean operator alone is not an extinction or branching-martingale convergence theorem.


\entrylabel{Dated literature check}

2026-09-07 bounded primary searches for EWT,\allowbreak{} La-Kabkab and bilateral preference degree results;\allowbreak{} 2025 connectivity theorem credited. No current-open or novelty certification.


\subsection*{P0453-A001 | Exact conditioning,\allowbreak{} spectral integral estimate,\allowbreak{} deterministic graph comparison or order-statistic asymptotics}

\entrylabel{Argument or result}

One spectral step only needs finite mean. Put g2(x)=\allowbreak{}exp(-x)*sum\_\allowbreak{}(j>=\allowbreak{}2)P(j)*x\textasciicircum{}(j-2)/\allowbreak{}(j-2)!,\allowbreak{} and consider Kf(x)=\allowbreak{}integral min(x,\allowbreak{}y)*f(y)*g2(y)dy on L2(g2 dx). Tonelli gives integral x*g2(x)dx=\allowbreak{}sum\_\allowbreak{}(j>=\allowbreak{}2)(j-1)P(j)=\allowbreak{}E[K\_\allowbreak{}degree]-1. Since min(x,\allowbreak{}y)\textasciicircum{}2<=\allowbreak{}x*y,\allowbreak{} the squared Hilbert-Schmidt norm is at most (E[K\_\allowbreak{}degree]-1)\textasciicircum{}2. Thus finite mean already makes this symmetric integral operator compact. The kernel is also positive semidefinite because min(x,\allowbreak{}y)=\allowbreak{}integral\_\allowbreak{}0\textasciicircum{}infinity 1\_\allowbreak{}(t<x)*1\_\allowbreak{}(t<y)dt. This removes an exponential moment from the compactness step,\allowbreak{} while leaving the stronger population conclusions untouched.


\entrylabel{Second check}

The gamma integral of x\textasciicircum{}(j-1)*exp(-x)/\allowbreak{}(j-2)! equals j-1. For compactly supported test functions,\allowbreak{} Fubini expresses the quadratic form as integral (integral\_\allowbreak{}(x>t) f(x)g2(x)dx)\textasciicircum{}2 dt;\allowbreak{} continuity extends positivity to L2.


\entrylabel{Remaining obstruction}

No assertion that every source growth/\allowbreak{}extinction theorem survives under finite mean. Identifying the branching eigenfunctions,\allowbreak{} integrability of the spine,\allowbreak{} and uniform integrability of martingales still require proof.


Verification:\allowbreak{} author checked partial | Independent review:\allowbreak{} none | Full target solved:\allowbreak{} False


\clearpage

\section*{P0456  The fine connectivity window for bilateral preference graphs}

\entrylabel{Primary source}

Dabirian and Subramanian:\allowbreak{} Connectivity of a family of bilateral preference random graphs;\allowbreak{} Open Problems;\allowbreak{} PDF2-3


\url{https://arxiv.org/pdf/2209.05677}


Source statement date:\allowbreak{} 2025-07-08. Archived access:\allowbreak{} 2026-09-07T07:\allowbreak{}25:\allowbreak{}41Z.


Historical:\allowbreak{} source\_\allowbreak{}says\_\allowbreak{}open | Current:\allowbreak{} unverified | Result:\allowbreak{} unsolved


\entrylabel{Question}

Locate the sharp connectivity window and limiting connectivity probability of LK(n,\allowbreak{}k) inside the source remaining band around log(n).


\entrylabel{Scope}

The source already settles the leading log(n) threshold. The unresolved band has width of order sqrt(log(n))*log(log(n)).


\entrylabel{Dated literature check}

2026-09-07 bounded primary searches for EWT,\allowbreak{} La-Kabkab and bilateral preference degree results;\allowbreak{} 2025 connectivity theorem credited. No current-open or novelty certification.


\subsection*{P0456-A001 | Exact conditioning,\allowbreak{} spectral integral estimate,\allowbreak{} deterministic graph comparison or order-statistic asymptotics}

\entrylabel{Argument or result}

At k=\allowbreak{}log(n)+\allowbreak{}a*sqrt(log(n)),\allowbreak{} the known mean-degree expansion gives E[D]=\allowbreak{}log(n)+\allowbreak{}(a-1/\allowbreak{}sqrt(pi))*sqrt(log(n))+\allowbreak{}lower-order terms. Matching this mean to log(n) therefore suggests an apparent positive shift,\allowbreak{} but it does not derive a connectivity threshold:\allowbreak{} isolated-vertex probabilities depend on the far lower tail of D,\allowbreak{} not its mean. Even the candidate central fluctuation law in the separate packet probes deficits of order sqrt(k);\allowbreak{} isolation requires deficit k,\allowbreak{} a moving far-tail event outside that limit. An actual sharp-window proof must analyze the relevant order-statistic large deviation and then exclude nontrivial small components.


\entrylabel{Second check}

Substitute the proposed k into sqrt(k/\allowbreak{}pi) and divide the remainder by sqrt(log(n)). Separately,\allowbreak{} weak convergence at fixed scaled thresholds gives no estimate at a threshold diverging like sqrt(k).


\entrylabel{Remaining obstruction}

No connectivity shift,\allowbreak{} window width or limiting probability is established;\allowbreak{} the mean-matching calculation is explicitly a failed inference to test.


Verification:\allowbreak{} author checked partial | Independent review:\allowbreak{} none | Full target solved:\allowbreak{} False


\clearpage

\section*{P0457  Diameter near the newly established bilateral connectivity threshold}

\entrylabel{Primary source}

Dabirian and Subramanian:\allowbreak{} Connectivity of a family of bilateral preference random graphs;\allowbreak{} Open Problems;\allowbreak{} PDF3


\url{https://arxiv.org/pdf/2209.05677}


Source statement date:\allowbreak{} 2025-07-08. Archived access:\allowbreak{} 2026-09-07T07:\allowbreak{}25:\allowbreak{}41Z.


Historical:\allowbreak{} source\_\allowbreak{}says\_\allowbreak{}open | Current:\allowbreak{} unverified | Result:\allowbreak{} unsolved


\entrylabel{Question}

Characterize LK(n,\allowbreak{}k) diameter in the additional near-logarithmic range where the source proves connectivity.


\entrylabel{Scope}

The old diameter result follows from a coarser embedded graph and does not automatically extend through the improved connectivity range.


\entrylabel{Dated literature check}

2026-09-07 bounded primary searches for EWT,\allowbreak{} La-Kabkab and bilateral preference degree results;\allowbreak{} 2025 connectivity theorem credited. No current-open or novelty certification.


\subsection*{P0457-A001 | Exact conditioning,\allowbreak{} spectral integral estimate,\allowbreak{} deterministic graph comparison or order-statistic asymptotics}

\entrylabel{Argument or result}

Any connected graph with maximum degree k>=\allowbreak{}3 has at most 1+\allowbreak{}k*((k-1)\textasciicircum{}r-1)/\allowbreak{}(k-2) vertices within radius r. Hence its diameter d satisfies n<=\allowbreak{}1+\allowbreak{}k*((k-1)\textasciicircum{}d-1)/\allowbreak{}(k-2),\allowbreak{} giving a deterministic lower bound of logarithmic order log(n)/\allowbreak{}log(k). For LK this uses its exact degree cap and requires no independence. A matching upper bound needs expansion of balls grown through mutual rankings;\allowbreak{} the Erdős-Rényi subgraph used by a coarse coupling may itself be disconnected in the newly covered range and then offers no diameter bound.


\entrylabel{Second check}

Count at most k first neighbors and at most k-1 new choices on each later level. This also distinguishes graph distance from a count of path walks with arbitrary backtracking.


\entrylabel{Remaining obstruction}

No matching upper bound or limiting constant is proved. The tree-volume inequality is classical and credited as a diagnostic.


Verification:\allowbreak{} author checked partial | Independent review:\allowbreak{} none | Full target solved:\allowbreak{} False


\clearpage

\section*{P0458  The standardized degree law in growing bilateral preference graphs}

\entrylabel{Primary source}

Dabirian and Subramanian:\allowbreak{} Connectivity of a family of bilateral preference random graphs;\allowbreak{} Open Problems;\allowbreak{} PDF3 first additional question


\url{https://arxiv.org/pdf/2209.05677}


Source statement date:\allowbreak{} 2025-07-08. Archived access:\allowbreak{} 2026-09-07T07:\allowbreak{}25:\allowbreak{}41Z.


Historical:\allowbreak{} source\_\allowbreak{}says\_\allowbreak{}open | Current:\allowbreak{} unverified | Result:\allowbreak{} candidate\_\allowbreak{}partial\_\allowbreak{}counterexample


\entrylabel{Question}

Determine the limiting law of (D-E[D])/\allowbreak{}sqrt(Var(D)) for a typical LK(n,\allowbreak{}k) degree as the preference budget grows.


\entrylabel{Scope}

A private packet covers k->infinity with k=\allowbreak{}o(n\textasciicircum{}(2/\allowbreak{}5)),\allowbreak{} including the source mean-asymptotic regime k=\allowbreak{}o(n\textasciicircum{}(1/\allowbreak{}3)). Other growth regimes and novelty are not settled.


\entrylabel{Dated literature check}

2026-09-07 bounded primary searches for EWT,\allowbreak{} La-Kabkab and bilateral preference degree results;\allowbreak{} 2025 connectivity theorem credited. No current-open or novelty certification.


\subsection*{P0458-A001 | Exact conditioning,\allowbreak{} spectral integral estimate,\allowbreak{} deterministic graph comparison or order-statistic asymptotics}

\entrylabel{Argument or result}

The packet uses the exact root (k+\allowbreak{}1)st order statistic X\textasciitilde{}Beta(k+\allowbreak{}1,\allowbreak{}n-k-1). Ignoring internal edges among the root preferred neighbors gives D0|X\textasciitilde{}Bin(k,\allowbreak{}1-q(X)),\allowbreak{} where q(X)=\allowbreak{}E[(X-T)+\allowbreak{}]/\allowbreak{}X and T\textasciitilde{}Beta(k,\allowbreak{}n-2k). The actual error satisfies E[(D0-D)\textasciicircum{}2/\allowbreak{}k]=\allowbreak{}O(k\textasciicircum{}2/\allowbreak{}n+\allowbreak{}k\textasciicircum{}5/\allowbreak{}n\textasciicircum{}2). Beta central limits and explicit moment controls then give (k-D)/\allowbreak{}sqrt(k)->g(Z),\allowbreak{} g(z)=\allowbreak{}phi(z)+\allowbreak{}z*Phi(z),\allowbreak{} with convergence of its first two moments. The variance constant is v=\allowbreak{}1/\allowbreak{}3+\allowbreak{}(sqrt(3)-2)/\allowbreak{}(2*pi),\allowbreak{} and the standardized degree limit is (1/\allowbreak{}sqrt(pi)-g(Z))/\allowbreak{}sqrt(v),\allowbreak{} which is non-Gaussian.


\entrylabel{Second check}

All 720 rank permutations on four vertices for k=\allowbreak{}1,\allowbreak{}2 satisfy D<=\allowbreak{}D0<=\allowbreak{}D+\allowbreak{}2J;\allowbreak{} degree means 3/\allowbreak{}5 and 8/\allowbreak{}5 agree with independent two-endpoint rank probabilities. The separate fixed-k EWT root mixture followed by a gamma limit yields the same g(Z). Exact conditioning,\allowbreak{} L2 comparison,\allowbreak{} uniform integrability and a polar calculation of v are detailed in the packet.


\entrylabel{Remaining obstruction}

Author-only partial candidate for the broader source question. No independent review,\allowbreak{} current-openness or novelty certificate,\allowbreak{} and no assertion for every k regime.


Verification:\allowbreak{} author checked partial | Independent review:\allowbreak{} none | Full target solved:\allowbreak{} False


\entrylabel{Full proof packet}

proofs/\allowbreak{}Bilateral\_\allowbreak{}degree\_\allowbreak{}fluctuations\_\allowbreak{}partial.md


\clearpage

\section*{P0459  Simultaneous concentration of bilateral preference degrees}

\entrylabel{Primary source}

Dabirian and Subramanian:\allowbreak{} Connectivity of a family of bilateral preference random graphs;\allowbreak{} Open Problems;\allowbreak{} PDF3 second additional question


\url{https://arxiv.org/pdf/2209.05677}


Source statement date:\allowbreak{} 2025-07-08. Archived access:\allowbreak{} 2026-09-07T07:\allowbreak{}25:\allowbreak{}41Z.


Historical:\allowbreak{} source\_\allowbreak{}says\_\allowbreak{}open | Current:\allowbreak{} unverified | Result:\allowbreak{} candidate\_\allowbreak{}partial\_\allowbreak{}counterexample


\entrylabel{Question}

Find growth conditions and quantitative bounds for LK(n,\allowbreak{}k) degrees to concentrate simultaneously around their mean.


\entrylabel{Scope}

The source asks for concentration conditions. The private packet supplies a sufficient superlogarithmic regime using the source existing ER-sandwich method,\allowbreak{} with no optimality or novelty claim.


\entrylabel{Dated literature check}

2026-09-07 bounded primary searches for EWT,\allowbreak{} La-Kabkab and bilateral preference degree results;\allowbreak{} 2025 connectivity theorem credited. No current-open or novelty certification.


\subsection*{P0459-A001 | Exact conditioning,\allowbreak{} spectral integral estimate,\allowbreak{} deterministic graph comparison or order-statistic asymptotics}

\entrylabel{Argument or result}

Put delta=\allowbreak{}sqrt(6*k*log(n)) and assume k>2*delta. Threshold edge weights at t=\allowbreak{}(k-delta)/\allowbreak{}(n-1). Each threshold degree B\_\allowbreak{}v is binomial with mean k-delta. Chernoff and a union bound give k-2*delta<=\allowbreak{}B\_\allowbreak{}v<=\allowbreak{}k for all vertices with probability at least 1-2/\allowbreak{}n\textasciicircum{}2. On that event every threshold edge is preferred at both endpoints,\allowbreak{} so k-2*delta<=\allowbreak{}D\_\allowbreak{}v<=\allowbreak{}k simultaneously. Thus k/\allowbreak{}log(n)->infinity suffices for relative concentration. The argument uses independence of the underlying edge weights,\allowbreak{} and never independence of vertex degrees.


\entrylabel{Second check}

Upper and lower Chernoff exponents are at least delta\textasciicircum{}2/\allowbreak{}(2*k)=\allowbreak{}3*log(n),\allowbreak{} followed by a union over n vertices. The deterministic threshold inclusion follows directly from having at most k incident edges below t at each endpoint.


\entrylabel{Remaining obstruction}

A sufficient-condition partial only. Sharp tails and the logarithmic boundary are not determined;\allowbreak{} below connectivity isolated vertices rule out concentration of every degree near k.


Verification:\allowbreak{} author checked partial | Independent review:\allowbreak{} none | Full target solved:\allowbreak{} False


\entrylabel{Full proof packet}

proofs/\allowbreak{}Bilateral\_\allowbreak{}degree\_\allowbreak{}fluctuations\_\allowbreak{}partial.md


\clearpage

\section*{P0461  Entropy concavity for determinantal kernels}

\entrylabel{Primary source}

Lyons:\allowbreak{} Determinantal probability,\allowbreak{} basic properties and conjectures,\allowbreak{} ICM 2014;\allowbreak{} Conjectures2.6 and5.3;\allowbreak{} PDF5,\allowbreak{}22


\url{https://arxiv.org/pdf/1406.2707}


Source statement date:\allowbreak{} 2014-06-10. Archived access:\allowbreak{} 2026-09-07T06:\allowbreak{}19:\allowbreak{}37Z.


Historical:\allowbreak{} source\_\allowbreak{}says\_\allowbreak{}open | Current:\allowbreak{} unverified | Result:\allowbreak{} unsolved


\entrylabel{Question}

Prove or disprove concavity of finite determinantal Shannon entropy in a positive-contraction kernel. Include the source-related stationary-symbol entropy conjecture as a grouped consequence.


\entrylabel{Scope}

Finite-kernel concavity implies the stationary version by restriction to finite boxes and entropy-rate limits. They are not asserted equivalent. Diagonal means diagonal in the coordinate basis defining the point process,\allowbreak{} not merely simultaneously diagonalizable.


\entrylabel{Dated literature check}

2026-09-07 bounded primary searches on entropy concavity,\allowbreak{} union coupling,\allowbreak{} block factors,\allowbreak{} Fuglede-Kadison domination and the linear distance bound. Known rank-one coupling and degree-one block-factor cases credited. No current-openness certificate.


\subsection*{P0461-A001 | Integrator-authored exact special-case argument and failure analysis after checking the incoming source}

\entrylabel{Argument or result}

The coordinate-diagonal case is a sum of binary entropies,\allowbreak{} whose second derivatives are -1/\allowbreak{}[q(1-q)]. There is also an explicit non-diagonal line:\allowbreak{} for Q(c)=\allowbreak{}[[1/\allowbreak{}2,\allowbreak{}c],\allowbreak{}[c,\allowbreak{}1/\allowbreak{}2]],\allowbreak{} |c|<1/\allowbreak{}2,\allowbreak{} the empty and doubleton probabilities are a=\allowbreak{}1/\allowbreak{}4-c\textasciicircum{}2 and the two singleton probabilities are b=\allowbreak{}1/\allowbreak{}4+\allowbreak{}c\textasciicircum{}2. Hence H(c)=\allowbreak{}-2*a*log(a)-2*b*log(b),\allowbreak{} and Hsecond(c)=\allowbreak{}4*log(a/\allowbreak{}b)-8*c\textasciicircum{}2*(1/\allowbreak{}a+\allowbreak{}1/\allowbreak{}b)<=\allowbreak{}0. Continuity includes c=\allowbreak{}+\allowbreak{}-1/\allowbreak{}2. This verifies concavity on this affine kernel line. In general the determinant probabilities depend nonlinearly on Q,\allowbreak{} so concavity of Shannon entropy in a probability vector cannot be applied directly to an arithmetic kernel midpoint.


\entrylabel{Second check}

Inclusion-exclusion gives p11=\allowbreak{}det(Q),\allowbreak{} p10=\allowbreak{}Q11-det(Q),\allowbreak{} p01=\allowbreak{}Q22-det(Q),\allowbreak{} p00=\allowbreak{}det(I-Q);\allowbreak{} these are a,\allowbreak{}b,\allowbreak{}b,\allowbreak{}a and sum to one. Differentiating H first gives 4*c*log(a/\allowbreak{}b),\allowbreak{} independently confirming the second derivative and its sign.


\entrylabel{Remaining obstruction}

No joint concavity proof,\allowbreak{} no extension to arbitrary off-diagonal or noncommuting kernels,\allowbreak{} and no numerical trials treated as certification.


Verification:\allowbreak{} author checked restricted argument | Independent review:\allowbreak{} none | Full target solved:\allowbreak{} False


\clearpage

\section*{P0462  Disjoint coupling of orthogonal projection processes}

\entrylabel{Primary source}

Lyons:\allowbreak{} Determinantal probability,\allowbreak{} basic properties and conjectures,\allowbreak{} ICM 2014;\allowbreak{} Question2.8;\allowbreak{} PDF6


\url{https://arxiv.org/pdf/1406.2707}


Source statement date:\allowbreak{} 2014-06-10. Archived access:\allowbreak{} 2026-09-07T06:\allowbreak{}19:\allowbreak{}37Z.


Historical:\allowbreak{} source\_\allowbreak{}says\_\allowbreak{}open | Current:\allowbreak{} unverified | Result:\allowbreak{} unsolved


\entrylabel{Question}

For H=\allowbreak{}H1 orthogonal-direct-sum H2 in l2(E),\allowbreak{} construct or rule out a coupling of their projection determinantal processes whose union has the projection law of H.


\entrylabel{Scope}

Finite-rank cardinalities force disjointness. The rank-one argument below is a credited consequence of stochastic domination,\allowbreak{} already present in Lyons2003;\allowbreak{} it is not a new solution.


\entrylabel{Dated literature check}

2026-09-07 bounded primary searches on entropy concavity,\allowbreak{} union coupling,\allowbreak{} block factors,\allowbreak{} Fuglede-Kadison domination and the linear distance bound. Known rank-one coupling and degree-one block-factor cases credited. No current-openness certificate.


\subsection*{P0462-A001 | Integrator-authored exact special-case argument and failure analysis after checking the incoming source}

\entrylabel{Argument or result}

If H2 is spanned by a unit vector u and H1 has finite rank,\allowbreak{} monotone-couple A with law P\_\allowbreak{}H1 inside C with law P\_\allowbreak{}H. Their deterministic sizes differ by one,\allowbreak{} so C minus A is a singleton B. At any coordinate e,\allowbreak{} P(B=\allowbreak{}\{e\})=\allowbreak{}P(e in C)-P(e in A)=\allowbreak{}(P\_\allowbreak{}H)ee-(P\_\allowbreak{}H1)ee=\allowbreak{}|u(e)|\textasciicircum{}2. Thus B has precisely P\_\allowbreak{}H2 and (A,\allowbreak{}B) gives the required coupling in this special case. For higher rank,\allowbreak{} the same subtraction determines only one-point intensities of C minus A;\allowbreak{} its higher inclusion probabilities are not fixed. The finite problem can be written as nonnegative weights w(A,\allowbreak{}B) on disjoint pairs with the three required marginals,\allowbreak{} but setting up this transportation problem does not prove it feasible.


\entrylabel{Second check}

The rank-one masses sum to ||u||\textasciicircum{}2=\allowbreak{}1 and each difference indicator equals 1\_\allowbreak{}(e in C)-1\_\allowbreak{}(e in A). A random pair of same-size subsets can have matching one-point intensities yet different two-point probabilities,\allowbreak{} showing why the rank-one calculation does not generalize by intensity matching alone.


\entrylabel{Remaining obstruction}

No general coupling construction. Earlier floating-point LP claims are retained only in the frozen incoming record,\allowbreak{} not counted as exact certificates.


Verification:\allowbreak{} author checked restricted argument | Independent review:\allowbreak{} none | Full target solved:\allowbreak{} False


\clearpage

\section*{P0463  Block factors for finite-range determinantal symbols}

\entrylabel{Primary source}

Lyons:\allowbreak{} Determinantal probability,\allowbreak{} basic properties and conjectures,\allowbreak{} ICM 2014;\allowbreak{} Question5.4;\allowbreak{} PDF22


\url{https://arxiv.org/pdf/1406.2707}


Source statement date:\allowbreak{} 2014-06-10. Archived access:\allowbreak{} 2026-09-07T06:\allowbreak{}19:\allowbreak{}37Z.


Historical:\allowbreak{} source\_\allowbreak{}says\_\allowbreak{}open | Current:\allowbreak{} unverified | Result:\allowbreak{} unsolved


\entrylabel{Question}

For a trigonometric polynomial f:\allowbreak{}T->[0,\allowbreak{}1] of degree m,\allowbreak{} is its stationary determinantal process an (m+\allowbreak{}1)-block factor of iid variables?


\entrylabel{Scope}

The source credits Broman2005 for degree one. A sublattice reduction proves new cases only when the lower-degree process already has a suitable block-factor construction;\allowbreak{} it is not an induction proving every even degree.


\entrylabel{Dated literature check}

2026-09-07 bounded primary searches on entropy concavity,\allowbreak{} union coupling,\allowbreak{} block factors,\allowbreak{} Fuglede-Kadison domination and the linear distance bound. Known rank-one coupling and degree-one block-factor cases credited. No current-openness certificate.


\subsection*{P0463-A001 | Integrator-authored exact special-case argument and failure analysis after checking the incoming source}

\entrylabel{Argument or result}

Let f(theta)=\allowbreak{}g(k*theta). Its Fourier coefficients vanish at lags not divisible by k,\allowbreak{} so the determinantal kernel is block diagonal on the k residue classes. Determinant factorization proves that these restrictions are independent copies of P\_\allowbreak{}g. If P\_\allowbreak{}g can be represented as F(U\_\allowbreak{}j,\allowbreak{}...,\allowbreak{}U\_\allowbreak{}(j+\allowbreak{}r)) with iid U,\allowbreak{} define Y\_\allowbreak{}i=\allowbreak{}F(U\_\allowbreak{}i,\allowbreak{}U\_\allowbreak{}(i+\allowbreak{}k),\allowbreak{}...,\allowbreak{}U\_\allowbreak{}(i+\allowbreak{}r*k)). This is a shift-equivariant function of the window U\_\allowbreak{}i,\allowbreak{}...,\allowbreak{}U\_\allowbreak{}(i+\allowbreak{}r*k),\allowbreak{} so is an (r*k+\allowbreak{}1)-block factor. The residue-class families use disjoint iid coordinates and have the correct laws. In particular any single-harmonic symbol a+\allowbreak{}b*cos(k*theta)+\allowbreak{}c*sin(k*theta) with range in [0,\allowbreak{}1] inherits the known degree-one construction and is a (k+\allowbreak{}1)-block factor.


\entrylabel{Second check}

For cylinder events on distinct residue classes,\allowbreak{} all mixed kernel entries vanish and every required inclusion determinant factors. Inclusion-exclusion upgrades those identities to independence of the complete finite-dimensional distributions. The proposed coding depends only on positions at distances k,\allowbreak{} making its full window length exactly r*k+\allowbreak{}1.


\entrylabel{Remaining obstruction}

General mixed-harmonic symbols remain untreated. Finite dependence by itself does not imply a block-factor representation.


Verification:\allowbreak{} author checked restricted argument | Independent review:\allowbreak{} none | Full target solved:\allowbreak{} False


\clearpage

\section*{P0464  Geometric-mean product bounds for equivariant determinantal laws}

\entrylabel{Primary source}

Lyons:\allowbreak{} Determinantal probability,\allowbreak{} basic properties and conjectures,\allowbreak{} ICM 2014;\allowbreak{} Conjecture5.7;\allowbreak{} PDF23


\url{https://arxiv.org/pdf/1406.2707}


Source statement date:\allowbreak{} 2014-06-10. Archived access:\allowbreak{} 2026-09-07T06:\allowbreak{}19:\allowbreak{}37Z.


Historical:\allowbreak{} source\_\allowbreak{}says\_\allowbreak{}open | Current:\allowbreak{} unverified | Result:\allowbreak{} unsolved


\entrylabel{Question}

For a Gamma-equivariant positive contraction Q on l2(Gamma),\allowbreak{} establish or refute the optimal stochastic product bounds with lower density FK(Q)=\allowbreak{}exp(tr log Q) and upper density 1-FK(I-Q).


\entrylabel{Scope}

The following classical conditioning argument gives weaker harmonic-mean bounds for finite groups and 0<Q<I. It is an author-checked supporting calculation,\allowbreak{} not a novelty claim or a solution of the geometric-mean conjecture.


\entrylabel{Dated literature check}

2026-09-07 bounded primary searches on entropy concavity,\allowbreak{} union coupling,\allowbreak{} block factors,\allowbreak{} Fuglede-Kadison domination and the linear distance bound. Known rank-one coupling and degree-one block-factor cases credited. No current-openness certificate.


\subsection*{P0464-A001 | Integrator-authored exact special-case argument and failure analysis after checking the incoming source}

\entrylabel{Argument or result}

Put J=\allowbreak{}Q(I-Q)\textasciicircum{}(-1). Given the occupied set B outside e,\allowbreak{} the conditional odds of presence are s\_\allowbreak{}B=\allowbreak{}J\_\allowbreak{}ee-J\_\allowbreak{}eB*(J\_\allowbreak{}BB)\textasciicircum{}(-1)*J\_\allowbreak{}Be. Schur-complement minimization makes s\_\allowbreak{}B decrease as B increases. Its minimum is s\_\allowbreak{}all=\allowbreak{}1/\allowbreak{}(J\textasciicircum{}(-1))ee=\allowbreak{}1/\allowbreak{}[(Q\textasciicircum{}(-1))ee-1]. The conditional probability,\allowbreak{} however,\allowbreak{} is s\_\allowbreak{}all/\allowbreak{}(1+\allowbreak{}s\_\allowbreak{}all)=\allowbreak{}1/\allowbreak{}(Q\textasciicircum{}(-1))ee. Equivariance makes (Q\textasciicircum{}(-1))ee=\allowbreak{}tr(Q\textasciicircum{}(-1)),\allowbreak{} with normalized trace,\allowbreak{} so this probability is the harmonic mean h of all eigenvalues counted with multiplicity. Sequential sampling with common independent uniforms couples the process above iid Bernoulli(h),\allowbreak{} since every conditional probability given a partial reveal is also at least h. Applying the same argument to the complement gives the upper density 1-HM(I-Q). AM-GM applied to the reciprocals gives HM(Q)<=\allowbreak{}FK(Q),\allowbreak{} exposing the exact remaining gap.


\entrylabel{Second check}

For a scalar Q=\allowbreak{}qI,\allowbreak{} s\_\allowbreak{}all=\allowbreak{}q/\allowbreak{}(1-q) but the conditional probability is q;\allowbreak{} this catches the odds/\allowbreak{}probability confusion in the incoming record. In a finite group,\allowbreak{} equal inverse diagonal entries average to the full matrix trace divided by |Gamma|,\allowbreak{} irrespective of whether the group is abelian. No sum over one-dimensional characters is substituted for the nonabelian spectrum.


\entrylabel{Remaining obstruction}

No geometric-mean domination theorem or sharpness certificate. Numerical coupling feasibility is not exact proof,\allowbreak{} and the infinite-group extension of this particular calculation is not asserted.


Verification:\allowbreak{} author checked restricted argument | Independent review:\allowbreak{} none | Full target solved:\allowbreak{} False


\clearpage

\section*{P0465  Linear coupling distance for noncommuting equivariant kernels}

\entrylabel{Primary source}

Lyons:\allowbreak{} Determinantal probability,\allowbreak{} basic properties and conjectures,\allowbreak{} ICM 2014;\allowbreak{} Section5.2 afterTheorem5.6;\allowbreak{} PDF23


\url{https://arxiv.org/pdf/1406.2707}


Source statement date:\allowbreak{} 2014-06-10. Archived access:\allowbreak{} 2026-09-07T06:\allowbreak{}19:\allowbreak{}37Z.


Historical:\allowbreak{} source\_\allowbreak{}says\_\allowbreak{}open | Current:\allowbreak{} unverified | Result:\allowbreak{} unsolved


\entrylabel{Question}

For sofic-group equivariant positive contractions Q,\allowbreak{}Qprime,\allowbreak{} does the invariant normalized Hamming coupling distance satisfy dbar(P\_\allowbreak{}Q,\allowbreak{}P\_\allowbreak{}Qprime)<=\allowbreak{}||Q-Qprime||\_\allowbreak{}1 without a commutativity assumption?


\entrylabel{Scope}

The source states the linear bound for commuting pairs and a weaker power bound in general. We examine why the natural common-upper-kernel proof fails;\allowbreak{} the example is not a counterexample to the actual coupling inequality.


\entrylabel{Dated literature check}

2026-09-07 bounded primary searches on entropy concavity,\allowbreak{} union coupling,\allowbreak{} block factors,\allowbreak{} Fuglede-Kadison domination and the linear distance bound. Known rank-one coupling and degree-one block-factor cases credited. No current-openness certificate.


\subsection*{P0465-A001 | Integrator-authored exact special-case argument and failure analysis after checking the incoming source}

\entrylabel{Argument or result}

Write D=\allowbreak{}Q-Qprime. A tempting common upper kernel is R=\allowbreak{}Qprime+\allowbreak{}Dplus=\allowbreak{}(Q+\allowbreak{}Qprime+\allowbreak{}|D|)/\allowbreak{}2. It dominates both in operator order,\allowbreak{} but need not be a contraction. In dimension two take Q=\allowbreak{}diag(1,\allowbreak{}0) and Qprime=\allowbreak{}(1/\allowbreak{}2)*[[1,\allowbreak{}1],\allowbreak{}[1,\allowbreak{}1]]. Then D\textasciicircum{}2=\allowbreak{}I/\allowbreak{}2 and |D|=\allowbreak{}I/\allowbreak{}sqrt(2). The largest eigenvalue of Q+\allowbreak{}Qprime is 1+\allowbreak{}1/\allowbreak{}sqrt(2),\allowbreak{} giving lambda\_\allowbreak{}max(R)=\allowbreak{}1/\allowbreak{}2+\allowbreak{}1/\allowbreak{}sqrt(2)>1. Thus R is not an admissible determinantal kernel and one cannot apply monotone-coupling bounds through it. This same obstruction can be placed in the two-dimensional irreducible block of the right group algebra of S3,\allowbreak{} with zeros on the other blocks;\allowbreak{} the resulting operators commute with left translations but not with one another.


\entrylabel{Second check}

Direct multiplication gives D=\allowbreak{}[[1/\allowbreak{}2,\allowbreak{}-1/\allowbreak{}2],\allowbreak{}[-1/\allowbreak{}2,\allowbreak{}-1/\allowbreak{}2]] and D\textasciicircum{}2=\allowbreak{}I/\allowbreak{}2. The right regular representation of S3 contains its two-dimensional irreducible representation,\allowbreak{} and equivariant operators can independently specify a matrix in that block. On abelian groups the equivariant convolution algebra is commutative,\allowbreak{} so that setting cannot test this obstruction.


\entrylabel{Remaining obstruction}

Failure of a common-upper-kernel construction does not refute the desired distance inequality. No claim of having executed a nonabelian coupling LP.


Verification:\allowbreak{} author checked restricted argument | Independent review:\allowbreak{} none | Full target solved:\allowbreak{} False


\clearpage

\section*{Excluded and resolved source candidates}

\subsection*{P0096  Four-edge-minimum lower bound for the FPP time constant}

Auffinger,\allowbreak{} Damron and Hanson:\allowbreak{} 50 years of first-passage percolation;\allowbreak{} Question12,\allowbreak{} PDF pp.32-33 boundary;\allowbreak{} repeated p.146;\allowbreak{} manuscript26September2016


\url{https://arxiv.org/pdf/1511.03262}


\entrylabel{current status}

resolved\_\allowbreak{}in\_\allowbreak{}literature


\entrylabel{result status}

literature\_\allowbreak{}counterexample\_\allowbreak{}rederived


\entrylabel{statement}

For iid nonnegative edge times on Z\textasciicircum{}2 with F(0)=\allowbreak{}0,\allowbreak{} is the axial time constant at least the expected minimum of four independent edge times?


\entrylabel{statement caveat}

E\_\allowbreak{}4 is expectation of a minimum,\allowbreak{} not a fourth moment. The target is the universal survey formulation;\allowbreak{} the motivating simulation paper reports selected families.


\entrylabel{resolution evidence}

Daniel Ahlberg,\allowbreak{} Tertiles and the time constant,\allowbreak{} arXiv1911.09949 (2019;\allowbreak{}v2May2020),\allowbreak{} Journal of Applied Probability57(2),\allowbreak{}407-408,\allowbreak{}2020. The full two-page paper was read. Theorem1 gives mu<=\allowbreak{}2t\_\allowbreak{}(2/\allowbreak{}3);\allowbreak{} final paragraph explicitly refutes Alm-Deijfen and surveyQuestion12. Published DOI10.1017/\allowbreak{}jpr.2020.9.


\entrylabel{resolution url}

\url{https://arxiv.org/pdf/1911.09949}


\entrylabel{reason}

Already resolved before this campaign;\allowbreak{} own construction is retained as a rederivation,\allowbreak{} excluded from1000 current-open target and original-solution counts.


\entrylabel{Retained derivation or attempt}

Take tau=\allowbreak{}eta+\allowbreak{}epsilon U with P(eta=\allowbreak{}0)=\allowbreak{}3/\allowbreak{}4 and U uniform[0,\allowbreak{}1]. For epsilon>0 this is atomless and bounded with F(0)=\allowbreak{}0. The four-edge minimum has expectation at least1/\allowbreak{}256. The eta-model time constant is0. For each n,\allowbreak{} passage times converge to their eta-model values and are bounded by2n;\allowbreak{} dominated convergence and mu=\allowbreak{}inf\_\allowbreak{}n E T(0,\allowbreak{}ne1)/\allowbreak{}n imply mu\_\allowbreak{}epsilon tends to0. All sufficiently small epsilon reverse the proposed inequality.


\entrylabel{Status limit}

No mathematical gap identified in the written candidate;\allowbreak{} later-literature and erratum search remain incomplete. No explicit numerical epsilon threshold or independent review is claimed.


\entrylabel{Retained derivation or attempt}

Ahlberg2019/\allowbreak{}2020 already refuted the question. Applying his mu<=\allowbreak{}2t\_\allowbreak{}(2/\allowbreak{}3) to the independently drafted tau=\allowbreak{}eta+\allowbreak{}epsilon U family gives t\_\allowbreak{}(2/\allowbreak{}3)=\allowbreak{}8epsilon/\allowbreak{}9,\allowbreak{} hence mu<=\allowbreak{}16epsilon/\allowbreak{}9. With epsilon=\allowbreak{}1/\allowbreak{}1024,\allowbreak{} mu<=\allowbreak{}1/\allowbreak{}576<1/\allowbreak{}256<=\allowbreak{}E minfour. This supplies an explicit atomless bounded counterexample using an established theorem.


\entrylabel{Status limit}

No current-open problem remains for this universal inequality. This is literature-established,\allowbreak{} not an original solution.


\subsection*{P0140  Random critical drift in the annealed thinned Levy process}

Van der Hofstad:\allowbreak{} Stochastic processes on random graphs,\allowbreak{} Saint-Flour2017;\allowbreak{} Open Problem4.1,\allowbreak{} PDF p.210;\allowbreak{} equations4.5.5 and4.5.20


\url{https://www.win.tue.nl/~rhofstad/SaintFlour_SPoRG_fin_Saint_Flour.pdf}


\entrylabel{current status}

resolved\_\allowbreak{}in\_\allowbreak{}literature


\entrylabel{result status}

literature\_\allowbreak{}result\_\allowbreak{}rederived\_\allowbreak{}with\_\allowbreak{}scope\_\allowbreak{}correction


\entrylabel{statement}

Explain directly how mixing the heavy-tailed degree weights yields the independent-increment iid-degree exploration limit.


\entrylabel{statement caveat}

The deterministic drift in the short2017 formulation must be replaced by the random critical shift Lambda\_\allowbreak{}0. Dhara et al2019 equation2.22 and Remark11 already identify the corrected law indirectly.


\entrylabel{later resolution}

\{"finding":\allowbreak{} "Corrected law is indirectly identified;\allowbreak{} direct marking proof supplied here is not counted as a new theorem",\allowbreak{} "locator":\allowbreak{} "Equation2.22 and Remark11,\allowbreak{} PDF pp.9-10",\allowbreak{} "url":\allowbreak{} "\url{https://arxiv.org/pdf/1612.00650"%7D}


\entrylabel{Retained derivation or attempt}

For PPP weights with intensity(tau-1)c*x\textasciicircum{}(-tau)dx and exponential marks of rate x/\allowbreak{}mu,\allowbreak{} the random drift Z/\allowbreak{}mu cancels the empirical quadratic compensator and leaves independent increments from a deterministic marked-Poisson intensity. With a deterministic drift,\allowbreak{} a mixed log-Laplace difference is positive for small a,\allowbreak{} so increments are dependent.


\entrylabel{Status limit}

No unproved theorem is claimed:\allowbreak{} the corrected law was already identified indirectly in2019. Exact source normalization and later availability of the direct proof are retained as scope/\allowbreak{}novelty limits.


\subsection*{P0206  Random-to-random insertion cutoff}

Levin-Peres:\allowbreak{} Markov Chains and Mixing Times,\allowbreak{} second edition;\allowbreak{} Chapter26 Question6(b)


\url{https://pages.uoregon.edu/dlevin/MARKOV/mcmt2e.pdf}


\entrylabel{current status}

resolved\_\allowbreak{}in\_\allowbreak{}literature


\entrylabel{result status}

excluded\_\allowbreak{}prior\_\allowbreak{}resolution


\entrylabel{resolution evidence}

Bernstein-Nestoridi prove cutoff at(3/\allowbreak{}4)*N*log N in Annals of Probability2019.


\entrylabel{resolution url}

\url{https://arxiv.org/abs/1703.06210}


\subsection*{P0207  Mixing of the supercritical hypercube giant}

Yuval Peres:\allowbreak{} related open Markov-chain problems;\allowbreak{} Related open problems Question3


\url{https://www.yuval-peres-books.com/markov-chains-and-mixing-times/}


\entrylabel{current status}

claimed\_\allowbreak{}resolved\_\allowbreak{}in\_\allowbreak{}literature


\entrylabel{result status}

excluded\_\allowbreak{}prior\_\allowbreak{}resolution


\entrylabel{resolution evidence}

Anastos-Diskin-Lichev-Zhukovskii arXiv2510.13348v4May2026 Theorem1(b) proves order d\textasciicircum{}2 for fixed c>1;\allowbreak{} abstract and full theorem read.


\entrylabel{resolution url}

\url{https://arxiv.org/abs/2510.13348}


\subsection*{P0237  Strictly fractal dimension of the Poisson-coloring frontier}

Aldous:\allowbreak{} Random partitions of the plane via Poissonian coloring and coalescing planar partitions;\allowbreak{} Section5.3.3


\url{https://arxiv.org/pdf/1701.00131}


\entrylabel{current status}

resolved\_\allowbreak{}in\_\allowbreak{}literature


\entrylabel{result status}

excluded\_\allowbreak{}prior\_\allowbreak{}resolution


\entrylabel{resolution evidence}

Basdevant-Blanc-Curien-Singh2023 Theorem1.1 proves d-1<dim\_\allowbreak{}H F<d and explicitly says it answers this Aldous conjecture.


\entrylabel{resolution url}

\url{https://alea.impa.br/articles/v20/20-59.pdf}


\subsection*{P0238  Zero area of the Poisson-coloring frontier}

Aldous:\allowbreak{} Random partitions of the plane via Poissonian coloring and coalescing planar partitions;\allowbreak{} Conjecture3


\url{https://arxiv.org/pdf/1701.00131}


\entrylabel{current status}

resolved\_\allowbreak{}in\_\allowbreak{}literature


\entrylabel{result status}

excluded\_\allowbreak{}prior\_\allowbreak{}resolution


\entrylabel{resolution evidence}

Preater2009 Theorem2 proves zero Lebesgue measure;\allowbreak{} the2023 primary paper explicitly identifies this as the earlier answer to Aldous Conjecture3.


\entrylabel{resolution url}

\url{https://www.stat.berkeley.edu/users/aldous/Research/OP/preater.pdf}


\subsection*{P0239  Exclusion of tree subnetworks in SIRSNs}

Aldous:\allowbreak{} Scale-invariant random spatial networks;\allowbreak{} Open Problem36,\allowbreak{}PDFp52


\url{https://arxiv.org/pdf/1204.0817}


\entrylabel{current status}

resolved\_\allowbreak{}in\_\allowbreak{}literature


\entrylabel{result status}

excluded\_\allowbreak{}prior\_\allowbreak{}resolution


\entrylabel{resolution evidence}

Aldous2021 Theorem2 forces infinite average route length for large bounded separations in every invariant Poisson tree network. A SIRSN has E length at separation s equal to Delta*s with finite Delta,\allowbreak{} giving finite radial average2*Delta*r/\allowbreak{}3,\allowbreak{} a contradiction. This is an application of the established theorem,\allowbreak{} not a new open-problem solution.


\entrylabel{resolution url}

\url{https://arxiv.org/abs/2103.00669}


\subsection*{P0240  Poisson-line construction of a SIRSN}

Aldous:\allowbreak{} Scale-invariant random spatial networks;\allowbreak{} Section4.1 candidate construction


\url{https://arxiv.org/pdf/1204.0817}


\entrylabel{current status}

resolved\_\allowbreak{}in\_\allowbreak{}literature


\entrylabel{result status}

excluded\_\allowbreak{}prior\_\allowbreak{}resolution


\entrylabel{resolution evidence}

Kahn2016 proves that the marked Poisson-line metric is a SIRSN in every dimension.


\entrylabel{resolution url}

\url{https://arxiv.org/abs/1503.03976}


\subsection*{P0241  Measure-theoretic compactification of finite reversible chains}

Aldous:\allowbreak{} Compactification notes for reversible Markov chains;\allowbreak{} Conjecture2 and refined uniqueness


\url{https://www.stat.berkeley.edu/~aldous/Talks/MCcompact.pdf}


\entrylabel{current status}

resolved\_\allowbreak{}in\_\allowbreak{}literature


\entrylabel{result status}

excluded\_\allowbreak{}prior\_\allowbreak{}resolution


\entrylabel{resolution evidence}

Towsner2015 proves existence and refined uniqueness of the measurable limit construction. The topological upgrade remains separately recorded as P0235.


\entrylabel{resolution url}

\url{https://arxiv.org/abs/1404.3815}


\subsection*{P0248  Diameter-squared optimized relaxation on Cayley graphs}

Aldous-Fill:\allowbreak{} Reversible Markov Chains and Random Walks on Graphs,\allowbreak{} 2002 recompiled2014;\allowbreak{} Open Problem7.18 and Theorem7.16,\allowbreak{}PDF253


\url{https://www.stat.berkeley.edu/~aldous/RWG/book.pdf}


\entrylabel{current status}

resolved\_\allowbreak{}in\_\allowbreak{}literature


\entrylabel{result status}

excluded\_\allowbreak{}prior\_\allowbreak{}resolution


\entrylabel{statement}

Given a finite group and a symmetric generating set of Cayley diameter D,\allowbreak{} can a symmetric step law supported on that set always have relaxation time at most a universal constant times D\textasciicircum{}2?


\entrylabel{statement caveat}

The source asks about optimizing the step law,\allowbreak{} not the uniform generating-set walk or the absolute eigenvalue gap. The proof below is a complete candidate;\allowbreak{} prior literature and novelty remain unverified.


\entrylabel{resolution evidence}

Pak SODA1999 Theorem3.1 explicitly resolves AF7.18. Author PDF dated2004;\allowbreak{} proof rederivation retained.


\entrylabel{resolution url}

\url{https://www.math.ucla.edu/~pak/papers/soda4.pdf}


\entrylabel{Later primary source}

Pak:\allowbreak{} Using stopping times to bound mixing times,\allowbreak{}SODA1999,\allowbreak{}author PDF dated2004;\allowbreak{} 


\url{https://www.math.ucla.edu/~pak/papers/soda4.pdf}


\entrylabel{Retained derivation or attempt}

Choose a shortest generator word for each group element x,\allowbreak{} and symmetrize these random word choices under x going to x\textasciicircum{}-1 and word reversal with inverse letters. Put a\_\allowbreak{}g=\allowbreak{}E\_\allowbreak{}x[length(word\_\allowbreak{}x)*number\_\allowbreak{}of\_\allowbreak{}g\_\allowbreak{}letters(word\_\allowbreak{}x)],\allowbreak{} x uniform,\allowbreak{} and A=\allowbreak{}sum\_\allowbreak{}g a\_\allowbreak{}g=\allowbreak{}E\_\allowbreak{}x length(word\_\allowbreak{}x)\textasciicircum{}2<=\allowbreak{}D\textasciicircum{}2. Then a\_\allowbreak{}g=\allowbreak{}a\_\allowbreak{}(g\textasciicircum{}-1),\allowbreak{} and mu(g)=\allowbreak{}a\_\allowbreak{}g/\allowbreak{}A is a symmetric probability law. The source comparison theorem with the uniform jump-to-any-group-element chain,\allowbreak{} whose relaxation time is1,\allowbreak{} gives tau2(mu)<=\allowbreak{}max\_\allowbreak{}g a\_\allowbreak{}g/\allowbreak{}mu(g)=\allowbreak{}A<=\allowbreak{}D\textasciicircum{}2.


\entrylabel{Status limit}

The mathematical bound is complete under the displayed standard definitions. Current-literature disposition and originality are not verified;\allowbreak{} the simplicity of optimizing the source comparison bound makes prior-art checking essential. No original full solution is promoted.


\entrylabel{Retained derivation or attempt}

Choose a shortest generator word for each group element x,\allowbreak{} and symmetrize these random word choices under x going to x\textasciicircum{}-1 and word reversal with inverse letters. Put a\_\allowbreak{}g=\allowbreak{}E\_\allowbreak{}x[length(word\_\allowbreak{}x)*number\_\allowbreak{}of\_\allowbreak{}g\_\allowbreak{}letters(word\_\allowbreak{}x)],\allowbreak{} x uniform,\allowbreak{} and A=\allowbreak{}sum\_\allowbreak{}g a\_\allowbreak{}g=\allowbreak{}E\_\allowbreak{}x length(word\_\allowbreak{}x)\textasciicircum{}2<=\allowbreak{}D\textasciicircum{}2. Then a\_\allowbreak{}g=\allowbreak{}a\_\allowbreak{}(g\textasciicircum{}-1),\allowbreak{} and mu(g)=\allowbreak{}a\_\allowbreak{}g/\allowbreak{}A is a symmetric probability law. The source comparison theorem with the uniform jump-to-any-group-element chain,\allowbreak{} whose relaxation time is1,\allowbreak{} gives tau2(mu)<=\allowbreak{}max\_\allowbreak{}g a\_\allowbreak{}g/\allowbreak{}mu(g)=\allowbreak{}A<=\allowbreak{}D\textasciicircum{}2.


\entrylabel{Status limit}

No mathematical obstruction remains in this finite-state question;\allowbreak{} the source question was already resolved by later literature. Independent proof review not claimed.


\subsection*{P0249  Trace inequality for additive reversibilization}

Aldous-Fill:\allowbreak{} Reversible Markov Chains and Random Walks on Graphs,\allowbreak{} 2002 recompiled2014;\allowbreak{} Conjecture9.22 and Corollary9.24,\allowbreak{}PDF327-329


\url{https://www.stat.berkeley.edu/~aldous/RWG/book.pdf}


\entrylabel{current status}

resolved\_\allowbreak{}in\_\allowbreak{}literature


\entrylabel{result status}

excluded\_\allowbreak{}prior\_\allowbreak{}resolution


\entrylabel{statement}

For an irreducible finite Markov chain with fundamental matrix Z and time reversal P*,\allowbreak{} prove or refute trace[Z\textasciicircum{}2*(P*-P)]>=\allowbreak{}0.


\entrylabel{statement caveat}

This concerns the mean-zero restriction of the fundamental matrix. A full analytic candidate is given,\allowbreak{} but the broad nonreversible variance-reduction literature has not been exhaustively reconciled.


\entrylabel{resolution evidence}

Huang-Mao2020 Theorem4.3 implies the trace monotonicity by summing asymptotic variances over a mean-zero orthonormal basis;\allowbreak{} one-half lazification removes periodicity. Exact reduction written in proof packet.


\entrylabel{resolution url}

\url{https://arxiv.org/abs/2012.13895}


\entrylabel{Later primary source}

Huang-Mao:\allowbreak{} Variational formulas of asymptotic variance for general discrete-time Markov chains;\allowbreak{} 


\url{https://arxiv.org/pdf/2012.13895}


\entrylabel{Retained derivation or attempt}

Work in real L2(pi) on the mean-zero subspace. Write A=\allowbreak{}I-P=\allowbreak{}H+\allowbreak{}K with H positive definite self-adjoint and K skew-adjoint. Put Q=\allowbreak{}H\textasciicircum{}-1/\allowbreak{}2*K*H\textasciicircum{}-1/\allowbreak{}2 and B=\allowbreak{}-Q\textasciicircum{}2>=\allowbreak{}0. For f(t)=\allowbreak{}trace[(H+\allowbreak{}tK)\textasciicircum{}-1],\allowbreak{} skew-adjointness gives f(t)=\allowbreak{}trace[H\textasciicircum{}-1*(I+\allowbreak{}t\textasciicircum{}2*B)\textasciicircum{}-1],\allowbreak{} because only the symmetric part of the inverse contributes to the trace. Differentiation yields fprime(t)=\allowbreak{}-2t*trace[H\textasciicircum{}-1*B*(I+\allowbreak{}t\textasciicircum{}2*B)\textasciicircum{}-2]<=\allowbreak{}0 for t>=\allowbreak{}0. On the other hand fprime(1)=\allowbreak{}-trace[A\textasciicircum{}-2*K]. Since P*-P=\allowbreak{}2K and Z=\allowbreak{}A\textasciicircum{}-1 on mean-zero vectors,\allowbreak{} the requested trace is -2*fprime(1)>=\allowbreak{}0.


\entrylabel{Status limit}

The proof is complete as a finite-dimensional matrix argument,\allowbreak{} but no independent review or novelty certificate has been obtained. Earlier additive-reversibilization and asymptotic-variance results may already imply it. Preserve candidate status until that is checked.


\entrylabel{Retained derivation or attempt}

Work in real L2(pi) on the mean-zero subspace. Write A=\allowbreak{}I-P=\allowbreak{}H+\allowbreak{}K with H positive definite self-adjoint and K skew-adjoint. Put Q=\allowbreak{}H\textasciicircum{}-1/\allowbreak{}2*K*H\textasciicircum{}-1/\allowbreak{}2 and B=\allowbreak{}-Q\textasciicircum{}2>=\allowbreak{}0. For f(t)=\allowbreak{}trace[(H+\allowbreak{}tK)\textasciicircum{}-1],\allowbreak{} skew-adjointness gives f(t)=\allowbreak{}trace[H\textasciicircum{}-1*(I+\allowbreak{}t\textasciicircum{}2*B)\textasciicircum{}-1],\allowbreak{} because only the symmetric part of the inverse contributes to the trace. Differentiation yields fprime(t)=\allowbreak{}-2t*trace[H\textasciicircum{}-1*B*(I+\allowbreak{}t\textasciicircum{}2*B)\textasciicircum{}-2]<=\allowbreak{}0 for t>=\allowbreak{}0. On the other hand fprime(1)=\allowbreak{}-trace[A\textasciicircum{}-2*K]. Since P*-P=\allowbreak{}2K and Z=\allowbreak{}A\textasciicircum{}-1 on mean-zero vectors,\allowbreak{} the requested trace is -2*fprime(1)>=\allowbreak{}0.


\entrylabel{Status limit}

No mathematical obstruction remains in this finite-state question;\allowbreak{} the source question was already resolved by later literature. Independent proof review not claimed.


\subsection*{P0260  Moving-target maximal hitting-time equality}

Aldous-Fill:\allowbreak{} Reversible Markov Chains and Random Walks on Graphs,\allowbreak{} 2002 recompiled2014;\allowbreak{} Open Problem4.20,\allowbreak{}PDF129


\url{https://www.stat.berkeley.edu/~aldous/RWG/book.pdf}


\entrylabel{current status}

resolved\_\allowbreak{}in\_\allowbreak{}literature


\entrylabel{result status}

excluded\_\allowbreak{}prior\_\allowbreak{}resolution


\entrylabel{resolution evidence}

Sousi-Winkler2014 Theorem1.7 provides a transitive counterexample for both continuous-time and lazy walks.


\entrylabel{resolution url}

\url{https://arxiv.org/abs/1210.5236}


\subsection*{P0261  Averaged versus continuous-time mixing}

Aldous-Fill:\allowbreak{} Reversible Markov Chains and Random Walks on Graphs,\allowbreak{} 2002 recompiled2014;\allowbreak{} Open Problem4.17,\allowbreak{}PDF127


\url{https://www.stat.berkeley.edu/~aldous/RWG/book.pdf}


\entrylabel{current status}

resolved\_\allowbreak{}in\_\allowbreak{}literature


\entrylabel{result status}

excluded\_\allowbreak{}prior\_\allowbreak{}resolution


\entrylabel{resolution evidence}

Hermon-Peres2017 proves the requested asymptotic comparison.


\entrylabel{resolution url}

\url{https://arxiv.org/abs/1508.04836}


\subsection*{P0262  Isoperimetric exponent greater than one half implies linear commute time}

Aldous-Fill:\allowbreak{} Reversible Markov Chains and Random Walks on Graphs,\allowbreak{} 2002 recompiled2014;\allowbreak{} Open Problem6.20,\allowbreak{}PDF224


\url{https://www.stat.berkeley.edu/~aldous/RWG/book.pdf}


\entrylabel{current status}

resolved\_\allowbreak{}in\_\allowbreak{}literature


\entrylabel{result status}

excluded\_\allowbreak{}prior\_\allowbreak{}resolution


\entrylabel{resolution evidence}

Benjamini-Kozma2005 resolves the stated Aldous-Fill mean commute-time problem through an isoperimetric resistance bound.


\entrylabel{resolution url}

\url{https://arxiv.org/abs/math/0212322}


\subsection*{P0263  Exact logarithmic Sobolev constant of finite cycles}

Aldous-Fill:\allowbreak{} Reversible Markov Chains and Random Walks on Graphs,\allowbreak{} 2002 recompiled2014;\allowbreak{} Open Problem8.33,\allowbreak{}PDF301


\url{https://www.stat.berkeley.edu/~aldous/RWG/book.pdf}


\entrylabel{current status}

resolved\_\allowbreak{}in\_\allowbreak{}literature


\entrylabel{result status}

excluded\_\allowbreak{}prior\_\allowbreak{}resolution


\entrylabel{resolution evidence}

Frank-Ivanisvili May2026 proves all n>=\allowbreak{}4;\allowbreak{} journal article scheduled October2026. The preprint predates this campaign.


\entrylabel{resolution url}

\url{https://arxiv.org/abs/2605.29035}


\subsection*{P0264  Sufficient normal-limit conditions for the antivoter model}

Aldous-Fill:\allowbreak{} Reversible Markov Chains and Random Walks on Graphs,\allowbreak{} 2002 recompiled2014;\allowbreak{} Open Problem14.25,\allowbreak{}PDF478


\url{https://www.stat.berkeley.edu/~aldous/RWG/book.pdf}


\entrylabel{current status}

resolved\_\allowbreak{}in\_\allowbreak{}literature


\entrylabel{result status}

excluded\_\allowbreak{}prior\_\allowbreak{}resolution


\entrylabel{resolution evidence}

Rinott-Rotar1997 gives CLT conditions and rates;\allowbreak{} Rollin2007 obtains stronger translated-Poisson approximation using the same setting.


\entrylabel{resolution url}

\url{https://doi.org/10.1214/aoap/1043862425}


\subsection*{P0265  Universal coalescence bound by maximal hitting time}

Aldous-Fill:\allowbreak{} Reversible Markov Chains and Random Walks on Graphs,\allowbreak{} 2002 recompiled2014;\allowbreak{} Open Problem14.13,\allowbreak{}PDF470


\url{https://www.stat.berkeley.edu/~aldous/RWG/book.pdf}


\entrylabel{current status}

resolved\_\allowbreak{}in\_\allowbreak{}literature


\entrylabel{result status}

excluded\_\allowbreak{}prior\_\allowbreak{}resolution


\entrylabel{resolution evidence}

Oliveira2012 establishes the universal-constant coalescence bound for reversible chains.


\entrylabel{resolution url}

\url{https://arxiv.org/abs/1009.0664}


\subsection*{P0304  Telek two-matrix RAP to MAP conjecture}

Peralta:\allowbreak{} Rational arrival processes with strictly positive densities need not be Markovian (2026);\allowbreak{} Conjecture1,\allowbreak{}PDF2;\allowbreak{}Theorem9,\allowbreak{}PDF9


\url{https://arxiv.org/pdf/2603.28047}


\entrylabel{current status}

resolved\_\allowbreak{}in\_\allowbreak{}literature


\entrylabel{result status}

literature\_\allowbreak{}counterexample


\entrylabel{statement}

Must a rational arrival process with strictly positive joint densities and a unique dominant real eigenvalue of its no-arrival matrix admit a finite Markovian arrival-process representation?


\entrylabel{statement caveat}

The2026 primary counterexample paper restates Telek2022 and refutes it. Excluded from the open target.


\entrylabel{resolution evidence}

Complete10-page primary preprint read;\allowbreak{} positive-density bound by summation by parts,\allowbreak{} normalization and irrational peripheral-pole obstruction checked. Not a new result here;\allowbreak{} independent human review not claimed.


\subsection*{P0319  Old Riesz number-rigidity threshold at s=\allowbreak{}d-1}

Lewin:\allowbreak{} Coulomb and Riesz gases,\allowbreak{} the known and the unknown,\allowbreak{}2022;\allowbreak{} SectionVI.C,\allowbreak{}PDF70;\allowbreak{}Thoma2023 Introduction PDF2


\url{https://arxiv.org/pdf/2202.09240}


\entrylabel{current status}

claimed\_\allowbreak{}resolved\_\allowbreak{}in\_\allowbreak{}literature


\entrylabel{result status}

literature\_\allowbreak{}counterexample\_\allowbreak{}claim\_\allowbreak{}not\_\allowbreak{}full\_\allowbreak{}proof\_\allowbreak{}reviewed


\entrylabel{statement}

The historical conjecture predicts number rigidity for Riesz gases at and below the threshold s=\allowbreak{}d-1.


\entrylabel{statement caveat}

Thoma explicitly reports disproving this conjecture by non-rigid Coulomb gas limits in d>=\allowbreak{}3. Excluded from the open target;\allowbreak{} the full15-page proof has not yet been read in this campaign.


\entrylabel{resolution evidence}

Primary introduction and Theorem1 read,\allowbreak{} plus2026 journal abstract;\allowbreak{} full proof not yet checked. New remaining interval recorded as P0320.


\subsection*{P0352  Maximum continuous height of the critical tree}

Aldous-Janson:\allowbreak{} Critical beta-splitting random tree II,\allowbreak{} overview and open problems;\allowbreak{} Open Problem3,\allowbreak{}PDF24


\url{https://arxiv.org/pdf/2303.02529}


\entrylabel{current status}

resolution\_\allowbreak{}claimed\_\allowbreak{}in\_\allowbreak{}literature


\entrylabel{result status}

literature\_\allowbreak{}resolution\_\allowbreak{}claimed\_\allowbreak{}not\_\allowbreak{}independently\_\allowbreak{}certified


\entrylabel{statement}

Maximum continuous height of the critical tree


\entrylabel{statement caveat}

The August18,\allowbreak{}2026 preprint proves a law of large numbers with constant2 and a mixed Gumbel law,\allowbreak{} contradicting the older suggested constant1.878. The complete maximum-height proof was read;\allowbreak{} its martingale integrability result uses an external theorem not independently rechecked here.


\entrylabel{Later primary source}

Chorbadzhiyska-Minchev-Savov:\allowbreak{} Beta-splitting trees via homogeneous fragmentations;\allowbreak{} Theorem2.12 and Remark2.13,\allowbreak{}PDF10-11;\allowbreak{}complete maximum-height argumentPDF20-25


\url{https://arxiv.org/pdf/2608.18320}


\subsection*{P0353  Fixed-clade count CLTs for critical beta-splitting trees}

Aldous-Janson:\allowbreak{} Critical beta-splitting random tree II,\allowbreak{} overview and open problems;\allowbreak{} Open Problem5,\allowbreak{}PDF37


\url{https://arxiv.org/pdf/2303.02529}


\entrylabel{current status}

resolution\_\allowbreak{}claimed\_\allowbreak{}in\_\allowbreak{}literature


\entrylabel{result status}

literature\_\allowbreak{}resolution\_\allowbreak{}claimed\_\allowbreak{}not\_\allowbreak{}independently\_\allowbreak{}certified


\entrylabel{statement}

Fixed-clade count CLTs for critical beta-splitting trees


\entrylabel{statement caveat}

The July2025 preprint explicitly answers the fixed-clade count question,\allowbreak{} with linear variance and a CLT for nontrivial clades. The proof text was read;\allowbreak{} several initial-index statements need care,\allowbreak{} including the literal variance-zero range for a size3 clade. This is a source-claimed resolution exclusion,\allowbreak{} not an independently certified proof.


\entrylabel{Later primary source}

Brandenberger-Chin-Mossel:\allowbreak{} Asymptotics for the harmonic descent chain and critical beta-splitting trees;\allowbreak{} Corollary1.4,\allowbreak{}PDF4;\allowbreak{}Sections3-4,\allowbreak{}PDF13-17


\url{https://arxiv.org/pdf/2505.24821}


\subsection*{P0354  Total continuous-time tree-length variance and CLT}

Aldous-Janson:\allowbreak{} Critical beta-splitting random tree II,\allowbreak{} overview and open problems;\allowbreak{} Open Problem6,\allowbreak{}PDF39


\url{https://arxiv.org/pdf/2303.02529}


\entrylabel{current status}

resolution\_\allowbreak{}claimed\_\allowbreak{}in\_\allowbreak{}literature


\entrylabel{result status}

literature\_\allowbreak{}resolution\_\allowbreak{}claimed\_\allowbreak{}not\_\allowbreak{}independently\_\allowbreak{}certified


\entrylabel{statement}

Total continuous-time tree-length variance and CLT


\entrylabel{statement caveat}

The July2025 preprint explicitly answers this question with linear variance and a CLT. The main argument and supporting displayed mean/\allowbreak{}variance estimates were read,\allowbreak{} but not every referenced contraction theorem or appendix independently verified. This is a source-claimed resolution exclusion,\allowbreak{} not a new campaign solution.


\entrylabel{Later primary source}

Brandenberger-Chin-Mossel:\allowbreak{} Asymptotics for the harmonic descent chain and critical beta-splitting trees;\allowbreak{} Corollary1.5,\allowbreak{}PDF4;\allowbreak{}Section5,\allowbreak{}PDF17-18


\url{https://arxiv.org/pdf/2505.24821}


\subsection*{P0407  FCFS characterization for minimum mean slowdown}

Aalto:\allowbreak{} Minimizing the mean slowdown in a single-server queue;\allowbreak{} Conjecture1;\allowbreak{} PDF4


\url{https://acris.aalto.fi/ws/portalfiles/portal/79466983/Aalto_Queueing_Systems_2022.pdf}


\entrylabel{current status}

resolved\_\allowbreak{}in\_\allowbreak{}literature


\entrylabel{result status}

known\_\allowbreak{}result\_\allowbreak{}not\_\allowbreak{}original


\entrylabel{statement}

FCFS minimizes mean slowdown among non-anticipating policies if and only if Htilde(x)>=\allowbreak{}Htilde(0) for every x.


\entrylabel{resolution evidence}

The later formula for H\_\allowbreak{}sld is identical to the source Htilde. Corollary3 states the iff claim and explicitly identifies the earlier conjecture. No new solution claimed.


\entrylabel{Later primary source}

Aalto and Scully:\allowbreak{} Minimizing the mean slowdown in the M/\allowbreak{}G/\allowbreak{}1 queue;\allowbreak{} Corollary3 equation18;\allowbreak{} PDF13


\url{https://aaltodoc.aalto.fi/server/api/core/bitstreams/a9d987aa-1806-45f5-a7be-7bce2d8cc8f0/content}


\subsection*{P0423  Ideal Voronoi face-volume intensities}

D Achille Curien Enriquez Lyons Unel:\allowbreak{} Ideal Poisson Voronoi tessellations on hyperbolic spaces;\allowbreak{} Question7.4,\allowbreak{} density portion;\allowbreak{} PDF48


\url{https://arxiv.org/pdf/2303.16831}


\entrylabel{current status}

resolved\_\allowbreak{}in\_\allowbreak{}literature


\entrylabel{result status}

known\_\allowbreak{}result\_\allowbreak{}not\_\allowbreak{}original


\entrylabel{statement}

Compute all ideal k-face volume densities.


\entrylabel{resolution evidence}

The theorem computes every density and Remark2.2 explicitly identifies D\_\allowbreak{}d,\allowbreak{}k with the original tilde-I\_\allowbreak{}d,\allowbreak{}k notation. The separate hole probabilities remain P0410.


\entrylabel{Later primary source}

D Achille and Thale:\allowbreak{} Face volume densities of positive intensity and ideal Poisson Voronoi tessellations in hyperbolic spaces;\allowbreak{} Corollary2.1 and Remark2.2;\allowbreak{} PDF6


\url{https://arxiv.org/pdf/2606.26049}


\subsection*{P0424  Indistinguishability of hyperbolic ideal Voronoi cells}

D Achille Curien Enriquez Lyons Unel:\allowbreak{} Ideal Poisson Voronoi tessellations on hyperbolic spaces;\allowbreak{} Question7.8;\allowbreak{} PDF48-49


\url{https://arxiv.org/pdf/2303.16831}


\entrylabel{current status}

resolved\_\allowbreak{}in\_\allowbreak{}literature


\entrylabel{result status}

known\_\allowbreak{}result\_\allowbreak{}not\_\allowbreak{}original


\entrylabel{statement}

Prove that ideal hyperbolic Voronoi cells admit no nontrivial equivariant selection.


\entrylabel{resolution evidence}

The source update explicitly credits Mellick. His Theorem1 includes real hyperbolic symmetric spaces and states indistinguishability.


\entrylabel{Later primary source}

Mellick:\allowbreak{} Indistinguishability of cells for the ideal Poisson Voronoi tessellation;\allowbreak{} Theorem1;\allowbreak{} PDF1


\url{https://arxiv.org/pdf/2408.10009}


\subsection*{P0448  Offline common-concave speedup allocation}

Harchol-Balter:\allowbreak{} Open problems in queueing theory inspired by datacenter computing;\allowbreak{} Section3.3 PDF10


\url{https://www.cs.cmu.edu/~harchol/Papers/QUESTA21.pdf}


\entrylabel{current status}

resolved\_\allowbreak{}in\_\allowbreak{}literature


\entrylabel{result status}

known\_\allowbreak{}result\_\allowbreak{}not\_\allowbreak{}original


\entrylabel{statement}

Offline common-concave speedup allocation


\entrylabel{resolution evidence}

The later paper provides SmartFill for known jobs all present at zero with a common strictly concave differentiable increasing speedup and compatible weights. This resolves the main offline smooth common-curve subproblem;\allowbreak{} no Poisson-arrival extension is claimed.


\entrylabel{Later primary source}

Li Ju Eryilmaz Shroff:\allowbreak{} Optimal Parallel Scheduling under Concave Speedup Functions;\allowbreak{} Sections2 and5;\allowbreak{} PDF3-4,\allowbreak{}14-15


\url{https://arxiv.org/pdf/2509.01811}


\subsection*{P0449  General FCFS multiserver-job stability representation}

Harchol-Balter:\allowbreak{} Open problems in queueing theory inspired by datacenter computing;\allowbreak{} Section2 PDF5


\url{https://www.cs.cmu.edu/~harchol/Papers/QUESTA21.pdf}


\entrylabel{current status}

resolved\_\allowbreak{}in\_\allowbreak{}literature


\entrylabel{result status}

known\_\allowbreak{}result\_\allowbreak{}not\_\allowbreak{}original


\entrylabel{statement}

General FCFS multiserver-job stability representation


\entrylabel{resolution evidence}

The later theorem gives lambda\_\allowbreak{}c=\allowbreak{}1/\allowbreak{}gamma using the saturated pile growth rate for the stated stationary input framework. This credits the representation;\allowbreak{} evaluation complexity remains P0444 and no blanket critical-boundary claim is added.


\entrylabel{Later primary source}

The Multiserver Job Stochastic Recurrence Equation for Cloud Computing Performance Evaluation;\allowbreak{} Theorem5.1;\allowbreak{} PDF11


\url{https://arxiv.org/pdf/2601.20653}


\subsection*{P0454  Iterated budget-filling bilateral graph limits}

Moharrami,\allowbreak{} Subramanian,\allowbreak{} Liu and Sundaresan:\allowbreak{} The Erlang Weighted Tree,\allowbreak{} A New Branching Process (v5,\allowbreak{} 27 May 2023);\allowbreak{} Section4 question5;\allowbreak{} PDF17


\url{https://arxiv.org/pdf/2002.03993}


\entrylabel{current status}

intake\_\allowbreak{}needs\_\allowbreak{}precise\_\allowbreak{}formulation


\entrylabel{result status}

not\_\allowbreak{}attempted\_\allowbreak{}as\_\allowbreak{}precise\_\allowbreak{}problem


\entrylabel{statement}

Iterated budget-filling bilateral graph limits


\entrylabel{reason}

The proposal requires an exact round/\allowbreak{}update protocol. Retain as an intake lead until that model is fixed;\allowbreak{} no extra admission for an arbitrary invented protocol.


\subsection*{P0455  A general spectral method for uncountable-type branching}

Moharrami,\allowbreak{} Subramanian,\allowbreak{} Liu and Sundaresan:\allowbreak{} The Erlang Weighted Tree,\allowbreak{} A New Branching Process (v5,\allowbreak{} 27 May 2023);\allowbreak{} Section4 question6;\allowbreak{} PDF17


\url{https://arxiv.org/pdf/2002.03993}


\entrylabel{current status}

intake\_\allowbreak{}needs\_\allowbreak{}precise\_\allowbreak{}formulation


\entrylabel{result status}

not\_\allowbreak{}attempted\_\allowbreak{}as\_\allowbreak{}precise\_\allowbreak{}problem


\entrylabel{statement}

A general spectral method for uncountable-type branching


\entrylabel{reason}

Methodological research direction without a sharp general target;\allowbreak{} hold rather than count a broad extension request as another precise problem.


\subsection*{P0460  Leading logarithmic connectivity threshold of bilateral preference graphs}

Dabirian and Subramanian:\allowbreak{} Connectivity of a family of bilateral preference random graphs;\allowbreak{} Introduction;\allowbreak{} Theorem1.1 and Theorems4.1,\allowbreak{}5.3


\url{https://arxiv.org/pdf/2209.05677}


\entrylabel{current status}

resolved\_\allowbreak{}in\_\allowbreak{}literature


\entrylabel{result status}

known\_\allowbreak{}result\_\allowbreak{}not\_\allowbreak{}original


\entrylabel{statement}

Determine the leading k=\allowbreak{}t*log(n) connectivity threshold originally conjectured by La and Kabkab.


\entrylabel{resolution evidence}

The archived v3 paper proves disconnectedness below log(n)-0.5*sqrt(log(n))*log(log(n)) and connectivity above the corresponding upper band,\allowbreak{} with strict inequalities in the coefficient. It settles the leading threshold;\allowbreak{} P0456 retains only the fine unresolved window.


\entrylabel{Later primary source}

Dabirian and Subramanian:\allowbreak{} Connectivity of a family of bilateral preference random graphs;\allowbreak{} 


\url{https://arxiv.org/pdf/2209.05677}


\subsection*{R0001  Complete bunkbed inequality}

Selected problems in probability theory;\allowbreak{} Conjecture 4.1,\allowbreak{} PDF page 8


\url{https://arxiv.org/pdf/2205.07318}


\entrylabel{current status}

resolved\_\allowbreak{}in\_\allowbreak{}literature


\entrylabel{result status}

literature\_\allowbreak{}counterexample


\entrylabel{resolution url}

\url{https://www.math.ucla.edu/~pak/papers/BunkBed4.pdf}


\subsection*{R0002  One-dimensional Brownian snails extinction}

Selected problems in probability theory;\allowbreak{} Question 6.1(i),\allowbreak{} PDF page 10


\url{https://arxiv.org/pdf/2205.07318}


\entrylabel{current status}

resolved\_\allowbreak{}in\_\allowbreak{}literature


\entrylabel{result status}

literature\_\allowbreak{}solution


\entrylabel{resolution url}

\url{https://www.normalesup.org/~hartarsky/files/research/publi/snails.pdf}


\subsection*{R0003  Schramm locality for critical percolation}

Probability on Trees and Networks,\allowbreak{} paperback electronic version 19 August 2026;\allowbreak{} Conjecture 7.45


\url{https://rdlyons.pages.iu.edu/prbtree/book_pb.pdf}


\entrylabel{current status}

resolved\_\allowbreak{}in\_\allowbreak{}literature


\entrylabel{result status}

literature\_\allowbreak{}resolution


\entrylabel{resolution evidence}

Theorem 1.1;\allowbreak{} abstract checked;\allowbreak{} 2023 preprint


\entrylabel{resolution url}

\url{https://arxiv.org/abs/2310.10983}


\subsection*{R0004  Generating-set invariance of free uniform-forest connectivity}

Probability on Trees and Networks,\allowbreak{} paperback electronic version 19 August 2026;\allowbreak{} Question 10.50


\url{https://rdlyons.pages.iu.edu/prbtree/book_pb.pdf}


\entrylabel{current status}

resolved\_\allowbreak{}in\_\allowbreak{}literature


\entrylabel{result status}

literature\_\allowbreak{}resolution


\entrylabel{resolution evidence}

Abstract and author institutional version;\allowbreak{} Annals of Probability 2022


\entrylabel{resolution url}

\url{https://arxiv.org/abs/2006.06387}


\subsection*{R0005  Analytic biased-walk speed on leafless Galton-Watson trees}

Probability on Trees and Networks,\allowbreak{} paperback electronic version 19 August 2026;\allowbreak{} Question 17.34


\url{https://rdlyons.pages.iu.edu/prbtree/book_pb.pdf}


\entrylabel{current status}

resolved\_\allowbreak{}in\_\allowbreak{}literature


\entrylabel{result status}

literature\_\allowbreak{}resolution


\entrylabel{resolution evidence}

Question 17.34 immediately following sentence;\allowbreak{} original theorem conditions still to verify


\entrylabel{resolution url}

\url{https://rdlyons.pages.iu.edu/prbtree/book_pb.pdf}


\subsection*{R0006  Harmonic dimension lower bound by expected log offspring}

Probability on Trees and Networks,\allowbreak{} paperback electronic version 19 August 2026;\allowbreak{} Question 17.28


\url{https://rdlyons.pages.iu.edu/prbtree/book_pb.pdf}


\entrylabel{current status}

resolved\_\allowbreak{}in\_\allowbreak{}literature


\entrylabel{result status}

literature\_\allowbreak{}resolution


\entrylabel{resolution evidence}

Theorem 1;\allowbreak{} primary theorem statement checked


\entrylabel{resolution url}

\url{https://perso.lpsm.paris/~slin/papers/harmonic_bias.pdf}


\subsection*{R0007  Exact Gaussian ellipsoid-fitting quarter threshold}

Randomstrasse101:\allowbreak{} Open Problems of 2024;\allowbreak{} Conjecture6,\allowbreak{} PDF pp.4-5


\url{https://arxiv.org/pdf/2504.20539v1}


\entrylabel{current status}

claimed\_\allowbreak{}resolution\_\allowbreak{}in\_\allowbreak{}recent\_\allowbreak{}primary\_\allowbreak{}preprints


\entrylabel{result status}

literature\_\allowbreak{}claim\_\allowbreak{}excluded\_\allowbreak{}pending\_\allowbreak{}review


\entrylabel{resolution evidence}

Misiakiewicz-Wen,\allowbreak{} Theorem1.2,\allowbreak{} PDF p.2,\allowbreak{} both exact feasibility below1/\allowbreak{}4 and unrestricted infeasibility above1/\allowbreak{}4;\allowbreak{} statement read,\allowbreak{} full47-page proof not reviewed


\entrylabel{resolution url}

\url{https://arxiv.org/pdf/2608.10184}


\entrylabel{additional resolution claims}

["\url{https://arxiv.org/abs/2608.12415",} "\url{https://arxiv.org/abs/2608.27372"]}


\clearpage

\section*{Primary-source archive registry}

\subsection*{Agresti-Sauerbrey-Veraar stochastic flow De Giorgi Nash Moser,\allowbreak{}2025}

agresti\_\allowbreak{}flow\_\allowbreak{}dgnm\_\allowbreak{}2025


\url{https://arxiv.org/pdf/2511.12692}


\entrylabel{access evidence}

Primary author/\allowbreak{}arXiv source used for later-status reconciliation,\allowbreak{} not automatic current-open certification


\entrylabel{reading status}

unread


\entrylabel{statement validation}

not started


\entrylabel{sha256}

e78da6d5\allowbreak{}18def19a\allowbreak{}4a350f5b\allowbreak{}3904e85f\allowbreak{}218bc2ab\allowbreak{}15e0ade6\allowbreak{}e20949a4\allowbreak{}5279a588


\entrylabel{acquired at}

2026-09-07T00:\allowbreak{}00:\allowbreak{}18Z


\subsection*{Agresti:\allowbreak{} Fractal dimension of singular times for SPDEs,\allowbreak{}2026}

agresti\_\allowbreak{}singular\_\allowbreak{}times\_\allowbreak{}2026


\url{https://arxiv.org/pdf/2602.21981}


\entrylabel{access evidence}

Primary later-literature check for endpoint energy questions


\entrylabel{reading status}

unread


\entrylabel{statement validation}

not started


\entrylabel{sha256}

c3b286ac\allowbreak{}85b040bd\allowbreak{}146e41ae\allowbreak{}dc29772c\allowbreak{}4b491f00\allowbreak{}8cc1c27b\allowbreak{}47c02726\allowbreak{}6e6109e2


\entrylabel{acquired at}

2026-09-06T23:\allowbreak{}44:\allowbreak{}24Z


\subsection*{Agresti-Veraar temporal regularity threshold,\allowbreak{}2026}

agresti\_\allowbreak{}temporal\_\allowbreak{}threshold\_\allowbreak{}2026


\url{https://arxiv.org/pdf/2509.07803}


\entrylabel{access evidence}

Primary author/\allowbreak{}arXiv source used for later-status reconciliation,\allowbreak{} not automatic current-open certification


\entrylabel{reading status}

unread


\entrylabel{statement validation}

not started


\entrylabel{sha256}

b837690c\allowbreak{}5077999a\allowbreak{}726ebe4e\allowbreak{}8037dddf\allowbreak{}7f17168c\allowbreak{}ea0099d1\allowbreak{}cb44f685\allowbreak{}f2c00e29


\entrylabel{acquired at}

2026-09-07T00:\allowbreak{}00:\allowbreak{}17Z


\subsection*{Agresti-Veraar:\allowbreak{} Nonlinear SPDEs and Maximal Regularity,\allowbreak{}2025 extended survey}

agresti\_\allowbreak{}veraar\_\allowbreak{}spde\_\allowbreak{}2025


\url{https://arxiv.org/pdf/2501.18561}


\entrylabel{access evidence}

Primary authors arXiv survey;\allowbreak{} final open-problem section


\entrylabel{reading status}

open problems section read PDF91 94 and relevant theorem checks


\entrylabel{statement validation}

admitted statements paraphrased against numbered questions


\entrylabel{sha256}

c2ef6839\allowbreak{}e5c3c78a\allowbreak{}9c72826b\allowbreak{}22cb509c\allowbreak{}9cb8af8a\allowbreak{}0753b705\allowbreak{}4a025217\allowbreak{}0cf9dae7


\entrylabel{acquired at}

2026-09-06T23:\allowbreak{}40:\allowbreak{}06Z


\subsection*{Agresti-Veraar:\allowbreak{} Nonlinear SPDEs and maximal regularity,\allowbreak{}published NoDEA2025}

agresti\_\allowbreak{}veraar\_\allowbreak{}spde\_\allowbreak{}journal\_\allowbreak{}2025


\url{https://repository.tudelft.nl/file/File_783e6459-5d3b-4957-bcd3-64da5b8d9025}


\entrylabel{access evidence}

Primary institutional repository published article;\allowbreak{} endpoint energy question source check


\entrylabel{reading status}

unread


\entrylabel{statement validation}

not started


\entrylabel{sha256}

a66ffa2c\allowbreak{}1d88f2b4\allowbreak{}6d563a3e\allowbreak{}4d3963bb\allowbreak{}c0f86234\allowbreak{}92e7e183\allowbreak{}7b4ae089\allowbreak{}008ad091


\entrylabel{acquired at}

2026-09-06T23:\allowbreak{}44:\allowbreak{}20Z


\subsection*{Ahlberg:\allowbreak{} Tertiles and the time constant (2019;\allowbreak{}published2020)}

ahlberg\_\allowbreak{}tertiles\_\allowbreak{}2020


\url{https://arxiv.org/pdf/1911.09949}


\entrylabel{access evidence}

Public primary arXiv manuscript;\allowbreak{} both pages including complete proof and counterexample read


\entrylabel{reading status}

unread


\entrylabel{statement validation}

not started


\entrylabel{sha256}

202f4455\allowbreak{}544198fc\allowbreak{}1a7cd832\allowbreak{}f4264ab0\allowbreak{}f8184d23\allowbreak{}f00e706a\allowbreak{}058ecaf5\allowbreak{}e327538b


\entrylabel{acquired at}

2026-09-06T19:\allowbreak{}17:\allowbreak{}14Z


\subsection*{AIM braid groups,\allowbreak{} clusters and free probability problems}

aim\_\allowbreak{}braid\_\allowbreak{}free\_\allowbreak{}probability


\url{https://aimath.org/WWN/braidgroups/braidgroups.pdf}


\entrylabel{access evidence}

Public primary author or workshop PDF;\allowbreak{} individual open-question statements and later status require separate validation


\entrylabel{reading status}

unread


\entrylabel{statement validation}

not started


\entrylabel{sha256}

88f54825\allowbreak{}edb11953\allowbreak{}2be13ece\allowbreak{}404988f8\allowbreak{}ed82dfef\allowbreak{}4716d243\allowbreak{}97e11074\allowbreak{}3c19c1e9


\entrylabel{acquired at}

2026-09-06T19:\allowbreak{}34:\allowbreak{}20Z


\subsection*{AIM free analysis problems,\allowbreak{} August2006}

aim\_\allowbreak{}free\_\allowbreak{}analysis\_\allowbreak{}2006


\url{https://aimath.org/WWN/freeanalysis/freeanalysis.pdf}


\entrylabel{access evidence}

Public primary author or workshop PDF;\allowbreak{} individual open-question statements and later status require separate validation


\entrylabel{reading status}

unread


\entrylabel{statement validation}

not started


\entrylabel{sha256}

a44e98be\allowbreak{}34192d10\allowbreak{}7d730b5b\allowbreak{}3f4b892d\allowbreak{}d4640563\allowbreak{}123b5097\allowbreak{}2d5b58ad\allowbreak{}a54d5b47


\entrylabel{acquired at}

2026-09-06T19:\allowbreak{}33:\allowbreak{}56Z


\subsection*{AIM All roads to KPZ universality:\allowbreak{} workshop open problems,\allowbreak{}March19 2025}

aim\_\allowbreak{}kpz\_\allowbreak{}2025


\url{https://lpetrov.cc/LaTeX/2025-03-19-AIM-KPZ/2025-03-19-AIM-KPZ.pdf}


\entrylabel{access evidence}

Organizer-hosted primary workshop problem list with named moderator and scribe


\entrylabel{reading status}

all question sections and references read


\entrylabel{statement validation}

seven concrete questions cross checked with followup


\entrylabel{sha256}

36276135\allowbreak{}20289f96\allowbreak{}22f9ab0c\allowbreak{}61a9b38d\allowbreak{}4fa21b37\allowbreak{}99fd4070\allowbreak{}c7b458c0\allowbreak{}ad99bdf6


\entrylabel{acquired at}

2026-09-06T23:\allowbreak{}41:\allowbreak{}31Z


\subsection*{AIM all roads to KPZ workshop follow-up report,\allowbreak{}2025}

aim\_\allowbreak{}kpz\_\allowbreak{}followup\_\allowbreak{}2025


\url{https://aimath.org/pastworkshops/roadtokpzrep.pdf}


\entrylabel{access evidence}

Primary workshop follow-up and author paper,\allowbreak{} bounded current-status reconciliation


\entrylabel{reading status}

complete six page followup read


\entrylabel{statement validation}

explicit no progress and model scope distinguished


\entrylabel{sha256}

cddeb8d6\allowbreak{}6a7cdf04\allowbreak{}c0fe7117\allowbreak{}fb4a5c39\allowbreak{}3dad9b65\allowbreak{}0a112925\allowbreak{}31863750\allowbreak{}3324e2e0


\entrylabel{acquired at}

2026-09-07T00:\allowbreak{}14:\allowbreak{}14Z


\subsection*{AIM stochastic methods for non-equilibrium dynamical systems,\allowbreak{} Ben Webb notes}

aim\_\allowbreak{}nonequilibrium


\url{https://aimath.org/pastworkshops/noneqdynsysproblems.pdf}


\entrylabel{access evidence}

Public primary author or workshop PDF;\allowbreak{} individual open-question statements and later status require separate validation


\entrylabel{reading status}

unread


\entrylabel{statement validation}

not started


\entrylabel{sha256}

276ff5e0\allowbreak{}93da6f81\allowbreak{}339a9749\allowbreak{}8b5bfabe\allowbreak{}a89acf83\allowbreak{}9ed754e5\allowbreak{}933bc5b0\allowbreak{}22de8ce4


\entrylabel{acquired at}

2026-09-06T19:\allowbreak{}32:\allowbreak{}44Z


\subsection*{AIM phase transitions workshop problems,\allowbreak{} compiled by Charles Radin}

aim\_\allowbreak{}phase\_\allowbreak{}transitions


\url{https://aimath.org/WWN/phasetransition/phasetransition.pdf}


\entrylabel{access evidence}

Public primary author or workshop PDF;\allowbreak{} individual open-question statements and later status require separate validation


\entrylabel{reading status}

unread


\entrylabel{statement validation}

not started


\entrylabel{sha256}

5a2f7236\allowbreak{}eb67c095\allowbreak{}fbc95c88\allowbreak{}213c28aa\allowbreak{}cf07d6fe\allowbreak{}434ba026\allowbreak{}abd9597c\allowbreak{}0669cea4


\entrylabel{acquired at}

2026-09-06T19:\allowbreak{}33:\allowbreak{}28Z


\subsection*{AIM random analytic functions workshop open problems,\allowbreak{} April2006}

aim\_\allowbreak{}random\_\allowbreak{}analytic\_\allowbreak{}2006


\url{https://aimath.org/WWN/randomzeros/arcc1.pdf}


\entrylabel{access evidence}

Public primary author or workshop PDF;\allowbreak{} individual open-question statements and later status require separate validation


\entrylabel{reading status}

unread


\entrylabel{statement validation}

not started


\entrylabel{sha256}

f415ecad\allowbreak{}26caf095\allowbreak{}b963ba9b\allowbreak{}7da03c2f\allowbreak{}1a5f85da\allowbreak{}aed16838\allowbreak{}bc388cc6\allowbreak{}7b026b4d


\entrylabel{acquired at}

2026-09-06T19:\allowbreak{}31:\allowbreak{}50Z


\subsection*{AIM random matrices workshop problem list}

aim\_\allowbreak{}random\_\allowbreak{}matrices


\url{https://aimath.org/WWN/randommatrices/randommatrices.pdf}


\entrylabel{access evidence}

Public primary author or workshop PDF;\allowbreak{} exact statements and later status separately checked


\entrylabel{reading status}

unread


\entrylabel{statement validation}

not started


\entrylabel{sha256}

5f3f0e00\allowbreak{}106140d5\allowbreak{}401fe671\allowbreak{}da093bbc\allowbreak{}2e23ab4f\allowbreak{}8faa4e2d\allowbreak{}b034ebc5\allowbreak{}b50e080d


\entrylabel{acquired at}

2026-09-06T19:\allowbreak{}50:\allowbreak{}04Z


\subsection*{Aldous open problems:\allowbreak{} Bacon and eggs batch scheduling}

aldous\_\allowbreak{}bacon\_\allowbreak{}20260906


\url{https://www.stat.berkeley.edu/~aldous/Research/OP/bacon.html}


\entrylabel{access evidence}

Public primary author source,\allowbreak{} accessed2026-09-06;\allowbreak{} mathematical content and current status checked separately


\entrylabel{reading status}

unread


\entrylabel{statement validation}

not started


\entrylabel{sha256}

ac9b73e5\allowbreak{}6230e8f4\allowbreak{}150301f8\allowbreak{}6bd9f583\allowbreak{}301476b4\allowbreak{}35070f56\allowbreak{}0262538b\allowbreak{}e3a22d67


\entrylabel{acquired at}

2026-09-06T22:\allowbreak{}04:\allowbreak{}58Z


\subsection*{A survey of max-type recursive distributional equations}

aldous\_\allowbreak{}bandyopadhyay\_\allowbreak{}2005


\url{https://arxiv.org/pdf/math/0401388}


\entrylabel{access evidence}

Public primary author website or arXiv manuscript


\entrylabel{reading status}

unread


\entrylabel{statement validation}

not started


\entrylabel{sha256}

a54121c0\allowbreak{}a0fe0663\allowbreak{}92edf3db\allowbreak{}a6652ef2\allowbreak{}fd22210b\allowbreak{}610c65aa\allowbreak{}9416aa71\allowbreak{}0f09631a


\entrylabel{acquired at}

2026-09-06T16:\allowbreak{}43:\allowbreak{}17Z


\subsection*{Aldous open problems:\allowbreak{} Metropolis on Cayley graphs}

aldous\_\allowbreak{}cayley\_\allowbreak{}20260906


\url{https://www.stat.berkeley.edu/~aldous/Research/OP/cayley.html}


\entrylabel{access evidence}

Public primary author source,\allowbreak{} accessed2026-09-06;\allowbreak{} mathematical content and current status checked separately


\entrylabel{reading status}

unread


\entrylabel{statement validation}

not started


\entrylabel{sha256}

9989446b\allowbreak{}f73a5d08\allowbreak{}579242e7\allowbreak{}0a89cb6d\allowbreak{}09cdc4b6\allowbreak{}0032ff9c\allowbreak{}bf29f090\allowbreak{}64eb8f31


\entrylabel{acquired at}

2026-09-06T22:\allowbreak{}05:\allowbreak{}06Z


\subsection*{Aldous open problems:\allowbreak{} Low-density coalescing branching random walk}

aldous\_\allowbreak{}cbrw\_\allowbreak{}20260906


\url{https://www.stat.berkeley.edu/~aldous/Research/OP/cbrw.html}


\entrylabel{access evidence}

Public primary author source,\allowbreak{} accessed2026-09-06;\allowbreak{} mathematical content and current status checked separately


\entrylabel{reading status}

unread


\entrylabel{statement validation}

not started


\entrylabel{sha256}

69f398a7\allowbreak{}ef7fc65e\allowbreak{}190b7382\allowbreak{}479bb108\allowbreak{}4f948067\allowbreak{}f93a3c25\allowbreak{}e734085b\allowbreak{}af817600


\entrylabel{acquired at}

2026-09-06T22:\allowbreak{}04:\allowbreak{}49Z


\subsection*{Aldous:\allowbreak{} A spatial model of city growth and formation}

aldous\_\allowbreak{}cities\_\allowbreak{}2012


\url{https://arxiv.org/pdf/1209.5120}


\entrylabel{access evidence}

Public primary source;\allowbreak{} precise open questions and later resolutions checked separately


\entrylabel{reading status}

unread


\entrylabel{statement validation}

not started


\entrylabel{sha256}

eeecb371\allowbreak{}b6e6c286\allowbreak{}0c01e1b7\allowbreak{}3d812468\allowbreak{}b619ddb7\allowbreak{}e3c4cd1d\allowbreak{}0e3affc0\allowbreak{}56e8a077


\entrylabel{acquired at}

2026-09-06T22:\allowbreak{}16:\allowbreak{}15Z


\subsection*{Aldous open problems:\allowbreak{} Spatial city growth and formation}

aldous\_\allowbreak{}cities\_\allowbreak{}20260906


\url{https://www.stat.berkeley.edu/~aldous/Research/OP/cities.html}


\entrylabel{access evidence}

Public primary author source,\allowbreak{} accessed2026-09-06;\allowbreak{} mathematical content and current status checked separately


\entrylabel{reading status}

unread


\entrylabel{statement validation}

not started


\entrylabel{sha256}

53f511b6\allowbreak{}a3424d07\allowbreak{}604eb19b\allowbreak{}cdc5814f\allowbreak{}ce6f710e\allowbreak{}b6c440a5\allowbreak{}09e49fcc\allowbreak{}d9a68999


\entrylabel{acquired at}

2026-09-06T22:\allowbreak{}05:\allowbreak{}59Z


\subsection*{Aldous open problems:\allowbreak{} Branch-rotation chain mixing}

aldous\_\allowbreak{}clad\_\allowbreak{}mix\_\allowbreak{}20260906


\url{https://www.stat.berkeley.edu/~aldous/Research/OP/clad-mix.html}


\entrylabel{access evidence}

Public primary author source,\allowbreak{} accessed2026-09-06;\allowbreak{} mathematical content and current status checked separately


\entrylabel{reading status}

unread


\entrylabel{statement validation}

not started


\entrylabel{sha256}

eb77b2a7\allowbreak{}bde90c9a\allowbreak{}fda28c7b\allowbreak{}8ed738a2\allowbreak{}f6db0c5d\allowbreak{}2c52813d\allowbreak{}586d6d75\allowbreak{}c2b21208


\entrylabel{acquired at}

2026-09-06T22:\allowbreak{}05:\allowbreak{}22Z


\subsection*{Aldous:\allowbreak{} Mixing times for branch rotation and triangulation walk}

aldous\_\allowbreak{}clad\_\allowbreak{}rotation\_\allowbreak{}source


\url{https://www.stat.berkeley.edu/~aldous/206-RWG/RWGpapers/clad-mix.pdf}


\entrylabel{access evidence}

Primary source statement and later-literature scope verification2026-09-06


\entrylabel{reading status}

unread


\entrylabel{statement validation}

not started


\entrylabel{sha256}

005420ad\allowbreak{}d4edf92c\allowbreak{}d7ad46df\allowbreak{}7799ef08\allowbreak{}2db1414a\allowbreak{}07f46d80\allowbreak{}818df109\allowbreak{}39316654


\entrylabel{acquired at}

2026-09-06T22:\allowbreak{}48:\allowbreak{}36Z


\subsection*{Aldous open problems:\allowbreak{} Coagulation-fragmentation mixing}

aldous\_\allowbreak{}coag\_\allowbreak{}frag\_\allowbreak{}20260906


\url{https://www.stat.berkeley.edu/~aldous/Research/OP/coag_frag.html}


\entrylabel{access evidence}

Public primary author source,\allowbreak{} accessed2026-09-06;\allowbreak{} mathematical content and current status checked separately


\entrylabel{reading status}

unread


\entrylabel{statement validation}

not started


\entrylabel{sha256}

4205247b\allowbreak{}098318c0\allowbreak{}b262384d\allowbreak{}f06c2efc\allowbreak{}93ad8b67\allowbreak{}3efaf821\allowbreak{}6e4f54c2\allowbreak{}4b7c9d47


\entrylabel{acquired at}

2026-09-06T22:\allowbreak{}04:\allowbreak{}42Z


\subsection*{Aldous-Fill unposted notes:\allowbreak{} coagulation fragmentation mixing}

aldous\_\allowbreak{}coag\_\allowbreak{}frag\_\allowbreak{}notes


\url{https://www.stat.berkeley.edu/~aldous/Research/OP/coag_frag_mix.pdf}


\entrylabel{access evidence}

Public primary source;\allowbreak{} precise open questions and later resolutions checked separately


\entrylabel{reading status}

unread


\entrylabel{statement validation}

not started


\entrylabel{sha256}

1ff70d6c\allowbreak{}2dba03cc\allowbreak{}80738b72\allowbreak{}3f22e476\allowbreak{}509f3255\allowbreak{}131343b6\allowbreak{}55726d0a\allowbreak{}6e86b7d7


\entrylabel{acquired at}

2026-09-06T22:\allowbreak{}16:\allowbreak{}12Z


\subsection*{Aldous open problems:\allowbreak{} Compactification of finite reversible chains}

aldous\_\allowbreak{}compact\_\allowbreak{}20260906


\url{https://www.stat.berkeley.edu/~aldous/Research/OP/compact.html}


\entrylabel{access evidence}

Public primary author source,\allowbreak{} accessed2026-09-06;\allowbreak{} mathematical content and current status checked separately


\entrylabel{reading status}

unread


\entrylabel{statement validation}

not started


\entrylabel{sha256}

96289780\allowbreak{}c33904b0\allowbreak{}f61c155e\allowbreak{}378d5d25\allowbreak{}08684845\allowbreak{}a355d77d\allowbreak{}b2494881\allowbreak{}719086e7


\entrylabel{acquired at}

2026-09-06T22:\allowbreak{}10:\allowbreak{}45Z


\subsection*{Aldous:\allowbreak{} Compactification notes for reversible Markov chains}

aldous\_\allowbreak{}compact\_\allowbreak{}notes


\url{https://www.stat.berkeley.edu/~aldous/Talks/MCcompact.pdf}


\entrylabel{access evidence}

Public primary source;\allowbreak{} precise open questions and later resolutions checked separately


\entrylabel{reading status}

unread


\entrylabel{statement validation}

not started


\entrylabel{sha256}

d556a455\allowbreak{}3f12d4a5\allowbreak{}c83fbe85\allowbreak{}d52a04bb\allowbreak{}833e58d8\allowbreak{}585c98ec\allowbreak{}7148146c\allowbreak{}47673ee0


\entrylabel{acquired at}

2026-09-06T22:\allowbreak{}16:\allowbreak{}10Z


\subsection*{Aldous open problems:\allowbreak{} Congestion and percolation-fragmentation}

aldous\_\allowbreak{}congestion\_\allowbreak{}20260906


\url{https://www.stat.berkeley.edu/~aldous/Research/OP/congestion.html}


\entrylabel{access evidence}

Public primary author source,\allowbreak{} accessed2026-09-06;\allowbreak{} mathematical content and current status checked separately


\entrylabel{reading status}

unread


\entrylabel{statement validation}

not started


\entrylabel{sha256}

341bf2e5\allowbreak{}f1e024c9\allowbreak{}e654a450\allowbreak{}cc084951\allowbreak{}8d1c3ad5\allowbreak{}2fcfd086\allowbreak{}702a7e4f\allowbreak{}d28baed8


\entrylabel{acquired at}

2026-09-06T22:\allowbreak{}04:\allowbreak{}19Z


\subsection*{Aldous open problems:\allowbreak{} Percolation and empires}

aldous\_\allowbreak{}empires\_\allowbreak{}20260906


\url{https://www.stat.berkeley.edu/~aldous/Research/OP/empires.html}


\entrylabel{access evidence}

Public primary author source,\allowbreak{} accessed2026-09-06;\allowbreak{} mathematical content and current status checked separately


\entrylabel{reading status}

unread


\entrylabel{statement validation}

not started


\entrylabel{sha256}

efb17f45\allowbreak{}74fd713e\allowbreak{}effdab4a\allowbreak{}fe7792e1\allowbreak{}21f8056d\allowbreak{}57b8f634\allowbreak{}f4855337\allowbreak{}9eda7919


\entrylabel{acquired at}

2026-09-06T22:\allowbreak{}05:\allowbreak{}50Z


\subsection*{Aldous open problems:\allowbreak{} Martingale for practical purposes}

aldous\_\allowbreak{}fields\_\allowbreak{}20260906


\url{https://www.stat.berkeley.edu/~aldous/Research/OP/fields.html}


\entrylabel{access evidence}

Public primary author source,\allowbreak{} accessed2026-09-06;\allowbreak{} mathematical content and current status checked separately


\entrylabel{reading status}

unread


\entrylabel{statement validation}

not started


\entrylabel{sha256}

a6e69e62\allowbreak{}ec5ca5c7\allowbreak{}7750245b\allowbreak{}07d1182d\allowbreak{}d44488ea\allowbreak{}b863cdc8\allowbreak{}3d8ca1a6\allowbreak{}122694e5


\entrylabel{acquired at}

2026-09-06T22:\allowbreak{}04:\allowbreak{}12Z


\subsection*{Aldous-Fill:\allowbreak{} Reversible Markov Chains and Random Walks on Graphs,\allowbreak{} 2002 recompiled2014}

aldous\_\allowbreak{}fill\_\allowbreak{}2014


\url{https://www.stat.berkeley.edu/~aldous/RWG/book.pdf}


\entrylabel{access evidence}

Public primary author source,\allowbreak{} accessed2026-09-06;\allowbreak{} mathematical content and current status checked separately


\entrylabel{reading status}

unread


\entrylabel{statement validation}

not started


\entrylabel{sha256}

8f9b9341\allowbreak{}11c30b3f\allowbreak{}8fd29d94\allowbreak{}a18794d4\allowbreak{}b86f824b\allowbreak{}19f7a3c2\allowbreak{}7b1153cd\allowbreak{}fac08e79


\entrylabel{acquired at}

2026-09-06T22:\allowbreak{}03:\allowbreak{}09Z


\subsection*{Aldous open problems:\allowbreak{} Bayesian graph Laplacian spectral gap}

aldous\_\allowbreak{}gap\_\allowbreak{}ion\_\allowbreak{}20260906


\url{https://www.stat.berkeley.edu/~aldous/Research/OP/gap_ion.html}


\entrylabel{access evidence}

Public primary author source,\allowbreak{} accessed2026-09-06;\allowbreak{} mathematical content and current status checked separately


\entrylabel{reading status}

unread


\entrylabel{statement validation}

not started


\entrylabel{sha256}

adc6b7cb\allowbreak{}91f88216\allowbreak{}642fbce3\allowbreak{}b0d9d56c\allowbreak{}f28be389\allowbreak{}728ec553\allowbreak{}31acf797\allowbreak{}daa1e870


\entrylabel{acquired at}

2026-09-06T22:\allowbreak{}05:\allowbreak{}10Z


\subsection*{Aldous open problems:\allowbreak{} Compression of sparse labelled graphs}

aldous\_\allowbreak{}graph\_\allowbreak{}compression\_\allowbreak{}20260906


\url{https://www.stat.berkeley.edu/~aldous/Research/OP/graph_compression.html}


\entrylabel{access evidence}

Public primary author source,\allowbreak{} accessed2026-09-06;\allowbreak{} mathematical content and current status checked separately


\entrylabel{reading status}

unread


\entrylabel{statement validation}

not started


\entrylabel{sha256}

d6aa44f2\allowbreak{}3bd408ec\allowbreak{}6a45d20e\allowbreak{}f59d0ff2\allowbreak{}5fa442f6\allowbreak{}3303ac3f\allowbreak{}164fd5c6\allowbreak{}12d3ff02


\entrylabel{acquired at}

2026-09-06T22:\allowbreak{}04:\allowbreak{}34Z


\subsection*{Aldous open problems:\allowbreak{} Discrete Hammersley process and oriented percolation flows}

aldous\_\allowbreak{}hammersley\_\allowbreak{}flow\_\allowbreak{}20260906


\url{https://www.stat.berkeley.edu/~aldous/Research/OP/hammersley_flow.html}


\entrylabel{access evidence}

Public primary author source,\allowbreak{} accessed2026-09-06;\allowbreak{} mathematical content and current status checked separately


\entrylabel{reading status}

unread


\entrylabel{statement validation}

not started


\entrylabel{sha256}

1c3078bb\allowbreak{}05b25554\allowbreak{}b7ab7bc1\allowbreak{}2e6d8567\allowbreak{}b777407b\allowbreak{}6c6ad24f\allowbreak{}ecfc8971\allowbreak{}5088a301


\entrylabel{acquired at}

2026-09-06T22:\allowbreak{}10:\allowbreak{}49Z


\subsection*{Aldous open problems:\allowbreak{} One-dimensional drift-jump particle process}

aldous\_\allowbreak{}left\_\allowbreak{}hamm\_\allowbreak{}20260906


\url{https://www.stat.berkeley.edu/~aldous/Research/OP/left_Hamm.html}


\entrylabel{access evidence}

Public primary author source,\allowbreak{} accessed2026-09-06;\allowbreak{} mathematical content and current status checked separately


\entrylabel{reading status}

unread


\entrylabel{statement validation}

not started


\entrylabel{sha256}

1a208ce3\allowbreak{}851d6502\allowbreak{}2ecc998b\allowbreak{}edcfe97f\allowbreak{}d8afe524\allowbreak{}0e52cec3\allowbreak{}172aef0d\allowbreak{}9a89b2b6


\entrylabel{acquired at}

2026-09-06T22:\allowbreak{}10:\allowbreak{}42Z


\subsection*{Bowen-Chapman-Lubotzky-Vidick:\allowbreak{} The Aldous-Lyons conjecture I,\allowbreak{} subgroup tests}

aldous\_\allowbreak{}lyons\_\allowbreak{}negative\_\allowbreak{}2024


\url{https://arxiv.org/pdf/2408.00110}


\entrylabel{access evidence}

Public primary source;\allowbreak{} precise open questions and later resolutions checked separately


\entrylabel{reading status}

unread


\entrylabel{statement validation}

not started


\entrylabel{sha256}

754f587a\allowbreak{}a5592c14\allowbreak{}7ded5964\allowbreak{}22b5c935\allowbreak{}0a1fc876\allowbreak{}7bc1c051\allowbreak{}7f054f2f\allowbreak{}19ccd2f5


\entrylabel{acquired at}

2026-09-06T22:\allowbreak{}16:\allowbreak{}18Z


\subsection*{Aldous:\allowbreak{} Limits for processes over general networks,\allowbreak{}2018}

aldous\_\allowbreak{}markov\_\allowbreak{}topology\_\allowbreak{}2018


\url{https://www.stat.berkeley.edu/~aldous/Research/OP/talk_towsner_2018.pdf}


\entrylabel{access evidence}

Primary URL discovered through focused literature search2026-09-06;\allowbreak{} substantive title and theorem verification follows acquisition.


\entrylabel{reading status}

unread


\entrylabel{statement validation}

not started


\entrylabel{sha256}

fdf0d2be\allowbreak{}a05dc981\allowbreak{}662d3362\allowbreak{}2fb6a104\allowbreak{}650efc92\allowbreak{}61d3d5bd\allowbreak{}cd118ad3\allowbreak{}694d6b8c


\entrylabel{acquired at}

2026-09-06T22:\allowbreak{}39:\allowbreak{}17Z


\subsection*{Aldous:\allowbreak{} The Nearest Unvisited Vertex Walk on Random Graphs,\allowbreak{}internal date26August2021}

aldous\_\allowbreak{}nuv\_\allowbreak{}arxiv


\url{https://arxiv.org/pdf/1912.13175}


\entrylabel{access evidence}

Primary author manuscript;\allowbreak{} title and question locations to be read


\entrylabel{reading status}

unread


\entrylabel{statement validation}

not started


\entrylabel{sha256}

df77b168\allowbreak{}14644797\allowbreak{}853f8748\allowbreak{}508f0f6b\allowbreak{}fe40506c\allowbreak{}9de17822\allowbreak{}7a80cf54\allowbreak{}d895e1e5


\entrylabel{acquired at}

2026-09-06T23:\allowbreak{}13:\allowbreak{}46Z


\subsection*{Aldous-Angel-Berestycki:\allowbreak{} Online random weight minimal spanning trees,\allowbreak{} draft2008}

aldous\_\allowbreak{}online\_\allowbreak{}mst\_\allowbreak{}2008


\url{https://www.stat.berkeley.edu/~aldous/Research/OP/Beres2008.pdf}


\entrylabel{access evidence}

Primary URL discovered through focused literature search2026-09-06;\allowbreak{} substantive title and theorem verification follows acquisition.


\entrylabel{reading status}

unread


\entrylabel{statement validation}

not started


\entrylabel{sha256}

250c6526\allowbreak{}9b510d59\allowbreak{}7de38a1b\allowbreak{}d2c8ce9b\allowbreak{}330b6414\allowbreak{}d5ab3f2d\allowbreak{}8a3cccc3\allowbreak{}fc2d0358


\entrylabel{acquired at}

2026-09-06T22:\allowbreak{}39:\allowbreak{}04Z


\subsection*{Aldous open problems:\allowbreak{} Online minimum spanning tree}

aldous\_\allowbreak{}online\_\allowbreak{}zeta3\_\allowbreak{}20260906


\url{https://www.stat.berkeley.edu/~aldous/Research/OP/online_zeta3.html}


\entrylabel{access evidence}

Public primary author source,\allowbreak{} accessed2026-09-06;\allowbreak{} mathematical content and current status checked separately


\entrylabel{reading status}

unread


\entrylabel{statement validation}

not started


\entrylabel{sha256}

8a0c5b22\allowbreak{}361b3bc8\allowbreak{}58867ab5\allowbreak{}e68dfe6c\allowbreak{}773e2f4f\allowbreak{}3816c1c5\allowbreak{}02a0d969\allowbreak{}28c7c4ee


\entrylabel{acquired at}

2026-09-06T22:\allowbreak{}06:\allowbreak{}20Z


\subsection*{David Aldous:\allowbreak{} Some Open Problems,\allowbreak{} May2003}

aldous\_\allowbreak{}open\_\allowbreak{}2003


\url{https://www.stat.berkeley.edu/~aldous/Research/OP/problems_old.pdf}


\entrylabel{access evidence}

Primary author publicly accessible source;\allowbreak{} numbered open questions require later-status verification


\entrylabel{reading status}

unread


\entrylabel{statement validation}

not started


\entrylabel{sha256}

3ddbd2d6\allowbreak{}b95c7bcf\allowbreak{}ccfd9767\allowbreak{}f3e9d4ec\allowbreak{}f221ddd6\allowbreak{}fb036d4e\allowbreak{}36e9700e\allowbreak{}c1e53418


\entrylabel{acquired at}

2026-09-06T21:\allowbreak{}44:\allowbreak{}57Z


\subsection*{David Aldous:\allowbreak{} current open-problem index and resolved-question registry}

aldous\_\allowbreak{}open\_\allowbreak{}index\_\allowbreak{}20260906


\url{https://www.stat.berkeley.edu/~aldous/Research/OP/index.html}


\entrylabel{access evidence}

Primary author publicly accessible source;\allowbreak{} numbered open questions require later-status verification


\entrylabel{reading status}

unread


\entrylabel{statement validation}

not started


\entrylabel{sha256}

69de0a61\allowbreak{}207bac04\allowbreak{}b9452263\allowbreak{}afbc7532\allowbreak{}1409adb7\allowbreak{}5e673d0f\allowbreak{}6dfb29a2\allowbreak{}441adb40


\entrylabel{acquired at}

2026-09-06T21:\allowbreak{}44:\allowbreak{}46Z


\subsection*{Aldous:\allowbreak{} Random partitions of the plane via Poissonian coloring and coalescing planar partitions}

aldous\_\allowbreak{}partition\_\allowbreak{}2017


\url{https://arxiv.org/pdf/1701.00131}


\entrylabel{access evidence}

Public primary source;\allowbreak{} precise open questions and later resolutions checked separately


\entrylabel{reading status}

unread


\entrylabel{statement validation}

not started


\entrylabel{sha256}

c8b365c6\allowbreak{}39a8611d\allowbreak{}b5c5941b\allowbreak{}d92a458a\allowbreak{}17baa4e9\allowbreak{}74699157\allowbreak{}e112fbdf\allowbreak{}73a2433d


\entrylabel{acquired at}

2026-09-06T22:\allowbreak{}16:\allowbreak{}08Z


\subsection*{Aldous open problems:\allowbreak{} Random plane partition topology}

aldous\_\allowbreak{}partition\_\allowbreak{}20260906


\url{https://www.stat.berkeley.edu/~aldous/Research/OP/partition.html}


\entrylabel{access evidence}

Public primary author source,\allowbreak{} accessed2026-09-06;\allowbreak{} mathematical content and current status checked separately


\entrylabel{reading status}

unread


\entrylabel{statement validation}

not started


\entrylabel{sha256}

8cce7766\allowbreak{}af016f82\allowbreak{}b00a40fe\allowbreak{}03ccfb20\allowbreak{}da692d12\allowbreak{}4fd92eeb\allowbreak{}ee48815d\allowbreak{}8b6928ff


\entrylabel{acquired at}

2026-09-06T22:\allowbreak{}06:\allowbreak{}12Z


\subsection*{Aldous-Diaconis:\allowbreak{} unpublished1993 patience-sorting notes}

aldous\_\allowbreak{}patience\_\allowbreak{}1993


\url{https://www.stat.berkeley.edu/~aldous/Research/OP/patience.pdf}


\entrylabel{access evidence}

Public primary source;\allowbreak{} precise open questions and later resolutions checked separately


\entrylabel{reading status}

unread


\entrylabel{statement validation}

not started


\entrylabel{sha256}

7d745742\allowbreak{}4b2f1422\allowbreak{}e3090069\allowbreak{}2e1626b9\allowbreak{}fe07d172\allowbreak{}e8b1cfe1\allowbreak{}2c3470f4\allowbreak{}18ac717c


\entrylabel{acquired at}

2026-09-06T22:\allowbreak{}16:\allowbreak{}01Z


\subsection*{Aldous-Diaconis:\allowbreak{} Longest increasing subsequences,\allowbreak{} from patience sorting to the Baik-Deift-Johansson theorem,\allowbreak{}1999}

aldous\_\allowbreak{}patience\_\allowbreak{}1999


\url{https://www.stat.berkeley.edu/~aldous/Papers/me86.pdf}


\entrylabel{access evidence}

Public primary source;\allowbreak{} precise open questions and later resolutions checked separately


\entrylabel{reading status}

unread


\entrylabel{statement validation}

not started


\entrylabel{sha256}

19350af5\allowbreak{}be8b46ac\allowbreak{}a281c32a\allowbreak{}9dc063e4\allowbreak{}6bf31c6f\allowbreak{}8e97ae91\allowbreak{}9ff9e8eb\allowbreak{}547e7317


\entrylabel{acquired at}

2026-09-06T22:\allowbreak{}16:\allowbreak{}05Z


\subsection*{Aldous open problems:\allowbreak{} Shortest routes in random proximity networks}

aldous\_\allowbreak{}random\_\allowbreak{}proximity\_\allowbreak{}20260906


\url{https://www.stat.berkeley.edu/~aldous/Research/OP/random_proximity.html}


\entrylabel{access evidence}

Public primary author source,\allowbreak{} accessed2026-09-06;\allowbreak{} mathematical content and current status checked separately


\entrylabel{reading status}

unread


\entrylabel{statement validation}

not started


\entrylabel{sha256}

dce18122\allowbreak{}09caaf03\allowbreak{}0c123fb8\allowbreak{}29130071\allowbreak{}c73aaea8\allowbreak{}b5d25856\allowbreak{}cc5f8625\allowbreak{}6a6a6073


\entrylabel{acquired at}

2026-09-06T22:\allowbreak{}05:\allowbreak{}16Z


\subsection*{Aldous and Bandyopadhyay:\allowbreak{} A survey of max-type recursive distributional equations}

aldous\_\allowbreak{}rde\_\allowbreak{}2005


\url{https://www.stat.berkeley.edu/~aldous/Papers/me107.pdf}


\entrylabel{access evidence}

Primary author-hosted collection or maintained status page


\entrylabel{reading status}

PDF25through30 full read other sections partially read


\entrylabel{statement validation}

Problem31 and Conjecture34 hypotheses checked literal degeneracy retained


\entrylabel{sha256}

63f611db\allowbreak{}22543fb7\allowbreak{}5b800b43\allowbreak{}b346a3af\allowbreak{}a5653441\allowbreak{}9c7eb474\allowbreak{}2cacde12\allowbreak{}05e778cf


\entrylabel{acquired at}

2026-09-07T01:\allowbreak{}59:\allowbreak{}45Z


\subsection*{Aldous author-maintained RDE survey updates}

aldous\_\allowbreak{}rde\_\allowbreak{}status\_\allowbreak{}page


\url{https://www.stat.berkeley.edu/~aldous/Research/RDE/index.html}


\entrylabel{access evidence}

Primary author-hosted collection or maintained status page


\entrylabel{reading status}

unread


\entrylabel{statement validation}

not started


\entrylabel{sha256}

6f8697c1\allowbreak{}e8bae36c\allowbreak{}a94a7b0f\allowbreak{}7039d9e0\allowbreak{}5a4fdc44\allowbreak{}d2fc10c4\allowbreak{}c4158794\allowbreak{}ab2f2044


\entrylabel{acquired at}

2026-09-07T01:\allowbreak{}59:\allowbreak{}49Z


\subsection*{Aldous open problems:\allowbreak{} Scale invariant random spatial networks}

aldous\_\allowbreak{}sirsn\_\allowbreak{}20260906


\url{https://www.stat.berkeley.edu/~aldous/Research/OP/sirsn.html}


\entrylabel{access evidence}

Public primary author source,\allowbreak{} accessed2026-09-06;\allowbreak{} mathematical content and current status checked separately


\entrylabel{reading status}

unread


\entrylabel{statement validation}

not started


\entrylabel{sha256}

0f2d77bc\allowbreak{}8c702dc5\allowbreak{}0c02f41f\allowbreak{}bc2cf1be\allowbreak{}1898d86e\allowbreak{}d82f174b\allowbreak{}e22f8701\allowbreak{}535ec748


\entrylabel{acquired at}

2026-09-06T22:\allowbreak{}06:\allowbreak{}08Z


\subsection*{Aldous:\allowbreak{} Open problems in spatial networks,\allowbreak{} Paris2022}

aldous\_\allowbreak{}spatial\_\allowbreak{}talk\_\allowbreak{}2022


\url{https://www.stat.berkeley.edu/~aldous/Talks/Paris_2022_talk.pdf}


\entrylabel{access evidence}

Primary URL discovered through focused literature search2026-09-06;\allowbreak{} substantive title and theorem verification follows acquisition.


\entrylabel{reading status}

unread


\entrylabel{statement validation}

not started


\entrylabel{sha256}

891dc17a\allowbreak{}a4273ba8\allowbreak{}7b645876\allowbreak{}2c94603d\allowbreak{}761309fd\allowbreak{}4d831197\allowbreak{}a074c408\allowbreak{}7ef09ac4


\entrylabel{acquired at}

2026-09-06T22:\allowbreak{}39:\allowbreak{}52Z


\subsection*{Aldous-Lando:\allowbreak{} The stretch-length tradeoff in geometric networks}

aldous\_\allowbreak{}stretch\_\allowbreak{}2015


\url{https://arxiv.org/pdf/1404.2653}


\entrylabel{access evidence}

Primary source statement and later-literature scope verification2026-09-06


\entrylabel{reading status}

unread


\entrylabel{statement validation}

not started


\entrylabel{sha256}

92ae1ef8\allowbreak{}a49209a9\allowbreak{}37bf12b5\allowbreak{}e7dcbc8c\allowbreak{}709545b2\allowbreak{}230490c2\allowbreak{}638e8719\allowbreak{}f5b11738


\entrylabel{acquired at}

2026-09-06T22:\allowbreak{}48:\allowbreak{}41Z


\subsection*{Aldous open problems:\allowbreak{} Spatial network stretch-length tradeoff}

aldous\_\allowbreak{}stretch\_\allowbreak{}20260906


\url{https://www.stat.berkeley.edu/~aldous/Research/OP/stretch.html}


\entrylabel{access evidence}

Public primary author source,\allowbreak{} accessed2026-09-06;\allowbreak{} mathematical content and current status checked separately


\entrylabel{reading status}

unread


\entrylabel{statement validation}

not started


\entrylabel{sha256}

f3f925f1\allowbreak{}12e37c8a\allowbreak{}27063d35\allowbreak{}c232135b\allowbreak{}e5d42943\allowbreak{}dae61658\allowbreak{}2e755fa8\allowbreak{}7065ff11


\entrylabel{acquired at}

2026-09-06T22:\allowbreak{}05:\allowbreak{}33Z


\subsection*{Aldous open problems:\allowbreak{} Largest common random substructures}

aldous\_\allowbreak{}substructures\_\allowbreak{}20260906


\url{https://www.stat.berkeley.edu/~aldous/Research/OP/substructures.html}


\entrylabel{access evidence}

Public primary author source,\allowbreak{} accessed2026-09-06;\allowbreak{} mathematical content and current status checked separately


\entrylabel{reading status}

unread


\entrylabel{statement validation}

not started


\entrylabel{sha256}

26952c61\allowbreak{}4e973684\allowbreak{}3f48855e\allowbreak{}b5e9fe79\allowbreak{}fcf627c5\allowbreak{}d54b2af3\allowbreak{}ed213614\allowbreak{}a3c93b5e


\entrylabel{acquired at}

2026-09-06T22:\allowbreak{}05:\allowbreak{}42Z


\subsection*{Aldous:\allowbreak{} Route Lengths in Invariant Spatial Tree Networks}

aldous\_\allowbreak{}tree\_\allowbreak{}lengths\_\allowbreak{}2021


\url{https://arxiv.org/pdf/2103.00669}


\entrylabel{access evidence}

Primary URL discovered through focused literature search2026-09-06;\allowbreak{} substantive title and theorem verification follows acquisition.


\entrylabel{reading status}

unread


\entrylabel{statement validation}

not started


\entrylabel{sha256}

942fb212\allowbreak{}91dba57e\allowbreak{}7214dcea\allowbreak{}915e6338\allowbreak{}b7907669\allowbreak{}d9d2f872\allowbreak{}940bfdec\allowbreak{}5b397bd7


\entrylabel{acquired at}

2026-09-06T22:\allowbreak{}38:\allowbreak{}07Z


\subsection*{Aldous open problems:\allowbreak{} Local weak limits and unimodularity}

aldous\_\allowbreak{}unimodular\_\allowbreak{}20260906


\url{https://www.stat.berkeley.edu/~aldous/Research/OP/unimodular.html}


\entrylabel{access evidence}

Public primary author source,\allowbreak{} accessed2026-09-06;\allowbreak{} mathematical content and current status checked separately


\entrylabel{reading status}

unread


\entrylabel{statement validation}

not started


\entrylabel{sha256}

9f7b9a4b\allowbreak{}2d4d6b4d\allowbreak{}364da371\allowbreak{}d1f741e8\allowbreak{}efc982c1\allowbreak{}964a6929\allowbreak{}7e873a60\allowbreak{}86902b1f


\entrylabel{acquired at}

2026-09-06T22:\allowbreak{}06:\allowbreak{}16Z


\subsection*{Hoffman Johnson Junge Meisel:\allowbreak{} Cutoff for activated random walk}

arw\_\allowbreak{}cutoff\_\allowbreak{}2025


\url{https://arxiv.org/pdf/2501.17938}


\entrylabel{access evidence}

Primary resolution of the one-dimensional cutoff conjecture


\entrylabel{reading status}

PDF1-2,\allowbreak{} Theorem1.1


\entrylabel{statement validation}

Total-variation cutoff on intervals credited


\entrylabel{sha256}

38ac26aa\allowbreak{}b9a57a5f\allowbreak{}26b5e4af\allowbreak{}9c79ea8b\allowbreak{}ea4355e9\allowbreak{}5e60c87e\allowbreak{}e7a723be\allowbreak{}f44f3057


\entrylabel{acquired at}

2026-09-07T05:\allowbreak{}11:\allowbreak{}32Z


\subsection*{Hoffman Johnson Junge:\allowbreak{} The density conjecture for activated random walk}

arw\_\allowbreak{}density\_\allowbreak{}2024


\url{https://arxiv.org/pdf/2406.01731}


\entrylabel{access evidence}

Later primary result identified in survey status check


\entrylabel{reading status}

PDF1-5 and Corollary8.9 statement;\allowbreak{} partial proof overview


\entrylabel{statement validation}

One-dimensional density and ball results credited;\allowbreak{} remaining dimensions preserved


\entrylabel{sha256}

6f8553f1\allowbreak{}297050e4\allowbreak{}e9da86ca\allowbreak{}bd81804d\allowbreak{}f0b80013\allowbreak{}f2351c71\allowbreak{}d4bb45c6\allowbreak{}bb765563


\entrylabel{acquired at}

2026-09-07T05:\allowbreak{}10:\allowbreak{}50Z


\subsection*{Hoffman Johnson Junge:\allowbreak{} The hockey-stick conjecture for activated random walk}

arw\_\allowbreak{}hockey\_\allowbreak{}2024


\url{https://arxiv.org/pdf/2411.02541}


\entrylabel{access evidence}

Later primary resolution on intervals


\entrylabel{reading status}

PDF1-3,\allowbreak{} Theorem1


\entrylabel{statement validation}

Interval density curve credited;\allowbreak{} d>=\allowbreak{}2 retained


\entrylabel{sha256}

a65971c6\allowbreak{}f0c5a3ed\allowbreak{}edf5f744\allowbreak{}6e5c7b78\allowbreak{}407872a2\allowbreak{}7bc01b6d\allowbreak{}5e58c7ae\allowbreak{}19b4aa2e


\entrylabel{acquired at}

2026-09-07T05:\allowbreak{}13:\allowbreak{}07Z


\subsection*{Hoffman Richey Son:\allowbreak{} Local density of activated random walk on Z}

arw\_\allowbreak{}local\_\allowbreak{}density\_\allowbreak{}2026


\url{https://arxiv.org/pdf/2601.07816}


\entrylabel{access evidence}

Later primary one-site and symmetry result


\entrylabel{reading status}

PDF1-4 including Theorem1.1 and Proposition1.2


\entrylabel{statement validation}

One-site convergence and subsequential shift invariance are not full joint convergence


\entrylabel{sha256}

fe4bb034\allowbreak{}de53a789\allowbreak{}c024dbd3\allowbreak{}e80dbbbe\allowbreak{}41cdf931\allowbreak{}e7db44e9\allowbreak{}721cd673\allowbreak{}cffc94c6


\entrylabel{acquired at}

2026-09-07T05:\allowbreak{}12:\allowbreak{}59Z


\subsection*{Levine and Silvestri:\allowbreak{} Universality conjectures for activated random walk}

arw\_\allowbreak{}universality\_\allowbreak{}2024


\url{https://arxiv.org/pdf/2306.01698}


\entrylabel{access evidence}

Author primary survey located by focused probability literature search


\entrylabel{reading status}

PDF3-16 statements and definitions read;\allowbreak{} Section9 pair-comparison discussion read


\entrylabel{statement validation}

Nineteen grouped questions;\allowbreak{} no counting of theorems or parameter-value splits


\entrylabel{sha256}

01f7fd60\allowbreak{}a4e31f47\allowbreak{}13fd1b59\allowbreak{}406b135a\allowbreak{}66599d84\allowbreak{}b8757e5e\allowbreak{}6c9bdee1\allowbreak{}4484cd24


\entrylabel{acquired at}

2026-09-07T05:\allowbreak{}09:\allowbreak{}51Z


\subsection*{Basdevant-Blanc-Curien-Singh:\allowbreak{} Fractal properties of the frontier in Poissonian coloring,\allowbreak{}2023}

basdevant\_\allowbreak{}coloring\_\allowbreak{}2023


\url{https://alea.impa.br/articles/v20/20-59.pdf}


\entrylabel{access evidence}

Primary URL discovered through focused literature search2026-09-06;\allowbreak{} substantive title and theorem verification follows acquisition.


\entrylabel{reading status}

unread


\entrylabel{statement validation}

title and relevant theorem read


\entrylabel{sha256}

6757c497\allowbreak{}9446f51b\allowbreak{}3f155d4a\allowbreak{}a40ecd62\allowbreak{}5b7e0a55\allowbreak{}9a255fdf\allowbreak{}8f9345e2\allowbreak{}c13839f0


\entrylabel{acquired at}

2026-09-06T22:\allowbreak{}39:\allowbreak{}35Z


\subsection*{Chorbadzhiyska-Minchev-Savov:\allowbreak{} Beta-splitting trees via homogeneous fragmentations}

beta\_\allowbreak{}fragmentation\_\allowbreak{}2026


\url{https://arxiv.org/pdf/2608.18320}


\entrylabel{access evidence}

Primary later-literature source found during dated status checks


\entrylabel{reading status}

PDF1through5 10 11 20through25 read full other pages unread


\entrylabel{statement validation}

maximum continuous height resolution matched no discrete maximum resolution


\entrylabel{sha256}

85c9fad6\allowbreak{}604d435f\allowbreak{}9f25f141\allowbreak{}f0d068ab\allowbreak{}c9ad9fd9\allowbreak{}b50c1744\allowbreak{}867dd263\allowbreak{}51ebcf2c


\entrylabel{acquired at}

2026-09-07T02:\allowbreak{}14:\allowbreak{}31Z


\subsection*{Brandenberger-Chin-Mossel:\allowbreak{} Asymptotics for the harmonic descent chain and critical beta-splitting trees}

beta\_\allowbreak{}harmonic\_\allowbreak{}asymptotics\_\allowbreak{}2025


\url{https://arxiv.org/pdf/2505.24821}


\entrylabel{access evidence}

Primary later-literature source found during dated status checks


\entrylabel{reading status}

PDF1through5 14through18 read full other pages unread


\entrylabel{statement validation}

explicit resolution claims 5 6 and partial8 matched no full certification


\entrylabel{sha256}

04ff3266\allowbreak{}f9898d46\allowbreak{}f2cfe995\allowbreak{}8d68d477\allowbreak{}ec4cc8ef\allowbreak{}9641f67c\allowbreak{}31a7607d\allowbreak{}a5c58631


\entrylabel{acquired at}

2026-09-07T02:\allowbreak{}14:\allowbreak{}27Z


\subsection*{Aldous-Janson:\allowbreak{} Critical beta-splitting random tree II,\allowbreak{} overview and open problems}

beta\_\allowbreak{}splitting\_\allowbreak{}overview


\url{https://arxiv.org/pdf/2303.02529}


\entrylabel{access evidence}

Primary author or arXiv source located in later-literature checks


\entrylabel{reading status}

PDF1 4 11through13 22through31 35through46 54through58 read full


\entrylabel{statement validation}

numbered questions matched three later claimed resolutions excluded sparse ansatz and drawn length counterchecks recorded


\entrylabel{sha256}

a93a520b\allowbreak{}c7e7f157\allowbreak{}9802bfb8\allowbreak{}7ac88887\allowbreak{}12aadd59\allowbreak{}06c0980c\allowbreak{}a15d2efd\allowbreak{}4fa7735d


\entrylabel{acquired at}

2026-09-07T02:\allowbreak{}02:\allowbreak{}03Z


\subsection*{Dabirian and Subramanian:\allowbreak{} Connectivity of a family of bilateral preference random graphs}

bilateral\_\allowbreak{}connectivity\_\allowbreak{}2025


\url{https://arxiv.org/pdf/2209.05677}


\entrylabel{access evidence}

Author primary preprint;\allowbreak{} verify revision date and conclusion;\allowbreak{} 2025 journal counterpart distinguishes logarithmic threshold from fixed-degree giant


\entrylabel{reading status}

Model and all open questions PDF1-3;\allowbreak{} mean-degree derivation PDF28-30 read


\entrylabel{statement validation}

Four remaining questions matched;\allowbreak{} leading connectivity theorem credited


\entrylabel{sha256}

7321bbc7\allowbreak{}a6431d62\allowbreak{}cdb4cd37\allowbreak{}aa180ac5\allowbreak{}6fc0b190\allowbreak{}5f3a9ab8\allowbreak{}7d0dbbf7\allowbreak{}bb8c172c


\entrylabel{acquired at}

2026-09-07T07:\allowbreak{}25:\allowbreak{}41Z


\subsection*{Baslingker-Krishnapur-Madiman log-concavity Coulomb gases,\allowbreak{}2026 revision}

bkm\_\allowbreak{}logconcavity\_\allowbreak{}2026


\url{https://arxiv.org/pdf/2412.15116}


\entrylabel{access evidence}

Primary workshop follow-up and author paper,\allowbreak{} bounded current-status reconciliation


\entrylabel{reading status}

unread


\entrylabel{statement validation}

not started


\entrylabel{sha256}

dd5a8b57\allowbreak{}8092aeef\allowbreak{}fe39514e\allowbreak{}a2bf305f\allowbreak{}501bad79\allowbreak{}a5f55ca1\allowbreak{}6f6cef1a\allowbreak{}1f5e3340


\entrylabel{acquired at}

2026-09-07T00:\allowbreak{}14:\allowbreak{}22Z


\subsection*{Boursier:\allowbreak{} Decay of correlations and thermodynamic limit for circular Riesz gas}

boursier\_\allowbreak{}riesz\_\allowbreak{}decay\_\allowbreak{}2022


\url{https://arxiv.org/pdf/2209.00396}


\entrylabel{access evidence}

Primary arXiv author preprint found in dated later-literature check


\entrylabel{reading status}

PDF1 through7 read full


\entrylabel{statement validation}

questions or later status matched with explicit scope


\entrylabel{sha256}

03c688f1\allowbreak{}db0d66c8\allowbreak{}4971e144\allowbreak{}dcaa8707\allowbreak{}6c904bc4\allowbreak{}dfdf1423\allowbreak{}f81d0067\allowbreak{}71ee653d


\entrylabel{acquired at}

2026-09-07T01:\allowbreak{}33:\allowbreak{}41Z


\subsection*{Yi Han:\allowbreak{} Brown measure convergence for the spectrum of polynomials in Ginibre matrices}

brown\_\allowbreak{}ginibre\_\allowbreak{}2026


\url{https://arxiv.org/pdf/2606.01664}


\entrylabel{access evidence}

Primary author or publisher public manuscript;\allowbreak{} no license or current-status inference from access


\entrylabel{reading status}

unread


\entrylabel{statement validation}

not started


\entrylabel{sha256}

cbce9cea\allowbreak{}23d33efc\allowbreak{}19032bba\allowbreak{}7a35d79c\allowbreak{}fd080b24\allowbreak{}222d85be\allowbreak{}fe79000f\allowbreak{}66b2f656


\entrylabel{acquired at}

2026-09-06T21:\allowbreak{}27:\allowbreak{}13Z


\subsection*{Berestycki and Sauzedde:\allowbreak{} Brownian paths as loop-decorated SLEs,\allowbreak{}2026}

brownian\_\allowbreak{}loop\_\allowbreak{}sle\_\allowbreak{}2026


\url{https://arxiv.org/pdf/2602.04673}


\entrylabel{access evidence}

Public primary author or workshop PDF;\allowbreak{} exact statements and later status separately checked


\entrylabel{reading status}

unread


\entrylabel{statement validation}

not started


\entrylabel{sha256}

992bf223\allowbreak{}1f5a814a\allowbreak{}8e45ff02\allowbreak{}2df405ef\allowbreak{}47e01ff3\allowbreak{}7743875f\allowbreak{}785265cc\allowbreak{}7d77f2f0


\entrylabel{acquired at}

2026-09-06T19:\allowbreak{}42:\allowbreak{}01Z


\subsection*{Krzysztof Burdzy:\allowbreak{} My favorite open problems}

burdzy\_\allowbreak{}open\_\allowbreak{}20260906


\url{https://sites.math.washington.edu/~burdzy/open_mathjax.php}


\entrylabel{source date}

undated;\allowbreak{} official page includes 2024 update


\entrylabel{access evidence}

Public primary author website or arXiv manuscript


\entrylabel{reading status}

all problem statements read


\entrylabel{statement validation}

author read exact locators


\entrylabel{sha256}

1620000b\allowbreak{}e15308b6\allowbreak{}996ed3c2\allowbreak{}15b10c36\allowbreak{}828e9132\allowbreak{}8bc78314\allowbreak{}6c789e02\allowbreak{}a58ecd86


\entrylabel{acquired at}

2026-09-06T16:\allowbreak{}41:\allowbreak{}36Z


\subsection*{Forward Brownian motion}

burdzy\_\allowbreak{}scheutzow\_\allowbreak{}2014


\url{https://math.washington.edu/sites/math/files/documents/research/1302.6958.pdf}


\entrylabel{source date}

2014;\allowbreak{} arXiv:\allowbreak{}1302.6958


\entrylabel{access evidence}

Public primary author website or arXiv manuscript


\entrylabel{reading status}

all problem statements read


\entrylabel{statement validation}

author read exact locators


\entrylabel{sha256}

8c18daa5\allowbreak{}4bac6ccc\allowbreak{}5da96a3a\allowbreak{}0d3d614c\allowbreak{}41cc4167\allowbreak{}8df671d1\allowbreak{}17b3a861\allowbreak{}68920b9e


\entrylabel{acquired at}

2026-09-06T16:\allowbreak{}41:\allowbreak{}55Z


\subsection*{Deift:\allowbreak{} Some open problems in random matrix theory and integrable systems}

deift\_\allowbreak{}problems\_\allowbreak{}2007


\url{https://arxiv.org/pdf/0712.0849}


\entrylabel{access evidence}

Public primary author or workshop PDF;\allowbreak{} exact statements and later status separately checked


\entrylabel{reading status}

unread


\entrylabel{statement validation}

not started


\entrylabel{sha256}

40d0f656\allowbreak{}4a81d322\allowbreak{}2e304309\allowbreak{}40d89358\allowbreak{}04778e5c\allowbreak{}5d350ce7\allowbreak{}e411a91b\allowbreak{}d3e2d947


\entrylabel{acquired at}

2026-09-06T19:\allowbreak{}48:\allowbreak{}44Z


\subsection*{Deift:\allowbreak{} Some open problems in random matrix theory and integrable systems II}

deift\_\allowbreak{}problems\_\allowbreak{}2017


\url{https://arxiv.org/pdf/1703.04931}


\entrylabel{access evidence}

Public primary author or workshop PDF;\allowbreak{} exact statements and later status separately checked


\entrylabel{reading status}

unread


\entrylabel{statement validation}

not started


\entrylabel{sha256}

675b61a5\allowbreak{}0b2a22a2\allowbreak{}0b2c3d29\allowbreak{}b32044cc\allowbreak{}dd9714b9\allowbreak{}51626220\allowbreak{}28880608\allowbreak{}45c7d5a1


\entrylabel{acquired at}

2026-09-06T19:\allowbreak{}47:\allowbreak{}17Z


\subsection*{Dereudre and Digneaux:\allowbreak{} Rigidity of one dimensional point processes via optimal transport}

dereudre\_\allowbreak{}digneaux\_\allowbreak{}rigidity\_\allowbreak{}2025


\url{https://arxiv.org/pdf/2510.17257}


\entrylabel{access evidence}

Primary arXiv author preprint found in dated later-literature check


\entrylabel{reading status}

PDF1 through3 and24 read full


\entrylabel{statement validation}

questions or later status matched with explicit scope


\entrylabel{sha256}

6f731047\allowbreak{}dbdc8e28\allowbreak{}800bf1dd\allowbreak{}123b3745\allowbreak{}4cadb9d3\allowbreak{}e50bbbf1\allowbreak{}9b18ebad\allowbreak{}be17c147


\entrylabel{acquired at}

2026-09-07T01:\allowbreak{}33:\allowbreak{}44Z


\subsection*{Biskup,\allowbreak{} Gufler,\allowbreak{} Louidor:\allowbreak{} Near-maxima of the two-dimensional DGFF}

dgff\_\allowbreak{}near\_\allowbreak{}maxima\_\allowbreak{}2024


\url{https://arxiv.org/pdf/2010.13939}


\entrylabel{access evidence}

Primary arXiv paper;\allowbreak{} published AIHP 2024;\allowbreak{} checked against Biskup Conjecture16.7


\entrylabel{reading status}

unread


\entrylabel{statement validation}

not started


\entrylabel{sha256}

18ff1110\allowbreak{}a623df7c\allowbreak{}d52fb9a6\allowbreak{}e0d74be3\allowbreak{}9e8ea95a\allowbreak{}407cf0b0\allowbreak{}3528b4e8\allowbreak{}74988979


\entrylabel{acquired at}

2026-09-07T08:\allowbreak{}11:\allowbreak{}57Z


\subsection*{Zhang:\allowbreak{} Boundedness of discounted branching random walks via generic chaining}

discounted\_\allowbreak{}chaining\_\allowbreak{}2026


\url{https://arxiv.org/pdf/2602.18919}


\entrylabel{access evidence}

Primary author or arXiv source located in later-literature checks


\entrylabel{reading status}

PDF1through4 full read remainder unread


\entrylabel{statement validation}

Theorem1 scope matched partial resolution of general Problem31


\entrylabel{sha256}

749db57e\allowbreak{}a84a1753\allowbreak{}f222684d\allowbreak{}acb88678\allowbreak{}582c7c5c\allowbreak{}638faf79\allowbreak{}e2b0b950\allowbreak{}eacfa98a


\entrylabel{acquired at}

2026-09-07T02:\allowbreak{}01:\allowbreak{}57Z


\subsection*{Aidekon-Hu-Shi:\allowbreak{} Boundedness of discounted tree sums}

discounted\_\allowbreak{}trees\_\allowbreak{}2024


\url{https://arxiv.org/pdf/2409.01048}


\entrylabel{access evidence}

Primary author or arXiv source located in later-literature checks


\entrylabel{reading status}

PDF1through6 full read


\entrylabel{statement validation}

Three numbered questions and general critical boundary matched


\entrylabel{sha256}

2b484ddf\allowbreak{}8e97bbf2\allowbreak{}08a69807\allowbreak{}c46aa40d\allowbreak{}3b5a4fdc\allowbreak{}35fbebd4\allowbreak{}5c3db4b4\allowbreak{}63a5cce7


\entrylabel{acquired at}

2026-09-07T02:\allowbreak{}01:\allowbreak{}54Z


\subsection*{Alekseev et al:\allowbreak{} Generalised Linial-Nisan Conjecture Is False for DNFs}

dnf\_\allowbreak{}almost\_\allowbreak{}independence\_\allowbreak{}2025


\url{https://drops.dagstuhl.de/storage/00lipics/lipics-vol339-ccc2025/LIPIcs.CCC.2025.29/LIPIcs.CCC.2025.29.pdf}


\entrylabel{access evidence}

Public primary mathematical source;\allowbreak{} exact status and scope verification


\entrylabel{reading status}

unread


\entrylabel{statement validation}

not started


\entrylabel{sha256}

aac710a1\allowbreak{}3425798b\allowbreak{}1fa0feb3\allowbreak{}b36472e6\allowbreak{}0ad6f988\allowbreak{}cb74dcd0\allowbreak{}f5b75734\allowbreak{}a050905a


\entrylabel{acquired at}

2026-09-06T20:\allowbreak{}47:\allowbreak{}59Z


\subsection*{Dumaz and Malvy:\allowbreak{} Long-Range Correlation of Sine beta Point Process (2026)}

dumaz\_\allowbreak{}malvy\_\allowbreak{}sine\_\allowbreak{}correlations\_\allowbreak{}2026


\url{https://arxiv.org/pdf/2603.15289}


\entrylabel{access evidence}

Primary arXiv author preprint found in dated later-literature check


\entrylabel{reading status}

PDF1 through3 read full


\entrylabel{statement validation}

questions or later status matched with explicit scope


\entrylabel{sha256}

5e32e707\allowbreak{}8fe8d252\allowbreak{}3012ef7b\allowbreak{}9622d619\allowbreak{}e0706518\allowbreak{}9a7584d3\allowbreak{}6c0a38b6\allowbreak{}f5e31010


\entrylabel{acquired at}

2026-09-07T01:\allowbreak{}33:\allowbreak{}42Z


\subsection*{Duminil-Copin thesis:\allowbreak{} Chapter17 Many questions and few answers}

duminil\_\allowbreak{}copin\_\allowbreak{}thesis


\url{https://www.ihes.fr/~duminil/publi/these.pdf}


\entrylabel{access evidence}

Public primary author or workshop PDF;\allowbreak{} exact statements and later status separately checked


\entrylabel{reading status}

unread


\entrylabel{statement validation}

not started


\entrylabel{sha256}

203a2abb\allowbreak{}e1237d63\allowbreak{}ac6018db\allowbreak{}908db1d7\allowbreak{}bf925754\allowbreak{}aa061404\allowbreak{}191bf465\allowbreak{}283018c4


\entrylabel{acquired at}

2026-09-06T19:\allowbreak{}45:\allowbreak{}27Z


\subsection*{Rick Durrett:\allowbreak{} Dynamics on Graphs (10 November 2025)}

durrett\_\allowbreak{}dynamics\_\allowbreak{}graphs\_\allowbreak{}2025


\url{https://sites.math.duke.edu/~rtd/DoG/DoGNov10.pdf}


\entrylabel{access evidence}

Author-hosted primary book manuscript;\allowbreak{} contact process and coevolving systems question intake


\entrylabel{reading status}

unread


\entrylabel{statement validation}

not started


\entrylabel{sha256}

daa94e20\allowbreak{}ea8c019d\allowbreak{}2584cab2\allowbreak{}1484f3d5\allowbreak{}c39f6976\allowbreak{}7b753ca4\allowbreak{}6478ba8d\allowbreak{}8c96f676


\entrylabel{acquired at}

2026-09-07T07:\allowbreak{}23:\allowbreak{}17Z


\subsection*{Moharrami,\allowbreak{} Subramanian,\allowbreak{} Liu and Sundaresan:\allowbreak{} The Erlang Weighted Tree,\allowbreak{} A New Branching Process (v5,\allowbreak{} 27 May 2023)}

erlang\_\allowbreak{}weighted\_\allowbreak{}tree\_\allowbreak{}2023


\url{https://arxiv.org/pdf/2002.03993}


\entrylabel{access evidence}

Author primary preprint;\allowbreak{} Section4 questions1-6;\allowbreak{} model and spectral assumptions checked


\entrylabel{reading status}

Model and main hypotheses PDF1-6;\allowbreak{} all Section4 questions PDF15-17 read


\entrylabel{statement validation}

Four precise questions admitted;\allowbreak{} two broad or protocol-dependent leads held


\entrylabel{sha256}

82971165\allowbreak{}5a26eb25\allowbreak{}fa989764\allowbreak{}58068506\allowbreak{}3fe96249\allowbreak{}a752d7ed\allowbreak{}939efd1a\allowbreak{}ba016985


\entrylabel{acquired at}

2026-09-07T07:\allowbreak{}23:\allowbreak{}15Z


\subsection*{Jardon Sanchez Mellick Poulin Wrobel:\allowbreak{} Exactness and topology of invariant random equivalence relations}

exactness\_\allowbreak{}ire\_\allowbreak{}2025


\url{https://antoinegpoulin.github.io/pdf/exac.pdf}


\entrylabel{access evidence}

Author-hosted primary paper Section6 names Bernoulli Voronoi questions


\entrylabel{reading status}

Section6 PDF17-18 fully read;\allowbreak{} general/\allowbreak{}Bernoulli distinctions checked


\entrylabel{statement validation}

See batch033 registry;\allowbreak{} later results credited and partial proof scoped


\entrylabel{sha256}

8b08b97b\allowbreak{}355a9782\allowbreak{}5f118026\allowbreak{}2c3a6ed8\allowbreak{}548a09d3\allowbreak{}4bae6b4f\allowbreak{}973cbb99\allowbreak{}25e4c70d


\entrylabel{acquired at}

2026-09-07T05:\allowbreak{}50:\allowbreak{}12Z


\subsection*{Gonzalez,\allowbreak{}MacManus,\allowbreak{}Pereyra:\allowbreak{} Further evidence towards the Fourier Entropy-Influence conjecture}

fei\_\allowbreak{}evidence\_\allowbreak{}2026


\url{https://arxiv.org/pdf/2606.00246}


\entrylabel{access evidence}

Public primary mathematical source;\allowbreak{} exact status and scope verification


\entrylabel{reading status}

unread


\entrylabel{statement validation}

not started


\entrylabel{sha256}

0cbdddd0\allowbreak{}c99e9aa6\allowbreak{}c525b60c\allowbreak{}ddfdf292\allowbreak{}b4abccd3\allowbreak{}4aafbf7e\allowbreak{}669d097f\allowbreak{}df11a8f4


\entrylabel{acquired at}

2026-09-06T20:\allowbreak{}45:\allowbreak{}45Z


\subsection*{Foley and McDonald:\allowbreak{} Yaglom limits can depend on the starting state}

foley\_\allowbreak{}mcdonald\_\allowbreak{}yaglom\_\allowbreak{}2017


\url{https://arxiv.org/pdf/1709.07578}


\entrylabel{access evidence}

Primary author or arXiv source;\allowbreak{} statement reading tracked separately


\entrylabel{reading status}

unread


\entrylabel{statement validation}

not started


\entrylabel{sha256}

aa4038df\allowbreak{}a41b3e6b\allowbreak{}95d24c0f\allowbreak{}511b1aec\allowbreak{}4fab0f2d\allowbreak{}4a798af9\allowbreak{}4b1a8ecb\allowbreak{}408e4ee9


\entrylabel{acquired at}

2026-09-07T01:\allowbreak{}49:\allowbreak{}08Z


\subsection*{Auffinger,\allowbreak{} Damron and Hanson:\allowbreak{} 50 years of first-passage percolation}

fpp\_\allowbreak{}survey\_\allowbreak{}2015


\url{https://arxiv.org/pdf/1511.03262}


\entrylabel{access evidence}

Public primary author preprint or workshop institutional host;\allowbreak{} exact statements and current status checked separately


\entrylabel{reading status}

unread


\entrylabel{statement validation}

not started


\entrylabel{sha256}

89ab8416\allowbreak{}bc42814c\allowbreak{}3d529475\allowbreak{}ad7e5608\allowbreak{}a7c67ba3\allowbreak{}68fcf7c6\allowbreak{}192c54eb\allowbreak{}82bf424d


\entrylabel{acquired at}

2026-09-06T18:\allowbreak{}34:\allowbreak{}26Z


\subsection*{Free probability and model theory of tracial W-star algebras}

free\_\allowbreak{}model\_\allowbreak{}theory\_\allowbreak{}2022


\url{https://arxiv.org/pdf/2208.13867}


\entrylabel{access evidence}

Public primary mathematical source;\allowbreak{} exact status and scope verification


\entrylabel{reading status}

unread


\entrylabel{statement validation}

not started


\entrylabel{sha256}

aa9031b2\allowbreak{}bb359626\allowbreak{}b9670abe\allowbreak{}01fc1e93\allowbreak{}01baa925\allowbreak{}adc0c62d\allowbreak{}ba3d041d\allowbreak{}b003c7f0


\entrylabel{acquired at}

2026-09-06T20:\allowbreak{}51:\allowbreak{}41Z


\subsection*{An Study of The Role of Software Project Manger in the Outcome of the Project}

free\_\allowbreak{}positive\_\allowbreak{}unimodality\_\allowbreak{}2020


\url{https://arxiv.org/pdf/2009.13869}


\entrylabel{access evidence}

Primary author or publisher public manuscript;\allowbreak{} no license or current-status inference from access


\entrylabel{reading status}

excluded off domain


\entrylabel{statement validation}

wrong arxiv identifier caught before problem admission


\entrylabel{sha256}

d522f76f\allowbreak{}caca29c5\allowbreak{}0e686a4e\allowbreak{}6d304b36\allowbreak{}2f75cccc\allowbreak{}8fa51ec9\allowbreak{}f4ff0231\allowbreak{}7567d21d


\entrylabel{acquired at}

2026-09-06T21:\allowbreak{}29:\allowbreak{}51Z


\subsection*{Hasebe-Ueda-Wang:\allowbreak{} Log-unimodality for free positive multiplicative Brownian motion}

free\_\allowbreak{}positive\_\allowbreak{}unimodality\_\allowbreak{}2020\_\allowbreak{}correct


\url{https://arxiv.org/pdf/2009.13848}


\entrylabel{access evidence}

Primary author or publisher public manuscript;\allowbreak{} no license or current-status inference from access


\entrylabel{reading status}

unread


\entrylabel{statement validation}

not started


\entrylabel{sha256}

cf993cc0\allowbreak{}fbd26af6\allowbreak{}cc88b15f\allowbreak{}78db0074\allowbreak{}01f61832\allowbreak{}d1cd68f6\allowbreak{}3847440d\allowbreak{}aec5842a


\entrylabel{acquired at}

2026-09-06T21:\allowbreak{}33:\allowbreak{}40Z


\subsection*{Barbieri-Jezernik:\allowbreak{} Strong convergence of random representations of free products of finite groups}

free\_\allowbreak{}products\_\allowbreak{}strong\_\allowbreak{}2026


\url{https://arxiv.org/pdf/2602.11905}


\entrylabel{access evidence}

Primary author or publisher public manuscript;\allowbreak{} no license or current-status inference from access


\entrylabel{reading status}

unread


\entrylabel{statement validation}

not started


\entrylabel{sha256}

362c34fa\allowbreak{}ffdb6dcb\allowbreak{}ed4f493f\allowbreak{}d2a90235\allowbreak{}bd4409eb\allowbreak{}963f949d\allowbreak{}26895020\allowbreak{}aa3e96bc


\entrylabel{acquired at}

2026-09-06T21:\allowbreak{}28:\allowbreak{}23Z


\subsection*{Hasebe,\allowbreak{}Ueda:\allowbreak{} Large time unimodality for classical and free Brownian motions}

free\_\allowbreak{}unimodality\_\allowbreak{}2018


\url{https://arxiv.org/pdf/1710.08240}


\entrylabel{access evidence}

Public primary mathematical source;\allowbreak{} exact status and scope verification


\entrylabel{reading status}

unread


\entrylabel{statement validation}

not started


\entrylabel{sha256}

b49e1f95\allowbreak{}81f220a4\allowbreak{}d0618c2e\allowbreak{}1598cbaf\allowbreak{}75015495\allowbreak{}ff26dd7d\allowbreak{}84d02f5a\allowbreak{}3acfea76


\entrylabel{acquired at}

2026-09-06T20:\allowbreak{}48:\allowbreak{}42Z


\subsection*{Hasebe-Ueda:\allowbreak{} Large time unimodality,\allowbreak{} published version}

free\_\allowbreak{}unimodality\_\allowbreak{}journal\_\allowbreak{}2018


\url{https://alea.math.cnrs.fr/articles/v15/15-15.pdf}


\entrylabel{access evidence}

Primary author or publisher public manuscript;\allowbreak{} no license or current-status inference from access


\entrylabel{reading status}

unread


\entrylabel{statement validation}

not started


\entrylabel{sha256}

4e58df7a\allowbreak{}b9aec6d7\allowbreak{}efc2172f\allowbreak{}89e251f9\allowbreak{}6af44e39\allowbreak{}ca80fd5b\allowbreak{}a724569d\allowbreak{}03e719a5


\entrylabel{acquired at}

2026-09-06T21:\allowbreak{}27:\allowbreak{}11Z


\subsection*{Hasebe-Ueda:\allowbreak{} Unimodality for free multiplicative convolution on the unit circle}

free\_\allowbreak{}unitary\_\allowbreak{}unimodality\_\allowbreak{}2019


\url{https://arxiv.org/pdf/1903.05327}


\entrylabel{access evidence}

Primary author or publisher public manuscript;\allowbreak{} no license or current-status inference from access


\entrylabel{reading status}

unread


\entrylabel{statement validation}

not started


\entrylabel{sha256}

8c8f3754\allowbreak{}5b05cd5a\allowbreak{}dfd3a45a\allowbreak{}f6959c14\allowbreak{}2fdfe617\allowbreak{}85d6f7bb\allowbreak{}2634438a\allowbreak{}296c95b1


\entrylabel{acquired at}

2026-09-06T21:\allowbreak{}29:\allowbreak{}17Z


\subsection*{Frozen percolation on the binary tree is nonendogenous}

frozen\_\allowbreak{}nonendogeny\_\allowbreak{}arxiv


\url{https://arxiv.org/pdf/1910.09213}


\entrylabel{access evidence}

Primary arXiv article located in later-literature search;\allowbreak{} alternate author TLS URL failed.


\entrylabel{reading status}

PDF3 9through17 29 30 full read other intro output partly truncated


\entrylabel{statement validation}

Questions1through10 matched distinct targets no binary nonendogeny recount


\entrylabel{sha256}

e1988718\allowbreak{}8122afc9\allowbreak{}b7da0dda\allowbreak{}2ff3edf4\allowbreak{}39001672\allowbreak{}11ae9577\allowbreak{}52d5f6f7\allowbreak{}795bd2f4


\entrylabel{acquired at}

2026-09-07T02:\allowbreak{}40:\allowbreak{}11Z


\subsection*{Rath-Swart-Szoke:\allowbreak{} A phase transition between endogeny and nonendogeny}

frozen\_\allowbreak{}transition\_\allowbreak{}2021


\url{https://arxiv.org/pdf/2103.14408}


\entrylabel{access evidence}

Primary author or arXiv source located in later-literature checks


\entrylabel{reading status}

unread


\entrylabel{statement validation}

not started


\entrylabel{sha256}

4bb8074c\allowbreak{}4a8d9c3d\allowbreak{}7a6f1346\allowbreak{}6cf0f5ef\allowbreak{}18a798c0\allowbreak{}78ad93ad\allowbreak{}e98c868f\allowbreak{}e783a07c


\entrylabel{acquired at}

2026-09-07T02:\allowbreak{}02:\allowbreak{}00Z


\subsection*{Fuhrman:\allowbreak{} The randomization method in stochastic optimal control,\allowbreak{}2025}

fuhrman\_\allowbreak{}randomization\_\allowbreak{}2025


\url{https://arxiv.org/pdf/2502.06356}


\entrylabel{access evidence}

Primary arXiv survey and open problems


\entrylabel{reading status}

unread


\entrylabel{statement validation}

not started


\entrylabel{sha256}

9f664b98\allowbreak{}ffe20638\allowbreak{}55d40d7f\allowbreak{}0be619b2\allowbreak{}304cef87\allowbreak{}834dc8c6\allowbreak{}02e4720d\allowbreak{}1c5ccf5a


\entrylabel{acquired at}

2026-09-06T23:\allowbreak{}40:\allowbreak{}11Z


\subsection*{Roberts-Rosenthal:\allowbreak{} General state space Markov chains and MCMC algorithms}

general\_\allowbreak{}state\_\allowbreak{}mcmc\_\allowbreak{}2004


\url{https://arxiv.org/pdf/math/0404033}


\entrylabel{access evidence}

Primary author publicly accessible source;\allowbreak{} numbered open questions require later-status verification


\entrylabel{reading status}

unread


\entrylabel{statement validation}

not started


\entrylabel{sha256}

0c7bddb7\allowbreak{}99b48b4c\allowbreak{}f8ce7f16\allowbreak{}4292a96d\allowbreak{}4ebb107b\allowbreak{}279bf6dd\allowbreak{}620a2d22\allowbreak{}6526c743


\entrylabel{acquired at}

2026-09-06T21:\allowbreak{}45:\allowbreak{}16Z


\subsection*{Delgosha-Anantharam:\allowbreak{} A universal low complexity compression algorithm for sparse marked graphs}

graph\_\allowbreak{}compression\_\allowbreak{}2023


\url{https://arxiv.org/pdf/2301.05742}


\entrylabel{access evidence}

Public primary source;\allowbreak{} precise open questions and later resolutions checked separately


\entrylabel{reading status}

unread


\entrylabel{statement validation}

not started


\entrylabel{sha256}

1f9cd89e\allowbreak{}8aab3f23\allowbreak{}e8aa6ce9\allowbreak{}2a92c911\allowbreak{}ffbccb69\allowbreak{}e253443a\allowbreak{}00395f74\allowbreak{}1b93ddb7


\entrylabel{acquired at}

2026-09-06T22:\allowbreak{}16:\allowbreak{}22Z


\subsection*{Ben Green:\allowbreak{}100open problems;\allowbreak{} select probability questions only}

green\_\allowbreak{}100


\url{https://people.maths.ox.ac.uk/greenbj/papers/open-problems.pdf}


\entrylabel{access evidence}

Public primary author or workshop PDF;\allowbreak{} individual open-question statements and later status require separate validation


\entrylabel{reading status}

unread


\entrylabel{statement validation}

not started


\entrylabel{sha256}

e0697124\allowbreak{}5914947f\allowbreak{}152550de\allowbreak{}e59bbb29\allowbreak{}fe0e798f\allowbreak{}0c51b2bc\allowbreak{}2557f824\allowbreak{}c2f9a44a


\entrylabel{acquired at}

2026-09-06T19:\allowbreak{}34:\allowbreak{}40Z


\subsection*{Selected problems in probability theory}

grimmett\_\allowbreak{}2022\_\allowbreak{}selected


\url{https://arxiv.org/pdf/2205.07318}


\entrylabel{access evidence}

Author-deposited open arXiv preprint,\allowbreak{} verified landing page


\entrylabel{reading status}

unread


\entrylabel{statement validation}

not started


\entrylabel{sha256}

21d15bcf\allowbreak{}b892c8b1\allowbreak{}9b4ca593\allowbreak{}f87ba8df\allowbreak{}c58504de\allowbreak{}3c02aaad\allowbreak{}896ec655\allowbreak{}51aee167


\entrylabel{acquired at}

2026-09-06T15:\allowbreak{}29:\allowbreak{}11Z


\subsection*{Jorritsma,\allowbreak{}Maillard,\allowbreak{}Moerters:\allowbreak{} The critical percolation window in growing random graphs}

growing\_\allowbreak{}critical\_\allowbreak{}2025


\url{https://arxiv.org/pdf/2512.18937}


\entrylabel{access evidence}

Public primary arXiv source;\allowbreak{} later-resolution and theorem-scope verification


\entrylabel{reading status}

unread


\entrylabel{statement validation}

not started


\entrylabel{sha256}

2992c568\allowbreak{}6f3ac711\allowbreak{}5a7a10fb\allowbreak{}42a8779b\allowbreak{}4fe44225\allowbreak{}aed7f5c3\allowbreak{}306bf43e\allowbreak{}1ee11212


\entrylabel{acquired at}

2026-09-06T20:\allowbreak{}03:\allowbreak{}18Z


\subsection*{Dhara,\allowbreak{}van der Hofstad,\allowbreak{}van Leeuwaarden,\allowbreak{}Sen:\allowbreak{} Heavy-tailed configuration models at criticality,\allowbreak{}2019version}

heavy\_\allowbreak{}tail\_\allowbreak{}cm\_\allowbreak{}2019


\url{https://arxiv.org/pdf/1612.00650}


\entrylabel{access evidence}

Public primary arXiv source;\allowbreak{} later-resolution and theorem-scope verification


\entrylabel{reading status}

unread


\entrylabel{statement validation}

not started


\entrylabel{sha256}

beff3f5f\allowbreak{}7c9b19c1\allowbreak{}0859bb0c\allowbreak{}5cad470a\allowbreak{}0e7c876e\allowbreak{}5a61544a\allowbreak{}e28b052c\allowbreak{}a6a7da63


\entrylabel{acquired at}

2026-09-06T20:\allowbreak{}03:\allowbreak{}12Z


\subsection*{van der Hofstad:\allowbreak{} Stochastic processes on random graphs,\allowbreak{}monograph draft4August2026}

hofstad\_\allowbreak{}saintflour\_\allowbreak{}grundlehre


\url{https://rhofstad.win.tue.nl/Saint_Flour_SPoRG_Grundlehre.pdf}


\entrylabel{access evidence}

Primary author manuscript;\allowbreak{} title and question locations to be read


\entrylabel{reading status}

unread


\entrylabel{statement validation}

not started


\entrylabel{sha256}

5c9817ae\allowbreak{}9da61be2\allowbreak{}abc4bf3a\allowbreak{}6287940e\allowbreak{}621d023a\allowbreak{}d20bcaf6\allowbreak{}621b4adb\allowbreak{}99c3b062


\entrylabel{acquired at}

2026-09-06T23:\allowbreak{}13:\allowbreak{}51Z


\subsection*{Houdre-Islak:\allowbreak{} A central limit theorem for longest common subsequences in random words,\allowbreak{}2023 revision}

houdre\_\allowbreak{}islak\_\allowbreak{}lcs\_\allowbreak{}2023


\url{https://arxiv.org/pdf/1408.1559}


\entrylabel{access evidence}

Primary arXiv abstract identifies variance lower bound assumption


\entrylabel{reading status}

unread


\entrylabel{statement validation}

not started


\entrylabel{sha256}

5f5f4ee2\allowbreak{}ebd146c5\allowbreak{}9a27615b\allowbreak{}9a44932c\allowbreak{}d206f220\allowbreak{}63622094\allowbreak{}5e5fca6c\allowbreak{}577b67e2


\entrylabel{acquired at}

2026-09-07T00:\allowbreak{}34:\allowbreak{}19Z


\subsection*{Huang-Mao:\allowbreak{} Variational formulas of asymptotic variance for general discrete-time Markov chains}

huang\_\allowbreak{}mao\_\allowbreak{}variance\_\allowbreak{}2020


\url{https://arxiv.org/pdf/2012.13895}


\entrylabel{access evidence}

Primary arXiv full manuscript;\allowbreak{} Theorem4.3 is prior-art lead for P0249


\entrylabel{reading status}

unread


\entrylabel{statement validation}

not started


\entrylabel{sha256}

d519be28\allowbreak{}0282f8ca\allowbreak{}6eaddf92\allowbreak{}49401f53\allowbreak{}50b79cf5\allowbreak{}4d64a826\allowbreak{}e1c4e8b4\allowbreak{}3c08e85e


\entrylabel{acquired at}

2026-09-06T23:\allowbreak{}35:\allowbreak{}27Z


\subsection*{Hyperbolic meets stochastic geometry,\allowbreak{} Oberwolfach57/\allowbreak{}2025}

hyperbolic\_\allowbreak{}stochastic\_\allowbreak{}2025


\url{https://homeweb.unifr.ch/kellerha/pub/OWR-2025.pdf}


\entrylabel{access evidence}

Primary institutional source located and opened by web research.


\entrylabel{reading status}

PDF16-19 survey statements and references read


\entrylabel{statement validation}

See batch033 registry;\allowbreak{} later results credited and partial proof scoped


\entrylabel{sha256}

5b9b9688\allowbreak{}0cd2957f\allowbreak{}cc04a17b\allowbreak{}c708b15d\allowbreak{}eea94300\allowbreak{}f0184c1b\allowbreak{}6bb89ba5\allowbreak{}0a5f887b


\entrylabel{acquired at}

2026-09-07T02:\allowbreak{}49:\allowbreak{}39Z


\subsection*{Buhler Hug Thale:\allowbreak{} Visibility and intersection density for Boolean models in hyperbolic space}

hyperbolic\_\allowbreak{}visibility\_\allowbreak{}2025


\url{https://arxiv.org/pdf/2501.13447}


\entrylabel{access evidence}

Primary article referenced by the 2025 hyperbolic stochastic geometry workshop


\entrylabel{reading status}

Targeted final theorem/\allowbreak{}remarks PDF22-23 read;\allowbreak{} no precise new open question admitted


\entrylabel{statement validation}

See batch033 registry;\allowbreak{} later results credited and partial proof scoped


\entrylabel{sha256}

7d6c0996\allowbreak{}b552fde8\allowbreak{}e1661058\allowbreak{}92ce4b7e\allowbreak{}2643aee8\allowbreak{}012fba9d\allowbreak{}42060f66\allowbreak{}f9d520a1


\entrylabel{acquired at}

2026-09-07T05:\allowbreak{}50:\allowbreak{}16Z


\subsection*{Anastos-Diskin-Lichev-Zhukovskii:\allowbreak{} Diameter and mixing time of the giant component in the percolated hypercube,\allowbreak{}v4May2026}

hypercube\_\allowbreak{}giant\_\allowbreak{}mixing\_\allowbreak{}2026


\url{https://arxiv.org/pdf/2510.13348}


\entrylabel{access evidence}

Public primary author source,\allowbreak{} accessed2026-09-06;\allowbreak{} mathematical content and current status checked separately


\entrylabel{reading status}

unread


\entrylabel{statement validation}

not started


\entrylabel{sha256}

15b7562d\allowbreak{}69b79480\allowbreak{}a7bcfdd2\allowbreak{}3b3245b7\allowbreak{}8cd77ee0\allowbreak{}8f918a5e\allowbreak{}09c851af\allowbreak{}f66eac93


\entrylabel{acquired at}

2026-09-06T22:\allowbreak{}03:\allowbreak{}54Z


\subsection*{D Achille and Khezeli:\allowbreak{} Convergence towards Ideal Poisson Voronoi tessellations with a focus on Diestel Leader graphs}

ideal\_\allowbreak{}convergence\_\allowbreak{}2026


\url{https://arxiv.org/pdf/2606.30517}


\entrylabel{access evidence}

Later primary convergence and distinguishability results identified by current literature search


\entrylabel{reading status}

PDF1-4,\allowbreak{} theorem statements on PDF5,\allowbreak{} Lemma4.9 PDF25,\allowbreak{} Proposition5.14 PDF38,\allowbreak{} PDF46-47 and questions PDF49 read;\allowbreak{} proofs partly read


\entrylabel{statement validation}

See batch033 registry;\allowbreak{} later results credited and partial proof scoped


\entrylabel{sha256}

43971265\allowbreak{}122412bb\allowbreak{}3188ffcb\allowbreak{}79cd0167\allowbreak{}d255e895\allowbreak{}018b102c\allowbreak{}72e7afd5\allowbreak{}ff71ea60


\entrylabel{acquired at}

2026-09-07T05:\allowbreak{}52:\allowbreak{}03Z


\subsection*{D Achille and Thale:\allowbreak{} Face volume densities of positive intensity and ideal Poisson Voronoi tessellations in hyperbolic spaces}

ideal\_\allowbreak{}face\_\allowbreak{}volumes\_\allowbreak{}2026


\url{https://arxiv.org/pdf/2606.26049}


\entrylabel{access evidence}

Later primary resolution of ideal face volume density question


\entrylabel{reading status}

Theorem1.1 PDF5 and Corollary2.1/\allowbreak{}Remark2.2 PDF6 read;\allowbreak{} full proof not reviewed


\entrylabel{statement validation}

See batch033 registry;\allowbreak{} later results credited and partial proof scoped


\entrylabel{sha256}

7a77d101\allowbreak{}713de1d5\allowbreak{}7c582458\allowbreak{}3b04d45c\allowbreak{}c06abd08\allowbreak{}7a8a22d6\allowbreak{}61027574\allowbreak{}51564995


\entrylabel{acquired at}

2026-09-07T05:\allowbreak{}52:\allowbreak{}06Z


\subsection*{Mellick:\allowbreak{} Indistinguishability of cells for the ideal Poisson Voronoi tessellation}

ideal\_\allowbreak{}indistinguishable\_\allowbreak{}2024


\url{https://arxiv.org/pdf/2408.10009}


\entrylabel{access evidence}

Primary resolution explicitly credited in updated source Question7.8


\entrylabel{reading status}

PDF1 Theorem1 and context read;\allowbreak{} full proof not reviewed


\entrylabel{statement validation}

See batch033 registry;\allowbreak{} later results credited and partial proof scoped


\entrylabel{sha256}

d5740d4a\allowbreak{}1db5eca5\allowbreak{}2933ea2f\allowbreak{}35bc5669\allowbreak{}c8e74e53\allowbreak{}94621952\allowbreak{}ed970bcd\allowbreak{}447c6ed8


\entrylabel{acquired at}

2026-09-07T05:\allowbreak{}52:\allowbreak{}12Z


\subsection*{D Achille Curien Enriquez Lyons Unel:\allowbreak{} Ideal Poisson Voronoi tessellations on hyperbolic spaces}

ideal\_\allowbreak{}voronoi\_\allowbreak{}2023


\url{https://arxiv.org/pdf/2303.16831}


\entrylabel{access evidence}

Primary article identified through the hyperbolic workshop references


\entrylabel{reading status}

PDF3,\allowbreak{}35-36,\allowbreak{}39,\allowbreak{}48-49 read;\allowbreak{} targeted Proposition4.4/\allowbreak{}4.6 and reconstruction discussion;\allowbreak{} remainder partly read


\entrylabel{statement validation}

See batch033 registry;\allowbreak{} later results credited and partial proof scoped


\entrylabel{sha256}

54a55bf0\allowbreak{}e27320f6\allowbreak{}23310010\allowbreak{}ac09eaa7\allowbreak{}e8b541c6\allowbreak{}3547f2bf\allowbreak{}66cc904d\allowbreak{}2eb37892


\entrylabel{acquired at}

2026-09-07T05:\allowbreak{}50:\allowbreak{}05Z


\subsection*{Usciati Guillarmou Rhodes Santachiara:\allowbreak{} Probabilistic construction of non compactified imaginary Liouville field theory (2025)}

imaginary\_\allowbreak{}lqg\_\allowbreak{}noncompact\_\allowbreak{}2025


\url{https://arxiv.org/pdf/2505.09390}


\entrylabel{access evidence}

Primary preprint or institution-hosted final PDF found and opened by web research 2026-09-07


\entrylabel{reading status}

PDF1 through3 read full


\entrylabel{statement validation}

questions or later status matched with explicit scope


\entrylabel{sha256}

a9196fd2\allowbreak{}68500e92\allowbreak{}06bd28e3\allowbreak{}089b906d\allowbreak{}80e0801e\allowbreak{}fadf73a7\allowbreak{}cdc7b43a\allowbreak{}c6472c80


\entrylabel{acquired at}

2026-09-07T01:\allowbreak{}25:\allowbreak{}44Z


\subsection*{Can,\allowbreak{}Giardina,\allowbreak{}Giberti,\allowbreak{}van der Hofstad:\allowbreak{} Annealed Ising model on configuration models}

ising\_\allowbreak{}cm\_\allowbreak{}2019


\url{https://arxiv.org/pdf/1904.03664}


\entrylabel{access evidence}

Public primary arXiv source;\allowbreak{} later-resolution and theorem-scope verification


\entrylabel{reading status}

unread


\entrylabel{statement validation}

not started


\entrylabel{sha256}

57d84e12\allowbreak{}9c70d875\allowbreak{}078c9274\allowbreak{}f075aca6\allowbreak{}36fc0b2e\allowbreak{}7325966d\allowbreak{}41feee5c\allowbreak{}9b3962ef


\entrylabel{acquired at}

2026-09-06T20:\allowbreak{}05:\allowbreak{}35Z


\subsection*{Rapid phase ordering of Ising dynamics on Z2,\allowbreak{}2026}

ising\_\allowbreak{}phase\_\allowbreak{}order\_\allowbreak{}2026


\url{https://arxiv.org/pdf/2605.08052}


\entrylabel{access evidence}

Public primary author source,\allowbreak{} accessed2026-09-06;\allowbreak{} mathematical content and current status checked separately


\entrylabel{reading status}

unread


\entrylabel{statement validation}

not started


\entrylabel{sha256}

9d039f66\allowbreak{}d423000d\allowbreak{}5bcb33e1\allowbreak{}4851a272\allowbreak{}489e215e\allowbreak{}3afb6c31\allowbreak{}dcdf6041\allowbreak{}47c87470


\entrylabel{acquired at}

2026-09-06T22:\allowbreak{}03:\allowbreak{}38Z


\subsection*{Kahn:\allowbreak{} Improper Poisson line process as SIRSN in any dimension}

kahn\_\allowbreak{}sirsn\_\allowbreak{}2016


\url{https://arxiv.org/pdf/1503.03976}


\entrylabel{access evidence}

Primary URL discovered through focused literature search2026-09-06;\allowbreak{} substantive title and theorem verification follows acquisition.


\entrylabel{reading status}

unread


\entrylabel{statement validation}

not started


\entrylabel{sha256}

84af8a30\allowbreak{}84755d70\allowbreak{}c20a16f5\allowbreak{}de08e061\allowbreak{}29542058\allowbreak{}a88ff7d2\allowbreak{}f16399e5\allowbreak{}bfafe18c


\entrylabel{acquired at}

2026-09-06T22:\allowbreak{}38:\allowbreak{}27Z


\subsection*{Rhodes-Vargas:\allowbreak{} Two decades of probabilistic approach to Liouville conformal field theory,\allowbreak{}2025}

lcft\_\allowbreak{}rhodes\_\allowbreak{}vargas\_\allowbreak{}2025


\url{https://arxiv.org/pdf/2509.21053}


\entrylabel{access evidence}

Primary arXiv author survey identifies open questions


\entrylabel{reading status}

PDF3 4 17 18 read full other pages not complete


\entrylabel{statement validation}

five explicit questions matched with later partial status limits


\entrylabel{sha256}

b9d1871b\allowbreak{}cf0df63d\allowbreak{}9dae43db\allowbreak{}b4ea0beb\allowbreak{}fd5704db\allowbreak{}c31794e6\allowbreak{}bbb8a492\allowbreak{}a548bd32


\entrylabel{acquired at}

2026-09-07T00:\allowbreak{}40:\allowbreak{}35Z


\subsection*{Leble:\allowbreak{} DLR equations number rigidity and translation invariance for 2DOCP limit points}

leble\_\allowbreak{}coulomb\_\allowbreak{}dlr\_\allowbreak{}2024


\url{https://arxiv.org/pdf/2410.04958}


\entrylabel{access evidence}

Primary arXiv preprint found in dated later-literature check


\entrylabel{reading status}

PDF1 through3 read full


\entrylabel{statement validation}

questions or later status matched with explicit scope


\entrylabel{sha256}

be792edf\allowbreak{}4928b09c\allowbreak{}953bcb28\allowbreak{}bab7abb9\allowbreak{}0e4ea52a\allowbreak{}b648d7a6\allowbreak{}b18c57a3\allowbreak{}886ee660


\entrylabel{acquired at}

2026-09-07T01:\allowbreak{}35:\allowbreak{}30Z


\subsection*{Leble:\allowbreak{} Les log-gas et leurs amis,\allowbreak{} HDR June2026}

leble\_\allowbreak{}hdr\_\allowbreak{}2026


\url{https://tleble.perso.math.cnrs.fr/assets/papers/HDR_Leble.pdf}


\entrylabel{access evidence}

Primary author or arXiv source;\allowbreak{} statement reading tracked separately


\entrylabel{reading status}

PDF21 24 through30 read full other pages only bounded leads


\entrylabel{statement validation}

recent open objectives and known hyperuniformity distinguished


\entrylabel{sha256}

211344a9\allowbreak{}4ab58df0\allowbreak{}784bf8ee\allowbreak{}7a6762fe\allowbreak{}04f23cbd\allowbreak{}5b5a443b\allowbreak{}41e36f31\allowbreak{}fe73de96


\entrylabel{acquired at}

2026-09-07T01:\allowbreak{}48:\allowbreak{}40Z


\subsection*{Lewin:\allowbreak{} Coulomb and Riesz gases,\allowbreak{} the known and the unknown,\allowbreak{}2022}

lewin\_\allowbreak{}riesz\_\allowbreak{}review\_\allowbreak{}2022


\url{https://arxiv.org/pdf/2202.09240}


\entrylabel{access evidence}

Primary author review and journal DOI10.1063/\allowbreak{}5.0086835


\entrylabel{reading status}

PDF13 17 27 28 59 through70 read full other pages not complete


\entrylabel{statement validation}

questions or later status matched with explicit scope


\entrylabel{sha256}

c3c70314\allowbreak{}375884c5\allowbreak{}5f43cf90\allowbreak{}786a0f57\allowbreak{}3484e343\allowbreak{}ae88bf99\allowbreak{}3fc12842\allowbreak{}e075c2b3


\entrylabel{acquired at}

2026-09-07T00:\allowbreak{}34:\allowbreak{}23Z


\subsection*{Bhatia Gwynne Harper:\allowbreak{} Dimension lower bounds in random geometry via Lipschitz functions (2026)}

lqg\_\allowbreak{}lipschitz\_\allowbreak{}2026


\url{https://arxiv.org/pdf/2606.10496}


\entrylabel{access evidence}

Primary preprint or institution-hosted final PDF found and opened by web research 2026-09-07


\entrylabel{reading status}

PDF1 through10 read full


\entrylabel{statement validation}

star family boundary conjecture and net dimension distinguished


\entrylabel{sha256}

4bd3086c\allowbreak{}ade53d05\allowbreak{}a3eafba4\allowbreak{}6b9f0b55\allowbreak{}7e7fb1bc\allowbreak{}a7d4e112\allowbreak{}6db1956e\allowbreak{}8cfba432


\entrylabel{acquired at}

2026-09-07T01:\allowbreak{}25:\allowbreak{}42Z


\subsection*{Lyons:\allowbreak{} Determinantal probability,\allowbreak{} basic properties and conjectures,\allowbreak{} ICM 2014}

lyons\_\allowbreak{}determinantal\_\allowbreak{}2014


\url{https://arxiv.org/pdf/1406.2707}


\entrylabel{source date}

2014-06-10


\entrylabel{access evidence}

arXiv 1406.2707


\entrylabel{reading status}

Primary statements and definitions PDF1,\allowbreak{}5-7,\allowbreak{}22-23 read by integrating author


\entrylabel{statement validation}

Five distinct grouped targets matched;\allowbreak{} implication,\allowbreak{} sublattice and odds/\allowbreak{}probability errors corrected


\entrylabel{sha256}

476170f8\allowbreak{}85edfe9b\allowbreak{}06710112\allowbreak{}5299cee8\allowbreak{}0b22a9e7\allowbreak{}39d082d5\allowbreak{}d7e871bb\allowbreak{}6728dea7


\entrylabel{acquired at}

2026-09-07T06:\allowbreak{}19:\allowbreak{}37Z


\subsection*{Probability on Trees and Networks,\allowbreak{} paperback electronic version 19 August 2026}

lyons\_\allowbreak{}peres\_\allowbreak{}20260819


\url{https://rdlyons.pages.iu.edu/prbtree/book_pb.pdf}


\entrylabel{access evidence}

Authors offer free online edition on their official book page


\entrylabel{reading status}

unread


\entrylabel{statement validation}

not started


\entrylabel{sha256}

5c9a799b\allowbreak{}2333b58c\allowbreak{}7f8b29a6\allowbreak{}08567bd1\allowbreak{}133713c9\allowbreak{}34ec6cf6\allowbreak{}54386dc1\allowbreak{}b01fbaf5


\entrylabel{acquired at}

2026-09-06T15:\allowbreak{}30:\allowbreak{}43Z


\subsection*{Kamath,\allowbreak{}Kumar,\allowbreak{}Manurangsi:\allowbreak{} When Majority Fails,\allowbreak{}2026}

majority\_\allowbreak{}fails\_\allowbreak{}2026


\url{https://arxiv.org/pdf/2604.06590}


\entrylabel{access evidence}

Public primary mathematical source;\allowbreak{} exact status and scope verification


\entrylabel{reading status}

unread


\entrylabel{statement validation}

not started


\entrylabel{sha256}

97ac67db\allowbreak{}723b6205\allowbreak{}edfe7e02\allowbreak{}ef0cf8fc\allowbreak{}dce45fdf\allowbreak{}d4ac0e66\allowbreak{}ec4e776c\allowbreak{}a878923b


\entrylabel{acquired at}

2026-09-06T20:\allowbreak{}45:\allowbreak{}41Z


\subsection*{Mallein:\allowbreak{} Three open problems,\allowbreak{}last modification September1 2017}

mallein\_\allowbreak{}open\_\allowbreak{}2017


\url{https://www.math.univ-toulouse.fr/~bmallein/en/openproblems.html}


\entrylabel{access evidence}

Primary author maintained page;\allowbreak{} author warns some questions may already be solved


\entrylabel{reading status}

unread


\entrylabel{statement validation}

not started


\entrylabel{sha256}

3f73f399\allowbreak{}c315345d\allowbreak{}6ef543fd\allowbreak{}aad8dce9\allowbreak{}aac2bf85\allowbreak{}817bd82d\allowbreak{}a6468565\allowbreak{}82500b4d


\entrylabel{acquired at}

2026-09-06T23:\allowbreak{}41:\allowbreak{}35Z


\subsection*{Levin-Peres:\allowbreak{} Markov Chains and Mixing Times,\allowbreak{} second edition}

markov\_\allowbreak{}mixing\_\allowbreak{}2017


\url{https://pages.uoregon.edu/dlevin/MARKOV/mcmt2e.pdf}


\entrylabel{access evidence}

Primary author publicly accessible source;\allowbreak{} numbered open questions require later-status verification


\entrylabel{reading status}

unread


\entrylabel{statement validation}

not started


\entrylabel{sha256}

9ef39f94\allowbreak{}67d9647f\allowbreak{}f3f5e874\allowbreak{}7b9ce24b\allowbreak{}7a90d13b\allowbreak{}e2f8156f\allowbreak{}bd827b95\allowbreak{}b661a772


\entrylabel{acquired at}

2026-09-06T21:\allowbreak{}43:\allowbreak{}30Z


\subsection*{Addario-Berry and Sen:\allowbreak{} Geometry of the minimal spanning tree of a random3-regular graph}

mst\_\allowbreak{}regular\_\allowbreak{}2018


\url{https://arxiv.org/pdf/1810.03802}


\entrylabel{access evidence}

Public primary arXiv source;\allowbreak{} later-resolution and theorem-scope verification


\entrylabel{reading status}

unread


\entrylabel{statement validation}

not started


\entrylabel{sha256}

c71ea6c0\allowbreak{}6334d753\allowbreak{}c32b8a92\allowbreak{}ed44fd86\allowbreak{}bedf7dcf\allowbreak{}988aba7f\allowbreak{}c16f006e\allowbreak{}2d917494


\entrylabel{acquired at}

2026-09-06T20:\allowbreak{}06:\allowbreak{}25Z


\subsection*{O Donnell:\allowbreak{} Open problems in analysis of Boolean functions}

odonnell\_\allowbreak{}2012\_\allowbreak{}open


\url{https://arxiv.org/pdf/1204.6447}


\entrylabel{access evidence}

Public author-hosted PDF or arXiv primary manuscript;\allowbreak{} open-question passages still require individual validation


\entrylabel{reading status}

unread


\entrylabel{statement validation}

not started


\entrylabel{sha256}

30bbb925\allowbreak{}bf8e05f8\allowbreak{}2d7653f9\allowbreak{}7e110a49\allowbreak{}ca34e980\allowbreak{}227ad5a6\allowbreak{}36ea8479\allowbreak{}4bd702a7


\entrylabel{acquired at}

2026-09-06T18:\allowbreak{}55:\allowbreak{}14Z


\subsection*{del Barrio et al:\allowbreak{} Distributional Limit Theory for Optimal Transport,\allowbreak{}July2026revision}

ot\_\allowbreak{}limit\_\allowbreak{}survey\_\allowbreak{}2026


\url{https://arxiv.org/pdf/2505.19104}


\entrylabel{access evidence}

Primary author or arXiv source;\allowbreak{} statement reading tracked separately


\entrylabel{reading status}

PDF7 8 10 12 20 21 27 through30 read full


\entrylabel{statement validation}

four numbered problems and three explicit optimizer questions matched


\entrylabel{sha256}

f367a677\allowbreak{}aa743d50\allowbreak{}71398827\allowbreak{}fa7ff1d3\allowbreak{}2f0547de\allowbreak{}4da0ddd7\allowbreak{}1a26e95c\allowbreak{}574e10b2


\entrylabel{acquired at}

2026-09-07T01:\allowbreak{}49:\allowbreak{}06Z


\subsection*{Hazra,\allowbreak{}van der Hofstad,\allowbreak{}Ray:\allowbreak{} Percolation on preferential attachment models}

pa\_\allowbreak{}percolation\_\allowbreak{}2023


\url{https://arxiv.org/pdf/2312.14085}


\entrylabel{access evidence}

Public primary arXiv source;\allowbreak{} later-resolution and theorem-scope verification


\entrylabel{reading status}

unread


\entrylabel{statement validation}

not started


\entrylabel{sha256}

931bc236\allowbreak{}2aab4d59\allowbreak{}cd0c8302\allowbreak{}12752d4e\allowbreak{}b612f025\allowbreak{}6e6303b7\allowbreak{}c35a33bc\allowbreak{}c5324698


\entrylabel{acquired at}

2026-09-06T20:\allowbreak{}03:\allowbreak{}14Z


\subsection*{Pak:\allowbreak{} Using stopping times to bound mixing times,\allowbreak{}SODA1999,\allowbreak{}author PDF dated2004}

pak\_\allowbreak{}stopping\_\allowbreak{}1999


\url{https://www.math.ucla.edu/~pak/papers/soda4.pdf}


\entrylabel{access evidence}

Primary author-hosted full paper;\allowbreak{} Theorem3.1 explicitly resolves AF7.18


\entrylabel{reading status}

unread


\entrylabel{statement validation}

not started


\entrylabel{sha256}

0dbd7e5e\allowbreak{}6ca77ffc\allowbreak{}b78e66c3\allowbreak{}26c19b8d\allowbreak{}fd4e8165\allowbreak{}2b163a47\allowbreak{}f4e64bfc\allowbreak{}1d57ffac


\entrylabel{acquired at}

2026-09-06T23:\allowbreak{}35:\allowbreak{}22Z


\subsection*{Pedrotti-Salez local product condition implies cutoff,\allowbreak{}2026}

pedrotti\_\allowbreak{}salez\_\allowbreak{}local\_\allowbreak{}product\_\allowbreak{}2026


\url{https://arxiv.org/pdf/2607.05345}


\entrylabel{access evidence}

Primary author/\allowbreak{}arXiv source used for later-status reconciliation,\allowbreak{} not automatic current-open certification


\entrylabel{reading status}

unread


\entrylabel{statement validation}

not started


\entrylabel{sha256}

add7dba4\allowbreak{}82e5f8d7\allowbreak{}0daee76d\allowbreak{}8279c338\allowbreak{}02442ba2\allowbreak{}2f7d36c7\allowbreak{}f73d5b4f\allowbreak{}8bb0dbf6


\entrylabel{acquired at}

2026-09-07T00:\allowbreak{}00:\allowbreak{}20Z


\subsection*{Pemantle:\allowbreak{} A survey of random processes with reinforcement}

pemantle\_\allowbreak{}2007


\url{https://www2.math.upenn.edu/~pemantle/papers/reinforce.pdf}


\entrylabel{access evidence}

Public author-hosted PDF or arXiv primary manuscript;\allowbreak{} open-question passages still require individual validation


\entrylabel{reading status}

unread


\entrylabel{statement validation}

not started


\entrylabel{sha256}

b1b3fc26\allowbreak{}30408ce2\allowbreak{}90fd1873\allowbreak{}369da86f\allowbreak{}58c3e17d\allowbreak{}2c89d9b8\allowbreak{}26ddef37\allowbreak{}85a13827


\entrylabel{acquired at}

2026-09-06T18:\allowbreak{}54:\allowbreak{}02Z


\subsection*{Pemantle:\allowbreak{} A survey of random processes with reinforcement}

pemantle\_\allowbreak{}reinforcement\_\allowbreak{}2007


\url{https://www2.math.upenn.edu/~pemantle/papers/reinforce.pdf}


\entrylabel{access evidence}

Primary author-hosted collection or maintained status page


\entrylabel{reading status}

unread


\entrylabel{statement validation}

not started


\entrylabel{sha256}

b1b3fc26\allowbreak{}30408ce2\allowbreak{}90fd1873\allowbreak{}369da86f\allowbreak{}58c3e17d\allowbreak{}2c89d9b8\allowbreak{}26ddef37\allowbreak{}85a13827


\entrylabel{acquired at}

2026-09-07T01:\allowbreak{}59:\allowbreak{}51Z


\subsection*{Peralta:\allowbreak{} Rational arrival processes with strictly positive densities need not be Markovian (2026)}

peralta\_\allowbreak{}rap\_\allowbreak{}2026


\url{https://arxiv.org/pdf/2603.28047}


\entrylabel{access evidence}

Primary author preprint or institutional publisher PDF,\allowbreak{} found by web search 2026-09-07


\entrylabel{reading status}

complete ten page preprint read


\entrylabel{statement validation}

historical conjecture and counterexample matched


\entrylabel{sha256}

3898bc8c\allowbreak{}2e96bc05\allowbreak{}8c3295a1\allowbreak{}3f233ce4\allowbreak{}044775aa\allowbreak{}b5dba9a0\allowbreak{}0c76ed5c\allowbreak{}0d36dc67


\entrylabel{acquired at}

2026-09-07T01:\allowbreak{}22:\allowbreak{}28Z


\subsection*{Yuval Peres:\allowbreak{} related open Markov-chain problems}

peres\_\allowbreak{}current\_\allowbreak{}open


\url{https://www.yuval-peres-books.com/markov-chains-and-mixing-times/}


\entrylabel{access evidence}

Primary author publicly accessible source;\allowbreak{} numbered open questions require later-status verification


\entrylabel{reading status}

unread


\entrylabel{statement validation}

not started


\entrylabel{sha256}

de00d475\allowbreak{}3d784616\allowbreak{}76088e84\allowbreak{}a3029163\allowbreak{}cb703203\allowbreak{}1d4ef95f\allowbreak{}c24338fc\allowbreak{}d636ebe8


\entrylabel{acquired at}

2026-09-06T21:\allowbreak{}45:\allowbreak{}37Z


\subsection*{Preater:\allowbreak{} A species of voter model driven by immigration}

preater\_\allowbreak{}voter\_\allowbreak{}2009


\url{https://www.stat.berkeley.edu/users/aldous/Research/OP/preater.pdf}


\entrylabel{access evidence}

Primary URL discovered through focused literature search2026-09-06;\allowbreak{} substantive title and theorem verification follows acquisition.


\entrylabel{reading status}

unread


\entrylabel{statement validation}

not started


\entrylabel{sha256}

eec8b012\allowbreak{}712b13c9\allowbreak{}2314fdfe\allowbreak{}db750549\allowbreak{}897b94e0\allowbreak{}80e85eec\allowbreak{}96831c8c\allowbreak{}b50b793a


\entrylabel{acquired at}

2026-09-06T22:\allowbreak{}38:\allowbreak{}50Z


\subsection*{Qu and Valko:\allowbreak{} On the pair correlation function of Sine beta process}

qu\_\allowbreak{}valko\_\allowbreak{}sine\_\allowbreak{}pair\_\allowbreak{}2025


\url{https://arxiv.org/pdf/2509.15446}


\entrylabel{access evidence}

Primary arXiv preprint found in dated later-literature check


\entrylabel{reading status}

unread


\entrylabel{statement validation}

not started


\entrylabel{sha256}

bd7c0709\allowbreak{}3b71f56f\allowbreak{}7705a4e3\allowbreak{}98b92ee1\allowbreak{}f5d478c3\allowbreak{}a7e5c696\allowbreak{}40a036bc\allowbreak{}540eb21e


\entrylabel{acquired at}

2026-09-07T01:\allowbreak{}35:\allowbreak{}32Z


\subsection*{Queueing Systems volume100 issue3-4:\allowbreak{}100 views on queues,\allowbreak{}primary issue index}

queues\_\allowbreak{}100\_\allowbreak{}issue\_\allowbreak{}2022


\url{https://link.springer.com/journal/11134/volumes-and-issues/100-3}


\entrylabel{access evidence}

Primary journal issue page lists102 contributions including editorial and notes


\entrylabel{reading status}

unread


\entrylabel{statement validation}

not started


\entrylabel{sha256}

37abcf1f\allowbreak{}886bf1d8\allowbreak{}716863eb\allowbreak{}949e3f17\allowbreak{}7763795f\allowbreak{}3c2f78ae\allowbreak{}6d0e127b\allowbreak{}05eeb520


\entrylabel{acquired at}

2026-09-07T00:\allowbreak{}40:\allowbreak{}31Z


\subsection*{100 views on queues:\allowbreak{}primary issue index page2}

queues\_\allowbreak{}100\_\allowbreak{}issue\_\allowbreak{}page2\_\allowbreak{}2022


\url{https://link.springer.com/journal/11134/volumes-and-issues/100-3?page=2}


\entrylabel{access evidence}

Observed primary issue pagination


\entrylabel{reading status}

unread


\entrylabel{statement validation}

not started


\entrylabel{sha256}

95e06ca0\allowbreak{}b50bd2ce\allowbreak{}4474575a\allowbreak{}2fb1529d\allowbreak{}ef2c58ed\allowbreak{}5c27f5e4\allowbreak{}7460b8b4\allowbreak{}0a9b2ce5


\entrylabel{acquired at}

2026-09-07T00:\allowbreak{}42:\allowbreak{}28Z


\subsection*{Harchol-Balter:\allowbreak{} Open problems in queueing theory inspired by datacenter computing}

queues\_\allowbreak{}computing\_\allowbreak{}frontiers\_\allowbreak{}2021


\url{https://www.cs.cmu.edu/~harchol/Papers/QUESTA21.pdf}


\entrylabel{access evidence}

Primary author-hosted journal PDF located through source search.


\entrylabel{reading status}

PDF1-17 and20-28 read;\allowbreak{} workload descriptive section18-19 and bibliography not fully read


\entrylabel{statement validation}

Batch034 source scopes and remaining questions matched


\entrylabel{sha256}

de92bbb3\allowbreak{}2c4ffd8d\allowbreak{}015f8d70\allowbreak{}2ee05306\allowbreak{}a31e9b86\allowbreak{}892ca18b\allowbreak{}d23e8f28\allowbreak{}de436049


\entrylabel{acquired at}

2026-09-07T02:\allowbreak{}54:\allowbreak{}43Z


\subsection*{Li Ju Eryilmaz Shroff:\allowbreak{} Optimal Parallel Scheduling under Concave Speedup Functions}

queues\_\allowbreak{}concave\_\allowbreak{}resolution\_\allowbreak{}2025


\url{https://arxiv.org/pdf/2509.01811}


\entrylabel{access evidence}

Primary author manuscript resolving a source-survey scheduling question


\entrylabel{reading status}

Model and assumptions PDF1-4,\allowbreak{} structural proof PDF5-6,\allowbreak{} algorithm discussion PDF14-16 read;\allowbreak{} full proof not independently reviewed


\entrylabel{statement validation}

Batch034 source scopes and remaining questions matched


\entrylabel{sha256}

070f9b54\allowbreak{}17702395\allowbreak{}fee75807\allowbreak{}be54c8cf\allowbreak{}c74ba185\allowbreak{}27e49aae\allowbreak{}d9cbe322\allowbreak{}28cdb1f4


\entrylabel{acquired at}

2026-09-07T06:\allowbreak{}25:\allowbreak{}09Z


\subsection*{Whitt-You:\allowbreak{} New decomposition approximations for queueing networks}

queues\_\allowbreak{}decomposition\_\allowbreak{}2022


\url{https://www.columbia.edu/~ww2040/view_RQNA_format_020222.pdf}


\entrylabel{access evidence}

Primary author or institutional PDF located through the100 views on queues issue registry and web search.


\entrylabel{reading status}

All three PDF pages read;\allowbreak{} four grouped targets from Sections2-3


\entrylabel{statement validation}

See batch032 registry;\allowbreak{} no present-openness certification


\entrylabel{sha256}

e99f4793\allowbreak{}059b926a\allowbreak{}382c4b9b\allowbreak{}7ee241e2\allowbreak{}0593518a\allowbreak{}c687c4e9\allowbreak{}f1842aa0\allowbreak{}016e768d


\entrylabel{acquired at}

2026-09-07T02:\allowbreak{}50:\allowbreak{}58Z


\subsection*{Moyal-Perry:\allowbreak{} Many server limits with dependent service and patience times,\allowbreak{}2022}

queues\_\allowbreak{}dependent\_\allowbreak{}patience\_\allowbreak{}2022


\url{https://users.iems.northwestern.edu/~perry/QueuesOpen.pdf}


\entrylabel{access evidence}

Primary author or institutional full-text archive found from journal collection


\entrylabel{reading status}

complete three page note read


\entrylabel{statement validation}

numbered or explicit problem matched with scope caveats


\entrylabel{sha256}

255ed3f7\allowbreak{}cd8a7be6\allowbreak{}1c1ac9de\allowbreak{}b0033726\allowbreak{}78fc76d1\allowbreak{}3f182a0b\allowbreak{}caa05fd6\allowbreak{}287eb39e


\entrylabel{acquired at}

2026-09-07T00:\allowbreak{}46:\allowbreak{}38Z


\subsection*{Williams Wang Harchol-Balter:\allowbreak{} An Upper Bound on the M M k Queue With Deterministic Setup Times}

queues\_\allowbreak{}deterministic\_\allowbreak{}setup\_\allowbreak{}2025


\url{https://arxiv.org/pdf/2512.16854}


\entrylabel{access evidence}

Later primary bounds and remaining queue setup questions


\entrylabel{reading status}

Model PDF8-9,\allowbreak{} related work PDF10-11,\allowbreak{} mean-bound statements PDF13-14,\allowbreak{} all concluding questions PDF39-40 read


\entrylabel{statement validation}

Batch034 source scopes and remaining questions matched


\entrylabel{sha256}

9e3b1d71\allowbreak{}90277e73\allowbreak{}db130970\allowbreak{}b87d0487\allowbreak{}a288ff83\allowbreak{}f1c2be80\allowbreak{}dceb4e20\allowbreak{}c890b516


\entrylabel{acquired at}

2026-09-07T06:\allowbreak{}25:\allowbreak{}33Z


\subsection*{Schol Vlasiou Zwart:\allowbreak{} Extreme values for the waiting time in large fork join queues}

queues\_\allowbreak{}forkjoin\_\allowbreak{}extremes\_\allowbreak{}2023


\url{https://arxiv.org/pdf/2309.08373}


\entrylabel{access evidence}

Later primary asymptotic fork-join theorem;\allowbreak{} finite-size scope checked separately


\entrylabel{reading status}

Main theorem assumptions and scope read;\allowbreak{} finite-n joint-law question retained


\entrylabel{statement validation}

Batch034 source scopes and remaining questions matched


\entrylabel{sha256}

c9ba3251\allowbreak{}5a4cf1fd\allowbreak{}019d4370\allowbreak{}7599369c\allowbreak{}0b0464b3\allowbreak{}ed64da75\allowbreak{}df4360cc\allowbreak{}2a4a2191


\entrylabel{acquired at}

2026-09-07T06:\allowbreak{}46:\allowbreak{}03Z


\subsection*{Ravner:\allowbreak{} Queue input estimation from discrete workload observations}

queues\_\allowbreak{}input\_\allowbreak{}estimation\_\allowbreak{}2022


\url{https://arxiv.org/pdf/2202.10750}


\entrylabel{access evidence}

Primary author or institutional PDF located through the100 views on queues issue registry and web search.


\entrylabel{reading status}

All three PDF pages read;\allowbreak{} efficient estimation distinguished from consistency


\entrylabel{statement validation}

See batch032 registry;\allowbreak{} no present-openness certification


\entrylabel{sha256}

0667f64a\allowbreak{}dc9adc80\allowbreak{}73412734\allowbreak{}5a93caef\allowbreak{}ce292a83\allowbreak{}338f7883\allowbreak{}261804a9\allowbreak{}e816de5c


\entrylabel{acquired at}

2026-09-07T02:\allowbreak{}50:\allowbreak{}55Z


\subsection*{Vlasiou:\allowbreak{} Queues on interacting networks}

queues\_\allowbreak{}interacting\_\allowbreak{}2022


\url{https://arxiv.org/pdf/2212.07945}


\entrylabel{access evidence}

Primary author or institutional PDF located through the100 views on queues issue registry and web search.


\entrylabel{reading status}

All three pages read;\allowbreak{} broad questions held for concrete model specification


\entrylabel{statement validation}

See batch032 registry;\allowbreak{} no present-openness certification


\entrylabel{sha256}

836ba3f8\allowbreak{}21f57076\allowbreak{}72048400\allowbreak{}1cd92b04\allowbreak{}076167b5\allowbreak{}2f713c86\allowbreak{}19ab13e6\allowbreak{}8cb3d79b


\entrylabel{acquired at}

2026-09-07T02:\allowbreak{}51:\allowbreak{}03Z


\subsection*{Boon and Winands:\allowbreak{} Open problems for critically loaded k-limited polling systems (2022)}

queues\_\allowbreak{}klimited\_\allowbreak{}uva\_\allowbreak{}2022


\url{https://pure.uva.nl/ws/files/120928955/Open_problems_for_critically_loaded_k_limited_polling_systems.pdf}


\entrylabel{access evidence}

Primary preprint or institution-hosted final PDF found and opened by web research 2026-09-07


\entrylabel{reading status}

complete three page note read plus cover


\entrylabel{statement validation}

Conjecture1 Problem1 grouped Problem2 distinct


\entrylabel{sha256}

08caacdd\allowbreak{}beea015b\allowbreak{}9b802178\allowbreak{}6b299d0c\allowbreak{}0cad8090\allowbreak{}067e9932\allowbreak{}c789a81f\allowbreak{}bcccab28


\entrylabel{acquired at}

2026-09-07T01:\allowbreak{}25:\allowbreak{}41Z


\subsection*{Wang Harchol-Balter Jiang Scheller-Wolf Srikant:\allowbreak{} Delay Asymptotics and Bounds for Multi Task Parallel Jobs}

queues\_\allowbreak{}limited\_\allowbreak{}forkjoin\_\allowbreak{}2018


\url{https://www.cs.cmu.edu/~harchol/Papers/Performance18b.pdf}


\entrylabel{access evidence}

Original primary source for the growing fork-join independence exponent


\entrylabel{reading status}

unread


\entrylabel{statement validation}

not started


\entrylabel{sha256}

15409983\allowbreak{}9f75dc2b\allowbreak{}7f80563f\allowbreak{}82b0af02\allowbreak{}2041d0c0\allowbreak{}40838b0d\allowbreak{}16e72c4c\allowbreak{}cb275c41


\entrylabel{acquired at}

2026-09-07T06:\allowbreak{}52:\allowbreak{}10Z


\subsection*{The Multiserver Job Stochastic Recurrence Equation for Cloud Computing Performance Evaluation}

queues\_\allowbreak{}msj\_\allowbreak{}recurrence\_\allowbreak{}2026


\url{https://arxiv.org/pdf/2601.20653}


\entrylabel{access evidence}

Later primary general FCFS stability and sampling results


\entrylabel{reading status}

PDF1-4,\allowbreak{} Theorem5.1 and discussion PDF11-12,\allowbreak{} concluding questions PDF22 read


\entrylabel{statement validation}

Batch034 source scopes and remaining questions matched


\entrylabel{sha256}

5eb058fb\allowbreak{}ca320627\allowbreak{}1b3fb04d\allowbreak{}dd96c284\allowbreak{}7742fc0b\allowbreak{}4ce5791b\allowbreak{}9dba05ac\allowbreak{}6222cb5a


\entrylabel{acquired at}

2026-09-07T06:\allowbreak{}25:\allowbreak{}16Z


\subsection*{Harchol-Balter:\allowbreak{} The multiserver job queueing model}

queues\_\allowbreak{}multiserver\_\allowbreak{}jobs\_\allowbreak{}2022


\url{https://www.cs.cmu.edu/~harchol/Papers/QUESTA22.pdf}


\entrylabel{access evidence}

Primary author-hosted journal PDF located through source search.


\entrylabel{reading status}

All three pages read;\allowbreak{} intake held pending later 2023-2026 results


\entrylabel{statement validation}

See batch032 registry;\allowbreak{} no present-openness certification


\entrylabel{sha256}

e32b7798\allowbreak{}bdd59906\allowbreak{}d84cc183\allowbreak{}64d6f2ec\allowbreak{}73e98dc8\allowbreak{}8766676d\allowbreak{}eda23f19\allowbreak{}6d07a9ce


\entrylabel{acquired at}

2026-09-07T02:\allowbreak{}54:\allowbreak{}38Z


\subsection*{Dong-Ibrahim:\allowbreak{} Shortest-job-first scheduling in many-server queues with impatient customers and noisy estimates}

queues\_\allowbreak{}noisy\_\allowbreak{}sjf\_\allowbreak{}2025


\url{https://www.roubaibrahim.com/SRPT_with_errors.pdf}


\entrylabel{access evidence}

Primary author or institutional full text found through source and later-literature lookup.


\entrylabel{reading status}

PDF1-5 introduction and contributions read;\allowbreak{} formal proofs not reviewed


\entrylabel{statement validation}

See batch032 registry;\allowbreak{} no present-openness certification


\entrylabel{sha256}

ea0555ad\allowbreak{}2d03d0aa\allowbreak{}749bbc2b\allowbreak{}717390e5\allowbreak{}4a04d410\allowbreak{}ed081bba\allowbreak{}0fd9cbae\allowbreak{}f39bfa60


\entrylabel{acquired at}

2026-09-07T02:\allowbreak{}54:\allowbreak{}07Z


\subsection*{Ibrahim:\allowbreak{} Personalized scheduling in service systems}

queues\_\allowbreak{}personalized\_\allowbreak{}2022


\url{https://discovery.ucl.ac.uk/10147615/1/problem_Ibrahim.pdf}


\entrylabel{access evidence}

Primary author or institutional PDF located through the100 views on queues issue registry and web search.


\entrylabel{reading status}

All source body read in resumed intake;\allowbreak{} later noisy-SJF theorem scope credited


\entrylabel{statement validation}

See batch032 registry;\allowbreak{} no present-openness certification


\entrylabel{sha256}

14a778e4\allowbreak{}6d36062c\allowbreak{}68666c18\allowbreak{}774ec257\allowbreak{}41a07190\allowbreak{}5afc1f71\allowbreak{}34b4b137\allowbreak{}5bd38785


\entrylabel{acquired at}

2026-09-07T02:\allowbreak{}50:\allowbreak{}59Z


\subsection*{Antunes-Jacinto-Pacheco:\allowbreak{} Statistical inference in queueing networks with probing information}

queues\_\allowbreak{}probing\_\allowbreak{}2022


\url{https://rdpc.uevora.pt/bitstream/10174/32450/3/paper_Q65.pdf}


\entrylabel{access evidence}

Primary author or institutional PDF located through the100 views on queues issue registry and web search.


\entrylabel{reading status}

PDF1-2 fully read;\allowbreak{} PDF3 partly read;\allowbreak{} exact n=\allowbreak{}1 identifiability limitation checked


\entrylabel{statement validation}

See batch032 registry;\allowbreak{} no present-openness certification


\entrylabel{sha256}

570e83f2\allowbreak{}28861542\allowbreak{}391b2d49\allowbreak{}708be993\allowbreak{}419bc225\allowbreak{}55f113f6\allowbreak{}cc6830b1\allowbreak{}21731b76


\entrylabel{acquired at}

2026-09-07T02:\allowbreak{}51:\allowbreak{}02Z


\subsection*{Pollett:\allowbreak{} Quasi stationary distributions for queueing and other models,\allowbreak{}2022}

queues\_\allowbreak{}qsd\_\allowbreak{}pollett\_\allowbreak{}2022


\url{https://arxiv.org/pdf/2202.05428}


\entrylabel{access evidence}

Primary author or institutional full-text archive found from journal collection


\entrylabel{reading status}

complete three page primary read


\entrylabel{statement validation}

conjecture ii checked against full hypotheses and discussion


\entrylabel{sha256}

12fa667b\allowbreak{}bfe2b2b1\allowbreak{}880aa7fd\allowbreak{}f0c3e110\allowbreak{}42fe4472\allowbreak{}3dc56114\allowbreak{}01ef32a6\allowbreak{}059d0ebd


\entrylabel{acquired at}

2026-09-07T00:\allowbreak{}46:\allowbreak{}35Z


\subsection*{Gardner:\allowbreak{} Correlation in redundancy systems}

queues\_\allowbreak{}redundancy\_\allowbreak{}2022


\url{https://kgardner.people.amherst.edu/papers/gardner_questa_100views.pdf}


\entrylabel{access evidence}

Primary author or institutional full text found through source and later-literature lookup.


\entrylabel{reading status}

All three PDF pages read;\allowbreak{} all five questions with prior special cases credited


\entrylabel{statement validation}

See batch032 registry;\allowbreak{} no present-openness certification


\entrylabel{sha256}

9da03e26\allowbreak{}8c67573d\allowbreak{}97b7ad29\allowbreak{}ca9f7c48\allowbreak{}365554e3\allowbreak{}09a3f81b\allowbreak{}054a5e9c\allowbreak{}f4380553


\entrylabel{acquired at}

2026-09-07T02:\allowbreak{}54:\allowbreak{}03Z


\subsection*{Hellemans-Van Houdt:\allowbreak{} Open problems in refined mean field approximations,\allowbreak{}2022}

queues\_\allowbreak{}refined\_\allowbreak{}meanfield\_\allowbreak{}2022


\url{https://link.springer.com/article/10.1007/s11134-022-09756-9}


\entrylabel{access evidence}

Primary issue article link


\entrylabel{reading status}

landing page only full text not available


\entrylabel{statement validation}

not started


\entrylabel{sha256}

d59060ac\allowbreak{}56cdf58c\allowbreak{}3ba8a79d\allowbreak{}c2e22588\allowbreak{}b07287b1\allowbreak{}56cbe6b6\allowbreak{}d4e2b9a9\allowbreak{}4917f56b


\entrylabel{acquired at}

2026-09-07T00:\allowbreak{}42:\allowbreak{}32Z


\subsection*{Hellemans-Van Houdt:\allowbreak{} Open problems in refined mean field approximations,\allowbreak{}2022}

queues\_\allowbreak{}refined\_\allowbreak{}meanfield\_\allowbreak{}pdf\_\allowbreak{}2022


\url{https://win.uantwerpen.be/~vanhoudt/papers/2022/Questa_100.pdf}


\entrylabel{access evidence}

Primary author or institutional full-text archive found from journal collection


\entrylabel{reading status}

complete three page note read


\entrylabel{statement validation}

numbered or explicit problem matched with scope caveats


\entrylabel{sha256}

4e04451c\allowbreak{}087f3e70\allowbreak{}8a5c3efa\allowbreak{}e24d316d\allowbreak{}809aa634\allowbreak{}33e49eea\allowbreak{}dfb7b974\allowbreak{}0c19ec94


\entrylabel{acquired at}

2026-09-07T00:\allowbreak{}46:\allowbreak{}29Z


\subsection*{Leskela:\allowbreak{} Ross second conjecture and supermodular stochastic ordering,\allowbreak{}2022}

queues\_\allowbreak{}ross\_\allowbreak{}second\_\allowbreak{}2022


\url{https://acris.aalto.fi/ws/portalfiles/portal/89616571/Ross_s_second_conjecture_and_supermodular_stochastic_ordering.pdf}


\entrylabel{access evidence}

Primary author or institutional full-text archive found from journal collection


\entrylabel{reading status}

complete three page note read


\entrylabel{statement validation}

numbered or explicit problem matched with scope caveats


\entrylabel{sha256}

feb092d9\allowbreak{}f72c91f0\allowbreak{}e6430668\allowbreak{}82613ab6\allowbreak{}6339666d\allowbreak{}341427af\allowbreak{}042d3aa6\allowbreak{}7436e7d3


\entrylabel{acquired at}

2026-09-07T00:\allowbreak{}46:\allowbreak{}33Z


\subsection*{Grosof Hurtado-Lange:\allowbreak{} Outperforming Multiserver SRPT at All Loads}

queues\_\allowbreak{}sek\_\allowbreak{}smod\_\allowbreak{}2025


\url{https://arxiv.org/pdf/2510.25963}


\entrylabel{access evidence}

Later primary resolution of the SRPT improvement-existence question


\entrylabel{reading status}

PDF1-3 and Theorem4.1 read;\allowbreak{} improvement credited,\allowbreak{} optimality not inferred


\entrylabel{statement validation}

Batch034 source scopes and remaining questions matched


\entrylabel{sha256}

8ca04bd3\allowbreak{}451a04d1\allowbreak{}402f7368\allowbreak{}35fd00f9\allowbreak{}fe002af8\allowbreak{}a0f00cda\allowbreak{}0c84839d\allowbreak{}183452fb


\entrylabel{acquired at}

2026-09-07T06:\allowbreak{}45:\allowbreak{}50Z


\subsection*{Aalto:\allowbreak{} Minimizing the mean slowdown in a single-server queue}

queues\_\allowbreak{}slowdown\_\allowbreak{}2022


\url{https://acris.aalto.fi/ws/portalfiles/portal/79466983/Aalto_Queueing_Systems_2022.pdf}


\entrylabel{access evidence}

Primary author or institutional full text found through source and later-literature lookup.


\entrylabel{reading status}

All four PDF pages read;\allowbreak{} Conjecture1 matched to later Corollary3


\entrylabel{statement validation}

See batch032 registry;\allowbreak{} no present-openness certification


\entrylabel{sha256}

b580842b\allowbreak{}95e365ab\allowbreak{}0f9c49af\allowbreak{}e2dc0f54\allowbreak{}1927855a\allowbreak{}7e1629dd\allowbreak{}4b74815d\allowbreak{}40709cf5


\entrylabel{acquired at}

2026-09-07T02:\allowbreak{}54:\allowbreak{}05Z


\subsection*{Aalto and Scully:\allowbreak{} Minimizing the mean slowdown in the M/\allowbreak{}G/\allowbreak{}1 queue}

queues\_\allowbreak{}slowdown\_\allowbreak{}resolution\_\allowbreak{}2023


\url{https://aaltodoc.aalto.fi/server/api/core/bitstreams/a9d987aa-1806-45f5-a7be-7bce2d8cc8f0/content}


\entrylabel{access evidence}

Primary institutional open article resolves 2022 conjecture;\allowbreak{} verify theorem correspondence


\entrylabel{reading status}

PDF12-15 fully read plus targeted theorem locations;\allowbreak{} remainder not fully reviewed


\entrylabel{statement validation}

See batch032 registry;\allowbreak{} no present-openness certification


\entrylabel{sha256}

6166b192\allowbreak{}f1b40fd5\allowbreak{}ed94f265\allowbreak{}9a1ec6ae\allowbreak{}ca06df71\allowbreak{}9d56c66a\allowbreak{}ba145d06\allowbreak{}e590506c


\entrylabel{acquired at}

2026-09-07T05:\allowbreak{}39:\allowbreak{}06Z


\subsection*{Zwart:\allowbreak{} Conjectures on Symmetric Queues in Heavy Traffic}

queues\_\allowbreak{}symmetric\_\allowbreak{}2022


\url{https://arxiv.org/pdf/2202.02366}


\entrylabel{access evidence}

Primary author or institutional PDF located through the100 views on queues issue registry and web search.


\entrylabel{reading status}

All three PDF pages read;\allowbreak{} two numbered problems and separately named busy-cycle tail question


\entrylabel{statement validation}

See batch032 registry;\allowbreak{} no present-openness certification


\entrylabel{sha256}

0c1d9c15\allowbreak{}dd9bf01e\allowbreak{}3f568a18\allowbreak{}eef8839b\allowbreak{}3251c750\allowbreak{}b5183722\allowbreak{}8a887cd7\allowbreak{}cbaac19f


\entrylabel{acquired at}

2026-09-07T02:\allowbreak{}50:\allowbreak{}54Z


\subsection*{Can and van der Hofstad:\allowbreak{} Random cluster models on random graphs}

random\_\allowbreak{}cluster\_\allowbreak{}graphs\_\allowbreak{}2025


\url{https://arxiv.org/pdf/2503.17636}


\entrylabel{access evidence}

Public primary arXiv source;\allowbreak{} later-resolution and theorem-scope verification


\entrylabel{reading status}

unread


\entrylabel{statement validation}

not started


\entrylabel{sha256}

01d395b7\allowbreak{}9764061d\allowbreak{}57e5d34c\allowbreak{}425e4194\allowbreak{}03fde0d3\allowbreak{}da27a6c3\allowbreak{}2fcae4d5\allowbreak{}18d523a0


\entrylabel{acquired at}

2026-09-06T20:\allowbreak{}04:\allowbreak{}32Z


\subsection*{Budzinski-Senizergues:\allowbreak{} Maximum agreement subtrees and Holder homeomorphisms between Brownian trees}

random\_\allowbreak{}tree\_\allowbreak{}agreement\_\allowbreak{}2024


\url{https://arxiv.org/pdf/2304.00905}


\entrylabel{access evidence}

Public primary source;\allowbreak{} precise open questions and later resolutions checked separately


\entrylabel{reading status}

unread


\entrylabel{statement validation}

not started


\entrylabel{sha256}

83a20c5f\allowbreak{}80d40d63\allowbreak{}e1d39427\allowbreak{}a418d69d\allowbreak{}3b7f776c\allowbreak{}c461fa30\allowbreak{}ad7ec967\allowbreak{}e0e8eac2


\entrylabel{acquired at}

2026-09-06T22:\allowbreak{}16:\allowbreak{}26Z


\subsection*{Randomstrasse101:\allowbreak{} Open Problems of 2024}

randomstrasse\_\allowbreak{}2024\_\allowbreak{}v1


\url{https://arxiv.org/pdf/2504.20539v1}


\entrylabel{access evidence}

Public primary author website or arXiv manuscript


\entrylabel{reading status}

unread


\entrylabel{statement validation}

not started


\entrylabel{sha256}

5299b7af\allowbreak{}f49b581c\allowbreak{}963b2221\allowbreak{}fdc16dac\allowbreak{}f3047408\allowbreak{}71ed6318\allowbreak{}8ae40e8f\allowbreak{}0c43aac2


\entrylabel{acquired at}

2026-09-06T16:\allowbreak{}42:\allowbreak{}33Z


\subsection*{Randomstrasse101:\allowbreak{} Open Problems of 2025}

randomstrasse\_\allowbreak{}2025\_\allowbreak{}v1


\url{https://arxiv.org/pdf/2603.29571}


\entrylabel{access evidence}

Author-deposited arXiv preprint under CC BY 4.0


\entrylabel{reading status}

unread


\entrylabel{statement validation}

not started


\entrylabel{sha256}

cd4e929a\allowbreak{}83143f91\allowbreak{}57283d2d\allowbreak{}9de7e6a9\allowbreak{}9bca40f5\allowbreak{}7e16d62a\allowbreak{}9b1f3ba3\allowbreak{}3070801c


\entrylabel{acquired at}

2026-09-06T15:\allowbreak{}32:\allowbreak{}10Z


\subsection*{Aletti et al.:\allowbreak{} Open problems in stochastic reinforcement models;\allowbreak{} March2023 workshop,\allowbreak{} published2025}

reinforcement\_\allowbreak{}matrix\_\allowbreak{}2023


\url{https://www.matrix-inst.org.au/wp_Matrix2016/wp-content/uploads/2023/Aletti.pdf}


\entrylabel{access evidence}

Public primary author preprint or workshop institutional host;\allowbreak{} exact statements and current status checked separately


\entrylabel{reading status}

unread


\entrylabel{statement validation}

not started


\entrylabel{sha256}

93277e98\allowbreak{}5f6048ba\allowbreak{}f322b243\allowbreak{}858f882a\allowbreak{}c56cbad0\allowbreak{}be0a259f\allowbreak{}a2f406fe\allowbreak{}dd3fb1db


\entrylabel{acquired at}

2026-09-06T18:\allowbreak{}00:\allowbreak{}32Z


\subsection*{Van der Hofstad:\allowbreak{} Stochastic processes on random graphs,\allowbreak{} Saint-Flour2017}

rhofstad\_\allowbreak{}saintflour\_\allowbreak{}2017


\url{https://www.win.tue.nl/~rhofstad/SaintFlour_SPoRG_fin_Saint_Flour.pdf}


\entrylabel{access evidence}

Public primary author or workshop PDF;\allowbreak{} individual open-question statements and later status require separate validation


\entrylabel{reading status}

unread


\entrylabel{statement validation}

not started


\entrylabel{sha256}

77d935eb\allowbreak{}e206ee81\allowbreak{}1d131b37\allowbreak{}64b8c51f\allowbreak{}aabb8188\allowbreak{}1423b21e\allowbreak{}638d151d\allowbreak{}f4e924bd


\entrylabel{acquired at}

2026-09-06T19:\allowbreak{}31:\allowbreak{}10Z


\subsection*{Fang-Jiang-Lin-Xu:\allowbreak{} Generalized universal irrational rotation algebras and u+\allowbreak{}v}

rotation\_\allowbreak{}brown\_\allowbreak{}2012


\url{https://arxiv.org/pdf/1210.4771}


\entrylabel{access evidence}

Primary author or publisher public manuscript;\allowbreak{} no license or current-status inference from access


\entrylabel{reading status}

unread


\entrylabel{statement validation}

not started


\entrylabel{sha256}

5430fde0\allowbreak{}dbc9b46b\allowbreak{}fcfbc618\allowbreak{}7cfa94be\allowbreak{}bf03b34b\allowbreak{}fbd85169\allowbreak{}745c9119\allowbreak{}a0a0178b


\entrylabel{acquired at}

2026-09-06T21:\allowbreak{}27:\allowbreak{}36Z


\subsection*{Salez:\allowbreak{} Modern aspects of Markov chains,\allowbreak{} entropy curvature and cutoff,\allowbreak{}2025}

salez\_\allowbreak{}markov\_\allowbreak{}2025


\url{https://arxiv.org/pdf/2508.21055}


\entrylabel{access evidence}

Public primary author source,\allowbreak{} accessed2026-09-06;\allowbreak{} mathematical content and current status checked separately


\entrylabel{reading status}

questions 2 1 through 2 9 context read PDF25 32 and other previously used passages


\entrylabel{statement validation}

admitted statements paraphrased against numbered questions


\entrylabel{sha256}

9b58f5e6\allowbreak{}c2b13103\allowbreak{}f9ca02dd\allowbreak{}d53b7890\allowbreak{}e259b2d9\allowbreak{}3f632c25\allowbreak{}d3414dd8\allowbreak{}7002e267


\entrylabel{acquired at}

2026-09-06T22:\allowbreak{}03:\allowbreak{}23Z


\subsection*{Schramm:\allowbreak{} Conformally invariant scaling limits,\allowbreak{} overview and collection of problems,\allowbreak{} 2006v2}

schramm\_\allowbreak{}2006\_\allowbreak{}v2


\url{https://arxiv.org/pdf/math/0602151}


\entrylabel{access evidence}

Public primary author preprint or workshop institutional host;\allowbreak{} exact statements and current status checked separately


\entrylabel{reading status}

unread


\entrylabel{statement validation}

not started


\entrylabel{sha256}

aed7e924\allowbreak{}69d5b0f4\allowbreak{}72c74e9c\allowbreak{}bdd5eb6f\allowbreak{}109d092a\allowbreak{}ddb7fcc7\allowbreak{}de432778\allowbreak{}0a116173


\entrylabel{acquired at}

2026-09-06T18:\allowbreak{}33:\allowbreak{}45Z


\subsection*{Serfaty:\allowbreak{} Lectures on Coulomb and Riesz gases,\allowbreak{}2024}

serfaty\_\allowbreak{}coulomb\_\allowbreak{}lectures\_\allowbreak{}2024


\url{https://arxiv.org/pdf/2407.21194}


\entrylabel{access evidence}

Primary author or arXiv source;\allowbreak{} statement reading tracked separately


\entrylabel{reading status}

PDF21 163 164 199 through203 223 224 read full other pages only bounded leads


\entrylabel{statement validation}

explicit conjectures matched and schematic normalization inconsistency preserved


\entrylabel{sha256}

492af6eb\allowbreak{}7f295cce\allowbreak{}48681246\allowbreak{}b385777c\allowbreak{}515531fa\allowbreak{}92149f1b\allowbreak{}b41f3863\allowbreak{}baa5b7da


\entrylabel{acquired at}

2026-09-07T01:\allowbreak{}48:\allowbreak{}41Z


\subsection*{Du,\allowbreak{}Gao,\allowbreak{}Li,\allowbreak{}Zhuang:\allowbreak{} Sharp asymptotics for arm probabilities in critical planar percolation}

sharp\_\allowbreak{}arms\_\allowbreak{}2022


\url{https://arxiv.org/pdf/2205.15901}


\entrylabel{access evidence}

Public primary author or workshop PDF;\allowbreak{} exact statements and later status separately checked


\entrylabel{reading status}

unread


\entrylabel{statement validation}

not started


\entrylabel{sha256}

99756d8e\allowbreak{}5958a049\allowbreak{}67172b68\allowbreak{}868b70eb\allowbreak{}36b4d3c4\allowbreak{}ff90a4d7\allowbreak{}f6e351b0\allowbreak{}448e866d


\entrylabel{acquired at}

2026-09-06T19:\allowbreak{}42:\allowbreak{}22Z


\subsection*{Heilman:\allowbreak{} Sharp Hardness for MAX-3-CUT and Quantum MAX-CUT}

simplex\_\allowbreak{}partial\_\allowbreak{}2026


\url{https://arxiv.org/pdf/2608.00333}


\entrylabel{access evidence}

Public primary mathematical source;\allowbreak{} exact status and scope verification


\entrylabel{reading status}

unread


\entrylabel{statement validation}

not started


\entrylabel{sha256}

225a538b\allowbreak{}41d36176\allowbreak{}f70f166d\allowbreak{}f6f412e6\allowbreak{}459826f0\allowbreak{}8048393b\allowbreak{}a317cc6d\allowbreak{}8c1ee25b


\entrylabel{acquired at}

2026-09-06T20:\allowbreak{}45:\allowbreak{}47Z


\subsection*{Aldous:\allowbreak{} Scale-invariant random spatial networks}

sirsn\_\allowbreak{}2014


\url{https://arxiv.org/pdf/1204.0817}


\entrylabel{access evidence}

Public primary source;\allowbreak{} precise open questions and later resolutions checked separately


\entrylabel{reading status}

unread


\entrylabel{statement validation}

not started


\entrylabel{sha256}

b5e0a0e4\allowbreak{}df824477\allowbreak{}2e0b4261\allowbreak{}2799d851\allowbreak{}beee0b34\allowbreak{}ef637570\allowbreak{}eb4e620e\allowbreak{}9e2a567c


\entrylabel{acquired at}

2026-09-06T22:\allowbreak{}15:\allowbreak{}59Z


\subsection*{Rezaei:\allowbreak{} Hausdorff measure of SLE curves}

sle\_\allowbreak{}hausdorff\_\allowbreak{}2012


\url{https://arxiv.org/pdf/1212.5847}


\entrylabel{access evidence}

Public primary author or workshop PDF;\allowbreak{} exact statements and later status separately checked


\entrylabel{reading status}

unread


\entrylabel{statement validation}

not started


\entrylabel{sha256}

7fc16233\allowbreak{}4e1466ad\allowbreak{}c50a8d1f\allowbreak{}e240dd31\allowbreak{}5b2d50b4\allowbreak{}1129088f\allowbreak{}90071e59\allowbreak{}3d4a9687


\entrylabel{acquired at}

2026-09-06T19:\allowbreak{}44:\allowbreak{}16Z


\subsection*{The strong convergence phenomenon,\allowbreak{}2025 survey}

strong\_\allowbreak{}convergence\_\allowbreak{}2025


\url{https://arxiv.org/pdf/2507.00346}


\entrylabel{access evidence}

Public primary mathematical source;\allowbreak{} exact status and scope verification


\entrylabel{reading status}

unread


\entrylabel{statement validation}

not started


\entrylabel{sha256}

78721e89\allowbreak{}ffb1817d\allowbreak{}0887a06a\allowbreak{}58d2ed1c\allowbreak{}516784f4\allowbreak{}90710bd6\allowbreak{}030307a3\allowbreak{}666cf0ad


\entrylabel{acquired at}

2026-09-06T20:\allowbreak{}51:\allowbreak{}12Z


\subsection*{Gaitonde-Mossel:\allowbreak{} Comparison theorems for the mixing times of systematic and random scan dynamics,\allowbreak{}v2July2025}

systematic\_\allowbreak{}scan\_\allowbreak{}2025


\url{https://arxiv.org/pdf/2410.11136}


\entrylabel{access evidence}

Public primary author source,\allowbreak{} accessed2026-09-06;\allowbreak{} mathematical content and current status checked separately


\entrylabel{reading status}

unread


\entrylabel{statement validation}

not started


\entrylabel{sha256}

77d17d6a\allowbreak{}114d9b68\allowbreak{}d7a4b23a\allowbreak{}0680dd93\allowbreak{}b057265b\allowbreak{}6821f4de\allowbreak{}791f697b\allowbreak{}164eaf19


\entrylabel{acquired at}

2026-09-06T22:\allowbreak{}04:\allowbreak{}07Z


\subsection*{Thoma:\allowbreak{} Non-rigidity Properties of the Coulomb Gas}

thoma\_\allowbreak{}coulomb\_\allowbreak{}nonrigidity\_\allowbreak{}2023


\url{https://arxiv.org/pdf/2303.11486}


\entrylabel{access evidence}

Primary arXiv preprint found in dated later-literature check


\entrylabel{reading status}

PDF1 through3 read full full proof not read


\entrylabel{statement validation}

questions or later status matched with explicit scope


\entrylabel{sha256}

c13025c7\allowbreak{}19718118\allowbreak{}5af8c75e\allowbreak{}5f720522\allowbreak{}e45f517b\allowbreak{}f07fb90e\allowbreak{}ef42ade8\allowbreak{}9dcd0da7


\entrylabel{acquired at}

2026-09-07T01:\allowbreak{}35:\allowbreak{}29Z


\subsection*{Towsner:\allowbreak{} Limits of Sequences of Markov Chains}

towsner\_\allowbreak{}markov\_\allowbreak{}limits\_\allowbreak{}2015


\url{https://arxiv.org/pdf/1404.3815}


\entrylabel{access evidence}

Primary URL discovered through focused literature search2026-09-06;\allowbreak{} substantive title and theorem verification follows acquisition.


\entrylabel{reading status}

unread


\entrylabel{statement validation}

not started


\entrylabel{sha256}

7e8d8419\allowbreak{}4659245c\allowbreak{}456b216b\allowbreak{}ae2b64ef\allowbreak{}fede229f\allowbreak{}2c488f8a\allowbreak{}a09661e9\allowbreak{}799000d8


\entrylabel{acquired at}

2026-09-06T22:\allowbreak{}38:\allowbreak{}10Z


\subsection*{Random tree recursions:\allowbreak{} which fixed points correspond to tangible sets of trees?}

tree\_\allowbreak{}automata\_\allowbreak{}2019


\url{https://arxiv.org/pdf/1808.03019v2}


\entrylabel{access evidence}

Public primary author website or arXiv manuscript


\entrylabel{reading status}

unread


\entrylabel{statement validation}

not started


\entrylabel{sha256}

65b10d5d\allowbreak{}474c8ef7\allowbreak{}4b3aceb0\allowbreak{}af983e32\allowbreak{}401a62e7\allowbreak{}582086dc\allowbreak{}5f85a599\allowbreak{}e60bc7f1


\entrylabel{acquired at}

2026-09-06T16:\allowbreak{}42:\allowbreak{}18Z


\subsection*{Alev-Frishberg-Sarantis-Tetali:\allowbreak{} Faster triangulation mixing via transport flows,\allowbreak{}ICALP2026}

triangulation\_\allowbreak{}transport\_\allowbreak{}2026


\url{https://drops.dagstuhl.de/storage/00lipics/lipics-vol374-icalp2026/LIPIcs.ICALP.2026.12/LIPIcs.ICALP.2026.12.pdf}


\entrylabel{access evidence}

Primary source statement and later-literature scope verification2026-09-06


\entrylabel{reading status}

unread


\entrylabel{statement validation}

title and relevant theorem read


\entrylabel{sha256}

50f80931\allowbreak{}82b1ae9f\allowbreak{}fe705836\allowbreak{}a8700314\allowbreak{}75c14d00\allowbreak{}7083177d\allowbreak{}2bad4c53\allowbreak{}3af62ec4


\entrylabel{acquired at}

2026-09-06T22:\allowbreak{}48:\allowbreak{}39Z


\subsection*{Burdzy and Pal:\allowbreak{} Twin peaks}

twin\_\allowbreak{}peaks\_\allowbreak{}2016


\url{https://arxiv.org/pdf/1606.08025}


\entrylabel{access evidence}

Public primary author preprint or workshop institutional host;\allowbreak{} exact statements and current status checked separately


\entrylabel{reading status}

unread


\entrylabel{statement validation}

not started


\entrylabel{sha256}

0c184034\allowbreak{}31de50d5\allowbreak{}3950c879\allowbreak{}7f008fd7\allowbreak{}6b5e74f1\allowbreak{}fb2c6135\allowbreak{}55fc66e3\allowbreak{}38fb606d


\entrylabel{acquired at}

2026-09-06T18:\allowbreak{}35:\allowbreak{}36Z


\clearpage

\section*{Proof appendix:\allowbreak{} ARW fixed mass window obstructions}

The current ledger disposition and the exact status limits in this packet govern its interpretation. Classical supporting facts and known-result rederivations are distinct from original open-problem solutions.


\subsection*{Fixed mass and the ARW window-variance questions}

Status:\allowbreak{} author-checked scope counterexample only. No independent review or novelty claim. The intended mesoscopic hyperuniformity problems remain unsolved.

Source:\allowbreak{} Levine and Silvestri,\allowbreak{} \emph{Universality Conjectures for Activated Random Walk},\allowbreak{} arXiv:\allowbreak{}2306.01698v2,\allowbreak{} Conjecture 27 and Question 28,\allowbreak{} PDF pages 15-16. Corpus entries P0389 and P0390. The two questions concern different laws,\allowbreak{} but share the following necessary correction to the observation-window quantifier.

\subsection*{Deterministic complement identity}

Let T be a finite set with N sites,\allowbreak{} and let I\_\allowbreak{}x be zero-one occupation indicators with exactly k occupied sites almost surely. For every B contained in T,\allowbreak{}

\[S_B=\sum_{x\in B}I_x=k-S_{T\setminus B},\qquad \operatorname{Var}(S_B)=\operatorname{Var}(S_{T\setminus B}).\]

Both counts are bounded by their respective window sizes. The variance bound for a variable in an interval of length a is a squared divided by four. Thus

\[\operatorname{Var}(S_B)\le\frac14\min\{|B|,N-|B|\}^2.\]

This proof uses only fixed total mass,\allowbreak{} which holds after every free-chain addition and throughout each fixed-particle-number wake chain.

\subsection*{Proper boxes already contradict the unrestricted lower bound}

Take d=\allowbreak{}1 and T the cycle with L vertices. Fix any positive subcritical density rho,\allowbreak{} and put k=\allowbreak{}floor(rho L). Take the proper interval B consisting of L-1 vertices. If x is the omitted vertex,\allowbreak{} S\_\allowbreak{}B=\allowbreak{}k-I\_\allowbreak{}x,\allowbreak{} so its variance is at most 1/\allowbreak{}4. It cannot be at least c(L-1) for any positive c independent of L.

Uniform free driving is translation equivariant. For the wake chain,\allowbreak{} the unique stationary law is translation invariant because the transition kernel commutes with translations. Therefore each single-site mean is rho\_\allowbreak{}L=\allowbreak{}k/\allowbreak{}L in both ensembles,\allowbreak{} and the variance of this proper box is exactly rho\_\allowbreak{}L(1-rho\_\allowbreak{}L). The whole torus gives the even simpler value zero,\allowbreak{} but the proper interval shows that requiring merely a strict subset does not repair the claim.

At zero density the configuration is empty and has variance zero for every window,\allowbreak{} giving another boundary exception. The positive fixed-density argument above does not rely on that exception.

\subsection*{Global cancellation is not mesoscopic hyperuniformity}

Under any fixed-mass law,\allowbreak{}

\[\sum_{y\in T}\operatorname{Cov}(I_x,I_y)=\operatorname{Cov}(I_x,k)=0.\]

Consider,\allowbreak{} only as a comparison model,\allowbreak{} a uniformly chosen k-subset of T. For x distinct from y,\allowbreak{}

\[\mathbb E I_x=\frac{k}{N},\qquad\mathbb E I_xI_y=\frac{k(k-1)}{N(N-1)}.\]

For a window with b sites,\allowbreak{} expansion of the variance gives

\[\operatorname{Var}(S_B)=b\,\rho_N(1-\rho_N)\frac{N-b}{N-1},\qquad\rho_N=k/N.\]

When b grows but b/\allowbreak{}N tends to zero and rho\_\allowbreak{}N tends to a number strictly between zero and one,\allowbreak{} this is asymptotic to b times rho(1-rho). Therefore exact finite-volume covariance cancellation is compatible with volume-order mesoscopic fluctuations. Passing a finite-torus zero sum to an infinite-volume covariance sum requires additional uniform tail control.

\subsection*{What remains}

The natural next version restricts windows away from full volume,\allowbreak{} or takes the thermodynamic limit before the window limit. Neither restriction is asserted to have been prescribed by the source. Proving the ARW transition and its exponent in such a regime remains the actual mathematical problem. These elementary identities do not show that either actual ARW law is a uniform subset law,\allowbreak{} negatively associated,\allowbreak{} or hyperuniform.

\subsection*{Second check}

An exact rational enumeration of all k-subsets for small N independently checks the comparison variance and complement formula. The enumeration is a finite arithmetic control,\allowbreak{} not evidence that an unbounded ARW conjecture is solved. Its receipt is in ARW\_\allowbreak{}window\_\allowbreak{}exact\_\allowbreak{}controls.json.

\clearpage

\section*{Proof appendix:\allowbreak{} Bilateral degree fluctuations partial}

The current ledger disposition and the exact status limits in this packet govern its interpretation. Classical supporting facts and known-result rederivations are distinct from original open-problem solutions.


\subsection*{Bilateral-preference degrees:\allowbreak{} a non-Gaussian fluctuation candidate}

Status:\allowbreak{} private author-checked candidate for the growing-degree regime in the degree-distribution question on PDF page 3 of Dabirian and Subramanian,\allowbreak{} \emph{Connectivity of a Family of Bilateral Preference Random Graphs},\allowbreak{} arXiv:\allowbreak{}2209.05677v3,\allowbreak{} 8 July 2025. No independent review or novelty certification. This packet does not settle every growth regime in the source's broadly stated question. The related simultaneous-concentration estimate is a sufficient condition,\allowbreak{} not an optimal characterization. Both remain partial in the campaign ledger.

The model is exactly the source model:\allowbreak{} assign independent uniform $(0,1)$ weights to all undirected edges of the complete graph on $n$ vertices. An edge survives in $LK(n,k)$ precisely when it is among the $k$ smallest incident weights at both endpoints. Let $D$ be the degree of a specified vertex;\allowbreak{} symmetry gives the same law for a uniformly chosen vertex. Ties have probability zero.

\subsection*{Candidate theorem and meaning}

Suppose $k=k(n)\to\infty$ and $k=o(n^{2/5})$. Let $Z$ be standard normal and write

\[g(z)=\phi(z)+z\Phi(z),\qquad v=\frac13+\frac{\sqrt3-2}{2\pi}>0.\]

The candidate conclusion is convergence in distribution and of the first two moments:\allowbreak{}

\[\frac{k-D}{\sqrt{k}}\ \Longrightarrow\ g(Z).\]

Consequently,\allowbreak{}

\[\frac{\operatorname{Var}(D)}{k}\longrightarrow v,\qquad
\frac{D-\mathbb ED}{\sqrt{\operatorname{Var}(D)}}\ \Longrightarrow\
\frac{1/\sqrt\pi-g(Z)}{\sqrt v}.\]

This includes the source's $k=o(n^{1/3})$ regime for its detailed mean asymptotic. The limiting distribution is not normal:\allowbreak{} $g$ is strictly increasing,\allowbreak{} positive,\allowbreak{} and has range $(0,\infty)$,\allowbreak{} so the standardized degree limit has a finite upper endpoint. Its negative skew is consistent with the deterministic ceiling $D\le k$.

\subsection*{Exact finite-graph comparison}

Fix the root $0$. Let $S$ be its $k$ preferred neighbors,\allowbreak{} let $O$ be the other $n-k-1$ vertices,\allowbreak{} and let $X$ be the root's $(k+1)$st smallest incident weight. Then

\[X\sim\operatorname{Beta}(k+1,n-k-1),\quad
\mathbb EX=\frac{k+1}{n},\quad
\mathbb EX^2=\frac{(k+1)(k+2)}{n(n+1)}.\]

Conditional on $X$ and the selected label set $S$,\allowbreak{} the $k$ unordered root-to-$S$ weights are independent uniforms on $(0,X)$. All edges not incident to the root retain their independent uniform laws. This conditional statement follows directly from the joint density of independent uniform order statistics;\allowbreak{} no independence of actual degrees is asserted.

Temporarily ignore edges whose two endpoints lie in $S$ when deciding whether a selected neighbor also prefers the root. Let $D^0$ count the resulting accepted root edges. Removing competitors can only help acceptance,\allowbreak{} so $D^0\ge D$.

Let $J$ count edges internal to $S$ with weight below $X$,\allowbreak{} and put $M=\binom{k}{2}$. If a neighbor's acceptance changes,\allowbreak{} it must touch such an edge:\allowbreak{} every root-to-$S$ weight is below $X$. Therefore

\[0\le D^0-D\le2J,\qquad J\mid X\sim\operatorname{Bin}(M,X).\]

In particular,\allowbreak{}

\[\mathbb E\frac{(D^0-D)^2}{k}
\le\frac4k\left(M\frac{k+1}{n}
+M(M-1)\frac{(k+1)(k+2)}{n(n+1)}\right)
=O\left(\frac{k^2}{n}+\frac{k^5}{n^2}\right)=o(1).\]

The last step is exactly where the sufficient restriction $k=o(n^{2/5})$ is used. This is an $L^2$ comparison;\allowbreak{} a bound only on the probability of a discrepancy would not by itself justify the variance conclusion.

There are $m=n-k-1$ outside edges at each selected neighbor. The outside-edge arrays for different selected neighbors are disjoint sets of independent edge variables. Conditional on $X$ and $S$,\allowbreak{} after also averaging over their independent root weights,\allowbreak{} the ideal acceptance indicators are thus independent and identically distributed. For large $n$,\allowbreak{} $m\ge k$. If $T$ is the $k$th order statistic of $m$ independent uniforms,\allowbreak{} independent of $X$,\allowbreak{} then

\[T\sim\operatorname{Beta}(k,n-2k),\qquad
D^0\mid X\sim\operatorname{Bin}(k,1-q_n(X)),\qquad
q_n(X)=\frac{\mathbb E[(X-T)_+\mid X]}{X}.\]

Indeed,\allowbreak{} for a root edge of weight $u$,\allowbreak{} acceptance means that at most $k-1$ outside edges have weight below $u$,\allowbreak{} equivalently $u<T$. Integrating $u$ uniformly over $(0,X)$ gives the displayed formula. These are conditional laws of the comparison variable,\allowbreak{} not of $D$.

\subsection*{Fluctuation and moment limits}

Set

\[A=nX,\quad B=nT,\quad U_n=\frac{A-k}{\sqrt k},\quad V_n=\frac{B-k}{\sqrt k}.\]

Both $U_n$ and $V_n$ converge to standard normal variables in their marginal laws,\allowbreak{} and their fourth moments are uniformly bounded. Here are the needed checks. Represent a beta variable as a ratio of independent gamma variables. For $X$,\allowbreak{} the numerator is $\operatorname{Gamma}(k+1,1)$ and the denominator has shape $n$;\allowbreak{} the denominator replacement error after scaling by $\sqrt k$ is $O_P(\sqrt{k/n})$. For $T$,\allowbreak{} the numerator has shape $k$ and the denominator has shape $n-k$;\allowbreak{} the additional centering error is $k^{3/2}/(n-k)=o(1)$. The gamma central limit theorem then applies. The moment assertion also follows directly from the beta moments $\mathbb E X^j=(k+1)^{\overline j}/n^{\overline j}$ and $\mathbb E T^j=k^{\overline j}/(n-k)^{\overline j}$ for $j\le4$. Equivalently,\allowbreak{} the beta variance and kurtosis formulas show bounded fourth moments after this centering and scaling. No exponentially growing degree moment is being assumed.

Define

\[g_n(u)=\mathbb E[(u-V_n)_+].\]

Convergence in law of $V_n$ and uniform integrability of $|V_n|$ give convergence to normal in Wasserstein distance of order one. Since the positive-part function is one-Lipschitz,\allowbreak{}

\[\sup_{u\in\mathbb R}|g_n(u)-g(u)|\longrightarrow0,\qquad
g(u)=\mathbb E[(u-Z)_+]=\phi(u)+u\Phi(u).\]

The conditional mean of the scaled degree deficit is

\[C_n:=\frac{kq_n(X)}{\sqrt k}=\frac{k}{A}g_n(U_n).\]

Since $k/A\to1$ in probability,\allowbreak{} the preceding facts imply $C_n\Longrightarrow g(Z)$.

To justify moments,\allowbreak{} split according to $A\ge k/2$. On that event,\allowbreak{}

\[C_n\le2\bigl(|U_n|+\mathbb E|V_n|\bigr).\]

On its complement,\allowbreak{} $C_n\le\sqrt k$ because $0\le q_n\le1$. Furthermore,\allowbreak{} the order-statistic/\allowbreak{}binomial identity and a Chernoff bound give

\[\mathbb P(A<k/2)
\le\mathbb P\{\operatorname{Bin}(n-1,k/(2n))\ge k\}
\le e^{-(\log2-1/2)k}.\]

For the last inequality,\allowbreak{} apply Markov's inequality to $2^{\operatorname{Bin}}$ and use that the binomial mean is at most $k/2$. Thus $\sup_n\mathbb E C_n^4<\infty$ along the asymptotic sequence. It follows that the first two moments of $C_n$ converge to those of $g(Z)$.

The remaining conditional binomial noise has vanishing second moment:\allowbreak{}

\[\mathbb E\left[\left(\frac{k-D^0}{\sqrt k}-C_n\right)^2\right]
=\mathbb E[q_n(X)(1-q_n(X))]
\le\frac{\mathbb EC_n}{\sqrt k}\longrightarrow0.\]

Combining this with the earlier $L^2$ comparison from $D^0$ to $D$ proves the candidate distribution and moment conclusions.

\subsection*{Two independent calculations of the constants}

Let $Z,Z_1,Z_2$ be independent standard normal variables. Conditioning on $Z$ gives

\[\mathbb E g(Z)=\mathbb E(Z-Z_1)_+=\frac1{\sqrt\pi},\]

and

\[\mathbb E g(Z)^2=\mathbb E[(Z-Z_1)_+(Z-Z_2)_+]
=\frac13+\frac{\sqrt3}{2\pi}.\]

For the last equality,\allowbreak{} the two differences have variances two and correlation $1/2$. Write them as $\sqrt2 R\cos\theta$ and $\sqrt2 R\cos(\theta-\pi/3)$ in polar coordinates of two independent standard normals. The angle is uniform,\allowbreak{} $\mathbb ER^2=2$,\allowbreak{} and the positive overlap is $(-\pi/6,\pi/2)$. Integrating the product of cosines on this interval gives $1/3+\sqrt3/(2\pi)$. Subtracting $1/\pi$ gives $v$ above,\allowbreak{} approximately $0.29068789486043867$.

A separate consistency calculation starts with the fixed-$k$ Erlang weighted tree root law in Moharrami et al. Its root threshold has law $\operatorname{Gamma}(k+1,1)$,\allowbreak{} and its conditional degree is binomial with retention probability $\mathbb E\min(\operatorname{Gamma}(k,1),X)/X$. Applying the gamma central limit theorem as $k\to\infty$ gives the same $g(Z)$ limit. This second lane does not replace the exact finite-$n$ comparison above:\allowbreak{} a sequential local limit alone would not prove the triangular-array theorem.

\subsection*{A separate sufficient condition for simultaneous concentration}

The same source asks for conditions under which degrees concentrate around their mean. A simple sufficient regime is $k/\log n\to\infty$. This uses the source's existing Erdős-Rényi sandwich idea;\allowbreak{} no novelty is claimed.

Let $\delta=\sqrt{6k\log n}$ and suppose $k>2\delta$. Threshold every uniform edge weight at $t=(k-\delta)/(n-1)$. Write $B_v$ for the resulting degree at vertex $v$,\allowbreak{} so $B_v\sim\operatorname{Bin}(n-1,t)$ and $\mathbb EB_v=k-\delta$. Chernoff bounds give,\allowbreak{} for each vertex,\allowbreak{}

\[\mathbb P(B_v>k)\le n^{-3},\qquad
\mathbb P(B_v<k-2\delta)\le n^{-3}.\]

The upper-tail exponent is at least $\delta^2/(2(k-\delta)+2\delta/3)\ge3\log n$;\allowbreak{} the lower-tail exponent is at least $\delta^2/(2(k-\delta))\ge3\log n$. A union bound,\allowbreak{} requiring no independence between vertex degrees,\allowbreak{} now shows that all threshold degrees lie in $[k-2\delta,k]$ with probability at least $1-2/n^2$.

On this event every threshold edge is among the $k$ preferred edges at both its endpoints. Hence the threshold graph is a subgraph of $LK(n,k)$ and

\[k-2\sqrt{6k\log n}\le D_v\le k\quad\hbox{for every vertex }v.\]

This proves simultaneous relative concentration when $k/\log n\to\infty$. It does not determine optimal tails or the critical logarithmic regime;\allowbreak{} below the connectivity threshold,\allowbreak{} isolated vertices already obstruct simultaneous concentration near $k$.

\subsection*{Checks and remaining limits}

An exhaustive independent rank enumeration of all $6!=720$ edge orders on four vertices,\allowbreak{} for $k=1,2$,\allowbreak{} found no violation of $D\le D^0\le D+2J$. The degree histograms were $(288,432,0,0)$ and $(48,192,480,0)$,\allowbreak{} giving means $3/5$ and $8/5$. These means also follow from direct two-endpoint ranking probabilities:\allowbreak{} a root's nearest edge is mutual with probability $3/5$,\allowbreak{} while for $k=2$ a given edge is kept with probability $1-2/3+1/5=8/15$. These finite controls are not a proof of the asymptotic theorem.

The fragile steps are the root's $k+1$ order statistic,\allowbreak{} the outside beta parameter $n-2k$,\allowbreak{} the $L^2$ rather than mere probability comparison,\allowbreak{} and moment convergence before standardization. They are exposed explicitly above. Bounded-degree regimes,\allowbreak{} larger growth regimes,\allowbreak{} optimal concentration,\allowbreak{} independent review,\allowbreak{} and novelty remain unresolved here. No original solution is promoted. The exact enumeration receipt and a retained proof plan are in \path{bilateral_degree_exact_checks.json} and \path{bilateral_degree_proof_plan.json}.

\clearpage

\section*{Proof appendix:\allowbreak{} IPVT boundary hole partial}

The current ledger disposition and the exact status limits in this packet govern its interpretation. Classical supporting facts and known-result rederivations are distinct from original open-problem solutions.


\subsection*{Boundary-hole monotonicity and the odd-dimensional polynomial}

Status:\allowbreak{} private author-checked partial candidate. No independent review,\allowbreak{} no novelty certification,\allowbreak{} and no claim that the whole source conjecture is solved.

Source:\allowbreak{} D'Achille,\allowbreak{} Curien,\allowbreak{} Enriquez,\allowbreak{} Lyons and Unel,\allowbreak{} \emph{Ideal Poisson--Voronoi tessellations on hyperbolic spaces},\allowbreak{} archived as \path{ideal_voronoi_2023},\allowbreak{} Proposition 5.1 and equation (5.1),\allowbreak{} PDF35;\allowbreak{} Conjecture 7.5,\allowbreak{} PDF48. The source already computes the dimension-three formula and observes rationality in odd dimensions. Those facts are credited. The claims attempted here are the all-dimensional monotonicity and positive coefficients in every odd dimension. The final addendum gives a candidate counterexample to the even-dimensional positive-power-series clause in its displayed distance variable.

\subsection*{Exact starting identity}

Let $n=d-1\ge1$ and let $H_d(u)$ be the probability that the hyperbolic ball of radius $u\ge0$ about the origin lies inside its ideal Voronoi cell. Write $J_d(u)=1/H_d(u)$. Equation (5.1) and Proposition 5.1 give

\[J_d(u)=\frac1{c_d}\int_{\mathbb R^n}
\left(\sqrt{\cosh^2u+|x|^2}-\sinh u\right)^{-2n}\,dx,
\qquad
c_d=\frac{\pi^{n/2}\Gamma(n/2)}{\Gamma(n)}.\]

The normalization follows independently by setting $u=0$ and integrating $(1+|x|^2)^{-n}$ in polar coordinates. In particular $J_d(0)=1$.

\subsection*{Monotonicity in every dimension}

Set $q=e^{2u}$ and substitute $x=e^{-u}z$. Then

\[e^{-nu}J_d(u)=\frac1{c_d}\int_{\mathbb R^n}D_q(z)^{-2n}\,dz,
\qquad
D_q(z)=\sqrt{(q+1)^2/4+|z|^2}-(q-1)/2.\]

The denominator is positive. For $z\ne0$,\allowbreak{}

\[\frac{\partial D_q(z)}{\partial q}
=\frac{q+1}{4\sqrt{(q+1)^2/4+|z|^2}}-\frac12<0.\]

Thus the integrand strictly increases with $q$ away from a null set. All integrals are finite by the exact source identity. Pointwise comparison of two values of $q$ proves that $e^{-nu}J_d(u)$ is strictly increasing;\allowbreak{} no interchange of differentiation and integration is needed. Taking reciprocals proves

\[u\longmapsto e^{(d-1)u}H_d(u)
\quad\hbox{is strictly decreasing on }[0,\infty).\]

The strictness follows because the difference of the positive finite integrands is strictly positive on a set of positive Lebesgue measure.

\subsection*{Positive rational polynomial for odd dimensions}

Let $d=2m+1$,\allowbreak{} so $n=2m$,\allowbreak{} with integer $m\ge1$. Put $a=e^{-u}$ and $s=\sinh u=(a^{-1}-a)/2$. In the radial integral make the substitution

\[y=\sqrt{\cosh^2u+r^2}-s.\]

Its lower endpoint is $a$,\allowbreak{} and

\[r^2=(y-a)(y+a^{-1}),\qquad r\,dr=(y+s)\,dy.\]

After setting $y=az$,\allowbreak{} the exact integral becomes

\[J_{2m+1}(u)=C_m a^{-4m}
\int_1^\infty (z-1)^{m-1}(a^2z+1)^{m-1}
\bigl[a^2(2z-1)+1\bigr]z^{-4m}\,dz,
\qquad C_m=\frac{\Gamma(2m)}{\Gamma(m)^2}.\]

Here the factor $C_m$ is half the ratio of the $(2m-1)$-sphere area to $c_d$. Expand the two bracketed factors as a polynomial in $a^2$. Every coefficient is a strictly positive integrable function of $z>1$:\allowbreak{} the only apparent subtraction is $2z-1$,\allowbreak{} which is positive on this domain. The degree is exactly $m$. The resulting integrals are rational because expanding $(z-1)^{m-1}$ reduces them to finite sums of powers of $z$ with integer exponents at most $-2$.

For an explicit formula,\allowbreak{} let $B$ denote the beta function,\allowbreak{} and interpret binomial coefficients outside $0,\ldots,m-1$ as zero. Define,\allowbreak{} for $0\le j\le m$,\allowbreak{}

\[b_{m,j}=\binom{m-1}{j}B(3m-j,m)
+\binom{m-1}{j-1}\left[2B(3m-j,m)-B(3m-j+1,m)\right].\]

This follows from $\int_1^\infty(z-1)^{m-1}z^{j-4m}\,dz=B(3m-j,m)$. All beta arguments are positive integers. Moreover $B(A+1,m)<B(A,m)$,\allowbreak{} so each displayed coefficient is positive. Hence

\[p_m(q)=C_m\sum_{j=0}^m b_{m,j}q^{m-j}\]

is a degree-$m$ polynomial with strictly positive rational coefficients,\allowbreak{} and

\[H_{2m+1}(u)=\frac{e^{-2mu}}{p_m(e^{2u})}.\]

The normalization gives $p_m(1)=1$. The first two cases are

\[p_1(q)=\frac{q+2}{3},\qquad
p_2(q)=\frac{q^2}{7}+\frac{16q}{35}+\frac25.\]

The first reproduces the already-known dimension-three formula in Remark 5.2. It is a control,\allowbreak{} not a newly solved special case.

\subsection*{An additional odd-dimensional consequence}

The polynomial representation makes $J_d(u)$ a positive sum of exponentials with rates in $[n,2n]$. Let $L=\log J_d$. Under the exponential weights,\allowbreak{} $L'$ is the mean rate and $L''$ its variance. Therefore $L'\ge n$ and $L''\le n^2/4$. The density $f=-H_d'=H_dL'$ satisfies

\[f'=H_d\bigl(L''-(L')^2\bigr)<0.\]

This proves the source's proposed decreasing-density property in odd dimensions. It also implies that the median is smaller than the mean. For completeness,\allowbreak{} if $M>0$ is the median and $v=f(M)>0$,\allowbreak{} convexity of the survival function gives

\[H_d(x)\ge\max\{1/2-v(x-M),0\},\qquad x\ge0.\]

Integrating yields $E X\ge(1/2+vM)^2/(2v)\ge M$,\allowbreak{} since the difference in the last inequality is $(vM-1/2)^2/(2v)$. Strict convexity makes the first inequality strict on a positive-length interval. The mean is finite by the source exponential tail bound.

\subsection*{Adversarial checks and remaining gap}

The load-bearing points to try to break are the factor from the radial Jacobian,\allowbreak{} the power $a^{-4m}$,\allowbreak{} and the distinction between polynomial positivity in $q=e^{2u}$ and power-series positivity in $u$. The check script compares the beta coefficients with an independent expansion and exact integration of the radial polynomial. It also compares the resulting formula with the original angular integral in Proposition 5.1,\allowbreak{} including even dimensions for the separate monotonicity calculation.

The even-dimensional positive-power-series assertion is not implied by monotonicity. The addendum below supplies a candidate counterexample to that clause in dimension two. No independent referee has checked this packet,\allowbreak{} and the bounded literature search is not a novelty certificate. Other even-dimensional density and median claims remain unproved here,\allowbreak{} so this grouped entry retains its partial-candidate classification.

\subsection*{Addendum:\allowbreak{} a negative sixth coefficient in dimension two}

Author-checked candidate counterexample,\allowbreak{} 2026-09-07. This uses the source's displayed distance variable $u$ and interprets "power series" as a convergent series in nonnegative integer powers of $u$ about zero. It makes no claim about series in a different variable.

The source's Remark 5.2,\allowbreak{} equation (5.2),\allowbreak{} credits its dimension-two formula to earlier work. We do not claim that formula as new:\allowbreak{}

\[H_2(u)=\frac{\pi}{4\arctan(e^u)\cosh^2u+2\sinh u}.\]

Conjecture 7.5 would require $H_2(u)=e^{-u}/f(u)$ with positive power-series coefficients. Thus $f$ is uniquely determined near zero by

\[f(u)=\frac{(e^u+2e^{-u}+e^{-3u})A(u)+1-e^{-2u}}{\pi},\qquad A(u)=\arctan(e^u).\]

Since $A(0)=\pi/4$ and $A'(u)=1/(2\cosh u)$,\allowbreak{}

\[A(u)=\frac\pi4+\frac u2-\frac{u^3}{12}+\frac{u^5}{48}+O(u^7).\]

Put $E(u)=e^u+2e^{-u}+e^{-3u}$. The coefficients needed for degree six are $[u^6]E=61/60$,\allowbreak{} $[u^5]E=-61/30$,\allowbreak{} $[u^3]E=-14/3$,\allowbreak{} and $[u]E=-4$. Consequently,\allowbreak{}

\[[u^6]f(u)=\frac{61}{240}+\frac1\pi\left(-\frac{61}{60}+\frac7{18}-\frac1{12}-\frac4{45}\right)=\frac{61\pi-192}{240\pi}<0.\]

The sign is exact:\allowbreak{} $\pi<22/7$ gives $61\pi-192<-2/7$. Numerically the coefficient is approximately $-0.000481242280365871$. A function analytic at zero has unique power-series coefficients,\allowbreak{} so no alternative positive-coefficient series in the same variable can represent this $f$ near zero. This contradicts the stated even-dimensional coefficient clause already at $d=2$.

\subsection*{Checks and strongest alternative explanation}

The coefficient was obtained by exact rational convolution using the recurrence for the reciprocal of $\cosh u$,\allowbreak{} by high-precision differentiation of the closed formula,\allowbreak{} and by Cauchy coefficient extraction directly from the original angular integral of Proposition 5.1. The last lane avoids the closed arctangent formula. Reproduction:\allowbreak{} \path{check_IPVT_even_coefficients.py};\allowbreak{} receipt:\allowbreak{} \path{IPVT_even_coefficient_checks.json}.

The strongest interpretation risk is that the authors intended a series in another transformed variable or a different normalization. The archived conjecture writes $f(u)$ in the distance variable and specifies no alternative. The calculation refutes that literal Maclaurin-series clause;\allowbreak{} it does not refute monotonicity,\allowbreak{} the odd-dimensional formula,\allowbreak{} or every possible reformulation. A source clarification or corrected conjecture could change the disposition. The dimension-two hole formula itself is already known;\allowbreak{} novelty of this coefficient observation has not been established.

\clearpage

\section*{Proof appendix:\allowbreak{} Limited forkjoin stationary exponential partial}

The current ledger disposition and the exact status limits in this packet govern its interpretation. Classical supporting facts and known-result rederivations are distinct from original open-problem solutions.


\subsection*{Limited fork-join:\allowbreak{} a stationary exponential-service partial candidate}

Status:\allowbreak{} private author-checked partial result for P0429. No independent review or novelty certification. The general finite-second-moment service question is not solved. The basic collision and empty-start coupling method is already in Wang et al.,\allowbreak{} \emph{Delay Asymptotics and Bounds for Multi-Task Parallel Jobs},\allowbreak{} Sections 3-5,\allowbreak{} archived as \path{queues_limited_forkjoin_2018}. We credit that method and its $k=o(n^{1/4})$ theorem. This packet sharpens the marginal emptying-time step under an additional exponential-service assumption and restores thinned arrival rates explicitly.

\subsection*{Model and candidate bound}

There are $n\ge2$ FIFO servers. Jobs arrive as a Poisson process of rate $\Lambda=n\lambda/k$,\allowbreak{} independently select uniform $k$-subsets of servers,\allowbreak{} and place one task in each selected queue. Service requirements of all tasks,\allowbreak{} including different tasks of the same job,\allowbreak{} are independent and identically distributed,\allowbreak{} independent of routing and arrivals. Fix the marginal rate $0<\lambda<\mu$,\allowbreak{} with mean service $1/\mu$. Let $\pi_{n,k}$ be the stationary joint workload law at any specified $k$ servers,\allowbreak{} and let $\nu_k$ be the product of their marginal stationary workload laws. Total variation uses $\sup_A|\pi(A)-\nu(A)|$.

Write $r(T)$ for the probability that a stationary marginal M/\allowbreak{}G/\allowbreak{}1 workload has not hit zero by time $T$. The modular coupling bound is

\[d_{\rm TV}(\pi_{n,k},\nu_k)\le 2k r(T)+\frac{3\lambda T k(k-1)^2}{2(n-1)},\qquad T>0.\]

With exponential service of rate $\mu$,\allowbreak{} put $\rho=\lambda/\mu$ and $c=(\sqrt\mu-\sqrt\lambda)^2>0$. Then

\[r(T)\le(1+\sqrt\rho)e^{-cT},\]

so the candidate explicit stationary bound is

\[d_{\rm TV}(\pi_{n,k},\nu_k)\le 2k(1+\sqrt\rho)e^{-cT}+\frac{3\lambda T k(k-1)^2}{2(n-1)}.\]

Either right-hand side may also be capped at one. The constants are uniform in $n,k$ at fixed $\lambda,\mu$.

\subsection*{Step 1:\allowbreak{} independent arrivals with exactly restored marginal rates}

Fix the tagged set $K$ with $|K|=k$. For one job let $N$ count selected servers in $K$. Counting ordered pairs gives

\[\mathbb E[N(N-1)]=\frac{k^2(k-1)^2}{n(n-1)}.\]

Poisson marking decomposes the arrival stream into independent streams indexed by its selected tagged subset. Retain every singleton-tagged arrival,\allowbreak{} discard all arrivals hitting two or more tagged queues,\allowbreak{} and add independent Poisson arrivals to each tagged queue at the marginal rate removed from that queue. The resulting $k$ arrival processes are independent Poisson processes of rate exactly $\lambda$. Give added tasks independent service requirements with the original law,\allowbreak{} and use identical marks on shared singleton arrivals.

Over $[0,T]$,\allowbreak{} the probability of any discrepancy is bounded by the expected number of discarded job events plus the expected total number of added task events. These are bounded by

\[\Lambda T\left(\mathbb P(N\ge2)+\mathbb E[N\mathbf1_{N\ge2}]\right)\le\frac32\Lambda T\mathbb E[N(N-1)].\]

Here $1\le N(N-1)/2$ and $N\le N(N-1)$ when $N\ge2$. Substitution of $\Lambda=n\lambda/k$ yields the collision term in the stated bound. This argument is nonasymptotic and does not require $2k\le n$.

Start the original and modified queues empty. Until the first discrepancy their tagged workload paths agree,\allowbreak{} because they have identical tagged arrival times and task marks. Thus the same bound controls total variation at time $T$.

\subsection*{Step 2:\allowbreak{} couple empty starts to stationary starts}

For the original system,\allowbreak{} couple one copy with stationary workloads at time zero to a copy with empty workloads,\allowbreak{} using the same future routing,\allowbreak{} arrivals and service marks. The stationary workload is coordinatewise above the empty-start workload. Whenever a stationary coordinate hits zero,\allowbreak{} both copies have zero workload there and then agree forever at that coordinate.

Each stationary coordinate is a marginal M/\allowbreak{}G/\allowbreak{}1 queue with the same $\lambda$ and service law,\allowbreak{} even though coordinates are dependent. By a union bound,\allowbreak{} the probability that any of the $k$ coordinates has failed to couple by $T$ is at most $kr(T)$. The independent comparison system has the same marginal law,\allowbreak{} so its own stationary-versus-empty coupling has the same bound. The triangle inequality between these two comparisons and Step 1 proves the modular estimate. Independence of the original stationary coordinates is neither assumed nor required.

\subsection*{Step 3:\allowbreak{} exponential-service emptying tail}

Consider the marginal M/\allowbreak{}M/\allowbreak{}1 queue count $Q$ until its first hit $\tau_0$ of zero. For $q\ge1$,\allowbreak{} its generator applied to $h(q)=a^q$,\allowbreak{} with $a=\sqrt{\mu/\lambda}>1$,\allowbreak{} is

\[\frac{{\cal L}h(q)}{h(q)}=\lambda(a-1)+\mu(a^{-1}-1)=-(\sqrt\mu-\sqrt\lambda)^2=-c.\]

Stopping first at zero and at a finite upper count level makes the standard exponential generator martingale legitimate. Nonnegativity and Fatou's lemma upon removal of the upper stopping level give

\[\mathbb E_q\left[e^{c(T\wedge\tau_0)}a^{Q(T\wedge\tau_0)}\right]\le a^q.\]

On $\{\tau_0>T\}$ the integrand is at least $e^{cT}$,\allowbreak{} so $\mathbb P_q(\tau_0>T)\le a^qe^{-cT}$. Averaging over the stationary geometric queue count yields

\[\mathbb E[a^{Q(0)}]=\frac{1-\rho}{1-\rho a}=\frac{1-\rho}{1-\sqrt\rho}=1+\sqrt\rho.\]

A marginal workload is zero exactly when its queue count is zero. The count process here is used only to bound this marginal stopping time;\allowbreak{} the pathwise workload coupling in Step 2 remains driven by shared task marks. No erroneous identification of arbitrary coupled count and workload paths is needed.

\subsection*{Consequences and limits}

Taking $T=2\log n/c$ gives

\[d_{\rm TV}(\pi_{n,k},\nu_k)\le\frac{2k(1+\sqrt\rho)}{n^2}+\frac{3\lambda k(k-1)^2\log n}{c(n-1)}.\]

Hence the distance tends to zero when $k=o((n/\log n)^{1/3})$ under exponential service. This is a sufficient regime,\allowbreak{} not an optimal threshold,\allowbreak{} and says nothing at the boundary $k\asymp n^{1/3}$ by itself. Poisson arrivals see the stationary pre-arrival workload law by PASTA. The arriving job selects its uniform target set independently of that law;\allowbreak{} exchangeability makes its selected workloads have the law of the specified tagged set. Appending independent tagged-job service requirements and taking a maximum is a Markov kernel,\allowbreak{} so total variation contracts;\allowbreak{} the same error controls the job-delay distribution against the independent comparison.

More generally,\allowbreak{} if $r(T)\le C T^{-p}$ for some fixed $p>0$,\allowbreak{} the modular estimate gives an error of order $k/T^p+Tk^3/n$. Balancing at $T\asymp(n/k^2)^{1/(p+1)}$ gives the sufficient regime $k=o(n^{p/(3p+1)})$. This conclusion is conditional on the stated emptying-tail bound;\allowbreak{} no unproved implication from a service moment is used. For $p=1$,\allowbreak{} the source's finite-second-moment estimate recovers the known exponent $1/4$.

\subsection*{Attempts to break the argument}

The vulnerable points are the growing number of tagged queues,\allowbreak{} simultaneous multi-hit arrivals,\allowbreak{} the marginal rate after removing those arrivals,\allowbreak{} and the passage to stationarity. The proof restores lost marginal rates rather than silently comparing to lower-load queues. It uses a union bound over all $k$ queues and retains both empty-start errors. The exponential tail constant tends to zero as load approaches one,\allowbreak{} so this packet does not claim a uniform heavy-traffic result. Correlated within-job service marks are excluded. Exponential control of the stationary emptying time has not been proved here for general service laws.

The original paper already has the same broad finite-time/\allowbreak{}emptying-time proof structure and numerical experiments at larger exponents. The observation here could be a known or straightforward corollary;\allowbreak{} bounded search is not a novelty certificate. No full original open-problem solution is promoted. Arithmetic controls are in \path{check_forkjoin_stationary.py} and \path{forkjoin_stationary_checks.json}.

\clearpage

\section*{Proof appendix:\allowbreak{} P0001 cactus forests}

The current ledger disposition and the exact status limits in this packet govern its interpretation. Classical supporting facts and known-result rederivations are distinct from original open-problem solutions.


\subsection*{P0001:\allowbreak{} negative dependence for random forests}

\textbf{Target remains unsolved here.} This note proves a special case and refutes an overly broad proof route. The special case is author checked;\allowbreak{} novelty is not claimed. No independent reviewer or proof assistant has checked this note.

The sourced target is Grimmett,\allowbreak{} \emph{Selected problems in probability theory} (2022),\allowbreak{} Conjecture 4.2,\allowbreak{} PDF page 8:\allowbreak{} for a finite simple graph,\allowbreak{} uniformly choose an acyclic edge subset F. Are distinct edge indicators always negatively correlated? Tang and Zhang,\allowbreak{} arXiv:\allowbreak{}2603.10738v1 (March 2026),\allowbreak{} Conjecture 1.3,\allowbreak{} still states this general question. Their complete-graph results do not settle arbitrary graphs.

\subsection*{A special case with an exact formula}

Let G be a finite cactus graph:\allowbreak{} each edge belongs to at most one simple cycle. Give every edge a finite positive weight w\_\allowbreak{}e. Choose an acyclic edge set F with probability proportional to the product of its selected edge weights. This includes the uniform forest measure when every w\_\allowbreak{}e=\allowbreak{}1.

\textbf{Proposition.} This measure is negatively associated:\allowbreak{} if A and B are disjoint edge sets and f and g are increasing real functions of the indicators in A and B respectively,\allowbreak{} then Cov(f,\allowbreak{}g) <=\allowbreak{} 0.

\textbf{Proof.} Begin with independent Bernoulli edge indicators X\_\allowbreak{}e with p\_\allowbreak{}e=\allowbreak{}w\_\allowbreak{}e/\allowbreak{}(1+\allowbreak{}w\_\allowbreak{}e). Conditioning on acyclicity gives exactly the desired measure,\allowbreak{} because the product probability of an edge set is a constant times its product weight.

The cactus cycles have disjoint edge sets. Acyclicity is therefore the intersection,\allowbreak{} over cycles C,\allowbreak{} of the events that not every edge of C is selected. These events depend on disjoint coordinate sets. Conditioning preserves independence between the cycle blocks and the remaining bridge edges.

Consider one cycle block. Write H for the event that every coordinate equals one and q=\allowbreak{}P(H)=\allowbreak{}product(p\_\allowbreak{}e). For disjoint coordinate sets A and B in that block,\allowbreak{} set a=\allowbreak{}E f,\allowbreak{} b=\allowbreak{}E g,\allowbreak{} f\_\allowbreak{}*=\allowbreak{}f(1\_\allowbreak{}A),\allowbreak{} and g\_\allowbreak{}*=\allowbreak{}g(1\_\allowbreak{}B). Independence gives E(fg)=\allowbreak{}ab. On H,\allowbreak{} f=\allowbreak{}f\_\allowbreak{}* and g=\allowbreak{}g\_\allowbreak{}*. Consequently

E(f | H\textasciicircum{}c) =\allowbreak{} (a-q f\_\allowbreak{}*)/\allowbreak{}(1-q),\allowbreak{}
    E(g | H\textasciicircum{}c) =\allowbreak{} (b-q g\_\allowbreak{}*)/\allowbreak{}(1-q),\allowbreak{}
    E(fg | H\textasciicircum{}c) =\allowbreak{} (ab-q f\_\allowbreak{}* g\_\allowbreak{}*)/\allowbreak{}(1-q).

Subtracting the product of the first two equations from the third yields

Cov(f,\allowbreak{}g | H\textasciicircum{}c)
      =\allowbreak{} -q (f\_\allowbreak{}*-a)(g\_\allowbreak{}*-b)/\allowbreak{}(1-q)\textasciicircum{}2 <=\allowbreak{} 0.

The inequality follows from monotonicity:\allowbreak{} f\_\allowbreak{}* >=\allowbreak{} a and g\_\allowbreak{}* >=\allowbreak{} b. The denominator is positive because every p\_\allowbreak{}e<1.

For completeness,\allowbreak{} independence products preserve negative association. Suppose X and Y are independent negatively associated blocks. Condition first on Y. The conditional covariance of f(X\_\allowbreak{}A,\allowbreak{}Y\_\allowbreak{}C) and g(X\_\allowbreak{}B,\allowbreak{}Y\_\allowbreak{}D) is nonpositive whenever A,\allowbreak{}B are disjoint and C,\allowbreak{}D are disjoint. The two conditional expectations are increasing functions of the disjoint Y-coordinate sets C,\allowbreak{}D;\allowbreak{} their covariance is also nonpositive. The total covariance identity proves the claim for (X,\allowbreak{}Y). Induct over the finitely many blocks,\allowbreak{} including independent bridge coordinates. This proves the proposition.

For two distinct edges e,\allowbreak{}f in the same cycle,\allowbreak{} the formula becomes

Cov(X\_\allowbreak{}e,\allowbreak{}X\_\allowbreak{}f | H\textasciicircum{}c)
      =\allowbreak{} -q (1-p\_\allowbreak{}e)(1-p\_\allowbreak{}f)/\allowbreak{}(1-q)\textasciicircum{}2.

For a uniform cycle with m edges this is

-2\textasciicircum{}(m-2)/\allowbreak{}(2\textasciicircum{}m-1)\textasciicircum{}2.

Edges in different blocks have covariance zero.

\subsection*{Second derivation and adversarial check}

For the uniform cycle,\allowbreak{} count subsets directly. There are 2\textasciicircum{}m-1 forests,\allowbreak{} 2\textasciicircum{}(m-1)-1 containing a given edge,\allowbreak{} and 2\textasciicircum{}(m-2)-1 containing a specified pair. Thus the covariance numerator is

(2\textasciicircum{}(m-2)-1)(2\textasciicircum{}m-1) - (2\textasciicircum{}(m-1)-1)\textasciicircum{}2 =\allowbreak{} -2\textasciicircum{}(m-2),\allowbreak{}

which agrees with the conditional-product calculation. For a triangle it gives -2/\allowbreak{}49. For a tree,\allowbreak{} all edge subsets are forests,\allowbreak{} and independent Bernoulli(1/\allowbreak{}2) indicators give zero covariance. These checks test both the sign and the degenerate no-cycle case.

The strongest reason this does not solve the target is that general cycles share edges. Their exclusion events no longer factor into independent blocks. One cannot simply say that conditioning independent bits on a decreasing event preserves negative association.

Here is a minimal counterexample to that proposed general principle. Take three fair independent bits (X1,\allowbreak{}X2,\allowbreak{}X3),\allowbreak{} conditioned on neither X1 X2 nor X2 X3 equaling one. The five allowed states are

000,\allowbreak{} 100,\allowbreak{} 010,\allowbreak{} 001,\allowbreak{} 101.

Then E X1=\allowbreak{}E X3=\allowbreak{}2/\allowbreak{}5 and E(X1 X3)=\allowbreak{}1/\allowbreak{}5,\allowbreak{} so Cov(X1,\allowbreak{}X3)=\allowbreak{}1/\allowbreak{}25>0. The conditioning event is decreasing,\allowbreak{} but creates positive correlation. This is a counterexample to a proof route,\allowbreak{} not to the forest conjecture:\allowbreak{} the forbidden pairs do not form the circuit family of a graphic matroid. Circuit elimination would require another circuit,\allowbreak{} which is absent.

\subsection*{Exact unresolved obligation}

To extend this argument one needs a valid inequality controlling the interaction of overlapping cycle constraints,\allowbreak{} using graphic-matroid structure. Arbitrary decreasing-event conditioning cannot provide it. The target is not promoted beyond a partial attempt.

\subsection*{Sources}

\begin{itemize}

\item \url{https://arxiv.org/abs/2205.07318} (Conjecture 4.2,\allowbreak{} page 8).

\item \url{https://arxiv.org/html/2603.10738v1} (Conjecture 1.3 and the discussion immediately following).

\end{itemize}

The calculations and proof above were derived during this campaign. They are not attributed to the cited papers.

\clearpage

\section*{Proof appendix:\allowbreak{} P0004 quenched limsup}

The current ledger disposition and the exact status limits in this packet govern its interpretation. Classical supporting facts and known-result rederivations are distinct from original open-problem solutions.


\subsection*{P0004:\allowbreak{} what can be proved about the limsup}

Status:\allowbreak{} author-checked elementary partial result. No novelty claim. The source's question about existence of a limit remains unsolved here.

Target source:\allowbreak{} Geoffrey Grimmett,\allowbreak{} \emph{Selected problems in probability theory},\allowbreak{} arXiv:\allowbreak{}2205.07318,\allowbreak{} Question 2.3,\allowbreak{} PDF pages 3-4. The target is the exponential growth limit of self-avoiding walks on the infinite supercritical bond-percolation cluster of the square lattice.

\subsection*{Deterministic root-independence lemma}

Let G be an infinite connected graph of maximum degree Delta < infinity. Let c\_\allowbreak{}n(v) count n-step self-avoiding walks starting at v,\allowbreak{} with c\_\allowbreak{}0(v)=\allowbreak{}1. Then the quantity

M(v) =\allowbreak{} limsup as n tends to infinity of c\_\allowbreak{}n(v)\textasciicircum{}(1/\allowbreak{}n)

does not depend on v.

Proof. First let v and w be adjacent. Partition the walks from v according to whether they visit w. A walk which avoids w becomes an (n+\allowbreak{}1)-step self-avoiding walk from w after prepending the edge wv. If a walk visits w,\allowbreak{} let k in \{1,\allowbreak{}...,\allowbreak{}n\} be its unique visiting time. Reverse its prefix from v to w and retain its suffix from w. These are respectively k-step and (n-k)-step self-avoiding walks from w. Forgetting their disjointness only enlarges the number of pairs. Thus

c\_\allowbreak{}n(v) <=\allowbreak{} c\_\allowbreak{}(n+\allowbreak{}1)(w) +\allowbreak{} sum\_\allowbreak{}(k=\allowbreak{}1)\textasciicircum{}n c\_\allowbreak{}k(w)c\_\allowbreak{}(n-k)(w).

An infinite connected locally finite graph contains a ray from every vertex,\allowbreak{} by the elementary finite-branching tree argument applied to finite self-avoiding walks. Hence 1 <=\allowbreak{} M(w) <=\allowbreak{} Delta. For every epsilon>0 there is a finite C>=\allowbreak{}1 such that c\_\allowbreak{}m(w)<=\allowbreak{}C(M(w)+\allowbreak{}epsilon)\textasciicircum{}m for every m>=\allowbreak{}0,\allowbreak{} by the definition of limsup and enlarging C for the finitely many exceptions. The displayed inequality gives

c\_\allowbreak{}n(v) <=\allowbreak{} [C(M(w)+\allowbreak{}epsilon) +\allowbreak{} n C\textasciicircum{}2](M(w)+\allowbreak{}epsilon)\textasciicircum{}n.

Take nth roots and limsup,\allowbreak{} then let epsilon decrease to zero. This gives M(v)<=\allowbreak{}M(w). Interchanging the roots gives equality. A finite path connects any two roots. QED.

Independent check:\allowbreak{} apply the argument to a tree. The visiting decomposition still holds although the exact counts can vary dramatically by root at finite n;\allowbreak{} the polynomial factor disappears at the exponential scale. The proof never uses transitivity or stationarity and therefore cannot,\allowbreak{} by itself,\allowbreak{} prove convergence of c\_\allowbreak{}n(v)\textasciicircum{}(1/\allowbreak{}n).

\subsection*{Consequences in percolation and an annealed upper bound}

Using uniqueness of the infinite cluster and translation ergodicity of independent supercritical bond percolation,\allowbreak{} the common limsup is translation-invariant and therefore deterministic. These two standard percolation theorems are dependencies,\allowbreak{} not proved here. This conclusion is already in the literature:\allowbreak{} Hubert Lacoin,\allowbreak{} \emph{Non-coincidence of quenched and annealed connective constants on the supercritical planar percolation cluster},\allowbreak{} arXiv:\allowbreak{}1203.6051.

There is also an elementary upper bound. For a fixed lattice root,\allowbreak{} let c\_\allowbreak{}n\textasciicircum{}omega count open n-step self-avoiding walks,\allowbreak{} including environments where the root has a finite cluster. Every lattice self-avoiding walk uses n distinct edges,\allowbreak{} so

E c\_\allowbreak{}n\textasciicircum{}omega =\allowbreak{} p\textasciicircum{}n c\_\allowbreak{}n\textasciicircum{}lattice.

Write mu for the lattice connective constant. For any epsilon>0,\allowbreak{} eventually c\_\allowbreak{}n\textasciicircum{}lattice<=\allowbreak{}(mu+\allowbreak{}epsilon)\textasciicircum{}n. Markov's inequality yields

P(c\_\allowbreak{}n\textasciicircum{}omega >=\allowbreak{} [p(mu+\allowbreak{}2epsilon)]\textasciicircum{}n)
        <=\allowbreak{} ((mu+\allowbreak{}epsilon)/\allowbreak{}(mu+\allowbreak{}2epsilon))\textasciicircum{}n.

The right side is summable. Borel-Cantelli,\allowbreak{} followed by a countable epsilon sequence tending to zero,\allowbreak{} proves limsup(c\_\allowbreak{}n\textasciicircum{}omega)\textasciicircum{}(1/\allowbreak{}n)<=\allowbreak{}p mu almost surely. It remains valid after conditioning the root to belong to the infinite cluster. Lacoin proves a strictly smaller bound in the planar supercritical regime p<1;\allowbreak{} that strict improvement is not derived here.

\subsection*{Why the obvious subadditivity route stops}

Cutting a walk at time m only gives

c\_\allowbreak{}(m+\allowbreak{}n)\textasciicircum{}omega(v) <=\allowbreak{} c\_\allowbreak{}m\textasciicircum{}omega(v) max\_\allowbreak{}(u reachable in m steps) c\_\allowbreak{}n\textasciicircum{}omega(u).

The endpoint maximum depends on the same environment and is not c\_\allowbreak{}n\textasciicircum{}omega(v),\allowbreak{} nor a single predetermined stationary translate. Fekete's lemma therefore does not apply to log c\_\allowbreak{}n\textasciicircum{}omega(v). Kingman's theorem needs a specified stationary subadditive process;\allowbreak{} no such process has been constructed by this inequality. Replacing the maximum with a typical-root value would require a uniform estimate over the growing set of endpoints and cannot be justified by an almost-sure statement at each fixed root alone.

Exact obstruction:\allowbreak{} prove that the liminf has the same exponential rate as the deterministic limsup. Root-independence,\allowbreak{} annealed expectations,\allowbreak{} and subexponential prefactors in the displayed comparison do not resolve it.

\clearpage

\section*{Proof appendix:\allowbreak{} P0032 brownian prefix candidate}

The current ledger disposition and the exact status limits in this packet govern its interpretation. Classical supporting facts and known-result rederivations are distinct from original open-problem solutions.


\subsection*{P0032:\allowbreak{} a Brownian-prefix embedding argument}

Written 2026-09-06. Status:\allowbreak{} complete candidate proof,\allowbreak{} undergoing author checks and a prior-literature search. No independent review and no novelty claim. Do not promote this file merely because a ledger validator accepts it.

Source problem:\allowbreak{} Burdzy and Scheutzow,\allowbreak{} \emph{Forward Brownian motion},\allowbreak{} arXiv:\allowbreak{}1302.6958,\allowbreak{} Problem 7.5,\allowbreak{} PDF page 26 (published 2014). The source asks whether a process on nonnegative time whose law is equivalent to Wiener measure can be preceded by a random finite path segment so that the concatenated process has Wiener law.

The interpretation here allows an extension of the probability space and dependence between prefix and suffix. The joining time is an arbitrary finite random time,\allowbreak{} not a stopping time for the ordinary Brownian filtration. Requiring a Brownian stopping time would immediately force the suffix itself to have Wiener law by the strong Markov property,\allowbreak{} so that different requirement cannot hold for a general equivalent law.

\subsection*{A general lemma}

Let (E,\allowbreak{}F,\allowbreak{}P) be a standard Borel probability space,\allowbreak{} and let theta:\allowbreak{}E->E be an invertible,\allowbreak{} measurable,\allowbreak{} measure-preserving,\allowbreak{} ergodic transformation with measurable inverse. Let Q have density f>=\allowbreak{}0 with respect to P,\allowbreak{} with integral f dP=\allowbreak{}1. On an extension supporting independent uniform random variables,\allowbreak{} there exists an almost surely finite integer N>=\allowbreak{}0 such that,\allowbreak{} for Z with law P,\allowbreak{} theta\textasciicircum{}N Z has law Q.

This is a filling construction related to classical work of Chacon-Ornstein,\allowbreak{} Rost,\allowbreak{} and Thorisson. The following proof is supplied in full;\allowbreak{} claiming its application as a new research result requires a separate literature check.

\subsection*{Construction of remaining densities}

For an integrable density h let Lh=\allowbreak{}h composed with theta\textasciicircum{}(-1). Thus pushing hP forward under theta gives (Lh)P. Set

a\_\allowbreak{}0=\allowbreak{}1,\allowbreak{}  b\_\allowbreak{}0=\allowbreak{}f,\allowbreak{}
    m\_\allowbreak{}n=\allowbreak{}min(a\_\allowbreak{}n,\allowbreak{}b\_\allowbreak{}n),\allowbreak{}
    a\_\allowbreak{}(n+\allowbreak{}1)=\allowbreak{}L(a\_\allowbreak{}n-m\_\allowbreak{}n),\allowbreak{}
    b\_\allowbreak{}(n+\allowbreak{}1)=\allowbreak{}b\_\allowbreak{}n-m\_\allowbreak{}n.

All these functions are measurable and nonnegative. Inductively 0<=\allowbreak{}a\_\allowbreak{}n<=\allowbreak{}1. The b\_\allowbreak{}n decrease and are bounded above by f,\allowbreak{} so they have an integrable pointwise limit b\_\allowbreak{}infinity. Since L preserves integrals,\allowbreak{} both remaining densities have the same mass r\_\allowbreak{}n:\allowbreak{}

integral a\_\allowbreak{}n dP =\allowbreak{} integral b\_\allowbreak{}n dP =\allowbreak{} r\_\allowbreak{}n,\allowbreak{}
    r\_\allowbreak{}n decreases to r =\allowbreak{} integral b\_\allowbreak{}infinity dP,\allowbreak{}
    sum\_\allowbreak{}(n>=\allowbreak{}0) integral m\_\allowbreak{}n dP =\allowbreak{} 1-r.

\subsection*{The remaining mass must be zero}

Iteration gives the exact identity

a\_\allowbreak{}n =\allowbreak{} 1 - sum\_\allowbreak{}(k=\allowbreak{}0)\textasciicircum{}(n-1) L\textasciicircum{}(n-k)m\_\allowbreak{}k.

For each fixed k,\allowbreak{} the ergodic theorem for the measure-preserving transformation theta\textasciicircum{}(-1) gives

(1/\allowbreak{}M) sum\_\allowbreak{}(n=\allowbreak{}k+\allowbreak{}1)\textasciicircum{}(M-1) L\textasciicircum{}(n-k)m\_\allowbreak{}k
       -> integral m\_\allowbreak{}k dP

in L1(P). The L1 assertion follows from the usual almost-sure ergodic theorem and equality of the integrals for nonnegative averages,\allowbreak{} or by approximation from bounded functions. Initial omitted terms do not affect the limit.

Truncate the sum in k at K. Its tail has L1 norm at most sum\_\allowbreak{}(k>=\allowbreak{}K) integral m\_\allowbreak{}k dP,\allowbreak{} uniformly in M. That tail tends to zero as K grows. Therefore

(1/\allowbreak{}M) sum\_\allowbreak{}(n=\allowbreak{}0)\textasciicircum{}(M-1) a\_\allowbreak{}n -> r

in L1(P),\allowbreak{} where r on the right is the constant function.

Suppose r>0. There is a delta>0 such that B=\allowbreak{}\{b\_\allowbreak{}infinity>=\allowbreak{}delta\} has positive P-measure. On B,\allowbreak{}

m\_\allowbreak{}n=\allowbreak{}min(a\_\allowbreak{}n,\allowbreak{}b\_\allowbreak{}n) >=\allowbreak{} min(a\_\allowbreak{}n,\allowbreak{}delta)
        >=\allowbreak{} min(1,\allowbreak{}delta) a\_\allowbreak{}n,\allowbreak{}

where the last step uses a\_\allowbreak{}n<=\allowbreak{}1. Averaging and integrating over B gives

liminf\_\allowbreak{}M (1/\allowbreak{}M) sum\_\allowbreak{}(n<M) integral\_\allowbreak{}B m\_\allowbreak{}n dP
        >=\allowbreak{} min(1,\allowbreak{}delta) r P(B) > 0.

But the entire sum of the integrals of m\_\allowbreak{}n is at most one,\allowbreak{} so its Cesaro average tends to zero. This contradiction proves r=\allowbreak{}0. In particular sum\_\allowbreak{}n m\_\allowbreak{}n=\allowbreak{}f in L1 and almost everywhere.

\subsection*{Realizing the matching as a random shift}

Sample Z with law P and independent U\_\allowbreak{}0,\allowbreak{}U\_\allowbreak{}1,\allowbreak{}... uniform on [0,\allowbreak{}1]. At stage n,\allowbreak{} if not previously accepted,\allowbreak{} inspect theta\textasciicircum{}n Z and accept when

U\_\allowbreak{}n <=\allowbreak{} h\_\allowbreak{}n(theta\textasciicircum{}n Z),\allowbreak{}
    h\_\allowbreak{}n=\allowbreak{}m\_\allowbreak{}n/\allowbreak{}a\_\allowbreak{}n on \{a\_\allowbreak{}n>0\},\allowbreak{} and h\_\allowbreak{}n=\allowbreak{}0 elsewhere.

By induction,\allowbreak{} the subprobability distribution of theta\textasciicircum{}n Z before this test is a\_\allowbreak{}n P. The mass accepted at stage n has distribution m\_\allowbreak{}n P,\allowbreak{} and the surviving mass pushed forward has distribution a\_\allowbreak{}(n+\allowbreak{}1)P. Hence P(N>=\allowbreak{}n)=\allowbreak{}r\_\allowbreak{}n tends to zero,\allowbreak{} while the law of theta\textasciicircum{}N Z is sum\_\allowbreak{}n m\_\allowbreak{}n P=\allowbreak{}fP=\allowbreak{}Q. Every U\_\allowbreak{}n is independent of Z and of preceding U's,\allowbreak{} so the induction is valid even though h\_\allowbreak{}n can inspect an entire trajectory. This proves the lemma.

\subsection*{Application to Brownian motion}

Let E=\allowbreak{}C\_\allowbreak{}0(R) carry two-sided Wiener measure P. Define

(theta w)(t)=\allowbreak{}w(t+\allowbreak{}1)-w(1),\allowbreak{} t in R.

This is an invertible measurable measure-preserving transformation. It is ergodic:\allowbreak{} the path is measurably equivalent to its bi-infinite sequence of unit-length Brownian increment segments,\allowbreak{} and theta shifts this independent identically distributed sequence by one. Ergodicity of that shift follows by approximating an invariant event by finite-coordinate cylinder events and using independence of sufficiently separated cylinders.

Let Q\_\allowbreak{}plus be the desired law of X on C\_\allowbreak{}0([0,\allowbreak{}infinity)),\allowbreak{} and let f\_\allowbreak{}plus=\allowbreak{}dQ\_\allowbreak{}plus/\allowbreak{}dW\_\allowbreak{}plus. Define f(w)=\allowbreak{}f\_\allowbreak{}plus(w restricted to [0,\allowbreak{}infinity)). Its integral under P is one. Apply the lemma to Q=\allowbreak{}fP. The forward part of theta\textasciicircum{}N W,\allowbreak{}

Y\_\allowbreak{}t=\allowbreak{}W\_\allowbreak{}(N+\allowbreak{}t)-W\_\allowbreak{}N,\allowbreak{} t>=\allowbreak{}0,\allowbreak{}

has exactly law Q\_\allowbreak{}plus,\allowbreak{} and N is finite and nonnegative. The path obtained by taking W on [0,\allowbreak{}N] and continuing with W\_\allowbreak{}N+\allowbreak{}Y\_\allowbreak{}(t-N) is identically the original path W on nonnegative time. It therefore has Wiener law.

If the requested suffix X already lives on a specified probability space,\allowbreak{} use the regular conditional distribution of (W,\allowbreak{}N) given Y in this coupling to extend that space. Standard Borel path spaces ensure this disintegration exists. In the resulting coupling Y=\allowbreak{}X almost surely,\allowbreak{} so the prescribed X receives a prefix without changing its law.

Thus the argument gives the requested embedding under the stated interpretation,\allowbreak{} and even uses only absolute continuity Q\_\allowbreak{}plus<<W\_\allowbreak{}plus rather than equivalence. It does not require a finite expectation for N and does not give a computable finite-lookahead rule.

\subsection*{Checks and strongest objections}

\begin{itemize}

\item Pushforward direction:\allowbreak{} Lh=\allowbreak{}h composed with theta\textasciicircum{}(-1),\allowbreak{} so removing m\_\allowbreak{}n and then moving forward really gives the stated a\_\allowbreak{}(n+\allowbreak{}1). Using h composed with theta would reverse the density transport.

\item The bound a\_\allowbreak{}n<=\allowbreak{}1 is indispensable in the contradiction. It follows from stationarity of the starting probability P;\allowbreak{} an arbitrary initial distribution would need a different argument.

\item No independence of the sampled future blocks is asserted. The density recursion replaces the invalid idea that infinitely long Brownian futures can be used as disjoint rejection trials.

\item The proof preserves the unconditional Wiener marginal of W. Selecting a suffix changes its law but does not condition the original W on the success of an event;\allowbreak{} acceptance occurs almost surely.

\item A two-state cyclic transformation with uniform P and target density (2,\allowbreak{}0) matches half the mass at n=\allowbreak{}0 and the other half after one shift. The density recursion yields Q exactly and checks its direction and normalization.

\item The mathematical and interpretive risk is that Problem 7.5 imposes an additional implicit restriction on the prefix. The displayed source wording contains no independence or stopping-time restriction. Check the surrounding paper and subsequent literature before labeling the problem solved.

\item A pre-existing embedding/\allowbreak{}filling theorem may already imply this conclusion. The argument would remain a valid explanation,\allowbreak{} but novelty and open-at-start status would then need to be withdrawn.

\end{itemize}

Prior-literature leads being checked:\allowbreak{} Pitman-Tang (2015),\allowbreak{} \emph{The Slepian zeros and Brownian bridge embedded in Brownian motion},\allowbreak{} Theorem 5.7 (a two-sided random shift),\allowbreak{} and Chacon-Ornstein/\allowbreak{}Rost filling schemes. Theorem 5.7 alone allows a shift of either sign;\allowbreak{} the nonnegative direction is established directly above.

\subsection*{A determinism obstruction and why the auxiliary uniforms matter}

A counter-check communicated by the parent agent identifies a real boundary of the lemma. If an integer shift n(z) is a deterministic measurable function of z alone,\allowbreak{} the full shifted law has density

g(y) =\allowbreak{} sum\_\allowbreak{}(k in Z) 1\_\allowbreak{}\{n(theta\textasciicircum{}(-k)y)=\allowbreak{}k\}

with respect to P. Indeed,\allowbreak{} for every measurable A,\allowbreak{} partitioning by n=\allowbreak{}k and applying invariance gives the integral of 1\_\allowbreak{}A times this sum. Tonelli applies because the summands are nonnegative. Thus g is integer-valued almost everywhere;\allowbreak{} its integral is one,\allowbreak{} so the sum is finite almost everywhere. Arbitrary noninteger f cannot arise this way. This is a counterexample to a stronger,\allowbreak{} deterministic version of the general lemma,\allowbreak{} not to the displayed randomized construction.

The present argument explicitly uses independent uniforms. Averaging their conditional acceptance probabilities yields real-valued densities m\_\allowbreak{}n,\allowbreak{} so the integer-valued obstruction does not apply. Alternatively,\allowbreak{} a future marginal may hide randomness in the unobserved past;\allowbreak{} a statement about that projection must be distinguished from a statement about the full two-sided law. No deterministic full-path transport is claimed here.

This check was independently suggested but does not constitute an independent review of the complete candidate proof. The author has re-derived the density formula above. Holroyd and Peres,\allowbreak{} \emph{Extra heads and invariant allocations} (2005),\allowbreak{} arXiv:\allowbreak{}math/\allowbreak{}0306402,\allowbreak{} Theorem 17 and the following remark,\allowbreak{} provide another prior-literature lead:\allowbreak{} their greedy transport can restrict displacements on Z to nonnegative integers. Their displayed setting uses binary target capacities;\allowbreak{} the density-matching proof here supplies its own general-density argument.

\clearpage

\section*{Proof appendix:\allowbreak{} P0033 stationary flow counterexample candidate}

The current ledger disposition and the exact status limits in this packet govern its interpretation. Classical supporting facts and known-result rederivations are distinct from original open-problem solutions.


\subsection*{P0033:\allowbreak{} a stationary Galton-Watson flow that is not ergodic}

Written 2026-09-06. Complete candidate counterexample;\allowbreak{} author checking and prior-literature search are in progress. No independent review or novelty claim. This file is not counted as a solved open problem merely because it exists.

Source:\allowbreak{} Lyons and Peres,\allowbreak{} \emph{Probability on Trees and Networks},\allowbreak{} version 19 August 2026,\allowbreak{} Question 17.18,\allowbreak{} printed page 575 (PDF page 592). The question asks whether a flow rule admitting a stationary measure equivalent to Galton-Watson measure must induce an ergodic Markov chain. Proposition 17.17 proves this under strict positivity of all child weights. Exercise 17.22 gives a nonergodic example with a merely absolutely continuous stationary measure;\allowbreak{} equivalence is the important extra requirement here.

\subsection*{Construction and claim}

Let GW be the law of a rooted Galton-Watson tree with independent offspring numbers D satisfying

P(D=\allowbreak{}2)=\allowbreak{}3/\allowbreak{}5,\allowbreak{}   P(D=\allowbreak{}3)=\allowbreak{}2/\allowbreak{}5.

The tree has no leaves and mean offspring 12/\allowbreak{}5. We construct an isomorphism-invariant measurable color c(T) in \{0,\allowbreak{}1\},\allowbreak{} each with probability 1/\allowbreak{}2,\allowbreak{} such that almost surely:\allowbreak{}

\begin{itemize}

\item at a vertex with two children,\allowbreak{} its color is the AND of its child colors;\allowbreak{}

\item at a vertex with three children,\allowbreak{} its color is the OR of its child colors.

\end{itemize}

Both Boolean rules preserve unanimity. Consequently every vertex has at least one child of its own color. At every vertex,\allowbreak{} the flow chooses uniformly among those children. The resulting flow is consistent under restriction to a descendant subtree. GW itself is stationary for the Markov chain that follows this flow;\allowbreak{} nevertheless the root color is constant along the chain and has probability 1/\allowbreak{}2 for each value. Thus the chain is not ergodic,\allowbreak{} despite a stationary measure equal to GW.

\subsection*{The color is a function of the uncolored tree}

For a finite truncation of T at level n,\allowbreak{} give the vertices at that level independent Bernoulli(1/\allowbreak{}2) colors and propagate colors upward using AND at two-child vertices and OR at three-child vertices. Let q\_\allowbreak{}n(T) be the probability,\allowbreak{} over these temporary boundary colors,\allowbreak{} that the root color is one. Thus q\_\allowbreak{}0=\allowbreak{}1/\allowbreak{}2. It is a measurable,\allowbreak{} isomorphism-invariant function of the uncolored tree through level n,\allowbreak{} with values in [0,\allowbreak{}1].

The mean output probability when the children independently have probability p of color one is

F(p)=\allowbreak{}(3/\allowbreak{}5)p\textasciicircum{}2+\allowbreak{}(2/\allowbreak{}5)[1-(1-p)\textasciicircum{}3].

Direct substitution gives F(1/\allowbreak{}2)=\allowbreak{}1/\allowbreak{}2. Let F\_\allowbreak{}n denote the sigma-field of the tree through level n,\allowbreak{} meaning that the offspring numbers at depths below n have been exposed. Conditional on F\_\allowbreak{}n,\allowbreak{} the offspring numbers of the level-n vertices are independent with the displayed law. The expected probability of a one at each such vertex after one additional level is 1/\allowbreak{}2. The root Boolean formula is a multilinear polynomial in the probabilities at its boundary vertices,\allowbreak{} because it is a read-once formula on a tree. Taking conditional expectations therefore gives

E[q\_\allowbreak{}(n+\allowbreak{}1) | F\_\allowbreak{}n]=\allowbreak{}q\_\allowbreak{}n.

Hence q\_\allowbreak{}n is a bounded martingale. It converges almost surely and in L2 to a measurable q\_\allowbreak{}infinity,\allowbreak{} with E q\_\allowbreak{}infinity=\allowbreak{}1/\allowbreak{}2.

We must show that q\_\allowbreak{}infinity is \{0,\allowbreak{}1\}-valued;\allowbreak{} otherwise a random boundary color would remain and would not define the required flow on uncolored trees. Put r\_\allowbreak{}n=\allowbreak{}E[q\_\allowbreak{}n\textasciicircum{}2]. Equivalently,\allowbreak{} use two independent boundary colorings of the same uncolored truncation:\allowbreak{} r\_\allowbreak{}n is the probability that both root outputs are one. Each output has marginal probability 1/\allowbreak{}2. For a two-child AND,\allowbreak{} the probability of two one outputs is r\_\allowbreak{}n\textasciicircum{}2. For a three-child OR,\allowbreak{} inclusion-exclusion gives

1-2(1/\allowbreak{}2)\textasciicircum{}3+\allowbreak{}r\_\allowbreak{}n\textasciicircum{}3=\allowbreak{}3/\allowbreak{}4+\allowbreak{}r\_\allowbreak{}n\textasciicircum{}3,\allowbreak{}

since for a pair with both marginals 1/\allowbreak{}2 the probabilities of 00 and 11 agree. Different child subtrees are independent. It follows that

r\_\allowbreak{}0=\allowbreak{}1/\allowbreak{}4,\allowbreak{}
    r\_\allowbreak{}(n+\allowbreak{}1)=\allowbreak{}3/\allowbreak{}10+\allowbreak{}(3/\allowbreak{}5)r\_\allowbreak{}n\textasciicircum{}2+\allowbreak{}(2/\allowbreak{}5)r\_\allowbreak{}n\textasciicircum{}3.

For 1/\allowbreak{}4 <=\allowbreak{} r < 1/\allowbreak{}2,\allowbreak{} the increment factors as

3/\allowbreak{}10+\allowbreak{}(3/\allowbreak{}5)r\textasciicircum{}2+\allowbreak{}(2/\allowbreak{}5)r\textasciicircum{}3-r
      =\allowbreak{}(r-1/\allowbreak{}2)[(2/\allowbreak{}5)r\textasciicircum{}2+\allowbreak{}(4/\allowbreak{}5)r-3/\allowbreak{}5] > 0.

The quadratic factor is negative throughout [1/\allowbreak{}4,\allowbreak{}1/\allowbreak{}2],\allowbreak{} and the recursion maps this interval into itself because it is increasing and fixes 1/\allowbreak{}2. The only fixed point in this interval is 1/\allowbreak{}2:\allowbreak{} the other positive root is (sqrt(10)-2)/\allowbreak{}2 > 1/\allowbreak{}2. Thus r\_\allowbreak{}n increases to 1/\allowbreak{}2. By L2 convergence,\allowbreak{}

E[q\_\allowbreak{}infinity(1-q\_\allowbreak{}infinity)] =\allowbreak{} 1/\allowbreak{}2-lim r\_\allowbreak{}n =\allowbreak{} 0.

Nonnegativity forces q\_\allowbreak{}infinity in \{0,\allowbreak{}1\} almost surely. Define c(T)=\allowbreak{}q\_\allowbreak{}infinity(T) on this probability-one set. The recursion for q\_\allowbreak{}n at the root is a continuous polynomial in the q\_\allowbreak{}(n-1) values of its child subtrees,\allowbreak{} so passage to the limit proves the AND/\allowbreak{}OR color recursion. Apply this simultaneously to all vertices,\allowbreak{} a countable set. For a global Borel definition on leafless locally finite rooted trees,\allowbreak{} set c(T)=\allowbreak{}0 outside the Borel set where the displayed limit exists and is binary (and outside the offspring-2-or-3 tree space). At each vertex apply the same local same-color rule;\allowbreak{} if there are no same-color children,\allowbreak{} use equal weights on all children. Defining ray probabilities by products of these local weights gives consistency everywhere,\allowbreak{} including the null exceptional set,\allowbreak{} and leaves all GW claims unchanged.

\subsection*{Stationarity with no density change}

Write A\_\allowbreak{}c=\allowbreak{}\{T:\allowbreak{}c(T)=\allowbreak{}c\} and GW\_\allowbreak{}c=\allowbreak{}GW(. | A\_\allowbreak{}c),\allowbreak{} for c=\allowbreak{}0,\allowbreak{}1. Both events have probability 1/\allowbreak{}2. Conditional on the root offspring number D=\allowbreak{}d and the vector of its child colors (c\_\allowbreak{}1,\allowbreak{}...,\allowbreak{}c\_\allowbreak{}d),\allowbreak{} the child subtrees are independent,\allowbreak{} and child i has law GW\_\allowbreak{}(c\_\allowbreak{}i). This follows directly from independent Galton-Watson child subtrees and the fact that each color is a measurable function of its own subtree.

The color of the parent is a deterministic function of d and this vector. Let I\_\allowbreak{}c=\allowbreak{}\{i:\allowbreak{}c\_\allowbreak{}i=\allowbreak{}c\} when the parent color is c. This set is nonempty by the AND/\allowbreak{}OR rules. Choosing uniformly from I\_\allowbreak{}c uses no further information about the internal shape of the chosen subtree. For every bounded measurable function h,\allowbreak{}

E[ sum\_\allowbreak{}(i in I\_\allowbreak{}c) h(T\_\allowbreak{}i)/\allowbreak{}|I\_\allowbreak{}c| | D=\allowbreak{}d,\allowbreak{} (c\_\allowbreak{}i) ]
       =\allowbreak{} integral h dGW\_\allowbreak{}c.

Averaging conditional on the parent color c proves GW\_\allowbreak{}c K=\allowbreak{}GW\_\allowbreak{}c for the child-selection kernel K. Finally,\allowbreak{}

GW K =\allowbreak{} (1/\allowbreak{}2) GW\_\allowbreak{}0 K +\allowbreak{} (1/\allowbreak{}2) GW\_\allowbreak{}1 K
         =\allowbreak{} (1/\allowbreak{}2) GW\_\allowbreak{}0 +\allowbreak{} (1/\allowbreak{}2) GW\_\allowbreak{}1 =\allowbreak{} GW.

Thus the stationary measure is exactly GW,\allowbreak{} not merely an absolutely continuous measure supported on a smaller set of trees.

\subsection*{Consistency and failure of ergodicity}

The local rule assigns probability 1/\allowbreak{}|I\_\allowbreak{}c| to each same-color child and zero to the other children. Products of these transition weights assign mass to all finite initial ray segments. Their consistency defines a Borel probability measure on the infinite boundary by the usual cylinder extension. At a descendant receiving positive mass,\allowbreak{} the conditional continuation is the identical rule applied to that descendant subtree. This is a flow rule in the sense of the source.

The event A\_\allowbreak{}1 has GW measure 1/\allowbreak{}2,\allowbreak{} and K(T,\allowbreak{}A\_\allowbreak{}1)=\allowbreak{}1\_\allowbreak{}(A\_\allowbreak{}1)(T) almost surely. Therefore it is a nontrivial invariant event for the stationary Markov chain. Equivalently,\allowbreak{} in its stationary path space the root-color event is invariant under the time shift and has probability 1/\allowbreak{}2. The chain is not ergodic.

\subsection*{Adversarial checks and remaining status work}

\begin{itemize}

\item A recursive distributional fixed point alone is insufficient:\allowbreak{} it might require extra randomness after the full tree is known. The second-moment calculation above proves this does not happen here.

\item The two boundary trials are independent conditional on the tree,\allowbreak{} so their joint probability really is E q\_\allowbreak{}n\textasciicircum{}2. Their unconditional correlation is retained in r\_\allowbreak{}n;\allowbreak{} treating them as independent would destroy the proof.

\item Merely conditioning GW on the existence of an n-ray does not give equivalence to GW,\allowbreak{} as the source's Exercise 17.22 illustrates. Here the two invariant classes partition a GW-full set and GW itself is stationary.

\item The selected child may depend on all child colors,\allowbreak{} but these are exactly the conditions used in the disintegration argument. No unconditional independence after selection is asserted.

\item Rooted isomorphisms permute children;\allowbreak{} both Boolean rules,\allowbreak{} q\_\allowbreak{}n,\allowbreak{} its limit,\allowbreak{} and uniform selection of eligible children are invariant under those permutations. No hidden labels are used.

\item Strict positivity of every child weight fails on some configurations,\allowbreak{} as it must in view of Proposition 17.17. At a two-child vertex with colors 0,\allowbreak{}1 only the zero child is chosen;\allowbreak{} at a three-child vertex with both colors only the one children are chosen.

\item All arithmetic is exact rational algebra. Independent symbolic or finite-depth checks are diagnostics only;\allowbreak{} the argument proves all depths and the infinite-tree limit.

\item Search prior work on endogeny and interpretable tree automata before calling this a new resolution. Johnson,\allowbreak{} Podder,\allowbreak{} and Skerman,\allowbreak{} \emph{Random tree recursions:\allowbreak{} which fixed points correspond to tangible sets of trees?},\allowbreak{} arXiv:\allowbreak{}1808.03019,\allowbreak{} is a relevant primary source. Its general endogeny criteria are not assumed in this self-contained argument.

\end{itemize}

Exact scope check:\allowbreak{} the definition in Section 17.5,\allowbreak{} printed pages 573-574,\allowbreak{} explicitly permits nonnegative child probabilities and imposes conditional-flow consistency only at children with positive mass. Equivariance is required and was checked above. Thus zeros in this construction are permitted by the source definition;\allowbreak{} they are not a silent relaxation of it.

\clearpage

\section*{Proof appendix:\allowbreak{} P0096 FPP minimum bound counterexample candidate}

The current ledger disposition and the exact status limits in this packet govern its interpretation. Classical supporting facts and known-result rederivations are distinct from original open-problem solutions.


\subsection*{P0096:\allowbreak{} continuous-law counterexample to the four-edge lower bound}

Status update19:\allowbreak{}18Z2026-09-06:\allowbreak{} ALREADY REFUTED IN THE LITERATURE. Ahlberg1911.09949 (2019;\allowbreak{}published2020) explicitly settles this question. The proof below is retained as an author-checked rederivation;\allowbreak{} no novelty or original solved-open credit. The two-page primary paper was read in full;\allowbreak{} no independent review of this document occurred.

\subsection*{Source and exact target}

Auffinger,\allowbreak{} Damron and Hanson,\allowbreak{} *50 years of first-passage percolation*,\allowbreak{} arXiv:\allowbreak{}1511.03262,\allowbreak{} manuscript26September2016,\allowbreak{} Question12 (PDF p.33;\allowbreak{} repeated p.146),\allowbreak{} asks whether iid nonnegative edge passage times on Z\textasciicircum{}2 with F(0)=\allowbreak{}0 always satisfy

\[\mu(e_1)\geq \mathbb E\min(\tau_1,\tau_2,\tau_3,\tau_4).\]

E\_\allowbreak{}4 denotes expectation of a minimum of four independent copies,\allowbreak{} not a fourth moment. The survey motivates this by Alm-Deijfen's simulation study,\allowbreak{} arXiv:\allowbreak{}1412.5924 (December2014). That paper describes an observed relation for selected distribution families. The target here is the survey's universal formulation,\allowbreak{} not a universal theorem attributed to the simulation authors.

Primary sources:\allowbreak{} \url{https://arxiv.org/pdf/1511.03262} and \url{https://arxiv.org/pdf/1412.5924.} The two standard time-constant inputs below are stated in the latter paper,\allowbreak{} equations(2),\allowbreak{}(3),\allowbreak{} PDF p.3.

\subsection*{Construction and lower bound}

Independently for each lattice edge e,\allowbreak{} take eta\_\allowbreak{}e in \{0,\allowbreak{}1\} with P(eta\_\allowbreak{}e=\allowbreak{}0)=\allowbreak{}3/\allowbreak{}4 and an independent U\_\allowbreak{}e uniform on [0,\allowbreak{}1]. For 0<epsilon<=\allowbreak{}1 put

\[\tau_e^{(\epsilon)}=\eta_e+\epsilon U_e.\]

The resulting iid law is bounded and absolutely continuous:\allowbreak{} its density is 3/\allowbreak{}(4epsilon) on (0,\allowbreak{}epsilon) and 1/\allowbreak{}(4epsilon) on (1,\allowbreak{}1+\allowbreak{}epsilon). It has no atoms,\allowbreak{} F\_\allowbreak{}epsilon(0)=\allowbreak{}0,\allowbreak{} and infimum of support zero. Every moment exists.

For four independent copies,\allowbreak{} the event that all eta values equal1 has probability (1/\allowbreak{}4)\textasciicircum{}4=\allowbreak{}1/\allowbreak{}256. Their minimum is at least1 on this event,\allowbreak{} so

\[\mathbb E\min_{1\le i\le4}\tau_i^{(\epsilon)}\geq1/256.\tag{1}\]

\subsection*{Time constants tend to zero}

For bounded iid nonnegative edge weights the established subadditive formula is

\[\mu(e_1)=\inf_{n\ge1}\mathbb E T(0,ne_1)/n.\tag{2}\]

The established positivity criterion says mu(e\_\allowbreak{}1)>0 iff F(0)<p\_\allowbreak{}c(Z\textasciicircum{}2)=\allowbreak{}1/\allowbreak{}2. Hence the eta-only model has mu\_\allowbreak{}0=\allowbreak{}0 because its zero mass is3/\allowbreak{}4. Each positive-epsilon model has mu\_\allowbreak{}epsilon>0.

Fix n and couple all epsilon values with the same eta,\allowbreak{}U. Trivially T\_\allowbreak{}epsilon(0,\allowbreak{}ne\_\allowbreak{}1)>=\allowbreak{}T\_\allowbreak{}0(0,\allowbreak{}ne\_\allowbreak{}1). Given delta>0,\allowbreak{} choose a finite path gamma whose eta-cost is at most T\_\allowbreak{}0(0,\allowbreak{}ne\_\allowbreak{}1)+\allowbreak{}delta. Such a path exists by the definition of an infimum;\allowbreak{} a minimizing path is not assumed. Since U\_\allowbreak{}e<=\allowbreak{}1,\allowbreak{}

\[T_\epsilon(0,ne_1)\le T_0(0,ne_1)+\delta+\epsilon|\gamma|.\]

First send epsilon to zero and then delta to zero. This proves T\_\allowbreak{}epsilon(0,\allowbreak{}ne\_\allowbreak{}1) tends almost surely to T\_\allowbreak{}0(0,\allowbreak{}ne\_\allowbreak{}1). The straight coordinate path gives 0<=\allowbreak{}T\_\allowbreak{}epsilon(0,\allowbreak{}ne\_\allowbreak{}1)<=\allowbreak{}2n. Dominated convergence therefore gives convergence of expectations.

For each fixed n,\allowbreak{} (2) now implies

\[\limsup_{\epsilon\downarrow0}\mu_\epsilon(e_1)
\leq \mathbb E T_0(0,ne_1)/n.\]

Take the infimum over n and use mu\_\allowbreak{}0=\allowbreak{}0. Nonnegativity yields

\[\lim_{\epsilon\downarrow0}\mu_\epsilon(e_1)=0.\tag{3}\]

By(1),\allowbreak{}(3),\allowbreak{} every sufficiently small positive epsilon satisfies

\[0<\mu_\epsilon(e_1)<1/256\leq\mathbb E\min_{i\le4}\tau_i^{(\epsilon)}.\]

Thus this explicit family eventually reverses the proposed universal inequality. The proof gives existence of a small-enough parameter,\allowbreak{} not a numerical cutoff. All sufficiently large integers k give examples with epsilon=\allowbreak{}1/\allowbreak{}k.

\subsection*{Adversarial checks and limits}

The argument never interchanges an infinite-path minimization with a limit. It uses a finite approximate minimizer for each fixed n,\allowbreak{} then an upper bound valid for every n. No uniform bound on geodesic length is needed.

The limiting distribution has an atom at zero;\allowbreak{} the actual counterexample distributions do not. Their support has a gap,\allowbreak{} which the survey's stated hypotheses do not prohibit. Thus this is stronger than an atomic-law loophole.

A local entrance cost remains positive even when long travel can exploit a supercritical cheap-edge network;\allowbreak{} dividing a single entrance cost by distance explains the failure of the proposed bound. As a second literature lane,\allowbreak{} Cox-Kesten continuity of the time constant,\allowbreak{} cited in Alm-Deijfen,\allowbreak{} gives the same limiting conclusion for this bounded weakly convergent family. The direct proof above does not need that general continuity theorem.

Searches for Question12,\allowbreak{} Alm-Deijfen lower bounds,\allowbreak{} and minimum-of-four counterexamples did not locate a matching prior construction in the inspected results. This does not establish novelty. Before promotion,\allowbreak{} recheck source notation and standard theorem hypotheses,\allowbreak{} compare the original simulation discussion and later errata,\allowbreak{} and obtain independent review. No mathematical gap has been identified in the current argument,\allowbreak{} but neither independent review nor current open status is claimed.

\subsection*{Later-literature correction and an explicit parameter}

Ahlberg,\allowbreak{} Tertiles and the time constant,\allowbreak{} \url{https://arxiv.org/pdf/1911.09949,} Theorem1 and final paragraph,\allowbreak{} already gives counterexamples. The publication is Journal of Applied Probability57(2),\allowbreak{}407-408(2020),\allowbreak{} DOI10.1017/\allowbreak{}jpr.2020.9. The earlier search statement above records the preliminary search only and is superseded by this finding.

For our family,\allowbreak{} its two-thirds quantile is8epsilon/\allowbreak{}9. Ahlberg proves mu<=\allowbreak{}2t\_\allowbreak{}(2/\allowbreak{}3),\allowbreak{} so epsilon=\allowbreak{}1/\allowbreak{}1024 gives mu<=\allowbreak{}1/\allowbreak{}576<1/\allowbreak{}256<=\allowbreak{}E minfour. This is an explicit atomless bounded example derived from an established theorem. The continuity argument above remains a separate qualitative derivation. No original solution is claimed.

\clearpage

\section*{Proof appendix:\allowbreak{} P0140 thinned levy random drift candidate}

The current ledger disposition and the exact status limits in this packet govern its interpretation. Classical supporting facts and known-result rederivations are distinct from original open-problem solutions.


\subsection*{P0140:\allowbreak{} Poisson marking explains the random critical drift}

Status:\allowbreak{} author-checked derivation,\allowbreak{} no independent review,\allowbreak{} no novelty claim. Excluded from the current-open target and original-solution counts. The filename preserves the initial candidate workflow. The later primary paper already states an indirect identification with Joseph's independent-increment limit. What is supplied here is a direct calculation and an explicit obstruction to holding the critical drift deterministic.

Sources:\allowbreak{} van der Hofstad,\allowbreak{} \emph{Stochastic processes on random graphs},\allowbreak{} July12,\allowbreak{}2017,\allowbreak{} equations4.5.5 and4.5.20,\allowbreak{} Open Problem4.1,\allowbreak{} PDF pp.206,\allowbreak{}210;\allowbreak{} \url{https://www.win.tue.nl/~rhofstad/SaintFlour_SPoRG_fin_Saint_Flour.pdf} . Dhara,\allowbreak{}van der Hofstad,\allowbreak{}van Leeuwaarden,\allowbreak{}Sen,\allowbreak{} \emph{Heavy-tailed configuration models at criticality},\allowbreak{} December2019 version,\allowbreak{} equation2.22 and Remark11,\allowbreak{} PDF pp.9-10;\allowbreak{} \url{https://arxiv.org/pdf/1612.00650} . Joseph's original independent-increment limit:\allowbreak{} \url{https://arxiv.org/abs/1012.2352} . Exact comparison with every convention in Joseph's characteristic function remains a separate check;\allowbreak{} the independence statement below is direct.

\subsection*{Setup and the distinction in the question}

Fix 3<tau<4,\allowbreak{} c>0,\allowbreak{} and mu>0. Let Gamma\_\allowbreak{}i be the arrival times of a rate-one Poisson process and set theta\_\allowbreak{}i=\allowbreak{}(c/\allowbreak{}Gamma\_\allowbreak{}i)\textasciicircum{}(1/\allowbreak{}(tau-1)). Thus the theta\_\allowbreak{}i form a Poisson point process on(0,\allowbreak{}infinity) with intensity

rho(dx)=\allowbreak{}(tau-1)c x\textasciicircum{}(-tau) dx.

Conditionally on these points,\allowbreak{} give each theta\_\allowbreak{}i an independent exponential clock T\_\allowbreak{}i with rate theta\_\allowbreak{}i/\allowbreak{}mu. The source's process with a deterministic real parameter lambda is

S\_\allowbreak{}lambda(t)=\allowbreak{}sum\_\allowbreak{}i theta\_\allowbreak{}i[1\{T\_\allowbreak{}i<=\allowbreak{}t\}-(theta\_\allowbreak{}i/\allowbreak{}mu)t]+\allowbreak{}lambda*t.

The ordered series is well defined. For small theta\_\allowbreak{}i,\allowbreak{} its conditional means are O(theta\_\allowbreak{}i\textasciicircum{}3) and its conditional variances are O(theta\_\allowbreak{}i\textasciicircum{}3);\allowbreak{} sum theta\_\allowbreak{}i\textasciicircum{}3 is finite almost surely. There are finitely many points above1,\allowbreak{} each finite. A process-level construction also follows from the compensated Poisson representation below.

The important parameter for iid degrees is the RANDOM variable Lambda\_\allowbreak{}0 in equation2.22 of the later paper,\allowbreak{} not a fixed lambda. Define

Z=\allowbreak{}lim\_\allowbreak{}(epsilon down0)[sum\_\allowbreak{}(theta\_\allowbreak{}i>epsilon) theta\_\allowbreak{}i\textasciicircum{}2 - integral\_\allowbreak{}(epsilon,\allowbreak{}infinity) x\textasciicircum{}2 rho(dx)].

The small-point part of this compensated Poisson integral converges in L2 and almost surely because integral\_\allowbreak{}(0,\allowbreak{}1) x\textasciicircum{}4 rho(dx)<infinity. The large-point part is finite almost surely and has finite first absolute moment because integral\_\allowbreak{}(1,\allowbreak{}infinity) x\textasciicircum{}2 rho(dx)<infinity. Its second moment need not be finite. Set Lambda\_\allowbreak{}0=\allowbreak{}Z/\allowbreak{}mu.

At the random cutoff epsilon=\allowbreak{}theta\_\allowbreak{}n,\allowbreak{} including the nth point changes the sum only by theta\_\allowbreak{}n\textasciicircum{}2,\allowbreak{} which tends to0. Since

integral\_\allowbreak{}(theta\_\allowbreak{}n,\allowbreak{}infinity) x\textasciicircum{}2 rho(dx)=\allowbreak{}c\textasciicircum{}(2/\allowbreak{}(tau-1))*Gamma\_\allowbreak{}n\textasciicircum{}(1-2/\allowbreak{}(tau-1))/\allowbreak{}(1-2/\allowbreak{}(tau-1)),\allowbreak{}

the definition of Z/\allowbreak{}mu is exactly equation2.22 after telescoping the integrals over[Gamma\_\allowbreak{}(i-1),\allowbreak{}Gamma\_\allowbreak{}i]. This checks the random-drift convention against the primary source,\allowbreak{} not merely by analogy.

\subsection*{Direct independent-increment representation}

For every epsilon>0 the cancellation is an identity:\allowbreak{}

S\_\allowbreak{}0,\allowbreak{}epsilon(t)+\allowbreak{}(t/\allowbreak{}mu)Z\_\allowbreak{}epsilon
=\allowbreak{} sum\_\allowbreak{}(theta\_\allowbreak{}i>epsilon) theta\_\allowbreak{}i 1\{T\_\allowbreak{}i<=\allowbreak{}t\} -(t/\allowbreak{}mu) integral\_\allowbreak{}(epsilon,\allowbreak{}infinity) x\textasciicircum{}2 rho(dx).

By independent Poisson marking,\allowbreak{} the pairs(theta\_\allowbreak{}i,\allowbreak{}T\_\allowbreak{}i) form a Poisson random measure with deterministic intensity

nu(dx,\allowbreak{}ds)=\allowbreak{}rho(dx)*(x/\allowbreak{}mu)*exp(-xs/\allowbreak{}mu) ds.

The jump sum in the last display has independent increments in t,\allowbreak{} since disjoint time intervals select disjoint sets of this Poisson measure. Its centering is deterministic. Passing to finite-dimensional limits preserves independence,\allowbreak{} so

L(t):\allowbreak{}=\allowbreak{}S\_\allowbreak{}0(t)+\allowbreak{}(t/\allowbreak{}mu)Z=\allowbreak{}S\_\allowbreak{}(Lambda\_\allowbreak{}0)(t)

has independent increments. They are generally NONSTATIONARY:\allowbreak{} the jump intensity contains exp(-xs/\allowbreak{}mu).

For completeness,\allowbreak{} this passage is not based on a formally divergent cancellation. Write the truncated expression as M\_\allowbreak{}epsilon(t)+\allowbreak{}D\_\allowbreak{}epsilon(t),\allowbreak{} where

M\_\allowbreak{}epsilon(t)=\allowbreak{}sum\_\allowbreak{}(theta\_\allowbreak{}i>epsilon) theta\_\allowbreak{}i1\{T\_\allowbreak{}i<=\allowbreak{}t\}
             -integral\_\allowbreak{}(epsilon,\allowbreak{}infinity) x(1-exp(-xt/\allowbreak{}mu))rho(dx),\allowbreak{}

D\_\allowbreak{}epsilon(t)=\allowbreak{}integral\_\allowbreak{}(epsilon,\allowbreak{}infinity) x[1-exp(-xt/\allowbreak{}mu)-xt/\allowbreak{}mu]rho(dx).

On bounded time intervals,\allowbreak{} the centered contribution from points x<=\allowbreak{}1 is an L2 martingale in t. The variance contributed by0<x<=\allowbreak{}epsilon is bounded by

(T/\allowbreak{}mu) integral\_\allowbreak{}(0,\allowbreak{}epsilon) x\textasciicircum{}3 rho(dx),\allowbreak{}

which tends to0 because tau<4. Doob's inequality gives uniform-in-time L2 convergence of this small-point martingale. The deterministic drift converges uniformly on bounded time intervals because its integrand near0 is O(x\textasciicircum{}3),\allowbreak{} uniformly for bounded t. The finitely many large points and their finite deterministic compensator cause no problem. Hence L has a cadlag version obtained by this construction,\allowbreak{} and S\_\allowbreak{}0=\allowbreak{}L-(t/\allowbreak{}mu)Z does too.

The deterministic drift can also be evaluated:\allowbreak{}

D(t)=\allowbreak{} - (tau-1)c * Gamma(4-tau)/\allowbreak{}[(tau-2)(tau-3)] * (t/\allowbreak{}mu)\textasciicircum{}(tau-2).

To check the sign and constant,\allowbreak{} substitute u=\allowbreak{}xt/\allowbreak{}mu and integrate twice by parts in the positive integral of u\textasciicircum{}(1-tau)(exp(-u)-1+\allowbreak{}u). All boundary terms vanish for3<tau<4. The remaining integral is Gamma(4-tau)/\allowbreak{}[(tau-2)(tau-3)]. In particular D(t)<0,\allowbreak{} consistent with depletion of previously discovered vertices.

\subsection*{A deterministic lambda does not have independent increments}

This part checks the strongest way the preceding conclusion could be misreported. Consider increments A=\allowbreak{}S\_\allowbreak{}0(mu) and B=\allowbreak{}S\_\allowbreak{}0(2mu)-S\_\allowbreak{}0(mu). A marked point of size x contributes

F1=\allowbreak{}x1\{T<=\allowbreak{}mu\}-x\textasciicircum{}2,\allowbreak{}
F2=\allowbreak{}x1\{mu<T<=\allowbreak{}2mu\}-x\textasciicircum{}2.

Let K(a,\allowbreak{}b)=\allowbreak{}log E exp(aA+\allowbreak{}bB),\allowbreak{} for a,\allowbreak{}b>=\allowbreak{}0. These Laplace transforms are finite. Indeed,\allowbreak{} large points have nonpositive F1,\allowbreak{}F2 once x>=\allowbreak{}1;\allowbreak{} for small points the linear terms have integrable expectation O(x\textasciicircum{}3) and the quadratic remainder has integrable expectation O(x\textasciicircum{}3). The compensated Poisson formula,\allowbreak{} or its finite-cutoff version followed by these bounds,\allowbreak{} gives

K(a,\allowbreak{}a)-K(a,\allowbreak{}0)-K(0,\allowbreak{}a)
=\allowbreak{}integral rho(dx) E\_\allowbreak{}T[(exp(aF1)-1)(exp(aF2)-1)].

For x>=\allowbreak{}1 both F1 and F2 are nonpositive,\allowbreak{} so the integrand is nonnegative. If0<a<=\allowbreak{}1/\allowbreak{}4 and x lies in[a\textasciicircum{}(-1/\allowbreak{}2),\allowbreak{}2a\textasciicircum{}(-1/\allowbreak{}2)],\allowbreak{} then x>=\allowbreak{}2 and F1,\allowbreak{}F2<=\allowbreak{}-x\textasciicircum{}2/\allowbreak{}2. Consequently the integrand is at least(1-exp(-1/\allowbreak{}2))\textasciicircum{}2 on this interval. Its rho-mass is

c*(1-2\textasciicircum{}(1-tau))*a\textasciicircum{}((tau-1)/\allowbreak{}2).

The potentially negative part is confined to0<x<1. The two indicators in F1,\allowbreak{}F2 are disjoint,\allowbreak{} and

E|F1F2| <=\allowbreak{} x\textasciicircum{}3 P(T<=\allowbreak{}2mu)+\allowbreak{}x\textasciicircum{}4 <=\allowbreak{}3x\textasciicircum{}4.

Since |exp(au)-1|<=\allowbreak{}a|u|exp(a|u|) and |Fj|<=\allowbreak{}2 on this range,\allowbreak{} the absolute contribution of this whole region is at most

3a\textasciicircum{}2 exp(4a) integral\_\allowbreak{}(0,\allowbreak{}1) x\textasciicircum{}4 rho(dx)=\allowbreak{}C*a\textasciicircum{}2 exp(4a).

Thus the mixed log-Laplace difference is bounded below by a positive constant times a\textasciicircum{}((tau-1)/\allowbreak{}2) minus C*a\textasciicircum{}2 exp(4a). Since (tau-1)/\allowbreak{}2 is strictly less than2,\allowbreak{} it is positive for all sufficiently small a>0. Independence would make this difference identically0. Therefore A and B are dependent. Adding any deterministic lambda changes only their deterministic means and cannot fix this dependence.

This is a direct counterexample to interpreting the 2017 question as a claim about deterministic lambda after mixing the weights. It does not refute the corrected iid-degree statement,\allowbreak{} whose random drift is explicitly supplied in the later paper.

\subsection*{Checks and remaining status limits}

The argument has two complementary checks:\allowbreak{} Poisson marking proves independence after the exact random-slope cancellation;\allowbreak{} the mixed Laplace estimate disproves independence without it. The latter avoids covariance calculations,\allowbreak{} which would be invalid because the random slope can have infinite variance. The PPP intensity follows independently from the tail count E number\{theta\_\allowbreak{}i>x\}=\allowbreak{}c*x\textasciicircum{}(-(tau-1)). The source's Gamma-coordinate expression telescopes to exactly the same compensated slope.

No simulation is used as proof. No independent referee has checked this text. The 2019 paper already identifies the corrected law indirectly,\allowbreak{} so this derivation does not establish an unproved theorem for the corpus and is not counted as an original solution. Later literature may also contain this direct marking calculation;\allowbreak{} that novelty search remains incomplete.

\clearpage

\section*{Proof appendix:\allowbreak{} P0174 free brownian unimodality constant candidate}

The current ledger disposition and the exact status limits in this packet govern its interpretation. Classical supporting facts and known-result rederivations are distinct from original open-problem solutions.


\subsection*{P0174:\allowbreak{} candidate sharp diameter threshold for free Brownian unimodality}

Status:\allowbreak{} complete candidate proof,\allowbreak{} author-checked in two analytic formulations;\allowbreak{} no independent review and no claim of established novelty. Current-literature search is incomplete. Do not count as a promoted original solution.

Source:\allowbreak{} Hasebe and Ueda,\allowbreak{} \emph{Large time unimodality for classical and free Brownian motions with initial distributions},\allowbreak{} ALEA 15 (2018),\allowbreak{} 353–374,\allowbreak{} Problem 3.3;\allowbreak{} archived arXiv1710.08240v2,\allowbreak{} PDF p.6. The PDF carries the arXiv revision date14 September2018 and a manuscript header4 November2021;\allowbreak{} retain that discrepancy. Primary links:\allowbreak{} \url{https://arxiv.org/abs/1710.08240} and \url{https://alea.math.cnrs.fr/articles/v15/15-15.pdf} . The source defines S(0,\allowbreak{}t) to have variance t,\allowbreak{} and D\_\allowbreak{}mu to be support diameter.

The source proves the sufficient constant4 and asks for the optimal universal C in t>=\allowbreak{}C D\_\allowbreak{}mu\textasciicircum{}2. Its paragraph says the symmetric Bernoulli law on\{-1,\allowbreak{}1\} has threshold t=\allowbreak{}4 and diameter2,\allowbreak{} then prints2<=\allowbreak{}C<=\allowbreak{}4. Those preceding numbers imply the lower bound C>=\allowbreak{}1. The candidate upper bound below is a separate argument,\allowbreak{} not a conclusion drawn from that arithmetic inconsistency.

\subsection*{Candidate theorem}

If mu is a compactly supported probability measure on R with support diameter D,\allowbreak{} then mu free-convolved with the centered semicircle law of variance t is unimodal whenever t>=\allowbreak{}D\textasciicircum{}2. The optimal universal constant in the source's normalization is therefore C=\allowbreak{}1.

When D=\allowbreak{}0 the measure is a point mass and the assertion is immediate. Henceforth D>0. No symmetry,\allowbreak{} density,\allowbreak{} or non-atomicity of mu is assumed.

\subsection*{An elementary moment inequality}

Let Q be a probability law of a centered real variable X supported in an interval of length D. Write a=\allowbreak{}E\_\allowbreak{}Q X\textasciicircum{}2 and b=\allowbreak{}E\_\allowbreak{}Q X\textasciicircum{}4. Then

\[
 a\le D^2/4,\qquad b+3a^2\le D^2a.
\]

For the first inequality,\allowbreak{} if the supporting interval is[l,\allowbreak{}h],\allowbreak{} then E[(X-l)(h-X)]>=\allowbreak{}0 gives a<=\allowbreak{}-lh,\allowbreak{} and -lh<=\allowbreak{}D\textasciicircum{}2/\allowbreak{}4 because l+\allowbreak{}h has a real square. For the second,\allowbreak{} let X' be an independent copy. Pointwise,\allowbreak{} (X-X')\textasciicircum{}4<=\allowbreak{}D\textasciicircum{}2(X-X')\textasciicircum{}2. Taking expectations and using E X=\allowbreak{}0 gives2b+\allowbreak{}6a\textasciicircum{}2<=\allowbreak{}2D\textasciicircum{}2a. This derivation checks both the coefficient3 and the use of the diameter rather than radius.

Equality a=\allowbreak{}D\textasciicircum{}2/\allowbreak{}4 forces l=\allowbreak{}-D/\allowbreak{}2,\allowbreak{}h=\allowbreak{}D/\allowbreak{}2 and Q to put equal mass on the two endpoints:\allowbreak{} equality in E[(X-l)(h-X)]>=\allowbreak{}0 forces endpoint support,\allowbreak{} and centering fixes the weights.

\subsection*{Cauchy-profile lemma}

For R>0 define

\[
 F_R(u)=\int\frac{d\mu(x)}{(x-u)^2+R^2}.
\]

Claim:\allowbreak{} F\_\allowbreak{}R has no local minimum with value at most1/\allowbreak{}D\textasciicircum{}2.

Take any stationary point u of F\_\allowbreak{}R with F\_\allowbreak{}R''(u)>=\allowbreak{}0. Put y=\allowbreak{}R\textasciicircum{}2>0 and

\[
 Z=\int ((x-u)^2+y)^{-2}\,d\mu(x),\qquad
 dQ(x)=Z^{-1}((x-u)^2+y)^{-2}\,d\mu(x).
\]

Differentiation under the integral is justified by R>0 and compact support. Stationarity says E\_\allowbreak{}Q(X-u)=\allowbreak{}0. With a=\allowbreak{}E\_\allowbreak{}Q(X-u)\textasciicircum{}2 and b=\allowbreak{}E\_\allowbreak{}Q(X-u)\textasciicircum{}4,\allowbreak{} the preceding moment inequalities apply. The second derivative is

\[
 \frac{F_R''(u)}{2Z}
 =E_Q\frac{3(X-u)^2-y}{(X-u)^2+y}
 =3-4y E_Q\frac1{(X-u)^2+y}.
\]

Jensen's inequality and F\_\allowbreak{}R''(u)>=\allowbreak{}0 imply

\[
 \frac1{a+y}\le E_Q\frac1{(X-u)^2+y}\le\frac3{4y},
 \qquad y\le3a.
\]

The normalization of Q gives an exact reciprocal-profile identity:\allowbreak{}

\[
 \frac1{F_R(u)}
 =\frac{E_Q[((X-u)^2+y)^2]}{E_Q[(X-u)^2+y]}
 =\frac{b+2ay+y^2}{a+y}.
\]

Consequently

\[
 \begin{aligned}
 \frac1{F_R(u)}-D^2
 &\le\frac{y^2+2ay-3a^2-D^2y}{a+y}\\
 &=\frac{(y-3a)(y+a)+(4a-D^2)y}{a+y}\le0.
 \end{aligned}
\]

The first inequality uses b<=\allowbreak{}D\textasciicircum{}2a-3a\textasciicircum{}2;\allowbreak{} the final two summands are nonpositive because y<=\allowbreak{}3a and4a<=\allowbreak{}D\textasciicircum{}2. In particular,\allowbreak{} F\_\allowbreak{}R(u)>=\allowbreak{}1/\allowbreak{}D\textasciicircum{}2.

The equality case matters for the endpoint t=\allowbreak{}D\textasciicircum{}2. If F\_\allowbreak{}R(u)=\allowbreak{}1/\allowbreak{}D\textasciicircum{}2 and F\_\allowbreak{}R''(u)>=\allowbreak{}0,\allowbreak{} then the two final nonpositive terms both vanish:\allowbreak{} y=\allowbreak{}3a and4a=\allowbreak{}D\textasciicircum{}2,\allowbreak{} since y>0. Equality in the variance bound forces Q to be symmetric Bernoulli on u plus/\allowbreak{}minus D/\allowbreak{}2. On these two points the reweighting denominator is equal,\allowbreak{} so mu is the same symmetric Bernoulli law. Writing c=\allowbreak{}D/\allowbreak{}2 and s for displacement from its midpoint,\allowbreak{} with R\textasciicircum{}2=\allowbreak{}3c\textasciicircum{}2,\allowbreak{} gives

\[
 F_R(u+s)=\frac{s^2+4c^2}{s^4+4c^2s^2+16c^4},
 \qquad
 F_R(u+s)-\frac1{4c^2}
 =-\frac{s^4}{4c^2(s^4+4c^2s^2+16c^4)}<0\quad(s\ne0).
\]

Thus the only possible equality point is a strict maximum,\allowbreak{} not a local minimum. Every local minimum therefore has value strictly greater than1/\allowbreak{}D\textasciicircum{}2,\allowbreak{} proving the claim including equality.

\subsection*{Unimodality via Biane's density parametrization}

For fixed t>=\allowbreak{}D\textasciicircum{}2 and fixed R>0,\allowbreak{} the equation F\_\allowbreak{}R(u)=\allowbreak{}1/\allowbreak{}t has at most two solutions. Indeed,\allowbreak{} if there were three solutions u1<u2<u3,\allowbreak{} then either F\_\allowbreak{}R dips below1/\allowbreak{}t somewhere in[u1,\allowbreak{}u3],\allowbreak{} producing an interior local minimum below that level,\allowbreak{} or F\_\allowbreak{}R>=\allowbreak{}1/\allowbreak{}t throughout that interval,\allowbreak{} making u2 a local minimum at that level. Both alternatives contradict the Cauchy-profile lemma because1/\allowbreak{}t<=\allowbreak{}1/\allowbreak{}D\textasciicircum{}2. The analytic function cannot be identically constant on an interval,\allowbreak{} since F\_\allowbreak{}R tends to0 at infinity.

Now use exactly the density framework stated in source Section2.2 and Lemma3.1. Define v\_\allowbreak{}t(u)>=\allowbreak{}0 by

\[
 \int\frac{d\mu(x)}{(x-u)^2+v_t(u)^2}=\frac1t
\]

where its positive solution exists,\allowbreak{} setting v\_\allowbreak{}t=\allowbreak{}0 elsewhere. Biane's real homeomorphism psi\_\allowbreak{}t satisfies

\[
 p_t(\psi_t(u))=\frac{v_t(u)}{\pi t},
\]

where p\_\allowbreak{}t is the continuous density of mu free-convolved with S(0,\allowbreak{}t),\allowbreak{} analytic wherever positive. Each positive density level p\_\allowbreak{}t=\allowbreak{}a corresponds to F\_\allowbreak{}(pi*a*t)(u)=\allowbreak{}1/\allowbreak{}t,\allowbreak{} and therefore has at most two points. Source Lemma2.2,\allowbreak{} equivalently Lemma3.1,\allowbreak{} now gives unimodality. All source assumptions match:\allowbreak{} mu is compactly supported,\allowbreak{} t>0,\allowbreak{} and the semicircle parameter is variance.

\subsection*{Sharpness}

Let mu be symmetric Bernoulli on\{-c,\allowbreak{}c\},\allowbreak{} c>0,\allowbreak{} so D=\allowbreak{}2c. For c\textasciicircum{}2<t<4c\textasciicircum{}2,\allowbreak{} the positive solution at the symmetry center has v\_\allowbreak{}t(0)\textasciicircum{}2=\allowbreak{}t-c\textasciicircum{}2. Implicit differentiation at0 yields

\[
 v_t''(0)=\frac{3c^2-v_t(0)^2}{v_t(0)(c^2+v_t(0)^2)}
 =\frac{4c^2-t}{t\sqrt{t-c^2}}>0.
\]

Thus v\_\allowbreak{}t has a positive strict local minimum at0,\allowbreak{} and the density under the homeomorphism psi\_\allowbreak{}t does as well. It is not unimodal. For any proposed C<1,\allowbreak{} choose t strictly between max(C,\allowbreak{}1/\allowbreak{}4)D\textasciicircum{}2 and D\textasciicircum{}2. Then t>=\allowbreak{}CD\textasciicircum{}2 but unimodality fails. This proves the lower bound C>=\allowbreak{}1 independently of the source's printed inequality.

\subsection*{Second analytic check and remaining review obligations}

The same moment computation can be applied directly at critical points of v\_\allowbreak{}t. At such a point,\allowbreak{} the profile's first derivative vanishes and the sign of v\_\allowbreak{}t'' equals the sign of F\_\allowbreak{}v-fixed''. The calculation above says a nonnegative second derivative would force t<=\allowbreak{}D\textasciicircum{}2. Hence for t>D\textasciicircum{}2 every critical point of v\_\allowbreak{}t is a strict maximum,\allowbreak{} giving the same threshold. The equality analysis in the Cauchy-profile proof closes the boundary without appealing to weak closure of unimodal laws.

The load-bearing points to challenge are the Q-normalization identity,\allowbreak{} the centered fourth-moment inequality,\allowbreak{} the direction of Jensen,\allowbreak{} and the equality case at t=\allowbreak{}D\textasciicircum{}2. Each is written explicitly;\allowbreak{} neither a numerical validator nor brain retrieval supplies independent mathematical review. Later literature might already contain this sharp threshold or the same argument. The statement remains a candidate until further proof checking and literature comparison are complete.

\subsection*{Parent-reported scoped review receipt}

The parent reports a scoped check recorded at \path{~/<private-repo>/12_cognitive_architecture/audits/reports/2026/09/06/last200-implementation/P0174_free_brownian_candidate_review.md},\allowbreak{} SHA256 \texttt{f7c00a07\allowbreak{}0ac27d95\allowbreak{}03d4b486\allowbreak{}e9c78dbe\allowbreak{}8ed1809f\allowbreak{}137836b2\allowbreak{}dc785ffd\allowbreak{}c100fd1a}. Its checked calculations and source assumptions survived review. This is not human certification,\allowbreak{} a novelty certificate,\allowbreak{} or proof of present openness. Canonical audit-node binding remains pending. For P0189 the parent contributed to the amended proof,\allowbreak{} so the amended argument still requires a different independent reviewer. For P0266/\allowbreak{}P0267 the checked proof hash is \texttt{3afdb8ad\allowbreak{}f2e2b58a\allowbreak{}7778f807\allowbreak{}a9028374\allowbreak{}1394c820\allowbreak{}3e9541e1\allowbreak{}60ca7cc6\allowbreak{}f55be870};\allowbreak{} later moment-order and scope clarifications do not change the displayed construction. Any earlier no-review statement in this evolving packet describes its initial draft,\allowbreak{} superseded only within this explicitly scoped receipt.

\clearpage

\section*{Proof appendix:\allowbreak{} P0189 brown measure strong convergence counterexample candidate}

The current ledger disposition and the exact status limits in this packet govern its interpretation. Classical supporting facts and known-result rederivations are distinct from original open-problem solutions.


\subsection*{P0189:\allowbreak{} candidate counterexample to Brown-measure continuity under strong convergence}

Status:\allowbreak{} complete candidate construction,\allowbreak{} checked by the author;\allowbreak{} no independent review and no novelty claim. This addresses the literal general conjecture in van Handel,\allowbreak{} \emph{The strong convergence phenomenon},\allowbreak{} arXiv2507.00346v2 (11 February2026),\allowbreak{} Section6.9,\allowbreak{} pp.69-70. It does not address the narrower conjectures for independent Ginibre matrices or free limiting families.

Primary source for the question:\allowbreak{} \url{https://arxiv.org/abs/2507.00346} . Primary source for the established limiting Brown measure:\allowbreak{} Fang,\allowbreak{} Jiang,\allowbreak{} Lin and Xu,\allowbreak{} \emph{On generalized universal irrational rotation algebras and the operator u+\allowbreak{}v},\allowbreak{} \url{https://arxiv.org/abs/1210.4771} ,\allowbreak{} Theorem12.2. Its introduction also states simplicity and uniqueness of trace of the irrational rotation algebra. These literature theorems are dependencies,\allowbreak{} not results proved here.

\subsection*{Finite matrices}

Fix an irrational theta in(0,\allowbreak{}1). Choose reduced rational approximants p\_\allowbreak{}j/\allowbreak{}q\_\allowbreak{}j tending to theta,\allowbreak{} with q\_\allowbreak{}j tending to infinity. Let omega\_\allowbreak{}j=\allowbreak{}exp(2*pi*i*p\_\allowbreak{}j/\allowbreak{}q\_\allowbreak{}j). On the standard basis indexed by k modulo q\_\allowbreak{}j define

\[
 U_j e_k=\omega_j^k e_k,\qquad
 S_j e_k=e_{k-1},\qquad
 V_j=e^{\pi i/q_j}S_j.
\]

Then U\_\allowbreak{}j and V\_\allowbreak{}j are unitary and satisfy

\[
 V_jU_j=\omega_jU_jV_j,\qquad U_j^{q_j}=I,\qquad V_j^{q_j}=-I.
\]

Put A\_\allowbreak{}j=\allowbreak{}U\_\allowbreak{}j+\allowbreak{}V\_\allowbreak{}j. Its characteristic polynomial is exactly z\textasciicircum{}\{q\_\allowbreak{}j\}. Indeed,\allowbreak{} in the determinant of zI-U\_\allowbreak{}j-V\_\allowbreak{}j only the all-diagonal permutation and the full cyclic permutation contribute (for q\_\allowbreak{}j>=\allowbreak{}2). The former contributes

\[
 \prod_{k=0}^{q_j-1}(z-\omega_j^k)=z^{q_j}-1,
\]

because omega\_\allowbreak{}j is primitive. The latter contributes minus the product of the cyclic weights,\allowbreak{} namely -exp(pi*i)=\allowbreak{}+\allowbreak{}1. Thus det(zI-A\_\allowbreak{}j)=\allowbreak{}z\textasciicircum{}\{q\_\allowbreak{}j\}. Cayley-Hamilton gives A\_\allowbreak{}j\textasciicircum{}\{q\_\allowbreak{}j\}=\allowbreak{}0,\allowbreak{} and every empirical eigenvalue distribution is delta\_\allowbreak{}0.

An equivalent algebraic check is the root-of-unity binomial identity (U\_\allowbreak{}j+\allowbreak{}V\_\allowbreak{}j)\textasciicircum{}\{q\_\allowbreak{}j\}=\allowbreak{}U\_\allowbreak{}j\textasciicircum{}\{q\_\allowbreak{}j\}+\allowbreak{}V\_\allowbreak{}j\textasciicircum{}\{q\_\allowbreak{}j\}=\allowbreak{}0. The determinant calculation supplies the conclusion without relying on that identity.

\subsection*{Strong convergence of the pair}

Let A\_\allowbreak{}theta be the universal unital irrational rotation C*-algebra generated by unitaries u,\allowbreak{}v satisfying vu=\allowbreak{}exp(2*pi*i*theta)uv. This algebra is simple and has a unique faithful tracial state tau,\allowbreak{} as recalled in the cited primary paper.

Consider the norm quotient

\[
 \mathcal Q=\prod_j M_{q_j}(\mathbb C)\big/\bigoplus_j M_{q_j}(\mathbb C),
\]

where the direct sum consists of sequences whose operator norms tend to zero. The classes of (U\_\allowbreak{}j) and (V\_\allowbreak{}j) are unitaries and satisfy the irrational relation exactly,\allowbreak{} since omega\_\allowbreak{}j tends to exp(2*pi*i*theta). Universality therefore gives a unital *-homomorphism A\_\allowbreak{}theta to Q. Simplicity makes it injective,\allowbreak{} hence isometric. Since the quotient norm of a bounded sequence is its limsup norm,\allowbreak{} for every noncommutative *-polynomial P,\allowbreak{}

\[
 \limsup_j\|P(U_j,V_j,U_j^*,V_j^*)\|
 =\|P(u,v,u^*,v^*)\|.
\]

The same argument applies to every subsequence. If the norms failed to converge to the displayed value,\allowbreak{} a subsequence uniformly separated from that value would contradict the same limsup identity. Consequently the pair converges strongly.

The normalized traces also converge to tau. Using the commutation relation,\allowbreak{} every fixed word is a phase times U\_\allowbreak{}j\textasciicircum{}a V\_\allowbreak{}j\textasciicircum{}b. When b is a fixed nonzero integer and j is large,\allowbreak{} the shift has no diagonal entries,\allowbreak{} so its trace is zero. If b=\allowbreak{}0 and a is fixed nonzero,\allowbreak{} the geometric sum of omega\_\allowbreak{}j\textasciicircum{}\{ak\} is zero for all q\_\allowbreak{}j>|a|. The identity word has trace one,\allowbreak{} and phases converge. These are exactly the canonical irrational-rotation trace values. Thus both norm and tracial versions of strong convergence are satisfied,\allowbreak{} with a faithful limiting trace.

In particular A\_\allowbreak{}j converges strongly,\allowbreak{} including all *-polynomials in A\_\allowbreak{}j alone,\allowbreak{} to a=\allowbreak{}u+\allowbreak{}v.

The survey Definition 1.1 also allows fixed finite-dimensional matrix coefficients. For each fixed D,\allowbreak{} amplify the quotient homomorphism to M\_\allowbreak{}D(A\_\allowbreak{}theta). An injective *-homomorphism is completely isometric,\allowbreak{} and M\_\allowbreak{}D(Q) identifies with the quotient of the product of M\_\allowbreak{}D(M\_\allowbreak{}qj) by norm-null sequences. Therefore the same limsup identity and subsequence argument apply to every D-by-D matrix-valued *-polynomial. The identification follows entrywise in finite D:\allowbreak{} a matrix sequence tends to zero in norm exactly when all its entries do,\allowbreak{} and the quotient norm is the limsup of the block-matrix norms. Thus the precise matrix-coefficient definition is covered.

A separate norm-convergence proof uses the trace computation above for the lower bound and needs no appeal to simplicity. Write P\_\allowbreak{}j for a fixed matrix-valued polynomial and P for its limit. For every positive integer m,\allowbreak{} with normalized traces on the matrix factors,\allowbreak{}
\[
 \|P_j\|\ge\big[(\mathrm{tr}_D\otimes\mathrm{tr}_{q_j})((P_j^*P_j)^m)\big]^{1/(2m)}.
\]
Taking the lower limit in j and then m to infinity gives liminf ||P\_\allowbreak{}j|| >=\allowbreak{} ||P||,\allowbreak{} because the limiting trace tr\_\allowbreak{}D tensor tau is faithful. Complete contractivity of the universal quotient homomorphism gives limsup ||P\_\allowbreak{}j|| <=\allowbreak{} ||P||. This independently proves the required full norm convergence. The parent suggested these two explicit extensions during its first bounded review on 6 September 2026;\allowbreak{} the author checked and incorporated them. The remaining review and novelty status is pending.

\subsection*{The Brown-measure mismatch}

The cited Theorem12.2 states that the Brown measure of u+\allowbreak{}v,\allowbreak{} in the tracial von Neumann algebra generated by the irrational rotation pair,\allowbreak{} is Haar probability measure on the unit circle. It is therefore different from delta\_\allowbreak{}0. The finite empirical measures are all delta\_\allowbreak{}0 and do not converge to the Brown measure of a,\allowbreak{} despite strong convergence.

The same example also prevents a universal lower spectral-radius conclusion:\allowbreak{} every A\_\allowbreak{}j has spectral radius zero,\allowbreak{} whereas the source Theorem11.1 gives spectrum of u+\allowbreak{}v equal to the closed unit disk and hence spectral radius one.

To make the matrices random rather than deterministic,\allowbreak{} take Haar unitary Q\_\allowbreak{}j and replace the pair by (Q\_\allowbreak{}j U\_\allowbreak{}j Q\_\allowbreak{}j*,\allowbreak{}Q\_\allowbreak{}j V\_\allowbreak{}j Q\_\allowbreak{}j*). For every outcome,\allowbreak{} norms,\allowbreak{} normalized traces,\allowbreak{} nilpotency,\allowbreak{} and eigenvalue distributions are unchanged. This supplies genuinely random families satisfying the same convergence. The two members are dependent;\allowbreak{} the construction does not claim otherwise. Equivalently,\allowbreak{} the single random matrix Q\_\allowbreak{}j A\_\allowbreak{}j Q\_\allowbreak{}j* is already a counterexample for a one-variable polynomial.

\subsection*{Optional realization in every integer dimension}

The dimensions q\_\allowbreak{}j suffice for a counterexample along a sequence tending to infinity. If an all-integer-dimension formulation is required,\allowbreak{} take consecutive continued-fraction denominators,\allowbreak{} which are coprime. For each sufficiently large N,\allowbreak{} choose the largest j with q\_\allowbreak{}j*q\_\allowbreak{}\{j+\allowbreak{}1\}<=\allowbreak{}N. The elementary coin theorem gives N=\allowbreak{}a*q\_\allowbreak{}j+\allowbreak{}b*q\_\allowbreak{}\{j+\allowbreak{}1\} for nonnegative integers a,\allowbreak{}b. Define the N-dimensional pair as the block direct sum of a copies of the j-th pair and b copies of the(j+\allowbreak{}1)-st pair. Both block dimensions tend to infinity as N does. Polynomial norms are the maximum of their block norms,\allowbreak{} traces are convex combinations of block traces,\allowbreak{} and the sum remains nilpotent. Thus the same counterexample exists with dimension exactly N for every sufficiently large integer N.

\subsection*{Scope and remaining checks}

This is a candidate refutation of the survey's literal implication from arbitrary strong convergence to empirical Brown-measure convergence. It is not a refutation for independent-entry ensembles,\allowbreak{} free limits,\allowbreak{} or holomorphic polynomials in independent Ginibre matrices. Yi Han's June2026 paper,\allowbreak{} arXiv2606.01664,\allowbreak{} claims the latter case and is compatible with this construction.

Load-bearing obligations are the precise survey scope,\allowbreak{} the primitive-root determinant sign,\allowbreak{} isometry of the quotient homomorphism,\allowbreak{} faithful-trace interpretation,\allowbreak{} and the cited Brown-measure theorem. All are explicit here. It remains possible that the construction or the general failure of continuity is already known;\allowbreak{} current-literature novelty is unverified. No validator,\allowbreak{} numerical eigenvalue computation,\allowbreak{} or brain ingestion is treated as proof review. In fact numerical eigenvalues of these nilpotent matrices would be an especially unreliable check because of their extreme nonnormality.

\subsection*{Parent-reported scoped review receipt}

The parent reports a scoped check recorded at \path{~/<private-repo>/12_cognitive_architecture/audits/reports/2026/09/06/last200-implementation/P0189_brown_measure_candidate_review.md},\allowbreak{} SHA256 \texttt{509f76b5\allowbreak{}e525eae2\allowbreak{}3af4b986\allowbreak{}87498844\allowbreak{}9d19eee3\allowbreak{}ae92e555\allowbreak{}0124936f\allowbreak{}6ee515d3}. Its checked calculations and source assumptions survived review. This is not human certification,\allowbreak{} a novelty certificate,\allowbreak{} or proof of present openness. Canonical audit-node binding remains pending. For P0189 the parent contributed to the amended proof,\allowbreak{} so the amended argument still requires a different independent reviewer. For P0266/\allowbreak{}P0267 the checked proof hash is \texttt{3afdb8ad\allowbreak{}f2e2b58a\allowbreak{}7778f807\allowbreak{}a9028374\allowbreak{}1394c820\allowbreak{}3e9541e1\allowbreak{}60ca7cc6\allowbreak{}f55be870};\allowbreak{} later moment-order and scope clarifications do not change the displayed construction. Any earlier no-review statement in this evolving packet describes its initial draft,\allowbreak{} superseded only within this explicitly scoped receipt.

\clearpage

\section*{Proof appendix:\allowbreak{} P0193 compact multiplicative nonunimodality candidate}

The current ledger disposition and the exact status limits in this packet govern its interpretation. Classical supporting facts and known-result rederivations are distinct from original open-problem solutions.


\subsection*{P0193:\allowbreak{} a compactly supported initial law with persistent failure of log-unimodality}

Status:\allowbreak{} complete author-checked candidate counterexample,\allowbreak{} 2026-09-06. No independent review and no novelty certification. This is a proposed negative answer to Hasebe,\allowbreak{} Ueda and Wang,\allowbreak{} \emph{Log-unimodality for free positive multiplicative Brownian motion},\allowbreak{} arXiv:\allowbreak{}2009.13848v1 (29 September 2020),\allowbreak{} Problem 4.7,\allowbreak{} PDF page 11. The published reference is Colloquium Mathematicum 169 (2022),\allowbreak{} 209-226;\allowbreak{} the archived version used here is the arXiv version. Source:\allowbreak{} \url{https://arxiv.org/abs/2009.13848.}

\subsection*{Claim and normalization}

Let \(\nu=(\delta_{1/10}+\delta_{10})/2\). In the source normalization
\[\Sigma_{\sigma_t}(z)=\exp\!\left(\frac t2\frac{z+1}{z-1}\right),\]
the law \(\sigma_t\boxtimes\nu\) is **not log-unimodal for every \(t\ge5\pi\)**. Its initial law has support in the compact interval \([1/10,10]\),\allowbreak{} so this disproves eventual log-unimodality for arbitrary compactly supported positive initial laws,\allowbreak{} if the argument and source match survive review.

Log-unimodality means unimodality of the image measure under the logarithm;\allowbreak{} ordinary unimodality on the positive line is a different claim. The initial law is multiplicatively symmetric but is not log-unimodal. Thus the source theorem for symmetric log-unimodal initial laws does not apply.

\subsection*{The exact source criterion}

Lemma 4.1(3),\allowbreak{} PDF page 8,\allowbreak{} says that \(\sigma_t\boxtimes\nu\) is log-unimodal if and only if,\allowbreak{} for every \(R\in(0,\pi)\),\allowbreak{} the equation
\[\frac{\sin R}{R}\int_0^\infty\frac{r\xi}{1+r^2\xi^2-2r\xi\cos R}\,d\nu(\xi)=\frac1t\tag{1}\]
has at most two solutions \(r>0\). We use its contrapositive by finding four distinct solutions for one value of \(R\),\allowbreak{} for each specified \(t\).

Put \(a=\log10\),\allowbreak{} \(u=\log r\),\allowbreak{} and
\[K_R(s)=\frac{1}{2(\cosh s-\cos R)},\qquad F_R(u)=\frac{K_R(u-a)+K_R(u+a)}2.\]
The integral in (1) equals \(F_R(u)\),\allowbreak{} because dividing its integrand's numerator and denominator by \(r\xi\) gives
\[\frac{r\xi}{1+r^2\xi^2-2r\xi\cos R}=\frac1{r\xi+(r\xi)^{-1}-2\cos R}=K_R(\log r+\log\xi).\]
Hence (1) is equivalent to
\[F_R(u)=\frac{R}{t\sin R}.\tag{2}\]
For fixed \(R\in(0,\pi)\),\allowbreak{} the function \(F_R\) is continuous,\allowbreak{} positive and even,\allowbreak{} and tends to zero at both infinities.

\subsection*{Selecting the level and time}

The function \(g(R)=R/\sin R\) is strictly increasing from 1 to infinity on \((0,\pi)\). Indeed
\[g'(R)=\frac{\sin R-R\cos R}{\sin^2R}>0,\]
because the numerator vanishes at zero and has derivative \(R\sin R>0\). Also \(g(\pi/2)=\pi/2\).

For each \(t\ge5\pi\),\allowbreak{} there is therefore a unique \(R_t\in[\pi/2,\pi)\) with
\[g(R_t)=\frac t{10}.\]
For this choice the right side of (2) is exactly \(1/10\),\allowbreak{} and \(\cos R_t\le0\).

\subsection*{Four strict crossings}

Since \(\cosh(\log10)=101/20\),\allowbreak{}
\[F_{R_t}(0)=K_{R_t}(a)=\frac1{2(101/20-\cos R_t)}\le\frac{10}{101}<\frac1{10}.\tag{3}\]
At \(u=a\),\allowbreak{} the contribution from the centered kernel alone gives
\[F_{R_t}(a)=\frac{K_{R_t}(0)+K_{R_t}(2a)}2\ge\frac1{4(1-\cos R_t)}\ge\frac18>\frac1{10}.\tag{4}\]
The same inequality holds at \(-a\) by evenness. Together with the limits zero at both infinities,\allowbreak{} (3)-(4) and the intermediate value theorem give a solution of \(F_{R_t}(u)=1/10\) in each of the four disjoint intervals
\[(-\infty,-a),\quad(-a,0),\quad(0,a),\quad(a,\infty).\]
Exponentiation is injective,\allowbreak{} so these give four distinct positive solutions of (1). Lemma 4.1 proves the claim.

\subsection*{Separate algebra check in the original variable}

The same signs can be checked without hyperbolic functions. At \(r=1\),\allowbreak{} both atom integrands equal
\[\frac{10}{101-20\cos R}\le\frac{10}{101}<\frac1{10}.\]
At \(r=10\),\allowbreak{} the atom \(\xi=1/10\) contributes,\allowbreak{} including its mass one half,\allowbreak{}
\[\frac1{2(2-2\cos R)}=\frac1{4(1-\cos R)}\ge\frac18.\]
At \(r=1/10\),\allowbreak{} the atom \(\xi=10\) gives the same lower bound. The integral tends to zero as \(r\downarrow0\) and as \(r\uparrow\infty\). Thus there are crossings in \((0,1/10),(1/10,1),(1,10),(10,\infty)\). These strict rational inequalities include the boundary time \(5\pi\),\allowbreak{} for which \(R_t=\pi/2\).

\subsection*{Adversarial checks and limits}

The crucial direction is the equivalence in Lemma 4.1,\allowbreak{} rather than a merely sufficient unimodality condition. Its full statement and proof must remain pinned to the source. The four crossings concern a positive density level,\allowbreak{} not just multiple extrema in an auxiliary transform. The time is selected separately for each \(t\),\allowbreak{} as the criterion permits. The upper bound on the central value is uniform for \(R\ge\pi/2\);\allowbreak{} the lower bound at the two flanks is uniform for all \(R\in(0,\pi)\). No numerical root finder or plot is used as proof. The argument only needs failure for arbitrarily large times;\allowbreak{} it proves failure for every time in the displayed half-line. It makes no assertion about smaller times or optimality of \(5\pi\).

The large-time mechanism can also be understood from \(K_R(s)\to[4\cosh^2(s/2)]^{-1}\) as \(R\uparrow\pi\):\allowbreak{} its width does not diverge. A separated two-point mixture of this limiting kernel remains bimodal. This observation motivated the construction;\allowbreak{} the finite-time rational inequalities above are the proof.

Focused primary-source searches on 6 September 2026 found the original paper,\allowbreak{} the authors' current publication pages,\allowbreak{} and Auer's 2026 paper on the law started at one. No matching later resolution of this exact compact-initial-law question was verified. This is not a certificate of current openness or originality. Author review and source checking are separate from independent review,\allowbreak{} which is pending.

\subsection*{Scoped parent review receipt}

The parent independently reviewed the historical question and the checked construction. Its report is at \path{~/<private-repo>/12_cognitive_architecture/audits/reports/2026/09/06/last200-implementation/P0193_compact_multiplicative_candidate_review.md},\allowbreak{} SHA256 \texttt{25f73461\allowbreak{}45df9adc\allowbreak{}8b33ff33\allowbreak{}d8029946\allowbreak{}13bd2434\allowbreak{}2b25825c\allowbreak{}6cbc7df0\allowbreak{}d8d9261a}. The report found no gap in its checked scope. Structural validation is not mathematical certification. Present openness,\allowbreak{} novelty,\allowbreak{} human review,\allowbreak{} and unique canonical audit-node binding remain unverified. This receipt does not promote the result to an original solution.

\clearpage

\section*{Proof appendix:\allowbreak{} P0208 wreath graph counterexample candidate}

The current ledger disposition and the exact status limits in this packet govern its interpretation. Classical supporting facts and known-result rederivations are distinct from original open-problem solutions.


\subsection*{P0208:\allowbreak{} the wreath graph and Aldous–Fill Problem 7.27}

Status:\allowbreak{} complete counterexample candidate,\allowbreak{} checked by the author through hitting-time and spectral calculations. No independent review or novelty certificate. The graph itself is an established construction,\allowbreak{} not an invention here. This is a response to a historically sourced question;\allowbreak{} present-day open status remains unverified.

Primary question:\allowbreak{} Aldous and Fill,\allowbreak{} \emph{Reversible Markov Chains and Random Walks on Graphs},\allowbreak{} 2002 unfinished monograph,\allowbreak{} 2014 recompilation,\allowbreak{} Open Problem 7.27,\allowbreak{} PDF p.259,\allowbreak{} \url{https://www.stat.berkeley.edu/~aldous/RWG/book.pdf} . Section 7.2 is internally dated January 31,\allowbreak{} 1994. It asks,\allowbreak{} for arc-transitive graphs other than cycles on n vertices,\allowbreak{} whether maximal mean commute time is O(n log n),\allowbreak{} relaxation time is O(n),\allowbreak{} and the ratio of maximal commute time to twice stationary-average hitting time tends to one. The three assertions are counted as one related question.

\subsection*{Source scope and a printed definition error}

Page 254 prints an inconsistent variable assignment in the definition of arc-transitivity:\allowbreak{} the images of two possible identical starting vertices would be allowed to be distinct. Read literally,\allowbreak{} that formula excludes any graph with a vertex having two neighbors. The surrounding discussion,\allowbreak{} examples of tori and Cayley graphs,\allowbreak{} and Lemma 7.19 unmistakably use the usual definition:\allowbreak{} the automorphism group is transitive on ordered adjacent pairs. Our result concerns that intended standard definition,\allowbreak{} and does not silently claim to satisfy the misprinted formula.

Conder's 2013 Rogla summer-school notes,\allowbreak{} p.8,\allowbreak{} independently describe the same established wreath/\allowbreak{}necklace graph and its arc-transitivity:\allowbreak{} \url{https://www.famnit.upr.si/sl/resources/files/konference/rogla2013/conderlecturenotes.pdf} . This is evidence that the construction is prior art. It does not establish when the probability question was answered.

\subsection*{Construction and symmetry}

Let m be an odd integer at least five,\allowbreak{} and put

\[
 V_m=(\mathbb Z/m\mathbb Z)\times\{0,1\},\qquad
 (i,b)\sim(i+1,c),\ (i,b)\sim(i-1,c)\quad(c=0,1).
\]

This is the lexicographic product of a cycle with an edgeless two-vertex graph. It is connected,\allowbreak{} simple,\allowbreak{} 4-regular,\allowbreak{} and has n=\allowbreak{}2m vertices,\allowbreak{} so it is not an n-cycle. Dihedral symmetries act on the first coordinate. Independently swapping the two vertices in any fiber preserves every edge. Given two directed edges,\allowbreak{} a dihedral symmetry sends their ordered base edges to each other,\allowbreak{} after which independent swaps in the two distinct endpoint fibers match both second coordinates. Thus the graph is arc-transitive.

Use discrete-time simple random walk,\allowbreak{} or its rate-one continuization. The book states on p.115 that the parameters at issue are unchanged by this continuization. The first-coordinate projection is simple random walk on the m-cycle,\allowbreak{} and the second coordinate is replaced at each step by a fresh fair bit independent of the base step.

\subsection*{Exact hitting times}

Write h(v,\allowbreak{}w) for the mean first hitting time,\allowbreak{} with h(w,\allowbreak{}w)=\allowbreak{}0. From a base position at cyclic distance j in \{1,\allowbreak{}...,\allowbreak{}m-1\},\allowbreak{} the mean first arrival at base fiber zero is

\[
 H_j=j(m-j).
\]

Indeed H\_\allowbreak{}0=\allowbreak{}H\_\allowbreak{}m=\allowbreak{}0 and H\_\allowbreak{}j=\allowbreak{}1+\allowbreak{}(H\_\allowbreak{}\{j-1\}+\allowbreak{}H\_\allowbreak{}\{j+\allowbreak{}1\})/\allowbreak{}2. At the first fiber arrival,\allowbreak{} the fair bit hits a specified vertex with probability one half.

Let h be the mean hitting time from its twin to the specified vertex. Starting in that fiber,\allowbreak{} the first positive return of the base walk has mean m. The arrival bit succeeds with probability one half;\allowbreak{} on failure the same twin state recurs. Hence h=\allowbreak{}m+\allowbreak{}h/\allowbreak{}2,\allowbreak{} giving h=\allowbreak{}2m. Therefore

\[
 h((j,b),(0,0))=j(m-j)+m\quad(j\ne0),\qquad
 h((0,1),(0,0))=2m.
\]

The formula is symmetric in source and target and gives 2m-1=\allowbreak{}n-1 for neighboring vertices,\allowbreak{} exactly as Lemma 7.19 requires. Its maximum is floor(m²/\allowbreak{}4)+\allowbreak{}m for m at least five. Thus

\[
 \tau^*=2\lfloor m^2/4\rfloor+2m.
\]

The stationary law is uniform. Summing the same hitting-time formula gives

\[
 \tau_0=\frac1{2m}\left[2m+2\sum_{j=1}^{m-1}
          \{j(m-j)+m\}\right]
       =\frac{m^2-1}{6}+m.
\]

On a vertex-transitive graph the stationary mean to a fixed target equals the source's stationary-average target parameter,\allowbreak{} so either equivalent convention yields this value.

\subsection*{Independent spectral route}

For functions constant on each two-vertex fiber,\allowbreak{} the transition operator is the cycle operator,\allowbreak{} with eigenvalues cos(2\(\pi\)k/\allowbreak{}m),\allowbreak{} k=\allowbreak{}0,\allowbreak{}...,\allowbreak{}m-1. For functions antisymmetric within every fiber it is zero. These orthogonal invariant subspaces both have dimension m and exhaust the state space. Consequently

\[
 \tau_2=\frac1{1-\cos(2\pi/m)}\sim\frac{m^2}{2\pi^2}.
\]

The fundamental-matrix trace is

\[
 \sum_{k=1}^{m-1}\frac1{1-\cos(2\pi k/m)}+m
 =\frac{m^2-1}{6}+m,
\]

which agrees with the direct hitting-time sum. The trigonometric sum can also be recovered from the already derived cycle hitting-time formula by stationary averaging;\allowbreak{} it is not needed as an unproved numerical identity for the counterexample.

\subsection*{Consequences and limitations}

Since n=\allowbreak{}2m,\allowbreak{} both \(\tau\)* and \(\tau\)2 are of order n². They violate the proposed n log n and n upper bounds. Moreover

\[
 \frac{\tau^*}{2\tau_0}
 =\frac{\lfloor m^2/4\rfloor+m}{(m^2-1)/6+m}
 \longrightarrow\frac32,
\]

which violates the third assertion. These are exact formulas,\allowbreak{} not simulation estimates.

The strongest reasons for caution are the old unfinished source,\allowbreak{} its definition typo,\allowbreak{} and the fact that wreath graphs have long been standard. Focused searches on September 6,\allowbreak{} 2026 located the original question and primary graph-theory descriptions but no verified later disposition of Problem 7.27. Absence from that search is not evidence of novelty. A future reviewer should check the intended arc-transitivity convention,\allowbreak{} rate-one normalization,\allowbreak{} and prior probability literature before any promotion.

\subsection*{Scoped parent review receipt}

The parent independently reviewed the historical question and the checked construction. Its report is at \path{~/<private-repo>/12_cognitive_architecture/audits/reports/2026/09/06/last200-implementation/P0208_historical_question_candidate_review.md},\allowbreak{} SHA256 \texttt{b968660d\allowbreak{}bd875880\allowbreak{}70d48502\allowbreak{}f0d93a24\allowbreak{}84ae90df\allowbreak{}2f368b02\allowbreak{}875c4694\allowbreak{}ec5ad409}. The report found no gap in its checked scope. Structural validation is not mathematical certification. Present openness,\allowbreak{} novelty,\allowbreak{} human review,\allowbreak{} and unique canonical audit-node binding remain unverified. This receipt does not promote the result to an original solution.

\clearpage

\section*{Proof appendix:\allowbreak{} P0209 uniform clt counterexample candidate}

The current ledger disposition and the exact status limits in this packet govern its interpretation. Classical supporting facts and known-result rederivations are distinct from original open-problem solutions.


\subsection*{P0209:\allowbreak{} a four-state obstruction to the proposed uniform normal approximation}

Status:\allowbreak{} complete negative-answer candidate,\allowbreak{} author checked by spectral covariance and a direct rare-event argument. No independent review or novelty certificate. The construction responds to the literal historical question,\allowbreak{} whose current status is unverified.

Source:\allowbreak{} Aldous and Fill,\allowbreak{} \emph{Reversible Markov Chains and Random Walks on Graphs},\allowbreak{} 2002 unfinished monograph,\allowbreak{} 2014 recompilation,\allowbreak{} Open Problem 4.30,\allowbreak{} PDF p.137,\allowbreak{} \url{https://www.stat.berkeley.edu/~aldous/RWG/book.pdf} . The chapter is internally dated October 11,\allowbreak{} 1994. With a stationary finite reversible continuous-time chain and a function g satisfying E\(\pi\)g=\allowbreak{}0 and E\(\pi\)g²=\allowbreak{}1,\allowbreak{} the source asks for a universal bound

\[
 \sup_x\left|\mathbb P_\pi\left\{\frac{\int_0^t g(X_u)\,du}
 {\sigma\sqrt t}\le x\right\}-\Phi(x)\right|
 \le\psi(\|g\|_\infty,t/\tau_2),
 \qquad \psi(b,s)\longrightarrow0\quad(s\longrightarrow\infty).
\]

Here \(\sigma\)² is the asymptotic variance rate and \(\tau\)2 is the relaxation time. We keep the first argument exactly equal to two,\allowbreak{} so there is no hidden assumption of uniformity or continuity in b.

\subsection*{Four-state reversible family}

For an integer s at least three,\allowbreak{} set \(\epsilon\)=\allowbreak{}s\textasciicircum{}\{-2\} and R=\allowbreak{}s\textasciicircum{}4. Let Y be a stationary two-state chain on \{0,\allowbreak{}1\},\allowbreak{} with rates 0\(\to\)1 equal to \(\epsilon\) and 1\(\to\)0 equal to 1-\(\epsilon\). Its stationary success probability is \(\epsilon\) and its nonzero generator eigenvalue is -1. Independently let Z be a fair stationary sign chain,\allowbreak{} flipping at rate R/\allowbreak{}2,\allowbreak{} with nonzero generator eigenvalue -R. The product X=\allowbreak{}(Y,\allowbreak{}Z) has four states,\allowbreak{} is irreducible and reversible,\allowbreak{} and has generator eigenvalues 0,\allowbreak{}-1,\allowbreak{}-R,\allowbreak{}-(R+\allowbreak{}1). Thus \(\tau\)2=\allowbreak{}1.

Define

\[
 A=2\varepsilon+\sqrt{(1-\varepsilon)(1-4\varepsilon)},\qquad
 B=\frac{2-A}{1-\varepsilon},\qquad
 g(Y,Z)=AZ+B(Y-\varepsilon).
\]

For 0<\(\epsilon\)\(\le\)1/\allowbreak{}9,\allowbreak{} A and B are positive. The positive root A solves

\[
 A^2-4\varepsilon A+5\varepsilon-1=0,
\]

which is equivalent to A²+\allowbreak{}B²\(\epsilon\)(1-\(\epsilon\))=\allowbreak{}1. Independence and centering therefore give E\(\pi\)g=\allowbreak{}0 and E\(\pi\)g²=\allowbreak{}1. Also A+\allowbreak{}B(1-\(\epsilon\))=\allowbreak{}2. Every possible absolute value of g is at most max\{A+\allowbreak{}B\(\epsilon\),\allowbreak{}A+\allowbreak{}B(1-\(\epsilon\))\}=\allowbreak{}2 because \(\epsilon\)<1/\allowbreak{}2,\allowbreak{} and the value 2 occurs when Y=\allowbreak{}1,\allowbreak{}Z=\allowbreak{}1. Hence ||g||\(\infty\)=\allowbreak{}2 exactly for every s. Finally A\(\to\)1 and B\(\to\)1.

\subsection*{Exact variance rate}

The stationary autocovariance is

\[
 \mathbb E[g(X_0)g(X_u)]
 =A^2e^{-Ru}+B^2\varepsilon(1-\varepsilon)e^{-u}.
\]

This follows directly from the two nonconstant coordinate eigenfunctions. Integrating twice over the covariance triangle gives

\[
 \sigma^2=2\int_0^\infty\mathbb E[g(X_0)g(X_u)]\,du
 =\frac{2A^2}{R}+2B^2\varepsilon(1-\varepsilon)>0.
\]

Each fixed member of the family has the ordinary asymptotic normal limit. The issue is uniformity over chains at a time depending on the member.

\subsection*{Collapse at time t=\allowbreak{}s}

Observe the chain until t=\allowbreak{}s. Write F=\allowbreak{}A\(\int\)\(_0\)ˢZ\_\allowbreak{}u du and H=\allowbreak{}B\(\int\)\(_0\)ˢ(Y\_\allowbreak{}u-\(\epsilon\))du. The probability that Y is ever one on this interval is at most

\[
 \mathbb P(Y_0=1)+\mathbb E[\#\{0\to1\hbox{ transitions in }[0,s]\}]
 \le\varepsilon+s\varepsilon=\frac1{s^2}+\frac1s.
\]

On the complementary event,\allowbreak{} H=\allowbreak{}-B\(\epsilon\)s,\allowbreak{} and

\[
 \left|\frac{H}{\sigma\sqrt s}\right|
 \le\sqrt{\frac{\varepsilon s}{2(1-\varepsilon)}}\longrightarrow0.
\]

The fast component is negligible in L²:\allowbreak{}

\[
 \mathbb E F^2
 =2A^2\int_0^s(s-u)e^{-Ru}\,du
 \le\frac{2A^2s}{R},
\qquad
 \mathbb E\left(\frac F{\sigma\sqrt s}\right)^2
 \le\frac{A^2}{RB^2\varepsilon(1-\varepsilon)}\longrightarrow0.
\]

Here R\(\epsilon\)=\allowbreak{}s². Combining the rare-event bound with Chebyshev yields

\[
 \frac{\int_0^s g(X_u)\,du}{\sigma\sqrt s}
 \longrightarrow0\quad\hbox{in probability}.
\]

For any fixed x>0 its distribution function thus tends to one. The limiting inferior of its Kolmogorov distance from a standard normal is at least 1-Φ(x). Letting x decrease to zero shows that this limiting inferior is at least one half. Yet ||g||\(\infty\)=\allowbreak{}2 and t/\allowbreak{}\(\tau\)2=\allowbreak{}s\(\to\)\(\infty\). No function \(\psi\) of the stated kind can therefore exist.

\subsection*{Second checks and rate normalization}

The apparent conflict with variance convergence is illusory. For a coordinate with covariance c e\textasciicircum{}\{-\(\lambda\)u\},\allowbreak{} the exact finite-time variance is 2c[t/\allowbreak{}\(\lambda\)-(1-e\textasciicircum{}\{-\(\lambda\)t\})/\allowbreak{}\(\lambda\)²]. Substitution shows Var(\(\int\)g)/\allowbreak{}(\(\sigma\)²s)\(\to\)1. Rare visits carry the variance while their probability tends to zero;\allowbreak{} the normalized squares are not uniformly integrable. Variance convergence alone cannot force a central limit theorem.

Page 26 explicitly permits nonnegative transition rates with no constraint on their sums. Even imposing total jump rate at most one would not rescue the proposed bound. Replace the generator Q by Q/\allowbreak{}R and observe for time Rs. Every outgoing rate is at most (R/\allowbreak{}2+\allowbreak{}1)/\allowbreak{}R\(\le\)1. The relaxation time becomes R,\allowbreak{} the variance rate becomes R\(\sigma\)²,\allowbreak{} and the time integral has the same law as R times the old integral. Both the standardized law and the ratio t/\allowbreak{}\(\tau\)2=\allowbreak{}s are unchanged.

The independent coordinate construction also makes clear the missing quantity:\allowbreak{} an asymptotic variance scale relative to the function and the slow rare-event mode. A gap and a bounded instantaneous variance do not prevent the asymptotic variance from concentrating in rare excursions. This observation is explanatory,\allowbreak{} not a claim that any particular modified bound has been proved here.

Focused primary-source searches on September 6,\allowbreak{} 2026 found the original statement and the author's chapter notes,\allowbreak{} which discuss fixed-chain Berry–Esseen constants. No later resolution of this precise uniform question was verified. This is not a novelty certificate. A future reviewer should recheck the exact fixed L\(\infty\) normalization,\allowbreak{} stationary initialization,\allowbreak{} spectral gap,\allowbreak{} covariance integral,\allowbreak{} and uniform-limit quantifiers.

\subsection*{Scoped parent review receipt}

The parent independently reviewed the historical question and the checked construction. Its report is at \path{~/<private-repo>/12_cognitive_architecture/audits/reports/2026/09/06/last200-implementation/P0209_historical_question_candidate_review.md},\allowbreak{} SHA256 \texttt{d0a58ea9\allowbreak{}1383b4ab\allowbreak{}cbff6b4d\allowbreak{}85c2515c\allowbreak{}e723250e\allowbreak{}2794cee4\allowbreak{}b76f2b76\allowbreak{}84d70f14}. The report found no gap in its checked scope. Structural validation is not mathematical certification. Present openness,\allowbreak{} novelty,\allowbreak{} human review,\allowbreak{} and unique canonical audit-node binding remain unverified. This receipt does not promote the result to an original solution.

\clearpage

\section*{Proof appendix:\allowbreak{} P0248 optimized cayley known rederivation}

The current ledger disposition and the exact status limits in this packet govern its interpretation. Classical supporting facts and known-result rederivations are distinct from original open-problem solutions.


\subsection*{P0248:\allowbreak{} optimized Cayley walk ? known-result rederivation}

Status:\allowbreak{} author-checked rederivation of a literature-resolved historical question. This is excluded from the current-open target and original-solution count. No independent review of this proof is claimed.

Aldous?Fill,\allowbreak{} Conjecture/\allowbreak{}Open Problem 7.18,\allowbreak{} PDF page 253,\allowbreak{} asks whether a finite Cayley graph admits a symmetric generator distribution with relaxation time at most a universal constant times its squared diameter. Igor Pak,\allowbreak{} \emph{Using stopping times to bound mixing times},\allowbreak{} SODA 1999,\allowbreak{} Theorem 3.1,\allowbreak{} explicitly states that it resolves this question. The author-hosted PDF is internally dated October 31,\allowbreak{} 2004:\allowbreak{} \url{https://www.math.ucla.edu/~pak/papers/soda4.pdf} . The date of the archived copy is distinguished from the 1999 publication date.

\subsection*{Construction and direct proof}

Let the finite group be $G$,\allowbreak{} with inverse-closed generating set $S$ excluding the identity,\allowbreak{} and diameter $D$. For every $x\in G$,\allowbreak{} choose a geodesic word $w_x$ of length $\ell_x$. Randomize the word for $x$ equally between $w_x$ and the reversed,\allowbreak{} inverted word chosen for $x^{-1}$. Both words represent $x$ and have the same length. Let $N_s(x)$ be the number of occurrences of $s$ in this randomized word. Define

\[a_s=\frac1{|G|}\sum_x \ell_x\,\mathbb E N_s(x),\qquad A=\sum_s a_s=\frac1{|G|}\sum_x\ell_x^2\le D^2.\]

Inversion of words shows $a_s=a_{s^{-1}}$. Each $a_s>0$,\allowbreak{} because the word representing $s$ has length one. Thus $\mu(s)=a_s/A$ is a symmetric probability measure whose support is exactly $S$.

For a real function $f$ and uniform $\pi$,\allowbreak{} the independent uniform-pair identity is

\[\operatorname{Var}_\pi f=\frac1{2|G|^2}\sum_{z,x\in G}(f(zx)-f(z))^2.\]

Telescope along the chosen word for $x$ and use Cauchy?Schwarz. For each occurrence of $s$,\allowbreak{} translating the summation variable $z$ by the preceding prefix leaves the uniform sum unchanged. Therefore

\[\operatorname{Var}_\pi f\le\frac12\sum_s a_s\,\mathbb E_{z\sim\pi}(f(zs)-f(z))^2=A\,\mathcal E_\mu(f,f).\]

The variational characterization of the gap gives $\tau_2\le A\le D^2$. This is the continuous-time/\allowbreak{}rate-one spectral convention $\tau_2=(1-\lambda_2)^{-1}$;\allowbreak{} it is not a claim about absolute discrete-time eigenvalues when the graph is bipartite. Adding one-half holding makes aperiodic discrete-time mixing and changes the relaxation bound by a factor two.

\subsection*{Second check and limits}

The Aldous?Fill comparison theorem 7.16,\allowbreak{} with comparison increments uniform on the whole group,\allowbreak{} gives $\tau_2\le\max_s a_s/\mu(s)=A$. It independently reproduces the direct Dirichlet-form argument. The inversion randomization is necessary to justify symmetry in a noncommutative group;\allowbreak{} a deterministically chosen word table need not give symmetric counts.

This derivation was written before the Pak source was found in the campaign's later-literature check. Discovery order does not imply novelty. The open-problem record is now excluded and the argument survives only as a provenance-labelled mathematical contribution.

\clearpage

\section*{Proof appendix:\allowbreak{} P0249 reversibilization trace known consequence}

The current ledger disposition and the exact status limits in this packet govern its interpretation. Classical supporting facts and known-result rederivations are distinct from original open-problem solutions.


\subsection*{P0249:\allowbreak{} reversibilization trace inequality ? known-result consequence}

Status:\allowbreak{} author-checked complete argument for a historical conjecture. A later primary theorem implies the result,\allowbreak{} so this entry is excluded from the current-open target and original-solution count. The direct matrix proof is not independently reviewed here.

Source question:\allowbreak{} Aldous?Fill,\allowbreak{} Conjecture 9.22,\allowbreak{} PDF pages 327?329,\allowbreak{} asks whether $\operatorname{tr}[Z^2(P^*-P)]\ge0$. Here $P$ is an irreducible finite-state transition matrix,\allowbreak{} $P^*$ its adjoint in $L^2(\pi)$,\allowbreak{} and $Z$ the centered fundamental matrix. The centered convention is $Z1=0$ and $Z=(I-P)^{-1}$ on $L^2_0(\pi)$. Adding the constant projection to $Z$ would not change the displayed product.

\subsection*{Direct matrix proof}

Restrict all operators to the finite-dimensional real Hilbert space $L^2_0(\pi)$. Put

\[A=I-P=H+K,\quad H=\frac{A+A^*}{2},\quad K=\frac{A-A^*}{2}.\]

The reversible Dirichlet form shows $H$ is positive definite:\allowbreak{} its only zero-energy functions on the full space are constants,\allowbreak{} because the underlying undirected chain is connected. Also $K^*=-K$ and $P^*-P=2K$.

Let $Q=H^{-1/2}KH^{-1/2}$,\allowbreak{} $B=-Q^2=Q^*Q\succeq0$,\allowbreak{} and define $F(t)=\operatorname{tr}(H+tK)^{-1}$ for real $t$. All these inverses exist since their symmetric part is positive definite. Using transpose invariance of trace and averaging the inverses with $t$ and $-t$ gives

\[F(t)=\operatorname{tr}\!\left[H^{-1}(I+t^2B)^{-1}\right].\]

Indeed $(I+tQ)^{-1}+(I-tQ)^{-1}=2(I-t^2Q^2)^{-1}$. Differentiate the last trace formula:\allowbreak{}

\[F'(t)=-2t\operatorname{tr}\!\left[H^{-1}B(I+t^2B)^{-2}\right]\le0\qquad(t\ge0).\]

The product inside the trace need not be symmetric. Its trace is nonnegative because $C=B(I+t^2B)^{-2}$ is positive semidefinite and $\operatorname{tr}(H^{-1}C)=\operatorname{tr}(C^{1/2}H^{-1}C^{1/2})\ge0$. This explicitly avoids an invalid commuting assumption.

Differentiating the original inverse instead gives $F'(1)=-\operatorname{tr}(A^{-1}KA^{-1})=-\operatorname{tr}(A^{-2}K)$. Consequently

\[\operatorname{tr}[Z^2(P^*-P)]=-2F'(1)=4\operatorname{tr}\!\left[H^{-1}B(I+B)^{-2}\right]\ge0.\]

In fact equality holds if and only if $B=0$,\allowbreak{} equivalently $K=0$,\allowbreak{} equivalently the original chain is reversible. This follows from positive definiteness of $H^{-1}$ and the fact that the positive eigenvalues of $B$ remain positive under $b\mapsto b/(1+b)^2$.

\subsection*{Independent route through established literature}

Huang and Mao,\allowbreak{} \emph{Variational formulas of asymptotic variance for general discrete-time Markov chains},\allowbreak{} arXiv:\allowbreak{}2012.13895 (December 27,\allowbreak{} 2020),\allowbreak{} Theorem 4.3 on PDF page 14,\allowbreak{} proves that the asymptotic variance for a reversible kernel plus an antisymmetric perturbation is even in the perturbation parameter and decreases with its absolute value. Primary source:\allowbreak{} \url{https://arxiv.org/abs/2012.13895} . This is a stronger observable-by-observable statement than the required trace monotonicity.

For finite irreducible chains,\allowbreak{} first add one-half holding:\allowbreak{} $R=(I+P)/2$. Every convex mixture of $R$ and $R^*$ is irreducible and aperiodic,\allowbreak{} so the theorem's spectral-radius assumption on the mean-zero space is satisfied. Apply it to $R_t=(R+R^*)/2+t(R-R^*)/2$,\allowbreak{} for $-1\le t\le1$.

For an orthonormal basis $(e_j)$ of $L^2_0(\pi)$,\allowbreak{} the Poisson-equation formula reads

\[\sum_j\widetilde\sigma^2(R_t,e_j)=2\operatorname{tr}(I-R_t)^{-1}-(n-1).\]

Theorem 4.3 therefore gives the same trace monotonicity. Differentiate from the left at $t=1$;\allowbreak{} analyticity identifies this with the full derivative and yields the conjectured inequality for $R$. Since $Z_R=2Z_P$ and $R^*-R=(P^*-P)/2$,\allowbreak{} its trace expression is twice that for $P$. The conclusion follows even if $P$ is periodic. The original series definition of $Z$ in the periodic case is interpreted through the centered inverse or Abel limit.

The source does not explicitly name Aldous?Fill Conjecture 9.22 in the checked passage. The implication just proved is the reason for exclusion;\allowbreak{} no claim of a new resolution is made. Earlier finite-state antisymmetric-perturbation work cited by Huang?Mao may contain the result still earlier.

\clearpage

\section*{Proof appendix:\allowbreak{} P0266 P0267 critical allen cahn counterexamples candidate}

The current ledger disposition and the exact status limits in this packet govern its interpretation. Classical supporting facts and known-result rederivations are distinct from original open-problem solutions.


\subsection*{P0266 and P0267:\allowbreak{} critical Allen-Cahn endpoint counterexamples}

Status:\allowbreak{} complete author-checked candidate counterexamples to two explicitly numbered questions in a 2025 primary survey. No independent proof review or originality certification is claimed. Current-literature searches are incomplete. This packet is research,\allowbreak{} not a promoted solution.

\subsection*{Source and exact target}

Agresti and Veraar,\allowbreak{} \emph{Nonlinear SPDEs and Maximal Regularity:\allowbreak{} An Extended Survey},\allowbreak{} NoDEA (2025),\allowbreak{} Problems 4 and 5,\allowbreak{} journal page 130 of 150,\allowbreak{} PDF page 131 (the repository adds a cover page). Primary published PDF:\allowbreak{} \url{https://repository.tudelft.nl/file/File_783e6459-5d3b-4957-bcd3-64da5b8d9025} . Archived SHA256:\allowbreak{} a66ffa2c\allowbreak{}1d88f2b4\allowbreak{}6d563a3e\allowbreak{}4d3963bb\allowbreak{}c0f86234\allowbreak{}92e7e183\allowbreak{}7b4ae089\allowbreak{}008ad091. The arXiv version \url{https://arxiv.org/abs/2501.18561} has the questions on PDF page 92. The published page was rendered and visually read in this campaign.

Problem 4 asks for an expected squared maximal-energy bound in the setting of Theorem 6.4. Its explicit model is the Ito equation on the one-dimensional torus

\[du=(\Delta u+u-u^3)\,dt+\gamma u^2\,dW,\]

with one scalar Brownian driver,\allowbreak{} and it singles out $\gamma^2=2$ as unresolved. Problem 5 asks whether convergence of initial conditions in $L^2(\Omega;H)$ implies convergence of solutions in $L^2(\Omega;L^2(0,T;V)\cap C([0,T];H))$ under the same general assumptions. These are distinct bounds and continuity questions,\allowbreak{} not parameter copies. Both fail for the same critical equation.

\subsection*{Scalar reduction and global construction}

Normalize the torus volume to one and take the deterministic constant initial condition $u_0(x)=a>0$. A spatially constant solution is $u(t,x)=U_t$,\allowbreak{} where

\[dU_t=(U_t-U_t^3)\,dt+\sqrt2 U_t^2\,dW_t,\qquad U_0=a.\]

Set $s(t)=(e^{2t}-1)/2$. Let $R_s$ be a four-dimensional Bessel process,\allowbreak{} starting at $r=1/(\sqrt2 a)$ and satisfying

\[dR_s=dB_s+\frac{3}{2R_s}\,ds.\]

It is the radial part of four-dimensional Brownian motion started at a point of length $r$ and never hits zero. Define

\[U_t=\frac{e^t}{\sqrt2 R_{s(t)}}.\]

Ito's formula for $z(R)=1/(\sqrt2 R)$ gives drift $-z^3$ and diffusion $-\sqrt2 z^2$. Deterministic time change followed by multiplication by $e^t$ yields drift $U-U^3$ and diffusion $-\sqrt2 U^2$;\allowbreak{} replace the scalar Brownian driver by its negative. This proves the displayed SDE. Positivity and nonexplosion follow from the Bessel construction:\allowbreak{} a continuous positive Bessel path has a positive minimum on each compact time interval. Local Lipschitz coefficients give pathwise uniqueness;\allowbreak{} therefore this construction gives the law of the unique global strong solution for any prescribed scalar Brownian driver.

\subsection*{A finite-time Bessel minimum lemma}

For $0<b<r$,\allowbreak{} define $\tau_b=\inf\{s\ge0:R_s=b\}$. Fix $S>0$. The function $h(x)=|x|^{-2}$ is harmonic away from zero in dimension four. Its stopped process at $\tau_b$ is bounded by $b^{-2}$. Optional stopping at $\tau_b\wedge S$ gives the exact identity

\[\frac{\mathbb P_r(\tau_b\le S)}{b^2}=\frac1{r^2}-\mathbb E_r[R_S^{-2};\tau_b>S].\]

Since the process never hits zero,\allowbreak{} $\mathbf1_{\{\tau_b>S\}}\to1$ almost surely as $b\downarrow0$. The required integrability and strict loss of harmonic expectation can be computed directly. If $x\in\mathbb R^4$ has length $r$ and $B_S^{(4)}\sim N(0,SI_4)$,\allowbreak{} then Tonelli's theorem and a Gaussian integral give

\[\begin{aligned}
\mathbb E_r[R_S^{-2}]&=\int_0^\infty(1+2S\lambda)^{-2}\exp\!\left(-\frac{\lambda r^2}{1+2S\lambda}\right)d\lambda\\
&=\frac1{2S}\int_0^1\exp\!\left(-\frac{r^2(1-v)}{2S}\right)dv\\
&=\frac{1-e^{-r^2/(2S)}}{r^2}.
\end{aligned}\]

Dominated convergence in the stopped identity therefore proves

\[\lim_{b\downarrow0}\frac{\mathbb P_r(\tau_b\le S)}{b^2}=\frac{e^{-r^2/(2S)}}{r^2}>0.\]

No infinite-horizon hitting probability has been substituted for a finite-horizon one. The strict positive constant is explicitly derived for every $S>0$.

\subsection*{Failure of Problem 4}

Let $M_S=\sup_{0\le s\le S}(2R_s^2)^{-1}$. Substituting $b=(2x)^{-1/2}$ in the preceding lemma shows

\[\lim_{x\to\infty}x\,\mathbb P_r(M_S\ge x)=\frac{e^{-r^2/(2S)}}{2r^2}>0.\]

Thus $\mathbb E M_S=\infty$ by integration of the tail. For $S=s(T)$,\allowbreak{} the scalar construction gives

\[\sup_{0\le t\le T}U_t^2\ \ge\ M_{s(T)},\qquad \mathbb E\sup_{0\le t\le T}U_t^2=\infty\quad(T>0).\]

Consequently,\allowbreak{} $\mathbb E\sup_{t\le T}\|u(t)\|_{L^2(\mathbb T)}^2=\infty$,\allowbreak{} although $u_0$ is a smooth bounded deterministic function. The proposed bound fails for each positive finite time,\allowbreak{} not merely asymptotically.

At each fixed positive time,\allowbreak{} by contrast,\allowbreak{}

\[\mathbb E U_t^2=e^{2t}a^2\left(1-\exp\!\left[-\frac1{4a^2s(t)}\right]\right)<\infty.\]

This is consistent with the existing theorem's pointwise second-moment bound. The rare finite-time maxima have a $1/x$ tail for squared amplitude;\allowbreak{} replacing a localized stochastic integral by a mean-zero integrable martingale would miss exactly this issue.

\subsection*{Verification of the abstract assumptions}

There is already a finite-dimensional counterexample:\allowbreak{} choose the Gelfand triple $V=H=V^*=\mathbb R$,\allowbreak{} $A=I$,\allowbreak{} $B=0$,\allowbreak{} $F(v)=2v-v^3$,\allowbreak{} and $G(v)=\sqrt2 v^2$. Then the evolution equation $du+Au\,dt=F(u)\,dt+G(u)\,dW$ is the scalar SDE above. The linear coercivity condition holds with $\theta=1,M=0$. The nonlinear condition is

\[v(Av-F(v))-\tfrac12|G(v)|^2=-v^2\ge |v|^2-2|v|^2-2,\]

so Theorem 6.4's coercivity holds with $\theta'=1,M'=2$. Polynomial local Lipschitz bounds satisfy Assumption 6.1 with $\rho=2$,\allowbreak{} $\beta=2/3$,\allowbreak{} since $(2\beta-1)(\rho+1)=1$. All interpolation norms on this one-dimensional triple are equivalent. No compactness or infinite-dimensionality is required in the source assumptions.

For the stated torus model,\allowbreak{} take $V=H^1(\mathbb T)$,\allowbreak{} $H=L^2(\mathbb T)$,\allowbreak{} $A=I-\Delta$,\allowbreak{} $B=0$,\allowbreak{} $F(v)=2v-v^3$,\allowbreak{} and the same one-dimensional noise. Then

\[\langle v,Av-F(v)\rangle-\tfrac12\|G(v)\|_{L^2}^2=\|\nabla v\|_2^2-\|v\|_2^2=\|v\|_{H^1}^2-2\|v\|_2^2.\]

The polynomial Lipschitz condition follows from $(H^{-1},H^1)_{2/3,1}\hookrightarrow H^{1/3}\hookrightarrow L^6$ in one dimension,\allowbreak{} using Holder for the cubic drift and the resulting $L^4$ embedding for the quadratic diffusion. Thus the same critical indices $\rho=2$,\allowbreak{} $\beta=2/3$ suffice. The unique solution from Theorem 6.4 agrees with the spatially constant strong solution above.

\subsection*{Failure of Problem 5}

Use a common scalar Brownian driver and initial conditions $u_0^{(n)}(x)=1/n$,\allowbreak{} with limiting initial condition $u_0=0$. The solutions exist strongly and can be constructed simultaneously for this countable family. Their initial conditions converge deterministically in $L^2$. The limiting solution is identically zero,\allowbreak{} since both drift and diffusion vanish at zero. For every $n$ and every $T>0$,\allowbreak{} the preceding calculation applied with $a=1/n$ gives

\[\mathbb E\|u^{(n)}-0\|_{C([0,T];H)}^2=\infty.\]

Therefore convergence in the claimed $L^2(\Omega;L^2(0,T;V)\cap C([0,T];H))$ space cannot hold under the general hypotheses. This does not contradict convergence in probability or the stated $L^p$ convergence for $p<2$. In particular,\allowbreak{} tiny tail coefficients for small initial conditions do not make an infinite second maximal moment finite.

\subsection*{Adversarial checks and remaining review}

The argument's load-bearing checks are the Ito rather than Stratonovich convention;\allowbreak{} the scalar spatially constant driver allowed by the published model;\allowbreak{} the factor $\sqrt2$ and Bessel dimension four;\allowbreak{} the finite-horizon stopping identity;\allowbreak{} and the distinction between a fixed-time moment and the expected supremum. All are explicit above. The zero initial condition alone is not a counterexample,\allowbreak{} and $\gamma^2<2$ is outside this construction's critical claim.

A first later-literature search found the 2026 paper \emph{Fractal dimension of singular times for SPDEs:\allowbreak{} Energy bounds,\allowbreak{} criticality,\allowbreak{} and weak-strong uniqueness},\allowbreak{} arXiv:\allowbreak{}2602.21981. Its abstract concerns singular-time dimensions and does not resolve these two endpoint questions. Its full archived text has not yet been completely reviewed. Searches and an author's publication page are not a novelty certificate. The Bessel and Gaussian facts themselves are classical;\allowbreak{} the claimed candidate contribution is their application to the two precisely sourced questions. Scoped independent review is recorded below;\allowbreak{} human review and novelty remain pending.

\subsection*{Moment-order and continuity-scope clarification}

Only the fixed-time second moment is asserted finite. The fixed-time fourth moment is infinite:\allowbreak{} a nondegenerate four-dimensional Gaussian has a strictly positive continuous density at the origin,\allowbreak{} so the expectation of its inverse fourth radial power diverges logarithmically. The continuity counterexample addresses the numbered general Problem 5. It does not assert a resolution of the narrower surrounding prose with additional mapping or growth hypotheses. The later scoped review receipts below govern the review status;\allowbreak{} they do not certify novelty or the narrower prose variants.

\subsection*{Parent-reported scoped review receipt}

The parent reports a scoped check recorded at \path{~/<private-repo>/12_cognitive_architecture/audits/reports/2026/09/06/last200-implementation/P0266_critical_variational_candidate_review.md},\allowbreak{} SHA256 \texttt{51bab8e9\allowbreak{}34ed3092\allowbreak{}ab984962\allowbreak{}7b7eccf6\allowbreak{}bdabaedb\allowbreak{}6aeae7d3\allowbreak{}0aeb6f6c\allowbreak{}bb53dada}. Its checked calculations and source assumptions survived review. This is not human certification,\allowbreak{} a novelty certificate,\allowbreak{} or proof of present openness. Canonical audit-node binding remains pending. For P0189 the parent contributed to the amended proof,\allowbreak{} so the amended argument still requires a different independent reviewer. For P0266/\allowbreak{}P0267 the checked proof hash is \texttt{3afdb8ad\allowbreak{}f2e2b58a\allowbreak{}7778f807\allowbreak{}a9028374\allowbreak{}1394c820\allowbreak{}3e9541e1\allowbreak{}60ca7cc6\allowbreak{}f55be870};\allowbreak{} later moment-order and scope clarifications do not change the displayed construction. Any earlier no-review statement in this evolving packet describes its initial draft,\allowbreak{} superseded only within this explicitly scoped receipt.

\subsection*{Parent-reported scoped review receipt}

The parent reports a scoped check recorded at \path{~/<private-repo>/12_cognitive_architecture/audits/reports/2026/09/06/last200-implementation/P0267_critical_variational_candidate_review.md},\allowbreak{} SHA256 \texttt{23cf03ae\allowbreak{}22293403\allowbreak{}2f483c29\allowbreak{}10676dcf\allowbreak{}e211d249\allowbreak{}0c2d62b6\allowbreak{}42d28522\allowbreak{}b1d3b682}. Its checked calculations and source assumptions survived review. This is not human certification,\allowbreak{} a novelty certificate,\allowbreak{} or proof of present openness. Canonical audit-node binding remains pending. For P0189 the parent contributed to the amended proof,\allowbreak{} so the amended argument still requires a different independent reviewer. For P0266/\allowbreak{}P0267 the checked proof hash is \texttt{3afdb8ad\allowbreak{}f2e2b58a\allowbreak{}7778f807\allowbreak{}a9028374\allowbreak{}1394c820\allowbreak{}3e9541e1\allowbreak{}60ca7cc6\allowbreak{}f55be870};\allowbreak{} later moment-order and scope clarifications do not change the displayed construction. Any earlier no-review statement in this evolving packet describes its initial draft,\allowbreak{} superseded only within this explicitly scoped receipt.

\clearpage

\section*{Proof appendix:\allowbreak{} P0276 unweighted abelian cutoff counterexample candidate}

The current ledger disposition and the exact status limits in this packet govern its interpretation. Classical supporting facts and known-result rederivations are distinct from original open-problem solutions.


\subsection*{P0276:\allowbreak{} unweighted Abelian Cayley graphs with the product condition and no cutoff}

Status:\allowbreak{} complete author-checked counterexample candidate to the literal unrestricted statement of Salez Problem 2.4. A scoped independent review is recorded below;\allowbreak{} novelty remains unverified. This construction has unbounded degree. It does not address a bounded-degree or random-generator version of the question.

\subsection*{Primary source and scope}

Justin Salez,\allowbreak{} \emph{Modern aspects of Markov chains:\allowbreak{} entropy,\allowbreak{} curvature and the cutoff phenomenon},\allowbreak{} arXiv:\allowbreak{}2508.21055,\allowbreak{} Problem 2.4,\allowbreak{} PDF page 27;\allowbreak{} the definitions are on pages 25-26. Primary PDF:\allowbreak{} \url{https://arxiv.org/pdf/2508.21055} . Archived SHA256:\allowbreak{} 9b58f5e6\allowbreak{}c2b13103\allowbreak{}f9ca02dd\allowbreak{}d53b7890\allowbreak{}e259b2d9\allowbreak{}3f632c25\allowbreak{}d3414dd8\allowbreak{}7002e267.

The source defines the unweighted Cayley walk by uniform choice from a generating set. It asks whether cutoff is equivalent to the product condition for Abelian Cayley graphs,\allowbreak{} without a degree bound in that question. The subsequent paragraph separately introduces weighted increment laws. Therefore a weighted reset walk alone would not answer it. The construction below is a connected simple undirected unweighted Cayley graph and controls its entire state-space law.

The reset mechanism is classical prior art (often associated with Pak counterexamples);\allowbreak{} novelty,\allowbreak{} if any,\allowbreak{} would concern this exact unweighted realization and its application to the literal source statement. No novelty claim is made here. Earlier and later formulations may impose additional restrictions that this construction does not satisfy.

\subsection*{The graph and its exact rates}

For integers $d\ge8$,\allowbreak{} set

\[m_d=1+\left\lceil\frac{2^d\log d}{4}\right\rceil,
\qquad G_d=(\mathbb Z/2\mathbb Z)^d\times(\mathbb Z/m_d\mathbb Z).\]

Write $H_d=(\mathbb Z/2\mathbb Z)^d$,\allowbreak{} $K_d=\mathbb Z/m_d\mathbb Z$,\allowbreak{} and let $e_i$ be the cube basis vectors. The generating set is the disjoint union

\[S_d=\{(e_i,k):1\le i\le d,\ k\in K_d\setminus\{0\}\}
\ \cup\ \{(x,0):x\in H_d\setminus\{0\}\}.\]

It contains no identity and is inverse closed. The second part generates the cube factor,\allowbreak{} and adding $(e_1,0)$ to $(e_1,k)$ generates $(0,k)$,\allowbreak{} so it generates the full group. The Cayley graph is consequently simple,\allowbreak{} connected,\allowbreak{} undirected and vertex transitive. Every edge has the same transition weight. Its degree is

\[D_d=d(m_d-1)+2^d-1.\]

Use the rate-one continuization of simple random walk,\allowbreak{} whose generator is $P-I$. Split its jumps into local-type and global-type events. Their rates are

\[\alpha_d=\frac{d(m_d-1)}{D_d},\qquad
\theta_d=\frac{2^d-1}{D_d},\qquad
\beta_d=\frac{2^d}{D_d}.\]

On a local event the cube flips a uniformly chosen coordinate and the second coordinate receives an independent uniform nonzero increment. On a global event the cube receives a uniform nonzero increment and the second coordinate is unchanged.

\subsection*{Exact projected semigroup}

Let $L_d$ be the rate-one cube generator,\allowbreak{} and $\Pi_d$ the projection onto the uniform law on $H_d$. The transition operator for a uniform nonzero cube increment is

\[Q_d=\frac{2^d\Pi_d-I}{2^d-1}.\]

Therefore the cube projection has generator

\[\alpha_dL_d+\theta_d(Q_d-I)
=\alpha_dL_d+\beta_d(\Pi_d-I).\]

Since $\Pi_d$ commutes with the cube semigroup,\allowbreak{} if $q_d(t)$ is its total-variation distance from uniform when started at zero and $c_d(t)$ is the corresponding distance for rate-one cube walk,\allowbreak{} then

\[q_d(t)=e^{-\beta_dt}\,c_d(\alpha_dt).\]

This is exact;\allowbreak{} no replacement of a nonzero increment by an approximate reset is used.

\subsection*{Control of the whole graph,\allowbreak{} not just its projection}

Let $v_d(t)$ denote the total-variation distance of the full walk from the uniform law on $G_d$. Projection gives $v_d(t)\ge q_d(t)$. Let $N_t$ count local-type events,\allowbreak{} so $N_t$ is Poisson of mean $\alpha_dt$.

Conditional on the event types,\allowbreak{} their times,\allowbreak{} and all cube increments,\allowbreak{} the second coordinate is the sum of $N_t$ independent uniform nonzero $K_d$ increments. When $N_t\ge1$,\allowbreak{} its distance from uniform is at most $1/m_d$:\allowbreak{} the first increment has exactly that distance and subsequent convolution contracts total variation. When $N_t=0$,\allowbreak{} use the bound one. Coupling this conditional law to an independent uniform coordinate yields

\[0\le v_d(t)-q_d(t)\le e^{-\alpha_dt}+\frac1{m_d}.\]

Conditioning on the cube path is legitimate because local increments in the two factors are chosen independently given the event type. Their common jump count does not invalidate the conditional bound. Uniformity of an auxiliary coordinate is not silently assumed.

\subsection*{A self-contained cube cutoff calculation}

Define the comparison time

\[T_d=\frac{d\log d}{4\alpha_d}.\]

The rate-one cube coordinates flip independently at rate $1/d$. Their signs have means $e^{-2t/d}$ from the all-zero start. Hence

\[\chi^2(\mathcal L(X_t),\operatorname{Unif}(H_d))
=(1+e^{-4t/d})^d-1.\]

At $t=sd\log d/4$ this tends to zero for $s>1$,\allowbreak{} proving $c_d(t)\to0$. For $0<s<1$,\allowbreak{} the sum of coordinate signs has mean $d^{1-s/2}$ and variance at most $d$,\allowbreak{} while its stationary mean is zero and stationary variance is $d$. Chebyshev applied at half the former mean under both laws gives

\[c_d(sd\log d/4)\ge1-8d^{s-1}\longrightarrow1.\]

Also

\[\beta_dT_d=\frac{2^d\log d}{4(m_d-1)}\longrightarrow1.\]

Combining the exact projected formula with the whole-space error bound proves the limiting profile away from $s=1$:\allowbreak{}

\[\lim_{d\to\infty}v_d(sT_d)=
\begin{cases}e^{-s},&0<s<1,\\0,&s>1.\end{cases}\]

The large drop near $s=1$ does not constitute cutoff:\allowbreak{} a nonvanishing fraction of total-variation decay has already occurred gradually over a time proportional to $T_d$.

\subsection*{Spectral gap and the product condition}

Characters of $G_d$ are indexed by a subset $A\subseteq\{1,\ldots,d\}$ and a character $\psi$ of $K_d$. Put $a=|A|$. The eigenvalues of the discrete jump operator are

\[\lambda_{A,\psi}=\frac{(d-2a)c_\psi+b_A}{D_d},\qquad
c_\psi=\begin{cases}m_d-1,&\psi=1,\\-1,&\psi\ne1,\end{cases}
\quad b_A=\begin{cases}2^d-1,&A=\varnothing,\\-1,&A\ne\varnothing.\end{cases}\]

The three nonconstant sectors have gap candidates

\[\frac{2a(m_d-1)+2^d}{D_d}\quad(a\ge1,\psi=1),\qquad
\frac{dm_d}{D_d}\quad(A=\varnothing,\psi\ne1),\qquad
\frac{D_d+d-2a+1}{D_d}\quad(A\ne\varnothing,\psi\ne1).\]

For all sufficiently large $d$ the minimum is in the first sector at $a=1$. Thus

\[\gamma_d=\frac{2(m_d-1)+2^d}{D_d}
=\frac{2\alpha_d}{d}+\beta_d\sim\frac2d.\]

This also follows by noting that every nonconstant auxiliary-factor character is contracted at a rate bounded away from zero,\allowbreak{} while the first cube character has the displayed slow rate. The full character computation checks that a hidden auxiliary bottleneck has not been overlooked.

Write $t_d(\varepsilon)=\inf\{t:v_d(t)\le\varepsilon\}$. Monotonicity and the limiting profile imply

\[\frac{t_d(1/4)}{T_d}\longrightarrow1,\qquad
\frac{t_d(3/4)}{T_d}\longrightarrow\log(4/3).\]

For the first assertion,\allowbreak{} $e^{-s}>e^{-1}>1/4$ for every $s<1$,\allowbreak{} whereas the limit is zero for $s>1$. For the second,\allowbreak{} $e^{-s}$ crosses $3/4$ at the interior point $s=\log(4/3)$,\allowbreak{} so bracketing on either side gives the limit. Therefore

\[\gamma_d t_d(1/4)\sim\frac12\log d\longrightarrow\infty,
\qquad \frac{t_d(3/4)}{t_d(1/4)}\longrightarrow\log(4/3)<1.\]

The product condition holds and total-variation cutoff fails. By transitivity these are also the worst-start mixing times.

\subsection*{Review limits}

The potentially surprising assertion is that uniform generators can encode the rare-reset mechanism. Its load-bearing checks are disjointness and inverse closure of the generator set;\allowbreak{} exact projected reset rate $\beta_d$,\allowbreak{} which is different from the actual global jump rate $\theta_d$;\allowbreak{} conditional mixing of the auxiliary factor;\allowbreak{} the full character spectrum;\allowbreak{} and the distinction between a discontinuity in the limiting profile and cutoff. Each is derived above.

The construction has very large degree and specially chosen generators. It says nothing against the bounded-degree transitive-expander problem or almost-sure cutoff for uniformly random generating sets. Searches found both the literal Salez statement and related nilpotent-group conjecture formulations,\allowbreak{} but no verified prior disposition of this exact construction. The simple mechanism makes prior-art checking essential. This packet remains a candidate,\allowbreak{} with no promoted original solution count.

\subsection*{Scoped independent review receipt}

Root independently reviewed the proof pinned at SHA256 80633cef\allowbreak{}95af2117\allowbreak{}2e4fc4f0\allowbreak{}f4c67c89\allowbreak{}845f5222\allowbreak{}e9806b74\allowbreak{}d8de721b\allowbreak{}4b64dbeb. Its report is at \path{~/<private-repo>/12_cognitive_architecture/audits/reports/2026/09/06/last200-implementation/P0276_unweighted_abelian_candidate_review.md},\allowbreak{} SHA256 ee332a2a\allowbreak{}a7f3ab31\allowbreak{}cb0ae91a\allowbreak{}0a745f84\allowbreak{}eb2b21e0\allowbreak{}c18f0288\allowbreak{}69602cfd\allowbreak{}f1b90e7b. The reviewer checked 8 graphs,\allowbreak{} 360 characters,\allowbreak{} 2720 projection entries,\allowbreak{} 16 conditional convolutions,\allowbreak{} 10 product-mass cases and 18 rejected bad controls,\allowbreak{} as well as the full total-variation coupling,\allowbreak{} spectrum and mixing-time brackets. No mathematical gap was found in that scope. The unbounded degree and chosen generators remain essential. Human review,\allowbreak{} unique live review node,\allowbreak{} present openness and novelty remain pending. This receipt is a scoped review,\allowbreak{} not original-solution promotion.

\clearpage

\section*{Proof appendix:\allowbreak{} P0288 pollett rank two transition counterexample candidate}

The current ledger disposition and the exact status limits in this packet govern its interpretation. Classical supporting facts and known-result rederivations are distinct from original open-problem solutions.


\subsection*{P0288:\allowbreak{} a rank-two polynomial transition asymptotic for a killed chain}

Status:\allowbreak{} full author-checked counterexample candidate to Pollett's Section 2 conjecture (ii). No independent review or novelty certificate is claimed. The chain has a countably infinite irreducible transient class and bounded rates. This does not contradict the rank-one theorem in the lambda-positive recurrent setting.

\subsection*{Primary source and precise target}

Phil Pollett,\allowbreak{} \emph{Quasi-stationary distributions for queueing and other models},\allowbreak{} Queueing Systems 100 (2022),\allowbreak{} 241-243;\allowbreak{} arXiv:\allowbreak{}2202.05428,\allowbreak{} Section 2,\allowbreak{} PDF page 2. Primary source:\allowbreak{} \url{https://arxiv.org/pdf/2202.05428} . The complete three-page preprint was read,\allowbreak{} including the discussion distinguishing lambda-null recurrence from lambda-transience.

The source assumes a conservative,\allowbreak{} stable and regular continuous-time chain on an absorbing state and an irreducible countable transient class,\allowbreak{} with absorption almost surely. Conjecture (ii) says that if the polynomially corrected transition probabilities have finite strictly positive limits,\allowbreak{} their limit is always of the form $A x_i m_j$,\allowbreak{} allowing invariant or subinvariant vectors and measures. We refute this factorization. The other numbered asymptotic claims are not separately asserted false.

\subsection*{Construction and source hypotheses}

Use transient labels $C=\mathbb Z$ and a cemetery label $\partial$. From every $i\in\mathbb Z$ jump to each nearest neighbor at rate one. Kill at rate $\kappa>0$ everywhere,\allowbreak{} and at an additional rate $v>0$ when $i=0$. Thus the generator on the transient class is

\[Q=\Delta-\kappa I-v|0\rangle\langle0|,
 \qquad (\Delta f)(i)=f(i+1)+f(i-1)-2f(i).\]

Give the cemetery state rate zero. The full rate matrix is conservative:\allowbreak{} its killing rate is an actual transition to the cemetery state. Exit rates are at most $2+\kappa+v$,\allowbreak{} so the chain is stable and nonexplosive,\allowbreak{} and its transition function is unique. Its transient class is irreducible because every finite nearest-neighbor path has positive probability of being followed before killing. Independent rate-$\kappa$ killing implies absorption almost surely from every transient state.

To use the source's literal labels,\allowbreak{} map $\partial$ to its absorbing label $0$,\allowbreak{} lattice $0$ to $1$,\allowbreak{} positive lattice $k$ to $2k$,\allowbreak{} and negative lattice $-k$ to $2k+1$. Thus $C$ is exactly countably infinite as required. The distinguished lattice site is not the absorbing state.

Put $K=-\Delta+v|0\rangle\langle0|$ on $\ell^2(\mathbb Z)$ and

\[H_t(i,j)=\langle\delta_i,e^{-tK}\delta_j\rangle,
 \qquad p_t(i,j)=e^{-\kappa t}H_t(i,j).\]

The operator $K$ is bounded,\allowbreak{} nonnegative and self-adjoint. Its semigroup is the nearest-neighbor walk killed at the additional rate $v$ at lattice zero,\allowbreak{} by the elementary killed-generator construction.

\subsection*{Exact resolvent}

For real $s>0$,\allowbreak{} define

\[D(s)=\sqrt{s(s+4)},\qquad
 r(s)=\frac{s+2-D(s)}2\in(0,1).\]

The free resolvent is

\[R_s^0(i,j)=\int_0^\infty e^{-st}e^{t\Delta}(i,j)\,dt
 =\frac{r(s)^{|i-j|}}{D(s)}.\]

Indeed $r+r^{-1}=s+2$ solves the homogeneous nearest-neighbor resolvent recurrence,\allowbreak{} and the jump of its first difference at $i=j$ is $D(s)$. Decay as $|i-j|\to\infty$ selects the displayed root. The rank-one resolvent identity then gives

\[R_s(i,j):=\langle\delta_i,(s+K)^{-1}\delta_j\rangle
 =\frac{r^{|i-j|}}D-\frac{v r^{|i|+|j|}}{D(D+v)}.                 \tag{1}\]

This formula can also be checked directly by substituting it in the recurrence with the extra diagonal term $v$ at zero.

\subsection*{The lower spectral edge and the time asymptotic}

Write $a=|i-j|$,\allowbreak{} $b=|i|+|j|$ and $q=\sqrt{s}$ on the principal branch. As $q\to0$,\allowbreak{}

\[D=2q+\tfrac14q^3+O(q^5),\qquad
 \log r=-q+\tfrac1{24}q^3+O(q^5).\]

Consequently

\[\frac{r^a}{D}
 =\frac1{2q}-\frac a2+\left(\frac{a^2}4-\frac1{16}\right)q+O(q^2),
 \qquad
 \frac v{D+v}=1-\frac{2q}v+\frac{4q^2}{v^2}+O(q^3).\]

Subtracting the two terms of (1) cancels the pole and yields

\[R_s(i,j)=\frac{b-a}{2}+\frac1v-L_{ij}\sqrt{s}+O(s),
 \qquad
 L_{ij}=\frac{b^2-a^2}{4}+\frac bv+\frac2{v^2}.                 \tag{2}\]

A bare Laplace-transform expansion is not,\allowbreak{} by itself,\allowbreak{} a pointwise Tauberian theorem for an off-diagonal transition probability. Here the explicit resolvent supplies the stronger spectral information that is needed.

For $0<E<4$,\allowbreak{} set $s=-E+i0$. Then

\[D=i\sqrt{E(4-E)},\qquad r=1-E/2-i\sqrt{E(4-E)}/2.\]

In a sufficiently small interval $0<E<\varepsilon$,\allowbreak{} formula (1) has smooth boundary values and no denominator zero. Stieltjes inversion for the finite spectral measures therefore gives an absolutely continuous density there. The coefficients of its convergent expansion in $q$ are real,\allowbreak{} so the $O(s)$ term of (2) is real on this interval and its first possible further imaginary contribution is $O(E^{3/2})$. Hence

\[\frac{d\mu_{ij}}{dE}(E)
 =-\frac1\pi\operatorname{Im} R_{-E+i0}(i,j)
 =\frac{L_{ij}}\pi E^{1/2}+O(E^{3/2}).                         \tag{3}\]

There is no atom at zero:\allowbreak{} $sR_s(i,j)\to0$,\allowbreak{} or equivalently the nonnegative quadratic form of $K$ has no nonzero square-summable null vector. Any spectral mass outside $[0,\varepsilon)$ contributes $O(e^{-\varepsilon t})$ to $H_t(i,j)$;\allowbreak{} the total variation of an off-diagonal spectral measure is finite by Cauchy-Schwarz for spectral projections. Possible high-energy atoms are consequently harmless.

Integrating (3) against $e^{-tE}$,\allowbreak{} substituting $u=tE$ and bounding the stated remainder gives

\[H_t(i,j)=\frac{L_{ij}}{2\sqrt\pi}t^{-3/2}+O(t^{-5/2}),\]

since $\Gamma(3/2)=\sqrt\pi/2$. Using
$b^2-a^2=2(|i||j|+ij)$,\allowbreak{} the leading kernel is

\[\boxed{\displaystyle
 \lim_{t\to\infty}t^{3/2}e^{\kappa t}p_t(i,j)
 = C_{ij}
 =\frac{(|i|+2/v)(|j|+2/v)+ij}{4\sqrt\pi}.}                  \tag{4}\]

Every $C_{ij}$ is finite and strictly positive,\allowbreak{} including opposite signs,\allowbreak{} because $(|i|+2/v)(|j|+2/v)>|ij|$. Thus the source decay parameter is exactly $\lambda=\kappa$.

\subsection*{Lambda-transience and failure of factorization}

Monotone convergence in the nonnegative integral defining the resolvent,\allowbreak{} followed by (2),\allowbreak{} gives

\[\int_0^\infty e^{\lambda t}p_t(i,j)\,dt
 =\lim_{s\downarrow0}R_s(i,j)
 =\frac{|i|+|j|-|i-j|}{2}+\frac1v<\infty.\]

The transient class is therefore lambda-transient in exactly the source definition. The power $3/2$ is the unique exponent yielding a finite strictly positive limit for each fixed pair,\allowbreak{} so it also satisfies the source's smallest-exponent wording.

Take $v=2$ and transient lattice labels $i,j\in\{0,1\}$. Equation (4) yields

\[(C_{ij})_{i,j\in\{0,1\}}
 =\frac1{4\sqrt\pi}\begin{pmatrix}1&2\\2&5\end{pmatrix},
 \qquad \det(C_{ij})=\frac1{16\pi}>0.\]

Every matrix of the form $A x_i m_j$ has all its two-by-two minors equal to zero. Allowing subinvariant vectors cannot alter this algebraic necessity. Thus conjecture (ii),\allowbreak{} under the stated general lambda-transient hypotheses,\allowbreak{} fails for this chain.

\subsection*{A second check from reflection and the odd sector}

Reflection about lattice zero commutes with $K$. On odd functions the killing potential vanishes. For positive integers $i,j$,\allowbreak{} cancellation of the defect terms in (1),\allowbreak{} or the killed-walk reflection principle,\allowbreak{} gives the exact identity

\[H_t(i,j)-H_t(i,-j)=k_t(i-j)-k_t(i+j),\]

where the free kernel has Fourier representation

\[k_t(n)=\frac1{2\pi}\int_{-\pi}^{\pi}
 e^{-2t(1-\cos\theta)}\cos(n\theta)\,d\theta.\]

The difference of cosines is $2\sin(i\theta)\sin(j\theta)=2ij\theta^2+O(\theta^4)$. Gaussian scaling around zero,\allowbreak{} with exponential suppression away from zero,\allowbreak{} therefore gives

\[\lim_{t\to\infty}t^{3/2}[H_t(i,j)-H_t(i,-j)]
 =\frac{ij}{2\sqrt\pi}.\]

This matches (4). It also supplies a separate structural obstruction:\allowbreak{} a strictly positive symmetric rank-one kernel invariant under simultaneous reflection has an even positive factor,\allowbreak{} so it would satisfy $C_{ij}=C_{i,-j}$. The nonzero odd-sector limit rules that out independently of the displayed minor calculation.

\subsection*{What this does and does not settle}

The source's common-exponent conjecture has an elementary conditional observation:\allowbreak{} if finite strictly positive pure-power limits already exist for every pair,\allowbreak{} irreducibility and the Chapman-Kolmogorov lower bound compare any two entries at times differing by a fixed constant. Comparison in both directions forces their exponents equal. This does not prove existence of all pairwise limits from a single pair and is not counted as another new solution.

This packet addresses the rank-one assertion (ii). It makes no claim to characterize all limiting conditional distributions or to refute the separate survival-exponent assertion (iii). Rank-one perturbations,\allowbreak{} killed nearest-neighbor kernels and two-end harmonic functions are classical tools. A later theorem,\allowbreak{} earlier example or prior formulation may already contain this counterexample. Initial searches of the title,\allowbreak{} author page and conjecture wording found no verified disposition;\allowbreak{} that is not a novelty or current-openness certificate. A scoped independent source and spectral review has since been supplied,\allowbreak{} as recorded below;\allowbreak{} it is not an originality certificate.

\subsection*{Attribution update:\allowbreak{} published two-end construction and reviewer-derived Poissonization}

The parent reviewer independently checked the pinned original packet (SHA256 \texttt{08d929d1\allowbreak{}70e30f91\allowbreak{}4ffbfcd3\allowbreak{}38b98b27\allowbreak{}0aad706a\allowbreak{}9bdee819\allowbreak{}fba1c0dd\allowbreak{}d2654d46}) and found no gap in the checked source hypotheses and spectral calculation. Its canonical report is \path{P0288_killed_chain_candidate_review.md},\allowbreak{} SHA256 \texttt{94a89d6e\allowbreak{}5bbadef6\allowbreak{}026a0ff1\allowbreak{}644f33bb\allowbreak{}ab4093f1\allowbreak{}736019bb\allowbreak{}e55f6439\allowbreak{}cd2bc150},\allowbreak{} under the canonical audits reports directory for September 6,\allowbreak{} 2026. This is a scoped independent review of the original packet;\allowbreak{} this attribution addition was written afterward. Live audit-node registration and human review remain pending.

More significantly,\allowbreak{} Foley and McDonald,\allowbreak{} Yaglom limits can depend on the starting state (\url{https://arxiv.org/abs/1709.07578}),\allowbreak{} equations (36),\allowbreak{} (45),\allowbreak{} (46),\allowbreak{} give a published discrete two-end construction from which the reviewer derived a continuous-time counterexample to the same general factorization assertion. Pollett cites their paper. The reduction below is explicitly attributed to the reviewer;\allowbreak{} it is not asserted to be a corollary explicitly printed by Foley and McDonald.

Let $J$ be their substochastic nearest-neighbor hub chain,\allowbreak{} with inward probability $a$,\allowbreak{} outward probability $b$,\allowbreak{} $a+b=1$,\allowbreak{} $0<b<1/2$,\allowbreak{} and absorption probability $a$ at the hub. Put $B=J^2$ on the even sites,\allowbreak{} $\mu=4ab$,\allowbreak{} and $r=\sqrt{a/b}$. Their ballot formulas yield,\allowbreak{} for fixed even $x,y$ and $y\ne0$,\allowbreak{}

\[B^n(x,y)\sim D_{xy}\mu^n n^{-3/2},\qquad
 D_{xy}=\frac{r^{|x|-|y|}}{2\sqrt\pi}
       \big[(1+|x|)(1+|y|)+xy\big].\]

The hub column is $D_{x0}=r^{|x|}(1+|x|)/\sqrt\pi$;\allowbreak{} all coefficients are positive. Rate-one Poissonization,\allowbreak{} with transient generator $B-I$ and the missing row rate sent to the cemetery,\allowbreak{} gives

\[p_t(x,y)=e^{-(1-\mu)t}\mathbb E[a_N],\qquad
 a_n=B^n(x,y)/\mu^n,\quad N\sim\operatorname{Poisson}(\mu t).\]

Since $a_n\sim D_{xy}n^{-3/2}$,\allowbreak{} there is a global bound $a_n\le C(1+n)^{-3/2}$. On $N\ge\mu t/2$,\allowbreak{} the variables $(\mu t)^{3/2}a_N$ are uniformly bounded and converge in probability to $D_{xy}$. On the complementary event a Poisson exponential lower-tail bound dominates the polynomial prefactor. Thus

\[t^{3/2}e^{(1-\mu)t}p_t(x,y)\longrightarrow\mu^{-3/2}D_{xy}.\]

For sites $2,4$,\allowbreak{} the polynomial bracket matrix is $\begin{pmatrix}13&23\\23&41\end{pmatrix}$,\allowbreak{} whose determinant is $4$. The diagonal powers of $r$ cancel in this principal minor,\allowbreak{} so the continuous-time limiting minor is $\mu^{-3}/\pi>0$. The chain has decay parameter $1-\mu$ and an integrable critical $t^{-3/2}$ tail,\allowbreak{} hence is lambda-transient. This is the same factorization obstruction,\allowbreak{} obtained from published discrete formulas plus the reviewer's proved Poissonization step.

This does not identify the two generators and does not establish when this exact continuous-time corollary was first published. It does rule out treating the present packet as a certified original discovery. The campaign's promoted-original count remains zero. The source-page checks and independent finite controls are preserved in the canonical review;\allowbreak{} they are not a substitute for a complete novelty search.

\clearpage

\section*{Proof appendix:\allowbreak{} P0289 ross modulation counterexample candidate}

The current ledger disposition and the exact status limits in this packet govern its interpretation. Classical supporting facts and known-result rederivations are distinct from original open-problem solutions.


\subsection*{P0289:\allowbreak{} faster Cox modulation can increase mean workload}

Status:\allowbreak{} author-checked counterexample candidate to the universal monotonicity clause of Leskela's Problem 1 (Ross's second conjecture). Its separate request for a necessary-and-sufficient condition remains unresolved. No novelty or independent-review certificate is claimed. The construction uses a finite-state stationary ergodic CTMC as its intensity environment and deterministic unit service times.

Primary source:\allowbreak{} Lasse Leskela,\allowbreak{} \emph{Ross's second conjecture and supermodular stochastic ordering},\allowbreak{} Queueing Systems 100 (2022),\allowbreak{} 213-215,\allowbreak{} Problem 1 on journal page 214;\allowbreak{} institutional PDF page 3. \url{https://acris.aalto.fi/ws/portalfiles/portal/89616571/Ross_s_second_conjecture_and_supermodular_stochastic_ordering.pdf} . The source specifies mean stationary workload,\allowbreak{} not a transient or arbitrarily time-weighted average. Its complete three-page article was read. The institution's PDF has a cover page.

\subsection*{A uniform second-order bound with an explicit remainder}

Let $F(t)$ be any stationary nonnegative intensity environment bounded by $M$,\allowbreak{} with mean $a$. Let the arrival process be Cox with intensity $\lambda_c(t)=\epsilon F(ct)$. Customers require exactly one unit of service from a rate-one FIFO server. Assume $\rho=\epsilon M<1$. Denote stationary workload by $V_c$ and its mean by $w(c)$.

Couple the arrivals by thinning a stationary homogeneous Poisson process of rate $\rho$. Its stationary M/\allowbreak{}D/\allowbreak{}1 workload $\widehat V$ dominates $V_c$ and has finite moments. For an arrival configuration $A$,\allowbreak{} define its parallel-service shot workload at time zero by

\[S(A)=\sum_{T\in A\cap(-1,0]}(1+T).\]

The single-server workload is

\[V(A)=\sup_{u\ge0}\{N_A(-u,0]-u\},\]

including zero in the supremum. Put $D(A)=V(A)-S(A)$. This difference is nonnegative and is nondecreasing when arrivals are added. To check the latter,\allowbreak{} add one arrival of age $v$. If $v\ge1$,\allowbreak{} its shot contribution is zero and workload cannot decrease. If $0<v<1$,\allowbreak{} any candidate $u>v$ in the supremum gains one. A candidate $u<v$ is dominated,\allowbreak{} after the addition,\allowbreak{} by candidates at $v+\eta$ as $\eta\downarrow0$,\allowbreak{} whose limiting value is at least its old value plus $1-v$. Thus $V$ increases by at least $1-v$,\allowbreak{} exactly the new shot contribution. Apply this argument first on a finite history and then pass to increasing histories.

Consequently

\[0\le D(A)\le D(\widehat A),\qquad
 E D(\widehat A)=\frac{\rho^2}{2(1-\rho)}.                 \tag{1}\]

For completeness,\allowbreak{} the identity follows from $E\widehat S=\rho/2$ by Campbell's formula and $E\widehat V=\rho/[2(1-\rho)]$. The latter can be derived directly from stationarity of $\widehat V^2$:\allowbreak{} between arrivals its derivative is $-2\widehat V$,\allowbreak{} while a unit-service arrival increases its square by $2\widehat V+1$. The stable Poisson queue has finite exponential moments,\allowbreak{} so this expectation identity is justified;\allowbreak{} alternatively one can stop and use its regenerative busy periods.

Apply the same square-workload balance to the Cox queue. The intensity is bounded,\allowbreak{} and domination supplies the required moments. It gives

\[w(c)=\frac{\epsilon a}{2}+E[\lambda_c(0)V_c(0)].\]

Conditioning on the environment and using Campbell's formula for the shot term yields the exact bounded expansion

\[w(c)=\frac{\epsilon a}{2}
 +\epsilon^2\int_0^1(1-u)E[F(0)F(-cu)]\,du+R_c,
 \qquad
 0\le R_c\le\frac{\epsilon^3M^3}{2(1-\epsilon M)}.          \tag{2}\]

The remainder bound uses (1) and $\lambda_c(0)\le\epsilon M$. It is uniform in $c$. In particular this is not a numerical small-traffic extrapolation or an uncontrolled formal power series.

\subsection*{A finite-state Markov intensity}

Take $m=1024$,\allowbreak{} $\theta=2\pi/m$,\allowbreak{} and let $X_t$ be the stationary CTMC on the cyclic group $\mathbb Z/m\mathbb Z$ with only clockwise jumps,\allowbreak{} each at rate

\[r=\frac1{\sin\theta}.\]

Its stationary law is uniform;\allowbreak{} it is irreducible and mixing in continuous time. Set

\[F(t)=2+\cos(\theta X_t+\theta/4),\qquad
 \delta=\frac{1-\cos\theta}{\sin\theta}=\tan(\pi/m).\]

Thus $1\le F\le3$ and $a=2$. The displayed cosine values are all distinct:\allowbreak{} equality for two different cyclic labels would require $k+l+1/2$ to be an integer multiple of $m$. Therefore $F$ itself,\allowbreak{} not merely its hidden environment,\allowbreak{} is a finite-state stationary ergodic CTMC after relabeling.

The elementary Poisson jump-count generating function gives

\[E[F(0)F(t)]=4+\tfrac12 e^{-\delta|t|}\cos t.              \tag{3}\]

Indeed the first Fourier character has eigenvalue
$r(e^{i\theta}-1)=-\delta+i$,\allowbreak{} and the uniform average removes the second spatial harmonic. This does not require reversibility or stochastic monotonicity;\allowbreak{} the directed cycle has neither of those extra properties in the intensity ordering.

Define

\[I_\delta(c)=\int_0^1(1-u)e^{-\delta cu}\cos(cu)\,du.\]

Equations (2)-(3) now read

\[w(c)=\epsilon+2\epsilon^2+\tfrac12\epsilon^2 I_\delta(c)+R_c,
 \qquad 0\le R_c\le\frac{27\epsilon^3}{2(1-3\epsilon)}.    \tag{4}\]

\subsection*{Explicit reversal of the proposed monotonicity}

Choose $c_1=2\pi$,\allowbreak{} $c_2=3\pi$ and $\epsilon=10^{-4}$. Without damping,\allowbreak{}

\[I_0(c)=\frac{1-\cos c}{c^2},\qquad
 I_0(2\pi)=0,\quad I_0(3\pi)=\frac2{9\pi^2}.\]

The inequality $1-e^{-x}\le x$ gives

\[|I_\delta(c)-I_0(c)|\le\frac{\delta c}{6}.\]

Hence

\[\tfrac12[I_\delta(3\pi)-I_\delta(2\pi)]
 \ge\frac1{9\pi^2}-\frac{5\pi\delta}{12}.\]

Since $\tan x<2x$ for $0<x<\pi/4$,\allowbreak{} here $\delta<\pi/512$. Using only $\pi^2<10$ gives the rational lower bound

\[\tfrac12[I_\delta(3\pi)-I_\delta(2\pi)]
 >\frac1{90}-\frac{50}{6144}=\frac{137}{46080}.\]

At the chosen $\epsilon$ the normalized remainder bound is

\[\frac{27\epsilon}{2(1-3\epsilon)}=\frac{27}{19994}<\frac1{700}.\]

Therefore (4) proves,\allowbreak{} with no simulation,\allowbreak{}

\[\boxed{\displaystyle
 w(3\pi)-w(2\pi)
 >\epsilon^2\left(\frac{137}{46080}-\frac1{700}\right)
 =\epsilon^2\frac{2491}{1612800}>0.}\]

The queue is stable since its intensity is at most $3\epsilon<1$. Thus faster modulation increases mean stationary workload for this fully specified finite-state Markov environment and iid deterministic service times.

\subsection*{Second derivation of the coefficient}

The second-order coefficient can also be recovered by adding two arrivals to an empty queue. Put their ages at $y+d$ and $y$,\allowbreak{} where $0\le d\le1$. The interaction increment,\allowbreak{} beyond the sum of their separate workloads,\allowbreak{} is $y$ for $0\le y\le1-d$,\allowbreak{} is $1-d$ for $1-d\le y\le1$,\allowbreak{} and is $2-d-y$ for $1\le y\le2-d$. Its integral over $y$ is exactly $1-d$. The ordered pair therefore contributes

\[\epsilon^2\int_0^1(1-d)E[F(0)F(-cd)]\,dd,\]

matching (2). This second calculation checks the kernel,\allowbreak{} absence of an extra factorial,\allowbreak{} and the unit-service time scale. The uniform bound in (2),\allowbreak{} rather than this formal coefficient check alone,\allowbreak{} makes the counterexample rigorous.

\subsection*{Limits and further work}

The universal assertion in Problem 1 is refuted if the argument survives independent review. The requested necessary-and-sufficient condition on the intensity process is not supplied. The construction does not refute the established doubly stochastically monotone case,\allowbreak{} and does not settle Problem 2 about a deterministic image of such a chain. It does not establish anything for three-state environments:\allowbreak{} the explicit choice here has1024 distinct intensity states.

Initial targeted searches found the original author note and older Ross-type literature,\allowbreak{} including Miyoshi-Rolski2004 and Heyman's older counterexample concerning nonstationary-arrival formulations. These must be compared in full before claiming originality. Mean stationary workload,\allowbreak{} arrival-Palm waiting time and a transient time average must not be conflated. For deterministic unit service the stationary workload balance also relates arrival-Palm waiting mean to $w(c)$,\allowbreak{} but the source's exact workload formulation is the target used here. No new paid service,\allowbreak{} numerical simulation or public publication was used.

\clearpage

\section*{Proof appendix:\allowbreak{} P0340 beta tree drawn length counterexample candidate}

The current ledger disposition and the exact status limits in this packet govern its interpretation. Classical supporting facts and known-result rederivations are distinct from original open-problem solutions.


\subsection*{P0340:\allowbreak{} The proposed drawn-length growth is incompatible with an exact tree bound}

Status:\allowbreak{} author-checked candidate counterexample to the literal proposed rate,\allowbreak{} not a certified original result. No independent reviewer has checked this packet. Later-literature search is bounded,\allowbreak{} not a novelty certificate. The argument below also derives a positive finite linear-order expectation limit. Its constant is an explicit convergent series of finite-tree mean heights,\allowbreak{} not a closed-form numerical value.

Primary source:\allowbreak{} Aldous and Janson,\allowbreak{} \emph{The Critical Beta-splitting Random Tree II:\allowbreak{} Overview and Open Problems},\allowbreak{} arXiv:\allowbreak{}2303.02529v3,\allowbreak{} manuscript dated April 17,\allowbreak{} 2025;\allowbreak{} PDF pages 4,\allowbreak{} 12,\allowbreak{} 28-29,\allowbreak{} 54-56. Source PDF SHA256:\allowbreak{} a93a520b\allowbreak{}c7e7f157\allowbreak{}9802bfb8\allowbreak{}7ac88887\allowbreak{}12aadd59\allowbreak{}06c0980c\allowbreak{}a15d2efd\allowbreak{}4fa7735d. URL:\allowbreak{} \url{https://arxiv.org/abs/2303.02529v3.}

Open Problem 14 proposes that the expected drawn length grows on the scale $n(\log n)^4$. This packet disproves that scale under the exact drawing convention and recursion in source equations (72)-(73). The paper's ordinary random-leaf depth and its expected drawn length both use variants of the symbol L;\allowbreak{} we distinguish them throughout.

\subsection*{Definitions and deterministic inequality}

For a finite full rooted binary tree $T$,\allowbreak{} let $n(T)$ be its number of leaves,\allowbreak{} $h(T)$ its graph height with leaf height zero,\allowbreak{} $S(T)$ the sum of graph depths of its leaves,\allowbreak{} and $\ell(T)$ the total vertical length in the source's minimal integer-height drawing. For children $A,B$,\allowbreak{} the source definition is

\[h(T)=1+\max(h(A),h(B)),\qquad
\ell(T)=\ell(A)+\ell(B)+2+|h(A)-h(B)|.\]

Also $S(T)=S(A)+S(B)+n(T)$ and all three functionals vanish on a single leaf. Since a full binary tree with $i$ leaves has height at most $i-1$,\allowbreak{} if the children have $i$ and $n-i$ leaves,\allowbreak{}

\[2+|h(A)-h(B)|\le 2+\max(i-1,n-i-1)\le n.\]

Induction therefore proves the exact,\allowbreak{} distribution-free inequality

\[2(n-1)\le \ell(T)\le S(T). \tag{1}\]

The lower bound follows from the same recursion. An independent telescoping derivation gives

\[\ell(T)=\sum_{v\text{ internal}}h(T_v)+h(T),\qquad
S(T)=\sum_{v\text{ internal}}n(T_v).\]

Indeed each child height is subtracted once and each internal parent height added twice when the drawn edge differences are summed;\allowbreak{} leaves contribute zero. Since $h(T_v)\le n(T_v)-1$ and there are $n-1$ internal vertices,\allowbreak{} the second derivation again gives $\ell(T)\le S(T)$ using $h(T)\le n-1$. For a comb the upper bound is attained;\allowbreak{} for a balanced four-leaf tree the two quantities are six and eight. Thus neither an inequality reversal nor confusion with ordinary total edge count is hidden in (1).

\subsection*{Consequence for the source random tree}

Let $T_n=\mathrm{DTCS}(n)$ and let $D_n$ denote the graph depth of a uniformly selected leaf,\allowbreak{} conditional on the realized tree. Then $\mathbb E S(T_n)=n\mathbb E D_n$ by conditional averaging. Source Theorem 3 on PDF page 12,\allowbreak{} an established Aldous-Pittel/\allowbreak{}Aldous-Janson theorem,\allowbreak{} states

\[\mathbb E D_n\sim\frac{3}{\pi^2}(\log n)^2.\]

Consequently

\[2(n-1)\le\mathbb E\ell(T_n)
\le\left(\frac{3}{\pi^2}+o(1)\right)n(\log n)^2,
\qquad
\frac{\mathbb E\ell(T_n)}{n(\log n)^4}\longrightarrow0. \tag{2}\]

This rules out any positive constant asymptotic or two-sided comparability with $n(\log n)^4$. It also gives $\ell(T_n)/[n(\log n)^4]\to0$ in probability by Markov's inequality. The published random-leaf theorem is credited,\allowbreak{} not claimed as a new solution.

\subsection*{Self-contained bound without the published asymptotic}

Even without that theorem,\allowbreak{} the required upper order follows directly from the split law. The clade size $M$ seen along a uniformly chosen leaf has transition law

\[\mathbb P(M_{t+1}=j\mid M_t=m)=\frac{1}{H_{m-1}(m-j)},\quad 1\le j<m,\]

where $H_k=\sum_{r=1}^k1/r$. This follows by multiplying the source split probability $m/[2H_{m-1}j(m-j)]$ by the size-bias factor $2j/m$. The next clade has size at most $m/2$ with probability at least $1/(2H_{m-1})$,\allowbreak{} since there are $\lfloor m/2\rfloor$ terms,\allowbreak{} each at least $1/[H_{m-1}(m-1)]$.

Use the nonnegative potential $V(m)=(\log m)^2$,\allowbreak{} with $V(1)=0$. On that event,\allowbreak{}

\[V(m)-V(M_{t+1})\ge2\log 2\log m-(\log2)^2\ge\log2\log m.\]

The decrease is nonnegative otherwise. Also $H_{m-1}\le1+\log(m-1)\le(1+1/\log2)\log m$ for $m\ge2$. Hence the expected potential drop before absorption is at least

\[c_0=\frac{(\log2)^2}{2(1+\log2)}>0.\]

The chain strictly decreases,\allowbreak{} so its absorption time is at most $n-1$. Summing conditional expected drops through that bounded time gives $c_0\mathbb E D_n\le(\log n)^2$. Together with (1),\allowbreak{} this independently proves the upper order needed to refute the proposed growth. It avoids deriving an expectation bound from a probability-only maximum-height bound.

\subsection*{The apparent source error and remaining question}

The source occupation identity on PDF page 29 is $\mathbb E N_n(m)=na(n,m)/m$. For any additive split toll $b(m)$,\allowbreak{} the exact expected total divided by n is therefore $\sum_m a(n,m)b(m)/m$. The first heuristic line of equation (75) on PDF page 55 omits this factor $1/m$. Further approximations of absolute height differences are also unproved,\allowbreak{} so inserting the factor alone does not establish a corrected sharp asymptotic. The following separate argument establishes a linear-order expectation limit,\allowbreak{} without relying on the heuristic replacement of an absolute difference by a difference of expectations.

Search on September 7,\allowbreak{} 2026 opened the current arXiv abstract and the author's publication/\allowbreak{}home pages,\allowbreak{} and queried drawn length,\allowbreak{} cladogram length,\allowbreak{} and critical beta-splitting corrections. The author page still identifies overview version 3 and mentions an intended 2026 update. No matching correction was found in that bounded search. This does not establish that the observation is unpublished or that every newer source has been checked.

\subsection*{Stronger conclusion:\allowbreak{} the expected drawn length is asymptotically linear}

Put $H_m=\mathbb E h(T_m)$,\allowbreak{} with $H_1=0$. We prove

\[\frac{\mathbb E\ell(T_n)}{n}\longrightarrow
c_{\ell}:=\sum_{m=2}^{\infty}\frac{a(m)H_m}{m}
=\frac6{\pi^2}\sum_{m=2}^{\infty}\frac{H_{m-1}^{\mathrm{harm}}H_m}{m(m-1)},
\qquad 2\le c_{\ell}<\infty. \tag{3}\]

Here $H_k^{\mathrm{harm}}=\sum_{j=1}^k1/j$ distinguishes harmonic numbers from the mean tree height. The proof has three steps.

First,\allowbreak{} $H_m=O((\log m)^2)$. For the tagged chain begun at m,\allowbreak{} at each unabsorbed step the conditional probability of reducing the current size by at least half is at least $p_m=1/(2H_{m-1}^{\mathrm{harm}})$. Use fresh uniform variables at successive transitions to couple successful halvings from below by independent Bernoulli variables of probability $p_m$. The absorption time is then bounded by the waiting time for $b_m=\lceil\log_2 m\rceil$ successes. This is a sum of $b_m$ independent geometric variables G supported on the positive integers,\allowbreak{} with success probability $p_m$.

For $0<p\le1$,\allowbreak{} convexity gives $e^{p/4}\le1+p/2$. Thus

\[\mathbb E e^{pG/4}=\frac{p e^{p/4}}{1-(1-p)e^{p/4}}\le3.\]

A union bound over all m leaves therefore yields,\allowbreak{} with no independence assumption among their depths,\allowbreak{}

\[\mathbb P(h(T_m)>t)\le\min\{1,m3^{b_m}e^{-p_m t/4}\}.\]

Integrating this bound proves

\[H_m\le\frac4{p_m}\bigl(\log m+b_m\log3+1\bigr)=O((\log m)^2). \tag{4}\]

Second,\allowbreak{} the tagged occupation measures obey a uniform dyadic-block bound. Once the decreasing chain enters the integer block $(b,2b]$,\allowbreak{} at each subsequent step in that block it exits to a size at most b with conditional probability at least $1/(2H_{2b-1}^{\mathrm{harm}})$. It cannot return after exiting. Its expected number of visits is consequently at most $2H_{2b-1}^{\mathrm{harm}}$,\allowbreak{} uniformly in its starting state n:\allowbreak{}

\[\sum_{j=b+1}^{2b}a(n,j)\le2H_{2b-1}^{\mathrm{harm}}. \tag{5}\]

This remains true if the chain skips the block. Combining (4) and (5),\allowbreak{} for dyadic b at least two,\allowbreak{}

\[\sum_{j=b+1}^{2b}a(n,j)\frac{H_j}{j}
\le C\frac{(1+\log b)^3}{b}.\]

The sum of these bounds over dyadic blocks above any cutoff tends to zero as the cutoff tends to infinity,\allowbreak{} uniformly in n. This supplies precisely the tail control missing from a bare pointwise limit. Since $a(n,j)\to a(j)$ for each fixed j,\allowbreak{} finite-sum convergence followed by this uniform tail bound proves

\[\sum_{j=2}^n a(n,j)\frac{H_j}{j}\longrightarrow\sum_{j=2}^{\infty}a(j)\frac{H_j}{j}.\]

Third,\allowbreak{} the exact telescoping identity already proved and the branching property give

\[\mathbb E\ell(T_n)=n\sum_{j=2}^n\frac{a(n,j)H_j}{j}+H_n. \tag{6}\]

To justify the expectation identity,\allowbreak{} reveal the ancestors of any clade until its size j is known;\allowbreak{} its internal random tree is an independent copy of DTCS(j). Sum the conditional mean height $H_j$ over all clades,\allowbreak{} using the exact source identity $\mathbb E N_n(j)=na(n,j)/j$. Equation (4) gives $H_n/n\to0$,\allowbreak{} proving (3). The deterministic lower bound in (1) gives $c_{\ell}\ge2$;\allowbreak{} the convergent positive series also makes its finiteness explicit.

This establishes an expectation asymptotic,\allowbreak{} not a law of large numbers for $\ell(T_n)/n$,\allowbreak{} a variance formula,\allowbreak{} or a CLT. Those stronger statements are not claimed. The pointwise occupation formula is established literature;\allowbreak{} the dyadic uniform-tail argument and its application are checked here but lack independent review. A later paper,\allowbreak{} Brandenberger-Chin-Mossel arXiv:\allowbreak{}2505.24821v2 (July 2025),\allowbreak{} states an appropriate weighted-convergence result for all polynomially decreasing test functions on PDF page 3. Together with (4),\allowbreak{} that provides a second literature-supported route to (3);\allowbreak{} the self-contained block argument above avoids assuming that general weighted theorem.

\clearpage

\section*{Proof appendix:\allowbreak{} P0341 beta tree sparse occupation counterexample candidate}

The current ledger disposition and the exact status limits in this packet govern its interpretation. Classical supporting facts and known-result rederivations are distinct from original open-problem solutions.


\subsection*{P0341:\allowbreak{} Sparse root-sensitive tests refute both parts of the occupation ansatz}

Status:\allowbreak{} author-checked partial counterexample to Open Problem 8,\allowbreak{} specifically its displayed Ansatz 20. Explicit useful bounds for the occupation probabilities remain unresolved here. No independent review or originality certificate.

Source:\allowbreak{} Aldous-Janson,\allowbreak{} arXiv:\allowbreak{}2303.02529v3,\allowbreak{} PDF pages 28-30 and 39-40,\allowbreak{} April 2025. \url{https://arxiv.org/abs/2303.02529v3.} The PDF is pinned by the source registry. Write $a(n,j)$ for the probability that the tagged harmonic descent chain starting at n visits j,\allowbreak{} and $a(j)=6H_{j-1}/[\pi^2(j-1)]$ for its established fixed-j limit. Thus $a(n,n)=1$ exactly and $a(n,j)\to a(j)$ for fixed j.

The source's ansatz quantifies over every nonnegative sequence with at most polynomial growth. Part (i) asserts convergence of weighted sums whenever the limiting weighted series is finite. Part (ii) asserts equivalence with the corresponding truncated limiting sum when the series is infinite.

\subsection*{Counterexample to the summable case}

Choose the bounded function $f(j)=\mathbf1_{\{j=2^r\text{ for some integer }r\ge1\}}$. Since $a(2^r)=O(r2^{-r})$,\allowbreak{} the series $S=\sum_{r\ge1}a(2^r)$ is finite. For n equal to $2^K$,\allowbreak{} nonnegativity and the certain root visit give,\allowbreak{} for every fixed M<K,\allowbreak{}

\[\sum_{j=2}^{2^K}a(2^K,j)f(j)
\ge1+\sum_{r=1}^{M}a(2^K,2^r).\]

Let K tend to infinity and then M tend to infinity. Pointwise occupation convergence and monotone convergence of the positive series yield

\[\liminf_{K\to\infty}\sum_{j=2}^{2^K}a(2^K,j)f(j)\ge1+S.\]

The proposed limit is S,\allowbreak{} so part (i) is false even for a bounded nonnegative test function. No rate of occupation convergence or numerical calculation is used.

\subsection*{Counterexample to the divergent case}

Choose $g(j)=j\mathbf1_{\{j=2^r\text{ for some integer }r\ge1\}}$. This satisfies $g(j)\le j$. The limiting weighted series diverges,\allowbreak{} and its truncation at $n=2^K$ is

\[\sum_{j=2}^{2^K}a(j)g(j)=\sum_{r=1}^K\frac{6\,2^r H_{2^r-1}}{\pi^2(2^r-1)}\asymp K^2.\]

The actual finite-chain weighted sum is at least its root contribution $a(2^K,2^K)2^K=2^K$. Hence the ratio of actual to proposed sums tends to infinity along powers of two,\allowbreak{} refuting part (ii).

\subsection*{Why this does not solve every occupation question}

The obstruction is a moving root contribution invisible to every fixed-j limit. Restricting the test functions to appropriately regular,\allowbreak{} non-sparse classes may lead to valid theorems;\allowbreak{} this argument decides none of those refined classes or useful two-parameter error bounds. It also does not contradict the published special cases already proved by separate methods. The requested general ansatz is false,\allowbreak{} while the broader numbered question is only partially answered.

Searches on September 7,\allowbreak{} 2026 for the exact ansatz,\allowbreak{} occupation asymptotics and newer critical beta-splitting papers were recorded as bounded checks. New 2025 and August 2026 papers are being read for precise scope;\allowbreak{} no priority claim is made.

\clearpage

\section*{Proof appendix:\allowbreak{} P0342 beta tree fringe transfer scope counterexample}

The current ledger disposition and the exact status limits in this packet govern its interpretation. Classical supporting facts and known-result rederivations are distinct from original open-problem solutions.


\subsection*{P0342:\allowbreak{} A global finite-tree perturbation obstructs an unrestricted fringe CLT transfer}

Status:\allowbreak{} author-checked scope counterexample and partial answer to Aldous-Janson Open Problem 12. The example has a zero-variance Gaussian limit on the infinite fringe process. If the intended meaning of CLT requires strictly positive limiting variance or restricts functionals to a fixed finite neighborhood,\allowbreak{} that narrower question is not settled by this example. No independent review or novelty claim.

Source:\allowbreak{} Aldous-Janson,\allowbreak{} arXiv:\allowbreak{}2303.02529v3,\allowbreak{} PDF pages 44-45,\allowbreak{} with split law on PDF page 4. \url{https://arxiv.org/abs/2303.02529v3.} The source asks whether a CLT for a given functional along the stationary infinite fringe process necessarily transfers to the corresponding sum over leaves of the finite tree. Its preceding text mentions suitable functionals but imposes no quantitative locality or uniform error condition in the numbered question.

\subsection*{Explicit bounded functional}

For a finite full binary tree with n leaves,\allowbreak{} its original root is the unique degree-two vertex when n is at least two. Thus the original root split is recoverable after choosing a leaf as the distinguished root. Let $R_n$ be the smaller of the two original root-clade sizes,\allowbreak{} let $r_n=\lfloor\sqrt n\rfloor$,\allowbreak{} and for n at least nine put

\[B_n=\mathbf1_{\{R_n\le r_n\}},\qquad p_n=\mathbb P(B_n=1).\]

Define a single deterministic functional on finite or infinite leaf-rooted binary trees by

\[\Phi(T)=n^{-1/4}(B_n(T)-p_n)\quad\hbox{if T is finite with n at least nine leaves},\]

and $\Phi(T)=0$ for infinite T or smaller finite trees. This does not use external randomness. It is the same function at every n;\allowbreak{} n is a measurable statistic of its input tree. It is bounded by one.

In the usual local topology on these bounded-degree rooted trees,\allowbreak{} this functional is continuous. Every finite tree is isolated after a large enough rooted ball is specified,\allowbreak{} including the degrees of its boundary leaves. If finite trees converge locally to an infinite tree,\allowbreak{} their leaf counts must tend to infinity:\allowbreak{} a full binary tree with n leaves has exactly $2n-1$ vertices. Consequently $|\Phi(T)|\le n^{-1/4}\to0$ along every such sequence. The global split indicator need not itself be continuous at infinity because its prefactor tends uniformly to zero.

\subsection*{The two sums behave differently}

The stationary fringe tree and all its leaf rerootings are infinite. Thus $\Phi(T_i)=0$ almost surely for all i. Every centered normalized fringe sum is exactly zero,\allowbreak{} a degenerate Gaussian limit.

For the finite tree,\allowbreak{} rerooting at a different leaf leaves n and its unique degree-two original root unchanged. Hence every leaf has the same functional value and

\[\frac1{\sqrt n}\sum_{i=1}^n\Phi(T_i^{(n)})
=n^{1/4}(B_n-p_n). \tag{1}\]

Each summand has expectation zero already,\allowbreak{} so both centering by the infinite-fringe mean and centering by the finite-tree expectation give exactly (1).

For n at least nine the small-left and small-right split events are disjoint. The exact split law yields

\[p_n=\frac1{H_{n-1}}\left(H_{r_n}+\sum_{i=1}^{r_n}\frac1{n-i}\right)\longrightarrow\frac12.\]

Indeed $H_{r_n}/H_{n-1}\to1/2$,\allowbreak{} and the second numerator sum is at most $r_n/(n-r_n)=O(n^{-1/2})$. Eventually both $p_n$ and $1-p_n$ exceed one third. Therefore the absolute value in (1) is at least $n^{1/4}/3$ with probability one for all sufficiently large n. The finite normalized sums are not tight. Nevertheless their unnormalized empirical averages have magnitude at most $n^{-1/4}$ and tend to zero deterministically,\allowbreak{} so this does not contradict first-order fringe convergence.

\subsection*{Positive-variance interpretation and remaining obstruction}

If a nondegenerate Gaussian limit is mandatory,\allowbreak{} the construction by itself is not a full counterexample. It still gives the following precise structural implication. Let a class of admissible functionals contain some $\Psi$ with a nondegenerate fringe CLT and be closed under addition of the bounded continuous perturbation $\Phi$ above. A universal finite-tree transfer theorem cannot hold for every member of that class. The infinite fringe sums for $\Psi$ and $\Psi+\Phi$ are identical. If both centered finite-tree normalized sums were tight,\allowbreak{} their difference would be tight;\allowbreak{} the difference is (1),\allowbreak{} which is not tight. Thus at least one fails,\allowbreak{} conditional on the existence and admissibility of such $\Psi$.

The useful unresolved question is a transfer theorem with quantitative locality,\allowbreak{} stabilization,\allowbreak{} mixing,\allowbreak{} or a finite/\allowbreak{}infinite approximation error that is negligible after summing at square-root scale. Mere bounded continuity and first-order convergence are insufficient. These scope limits are retained in the problem ledger;\allowbreak{} this packet is not counted as a full solution of all intended versions of Open Problem 12.

\clearpage

\section*{Proof appendix:\allowbreak{} P0355 discounted RTP degenerate scope counterexample}

The current ledger disposition and the exact status limits in this packet govern its interpretation. Classical supporting facts and known-result rederivations are distinct from original open-problem solutions.


\subsection*{P0355:\allowbreak{} A degenerate case omitted by the discounted-tree nonendogeny conjecture}

Status:\allowbreak{} author-checked literal-scope counterexample,\allowbreak{} retained as a partial candidate because strict discounting and nondegeneracy would exclude it. No current-open or originality certification. It is a sourced conjecture,\allowbreak{} not a newly invented problem;\allowbreak{} its substantive intended nondegenerate regime remains unsolved here.

Source:\allowbreak{} Aldous and Bandyopadhyay,\allowbreak{} \emph{A survey of max-type recursive distributional equations},\allowbreak{} Annals of Applied Probability 15 (2005),\allowbreak{} 1047-1110. Conjecture 34 and Corollary 33,\allowbreak{} journal page 1075,\allowbreak{} PDF page 29;\allowbreak{} definitions and preceding sufficient theorem on PDF pages 26-28. Primary author PDF:\allowbreak{} \url{https://www.stat.berkeley.edu/~aldous/Papers/me107.pdf} . The exact archived bytes are pinned in the source registry under \path{aldous_rde_2005}.

\subsection*{Exact scope}

The displayed equation is

\[X\overset d=\eta+\max_i\xi_iX_i,\qquad X\in\mathbb R_+.\]

Corollary 33 assumes that the supremum of discounted ray sums is finite,\allowbreak{} and that the homogeneous equation obtained by removing $\eta$ has a nonzero solution $Y$. Condition (39) adds independence of $\eta$ and the weights,\allowbreak{} and that zero belongs to the support of $\eta$. Its parenthetical reference says the hypotheses of Theorem 32 are one sufficient way to obtain the finite supremum;\allowbreak{} it does not state that they are mandatory. In particular,\allowbreak{} the strict inequalities $\xi_i<1$ in Theorem 32 are not repeated as assumptions of Corollary 33. The state space includes zero,\allowbreak{} as confirmed by the source's repeated use of $\delta_0$ and $\eta\ge0$.

Conjecture 34 asserts that the invariant recursive tree process with marginal $\mu_a$,\allowbreak{} obtained by iterating the equation from the law of $aY$,\allowbreak{} is nonendogenous whenever $a>0$.

\subsection*{Construction and two checks}

Use the deterministic full binary tree and,\allowbreak{} at every vertex,\allowbreak{} take

\[N=2,\qquad \eta=0,\qquad \xi_1=\xi_2=1,\qquad Y=1.\]

Every discounted ray sum is zero,\allowbreak{} so its supremum is the finite random variable $X=0$. The homogeneous equation holds exactly because $\max(1,1)=1$. The variable $\eta$ is independent of the deterministic weight vector and its support is $\{0\}$. For every $a\ge0$,\allowbreak{} the distributional operator fixes $\delta_a$:\allowbreak{}

\[T(\delta_a)=\delta_{\max(a,a)}=\delta_a.\]

Consequently $T^n(\operatorname{law}(aY))=\delta_a$ for every $n$,\allowbreak{} so the source's $\mu_a$ is exactly $\delta_a$,\allowbreak{} and these measures are distinct as $a$ varies.

On the entire tree set every vertex variable equal to $a$. This is an invariant recursive tree process:\allowbreak{} its recursion holds,\allowbreak{} variables at any finite boundary are independent constants with the required common marginal,\allowbreak{} and they are independent of all ancestral innovations. The root is a constant,\allowbreak{} hence measurable with respect to the innovation sigma-field. The process is endogenous.

An independent formulation uses conditional laws. Given all innovations,\allowbreak{} the conditional root law is $\delta_a$;\allowbreak{} two conditionally independent copies therefore agree with probability one. Equivalently the only coupling with both marginals $\delta_a$ is the diagonal point mass $\delta_{(a,a)}$. Both checks contradict the literal nonendogeny conclusion.

\subsection*{Why the restriction matters}

This example uses both zero additive marks and unit edge weights. Replacing the marks by the positive constant one makes every ray sum infinite,\allowbreak{} violating the finite-supremum hypothesis. Replacing both weights by a deterministic $c<1$ makes the constant positive $Y$ cease to solve the homogeneous equation,\allowbreak{} since $\max(cY,cY)=cY\ne Y$. These controls expose the two load-bearing boundary choices.

The example does not answer a revised conjecture with strict discounts,\allowbreak{} nonzero additive marks,\allowbreak{} or suitable nondegeneracy assumptions. Later work on bounded discounted tree sums often explicitly assumes positive additive marks,\allowbreak{} so this observation does not contradict those theorems. A bounded search for the precise conjecture did not identify a published correction;\allowbreak{} failure to find one is not evidence of novelty.

\end{document}